% Auto-generated from data/segments.json + layout.json (+ transforms) — do not edit by hand.
% volume: 25Khu08
% mode: printing
% segments: 2111
% transform_rules: 65
% toc_marks: 82

% Volume layout defaults (layout.json ``layout``).
\csromanlayoutapply{1}{1.5}{21.6pt}{8pt}{8pt}{65pt}{2.5em}%

\setcsromanpitaka{สุตฺตนฺตปิฏก}
\setcsromannikaya{ขุทฺทกนิกาย}
\csromanheader{2}{\textbf{ขุทฺทกนิกาย}}
\csromanmatikaanchor{mtk.0}
\csromantocmarknopage{chapter}{จูฬนิทฺเทสปาฬิ}{mtk.0}
\bookmark[dest={mtk.0},level=0]{จูฬนิทฺเทสปาฬิ}
\gambhira{\textbf{จูฬนิทฺเทส}\-\textbf{ปาฬิ}}
\namakkaram{นโม ตสฺส ภควโต อรหโต สมฺมา\-สมฺพุทฺธสฺส.}
\csromanmatikaanchor{mtk.1}
\csromantocmarknopage{section}{ปารายนวคฺค}{mtk.1}
\bookmark[dest={mtk.1},level=1]{ปารายนวคฺค}
\csromanheader{1}{\textbf{ปารายนวคฺค}}
\csromanmatikaanchor{mtk.2}
\csromantocmarkref{subsection}{1. วตฺถุคาถา}{mtk.2}
\bookmark[dest={mtk.2},level=2]{1. วตฺถุคาถา}
\proseitem{1}{\csromanfolio{1}โกสลานํ ปุรา รมฺมา, อคมา ทกฺขิณาปถํ.}

\prose{อากิญฺจญฺญํ ปตฺถยาโน, พฺราหฺมโณ มนฺตปารคู.}

\proseitem{2}{โส อสฺสกสฺส วิสเย, มฬกสฺส\csromansharedfootnote{cab20d234ca5bdcc}{อฬกสฺส (ขุ 1. 429 ปิฏฺเฐ), มุฬกสฺส (สฺยา), มูฬฺหกสฺส (ก)} สมาสเน\csromansharedfootnote{f206d2e1dcc3771f}{สมาสนฺเน (ก)}.}

\prose{วสิ โคธาวรีกูเล, อุญฺเฉน จ ผเลน จ.}

\proseitem{3}{ตสฺเสว\csromansharedfootnote{5f42d39a1ffb960c}{ตํเยว (ก) อฏฺฐกถา โอโลเกตพฺพา.} อุปนิสฺสาย, คาโม จ วิปุโล อหุ.}

\prose{ตโต ชาเตน อาเยน, มหายญฺญม\-กปฺปยิ.}

\proseitem{4}{มหายญฺญํ ยชิตฺวาน, ปุน ปาวิสิ อสฺสมํ.}

\prose{ตสฺมิํ ปฏิปวิฏฺฐมฺหิ, อญฺโญ อาคญฺฉิ พฺราหฺมโณ.}

\setlength{\csromangathapagewidth}{0pt}
\setlength{\csromangathawakwidth}{0pt}
\csromangathameasuregroup{อุคฺฆฏฺฏปาโท ตสิโต\textsuperscript{1}, \\ โส จ นํ อุปสงฺกมฺม, \\ ตเมนํ พาวรี ทิสฺวา, \\ สุขญฺจ กุสลํ ปุจฺฉิ, \\ ยํ โข มม เทยฺยธมฺมํ, \\ อนุชานาหิ เม พฺรหฺเม, \\ สเจ เม ยาจมานสฺส, \\ สตฺตเม ทิวเส ตุยฺหํ, \\ อภิสงฺขริตฺวา กุหโก, \\ ตสฺส ตํ วจนํ สุตฺวา, \\ อุสฺสุสฺสติ อนาหาโร, \\ อโถปิ เอวํ จิตฺตสฺส, \\ อุตฺรสฺตํ ทุกฺขิตํ ทิสฺวา, \\ พาวริํ อุปสงฺกมฺม, \\ น โส มุทฺธํ ปชานาติ, \\ มุทฺธนิ มุทฺธปาเต\textsuperscript{1} วา, \\ โภตี\textsuperscript{1} จรหิ ชานาติ, \\ มุทฺธํ มุทฺธาธิปาตญฺจ\textsuperscript{1}, \\ อหมฺเปตํ น ชานามิ, \\ มุทฺธนิ มุทฺธาธิปาเต จ, \\ อถ โก จรหิ\textsuperscript{1} ชานาติ, \\ มุทฺธํ มุทฺธาธิปาตญฺจ, \\ ปุรา กปิลวตฺถุมฺหา, \\ อปจฺโจ โอกฺกากราชสฺส, \\ โส หิ พฺราหฺมณ สมฺพุทฺโธ, \\ สพฺพาภิญฺญาพลปฺปตฺโต\textsuperscript{1}, \\ สพฺพกมฺมกฺขยํ ปตฺโต, \\ พุทฺโธ โส ภควา โลเก, \\ ตํ ตฺวํ คนฺตฺวาน ปุจฺฉสฺสุ, \\ สมฺพุทฺโธติ วโจ สุตฺวา, \\ โสก’สฺส ตนุโก อาสิ, \\ ตโต อามนฺตยี สิสฺเส, \\ “เอถ มาณวา อกฺขิสฺสํ, \\ ยสฺเสโส ทุลฺลโภ โลเก, \\ สฺวาชฺช โลกมฺหิ อุปฺปนฺโน, \\ ขิปฺปํ คนฺตฺวาน สาวตฺถิํ, \\ กถํ จรหิ ชาเนมุ, \\ อชานตํ โน ปพฺรูหิ, \\ อาคตานิ หิ มนฺเตสุ, \\ ทฺวตฺติํสานิ จ พฺยากฺขาตา, \\ ยสฺเสเต โหนฺติ คตฺเตสุ, \\ เทฺวเยว ตสฺส คติโย, \\ สเจ อคารํ อาวสติ, \\ อทณฺเฑน อสตฺเถน, \\ สเจ จ โส ปพฺพชติ, \\ วิวฏฺฏจฺฉโท\textsuperscript{1} สมฺพุทฺโธ, \\ ชาติํ โคตฺตญฺจ ลกฺขณํ, \\ มุทฺธํ มุทฺธาธิปาตญฺจ, \\ อนาวรณทสฺสาวี, \\ มนสา ปุจฺฉิเต ปเญฺห, \\ พาวริสฺส วโจ สุตฺวา, \\ อชิโต ติสฺสเมตฺเตยฺโย, \\ โธตโก อุปสีโว จ, \\ โตเทยฺยกปฺปา ทุภโย, \\ ภทฺราวุโธ อุทโย จ, \\ โมฆราชา จ เมธาวี, \\ ปจฺเจกคณิโน สพฺเพ, \\ ฌายี ฌานรตา ธีรา, \\ พาวริํ อภิวาเทตฺวา, \\ ชฏาชินธรา สพฺเพ, \\ มฬกสฺส ปติฏฺฐานํ, \\ อุชฺเชนิญฺจาปิ โคนทฺธํ, \\ โกสมฺพิญฺจาปิ สาเกตํ, \\ เสตพฺยํ กปิลวตฺถุํ, \\ ปาวญฺจ โภคนครํ, \\ ปาสาณกํ เจติยญฺจ, \\ ตสิโตวุทกํ สีตํ, \\ ฉายํ ฆมฺมาภิตตฺโตว, \\ ภควา ตมฺหิ สมเย, \\ ภิกฺขูนํ ธมฺมํ เทเสติ, \\ อชิโต อทฺทส พุทฺธํ, \\ จนฺทํ ยถา ปนฺนรเส, \\ อถสฺส คตฺเต ทิสฺวาน, \\ เอกมนฺตํ ฐิโต หฏฺโฐ, \\ อาทิสฺส ชมฺมนํ พฺรูหิ, \\ มนฺเตสุ ปารมิํ พฺรูหิ, \\ วีสํ วสฺสสตํ อายุ, \\ ตีณิสฺส ลกฺขณา คตฺเต, \\ ลกฺขเณ อิติหาเส จ, \\ ปญฺจสตานิ วาเจติ, \\ ลกฺขณานํ ปวิจยํ, \\ ตณฺหจฺฉิท\textsuperscript{1} ปกาเสหิ, \\ มุขํ ชิวฺหาย ฉาเทติ, \\ โกโสหิตํ วตฺถคุยฺหํ,}{\csromangathabat{อุคฺฆฏฺฏปาโท ตสิโต\textsuperscript{1},}{ปงฺกทนฺโต รชสฺสิโร.} \\ \csromangathabat{โส จ นํ อุปสงฺกมฺม,}{สตานิ ปญฺจ ยาจติ.} \\ \csromangathabat{ตเมนํ พาวรี ทิสฺวา,}{อาสเนน นิมนฺตยิ.} \\ \csromangathabat{สุขญฺจ กุสลํ ปุจฺฉิ,}{อิทํ วจนมพฺรวิ\textsuperscript{1}.} \\ \csromangathabat{ยํ โข มม เทยฺยธมฺมํ,}{สพฺพํ วิสชฺชิตํ มยา.} \\ \csromangathabat{อนุชานาหิ เม พฺรหฺเม,}{นตฺถิ ปญฺจสตานิ เม.} \\ \csromangathabat{สเจ เม ยาจมานสฺส,}{ภวํ นานุปทสฺสติ\textsuperscript{1}.} \\ \csromangathabat{สตฺตเม ทิวเส ตุยฺหํ,}{มุทฺธา ผลตุ สตฺตธา.} \\ \csromangathabat{อภิสงฺขริตฺวา กุหโก,}{เภรวํ โส อกิตฺตยิ.} \\ \csromangathabat{ตสฺส ตํ วจนํ สุตฺวา,}{พาวรี ทุกฺขิโต อหุ.} \\ \csromangathabat{อุสฺสุสฺสติ อนาหาโร,}{โสกสลฺลสมปฺปิโต.} \\ \csromangathabat{อโถปิ เอวํ จิตฺตสฺส,}{ฌาเน น รมตี มโน.} \\ \csromangathabat{อุตฺรสฺตํ ทุกฺขิตํ ทิสฺวา,}{เทวตา อตฺถกามินี.} \\ \csromangathabat{พาวริํ อุปสงฺกมฺม,}{อิทํ วจนมพฺรวิ.} \\ \csromangathabat{น โส มุทฺธํ ปชานาติ,}{กุหโก โส ธนตฺถิโก.} \\ \csromangathabat{มุทฺธนิ มุทฺธปาเต\textsuperscript{1} วา,}{ญาณํ ตสฺส น วิชฺชติ.} \\ \csromangathabat{โภตี\textsuperscript{1} จรหิ ชานาติ,}{ตํ เม อกฺขาหิ ปุจฺฉิตา.} \\ \csromangathabat{มุทฺธํ มุทฺธาธิปาตญฺจ\textsuperscript{1},}{ตํ สุโณม วโจ ตว.} \\ \csromangathabat{อหมฺเปตํ น ชานามิ,}{ญาณํ เมตฺถ น วิชฺชติ.} \\ \csromangathabat{มุทฺธนิ มุทฺธาธิปาเต จ,}{ชินานํ เหตฺถ\textsuperscript{1} ทสฺสนํ.} \\ \csromangathabat{อถ โก จรหิ\textsuperscript{1} ชานาติ,}{อสฺมิํ ปถวิมณฺฑเล\textsuperscript{1}.} \\ \csromangathabat{มุทฺธํ มุทฺธาธิปาตญฺจ,}{ตํ เม อกฺขาหิ เทวเต.} \\ \csromangathabat{ปุรา กปิลวตฺถุมฺหา,}{นิกฺขนฺโต โลกนายโก.} \\ \csromangathabat{อปจฺโจ โอกฺกากราชสฺส,}{สกฺยปุตฺโต ปภงฺกโร.} \\ \csromangathabat{โส หิ พฺราหฺมณ สมฺพุทฺโธ,}{สพฺพธมฺมาน ปารคู.} \\ \csromangathabat{สพฺพาภิญฺญาพลปฺปตฺโต\textsuperscript{1},}{สพฺพธมฺเมสุ จกฺขุมา.} \\ \csromangathabat{สพฺพกมฺมกฺขยํ ปตฺโต,}{วิมุตฺโต อุปธิกฺขเย.} \\ \csromangathabat{พุทฺโธ โส ภควา โลเก,}{ธมฺมํ เทเสติ จกฺขุมา.} \\ \csromangathabat{ตํ ตฺวํ คนฺตฺวาน ปุจฺฉสฺสุ,}{โส เต ตํ พฺยากริสฺสติ.} \\ \csromangathabat{สมฺพุทฺโธติ วโจ สุตฺวา,}{อุทคฺโค พาวรี อหุ.} \\ \csromangathabat{โสก’สฺส ตนุโก อาสิ,}{ปีติญฺจ วิปุลํ ลภิ.} \\ โส พาวรี อตฺตมโน อุทคฺโค, \\ ตํ เทวตํ ปุจฺฉติ เวทชาโต. \\ กตมมฺหิ คาเม นิคมมฺหิ วา ปน, \\ กตมมฺหิ วา ชนปเท โลกนาโถ. \\ ยตฺถ คนฺตฺวาน ปสฺเสมุ, สมฺพุทฺธํ ทฺวิปทุตฺตมํ. \\ สาวตฺถิยํ โกสลมนฺทิเร ชิโน, \\ ปหูตปญฺโญ วรภูริเมธโส. \\ โส สกฺยปุตฺโต วิธุโร อนาสโว, \\ มุทฺธาธิปาตสฺส วิทู นราสโภ. \\ \csromangathabat{ตโต อามนฺตยี สิสฺเส,}{พฺราหฺมเณ มนฺตปารคู\textsuperscript{1}.} \\ \csromangathabat{“เอถ มาณวา อกฺขิสฺสํ,}{สุณาถ วจนํ มม.} \\ \csromangathabat{ยสฺเสโส ทุลฺลโภ โลเก,}{ปาตุภาโว อภิณฺหโส.} \\ \csromangathabat{สฺวาชฺช โลกมฺหิ อุปฺปนฺโน,}{สมฺพุทฺโธ อิติ วิสฺสุโต.} \\ \csromangathabat{ขิปฺปํ คนฺตฺวาน สาวตฺถิํ,}{ปสฺสโวฺห ทฺวิปทุตฺตมํ”.} \\ \csromangathabat{กถํ จรหิ ชาเนมุ,}{ทิสฺวา “พุทฺโธ”ติ พฺราหฺมณ.} \\ \csromangathabat{อชานตํ โน ปพฺรูหิ,}{ยถา ชาเนมุ ตํ มยํ.} \\ \csromangathabat{อาคตานิ หิ มนฺเตสุ,}{มหาปุริสลกฺขณา.} \\ \csromangathabat{ทฺวตฺติํสานิ จ พฺยากฺขาตา,}{สมตฺตา อนุปุพฺพโส.} \\ \csromangathabat{ยสฺเสเต โหนฺติ คตฺเตสุ,}{มหาปุริสลกฺขณา.} \\ \csromangathabat{เทฺวเยว ตสฺส คติโย,}{ตติยา หิ น วิชฺชติ.} \\ \csromangathabat{สเจ อคารํ อาวสติ,}{วิเชยฺย ปถวิํ อิมํ.} \\ \csromangathabat{อทณฺเฑน อสตฺเถน,}{ธมฺเมน อนุสาสติ.} \\ \csromangathabat{สเจ จ โส ปพฺพชติ,}{อคารา อนคาริยํ.} \\ \csromangathabat{วิวฏฺฏจฺฉโท\textsuperscript{1} สมฺพุทฺโธ,}{อรหา ภวติ อนุตฺตโร.} \\ \csromangathabat{ชาติํ โคตฺตญฺจ ลกฺขณํ,}{มนฺเต สิสฺเส ปุนาปเร.} \\ \csromangathabat{มุทฺธํ มุทฺธาธิปาตญฺจ,}{มนสาเยว ปุจฺฉถ.} \\ \csromangathabat{อนาวรณทสฺสาวี,}{ยทิ พุทฺโธ ภวิสฺสติ.} \\ \csromangathabat{มนสา ปุจฺฉิเต ปเญฺห,}{วาจาย วิสชฺชิสฺสติ\textsuperscript{1}.} \\ \csromangathabat{พาวริสฺส วโจ สุตฺวา,}{สิสฺสา โสฬส พฺราหฺมณา.} \\ \csromangathabat{อชิโต ติสฺสเมตฺเตยฺโย,}{ปุณฺณโก อถ เมตฺตคู.} \\ \csromangathabat{โธตโก อุปสีโว จ,}{นนฺโท จ อถ เหมโก.} \\ \csromangathabat{โตเทยฺยกปฺปา ทุภโย,}{ชตุกณฺณี จ ปณฺฑิโต.} \\ \csromangathabat{ภทฺราวุโธ อุทโย จ,}{โปสาโล จาปิ พฺราหฺมโณ.} \\ \csromangathabat{โมฆราชา จ เมธาวี,}{ปิงฺคิโย จ มหาอิสิ.} \\ \csromangathabat{ปจฺเจกคณิโน สพฺเพ,}{สพฺพโลกสฺส วิสฺสุตา.} \\ \csromangathabat{ฌายี ฌานรตา ธีรา,}{ปุพฺพวาสนวาสิตา.} \\ \csromangathabat{พาวริํ อภิวาเทตฺวา,}{กตฺวา จ นํ ปทกฺขิณํ.} \\ \csromangathabat{ชฏาชินธรา สพฺเพ,}{ปกฺกามุํ อุตฺตรามุขา.} \\ \csromangathabat{มฬกสฺส ปติฏฺฐานํ,}{ปุรมาหิสฺสติํ\textsuperscript{1} ตทา\textsuperscript{1}.} \\ \csromangathabat{อุชฺเชนิญฺจาปิ โคนทฺธํ,}{เวทิสํ วนสวฺหยํ.} \\ \csromangathabat{โกสมฺพิญฺจาปิ สาเกตํ,}{สาวตฺถิญฺจ ปุรุตฺตมํ.} \\ \csromangathabat{เสตพฺยํ กปิลวตฺถุํ,}{กุสินารญฺจ มนฺทิรํ.} \\ \csromangathabat{ปาวญฺจ โภคนครํ,}{เวสาลิํ มาคธํ ปุรํ.} \\ \csromangathabat{ปาสาณกํ เจติยญฺจ,}{รมณียํ มโนรมํ.} \\ \csromangathabat{ตสิโตวุทกํ สีตํ,}{มหาลาภํว วาณิโช.} \\ \csromangathabat{ฉายํ ฆมฺมาภิตตฺโตว,}{ตุริตา ปพฺพตมารุหุํ.} \\ \csromangathabat{ภควา ตมฺหิ สมเย,}{ภิกฺขุสํฆปุรกฺขโต.} \\ \csromangathabat{ภิกฺขูนํ ธมฺมํ เทเสติ,}{สีโหว นทตี วเน.} \\ \csromangathabat{อชิโต อทฺทส พุทฺธํ,}{ปีตรํสิํว\textsuperscript{1} ภาณุมํ.} \\ \csromangathabat{จนฺทํ ยถา ปนฺนรเส,}{ปริปูรํ\textsuperscript{1} อุปาคตํ.} \\ \csromangathabat{อถสฺส คตฺเต ทิสฺวาน,}{ปริปูรญฺจ พฺยญฺชนํ.} \\ \csromangathabat{เอกมนฺตํ ฐิโต หฏฺโฐ,}{มโนปเญฺห อปุจฺฉถ.} \\ \csromangathabat{อาทิสฺส ชมฺมนํ พฺรูหิ,}{โคตฺตํ พฺรูหิ สลกฺขณํ.} \\ \csromangathabat{มนฺเตสุ ปารมิํ พฺรูหิ,}{กติ วาเจติ พฺราหฺมโณ.} \\ \csromangathabat{วีสํ วสฺสสตํ อายุ,}{โส จ โคตฺเตน พาวรี.} \\ \csromangathabat{ตีณิสฺส ลกฺขณา คตฺเต,}{ติณฺณํ เวทาน ปารคู.} \\ \csromangathabat{ลกฺขเณ อิติหาเส จ,}{สนิฆณฺฑุสเกฏุเภ.} \\ \csromangathabat{ปญฺจสตานิ วาเจติ,}{สธมฺเม ปารมิํ คโต.} \\ \csromangathabat{ลกฺขณานํ ปวิจยํ,}{พาวริสฺส นรุตฺตม.} \\ \csromangathabat{ตณฺหจฺฉิท\textsuperscript{1} ปกาเสหิ,}{มา โน กงฺขายิตํ อหุ.} \\ \csromangathabat{มุขํ ชิวฺหาย ฉาเทติ,}{อุณฺณ’สฺส ภมุกนฺตเร.} \\ \csromangathabat{โกโสหิตํ วตฺถคุยฺหํ,}{เอวํ ชานาหิ มาณว.}}
\csromangathameasurewak{โส พาวรี อตฺตมโน อุทคฺโค, \\ ตํ เทวตํ ปุจฺฉติ เวทชาโต. \\ กตมมฺหิ คาเม นิคมมฺหิ วา ปน, \\ กตมมฺหิ วา ชนปเท โลกนาโถ. \\ ยตฺถ คนฺตฺวาน ปสฺเสมุ, สมฺพุทฺธํ ทฺวิปทุตฺตมํ. \\ สาวตฺถิยํ โกสลมนฺทิเร ชิโน, \\ ปหูตปญฺโญ วรภูริเมธโส. \\ โส สกฺยปุตฺโต วิธุโร อนาสโว, \\ มุทฺธาธิปาตสฺส วิทู นราสโภ.}
\csromangathasetleft{อุคฺฆฏฺฏปาโท ตสิโต\textsuperscript{1}, \\ โส จ นํ อุปสงฺกมฺม, \\ ตเมนํ พาวรี ทิสฺวา, \\ สุขญฺจ กุสลํ ปุจฺฉิ, \\ ยํ โข มม เทยฺยธมฺมํ, \\ อนุชานาหิ เม พฺรหฺเม, \\ สเจ เม ยาจมานสฺส, \\ สตฺตเม ทิวเส ตุยฺหํ, \\ อภิสงฺขริตฺวา กุหโก, \\ ตสฺส ตํ วจนํ สุตฺวา, \\ อุสฺสุสฺสติ อนาหาโร, \\ อโถปิ เอวํ จิตฺตสฺส, \\ อุตฺรสฺตํ ทุกฺขิตํ ทิสฺวา, \\ พาวริํ อุปสงฺกมฺม, \\ น โส มุทฺธํ ปชานาติ, \\ มุทฺธนิ มุทฺธปาเต\textsuperscript{1} วา, \\ โภตี\textsuperscript{1} จรหิ ชานาติ, \\ มุทฺธํ มุทฺธาธิปาตญฺจ\textsuperscript{1}, \\ อหมฺเปตํ น ชานามิ, \\ มุทฺธนิ มุทฺธาธิปาเต จ, \\ อถ โก จรหิ\textsuperscript{1} ชานาติ, \\ มุทฺธํ มุทฺธาธิปาตญฺจ, \\ ปุรา กปิลวตฺถุมฺหา, \\ อปจฺโจ โอกฺกากราชสฺส, \\ โส หิ พฺราหฺมณ สมฺพุทฺโธ, \\ สพฺพาภิญฺญาพลปฺปตฺโต\textsuperscript{1}, \\ สพฺพกมฺมกฺขยํ ปตฺโต, \\ พุทฺโธ โส ภควา โลเก, \\ ตํ ตฺวํ คนฺตฺวาน ปุจฺฉสฺสุ, \\ สมฺพุทฺโธติ วโจ สุตฺวา, \\ โสก’สฺส ตนุโก อาสิ, \\ ตโต อามนฺตยี สิสฺเส, \\ “เอถ มาณวา อกฺขิสฺสํ, \\ ยสฺเสโส ทุลฺลโภ โลเก, \\ สฺวาชฺช โลกมฺหิ อุปฺปนฺโน, \\ ขิปฺปํ คนฺตฺวาน สาวตฺถิํ, \\ กถํ จรหิ ชาเนมุ, \\ อชานตํ โน ปพฺรูหิ, \\ อาคตานิ หิ มนฺเตสุ, \\ ทฺวตฺติํสานิ จ พฺยากฺขาตา, \\ ยสฺเสเต โหนฺติ คตฺเตสุ, \\ เทฺวเยว ตสฺส คติโย, \\ สเจ อคารํ อาวสติ, \\ อทณฺเฑน อสตฺเถน, \\ สเจ จ โส ปพฺพชติ, \\ วิวฏฺฏจฺฉโท\textsuperscript{1} สมฺพุทฺโธ, \\ ชาติํ โคตฺตญฺจ ลกฺขณํ, \\ มุทฺธํ มุทฺธาธิปาตญฺจ, \\ อนาวรณทสฺสาวี, \\ มนสา ปุจฺฉิเต ปเญฺห, \\ พาวริสฺส วโจ สุตฺวา, \\ อชิโต ติสฺสเมตฺเตยฺโย, \\ โธตโก อุปสีโว จ, \\ โตเทยฺยกปฺปา ทุภโย, \\ ภทฺราวุโธ อุทโย จ, \\ โมฆราชา จ เมธาวี, \\ ปจฺเจกคณิโน สพฺเพ, \\ ฌายี ฌานรตา ธีรา, \\ พาวริํ อภิวาเทตฺวา, \\ ชฏาชินธรา สพฺเพ, \\ มฬกสฺส ปติฏฺฐานํ, \\ อุชฺเชนิญฺจาปิ โคนทฺธํ, \\ โกสมฺพิญฺจาปิ สาเกตํ, \\ เสตพฺยํ กปิลวตฺถุํ, \\ ปาวญฺจ โภคนครํ, \\ ปาสาณกํ เจติยญฺจ, \\ ตสิโตวุทกํ สีตํ, \\ ฉายํ ฆมฺมาภิตตฺโตว, \\ ภควา ตมฺหิ สมเย, \\ ภิกฺขูนํ ธมฺมํ เทเสติ, \\ อชิโต อทฺทส พุทฺธํ, \\ จนฺทํ ยถา ปนฺนรเส, \\ อถสฺส คตฺเต ทิสฺวาน, \\ เอกมนฺตํ ฐิโต หฏฺโฐ, \\ อาทิสฺส ชมฺมนํ พฺรูหิ, \\ มนฺเตสุ ปารมิํ พฺรูหิ, \\ วีสํ วสฺสสตํ อายุ, \\ ตีณิสฺส ลกฺขณา คตฺเต, \\ ลกฺขเณ อิติหาเส จ, \\ ปญฺจสตานิ วาเจติ, \\ ลกฺขณานํ ปวิจยํ, \\ ตณฺหจฺฉิท\textsuperscript{1} ปกาเสหิ, \\ มุขํ ชิวฺหาย ฉาเทติ, \\ โกโสหิตํ วตฺถคุยฺหํ,}
\csromangathagroup{\csromangathastanza{\csromangathaitembat{5}{อุคฺฆฏฺฏปาโท ตสิโต\csromansharedfootnote{bf81f55ab7d0f647}{ตสฺสิโต (ก)},}{ปงฺกทนฺโต รชสฺสิโร.} \\ \csromangathacontbat{โส จ นํ อุปสงฺกมฺม,}{สตานิ ปญฺจ ยาจติ.}}\csromangathastanza{\csromangathaitembat{6}{ตเมนํ พาวรี ทิสฺวา,}{อาสเนน นิมนฺตยิ.} \\ \csromangathacontbat{สุขญฺจ กุสลํ ปุจฺฉิ,}{อิทํ วจนมพฺรวิ\csromansharedfootnote{8adfe6e289d8d21b}{วจนมพฺรุวิ (สี)}.}}\csromangathastanza{\csromangathaitembat{7}{ยํ โข มม เทยฺยธมฺมํ,}{สพฺพํ วิสชฺชิตํ มยา.} \\ \csromangathacontbat{อนุชานาหิ เม พฺรหฺเม,}{นตฺถิ ปญฺจสตานิ เม.}}\csromangathastanza{\csromanfolio{2}\csromangathaitembat{8}{สเจ เม ยาจมานสฺส,}{ภวํ นานุปทสฺสติ\csromansharedfootnote{4b8c298a8f39c02b}{ปเทสฺสติ (ก)}.} \\ \csromangathacontbat{สตฺตเม ทิวเส ตุยฺหํ,}{มุทฺธา ผลตุ สตฺตธา.}}\csromangathastanza{\csromangathaitembat{9}{อภิส\-งฺขริ\-ตฺวา กุหโก,}{เภรวํ โส อกิตฺตยิ.} \\ \csromangathacontbat{ตสฺส ตํ วจนํ สุตฺวา,}{พาวรี ทุกฺขิโต อหุ.}}\csromangathastanza{\csromangathaitembat{10}{อุสฺสุสฺสติ อนาหาโร,}{โสกสลฺล\-สมปฺปิโต.} \\ \csromangathacontbat{อโถปิ เอวํ จิตฺตสฺส,}{ฌาเน น รมตี มโน.}}\csromangathastanza{\csromangathaitembat{11}{อุตฺรสฺตํ ทุกฺขิตํ ทิสฺวา,}{เทวตา อตฺถกามินี.} \\ \csromangathacontbat{พาวริํ อุปสงฺกมฺม,}{อิทํ วจนมพฺรวิ.}}\csromangathastanza{\csromangathaitembat{12}{น โส มุทฺธํ ปชานาติ,}{กุหโก โส ธนตฺถิโก.} \\ \csromangathacontbat{มุทฺธนิ มุทฺธปาเต\csromansharedfootnote{7dc37cf328c4723d}{มุทฺธนิมฺมุทฺธปาเต (ก)} วา,}{ญาณํ ตสฺส น วิชฺชติ.}}\csromangathastanza{\csromangathaitembat{13}{โภตี\csromansharedfootnote{285d0d004e74f1cb}{โภติ (ก)} จรหิ ชานาติ,}{ตํ เม อกฺขาหิ ปุจฺฉิตา.} \\ \csromangathacontbat{มุทฺธํ มุทฺธาธิปาตญฺจ\csromansharedfootnote{ba6c9221fd6a6eae}{มุทฺธาติปาตญฺจ (ก)},}{ตํ สุโณม วโจ ตว.}}\csromangathastanza{\csromangathaitembat{14}{อหมฺเปตํ น ชานามิ,}{ญาณํ เมตฺถ น วิชฺชติ.} \\ \csromangathacontbat{มุทฺธนิ มุทฺธาธิปาเต จ,}{ชินานํ เหตฺถ\csromansharedfootnote{bf47e1a3a2ea43bf}{ชนานเญฺหตฺถ (ก)} ทสฺสนํ.}}\csromangathastanza{\csromangathaitembat{15}{อถ โก จรหิ\csromansharedfootnote{f056a74b547dcddc}{โย จรติ (ก)} ชานาติ,}{อสฺมิํ ปถวิมณฺฑเล\csromansharedfootnote{4a0a6ae24372b0c1}{ปุถวิมณฺฑเล (สี)}.} \\ \csromangathacontbat{มุทฺธํ มุทฺธาธิปาตญฺจ,}{ตํ เม อกฺขาหิ เทวเต.}}\csromangathastanza{\csromangathaitembat{16}{ปุรา กปิล\-วตฺถุมฺหา,}{นิกฺขนฺโต โลกนายโก.} \\ \csromangathacontbat{อปจฺโจ โอกฺกากราชสฺส,}{สกฺยปุตฺโต ปภงฺกโร.}}\csromangathastanza{\csromangathaitembat{17}{โส หิ พฺราหฺมณ สมฺพุทฺโธ,}{สพฺพธมฺมาน ปารคู.} \\ \csromangathacontbat{สพฺพาภิญฺญาพลปฺปตฺโต\csromansharedfootnote{ac29352b6146bea2}{ผลปฺปตฺโต (ก)},}{สพฺพธมฺเมสุ จกฺขุมา.} \\ \csromangathacontbat{สพฺพกมฺม\-กฺขยํ ปตฺโต,}{วิมุตฺโต อุปธิกฺขเย.}}\csromangathastanza{\csromangathaitembat{18}{พุทฺโธ โส ภควา โลเก,}{ธมฺมํ เทเสติ จกฺขุมา.} \\ \csromangathacontbat{ตํ ตฺวํ คนฺตฺวาน ปุจฺฉสฺสุ,}{โส เต ตํ พฺยากริสฺสติ.}}\csromangathastanza{\csromangathaitembat{19}{สมฺพุทฺโธติ วโจ สุตฺวา,}{อุทคฺโค พาวรี อหุ.} \\ \csromangathacontbat{โสก’สฺส ตนุโก อาสิ,}{ปีติญฺจ วิปุลํ ลภิ.}}\csromangathastanza{\csromangathawak{\csromanfolio{3}\csromangathaitemline{20}{โส พาวรี อตฺตมโน อุทคฺโค,} \\ \csromangathacontline{ตํ เทวตํ ปุจฺฉติ เวทชาโต.} \\ \csromangathacontline{กตมมฺหิ คาเม นิคมมฺหิ วา ปน,} \\ \csromangathacontline{กตมมฺหิ วา ชนปเท โลกนาโถ.} \\ \csromangathacontline{ยตฺถ คนฺตฺวาน ปสฺเสมุ, สมฺพุทฺธํ ทฺวิปทุตฺตมํ.}}}\csromangathastanza{\csromangathawak{\csromangathaitemline{21}{สาวตฺถิยํ โกสลมนฺทิเร ชิโน,} \\ \csromangathacontline{ปหูตปญฺโญ วรภูริ\-เมธโส.} \\ \csromangathacontline{โส สกฺยปุตฺโต วิธุโร อนาสโว,} \\ \csromangathacontline{มุทฺธา\-ธิปาตสฺส วิทู นราสโภ.}}}\csromangathastanza{\csromangathaitembat{22}{ตโต อามนฺตยี สิสฺเส,}{พฺราหฺมเณ มนฺตปารคู\csromansharedfootnote{912b89a5bcb931be}{ปารเค (สฺยา)}.} \\ \csromangathacontbat{“เอถ มาณวา อกฺขิสฺสํ,}{สุณาถ วจนํ มม.}}\csromangathastanza{\csromangathaitembat{23}{ยสฺเสโส ทุลฺลโภ โลเก,}{ปาตุภาโว อภิณฺหโส.} \\ \csromangathacontbat{สฺวาชฺช โลกมฺหิ อุปฺปนฺโน,}{สมฺพุทฺโธ อิติ วิสฺสุโต.} \\ \csromangathacontbat{ขิปฺปํ คนฺตฺวาน สาวตฺถิํ,}{ปสฺสโวฺห ทฺวิปทุตฺตมํ”.}}\csromangathastanza{\csromangathaitembat{24}{กถํ จรหิ ชาเนมุ,}{ทิสฺวา “พุทฺโธ”ติ พฺราหฺมณ.} \\ \csromangathacontbat{อชานตํ โน ปพฺรูหิ,}{ยถา ชาเนมุ ตํ มยํ.}}\csromangathastanza{\csromangathaitembat{25}{อาคตานิ หิ มนฺเตสุ,}{มหา\-ปุริส\-ลกฺขณา.} \\ \csromangathacontbat{ทฺวตฺติํสานิ จ พฺยากฺขาตา,}{สมตฺตา อนุปุพฺพโส.}}\csromangathastanza{\csromangathaitembat{26}{ยสฺเสเต โหนฺติ คตฺเตสุ,}{มหา\-ปุริส\-ลกฺขณา.} \\ \csromangathacontbat{เทฺวเยว ตสฺส คติโย,}{ตติยา หิ น วิชฺชติ.}}\csromangathastanza{\csromangathaitembat{27}{สเจ อคารํ อาวสติ,}{วิเชยฺย ปถวิํ อิมํ.} \\ \csromangathacontbat{อทณฺเฑน อสตฺเถน,}{ธมฺเมน อนุสาสติ.}}\csromangathastanza{\csromangathaitembat{28}{สเจ จ โส ปพฺพชติ,}{อคารา อนคาริยํ.} \\ \csromangathacontbat{วิวฏฺฏจฺฉโท\csromansharedfootnote{c4bb035648538c9d}{วิวตฺตจฺฉทฺโท (สี)} สมฺพุทฺโธ,}{อรหา ภวติ อนุตฺตโร.}}\csromangathastanza{\csromangathaitembat{29}{ชาติํ โคตฺตญฺจ ลกฺขณํ,}{มนฺเต สิสฺเส ปุนาปเร.} \\ \csromangathacontbat{มุทฺธํ มุทฺธาธิปาตญฺจ,}{มนสาเยว ปุจฺฉถ.}}\csromangathastanza{\csromangathaitembat{30}{อนาวรณทสฺสาวี,}{ยทิ พุทฺโธ ภวิสฺสติ.} \\ \csromangathacontbat{มนสา ปุจฺฉิเต ปเญฺห,}{วาจาย วิสชฺชิสฺสติ\csromansharedfootnote{f93c7f45d8b8a93a}{วิสฺสชิสฺสติ (ก)}.}}\csromangathastanza{\csromanfolio{4}\csromangathaitembat{31}{พาวริสฺส วโจ สุตฺวา,}{สิสฺสา โสฬส พฺราหฺมณา.} \\ \csromangathacontbat{อชิโต ติสฺสเมตฺเตยฺโย,}{ปุณฺณโก อถ เมตฺตคู.}}\csromangathastanza{\csromangathaitembat{32}{โธตโก อุปสีโว จ,}{นนฺโท จ อถ เหมโก.} \\ \csromangathacontbat{โตเทยฺยกปฺปา ทุภโย,}{ชตุกณฺณี จ ปณฺฑิโต.}}\csromangathastanza{\csromangathaitembat{33}{ภทฺราวุโธ อุทโย จ,}{โปสาโล จาปิ พฺราหฺมโณ.} \\ \csromangathacontbat{โมฆราชา จ เมธาวี,}{ปิงฺคิโย จ มหาอิสิ.}}\csromangathastanza{\csromangathaitembat{34}{ปจฺเจกคณิโน สพฺเพ,}{สพฺพโลกสฺส วิสฺสุตา.} \\ \csromangathacontbat{ฌายี ฌานรตา ธีรา,}{ปุพฺพ\-วาสน\-วาสิตา.}}\csromangathastanza{\csromangathaitembat{35}{พาวริํ อภิวาเทตฺวา,}{กตฺวา จ นํ ปทกฺขิณํ.} \\ \csromangathacontbat{ชฏาชินธรา สพฺเพ,}{ปกฺกามุํ อุตฺตรามุขา.}}\csromangathastanza{\csromangathaitembat{36}{มฬกสฺส ปติฏฺฐานํ,}{ปุรมาหิสฺสติํ\csromansharedfootnote{2864ade0864c28fd}{ปุรมาหิยติ (ก)} ตทา\csromansharedfootnote{d970937995bd57f6}{สทา (ก)}.} \\ \csromangathacontbat{อุชฺเชนิญฺจาปิ โคนทฺธํ,}{เวทิสํ วนสวฺหยํ.}}\csromangathastanza{\csromangathaitembat{37}{โกสมฺพิญฺจาปิ สาเกตํ,}{สาวตฺถิญฺจ ปุรุตฺตมํ.} \\ \csromangathacontbat{เสตพฺยํ กปิลวตฺถุํ,}{กุสินารญฺจ มนฺทิรํ.}}\csromangathastanza{\csromangathaitembat{38}{ปาวญฺจ โภคนครํ,}{เวสาลิํ มาคธํ ปุรํ.} \\ \csromangathacontbat{ปาสาณกํ เจติยญฺจ,}{รมณียํ มโนรมํ.}}\csromangathastanza{\csromangathaitembat{39}{ตสิโตวุทกํ สีตํ,}{มหาลาภํว วาณิโช.} \\ \csromangathacontbat{ฉายํ ฆมฺมา\-ภิตตฺโตว,}{ตุริตา ปพฺพตมารุหุํ.}}\csromangathastanza{\csromangathaitembat{40}{ภควา ตมฺหิ สมเย,}{ภิกฺขุสํฆ\-ปุรกฺขโต.} \\ \csromangathacontbat{ภิกฺขูนํ ธมฺมํ เทเสติ,}{สีโหว นทตี วเน.}}\csromangathastanza{\csromangathaitembat{41}{อชิโต อทฺทส พุทฺธํ,}{ปีตรํสิํว\csromansharedfootnote{59c481b8fb0e9828}{ชิตรํสิํ สีตรํสิํ (ก), วีตรํสิํ (สี, สฺยา)} ภาณุมํ.} \\ \csromangathacontbat{จนฺทํ ยถา ปนฺนรเส,}{ปริปูรํ\csromansharedfootnote{7cef65ec6fba6663}{ปาริปูริํ (สี, สฺยา)} อุปาคตํ.}}\csromangathastanza{\csromangathaitembat{42}{อถสฺส คตฺเต ทิสฺวาน,}{ปริปูรญฺจ พฺยญฺชนํ.} \\ \csromangathacontbat{เอกมนฺตํ ฐิโต หฏฺโฐ,}{มโนปเญฺห อปุจฺฉถ.}}\csromangathastanza{\csromangathaitembat{43}{อาทิสฺส ชมฺมนํ พฺรูหิ,}{โคตฺตํ พฺรูหิ สลกฺขณํ.} \\ \csromangathacontbat{มนฺเตสุ ปารมิํ พฺรูหิ,}{กติ วาเจติ พฺราหฺมโณ.}}\csromangathastanza{\csromanfolio{5}\csromangathaitembat{44}{วีสํ วสฺสสตํ อายุ,}{โส จ โคตฺเตน พาวรี.} \\ \csromangathacontbat{ตีณิสฺส ลกฺขณา คตฺเต,}{ติณฺณํ เวทาน ปารคู.}}\csromangathastanza{\csromangathaitembat{45}{ลกฺขเณ อิติหาเส จ,}{สนิฆณฺฑุ\-สเกฏุเภ.} \\ \csromangathacontbat{ปญฺจสตานิ วาเจติ,}{สธมฺเม ปารมิํ คโต.}}\csromangathastanza{\csromangathaitembat{46}{ลกฺขณานํ ปวิจยํ,}{พาวริสฺส นรุตฺตม.} \\ \csromangathacontbat{ตณฺหจฺฉิท\csromansharedfootnote{f31d542f9ba8876c}{กงฺขจฺฉิท (ก)} ปกาเสหิ,}{มา โน กงฺขายิตํ อหุ.}}\csromangathastanza{\csromangathaitembat{47}{มุขํ ชิวฺหาย ฉาเทติ,}{อุณฺณ’สฺส ภมุกนฺตเร.} \\ \csromangathacontbat{โกโสหิตํ วตฺถคุยฺหํ,}{เอวํ ชานาหิ มาณว.}}}

\csromanhangingbegin
\hangingheaditem{48}{ปุจฺฉญฺหิ กิญฺจิ อสุณนฺโต, สุตฺวา ปเญฺห วิยากเต.}
\hangingbody{วิจินฺเตติ ชโน สพฺโพ, เวทชาโต กตญฺชลี.’}
\csromanhangingend

\setlength{\csromangathapagewidth}{0pt}
\setlength{\csromangathawakwidth}{0pt}
\csromangathameasuregroup{โก นุ เทโว วา พฺรหฺมา วา, \\ มนสา ปุจฺฉิเต ปเญฺห, \\ มุทฺธํ มุทฺธาธิปาตญฺจ, \\ ตํ พฺยากโรหิ ภควา, \\ อวิชฺชา มุทฺธาติ ชานาหิ, \\ สทฺธาสติสมาธีหิ, \\ ตโต เวเทน มหตา, \\ เอกํสํ อชินํ กตฺวา, \\ พาวรี พฺราหฺมโณ โภโต, \\ อุทคฺคจิตฺโต สุมโน, \\ สุขิโต พาวรี โหตุ, \\ ตฺวญฺจาปิ สุขิโต โหหิ, \\ พาวริสฺส จ ตุยฺหํ วา, \\ กตาวกาสา ปุจฺฉโวฺห, \\ สมฺพุทฺเธน กโตกาโส, \\ อชิโต ปฐมํ ปญฺหํ,}{\csromangathabat{โก นุ เทโว วา พฺรหฺมา วา,}{อินฺโท วาปิ สุชมฺปติ.} \\ \csromangathabat{มนสา ปุจฺฉิเต ปเญฺห,}{กเมตํ ปฏิภาสติ.} \\ \csromangathabat{มุทฺธํ มุทฺธาธิปาตญฺจ,}{พาวรี ปริปุจฺฉติ.} \\ \csromangathabat{ตํ พฺยากโรหิ ภควา,}{กงฺขํ วินย โน อิเส.} \\ \csromangathabat{อวิชฺชา มุทฺธาติ ชานาหิ,}{วิชฺชา มุทฺธาธิปาตินี.} \\ \csromangathabat{สทฺธาสติสมาธีหิ,}{ฉนฺทวีริเยน สํยุตา.} \\ \csromangathabat{ตโต เวเทน มหตา,}{สนฺถมฺเภตฺวาน มาณโว.} \\ \csromangathabat{เอกํสํ อชินํ กตฺวา,}{ปาเทสุ สิรสา ปติ.} \\ \csromangathabat{พาวรี พฺราหฺมโณ โภโต,}{สห สิสฺเสหิ มาริส.} \\ \csromangathabat{อุทคฺคจิตฺโต สุมโน,}{ปาเท วนฺทติ จกฺขุม.} \\ \csromangathabat{สุขิโต พาวรี โหตุ,}{สห สิสฺเสหิ พฺราหฺมโณ.} \\ \csromangathabat{ตฺวญฺจาปิ สุขิโต โหหิ,}{จิรํ ชีวาหิ มาณว.} \\ \csromangathabat{พาวริสฺส จ ตุยฺหํ วา,}{สพฺเพสํ สพฺพสํสยํ.} \\ \csromangathabat{กตาวกาสา ปุจฺฉโวฺห,}{ยํ กิญฺจิ มนสิจฺฉถ.} \\ \csromangathabat{สมฺพุทฺเธน กโตกาโส,}{นิสีทิตฺวาน ปญฺชลี.} \\ \csromangathabat{อชิโต ปฐมํ ปญฺหํ,}{ตตฺถ ปุจฺฉิ ตถาคตํ.}}
\csromangathasetleft{โก นุ เทโว วา พฺรหฺมา วา, \\ มนสา ปุจฺฉิเต ปเญฺห, \\ มุทฺธํ มุทฺธาธิปาตญฺจ, \\ ตํ พฺยากโรหิ ภควา, \\ อวิชฺชา มุทฺธาติ ชานาหิ, \\ สทฺธาสติสมาธีหิ, \\ ตโต เวเทน มหตา, \\ เอกํสํ อชินํ กตฺวา, \\ พาวรี พฺราหฺมโณ โภโต, \\ อุทคฺคจิตฺโต สุมโน, \\ สุขิโต พาวรี โหตุ, \\ ตฺวญฺจาปิ สุขิโต โหหิ, \\ พาวริสฺส จ ตุยฺหํ วา, \\ กตาวกาสา ปุจฺฉโวฺห, \\ สมฺพุทฺเธน กโตกาโส, \\ อชิโต ปฐมํ ปญฺหํ,}
\csromangathagroupwithcloser{\csromangathastanza{\csromangathaitembat{49}{โก นุ เทโว วา พฺรหฺมา วา,}{อินฺโท วาปิ สุชมฺปติ.} \\ \csromangathacontbat{มนสา ปุจฺฉิเต ปเญฺห,}{กเมตํ ปฏิภาสติ.}}\csromangathastanza{\csromangathaitembat{50}{มุทฺธํ มุทฺธาธิปาตญฺจ,}{พาวรี ปริปุจฺฉติ.} \\ \csromangathacontbat{ตํ พฺยากโรหิ ภควา,}{กงฺขํ วินย โน อิเส.}}\csromangathastanza{\csromangathaitembat{51}{อวิชฺชา มุทฺธาติ ชานาหิ,}{วิชฺชา มุทฺธา\-ธิปาตินี.} \\ \csromangathacontbat{สทฺธา\-สติ\-สมาธีหิ,}{ฉนฺทวีริเยน สํยุตา.}}\csromangathastanza{\csromangathaitembat{52}{ตโต เวเทน มหตา,}{สนฺถมฺ\-เภตฺวาน มาณโว.} \\ \csromangathacontbat{เอกํสํ อชินํ กตฺวา,}{ปาเทสุ สิรสา ปติ.}}\csromangathastanza{\csromangathaitembat{53}{พาวรี พฺราหฺมโณ โภโต,}{สห สิสฺเสหิ มาริส.} \\ \csromangathacontbat{อุทคฺคจิตฺโต สุมโน,}{ปาเท วนฺทติ จกฺขุม.}}\csromangathastanza{\csromangathaitembat{54}{สุขิโต พาวรี โหตุ,}{สห สิสฺเสหิ พฺราหฺมโณ.} \\ \csromangathacontbat{ตฺวญฺจาปิ สุขิโต โหหิ,}{จิรํ ชีวาหิ มาณว.}}\csromangathastanza{\csromangathaitembat{55}{พาวริสฺส จ ตุยฺหํ วา,}{สพฺเพสํ สพฺพสํสยํ.} \\ \csromangathacontbat{กตาวกาสา ปุจฺฉโวฺห,}{ยํ กิญฺจิ มนสิจฺฉถ.}}\csromangathastanza{\csromangathaitembat{56}{สมฺพุทฺเธน กโตกาโส,}{นิสีทิตฺวาน ปญฺชลี.} \\ \csromangathacontbat{อชิโต ปฐมํ ปญฺหํ,}{ตตฺถ ปุจฺฉิ ตถาคตํ.}}}{วตฺถุคาถา นิฏฺฐิตา.}
\csromanmatikaanchor{mtk.3}
\csromantocmarkref{subsection}{1. อชิตมาณวปุจฺฉา}{mtk.3}
\bookmark[dest={mtk.3},level=2]{1. อชิตมาณวปุจฺฉา}
\csromanheader{2}{1. \textbf{อชิตมาณวปุจฺฉา}}
\csromanhangingbegin
\hangingheaditem{57}{\csromanfolio{6}เกนสฺสุ นิวุโต โลโก, (อิจฺจายสฺมา อชิโต,)}
\hangingbody{เกนสฺสุ นปฺปกาสติ.}
\hangingbody{กิสฺสาภิเลปนํ พฺรูสิ, กิํสุ ตสฺส มหพฺภยํ.}
\csromanhangingend

\csromanhangingbegin
\hangingheaditem{58}{อวิชฺชาย นิวุโต โลโก, (อชิตาติ ภควา,)}
\hangingbody{เววิจฺฉา ปมาทา นปฺปกาสติ,}
\hangingbody{ชปฺปาภิเลปนํ พฺรูมิ, ทุกฺข’มสฺส มหพฺภยํ.}
\csromanhangingend

\csromanhangingbegin
\hangingheaditem{59}{สวนฺติ สพฺพธิ โสตา, (อิจฺจายสฺมา อชิโต,)}
\hangingbody{โสตานํ กิํ นิวารณํ,}
\hangingbody{โสตานํ สํวรํ พฺรูหิ, เกน โสตา ปิธิยฺยเร.’}
\csromanhangingend

\csromanhangingbegin
\hangingheaditem{60}{ยานิ โสตานิ โลกสฺมิํ, (อชิตาติ ภควา,)}
\hangingbody{สติ เตสํ นิวารณํ.}
\hangingbody{โสตานํ สํวรํ พฺรูมิ, ปญฺญาเยเต ปิธิยฺยเร.}
\csromanhangingend

\csromanhangingbegin
\hangingheaditem{61}{ปญฺญา เจว สติ จาปิ\csromansharedfootnote{602a9c076bc5ca39}{สตี เจว (สี)}, (อิจฺจายสฺมา อชิโต,)}
\hangingbody{นามรูปญฺจ มาริส.}
\hangingbody{เอตํ เม ปุฏฺโฐ ปพฺรูหิ, กตฺเถตํ อุปรุชฺฌติ.}
\csromanhangingend

\setlength{\csromangathapagewidth}{0pt}
\setlength{\csromangathawakwidth}{0pt}
\csromangathameasuregroup{ยเมตํ ปญฺหํ อปุจฺฉิ, \\ ยตฺถ นามญฺจ รูปญฺจ, \\ วิญฺญาณสฺส นิโรเธน, \\ เย จ สงฺขาตธมฺมาเส, \\ เตสํ เม นิปโก อิริยํ, \\ กาเมสุ นาภิคิชฺเฌยฺย, \\ กุสโล สพฺพธมฺมานํ,}{\csromangathabat{ยเมตํ ปญฺหํ อปุจฺฉิ,}{อชิต ตํ วทามิ เต.} \\ \csromangathabat{ยตฺถ นามญฺจ รูปญฺจ,}{อเสสํ อุปรุชฺฌติ.} \\ \csromangathabat{วิญฺญาณสฺส นิโรเธน,}{เอตฺเถตํ อุปรุชฺฌติ.} \\ \csromangathabat{เย จ สงฺขาตธมฺมาเส,}{เย จ เสขา\textsuperscript{1} ปุถู อิธ.} \\ \csromangathabat{เตสํ เม นิปโก อิริยํ,}{ปุฏฺโฐ ปพฺรูหิ มาริส.} \\ \csromangathabat{กาเมสุ นาภิคิชฺเฌยฺย,}{มนสา’นาวิโล สิยา.} \\ \csromangathabat{กุสโล สพฺพธมฺมานํ,}{สโต ภิกฺขุ ปริพฺพเชติ.}}
\csromangathasetleft{ยเมตํ ปญฺหํ อปุจฺฉิ, \\ ยตฺถ นามญฺจ รูปญฺจ, \\ วิญฺญาณสฺส นิโรเธน, \\ เย จ สงฺขาตธมฺมาเส, \\ เตสํ เม นิปโก อิริยํ, \\ กาเมสุ นาภิคิชฺเฌยฺย, \\ กุสโล สพฺพธมฺมานํ,}
\csromangathagroup{\csromangathastanza{\csromangathaitembat{62}{ยเมตํ ปญฺหํ อปุจฺฉิ,}{อชิต ตํ วทามิ เต.} \\ \csromangathacontbat{ยตฺถ นามญฺจ รูปญฺจ,}{อเสสํ อุปรุชฺฌติ.} \\ \csromangathacontbat{วิญฺญาณสฺส นิโรเธน,}{เอตฺเถตํ อุปรุชฺฌติ.}}\csromangathastanza{\csromangathaitembat{63}{เย จ สงฺขาต\-ธมฺมาเส,}{เย จ เสขา\csromansharedfootnote{f53d5e3e9498c57d}{เสกฺขา (ก)} ปุถู อิธ.} \\ \csromangathacontbat{เตสํ เม นิปโก อิริยํ,}{ปุฏฺโฐ ปพฺรูหิ มาริส.}}\csromangathastanza{\csromangathaitembat{64}{กาเมสุ นาภิคิชฺเฌยฺย,}{มนสา’นาวิโล สิยา.} \\ \csromangathacontbat{กุสโล สพฺพธมฺมานํ,}{สโต ภิกฺขุ ปริพฺพเชติ.}}}

\vspace{\tipitakaparskip}
\csromancenter{อชิตมาณวปุจฺฉา ปฐมา.}
\csromanmatikaanchor{mtk.4}
\csromanmatikaanchor{mtk.5}
\csromantocmarkref{subsection}{2. ติสฺสเมตฺเตยฺยมาณวปุจฺฉา}{mtk.4}
\bookmark[dest={mtk.4},level=2]{2. ติสฺสเมตฺเตยฺยมาณวปุจฺฉา}
\csromantocmarkref{subsection}{3. ปุณฺณกมาณวปุจฺฉา}{mtk.5}
\bookmark[dest={mtk.5},level=2]{3. ปุณฺณกมาณวปุจฺฉา}
\csromanmarkh{3. ปุณฺณก\-มาณว\-ปุจฺฉา}\csromanmarkhii{2. ติสฺสเมตฺเตยฺยมาณวปุจฺฉา}\csromanheadernomark{2}{3. ปุณฺณก\-มาณว\-ปุจฺฉา \textbf{2. ติสฺสเมตฺเตยฺยมาณวปุจฺฉา}}
\csromanhangingbegin
\hangingheaditem{65}{\csromanfolio{7}โกธ สนฺตุสิโต โลเก, (อิจฺจายสฺมา ติสฺสเมตฺเตยฺโย,)}
\hangingbody{กสฺส โน สนฺติ อิญฺชิตา.}
\hangingbody{โก อุภนฺตมภิญฺญาย, มชฺเฌ มนฺตา น ลิปฺปติ1.}
\hangingbody{กํ พฺรูสิ มหาปุริโสติ, โก อิธ สิพฺพินิมจฺจคาติ2.}
\csromanhangingend

\csromanhangingbegin
\hangingheaditem{66}{กาเมสุ พฺรหฺมจริยวา, (เมตฺเตยฺยาติ ภควา,)}
\hangingbody{วีตตโณฺห สทา สโต.}
\hangingbody{สงฺขาย นิพฺพุโต ภิกฺขุ, ตสฺส โน สนฺติ อิญฺชิตา.}
\csromanhangingend

\setlength{\csromangathapagewidth}{0pt}
\setlength{\csromangathawakwidth}{0pt}
\csromangathameasuregroup{โส อุภนฺตมภิญฺญาย, \\ ตํ พฺรูมิ มหาปุริโสติ,}{\csromangathabat{โส อุภนฺตมภิญฺญาย,}{มชฺเฌ มนฺตา น ลิปฺปติ.} \\ \csromangathabat{ตํ พฺรูมิ มหาปุริโสติ,}{โส อิธ สิพฺพินิมจฺจคาติ.}}
\csromangathasetleft{โส อุภนฺตมภิญฺญาย, \\ ตํ พฺรูมิ มหาปุริโสติ,}
\csromangathagroup{\csromangathastanza{\csromangathaitembat{67}{โส อุภนฺตม\-ภิญฺญาย,}{มชฺเฌ มนฺตา น ลิปฺปติ.} \\ \csromangathacontbat{ตํ พฺรูมิ มหาปุริโสติ,}{โส อิธ สิพฺพินิม\-จฺจคาติ.}}}

\vspace{\tipitakaparskip}
\csromancenter{ติสฺสเมตฺเตยฺยมาณวปุจฺฉา ทุติยา.}
\csromanheader{2}{3. ปุณฺณก\-มาณว\-ปุจฺฉา}
\proseitem{68}{อเนชํ มูลทสฺสาวิํ, (อิจฺจายสฺมา ปุณฺณโก,)}

\prose{อตฺถิ ปเญฺหน อาคมํ.}

\prose{กิํ นิสฺสิตา อิสโย มนุชา,}

\proseitem{69}{เย เกจิเม อิสโย มนุชา, (ปุณฺณกาติ ภควา,)}

\prose{อาสีสมานา ปุณฺณก อิตฺตตฺตํ.}

\prose{ชรํ สิตา ยญฺญม\-กปฺปยิํสุ.}

\proseitem{70}{\csromanfolio{8}เย เกจิเม อิสโย มนุชา, (อิจฺจายสฺมา ปุณฺณโก,)}

\prose{กจฺจิสุ เต ภควา ยญฺญปเถ อปฺปมตฺตา.}

\prose{อตารุํ ชาติญฺจ ชรญฺจ มาริส,}

\csromanhangingbegin
\hangingheaditem{71}{อาสีสนฺติ โถมยนฺติ อภิชปฺปนฺติ ชุหนฺติ, (ปุณฺณกาติ ภควา,)}
\hangingbody{กามาภิชปฺปนฺติ ปฏิจฺจ ลาภํ.}
\hangingbody{เต ยาชโยคา ภวราครตฺตา,}
\hangingbody{นาตริํสุ ชาติชรนฺติ พฺรูมิ.}
\csromanhangingend

\proseitem{72}{เต เจ นาตริํสุ ยาชโยคา, (อิจฺจายสฺมา ปุณฺณโก,)}

\prose{ยญฺเญหิ ชาติญฺจ ชรญฺจ มาริส.}

\setlength{\csromangathapagewidth}{0pt}
\setlength{\csromangathawakwidth}{0pt}
\csromangathameasuregroup{}{อถ โก จรหิ เทวมนุสฺสโลเก, \\ อตาริ ชาติญฺจ ชรญฺจ มาริส.}
\csromangathameasurewak{อถ โก จรหิ เทวมนุสฺสโลเก, \\ อตาริ ชาติญฺจ ชรญฺจ มาริส.}
\setlength{\csromangathaleftcol}{0pt}
\csromangathagroup{\csromangathastanza{\csromangathawak{อถ โก จรหิ เทว\-มนุสฺส\-โลเก, \\ อตาริ ชาติญฺจ ชรญฺจ มาริส.}}}

\csromanhangingbegin
\hangingheaditem{73}{สงฺขาย โลกสฺมิ ปโรปรานิ, (ปุณฺณกาติ ภควา,)}
\hangingbody{ยสฺสิญฺชิตํ นตฺถิ กุหิญฺจิ โลเก.}
\hangingbody{สนฺโต วิธูโม อนีโฆ นิราโส,}
\hangingbody{อตาริ โส ชาติชรนฺติ พฺรูมีติ.}
\csromanhangingend

\vspace{\tipitakaparskip}
\csromancenter{ปุณฺณก\-มาณว\-ปุจฺฉา ตติยา.}
\csromanmatikaanchor{mtk.6}
\csromantocmarkref{subsection}{4. เมตฺตคูมาณวปุจฺฉา}{mtk.6}
\bookmark[dest={mtk.6},level=2]{4. เมตฺตคูมาณวปุจฺฉา}
\csromanheader{2}{4. \textbf{เมตฺตคูมาณวปุจฺฉา}}
\csromanhangingbegin
\hangingheaditem{74}{ปุจฺฉามิ ตํ ภควา พฺรูหิ เม’ตํ, (อิจฺจายสฺมา เมตฺตคู,)}
\hangingbody{มญฺญามิ ตํ เวทคุํ ภาวิตตฺตํ.}
\hangingbody{กุโต นุ ทุกฺขา สมุทาคตา อิเม,}
\hangingbody{เย เกจิ โลกสฺมิ’มเนกรูปา.}
\csromanhangingend

\csromanheader{2}{4. เมตฺตคูมาณวปุจฺฉา}
\csromanhangingbegin
\hangingheaditem{75}{\csromanfolio{9}ทุกฺขสฺส เว มํ ปภวํ อปุจฺฉสิ, (เมตฺตคูติ ภควา,)}
\hangingbody{ตํ เต ปวกฺขามิ ยถา ปชานํ.}
\hangingbody{อุปธินิทานา ปภวนฺติ ทุกฺขา,}
\hangingbody{เย เกจิ โลกสฺมิ’มเนกรูปา.}
\csromanhangingend

\setlength{\csromangathapagewidth}{0pt}
\setlength{\csromangathawakwidth}{0pt}
\csromangathameasuregroup{}{โย เว อวิทฺวา อุปธิํ กโรติ, \\ ปุนปฺปุนํ ทุกฺข’มุเปติ มนฺโท. \\ ตสฺมา ปชานํ อุปธิํ น กยิรา, \\ ทุกฺขสฺส ชาติปฺปภวานุปสฺสี.}
\csromangathameasurewak{โย เว อวิทฺวา อุปธิํ กโรติ, \\ ปุนปฺปุนํ ทุกฺข’มุเปติ มนฺโท. \\ ตสฺมา ปชานํ อุปธิํ น กยิรา, \\ ทุกฺขสฺส ชาติปฺปภวานุปสฺสี.}
\setlength{\csromangathaleftcol}{0pt}
\csromangathagroup{\csromangathastanza{\csromangathawak{\csromangathaitemline{76}{โย เว อวิทฺวา อุปธิํ กโรติ,} \\ \csromangathacontline{ปุนปฺปุนํ ทุกฺข’มุเปติ มนฺโท.} \\ \csromangathacontline{ตสฺมา ปชานํ อุปธิํ น กยิรา,} \\ \csromangathacontline{ทุกฺขสฺส ชาติปฺปภวา\-นุปสฺสี.}}}}

\proseitem{77}{ยํ ตํ อปุจฺฉิมฺห อกิตฺตยี โน,}

\prose{อญฺญํ ตํ ปุจฺฉาม ตทิงฺฆ พฺรูหิ.}

\setlength{\csromangathapagewidth}{0pt}
\setlength{\csromangathawakwidth}{0pt}
\csromangathameasuregroup{}{กถํ นุ ธีรา วิตรนฺติ โอฆํ, \\ ชาติํ ชรํ โสกปริทฺทวญฺจ. \\ ตํ เม มุนิ สาธุ วิยากโรหิ,}
\csromangathameasurewak{กถํ นุ ธีรา วิตรนฺติ โอฆํ, \\ ชาติํ ชรํ โสกปริทฺทวญฺจ. \\ ตํ เม มุนิ สาธุ วิยากโรหิ,}
\setlength{\csromangathaleftcol}{0pt}
\csromangathagroup{\csromangathastanza{\csromangathawak{กถํ นุ ธีรา วิตรนฺติ โอฆํ, \\ ชาติํ ชรํ โสกปริทฺทวญฺจ. \\ ตํ เม มุนิ สาธุ วิยากโรหิ,}}}

\csromanhangingbegin
\hangingheaditem{78}{กิตฺตยิสฺสามิ เต ธมฺมํ, (เมตฺตคูติ ภควา,)}
\hangingbody{ทิฏฺเฐ ธมฺเม อนีติหํ.}
\hangingbody{ยํ วิทิตฺวา สโต จรํ, ตเร โลเก วิสตฺติกํ.}
\csromanhangingend

\setlength{\csromangathapagewidth}{0pt}
\setlength{\csromangathawakwidth}{0pt}
\csromangathameasuregroup{ตญฺจาหํ อภินนฺทามิ, \\ ยํ วิทิตฺวา สโต จรํ,}{\csromangathabat{ตญฺจาหํ อภินนฺทามิ,}{มเหสิ ธมฺมมุตฺตมํ.} \\ \csromangathabat{ยํ วิทิตฺวา สโต จรํ,}{ตเร โลเก วิสตฺติกํ.}}
\csromangathasetleft{ตญฺจาหํ อภินนฺทามิ, \\ ยํ วิทิตฺวา สโต จรํ,}
\csromangathagroup{\csromangathastanza{\csromangathaitembat{79}{ตญฺจาหํ อภินนฺทามิ,}{มเหสิ ธมฺมมุตฺตมํ.} \\ \csromangathacontbat{ยํ วิทิตฺวา สโต จรํ,}{ตเร โลเก วิสตฺติกํ.}}}

\csromanhangingbegin
\hangingheaditem{80}{ยํ กิญฺจิ สมฺปชานาสิ, (เมตฺตคูติ ภควา,)}
\hangingbody{อุทฺธํ อโธ ติริยญฺจาปิ มชฺเฌ.}
\hangingbody{เอเตสุ นนฺทิญฺจ นิเวสนญฺจ,}
\hangingbody{ปนุชฺช วิญฺญาณํ ภเว น ติฏฺเฐ.}
\csromanhangingend

\setlength{\csromangathapagewidth}{0pt}
\setlength{\csromangathawakwidth}{0pt}
\csromangathameasuregroup{}{เอวํวิหารี สโต อปฺปมตฺโต, \\ ภิกฺขุ จรํ หิตฺวา มมายิตานิ. \\ ชาติํ ชรํ โสกปริทฺทวญฺจ, \\ อิเธว วิทฺวา ปชเหยฺย ทุกฺขํ.}
\csromangathameasurewak{เอวํวิหารี สโต อปฺปมตฺโต, \\ ภิกฺขุ จรํ หิตฺวา มมายิตานิ. \\ ชาติํ ชรํ โสกปริทฺทวญฺจ, \\ อิเธว วิทฺวา ปชเหยฺย ทุกฺขํ.}
\setlength{\csromangathaleftcol}{0pt}
\csromangathagroup{\csromangathastanza{\csromangathawak{\csromangathaitemline{81}{เอวํวิหารี สโต อปฺปมตฺโต,} \\ \csromangathacontline{ภิกฺขุ จรํ หิตฺวา มมายิตานิ.} \\ \csromangathacontline{ชาติํ ชรํ โสกปริทฺทวญฺจ,} \\ \csromangathacontline{อิเธว วิทฺวา ปชเหยฺย ทุกฺขํ.}}}}

\proseitem{82}{\csromanfolio{10}เอตา’ภินนฺทามิ วโจ มเหสิโน,}

\prose{สุกิตฺติตํ โคตม’นูปธีกํ.}

\prose{อทฺธา หิ ภควา ปหาสิ ทุกฺขํ,}

\prose{(9)}

\proseitem{83}{เต จาปิ นูนปฺปชเหยฺยุ ทุกฺขํ,}

\prose{เย ตฺวํ มุนิ อฏฺฐิตํ โอวเทยฺย.}

\prose{ตํ ตํ นมสฺสามิ สเมจฺจ นาค,}

\prose{อปฺเปว มํ ภควา อฏฺฐิตํ โอวเทยฺย. (10)}

\proseitem{84}{ยํ พฺราหฺมณํ เวทคุ’มาภิชญฺญา,}

\prose{อกิญฺจนํ กามภเว อสตฺตํ.}

\prose{อทฺธา หิ โส โอฆมิมํ อตาริ,}

\prose{ติณฺโณ จ ปารํ อขิโล อกงฺโข. (11)}

\proseitem{85}{วิทฺวา จ โย เวทคู นโร อิธ,}

\prose{ภวาภเว สงฺคมิมํ วิสชฺช.}

\prose{โส วีตตโณฺห อนีโฆ นิราโส,}

\prose{อตาริ โส ชาติชรนฺติ พฺรูมีติ. (12)}

\vspace{\tipitakaparskip}
\csromancenter{เมตฺตคูมาณวปุจฺฉา จตุตฺถี.}
\csromanmatikaanchor{mtk.7}
\csromantocmarkref{subsection}{5. โธตกมาณวปุจฺฉา}{mtk.7}
\bookmark[dest={mtk.7},level=2]{5. โธตกมาณวปุจฺฉา}
\csromanheader{2}{5. \textbf{โธตกมาณวปุจฺฉา}}
\csromanhangingbegin
\hangingheaditem{86}{ปุจฺฉามิ ตํ ภควา พฺรูหิ เม’ตํ, (อิจฺจายสฺมา โธตโก,)}
\hangingbody{วาจา’ภิกงฺขามิ มเหสิ ตุยฺหํ.}
\hangingbody{ตว สุตฺวาน นิคฺโฆสํ, สิกฺเข นิพฺพานมตฺตโน.}
\csromanhangingend

\csromanhangingbegin
\hangingheaditem{87}{เตนหา’ตปฺปํ กโรหิ, (โธตกาติ ภควา,)}
\hangingbody{อิเธว นิปโก สโต.}
\hangingbody{อิโต สุตฺวาน นิคฺโฆสํ, สิกฺเข นิพฺพานมตฺตโน.}
\csromanhangingend

\setlength{\csromangathapagewidth}{0pt}
\setlength{\csromangathawakwidth}{0pt}
\csromangathameasuregroup{}{ปสฺสามหํ เทวมนุสฺสโลเก, \\ อกิญฺจนํ พฺราหฺมณ’มิริยมานํ. \\ ตํ ตํ นมสฺสามิ สมนฺตจกฺขุ, \\ ปมุญฺจ มํ สกฺก กถํกถาหิ.}
\csromangathameasurewak{ปสฺสามหํ เทวมนุสฺสโลเก, \\ อกิญฺจนํ พฺราหฺมณ’มิริยมานํ. \\ ตํ ตํ นมสฺสามิ สมนฺตจกฺขุ, \\ ปมุญฺจ มํ สกฺก กถํกถาหิ.}
\setlength{\csromangathaleftcol}{0pt}
\csromangathagroup{\csromangathastanza{\csromangathawak{\csromangathaitemline{88}{ปสฺสามหํ เทว\-มนุสฺส\-โลเก,} \\ \csromangathacontline{อกิญฺจนํ พฺราหฺมณ’มิริยมานํ.} \\ \csromangathacontline{ตํ ตํ นมสฺสามิ สมนฺตจกฺขุ,} \\ \csromangathacontline{ปมุญฺจ มํ สกฺก กถํกถาหิ.}}}}

\csromanmatikaanchor{mtk.8}
\csromantocmarkref{subsection}{6. อุปสีวมาณวปุจฺฉา}{mtk.8}
\bookmark[dest={mtk.8},level=2]{6. อุปสีวมาณวปุจฺฉา}
\csromanheader{2}{6. อุปสีวมาณวปุจฺฉา}
\setlength{\csromangathapagewidth}{0pt}
\setlength{\csromangathawakwidth}{0pt}
\csromangathameasuregroup{}{นาหํ สหิสฺสามิ ปโมจนาย, \\ กถํกถิํ โธตก กญฺจิ โลเก. \\ ธมฺมญฺจ เสฏฺฐํ อภิชานมาโน\textsuperscript{1}, \\ เอวํ ตุวํ โอฆมิมํ ตเรสิ. \\ อนุสาส พฺรหฺเม กรุณายมาโน, \\ วิเวกธมฺมํ ยมหํ วิชญฺญํ. \\ ยถาหํ อากาโสว อพฺยาปชฺชมาโน, \\ อิเธว สนฺโต อสิโต จเรยฺยํ.}
\csromangathameasurewak{นาหํ สหิสฺสามิ ปโมจนาย, \\ กถํกถิํ โธตก กญฺจิ โลเก. \\ ธมฺมญฺจ เสฏฺฐํ อภิชานมาโน\textsuperscript{1}, \\ เอวํ ตุวํ โอฆมิมํ ตเรสิ. \\ อนุสาส พฺรหฺเม กรุณายมาโน, \\ วิเวกธมฺมํ ยมหํ วิชญฺญํ. \\ ยถาหํ อากาโสว อพฺยาปชฺชมาโน, \\ อิเธว สนฺโต อสิโต จเรยฺยํ.}
\setlength{\csromangathaleftcol}{0pt}
\csromangathagroup{\csromangathastanza{\csromangathawak{\csromanfolio{11}\csromangathaitemline{89}{นาหํ สหิสฺสามิ ปโมจนาย,} \\ \csromangathacontline{กถํกถิํ โธตก กญฺจิ โลเก.} \\ \csromangathacontline{ธมฺมญฺจ เสฏฺฐํ อภิชานมาโน\csromansharedfootnote{139df5526e3959e6}{อาชานมาโน (สี, สฺยา, อิ)},} \\ \csromangathacontline{เอวํ ตุวํ โอฆมิมํ ตเรสิ.}}}\csromangathastanza{\csromangathawak{\csromangathaitemline{90}{อนุสาส พฺรหฺเม กรุณายมาโน,} \\ \csromangathacontline{วิเวกธมฺมํ ยมหํ วิชญฺญํ.} \\ \csromangathacontline{ยถาหํ อากาโสว อพฺยาปชฺชมาโน,} \\ \csromangathacontline{อิเธว สนฺโต อสิโต จเรยฺยํ.}}}}

\csromanhangingbegin
\hangingheaditem{91}{กิตฺตยิสฺสามิ เต สนฺติํ, (โธตกาติ ภควา,)}
\hangingbody{ทิฏฺเฐ ธมฺเม อนีติหํ.}
\hangingbody{ยํ วิทิตฺวา สโต จรํ,}
\hangingbody{ตเร โลเก วิสตฺติกํ.}
\csromanhangingend

\setlength{\csromangathapagewidth}{0pt}
\setlength{\csromangathawakwidth}{0pt}
\csromangathameasuregroup{ตญฺจาหํ อภินนฺทามิ, \\ ยํ วิทิตฺวา สโต จรํ,}{\csromangathabat{ตญฺจาหํ อภินนฺทามิ,}{มเหสิ สนฺติ’มุตฺตมํ.} \\ \csromangathabat{ยํ วิทิตฺวา สโต จรํ,}{ตเร โลเก วิสตฺติกํ.}}
\csromangathasetleft{ตญฺจาหํ อภินนฺทามิ, \\ ยํ วิทิตฺวา สโต จรํ,}
\csromangathagroup{\csromangathastanza{\csromangathaitembat{92}{ตญฺจาหํ อภินนฺทามิ,}{มเหสิ สนฺติ’มุตฺตมํ.} \\ \csromangathacontbat{ยํ วิทิตฺวา สโต จรํ,}{ตเร โลเก วิสตฺติกํ.}}}

\csromanhangingbegin
\hangingheaditem{93}{ยํ กิญฺจิ สมฺปชานาสิ, (โธตกาติ ภควา,)}
\hangingbody{อุทฺธํ อโธ ติริยญฺจาปิ มชฺเฌ,}
\hangingbody{เอตํ วิทิตฺวา สงฺโคติ โลเก.}
\hangingbody{ภวาภวาย มากาสิ ตณฺหนฺติ.}
\csromanhangingend

\vspace{\tipitakaparskip}
\csromancenter{โธตกมาณวปุจฺฉา ปญฺจมี.}
\csromanheader{2}{6. อุปสีวมาณวปุจฺฉา}
\csromanhangingbegin
\hangingheaditem{94}{เอโก อหํ สกฺก มหนฺตโมฆํ, (อิจฺจายสฺมา อุปสีโว,)}
\hangingbody{อนิสฺสิโต โน วิสหามิ ตาริตุํ.}
\hangingbody{อารมฺมณํ พฺรูหิ สมนฺตจกฺขุ,}
\hangingbody{ยํ นิสฺสิโต โอฆมิมํ ตเรยฺยํ.}
\csromanhangingend

\csromanhangingbegin
\hangingheaditem{95}{\csromanfolio{12}อากิญฺจญฺญํ เปกฺขมาโน สติมา, (อุปสีวาติ ภควา,)}
\hangingbody{นตฺถีติ นิสฺสาย ตรสฺสุ โอฆํ.}
\hangingbody{กาเม ปหาย วิรโต กถาหิ,}
\hangingbody{ตณฺหกฺขยํ นตฺตมหา’ภิปสฺส.}
\csromanhangingend

\csromanhangingbegin
\hangingheaditem{96}{สพฺเพสุ กาเมสุ โย วีตราโค, (อิจฺจายสฺมา อุปสีโว,)}
\hangingbody{อากิญฺจญฺญํ นิสฺสิโต หิตฺวา มญฺญํ.}
\hangingbody{สญฺญาวิโมกฺเข ปรเม วิมุตฺโต1,}
\hangingbody{ติฏฺเฐ นุ โส ตตฺถ อนานุยายี2.}
\csromanhangingend

\csromanhangingbegin
\hangingheaditem{97}{สพฺเพสุ กาเมสุ โย วีตราโค, (อุปสีวาติ ภควา,)}
\hangingbody{อากิญฺจญฺญํ นิสฺสิโต หิตฺวา มญฺญํ.}
\hangingbody{สญฺญาวิโมกฺเข ปรเม วิมุตฺโต,}
\hangingbody{ติฏฺเฐยฺย โส ตตฺถ อนานุยายี.}
\csromanhangingend

\setlength{\csromangathapagewidth}{0pt}
\setlength{\csromangathawakwidth}{0pt}
\csromangathameasuregroup{}{ติฏฺเฐ เจ โส ตตฺถ อนานุยายี, \\ ปูคมฺปิ วสฺสานํ สมนฺตจกฺขุ. \\ ตตฺเถว โส สีติสิยา วิมุตฺโต, \\ จเวถ วิญฺญาณํ ตถาวิธสฺส.}
\csromangathameasurewak{ติฏฺเฐ เจ โส ตตฺถ อนานุยายี, \\ ปูคมฺปิ วสฺสานํ สมนฺตจกฺขุ. \\ ตตฺเถว โส สีติสิยา วิมุตฺโต, \\ จเวถ วิญฺญาณํ ตถาวิธสฺส.}
\setlength{\csromangathaleftcol}{0pt}
\csromangathagroup{\csromangathastanza{\csromangathawak{\csromangathaitemline{98}{ติฏฺเฐ เจ โส ตตฺถ อนานุยายี,} \\ \csromangathacontline{ปูคมฺปิ วสฺสานํ สมนฺตจกฺขุ.} \\ \csromangathacontline{ตตฺเถว โส สีติสิยา วิมุตฺโต,} \\ \csromangathacontline{จเวถ วิญฺญาณํ ตถาวิธสฺส.}}}}

\csromanhangingbegin
\hangingheaditem{99}{อจฺจิ ยถา วาตเวเคน ขิตฺตา, (อุปสีวาติ ภควา,)}
\hangingbody{อตฺถํ ปเลติ น อุเปติ สงฺขํ.}
\hangingbody{เอวํ มุนี นามกายา วิมุตฺโต,}
\hangingbody{อตฺถํ ปเลติ น อุเปติ สงฺขํ.}
\csromanhangingend

\proseitem{100}{อตฺถงฺคโต โส อุท วา โส นตฺถิ,}

\prose{อุทาหุ เว สสฺสติยา อโรโค.}

\prose{ตํ เม มุนี สาธุ วิยากโรหิ,}

\csromanhangingbegin
\hangingheaditem{101}{อตฺถงฺคตสฺส น ปมาณ’มตฺถิ, (อุปสีวาติ ภควา,)}
\hangingbody{เยน นํ วชฺชุํ ตํ ตสฺส นตฺถิ.}
\hangingbody{สพฺเพสุ ธมฺเมสุ สมูหเตสุ,}
\hangingbody{สมูหตา วาทปถาปิ สพฺเพติ.}
\csromanhangingend

\vspace{\tipitakaparskip}
\csromancenter{อุปสีวมาณวปุจฺฉา ฉฏฺฐี.}
\csromanmatikaanchor{mtk.9}
\csromantocmarkref{subsection}{7. นนฺทมาณวปุจฺฉา}{mtk.9}
\bookmark[dest={mtk.9},level=2]{7. นนฺทมาณวปุจฺฉา}
\csromanheader{2}{7. นนฺทมาณวปุจฺฉา}
\csromanheader{2}{7. นนฺทมาณวปุจฺฉา}
\csromanhangingbegin
\hangingheaditem{102}{\csromanfolio{13}“สนฺติ โลเก มุนโย”, (อิจฺจายสฺมา นนฺโท,)}
\hangingbody{ชนา วทนฺติ ตยิทํ กถํสุ.}
\hangingbody{ญาณูปปนฺนํ มุนิ โน วทนฺติ,}
\hangingbody{อุทาหุ เว ชีวิเตนูปปนฺนํ.}
\csromanhangingend

\setlength{\csromangathapagewidth}{0pt}
\setlength{\csromangathawakwidth}{0pt}
\csromangathameasuregroup{น ทิฏฺฐิยา น สุติยา น ญาเณน,}{\csromangathabat{น ทิฏฺฐิยา น สุติยา น ญาเณน,}{มุนีธ นนฺท กุสลา วทนฺติ.} \\ วิเสนิกตฺวา อนีฆา นิราสา. \\ จรนฺติ เย เต มุนโยติ พฺรูมิ.}
\csromangathasetleft{น ทิฏฺฐิยา น สุติยา น ญาเณน,}
\csromangathagroup{\csromangathastanza{\csromangathaitembat{103}{น ทิฏฺฐิยา น สุติยา น ญาเณน,}{มุนีธ นนฺท กุสลา วทนฺติ.} \\ \csromangathacontline{วิเสนิกตฺวา อนีฆา นิราสา.} \\ \csromangathacontline{จรนฺติ เย เต มุนโยติ พฺรูมิ.}}}

\proseitem{104}{เย เกจิเม สมณ\-พฺราหฺมณาเส, (อิจฺจายสฺมา นนฺโท,)}

\prose{ทิฏฺฐสฺสุเตนาปิ วทนฺติ สุทฺธิํ.}

\setlength{\csromangathapagewidth}{0pt}
\setlength{\csromangathawakwidth}{0pt}
\csromangathameasuregroup{}{สีลพฺพเตนาปิ วทนฺติ สุทฺธิํ, \\ อเนกรูเปน วทนฺติ สุทฺธิํ. \\ กจฺจิสฺสุ เต ภควา ตตฺถ ยตา จรนฺตา, \\ อตารุ ชาติญฺจ ชรญฺจ มาริส.}
\csromangathameasurewak{สีลพฺพเตนาปิ วทนฺติ สุทฺธิํ, \\ อเนกรูเปน วทนฺติ สุทฺธิํ. \\ กจฺจิสฺสุ เต ภควา ตตฺถ ยตา จรนฺตา, \\ อตารุ ชาติญฺจ ชรญฺจ มาริส.}
\setlength{\csromangathaleftcol}{0pt}
\csromangathagroup{\csromangathastanza{\csromangathawak{สีลพฺพเตนาปิ วทนฺติ สุทฺธิํ, \\ อเนกรูเปน วทนฺติ สุทฺธิํ. \\ กจฺจิสฺสุ เต ภควา ตตฺถ ยตา จรนฺตา, \\ อตารุ ชาติญฺจ ชรญฺจ มาริส.}}}

\csromanhangingbegin
\hangingheaditem{105}{เย เกจิเม สมณ\-พฺราหฺมณาเส, (นนฺทาติ ภควา,)}
\hangingbody{ทิฏฺฐสฺสุเตนาปิ วทนฺติ สุทฺธิํ.}
\hangingbody{สีลพฺพเตนาปิ วทนฺติ สุทฺธิํ,}
\hangingbody{อเนกรูเปน วทนฺติ สุทฺธิํ.}
\hangingbody{กิญฺจาปิ เต ตตฺถ ยตา จรนฺติ,}
\hangingbody{นาตริํสุ ชาติชรนฺติ พฺรูมิ.}
\csromanhangingend

\proseitem{106}{เย เกจิเม สมณ\-พฺราหฺมณาเส, (อิจฺจายสฺมา นนฺโท,)}

\prose{ทิฏฺฐสฺสุเตนาปิ วทนฺติ สุทฺธิํ.}

\setlength{\csromangathapagewidth}{0pt}
\setlength{\csromangathawakwidth}{0pt}
\csromangathameasuregroup{}{สีลพฺพเตนาปิ วทนฺติ สุทฺธิํ, \\ อเนกรูเปน วทนฺติ สุทฺธิํ. \\ เต เจ มุนิ พฺรูสิ อโนฆติณฺเณ, \\ อถ โก จรหิ เทวมนุสฺสโลเก. \\ อตาริ ชาติญฺจ ชรญฺจ มาริส,}
\csromangathameasurewak{สีลพฺพเตนาปิ วทนฺติ สุทฺธิํ, \\ อเนกรูเปน วทนฺติ สุทฺธิํ. \\ เต เจ มุนิ พฺรูสิ อโนฆติณฺเณ, \\ อถ โก จรหิ เทวมนุสฺสโลเก. \\ อตาริ ชาติญฺจ ชรญฺจ มาริส,}
\setlength{\csromangathaleftcol}{0pt}
\csromangathagroup{\csromangathastanza{\csromangathawak{สีลพฺพเตนาปิ วทนฺติ สุทฺธิํ, \\ อเนกรูเปน วทนฺติ สุทฺธิํ. \\ เต เจ มุนิ พฺรูสิ อโนฆติณฺเณ, \\ อถ โก จรหิ เทว\-มนุสฺส\-โลเก.}}\csromangathastanza{\csromangathawak{อตาริ ชาติญฺจ ชรญฺจ มาริส,}}}

\csromanhangingbegin
\hangingheaditem{107}{\csromanfolio{14}นาหํ สพฺเพ สมณ\-พฺราหฺมณาเส, (นนฺทาติ ภควา,)}
\hangingbody{ชาติชราย นิวุตาติ พฺรูมิ.}
\hangingbody{เย สีธ ทิฏฺฐํ ว สุตํ มุตํ วา,}
\hangingbody{สีลพฺพตํ วาปิ ปหาย สพฺพํ.}
\hangingbody{อเนกรูปมฺปิ ปหาย สพฺพํ,}
\hangingbody{ตณฺหํ ปริญฺญาย อนาสวาเส.}
\hangingbody{เต เว นรา โอฆติณฺณาติ พฺรูมิ.}
\csromanhangingend

\setlength{\csromangathapagewidth}{0pt}
\setlength{\csromangathawakwidth}{0pt}
\csromangathameasuregroup{}{เอตา’ภินนฺทามิ วโจ มเหสิโน, \\ สุกิตฺติตํ โคตม’นูปธีกํ. \\ เย สีธ ทิฏฺฐํ ว สุตํ มุตํ วา, \\ สีลพฺพตํ วาปิ ปหาย สพฺพํ. \\ อเนกรูปมฺปิ ปหาย สพฺพํ, \\ ตณฺหํ ปริญฺญาย อนาสวาเส. \\ อหมฺปิ เต โอฆติณฺณาติ พฺรูมีติ.}
\csromangathameasurewak{เอตา’ภินนฺทามิ วโจ มเหสิโน, \\ สุกิตฺติตํ โคตม’นูปธีกํ. \\ เย สีธ ทิฏฺฐํ ว สุตํ มุตํ วา, \\ สีลพฺพตํ วาปิ ปหาย สพฺพํ. \\ อเนกรูปมฺปิ ปหาย สพฺพํ, \\ ตณฺหํ ปริญฺญาย อนาสวาเส. \\ อหมฺปิ เต โอฆติณฺณาติ พฺรูมีติ.}
\setlength{\csromangathaleftcol}{0pt}
\csromangathagroup{\csromangathastanza{\csromangathawak{\csromangathaitemline{108}{เอตา’ภินนฺทามิ วโจ มเหสิโน,} \\ \csromangathacontline{สุกิตฺติตํ โคตม’นูปธีกํ.} \\ \csromangathacontline{เย สีธ ทิฏฺฐํ ว สุตํ มุตํ วา,} \\ \csromangathacontline{สีลพฺพตํ วาปิ ปหาย สพฺพํ.} \\ \csromangathacontline{อเนกรูปมฺปิ ปหาย สพฺพํ,} \\ \csromangathacontline{ตณฺหํ ปริญฺญาย อนาสวาเส.} \\ \csromangathacontline{อหมฺปิ เต โอฆติณฺณาติ พฺรูมีติ.}}}}

\vspace{\tipitakaparskip}
\csromancenter{นนฺทมาณวปุจฺฉา สตฺตมา.}
\csromanmatikaanchor{mtk.10}
\csromantocmarkref{subsection}{8. เหมกมาณวปุจฺฉา}{mtk.10}
\bookmark[dest={mtk.10},level=2]{8. เหมกมาณวปุจฺฉา}
\csromanheader{2}{8. \textbf{เหมกมาณวปุจฺฉา}}
\csromanhangingbegin
\hangingheaditem{109}{เย เม ปุพฺเพ วิยากํสุ, (อิจฺจายสฺมา เหมโก,)}
\hangingbody{หุรํ โคตมสาสนา.}
\hangingbody{“อิจฺจาสิ อิติ ภวิสฺสติ”, สพฺพํ ตํ อิติหีติหํ.}
\hangingbody{สพฺพํ ตํ ตกฺกวฑฺฒนํ, นาหํ ตตฺถ อภิรมิํ. (1)}
\csromanhangingend

\setlength{\csromangathapagewidth}{0pt}
\setlength{\csromangathawakwidth}{0pt}
\csromangathameasuregroup{ตฺวญฺจ เม ธมฺมมกฺขาหิ, \\ ยํ วิทิตฺวา สโต จรํ, \\ อิธ ทิฏฺฐสุตมุตวิญฺญาเตสุ, \\ ฉนฺทราควิโนทนํ, \\ เอตทญฺญาย เย สตา, \\ อุปสนฺตา จ เต สทา,}{\csromangathabat{ตฺวญฺจ เม ธมฺมมกฺขาหิ,}{ตณฺหานิคฺฆาตนํ มุนิ.} \\ \csromangathabat{ยํ วิทิตฺวา สโต จรํ,}{ตเร โลเก วิสตฺติกํ.} \\ \csromangathabat{อิธ ทิฏฺฐสุตมุตวิญฺญาเตสุ,}{ปิยรูเปสุ เหมก.} \\ \csromangathabat{ฉนฺทราควิโนทนํ,}{นิพฺพานปท’มจฺจุตํ.} \\ \csromangathabat{เอตทญฺญาย เย สตา,}{ทิฏฺฐธมฺมา’ภินิพฺพุตา.} \\ \csromangathabat{อุปสนฺตา จ เต สทา,}{ติณฺณา โลเก วิสตฺติกนฺติ.}}
\csromangathasetleft{ตฺวญฺจ เม ธมฺมมกฺขาหิ, \\ ยํ วิทิตฺวา สโต จรํ, \\ อิธ ทิฏฺฐสุตมุตวิญฺญาเตสุ, \\ ฉนฺทราควิโนทนํ, \\ เอตทญฺญาย เย สตา, \\ อุปสนฺตา จ เต สทา,}
\csromangathagroup{\csromangathastanza{\csromangathaitembat{110}{ตฺวญฺจ เม ธมฺมมกฺขาหิ,}{ตณฺหา\-นิคฺฆาตนํ มุนิ.} \\ \csromangathacontbat{ยํ วิทิตฺวา สโต จรํ,}{ตเร โลเก วิสตฺติกํ.}}\csromangathastanza{\csromangathaitembat{111}{อิธ ทิฏฺฐ\-สุต\-มุต\-วิญฺญาเตสุ,}{ปิยรูเปสุ เหมก.} \\ \csromangathacontbat{ฉนฺทราค\-วิโนทนํ,}{นิพฺพานปท’มจฺจุตํ.}}\csromangathastanza{\csromangathaitembat{112}{เอตทญฺญาย เย สตา,}{ทิฏฺฐธมฺมา’ภินิพฺพุตา.} \\ \csromangathacontbat{อุปสนฺตา จ เต สทา,}{ติณฺณา โลเก วิสตฺติกนฺติ.}}}

\vspace{\tipitakaparskip}
\csromancenter{เหมกมาณวปุจฺฉา อฏฺฐมา.}
\csromanmatikaanchor{mtk.11}
\csromanmatikaanchor{mtk.12}
\csromantocmarkref{subsection}{9. โตเทยฺยมาณวปุจฺฉา}{mtk.11}
\bookmark[dest={mtk.11},level=2]{9. โตเทยฺยมาณวปุจฺฉา}
\csromantocmarkref{subsection}{10. กปฺปมาณวปุจฺฉา}{mtk.12}
\bookmark[dest={mtk.12},level=2]{10. กปฺปมาณวปุจฺฉา}
\csromanmarkh{10. กปฺปมาณวปุจฺฉา}\csromanmarkhii{9. โตเทยฺยมาณวปุจฺฉา}\csromanheadernomark{2}{10. กปฺปมาณวปุจฺฉา \textbf{9. โตเทยฺยมาณวปุจฺฉา}}
\csromanhangingbegin
\hangingheaditem{113}{\csromanfolio{15}ยสฺมิํ กามา น วสนฺติ, (อิจฺจายสฺมา โตเทยฺโย,)}
\hangingbody{ตณฺหา ยสฺส น วิชฺชติ.}
\hangingbody{กถํกถา จ โย ติณฺโณ,}
\hangingbody{วิโมกฺโข ตสฺส กีทิโส.}
\csromanhangingend

\csromanhangingbegin
\hangingheaditem{114}{ยสฺมิํ กามา น วสนฺติ, (โตเทยฺยาติ ภควา,)}
\hangingbody{ตณฺหา ยสฺส น วิชฺชติ.}
\hangingbody{กถํกถา จ โย ติณฺโณ,}
\hangingbody{วิโมกฺโข ตสฺส นา’ปโร.}
\csromanhangingend

\setlength{\csromangathapagewidth}{0pt}
\setlength{\csromangathawakwidth}{0pt}
\csromangathameasuregroup{}{นิราสโส โส อุท อาสสาโน\textsuperscript{1}, \\ ปญฺญาณวา โส อุท ปญฺญกปฺปี. \\ มุนิํ อหํ สกฺก ยถา วิชญฺญํ, \\ ตํ เม วิยาจิกฺข สมนฺตจกฺขุ. \\ นิราสโส โส น จ อาสสาโน, \\ ปญฺญาณวา โส น จ ปญฺญกปฺปี. \\ เอวมฺปิ โตเทยฺย มุนิํ วิชาน, \\ อกิญฺจนํ กามภเว อสตฺตนฺติ.}
\csromangathameasurewak{นิราสโส โส อุท อาสสาโน\textsuperscript{1}, \\ ปญฺญาณวา โส อุท ปญฺญกปฺปี. \\ มุนิํ อหํ สกฺก ยถา วิชญฺญํ, \\ ตํ เม วิยาจิกฺข สมนฺตจกฺขุ. \\ นิราสโส โส น จ อาสสาโน, \\ ปญฺญาณวา โส น จ ปญฺญกปฺปี. \\ เอวมฺปิ โตเทยฺย มุนิํ วิชาน, \\ อกิญฺจนํ กามภเว อสตฺตนฺติ.}
\setlength{\csromangathaleftcol}{0pt}
\csromangathagroup{\csromangathastanza{\csromangathawak{\csromangathaitemline{115}{นิราสโส โส อุท อาสสาโน\csromansharedfootnote{b5e9a479eebc47be}{อาสยาโน (ก)},} \\ \csromangathacontline{ปญฺญาณวา โส อุท ปญฺญกปฺปี.} \\ \csromangathacontline{มุนิํ อหํ สกฺก ยถา วิชญฺญํ,} \\ \csromangathacontline{ตํ เม วิยาจิกฺข สมนฺตจกฺขุ.}}}\csromangathastanza{\csromangathawak{\csromangathaitemline{116}{นิราสโส โส น จ อาสสาโน,} \\ \csromangathacontline{ปญฺญาณวา โส น จ ปญฺญกปฺปี.} \\ \csromangathacontline{เอวมฺปิ โตเทยฺย มุนิํ วิชาน,} \\ \csromangathacontline{อกิญฺจนํ กามภเว อสตฺตนฺติ.}}}}

\vspace{\tipitakaparskip}
\csromancenter{โตเทยฺยมาณวปุจฺฉา นวมา.}
\csromanheader{2}{10. กปฺปมาณวปุจฺฉา}
\csromanhangingbegin
\hangingheaditem{117}{มชฺเฌ สรสฺมิํ ติฏฺฐตํ, (อิจฺจายสฺมา กปฺโป,)}
\hangingbody{โอเฆ ชาเต มหพฺภเย.}
\hangingbody{ชรามจฺจุปเรตานํ, ทีปํ ปพฺรูหิ มาริส.}
\hangingbody{ตฺวญฺจ เม ทีปมกฺขาหิ, ยถายิทํ นาปรํ สิยา.}
\csromanhangingend

\csromanhangingbegin
\hangingheaditem{118}{มชฺเฌ สรสฺมิํ ติฏฺฐตํ, (กปฺปาติ ภควา,)}
\hangingbody{โอเฆ ชาเต มหพฺภเย.}
\hangingbody{ชรามจฺจุปเรตานํ, ทีปํ ปพฺรูมิ กปฺป เต.}
\csromanhangingend

\setlength{\csromangathapagewidth}{0pt}
\setlength{\csromangathawakwidth}{0pt}
\csromangathameasuregroup{อกิญฺจนํ อนาทานํ, \\ นิพฺพานํ อิติ นํ พฺรูมิ, \\ เอตทญฺญาย เย สตา, \\ น เต มารวสานุคา,}{\csromangathabat{อกิญฺจนํ อนาทานํ,}{เอตํ ทีปํ อนาปรํ.} \\ \csromangathabat{นิพฺพานํ อิติ นํ พฺรูมิ,}{ชรามจฺจุปริกฺขยํ.} \\ \csromangathabat{เอตทญฺญาย เย สตา,}{ทิฏฺฐธมฺมาภินิพฺพุตา.} \\ \csromangathabat{น เต มารวสานุคา,}{น เต มารสฺส ปฏฺฐคูติ\textsuperscript{1}.}}
\csromangathasetleft{อกิญฺจนํ อนาทานํ, \\ นิพฺพานํ อิติ นํ พฺรูมิ, \\ เอตทญฺญาย เย สตา, \\ น เต มารวสานุคา,}
\csromangathagroup{\csromangathastanza{\csromanfolio{16}\csromangathaitembat{119}{อกิญฺจนํ อนาทานํ,}{เอตํ ทีปํ อนาปรํ.} \\ \csromangathacontbat{นิพฺพานํ อิติ นํ พฺรูมิ,}{ชรามจฺจุ\-ปริกฺขยํ.}}\csromangathastanza{\csromangathaitembat{120}{เอตทญฺญาย เย สตา,}{ทิฏฺฐธมฺมา\-ภินิพฺพุตา.} \\ \csromangathacontbat{น เต มารวสานุคา,}{น เต มารสฺส ปฏฺฐคูติ\csromansharedfootnote{2c6432e3c4a1fe15}{ปทฺธคู (สี)}.}}}

\vspace{\tipitakaparskip}
\csromancenter{กปฺปมาณวปุจฺฉา ทสมา.}
\csromanmatikaanchor{mtk.13}
\csromantocmarkref{subsection}{11. ชตุกณฺณิมาณวปุจฺฉา}{mtk.13}
\bookmark[dest={mtk.13},level=2]{11. ชตุกณฺณิมาณวปุจฺฉา}
\csromanheader{2}{11. \textbf{ชตุกณฺณิมาณวปุจฺฉา}}
\csromanhangingbegin
\hangingheaditem{121}{สุตฺวานหํ วีรม\-กามกามิํ, (อิจฺจายสฺมา ชตุกณฺณิ,)}
\hangingbody{โอฆาติคํ ปุฏฺฐุมกามมาคมํ.}
\hangingbody{สนฺติปทํ พฺรูหิ สหชเนตฺต,}
\hangingbody{ยถาตจฺฉํ ภควา พฺรูหิ เม’ตํ.}
\csromanhangingend

\setlength{\csromangathapagewidth}{0pt}
\setlength{\csromangathawakwidth}{0pt}
\csromangathameasuregroup{}{ภควา หิ กาเม อภิภุยฺย อิริยติ, \\ อาทิจฺโจว ปถวิํ เตชี เตชสา. \\ ปริตฺตปญฺญสฺส เม ภูริปญฺญ, \\ อาจิกฺข ธมฺมํ ยมหํ วิชญฺญํ. \\ ชาติชราย อิธ วิปฺปหานํ.}
\csromangathameasurewak{ภควา หิ กาเม อภิภุยฺย อิริยติ, \\ อาทิจฺโจว ปถวิํ เตชี เตชสา. \\ ปริตฺตปญฺญสฺส เม ภูริปญฺญ, \\ อาจิกฺข ธมฺมํ ยมหํ วิชญฺญํ. \\ ชาติชราย อิธ วิปฺปหานํ.}
\setlength{\csromangathaleftcol}{0pt}
\csromangathagroup{\csromangathastanza{\csromangathawak{\csromangathaitemline{122}{ภควา หิ กาเม อภิภุยฺย อิริยติ,} \\ \csromangathacontline{อาทิจฺโจว ปถวิํ เตชี เตชสา.} \\ \csromangathacontline{ปริตฺต\-ปญฺญสฺส เม ภูริปญฺญ,} \\ \csromangathacontline{อาจิกฺข ธมฺมํ ยมหํ วิชญฺญํ.} \\ \csromangathacontline{ชาติชราย อิธ วิปฺปหานํ.}}}}

\csromanhangingbegin
\hangingheaditem{123}{กาเมสุ วินย เคธํ, (ชตุกณฺณีติ ภควา,)}
\hangingbody{เนกฺขมฺมํ ทฏฺฐุ เขมโต.}
\hangingbody{อุคฺคหิตํ นิรตฺตํ วา, มา เต วิชฺชิตฺถ กิญฺจนํ.}
\csromanhangingend

\setlength{\csromangathapagewidth}{0pt}
\setlength{\csromangathawakwidth}{0pt}
\csromangathameasuregroup{ยํ ปุพฺเพ ตํ วิโสเสหิ, \\ มชฺเฌ เจ โน คเหสฺสสิ, \\ สพฺพโส นามรูปสฺมิํ, \\ อาสวาสฺส น วิชฺชนฺติ,}{\csromangathabat{ยํ ปุพฺเพ ตํ วิโสเสหิ,}{ปจฺฉา เต มาหุ กิญฺจนํ.} \\ \csromangathabat{มชฺเฌ เจ โน คเหสฺสสิ,}{อุปสนฺโต จริสฺสสิ.} \\ \csromangathabat{สพฺพโส นามรูปสฺมิํ,}{วีตเคธสฺส พฺราหฺมณ.} \\ \csromangathabat{อาสวาสฺส น วิชฺชนฺติ,}{เยหิ มจฺจุวสํ วเชติ.}}
\csromangathasetleft{ยํ ปุพฺเพ ตํ วิโสเสหิ, \\ มชฺเฌ เจ โน คเหสฺสสิ, \\ สพฺพโส นามรูปสฺมิํ, \\ อาสวาสฺส น วิชฺชนฺติ,}
\csromangathagroup{\csromangathastanza{\csromangathaitembat{124}{ยํ ปุพฺเพ ตํ วิโสเสหิ,}{ปจฺฉา เต มาหุ กิญฺจนํ.} \\ \csromangathacontbat{มชฺเฌ เจ โน คเหสฺสสิ,}{อุปสนฺโต จริสฺสสิ.}}\csromangathastanza{\csromangathaitembat{125}{สพฺพโส นามรูปสฺมิํ,}{วีตเคธสฺส พฺราหฺมณ.} \\ \csromangathacontbat{อาสวาสฺส น วิชฺชนฺติ,}{เยหิ มจฺจุวสํ วเชติ.}}}

\vspace{\tipitakaparskip}
\csromancenter{ชตุกณฺณิมาณวปุจฺฉา เอกาทสมา.}
\csromanmatikaanchor{mtk.14}
\csromanmatikaanchor{mtk.15}
\csromantocmarkref{subsection}{12. ภทฺราวุธมาณวปุจฺฉา}{mtk.14}
\bookmark[dest={mtk.14},level=2]{12. ภทฺราวุธมาณวปุจฺฉา}
\csromantocmarkref{subsection}{13. อุทยมาณวปุจฺฉา}{mtk.15}
\bookmark[dest={mtk.15},level=2]{13. อุทยมาณวปุจฺฉา}
\csromanmarkh{13. อุทย\-มาณว\-ปุจฺฉา}\csromanmarkhii{12. ภทฺราวุธมาณวปุจฺฉา}\csromanheadernomark{2}{13. อุทย\-มาณว\-ปุจฺฉา \textbf{12. ภทฺราวุธมาณวปุจฺฉา}}
\csromanhangingbegin
\hangingheaditem{126}{\csromanfolio{17}โอกญฺชหํ ตณฺหจฺฉิทํ อเนชํ, (อิจฺจายสฺมา ภทฺราวุโธ,)}
\hangingbody{นนฺทิญฺชหํ โอฆติณฺณํ วิมุตฺตํ.}
\hangingbody{กปฺปญฺชหํ อภิยาเจ สุเมธํ,}
\hangingbody{สุตฺวาน นาคสฺส อปนมิสฺสนฺติ อิโต.}
\csromanhangingend

\proseitem{127}{นานาชนา ชนปเทหิ สงฺคตา,}

\prose{ตว วีร วากฺยํ อภิกงฺขมานา.}

\prose{เตสํ ตุวํ สาธุ วิยากโรหิ,}

\csromanhangingbegin
\hangingheaditem{128}{อาทานตณฺหํ วินเยถ สพฺพํ, (ภทฺราวุธาติ ภควา,)}
\hangingbody{อุทฺธํ อโธ ติริยญฺจาปิ มชฺเฌ.}
\hangingbody{ยํ ยญฺหิ โลกสฺมิมุปาทิยนฺติ,}
\hangingbody{เตเนว มาโร อเนฺวติ ชนฺตุํ.}
\csromanhangingend

\setlength{\csromangathapagewidth}{0pt}
\setlength{\csromangathawakwidth}{0pt}
\csromangathameasuregroup{}{ตสฺมา ปชานํ น อุปาทิเยถ, \\ ภิกฺขุ สโต กิญฺจนํ สพฺพโลเก. \\ อาทานสตฺเต อิติ เปกฺขมาโน, \\ ปชํ อิมํ มจฺจุเธยฺเย วิสตฺตนฺติ.}
\csromangathameasurewak{ตสฺมา ปชานํ น อุปาทิเยถ, \\ ภิกฺขุ สโต กิญฺจนํ สพฺพโลเก. \\ อาทานสตฺเต อิติ เปกฺขมาโน, \\ ปชํ อิมํ มจฺจุเธยฺเย วิสตฺตนฺติ.}
\setlength{\csromangathaleftcol}{0pt}
\csromangathagroup{\csromangathastanza{\csromangathawak{\csromangathaitemline{129}{ตสฺมา ปชานํ น อุปาทิเยถ,} \\ \csromangathacontline{ภิกฺขุ สโต กิญฺจนํ สพฺพโลเก.} \\ \csromangathacontline{อาทานสตฺเต อิติ เปกฺขมาโน,} \\ \csromangathacontline{ปชํ อิมํ มจฺจุเธยฺเย วิสตฺตนฺติ.}}}}

\vspace{\tipitakaparskip}
\csromancenter{ภทฺราวุธมาณวปุจฺฉา ทฺวาทสมา.}
\csromanheader{2}{13. อุทย\-มาณว\-ปุจฺฉา}
\csromanhangingbegin
\hangingheaditem{130}{ฌายิํ วิรชมาสีนํ, (อิจฺจายสฺมา อุทโย)}
\hangingbody{กตกิจฺจํ อนาสวํ.}
\hangingbody{ปารคุํ สพฺพธมฺมานํ, อตฺถิ ปเญฺหน อาคมํ.}
\hangingbody{อญฺญาวิโมกฺขํ ปพฺรูหิ, อวิชฺชาย ปเภทนํ.}
\csromanhangingend

\csromanhangingbegin
\hangingheaditem{131}{ปหานํ กามจฺฉนฺทานํ, (อุทยาติ ภควา)}
\hangingbody{โทมนสฺสาน จูภยํ.}
\hangingbody{ถินสฺส จ ปนูทนํ, กุกฺกุจฺจานํ นิวารณํ.}
\csromanhangingend

\setlength{\csromangathapagewidth}{0pt}
\setlength{\csromangathawakwidth}{0pt}
\csromangathameasuregroup{อุเปกฺขาสติสํสุทฺธํ, \\ อญฺญาวิโมกฺขํ ปพฺรูมิ, \\ กิํสุ สํโยชโน โลโก, \\ กิสฺสสฺส วิปฺปหาเนน, \\ นนฺทิสํโยชโน โลโก, \\ ตณฺหาย วิปฺปหาเนน,}{\csromangathabat{อุเปกฺขาสติสํสุทฺธํ,}{ธมฺมตกฺกปุเรชวํ.} \\ \csromangathabat{อญฺญาวิโมกฺขํ ปพฺรูมิ,}{อวิชฺชาย ปเภทนํ.} \\ \csromangathabat{กิํสุ สํโยชโน โลโก,}{กิํสุ ตสฺส วิจารณํ.} \\ \csromangathabat{กิสฺสสฺส วิปฺปหาเนน,}{นิพฺพานํ อิติ วุจฺจติ.} \\ \csromangathabat{นนฺทิสํโยชโน โลโก,}{วิตกฺกสฺส วิจารณํ.} \\ \csromangathabat{ตณฺหาย วิปฺปหาเนน,}{นิพฺพานํ อิติ วุจฺจติ.}}
\csromangathasetleft{อุเปกฺขาสติสํสุทฺธํ, \\ อญฺญาวิโมกฺขํ ปพฺรูมิ, \\ กิํสุ สํโยชโน โลโก, \\ กิสฺสสฺส วิปฺปหาเนน, \\ นนฺทิสํโยชโน โลโก, \\ ตณฺหาย วิปฺปหาเนน,}
\csromangathagroup{\csromangathastanza{\csromanfolio{18}\csromangathaitembat{132}{อุเปกฺขา\-สติ\-สํสุทฺธํ,}{ธมฺมตกฺก\-ปุเรชวํ.} \\ \csromangathacontbat{อญฺญาวิโมกฺขํ ปพฺรูมิ,}{อวิชฺชาย ปเภทนํ.}}\csromangathastanza{\csromangathaitembat{133}{กิํสุ สํโยชโน โลโก,}{กิํสุ ตสฺส วิจารณํ.} \\ \csromangathacontbat{กิสฺสสฺส วิปฺปหาเนน,}{นิพฺพานํ อิติ วุจฺจติ.}}\csromangathastanza{\csromangathaitembat{134}{นนฺทิสํโยชโน โลโก,}{วิตกฺกสฺส วิจารณํ.} \\ \csromangathacontbat{ตณฺหาย วิปฺปหาเนน,}{นิพฺพานํ อิติ วุจฺจติ.}}}

\csromanhangingbegin
\hangingheaditem{135}{กถํ สตสฺส จรโต, วิญฺญาณํ อุปรุชฺฌติ.}
\hangingbody{ภควนฺตํ ปุฏฺฐุมาคมฺม, ตํ สุโณม วโจ ตว. (6)}
\csromanhangingend

\setlength{\csromangathapagewidth}{0pt}
\setlength{\csromangathawakwidth}{0pt}
\csromangathameasuregroup{อชฺฌตฺตญฺจ พหิทฺธา จ, \\ เอวํ สตสฺส จรโต,}{\csromangathabat{อชฺฌตฺตญฺจ พหิทฺธา จ,}{เวทนํ นาภินนฺทโต.} \\ \csromangathabat{เอวํ สตสฺส จรโต,}{วิญฺญาณํ อุปรุชฺฌตีติ.}}
\csromangathasetleft{อชฺฌตฺตญฺจ พหิทฺธา จ, \\ เอวํ สตสฺส จรโต,}
\csromangathagroup{\csromangathastanza{\csromangathaitembat{136}{อชฺฌตฺตญฺจ พหิทฺธา จ,}{เวทนํ นาภินนฺทโต.} \\ \csromangathacontbat{เอวํ สตสฺส จรโต,}{วิญฺญาณํ อุปรุชฺฌตีติ.}}}

\vspace{\tipitakaparskip}
\csromancenter{อุทย\-มาณว\-ปุจฺฉา เตรสมา.}
\csromanmatikaanchor{mtk.16}
\csromantocmarkref{subsection}{14. โปสาลมาณวปุจฺฉา}{mtk.16}
\bookmark[dest={mtk.16},level=2]{14. โปสาลมาณวปุจฺฉา}
\csromanheader{2}{14. \textbf{โปสาลมาณวปุจฺฉา}}
\csromanhangingbegin
\hangingheaditem{137}{โย อตีตํ อาทิสติ, (อิจฺจายสฺมา โปสาโล,)}
\hangingbody{อเนโช ฉินฺนสํสโย.}
\hangingbody{ปารคุํ สพฺพธมฺมานํ, อตฺถิ ปเญฺหน อาคมํ.}
\csromanhangingend

\setlength{\csromangathapagewidth}{0pt}
\setlength{\csromangathawakwidth}{0pt}
\csromangathameasuregroup{วิภูตรูปสญฺญิสฺส, \\ อชฺฌตฺตญฺจ พหิทฺธา จ, \\ ญาณํ สกฺกา’นุปุจฺฉามิ,}{\csromangathabat{วิภูตรูปสญฺญิสฺส,}{สพฺพกายปฺปหายิโน.} \\ \csromangathabat{อชฺฌตฺตญฺจ พหิทฺธา จ,}{นตฺถิ กิญฺจีติ ปสฺสโต.} \\ \csromangathabat{ญาณํ สกฺกา’นุปุจฺฉามิ,}{กถํ เนยฺโย ตถาวิโธ.}}
\csromangathasetleft{วิภูตรูปสญฺญิสฺส, \\ อชฺฌตฺตญฺจ พหิทฺธา จ, \\ ญาณํ สกฺกา’นุปุจฺฉามิ,}
\csromangathagroup{\csromangathastanza{\csromangathaitembat{138}{วิภูต\-รูป\-สญฺญิสฺส,}{สพฺพกาย\-ปฺปหายิโน.} \\ \csromangathacontbat{อชฺฌตฺตญฺจ พหิทฺธา จ,}{นตฺถิ กิญฺจีติ ปสฺสโต.} \\ \csromangathacontbat{ญาณํ สกฺกา’นุปุจฺฉามิ,}{กถํ เนยฺโย ตถาวิโธ.}}}

\csromanhangingbegin
\hangingheaditem{139}{วิญฺญาณ\-ฏฺฐิติโย สพฺพา, (โปสาลาติ ภควา,)}
\hangingbody{อภิชานํ ตถาคโต.}
\hangingbody{ติฏฺฐนฺตเมนํ ชานาติ, วิมุตฺตํ ตปฺปรายณํ.}
\csromanhangingend

\setlength{\csromangathapagewidth}{0pt}
\setlength{\csromangathawakwidth}{0pt}
\csromangathameasuregroup{อากิญฺจญฺญสมฺภวํ ญ ตฺวา, \\ เอวเมตํ อภิญฺญาย, \\ เอตํ\textsuperscript{1} ญาณํ ตถํ ตสฺส,}{\csromangathabat{อากิญฺจญฺญสมฺภวํ ญ ตฺวา,}{นนฺที สํโยชนํ อิติ.} \\ \csromangathabat{เอวเมตํ อภิญฺญาย,}{ตโต ตตฺถ วิปสฺสติ.} \\ \csromangathabat{เอตํ\textsuperscript{1} ญาณํ ตถํ ตสฺส,}{พฺราหฺมณสฺส วุสีมโตติ.}}
\csromangathasetleft{อากิญฺจญฺญสมฺภวํ ญ ตฺวา, \\ เอวเมตํ อภิญฺญาย, \\ เอตํ\textsuperscript{1} ญาณํ ตถํ ตสฺส,}
\csromangathagroup{\csromangathastanza{\csromangathaitembat{140}{อากิญฺจญฺญ\-สมฺภวํ ญ ตฺวา,}{นนฺที สํโยชนํ อิติ.} \\ \csromangathacontbat{เอวเมตํ อภิญฺญาย,}{ตโต ตตฺถ วิปสฺสติ.} \\ \csromangathacontbat{เอตํ\csromansharedfootnote{923f2d999811275a}{เอวํ (สฺยา, ก)} ญาณํ ตถํ ตสฺส,}{พฺราหฺมณสฺส วุสีมโตติ.}}}

\vspace{\tipitakaparskip}
\csromancenter{โปสาลมาณวปุจฺฉา จุทฺทสมา.}
\csromanmatikaanchor{mtk.17}
\csromanmatikaanchor{mtk.18}
\csromantocmarkref{subsection}{15. โมฆราชมาณวปุจฺฉา}{mtk.17}
\bookmark[dest={mtk.17},level=2]{15. โมฆราชมาณวปุจฺฉา}
\csromantocmarkref{subsection}{16. ปิงฺคิยมาณวปุจฺฉา}{mtk.18}
\bookmark[dest={mtk.18},level=2]{16. ปิงฺคิยมาณวปุจฺฉา}
\csromanmarkh{16. ปิงฺคิยมาณวปุจฺฉา}\csromanmarkhii{15. โมฆราชมาณวปุจฺฉา}\csromanheadernomark{2}{16. ปิงฺคิยมาณวปุจฺฉา \textbf{15. โมฆราชมาณวปุจฺฉา}}
\csromanhangingbegin
\hangingheaditem{141}{\csromanfolio{19}ทฺวาหํ สกฺกํ อปุจฺฉิสฺสํ, (อิจฺจายสฺมา โมฆราชา,)}
\hangingbody{น เม พฺยากาสิ จกฺขุมา.}
\hangingbody{ยาวตติยญฺจ เทวีสิ,}
\hangingbody{พฺยากโรตีติ เม สุตํ.}
\csromanhangingend

\setlength{\csromangathapagewidth}{0pt}
\setlength{\csromangathawakwidth}{0pt}
\csromangathameasuregroup{เอวํ อภิกฺกนฺตทสฺสาวิํ, \\ กถํ โลกํ อเวกฺขนฺตํ, \\ สุญฺญโต โลกํ อเวกฺขสฺสุ, \\ อตฺตานุทิฏฺฐิํ อูหจฺจ, \\ เอวํ โลกํ อเวกฺขนฺตํ,}{อยํ โลโก ปโร โลโก, \\ พฺรหฺมโลโก สเทวโก. \\ ทิฏฺฐิํ เต นาภิชานาติ, \\ โคตมสฺส ยสสฺสิโน. \\ \csromangathabat{เอวํ อภิกฺกนฺตทสฺสาวิํ,}{อตฺถิ ปเญฺหน อาคมํ.} \\ \csromangathabat{กถํ โลกํ อเวกฺขนฺตํ,}{มจฺจุราชา น ปสฺสติ.} \\ \csromangathabat{สุญฺญโต โลกํ อเวกฺขสฺสุ,}{โมฆราช สทา สโต.} \\ \csromangathabat{อตฺตานุทิฏฺฐิํ อูหจฺจ,}{เอวํ มจฺจุตโร สิยา.} \\ \csromangathabat{เอวํ โลกํ อเวกฺขนฺตํ,}{มจฺจุราชา น ปสฺสตีติ.}}
\csromangathameasurewak{อยํ โลโก ปโร โลโก, \\ พฺรหฺมโลโก สเทวโก. \\ ทิฏฺฐิํ เต นาภิชานาติ, \\ โคตมสฺส ยสสฺสิโน.}
\csromangathasetleft{เอวํ อภิกฺกนฺตทสฺสาวิํ, \\ กถํ โลกํ อเวกฺขนฺตํ, \\ สุญฺญโต โลกํ อเวกฺขสฺสุ, \\ อตฺตานุทิฏฺฐิํ อูหจฺจ, \\ เอวํ โลกํ อเวกฺขนฺตํ,}
\csromangathagroup{\csromangathastanza{\csromangathawak{\csromangathaitemline{142}{อยํ โลโก ปโร โลโก,} \\ \csromangathacontline{พฺรหฺมโลโก สเทวโก.} \\ \csromangathacontline{ทิฏฺฐิํ เต นาภิชานาติ,} \\ \csromangathacontline{โคตมสฺส ยสสฺสิโน.}}}\csromangathastanza{\csromangathaitembat{143}{เอวํ อภิกฺกนฺต\-ทสฺสาวิํ,}{อตฺถิ ปเญฺหน อาคมํ.} \\ \csromangathacontbat{กถํ โลกํ อเวกฺขนฺตํ,}{มจฺจุราชา น ปสฺสติ.}}\csromangathastanza{\csromangathaitembat{144}{สุญฺญโต โลกํ อเวกฺขสฺสุ,}{โมฆราช สทา สโต.} \\ \csromangathacontbat{อตฺตานุทิฏฺฐิํ อูหจฺจ,}{เอวํ มจฺจุตโร สิยา.} \\ \csromangathacontbat{เอวํ โลกํ อเวกฺขนฺตํ,}{มจฺจุราชา น ปสฺสตีติ.}}}

\vspace{\tipitakaparskip}
\csromancenter{โมฆราชมาณวปุจฺฉา ปนฺนรสมา.}
\csromanheader{2}{16. ปิงฺคิยมาณวปุจฺฉา}
\csromanhangingbegin
\hangingheaditem{145}{ชิณฺโณหมสฺมิ อพโล วีตวณฺโณ, (อิจฺจายสฺมา ปิงฺคิโย,)}
\hangingbody{เนตฺตา น สุทฺธา สวนํ น ผาสุ.}
\hangingbody{มา’หํ นสฺสํ โมมุโห อนฺตราว,}
\hangingbody{อาจิกฺข ธมฺมํ ยมหํ วิชญฺญํ.}
\hangingbody{ชาติชราย อิธ วิปฺปหานํ.}
\csromanhangingend

\csromanhangingbegin
\hangingheaditem{146}{ทิสฺวาน รูเปสุ วิหญฺญมาเน, (ปิงฺคิยาติ ภควา,)}
\hangingbody{รุปฺปนฺติ รูเปสุ ชนา ปมตฺตา.}
\hangingbody{ตสฺมา ตุวํ ปิงฺคิย อปฺปมตฺโต,}
\hangingbody{ชหสฺสุ รูปํ อปุนพฺภวาย.}
\csromanhangingend

\setlength{\csromangathapagewidth}{0pt}
\setlength{\csromangathawakwidth}{0pt}
\csromangathameasuregroup{}{ทิสา จตสฺโส วิทิสา จตสฺโส, \\ อุทฺธํ อโธ ทส ทิสา อิมาโย. \\ น ตุยฺหํ อทิฏฺฐํ อสุตํ อมุตํ\textsuperscript{1}, \\ อโถ อวิญฺญาตํ กิญฺจน’มตฺถิ\textsuperscript{1} โลเก. \\ อาจิกฺข ธมฺมํ ยมหํ วิชญฺญํ, \\ ชาติชราย อิธ วิปฺปหานํ.}
\csromangathameasurewak{ทิสา จตสฺโส วิทิสา จตสฺโส, \\ อุทฺธํ อโธ ทส ทิสา อิมาโย. \\ น ตุยฺหํ อทิฏฺฐํ อสุตํ อมุตํ\textsuperscript{1}, \\ อโถ อวิญฺญาตํ กิญฺจน’มตฺถิ\textsuperscript{1} โลเก. \\ อาจิกฺข ธมฺมํ ยมหํ วิชญฺญํ, \\ ชาติชราย อิธ วิปฺปหานํ.}
\setlength{\csromangathaleftcol}{0pt}
\csromangathagroup{\csromangathastanza{\csromangathawak{\csromanfolio{20}\csromangathaitemline{147}{ทิสา จตสฺโส วิทิสา จตสฺโส,} \\ \csromangathacontline{อุทฺธํ อโธ ทส ทิสา อิมาโย.} \\ \csromangathacontline{น ตุยฺหํ อทิฏฺฐํ อสุตํ อมุตํ\csromansharedfootnote{ef913ad043c9d4a1}{อสุตํ อมุตํ วา (สี), อสุตามุตํ วา (สฺยา), อสุตํ’มุตํ วา (อิ)},} \\ \csromangathacontline{อโถ อวิญฺญาตํ กิญฺจน’มตฺถิ\csromansharedfootnote{2ddec0f7fa08c0d1}{กญฺจิ มตฺถิ (สฺยา), กิญฺจิ นตฺถิ (อิ), กิญฺจินมตฺถิ (ก)} โลเก.} \\ \csromangathacontline{อาจิกฺข ธมฺมํ ยมหํ วิชญฺญํ,} \\ \csromangathacontline{ชาติชราย อิธ วิปฺปหานํ.}}}}

\csromanhangingbegin
\hangingheaditem{148}{ตณฺหาธิปนฺเน มนุเช เปกฺขมาโน, (ปิงฺคิยาติ ภควา,)}
\hangingbody{สนฺตาปชาเต ชรสา ปเรเต,}
\hangingbody{ตสฺมา ตุวํ ปิงฺคิย อปฺปมตฺโต.}
\hangingbody{ชหสฺสุ ตณฺหํ อปุนพฺภวายาติ.}
\csromanhangingend

\vspace{\tipitakaparskip}
\csromancenter{ปิงฺคิยมาณวปุจฺฉา โสฬสมา.}
\csromanmatikaanchor{mtk.19}
\csromantocmarkref{subsubsection}{ปารายนตฺถุติคาถา}{mtk.19}
\bookmark[dest={mtk.19},level=3]{ปารายนตฺถุติคาถา}
\csromanheader{3}{ปารายนานุคีติคาถา \textbf{ปารายนตฺถุติคาถา}}
\prose{\csromanfolio{21}อิทมโวจ ภควา มคเธสุ วิหรนฺโต ปาสาณเก เจติเย, ปริจารก\-โสฬสานํ\csromansharedfootnote{d94b07cfe7229ebf}{ปริจารกโสฬสนฺนํ (สฺยา, ก)} พฺราหฺมณานํ อชฺฌิฏฺโฐ ปุฏฺโฐ ปุฏฺโฐ ปญฺหํ\csromansharedfootnote{5a3f20abd35a94a6}{ปเญฺห (สี, อิ)} พฺยากาสิ. เอกเมกสฺส เจปิ ปญฺหสฺส อตฺถมญฺญาย ธมฺมมญฺญาย ธมฺมา\-นุธมฺมํ ปฏิปชฺเชยฺย, คจฺเฉยฺเยว ชรามรณสฺส ปารํ, ปารงฺคมนียา อิเม ธมฺมาติ, ตสฺมา อิมสฺส ธมฺม\-ปริยายสฺส ปารายนนฺเตว\csromansharedfootnote{125df8d59286f841}{ปารายณํเตฺวว (สีฏฺฐ)} อธิวจนํ.}

\setlength{\csromangathapagewidth}{0pt}
\setlength{\csromangathawakwidth}{0pt}
\csromangathameasuregroup{อชิโต ติสฺสเมตฺเตยฺโย, \\ โธตโก อุปสีโว จ, \\ โตเทยฺยกปฺปา ทุภโย, \\ ภทฺราวุโธ อุทโย จ, \\ โมฆราชา จ เมธาวี, \\ เอเต พุทฺธํ อุปาคจฺฉุํ, \\ ปุจฺฉนฺตา นิปุเณ ปเญฺห, \\ เตสํ พุทฺโธ ปพฺยากาสิ, \\ ปญฺหานํ เวยฺยากรเณน, \\ เต โตสิตา จกฺขุมตา, \\ พฺรหฺมจริยมจริํสุ, \\ เอกเมกสฺส ปญฺหสฺส, \\ ตถา โย ปฏิปชฺเชยฺย, \\ อปารา ปารํ คจฺเฉยฺย, \\ มคฺโค โส ปารํ คมนาย,}{\csromangathabat{อชิโต ติสฺสเมตฺเตยฺโย,}{ปุณฺณโก อถ เมตฺตคู.} \\ \csromangathabat{โธตโก อุปสีโว จ,}{นนฺโท จ อถ เหมโก.} \\ \csromangathabat{โตเทยฺยกปฺปา ทุภโย,}{ชตุกณฺณี จ ปณฺฑิโต.} \\ \csromangathabat{ภทฺราวุโธ อุทโย จ,}{โปสาโล จาปิ พฺราหฺมโณ.} \\ \csromangathabat{โมฆราชา จ เมธาวี,}{ปิงฺคิโย จ มหาอิสิ.} \\ \csromangathabat{เอเต พุทฺธํ อุปาคจฺฉุํ,}{สมฺปนฺนจรณํ อิสิํ.} \\ \csromangathabat{ปุจฺฉนฺตา นิปุเณ ปเญฺห,}{พุทฺธเสฏฺฐํ อุปาคมุํ.} \\ \csromangathabat{เตสํ พุทฺโธ ปพฺยากาสิ,}{ปเญฺห ปุฏฺโฐ ยถาตถํ.} \\ \csromangathabat{ปญฺหานํ เวยฺยากรเณน,}{โตเสสิ พฺราหฺมเณ มุนิ.} \\ \csromangathabat{เต โตสิตา จกฺขุมตา,}{พุทฺเธนาทิจฺจพนฺธุนา.} \\ \csromangathabat{พฺรหฺมจริยมจริํสุ,}{วรปญฺญสฺส สนฺติเก.} \\ \csromangathabat{เอกเมกสฺส ปญฺหสฺส,}{ยถา พุทฺเธน เทสิตํ.} \\ \csromangathabat{ตถา โย ปฏิปชฺเชยฺย,}{คจฺเฉ ปารํ อปารโต.} \\ \csromangathabat{อปารา ปารํ คจฺเฉยฺย,}{ภาเวนฺโต มคฺคมุตฺตมํ.} \\ \csromangathabat{มคฺโค โส ปารํ คมนาย,}{ตสฺมา ปารายนํ อิติ.}}
\csromangathasetleft{อชิโต ติสฺสเมตฺเตยฺโย, \\ โธตโก อุปสีโว จ, \\ โตเทยฺยกปฺปา ทุภโย, \\ ภทฺราวุโธ อุทโย จ, \\ โมฆราชา จ เมธาวี, \\ เอเต พุทฺธํ อุปาคจฺฉุํ, \\ ปุจฺฉนฺตา นิปุเณ ปเญฺห, \\ เตสํ พุทฺโธ ปพฺยากาสิ, \\ ปญฺหานํ เวยฺยากรเณน, \\ เต โตสิตา จกฺขุมตา, \\ พฺรหฺมจริยมจริํสุ, \\ เอกเมกสฺส ปญฺหสฺส, \\ ตถา โย ปฏิปชฺเชยฺย, \\ อปารา ปารํ คจฺเฉยฺย, \\ มคฺโค โส ปารํ คมนาย,}
\csromangathagroup{\csromangathastanza{\csromangathaitembat{149}{อชิโต ติสฺสเมตฺเตยฺโย,}{ปุณฺณโก อถ เมตฺตคู.} \\ \csromangathacontbat{โธตโก อุปสีโว จ,}{นนฺโท จ อถ เหมโก.}}\csromangathastanza{\csromangathaitembat{150}{โตเทยฺยกปฺปา ทุภโย,}{ชตุกณฺณี จ ปณฺฑิโต.} \\ \csromangathacontbat{ภทฺราวุโธ อุทโย จ,}{โปสาโล จาปิ พฺราหฺมโณ.} \\ \csromangathacontbat{โมฆราชา จ เมธาวี,}{ปิงฺคิโย จ มหาอิสิ.}}\csromangathastanza{\csromangathaitembat{151}{เอเต พุทฺธํ อุปาคจฺฉุํ,}{สมฺปนฺน\-จรณํ อิสิํ.} \\ \csromangathacontbat{ปุจฺฉนฺตา นิปุเณ ปเญฺห,}{พุทฺธเสฏฺฐํ อุปาคมุํ.}}\csromangathastanza{\csromangathaitembat{152}{เตสํ พุทฺโธ ปพฺยากาสิ,}{ปเญฺห ปุฏฺโฐ ยถาตถํ.} \\ \csromangathacontbat{ปญฺหานํ เวยฺยากรเณน,}{โตเสสิ พฺราหฺมเณ มุนิ.}}\csromangathastanza{\csromangathaitembat{153}{เต โตสิตา จกฺขุมตา,}{พุทฺเธนา\-ทิจฺจพนฺธุนา.} \\ \csromangathacontbat{พฺรหฺมจริยม\-จริํสุ,}{วรปญฺญสฺส สนฺติเก.}}\csromangathastanza{\csromangathaitembat{154}{เอกเมกสฺส ปญฺหสฺส,}{ยถา พุทฺเธน เทสิตํ.} \\ \csromangathacontbat{ตถา โย ปฏิปชฺเชยฺย,}{คจฺเฉ ปารํ อปารโต.}}\csromangathastanza{\csromangathaitembat{155}{อปารา ปารํ คจฺเฉยฺย,}{ภาเวนฺโต มคฺคมุตฺตมํ.} \\ \csromangathacontbat{มคฺโค โส ปารํ คมนาย,}{ตสฺมา ปารายนํ อิติ.}}}

\csromanmatikaanchor{mtk.20}
\csromantocmarkref{subsubsection}{ปารายนานุคีติคาถา}{mtk.20}
\bookmark[dest={mtk.20},level=3]{ปารายนานุคีติคาถา}
\csromanheader{3}{ปารายนานุคีติคาถา}
\csromanhangingbegin
\hangingheaditem{156}{ปารายนม\-นุคายิสฺสํ, (อิจฺจายสฺมา ปิงฺคิโย,)}
\hangingbody{ยถา’ทฺทกฺขิ ตถา’กฺขาสิ, วิมโล ภูริเมธโส.}
\hangingbody{นิกฺกาโม นิพฺพโน4 นาโค, กิสฺส เหตุ มุสา ภเณ.(1)}
\csromanhangingend

\setlength{\csromangathapagewidth}{0pt}
\setlength{\csromangathawakwidth}{0pt}
\csromangathameasuregroup{ปหีนมลโมหสฺส, \\ หนฺทาหํ กิตฺตยิสฺสามิ, \\ เยเม ปุพฺเพ วิยากํสุ, \\ “อิจฺจาสิ อิติ ภวิสฺสติ”, \\ เอโก ตมนุทาสิโน, \\ โคตโม ภูริปญฺญาโณ, \\ โย เม ธมฺมมเทเสสิ, \\ ตณฺหกฺขยมนีติกํ, \\ กิํ นุ ตมฺหา วิปฺปวสสิ, \\ โคตมา ภูริปญฺญาณา,}{\csromangathabat{ปหีนมลโมหสฺส,}{มานมกฺขปฺปหายิโน.} \\ \csromangathabat{หนฺทาหํ กิตฺตยิสฺสามิ,}{คิรํ วณฺณูปสญฺหิตํ.} \\ ตโมนุโท พุทฺโธ สมนฺตจกฺขุ, \\ โลกนฺตคู สพฺพภวาติวตฺโต. \\ อนาสโว สพฺพทุกฺขปฺปหีโน, \\ สจฺจวฺหโย พฺรหฺเม อุปาสิโต เม. \\ ทิโช ยถา กุพฺพนกํ ปหาย, \\ พหุปฺผลํ กานนมาวเสยฺย. \\ เอวมฺปหํ อปฺปทสฺเส ปหาย, \\ มโหทธิํ หํโสริว อชฺฌปตฺโต. \\ \csromangathabat{เยเม ปุพฺเพ วิยากํสุ,}{หุรํ โคตมสาสนา.} \\ \csromangathabat{“อิจฺจาสิ อิติ ภวิสฺสติ”,}{สพฺพํ ตํ อิติหีติหํ,} \\ สพฺพํ ตํ ตกฺกวฑฺฒนํ. \\ \csromangathabat{เอโก ตมนุทาสิโน,}{ชุติมา โส ปภงฺกโร.} \\ \csromangathabat{โคตโม ภูริปญฺญาโณ,}{โคตโม ภูริเมธโส.} \\ \csromangathabat{โย เม ธมฺมมเทเสสิ,}{สนฺทิฏฺฐิกมกาลิกํ.} \\ \csromangathabat{ตณฺหกฺขยมนีติกํ,}{ยสฺส นตฺถิ อุปมา กฺวจิ.} \\ \csromangathabat{กิํ นุ ตมฺหา วิปฺปวสสิ,}{มุหุตฺตมปิ ปิงฺคิย.} \\ \csromangathabat{โคตมา ภูริปญฺญาณา,}{โคตมา ภูริเมธสา.}}
\csromangathameasurewak{ตโมนุโท พุทฺโธ สมนฺตจกฺขุ, \\ โลกนฺตคู สพฺพภวาติวตฺโต. \\ อนาสโว สพฺพทุกฺขปฺปหีโน, \\ สจฺจวฺหโย พฺรหฺเม อุปาสิโต เม. \\ ทิโช ยถา กุพฺพนกํ ปหาย, \\ พหุปฺผลํ กานนมาวเสยฺย. \\ เอวมฺปหํ อปฺปทสฺเส ปหาย, \\ มโหทธิํ หํโสริว อชฺฌปตฺโต.}
\csromangathasetleft{ปหีนมลโมหสฺส, \\ หนฺทาหํ กิตฺตยิสฺสามิ, \\ เยเม ปุพฺเพ วิยากํสุ, \\ “อิจฺจาสิ อิติ ภวิสฺสติ”, \\ เอโก ตมนุทาสิโน, \\ โคตโม ภูริปญฺญาโณ, \\ โย เม ธมฺมมเทเสสิ, \\ ตณฺหกฺขยมนีติกํ, \\ กิํ นุ ตมฺหา วิปฺปวสสิ, \\ โคตมา ภูริปญฺญาณา,}
\csromangathagroup{\csromangathastanza{\csromangathaitembat{157}{ปหีน\-มล\-โมหสฺส,}{มาน\-มกฺข\-ปฺปหายิโน.} \\ \csromangathacontbat{หนฺทาหํ กิตฺตยิสฺสามิ,}{คิรํ วณฺณู\-ปสญฺหิตํ.}}\csromangathastanza{\csromangathawak{\csromanfolio{22}\csromangathaitemline{158}{ตโมนุโท พุทฺโธ สมนฺตจกฺขุ,} \\ \csromangathacontline{โลกนฺตคู สพฺพภวาติวตฺโต.} \\ \csromangathacontline{อนาสโว สพฺพทุกฺข\-ปฺปหีโน,} \\ \csromangathacontline{สจฺจวฺหโย พฺรหฺเม อุปาสิโต เม.}}}\csromangathastanza{\csromangathawak{\csromangathaitemline{159}{ทิโช ยถา กุพฺพนกํ ปหาย,} \\ \csromangathacontline{พหุปฺผลํ กานนมา\-วเสยฺย.} \\ \csromangathacontline{เอวมฺปหํ อปฺปทสฺเส ปหาย,} \\ \csromangathacontline{มโหทธิํ หํโสริว อชฺฌปตฺโต.}}}\csromangathastanza{\csromangathaitembat{160}{เยเม ปุพฺเพ วิยากํสุ,}{หุรํ โคตมสาสนา.} \\ \csromangathacontbat{“อิจฺจาสิ อิติ ภวิสฺสติ”,}{สพฺพํ ตํ อิติหีติหํ,}}\csromangathastanza{\csromangathaitemline{160}{สพฺพํ ตํ ตกฺกวฑฺฒนํ.}}\csromangathastanza{\csromangathaitembat{161}{เอโก ตมนุทาสิโน,}{ชุติมา โส ปภงฺกโร.} \\ \csromangathacontbat{โคตโม ภูริปญฺญาโณ,}{โคตโม ภูริเมธโส.}}\csromangathastanza{\csromangathaitembat{162}{โย เม ธมฺมมเทเสสิ,}{สนฺทิฏฺฐิกม\-กาลิกํ.} \\ \csromangathacontbat{ตณฺหกฺขยม\-นีติกํ,}{ยสฺส นตฺถิ อุปมา กฺวจิ.}}\csromangathastanza{\csromangathaitembat{163}{กิํ นุ ตมฺหา วิปฺปวสสิ,}{มุหุตฺตมปิ ปิงฺคิย.} \\ \csromangathacontbat{โคตมา ภูริปญฺญาณา,}{โคตมา ภูริเมธสา.}}}

\proseitem{164}{โย เต ธมฺมมเทเสสิ, สนฺทิฏฺฐิกม\-กาลิกํ.}

\prose{ตณฺหกฺขยม\-นีติกํ, ยสฺส นตฺถิ อุปมา กฺวจิ. (9)}

\proseitem{165}{นาหํ ตมฺหา วิปฺปวสามิ, มุหุตฺตมปิ พฺราหฺมณ.}

\prose{โคตมา ภูริปญฺญาณา, โคตมา ภูริเมธสา. (10)}

\proseitem{166}{โย เม ธมฺมมเทเสสิ, สนฺทิฏฺฐิกม\-กาลิกํ.}

\prose{ตณฺหกฺขยม\-นีติกํ, ยสฺส นตฺถิ อุปมา กฺวจิ. (11)}

\proseitem{167}{ปสฺสามิ นํ มนสา จกฺขุนาว,}

\prose{รตฺตินฺทิวํ พฺราหฺมณ อปฺปมตฺโต.}

\prose{นมสฺสมาโน วิวเสมิ รตฺติํ,}

\prose{เตเนว มญฺญามิ อวิปฺปวาสํ. (12)}

\csromanheader{2}{ปารายนานุคีติคาถา}
\proseitem{168}{\csromanfolio{23}สทฺธา จ ปีติ จ มโน สติ จ,}

\prose{นา’เปนฺติ’เม โคตม\-สาสนมฺหา.}

\prose{ยํ ยํ ทิสํ วชติ ภูริปญฺโญ,}

\prose{ส เตน เตเนว นโตหมสฺมิ. (13)}

\proseitem{169}{ชิณฺณสฺส เม ทุพฺพล\-ถามกสฺส,}

\prose{เตเนว กาโย น ปเลติ ตตฺถ.}

\prose{สํกปฺปย\-นฺตาย\csromansharedfootnote{dee1f105cc886b24}{สํกปฺปยตฺตาย (สี)} วชามิ นิจฺจํ,}

\prose{มโน หิ เม พฺราหฺมณ เตน ยุตฺโต. (14)}

\proseitem{170}{ปงฺเก สยาโน ปริผนฺทมาโน,}

\prose{ทีปา ทีปํ อุปลฺลวิํ.}

\prose{อถทฺทสาสิํ สมฺพุทฺธํ, โอฆติณฺณม\-นาสวํ. (15)}

\proseitem{171}{ยถา อหู วกฺกลิ มุตฺตสทฺโธ,}

\prose{ภทฺราวุโธ อาฬวิ โคตโม จ.}

\prose{เอวเมว ตฺวมฺปิ ปมุญฺจสฺสุ สทฺธํ,}

\prose{คมิสฺสสิ ตฺวํ ปิงฺคิย มจฺจุเธยฺยสฺส ปารํ\csromansharedfootnote{d42d393445cf7e16}{มจฺจุเธยฺยปารํ (สี)}. (16)}

\proseitem{172}{เอส ภิยฺโย ปสีทามิ, สุตฺวาน มุนิโน วโจ.}

\prose{วิวฏฺฏจฺฉโท สมฺพุทฺโธ, อขิโล ปฏิภานวา. (17)}

\proseitem{173}{อธิเทเว อภิญฺญาย, สพฺพํ เวทิ ปโรปรํ.}

\prose{ปญฺหานนฺตกโร สตฺถา, กงฺขีนํ ปฏิชานตํ. (18)}

\proseitem{174}{อสํหีรํ อสํกุปฺปํ,}

\prose{ยสฺส นตฺถิ อุปมา กฺวจิ.}

\prose{อทฺธา คมิสฺสามิ น เมตฺถ กงฺขา,}

\prosewithcloser{เอวํ มํ ธาเรหิ อธิมุตฺต\-จิตฺตนฺติ\csromansharedfootnote{be2466eff0246208}{อชิตมาณวปุจฺฉาย ปฏฺฐาย ยาวปารายนานุคีติคาถาปริโยสานา สฺยา...โปตฺถเก นตฺถิ.}. (19)}{\textbf{ปารายนานุคีติคาถา นิฏฺฐิตา.}}
\chapterheadpageread{\textbf{ปารายนวคฺค}}
\csromanmatikaanchor{mtk.21}
\csromanmatikaanchor{mtk.22}
\csromantocmarknopage{section}{ปารายนวคฺค}{mtk.21}
\bookmark[dest={mtk.21},level=1]{ปารายนวคฺค}
\csromantocmarkref{subsection}{1. อชิตมาณวปุจฺฉานิทฺเทส}{mtk.22}
\bookmark[dest={mtk.22},level=2]{1. อชิตมาณวปุจฺฉานิทฺเทส}
\proseitem{1}{\csromanfolio{24}\textbf{เกนสฺสุ นิวุโต โลโก, (อิจฺจายสฺมา อชิโต,)}}

\prose{\textbf{เกนสฺสุ นปฺปกาสติ.}}

\setlength{\csromangathapagewidth}{0pt}
\setlength{\csromangathawakwidth}{0pt}
\csromangathameasuregroup{}{\textbf{กิสฺสาภิเลปนํ พฺรูสิ}\textsuperscript{1}\textbf{,} \\ \textbf{กิํสุ ตสฺส มหพฺภยํ.}}
\csromangathameasurewak{\textbf{กิสฺสาภิเลปนํ พฺรูสิ}\textsuperscript{1}\textbf{,} \\ \textbf{กิํสุ ตสฺส มหพฺภยํ.}}
\setlength{\csromangathaleftcol}{0pt}
\csromangathagroup{\csromangathastanza{\csromangathawak{\textbf{กิสฺสาภิเลปนํ พฺรูสิ}\csromansharedfootnote{085d9888983879cb}{พฺรูหิ (สฺยา)}\textbf{,} \\ \textbf{กิํสุ ตสฺส มหพฺภยํ.}}}}

\prose{\textbf{เกนสฺสุ นิวุโต โลโก}ติ โลโกติ นิรย\textbf{โลโก} ติรจฺฉานโลโก เปตฺติวิสย\-โลโก มนุสฺสโลโก เทวโลโก ขนฺธโลโก ธาตุโลโก อายตนโลโก อยํ โลโก ปโร โลโก พฺรหฺมโลโก เทวโลโก. อยํ วุจฺจติ โลโก. อยํ โลโก เกน อาวุโต นิวุโต โอวุโต\csromansharedfootnote{b482540df50e0521}{โอผุโต (สฺยา)} ปิหิโต ปฏิจฺฉนฺโน ปฏิกุชฺชิโตติ เกนสฺสุ นิวุโต โลโก.}

\prose{\textbf{อิจฺจายสฺมา อชิโต}ติ \textbf{อิจฺจา}ติ ปทสนฺธิ ปทสํสคฺโค ปทปาริปูรี อกฺขร\-สมวาโย พฺยญฺชน\-สิลิฏฺฐตา ปทา\-นุปุพฺพตา\-เปตํ\csromansharedfootnote{fde928508f82525a}{ปทานุปุพฺพตาเมตํ (พหูสุ)} อิจฺจาติ. \textbf{อายสฺมา}ติ ปิยวจนํ ครุวจนํ สคารวสปฺปติสฺสาธิวจนเมตํ อายสฺมาติ. \textbf{อชิโต}ติ ตสฺส พฺราหฺมณสฺส นามํ สงฺขา สมญฺญา ปญฺญตฺติ โวหาโร นามํ นามกมฺมํ นามเธยฺยํ นิรุตฺติ พฺยญฺชนํ อภิลาโปติ อิจฺจายสฺมา อชิโต.}

\prose{\textbf{เกนสฺสุ นปฺปกาสตี}ติ เกน โลโก นปฺปกาสติ น ภาสติ น ตปติ น วิโรจติ น ญายติ น ปญฺญายตีติ เกนสฺสุ นปฺปกาสติ.}

\prose{\textbf{กิสฺสาภิเลปนํ พฺรูสี}ติ กิํ โลกสฺส เลปนํ ลคฺคนํ พนฺธนํ อุปกฺกิเลโส. เกน โลโก ลิตฺโต สํลิตฺโต อุปลิตฺโต กิลิฏฺโฐ สํกิลิฏฺโฐ มกฺขิโต สํสฏฺโฐ ลคฺโค ลคฺคิโต ปลิพุทฺโธ พฺรูสิ อาจิกฺขสิ เทเสสิ ปญฺญเปสิ\csromansharedfootnote{a48e870bac758108}{ปญฺญาเปสิ (ก)} ปฏฺฐเปสิ วิวรสิ วิภชสิ อุตฺตานีกโรสิ\csromansharedfootnote{c3b7befca29cd97f}{อุตฺตานิํ กโรสิ (ก)} ปกาเสสีติ กิสฺสาภิเลปนํ พฺรูสิ.}

\prose{\csromanfolio{25}\textbf{กิํสุ ตสฺส มหพฺภยนฺ}ติ กิํ โลกสฺส ภยํ มหพฺภยํ ปีฬนํ ฆฏฺฏนํ อุปทฺทโว อุปสคฺโคติ กิํสุ ตสฺส มหพฺภยํ. เตนาห โส พฺราหฺมโณ\csromandash{}}

\prose{เกนสฺสุ นิวุโต โลโก, (อิจฺจายสฺมา อชิโต.)}

\prose{เกนสฺสุ นปฺปกาสติ.}

\setlength{\csromangathapagewidth}{0pt}
\setlength{\csromangathawakwidth}{0pt}
\csromangathameasuregroup{}{กิสฺสาภิเลปนํ พฺรูสิ, \\ กิํสุ ตสฺส มหพฺภยนฺติ.}
\csromangathameasurewak{กิสฺสาภิเลปนํ พฺรูสิ, \\ กิํสุ ตสฺส มหพฺภยนฺติ.}
\setlength{\csromangathaleftcol}{0pt}
\csromangathagroup{\csromangathastanza{\csromangathawak{กิสฺสาภิเลปนํ พฺรูสิ, \\ กิํสุ ตสฺส มหพฺภยนฺติ.}}}

\proseitem{2}{\textbf{อวิชฺชาย นิวุโต โลโก, (อชิตาติ ภควา,)}}

\prose{\textbf{เววิจฺฉา ปมาทา นปฺปกาสติ.}}

\setlength{\csromangathapagewidth}{0pt}
\setlength{\csromangathawakwidth}{0pt}
\csromangathameasuregroup{}{\textbf{ชปฺปา}\textbf{ภิเลปนํ พฺรูมิ,} \\ \textbf{ทุกฺข’มสฺส มหพฺภยํ.}}
\csromangathameasurewak{\textbf{ชปฺปา}\textbf{ภิเลปนํ พฺรูมิ,} \\ \textbf{ทุกฺข’มสฺส มหพฺภยํ.}}
\setlength{\csromangathaleftcol}{0pt}
\csromangathagroup{\csromangathastanza{\csromangathawak{\textbf{ชปฺปา}\-\textbf{ภิเลปนํ พฺรูมิ,} \\ \textbf{ทุกฺข’มสฺส มหพฺภยํ.}}}}

\prose{\textbf{อวิชฺชาย นิวุโต โลโก}ติ \textbf{อวิชฺชา}ติ ทุกฺเข อญฺญาณํ ทุกฺขสมุทเย อญฺญาณํ ทุกฺขนิโรเธ อญฺญาณํ ทุกฺขนิโรธ\-คามินิยา ปฏิปทาย อญฺญาณํ ปุพฺพนฺเต อญฺญาณํ อปรนฺเต อญฺญาณํ ปุพฺพนฺตา\-ปรนฺเต อญฺญาณํ อิทปฺปจฺจยตา\-ปฏิจฺจ\-สมุปฺปนฺเนสุ ธมฺเมสุ อญฺญาณํ, ยํ เอวรูปํ อญฺญาณํ อทสฺสนํ อนภิสมโย อนนุโพโธ อสมฺโพโธ อปฺปฏิเวโธ อสํคาหนา อปริโยคาหนา อสมเปกฺขนา อปจฺจเวกฺขณา\csromansharedfootnote{dffd0efebee9f97c}{อปจฺจเวกฺขนา (สฺยา)} อปจฺจเวกฺขณ\-กมฺมํ ทุมฺเมชฺฌํ พาลฺยํ อสมฺปชญฺญํ โมโห ปโมโห สมฺโมโห อวิชฺชา อวิชฺโชโฆ อวิชฺชาโยโค อวิชฺชานุสโย อวิชฺชา\-ปริยุฏฺฐานํ อวิชฺชาลงฺคี โมโห อกุสลมูลํ. อยํ วุจฺจติ อวิชฺชา.}

\prose{\textbf{โลโก}ติ นิรยโลโก ติรจฺฉานโลโก เปตฺติวิสย\-โลโก มนุสฺสโลโก เทวโลโก ขนฺธโลโก ธาตุโลโก อายตนโลโก อยํ โลโก ปโร โลโก พฺรหฺมโลโก เทวโลโก. อยํ วุจฺจติ โลโก. อยํ โลโก อิมาย อวิชฺชาย อาวุโต นิวุโต โอวุโต ปิหิโต ปฏิจฺฉนฺโน ปฏิกุชฺชิโตติ อวิชฺชาย นิวุโต โลโก.}

\prose{\textbf{อชิตา}ติ ภควา ตํ พฺราหฺมณํ นาเมน อาลปติ. \textbf{ภควา}ติ คารวา\-ธิวจนํ. อปิ จ ภคฺคราโคติ ภควา, ภคฺคโทโสติ ภควา, \csromanfolio{26}ภคฺคโมโหติ ภควา, ภคฺคมาโนติ ภควา, ภคฺคทิฏฺฐีติ ภควา, ภคฺคกณฺฏโกติ ภควา, ภคฺคกิเลโสติ ภควา, ภชิ วิภชิ ปวิภชิ ธมฺม\-รตนนฺติ ภควา, ภวานํ อนฺตกโรติ ภควา, ภวิตกาโย, ภาวิตสีโล, ภาวิตจิตฺโต\csromansharedfootnote{f6c8cb5fde61abf5}{ภาวิตกาโยติ ภควา, ภาวิตสีโลติ ภาวิตจิตฺโตติ (สฺยา)}, ภาวิตปญฺโญติ ภควา, ภชิ วา ภควา อรญฺญ\-วนปตฺถานิ ปนฺตานิ เสนาสนานิ อปฺปสทฺทานิ อปฺปนิคฺโฆสานิ วิชนวาตานิ มนุสฺส\-ราหสฺเสยฺยกานิ\csromansharedfootnote{b8b3d5dc9a48ce74}{มนุสฺสราหเสยฺยกานิ (สฺยา)} ปฏิสลฺลานสารุปฺปานีติ ภควา, ภาคี วา ภควา จีวร\-ปิณฺฑปาต\-เสนาสน\-คิลานปจฺจย\-เภสชฺช\-ปริกฺขารานนฺติ ภควา, ภาคี วา ภควา อตฺถรสสฺส ธมฺมรสสฺส วิมุตฺติรสสฺส อธิสีลสฺส อธิจิตฺตสฺส อธิปญฺญายาติ ภควา, ภาคี วา ภควา จตุนฺนํ ฌานานํ จตุนฺนํ อปฺปมญฺญานํ จตุนฺนํ อรูป\-สมาปตฺตีนนฺติ ภควา, ภาคี วา ภควา อฏฺฐนฺนํ วิโมกฺขานํ อฏฺฐนฺนํ อภิภา\-ยตนานํ นวนฺนํ อนุปุพฺพ\-สมาปตฺตีนนฺติ ภควา, ภาคี วา ภควา ทสนฺนํ สญฺญา\-ภาวนานํ กสิณ\-สมาปตฺตีนํ อานาปาน\-สฺสติ\-สมาธิสฺส อสุภ\-สมาปตฺติยาติ ภควา, ภาคี วา ภควา จตุนฺนํ สติปฏฺฐานานํ จตุนฺนํ สมฺม\-ปฺปธานานํ จตุนฺนํ อิทฺธิปาทานํ ปญฺจนฺนํ อินฺทฺริยานํ ปญฺจนฺนํ พลานํ สตฺตนฺนํ โพชฺฌงฺคานํ อริยสฺส อฏฺฐงฺคิกสฺส มคฺคสฺสาติ ภควา, ภาคี วา ภควา ทสนฺนํ ตถาคต\-พลานํ จตุนฺนํ เวสารชฺชานํ จตุนฺนํ ปฏิสมฺภิทานํ ฉนฺนํ อภิญฺญานํ ฉนฺนํ พุทฺธธมฺมานนฺติ ภควา. ภควาติ เนตํ นามํ มาตรา กตํ, น ปิตรา กตํ, น ภาตรา กตํ, น ภคินิยา กตํ, น มิตฺตามจฺเจหิ กตํ, น ญาติสาโลหิเตหิ กตํ, น สมณ\-พฺราหฺมเณหิ กตํ, น เทวตาหิ กตํ, วิโมกฺข\-นฺติกเมตํ พุทฺธานํ ภควนฺตานํ โพธิยา มูเล สห สพฺพญฺญุต\-ญาณสฺส ปฏิลาภา สจฺฉิกา ปญฺญตฺติ, ยทิทํ ภควาติ อชิตาติ ภควา.}

\prose{\textbf{เววิจฺฉา ปมาทา นปฺปกาสตี}ติ เววิจฺฉํ วุจฺจติ ปญฺจ มจฺฉริยานิ อาวาส\-มจฺฉริยํ กุลมจฺฉริยํ ลาภ\-มจฺฉริยํ, วณฺณ\-มจฺฉริยํ ธมฺม\-มจฺฉริยํ, ยํ เอวรูปํ มจฺเฉรํ มจฺฉรายนา มจฺฉรายิ\-ตตฺตํ เววิจฺฉํ กทริยํ กฏุกญฺจุกตา อคฺคหิตตฺตํ จิตฺตสฺส, อิทํ วุจฺจติ มจฺฉริยํ. อปิ จ ขนฺธมจฺฉริยมฺปิ มจฺฉริยํ, ธาตุมจฺฉริยมฺปิ มจฺฉริยํ, อายตนมจฺฉริยมฺปิ มจฺฉริยํ, คาโห วุจฺจติ มจฺฉริยํ. ปมาโท วตฺตพฺโพ, กายทุจฺจริเต วา วจีทุจฺจริเต วา \csromanfolio{27}\textbf{มโนทุจฺจริเต วา ปญฺจสุ กามคุเณสุ วา จิตฺตสฺส โวสคฺโค}\csromansharedfootnote{0dd9c214b732f03e}{โวสฺสคฺโค (พหูสุ)} โวสคฺคา\-นุปฺปทานํ กุสลานํ ธมฺมานํ ภาวนาย อสกฺก\-จฺจ\-กิริ\-ยตา อสาตจฺจกิริยตา อนฏฺฐิต\-กิริยตา\csromansharedfootnote{32d67b1cc99055d5}{อนิฏฺฐิตกิริยตา (ก) อภิ 2. 363 ปิฏฺเฐ.} โอลีนวุตฺติตา นิกฺขิตฺต\-จฺฉนฺทตา นิกฺขิตฺตธุรตา อนาเสวนา อภาวนา อพหุลีกมฺมํ อนธิฏฺฐานํ อนนุโยโค ปมาโท, โย เอวรูโป ปมาโท ปมชฺชนา ปมชฺชิตตฺตํ. อยํ วุจฺจติ ปมาโท. เววิจฺฉา ปมาทา นปฺปกาสตีติ อิมินา จ มจฺฉริเยน อิมินา จ ปมาเทน โลโก นปฺปกาสติ น ภาสติ น ตปติ น วิโรจติ น ญายติ น ปญฺญายตีติ เววิจฺฉา ปมาทา นปฺปกาสติ.}

\prose{\textbf{ชปฺปา}\-\textbf{ภิเลปนํ พฺรูมี}ติ ชปฺปา วุจฺจติ ตณฺหา, โย ราโค สาราโค อนุนโย อนุโรโธ นนฺที\csromansharedfootnote{04b0b619220ae621}{นนฺทิ (สฺยา)} นนฺทิราโค จิตฺตสฺส สาราโค อิจฺฉา มุจฺฉา อชฺโฌสานํ เคโธ ปลิเคโธ สงฺโค ปงฺโก เอชา มายา ชนิกา สญฺชนนี สิพฺพินี ชาลินี สริตา วิสตฺติกา สุตฺตํ วิสฏา\csromansharedfootnote{bc9d24d35143b10c}{โสตฺตํ วิสตา (สฺยา)} อายูหนี ทุติยา ปณิธิ ภวเนตฺติ วนํ วนโถ สนฺถโว\csromansharedfootnote{54f27cc4a0a2d719}{สนฺธโว (ก) อภิ 2. 376 ปิฏฺเฐ.} สิเนโห อเปกฺขา ปฏิพนฺธุ อาสา อาสีสนา\csromansharedfootnote{82130b595aa05460}{อาสิํสนา (สฺยา)} อาสีสิตตฺตํ รูปาสา สทฺทาสา คนฺธาสา รสาสา โผฏฺฐพฺพาสา ลาภาสา ธนาสา ปุตฺตาสา ชีวิตาสา ชปฺปา ปชปฺปา อภิชปฺปา ชปฺปนา ชปฺปิตตฺตํ โลลุปฺปํ โลลุปฺปายนา โลลุปฺปา\-ยิตตฺตํ ปุจฺฉญฺชิกตา สาธุกมฺยตา อธมฺมราโค วิสมโลโภ นิกนฺติ นิกามนา ปตฺถนา ปิหนา สมฺปตฺถนา กามตณฺหา ภวตณฺหา วิภวตณฺหา รูปตณฺหา อรูปตณฺหา นิโรธตณฺหา รูปตณฺหา สทฺทตณฺหา คนฺธตณฺหา รสตณฺหา โผฏฺฐพฺพ\-ตณฺหา ธมฺมตณฺหา โอโฆ โยโค คนฺโถ อุปาทานํ อาวรณํ นีวรณํ ฉทนํ พนฺธนํ อุปกฺกิเลโส อนุสโย ปริยุฏฺฐานํ ลตา เววิจฺฉํ ทุกฺขมูลํ ทุกฺขนิทานํ ทุกฺขปฺปภโว มารปาโส มารพฬิสํ มารามิสํ มารวิสโย มารนิวาโส มารโคจโร มารพนฺธนํ ตณฺหานที ตณฺหาชาลํ ตณฺหาคทฺทุลํ ตณฺหาสมุทฺโท อภิชฺฌา โลโภ อกุสลมูลํ, อยํ วุจฺจติ ชปฺปา. โลกสฺส เลปนํ ลคฺคนํ พนฺธนํ อุปกฺกิเลโส อิมาย ชปฺปาย โลโก ลิตฺโต สํลิตฺโต อุปลิตฺโต กิลิฏฺโฐ สํกิลิฏฺโฐ มกฺขิโต สํสฏฺโฐ ลคฺโค ลคฺคิโต ปลิพุทฺโธติ พฺรูมิ อาจิกฺขามิ เทเสมิ ปญฺญเปมิ ปฏฺฐเปมิ วิวรามิ วิภชามิ อุตฺตานีกโรมิ ปกาเสมีติ ชปฺปา\-ภิเลปนํ พฺรูมิ.}

\proseitem{1}{\csromanfolio{28}\textbf{ทุกฺข’มสฺส มหพฺภยนฺ}ติ \textbf{ทุกฺขนฺ}ติ ชาติทุกฺขํ ชราทุกฺขํ พฺยาธิทุกฺขํ มรณทุกฺขํ โสกปริเทว\-ทุกฺข- โทมนสฺสุ\-ปายา\-สทุกฺขํ เนรยิกํ ทุกฺขํ ติรจฺฉาน\-โยนิกํ ทุกฺขํ เปตฺติวิสยิกํ ทุกฺขํ มานุสิกํ ทุกฺขํ คพฺโภ\-กฺกนฺติ\-มูลกํ ทุกฺขํ คพฺภ\-ฏฺฐิติ\-มูลกํ\csromansharedfootnote{badcb8d532c1eb7b}{คพฺเภฐิติมูลกํ (สฺยา, ก)} ทุกฺขํ คพฺภวุฏฺฐาน\-มูลกํ ทุกฺขํ ชาตสฺสู\-ปนิพนฺธกํ ทุกฺขํ ชาตสฺส ปราเธยฺยกํ ทุกฺขํ อตฺตู\-ปกฺกม\-ทุกฺขํ ปรูปกฺกม\-ทุกฺขํ สงฺขาร\-ทุกฺขํ วิปริณาม\-ทุกฺขํ, จกฺขุโรโค โสตโรโค ฆานโรโค ชิวฺหาโรโค กายโรโค สีสโรโค กณฺณโรโค มุขโรโค ทนฺตโรโค กาโส สาโส ปินาโส ฑาโห\csromansharedfootnote{9ef52c474c08f49b}{qอโห (สฺยา)} ชโร กุจฺฉิโรโค มุจฺฉา ปกฺขนฺทิกา สูลา วิสูจิกา กุฏฺฐํ คณฺโฑ กิลาโส โสโส อปมาโร ททฺทุ กณฺฑุ กจฺฉุ รขสา\csromansharedfootnote{88798fb5b338b5d3}{รกฺขสา (ก)} วิตจฺฉิกา โลหิตปิตฺตํ\csromansharedfootnote{89f663f1fb28caa0}{โลหิตํ ปิตฺตํ (พหูสุ)} มธุเมโห อํสา ปิฬกา ภคนฺทลา ปิตฺต\-สมุฏฺฐานา อาพาธา เสมฺห\-สมุฏฺฐานา อาพาธา วาตสมุฏฺฐานา อาพาธา สนฺนิปาติกา อาพาธา อุตุปริณามชา อาพาธา วิสม\-ปริหารชา อาพาธา โอปกฺกมิกา อาพาธา กมฺมวิปากชา อาพาธา สีตํ อุณฺหํ ชิฆจฺฉา ปิปาสา อุจฺจาโร ปสฺสาโว ฑํสมกส\-วาตา\-ตป\-สรีสป\-สมฺผสฺสํ ทุกฺขํ มาตุมรณํ ทุกฺขํ ปิตุมรณํ ทุกฺขํ ภาตุมรณํ ทุกฺขํ ภคินิมรณํ ทุกฺขํ ปุตฺตมรณํ ทุกฺขํ ธีตุมรณํ ทุกฺขํ ญาติพฺยสนํ ทุกฺขํ โรคพฺยสนํ ทุกฺขํ โภคพฺยสนํ ทุกฺขํ สีลพฺยสนํ ทุกฺขํ ทิฏฺฐิพฺยสนํ ทุกฺขํ เยสํ ธมฺมานํ อาทิโต สมุทาคมนํ ปญฺญายติ. อตฺถงฺคมโต นิโรโธ ปญฺญายติ. กมฺม\-สนฺนิสฺสิโต วิปาโก, วิปาก\-สนฺนิสฺสิตํ กมฺมํ, นาม\-สนฺนิสฺสิตํ รูปํ รูปสนฺนิสฺสิตํ นามํ, ชาติยา อนุคตํ ชราย อนุสฏํ พฺยาธินา อภิภูตํ มรเณน อพฺภาหตํ ทุกฺเข ปติฏฺฐิตํ อตาณํ อเลณํ อสรณํ อสรณีภูตํ. อิทํ วุจฺจติ ทุกฺขํ, อิทํ ทุกฺขํ โลกสฺส ภยํ มหาภยํ ปีฬนํ ฆฏฺฏนํ อุปทฺทโว อุปสคฺโคติ ทุกฺขมสฺส มหพฺภยํ. เตนาห ภควา\csromandash{}}

\prose{อวิชฺชาย นิวุโต โลโก, (อชิตาติ ภควา,)}

\prose{เววิจฺฉา ปมาทา นปฺปกาสติ.}

\setlength{\csromangathapagewidth}{0pt}
\setlength{\csromangathawakwidth}{0pt}
\csromangathameasuregroup{ชปฺปาภิเลปนํ พฺรูมิ,}{\csromangathabat{ชปฺปาภิเลปนํ พฺรูมิ,}{ทุกฺข’มสฺส มหพฺภยนฺติ.}}
\csromangathasetleft{ชปฺปาภิเลปนํ พฺรูมิ,}
\csromangathagroup{\csromangathastanza{\csromangathabat{ชปฺปา\-ภิเลปนํ พฺรูมิ,}{ทุกฺข’มสฺส มหพฺภยนฺติ.}}}

\proseitem{3}{\csromanfolio{29}\textbf{สวนฺติ สพฺพธิ โสตา, (อิจฺจายสฺมา อชิโต,)}}

\prose{\textbf{โสตานํ กิํ นิวารณํ.}}

\setlength{\csromangathapagewidth}{0pt}
\setlength{\csromangathawakwidth}{0pt}
\csromangathameasuregroup{\textbf{โสตานํ สํวรํ} พฺรูหิ,}{\csromangathabat{\textbf{โสตานํ สํวรํ} พฺรูหิ,}{เกน \textbf{โสตา} ปิธียเร\textsuperscript{1}\textbf{.}}}
\csromangathasetleft{\textbf{โสตานํ สํวรํ} พฺรูหิ,}
\csromangathagroup{\csromangathastanza{\csromangathabat{\textbf{โสตานํ สํวรํ} พฺรูหิ,}{เกน \textbf{โสตา} ปิธียเร\csromansharedfootnote{47ae57675ed574a1}{ปิถิยฺยเร (สฺยา), ปิถียเร (สีฏฺฐ)}\textbf{.}}}}

\prose{\textbf{สวนฺติ สพฺพธิ โสตา}ติ โสตาติ ตณฺหาโสโต ทิฏฺฐิโสโต กิเลสโสโต ทุจฺจริตโสโต อวิชฺชาโสโต. \textbf{สพฺพธี}ติ สพฺเพสุ อายตเนสุ. \textbf{สวนฺตี}ติ สวนฺติ อาสวนฺติ สนฺทนฺติ ปวตฺตนฺติ. จกฺขุโต รูเป สวนฺติ อาสวนฺติ สนฺทนฺติ ปวตฺตนฺติ. โสตโต สทฺเท สวนฺติ. ฆานโต คนฺเธ สวนฺติ. ชิวฺหาโต รเส สวนฺติ. กายโต โผฏฺฐพฺเพ สวนฺติ. มนโต ธมฺเม สวนฺติ อาสวนฺติ สนฺทนฺติ ปวตฺตนฺติ. จกฺขุโต รูปตณฺหา สวนฺติ อาสวนฺติ สนฺทนฺติ ปวตฺตนฺติ. โสตโต สทฺทตณฺหา สวนฺติ อาสวนฺติ สนฺทนฺติ ปวตฺตนฺติ. ฆานโต คนฺธตณฺหา สวนฺติ. ชิวฺหาโต รสตณฺหา สวนฺติ. กายโต โผฏฺฐพฺพ\-ตณฺหา สวนฺติ. มนโต ธมฺมตณฺหา สวนฺติ อาสวนฺติ สนฺทนฺติ ปวตฺตนฺตีติ สวนฺติ สพฺพธิ โสตา.}

\prose{\textbf{อิจฺจายสฺมา อชิโต}ติ \textbf{อิจฺจา}ติ ปทสนฺธิ ฯเปฯ ปทา\-นุปุพฺพตา\-เปตํ อิจฺจาติ ฯเปฯ อิจฺจายสฺมา อชิโต.}

\prose{\textbf{โสตานํ กิํ นิวารณนฺ}ติ โสตานํ กิํ อาวรณํ นีวรณํ สํวรณํ รกฺขนํ โคปนนฺติ โสตานํ กิํ นิวารณํ.}

\prose{\textbf{โสตานํ สํวรํ พฺรูหี}ติ โสตานํ อาวรณํ นีวรณํ สํวรณํ รกฺขนํ โคปนํ พฺรูหิ อาจิกฺข เทเสหิ ปญฺญเปหิ ปฏฺฐเปหิ วิวราหิ วิภชาหิ อุตฺตานีกโรหิ ปกาเสหีติ โสตานํ สํวรํ พฺรูหิ.}

\prose{\textbf{เกน โสตา ปิธียเร}ติ เกน โสตา ปิธียนฺติ ปจฺฉิชฺชนฺติ น สวนฺติ น อาสวนฺติ น สนฺทนฺติ นปฺป\-วตฺต\-นฺตีติ เกน โสตา ปิธียเร. เตนาห โส พฺราหฺมโณ\csromandash{}}

\prose{\textbf{สวนฺติ สพฺพธิ โสตา, (อิจฺจายสฺมา อชิโต,)}}

\prose{\textbf{โสตานํ กิํ นิวารณํ.}}

\setlength{\csromangathapagewidth}{0pt}
\setlength{\csromangathawakwidth}{0pt}
\csromangathameasuregroup{\textbf{โสตานํ สํวรํ} พฺรูหิ,}{\csromangathabat{\textbf{โสตานํ สํวรํ} พฺรูหิ,}{เกน \textbf{โสตา} ปิธียเร\textbf{.}}}
\csromangathasetleft{\textbf{โสตานํ สํวรํ} พฺรูหิ,}
\csromangathagroup{\csromangathastanza{\csromangathabat{\textbf{โสตานํ สํวรํ} พฺรูหิ,}{เกน \textbf{โสตา} ปิธียเร\textbf{.}}}}

\proseitem{4}{\csromanfolio{30}\textbf{ยานิ โสตานิ โลกสฺมิํ, (อชิตาติ ภควา,)}}

\prose{\textbf{สติ เตสํ นิวารณํ.}}

\setlength{\csromangathapagewidth}{0pt}
\setlength{\csromangathawakwidth}{0pt}
\csromangathameasuregroup{\textbf{โสตานํ สํวรํ} พฺรูมิ,}{\csromangathabat{\textbf{โสตานํ สํวรํ} พฺรูมิ,}{\textbf{ปญฺญา}เยเต ปิธียเร.}}
\csromangathasetleft{\textbf{โสตานํ สํวรํ} พฺรูมิ,}
\csromangathagroup{\csromangathastanza{\csromangathabat{\textbf{โสตานํ สํวรํ} พฺรูมิ,}{\textbf{ปญฺญา}เยเต ปิธียเร.}}}

\prose{\textbf{ยานิ โสตานิ โลกสฺมินฺ}ติ ยานิ เอตานิ โสตานิ มยา กิตฺติตานิ ปกิตฺติตานิ อาจิกฺขิตานิ เทสิตานิ ปญฺญปิตานิ ปฏฺฐปิตานิ วิวริตานิ วิภชิตานิ\csromansharedfootnote{fc733b481d98728c}{วิภตฺตานิ (ก)} อุตฺตานีกตานิ ปกาสิตานิ. เสยฺยถิทํ\csromansharedfootnote{91dcd3a9f4a6e543}{เสยฺยถีทํ (สฺยา)}, ตณฺหาโสโต ทิฏฺฐิโสโต กิเลสโสโต ทุจฺจริตโสโต อวิชฺชาโสโต. \textbf{โลกสฺมินฺ}ติ อปายโลเก มนุสฺสโลเก เทวโลเก ขนฺธโลเก ธาตุโลเก อายตนโลเกติ ยานิ โสตานิ โลกสฺมิํ. \textbf{อชิตา}ติ ภควา ตํ พฺราหฺมณํ นาเมน อาลปติ.}

\prose{\textbf{สติ เตสํ นิวารณนฺ}ติ \textbf{สตี}ติ ยา สติ อนุสฺสติ ปฏิสฺสติ สติ สรณตา ธารณตา อปิลาปนตา อสมฺมุสฺสนตา, สติ สตินฺทฺริยํ สติพลํ สมฺมาสติ สติสมฺโพชฺฌงฺโค เอกายนมคฺโค. อยํ วุจฺจติ สติ. \textbf{นิวารณนฺ}ติ อาวรณํ นีวรณํ สํวรณํ รกฺขนํ โคปนนฺติ สติ เตสํ นิวารณํ.}

\prose{\textbf{โสตานํ สํวรํ พฺรูมี}ติ โสตานํ อาวรณํ นีวรณํ สํวรณํ รกฺขนํ โคปนํ พฺรูมิ อาจิกฺขามิ ฯเปฯ อุตฺตานีกโรมิ ปกาเสมีติ โสตานํ สํวรํ พฺรูมิ.}

\prose{\textbf{ปญฺญาเยเต ปิธียเร}ติ \textbf{ปญฺญา}ติ ยา ปญฺญา ปชานนา ฯเปฯ อโมโห ธมฺมวิจโย สมฺมาทิฏฺฐิ. ปญฺญาเยเต ปิธียเรติ ปญฺญาเยเต โสตา ปิธียนฺติ ปจฺฉิชฺชนฺติ น สวนฺติ น อาสวนฺติ น สนฺทนฺติ นปฺปวตฺตนฺติ. “สพฺเพ สงฺขารา อนิจฺจา”ติ ชานโต ปสฺสโต ปญฺญาเยเต โสตา ปิธียนฺติ ปจฺฉิชฺชนฺติ น สวนฺติ น อาสวนฺติ น สนฺทนฺติ นปฺปวตฺตนฺติ. “สพฺเพ สงฺขารา ทุกฺขา”ติ ชานโต ปสฺสโต ปญฺญาเยเต โสตา ปิธียนฺติ ปจฺฉิชฺชนฺติ น สวนฺติ น อาสวนฺติ น สนฺทนฺติ นปฺปวตฺตนฺติ. “สพฺเพ สงฺขารา อนตฺตา”ติ ชานโต ปสฺสโต ปญฺญาเยเต โสตา ปิธียนฺติ ปจฺฉิชฺชนฺติ น สวนฺติ น อาสวนฺติ น สนฺทนฺติ นปฺปวตฺตนฺติ. “อวิชฺชาปจฺจยา สงฺขารา”ติ ชานโต ปสฺสโต ปญฺญาเยเต โสตา ปิธียนฺติ \csromanfolio{31}ปจฺฉิชฺชนฺติ น สวนฺติ น อาสวนฺติ น สนฺทนฺติ นปฺปวตฺตนฺติ. “สงฺขาร\-ปจฺจยา วิญฺญาณนฺ”ติ. “วิญฺญาณปจฺจยา นามรูปนฺ”ติ. “นามรูป\-ปจฺจยา สฬายตนนฺ”ติ. “สฬายตน\-ปจฺจยา ผสฺโส”ติ. “ผสฺสปจฺจยา เวทนา”ติ. “เวทนาปจฺจยา ตณฺหา”ติ. “ตณฺหาปจฺจยา อุปาทานนฺ”ติ. “อุปาทานปจฺจยา ภโว”ติ. “ภวปจฺจยา ชาตี”ติ. “ชาติปจฺจยา ชรามรณนฺ”ติ ชานโต ปสฺสโต ปญฺญาเยเต โสตา ปิธียนฺติ \textbf{ปจฺฉิชฺชนฺติ น สวนฺติ น อาสวนฺติ น สนฺทนฺติ นปฺปวตฺตนฺติ. “อวิชฺชานิโรธา} สงฺขาร\-นิโรโธ”ติ. “สงฺขาร\-นิโรธา วิญฺญาณนิโรโธ”ติ. “วิญฺญาณนิโรธา \textbf{นามรูป}\-\textbf{นิโรโธ”ติ. “นามรูป}\-\textbf{นิโรธา สฬายตน}\-\textbf{นิโรโธ”ติ} “\textbf{สฬายตน}\-\textbf{นิโรธา ผสฺสนิโรโธ”ติ. “ผสฺสนิโรธา เวทนานิโรโธ”ติ.} “เวทนานิโรธา ตณฺหานิโรโธ”ติ. “ตณฺหานิโรธา อุปาทานนิโรโธ”ติ. “อุปาทานนิโรธา ภวนิโรโธ”ติ. “ภวนิโรธา ชาตินิโรโธ”ติ. “ชาตินิโรธา ชรามรณ\-นิโรโธ”ติ ชานโต ปสฺสโต ปญฺญาเยเต โสตา ปิธียนฺติ ปจฺฉิชฺชนฺติ น สวนฺติ น อาสวนฺติ น สนฺทนฺติ นปฺปวตฺตนฺติ. “อิทํ ทุกฺขนฺ”ติ. “อยํ ทุกฺขสมุทโย”ติ. “อยํ ทุกฺขนิโรโธ”ติ. “อยํ ทุกฺขนิโรธ\-คามินี ปฏิปทา”ติ ชานโต ปสฺสโต ปญฺญาเยเต โสตา ปิธียนฺติ ปจฺฉิชฺชนฺติ น สวนฺติ น อาสวนฺติ น สนฺทนฺติ นปฺปวตฺตนฺติ. “อิเม ธมฺมา อาสวา”ติ. “อยํ อาสวสมุทโย”ติ. “อยํ อาสวนิโรโธ”ติ. “อยํ อาสว\-นิโรธ\-คามินี ปฏิปทา”ติ ชานโต ปสฺสโต ปญฺญาเยเต โสตา ปิธียนฺติ ปจฺฉิชฺชนฺติ น สวนฺติ น อาสวนฺติ น สนฺทนฺติ นปฺปวตฺตนฺติ. “อิเม ธมฺมา อภิญฺเญยฺยา”ติ. “อิเม ธมฺมา ปริญฺเญยฺยา”ติ. “อิเม ธมฺมา ปหาตพฺพา”ติ. “อิเม ธมฺมา ภาเวตพฺพา”ติ. “อิเม ธมฺมา สจฺฉิกาตพฺพา”ติ ชานโต ปสฺสโต ปญฺญาเยเต โสตา ปิธิยนฺติ ปจฺฉิชฺชนฺติ น สวนฺติ น อาสวนฺติ น สนฺทนฺติ นปฺปวตฺตนฺติ. ฉนฺนํ ผสฺสา\-ยตนานํ สมุทยญฺจ อตฺถงฺคมญฺจ อสฺสาทญฺจ อาทีนวญฺจ นิสฺสรณญฺจ ชานโต ปสฺสโต ปญฺญาเยเต โสตา ปิธียนฺติ ปจฺฉิชฺชนฺติ น สวนฺติ น อาสวนฺติ น สนฺทนฺติ นปฺปวตฺตนฺติ. ปญฺจนฺนํ อุปาทาน\-กฺขนฺธานํ สมุทยญฺจ อตฺถงฺคมญฺจ อสฺสาทญฺจ อาทีนวญฺจ นิสฺสรณญฺจ ชานโต ปสฺสโต. จตุนฺนํ มหาภูตานํ สมุทยญฺจ อตฺถงฺคมญฺจ อสฺสาทญฺจ อาทีนวญฺจ นิสฺสรณญฺจ ชานโต ปสฺสโต. “ยํ กิญฺจิ สมุทย\-ธมฺมํ, สพฺพํ ตํ นิโรธธมฺมนฺ”ติ ชานโต ปสฺสโต ปญฺญาเยเต โสตา ปิธียนฺติ ปจฺฉิชฺชนฺติ น สวนฺติ น \csromanfolio{32}อาสวนฺติ น สนฺทนฺติ นปฺป\-วตฺต\-นฺตีติ ปญฺญาเยเต ปิธียเร. เตนาห ภควา\csromandash{}}

\proseitem{1}{ยานิ โสตานิ โลกสฺมิํ, (อชิตาติ ภควา,)}

\prose{สติ เตสํ นิวารณํ.}

\setlength{\csromangathapagewidth}{0pt}
\setlength{\csromangathawakwidth}{0pt}
\csromangathameasuregroup{โสตานํ สํวรํ พฺรูมิ,}{\csromangathabat{โสตานํ สํวรํ พฺรูมิ,}{\textbf{ปญฺญา}เยเต ปิธียเรติ.}}
\csromangathasetleft{โสตานํ สํวรํ พฺรูมิ,}
\csromangathagroup{\csromangathastanza{\csromangathabat{โสตานํ สํวรํ พฺรูมิ,}{\textbf{ปญฺญา}เยเต ปิธียเรติ.}}}

\prose{\textbf{ปญฺญา เจว สติ จาปิ, (อิจฺจายสฺมา อชิโต,)}}

\prose{\textbf{นามรูปญฺจ มาริส.}}

\setlength{\csromangathapagewidth}{0pt}
\setlength{\csromangathawakwidth}{0pt}
\csromangathameasuregroup{\textbf{เอตํ เม ปุฏฺโฐ ปพฺ}รูหิ,}{\csromangathabat{\textbf{เอตํ เม ปุฏฺโฐ ปพฺ}รูหิ,}{กตฺเถตํ อุปรุชฺฌติ.}}
\csromangathasetleft{\textbf{เอตํ เม ปุฏฺโฐ ปพฺ}รูหิ,}
\csromangathagroup{\csromangathastanza{\csromangathabat{\textbf{เอตํ เม ปุฏฺโฐ ปพฺ}รูหิ,}{กตฺเถตํ อุปรุชฺฌติ.}}}

\prose{\textbf{ปญฺญา เจว สติ จาปี}ติ \textbf{ปญฺญา}ติ ยา ปญฺญา ปชานนา วิจโย ปวิจโย ธมฺมวิจโย สลฺลกฺขณา อุปลกฺขณา ปจฺจุ\-ปลกฺขณา ปณฺฑิจฺจํ โกสลฺลํ เนปุญฺญํ เวภพฺยา จินฺตา อุปปริกฺขา ภูรี\csromansharedfootnote{e6245be23e5483fd}{ภูริ (ก)} เมธา ปริณายิกา วิปสฺสนา สมฺปชญฺญํ ปโตโท ปญฺญา ปญฺญินฺทฺริยํ ปญฺญาพลํ ปญฺญาสตฺถํ ปญฺญาปาสาโท ปญฺญาอาโลโก ปญฺญาโอภาโส ปญฺญาปชฺโชโต ปญฺญารตนํ อโมโห ธมฺมวิจโย สมฺมาทิฏฺฐิ, \textbf{สตี}ติ ยา สติ อนุสฺสติ ฯเปฯ สมฺมาสตีติ ปญฺญา เจว สติจาปิ, อิจฺจายสฺมา อชิโต.}

\prose{\textbf{นามรูปญฺจ มาริสา}ติ \textbf{นามนฺ}ติ จตฺตาโร อรูปิโน ขนฺธา. \textbf{รูปนฺ}ติ จตฺตาโร จ มหาภูตา จตุนฺนญฺจ มหาภูตานํ อุปาทายรูปํ. \textbf{มาริสา}ติ ปิยวจนํ ครุวจนํ สคารวสปฺปติสฺสาธิวจนเมตํ มาริสาติ นามรูปญฺจ มาริส.}

\prose{\textbf{เอตํ เม ปุฏฺโฐ ปพฺรูหี}ติ \textbf{เอตํ เม}ติ ยํ ปุจฺฉามิ ยํ ยาจามิ ยํ อชฺเฌสามิ ยํ ปสาเทมิ. \textbf{ปุฏฺโฐ}ติ ปุจฺฉิโต ยาจิโต อชฺเฌสิโต ปสาทิโต. \textbf{ปพฺรูหี}ติ พฺรูหิ อาจิกฺขาหิ เทเสหิ ปญฺญเปหิ ปฏฺฐเปหิ วิวราหิ วิภชาหิ\csromansharedfootnote{1b9ab0a19d52c107}{วิวเรหิ วิภเชหิ (ก)} อุตฺตานีกโรหิ ปกาเสหีติ เอตํ เม ปุฏฺโฐ ปพฺรูหิ.}

\prose{\textbf{กตฺเถตํ อุปรุชฺฌตี}ติ กตฺเถตํ นิรุชฺฌติ วูปสมฺมติ อตฺถํ คจฺฉติ ปฏิปฺปสฺสมฺภตีติ กตฺเถตํ อุปรุชฺฌติ. เตนาห โส พฺราหฺมโณ\csromandash{}}

\prose{\csromanfolio{33}ปญฺญา เจว สติ จาปิ, (อิจฺจายสฺมา อชิโต,)}

\prose{นามรูปญฺจ มาริส.}

\setlength{\csromangathapagewidth}{0pt}
\setlength{\csromangathawakwidth}{0pt}
\csromangathameasuregroup{เอตํ เม ปุฏฺโฐ ปพฺรูหิ,}{\csromangathabat{เอตํ เม ปุฏฺโฐ ปพฺรูหิ,}{กตฺเถตํ อุปรุชฺฌตีติ.}}
\csromangathasetleft{เอตํ เม ปุฏฺโฐ ปพฺรูหิ,}
\csromangathagroup{\csromangathastanza{\csromangathabat{เอตํ เม ปุฏฺโฐ ปพฺรูหิ,}{กตฺเถตํ อุปรุชฺฌตีติ.}}}

\proseitem{6}{\textbf{ยเมตํ ปญฺหํ อปุจฺฉิ, อชิต ตํ วทามิ เต.}}

\setlength{\csromangathapagewidth}{0pt}
\setlength{\csromangathawakwidth}{0pt}
\csromangathameasuregroup{\textbf{ยตฺถ นามญฺจ รู}ปญฺจ, \\ \textbf{วิญฺญาณสฺส นิโรธฺ}เอน,}{\csromangathabat{\textbf{ยตฺถ นามญฺจ รู}ปญฺจ,}{อเสสํ อุปรุชฺฌติ.} \\ \csromangathabat{\textbf{วิญฺญาณสฺส นิโรธฺ}เอน,}{เอตฺเถตํ อุปรุชฺฌติ.}}
\csromangathasetleft{\textbf{ยตฺถ นามญฺจ รู}ปญฺจ, \\ \textbf{วิญฺญาณสฺส นิโรธฺ}เอน,}
\csromangathagroup{\csromangathastanza{\csromangathabat{\textbf{ยตฺถ นามญฺจ รู}ปญฺจ,}{อเสสํ อุปรุชฺฌติ.} \\ \csromangathabat{\textbf{วิญฺญาณสฺส นิโรธฺ}เอน,}{เอตฺเถตํ อุปรุชฺฌติ.}}}

\prose{\textbf{ยเมตํ ปญฺหํ อปุจฺฉี}ติ \textbf{ยเมตนฺ}ติ ปญฺญญฺจ สติญฺจ นามรูปญฺจ. \textbf{อปุจฺฉี}ติ อปุจฺฉสิ ยาจสิ อชฺเฌสสิ\csromansharedfootnote{58a9c0351855852d}{อชฺเฌสิ (ก)} ปสาเทสีติ ยเมตํ ปญฺหํ อปุจฺฉิ.}

\prose{\textbf{อชิต ตํ วทามิ เต}ติ \textbf{อชิตา}ติ ภควา ตํ พฺราหฺมณํ นาเมน อาลปติ. \textbf{ตนฺ}ติ ปญฺญญฺจ สติญฺจ นามรูปญฺจ. \textbf{วทามี}ติ วทามิ อาจิกฺขามิ เทเสมิ ปญฺญเปมิ ปฏฺฐเปมิ วิวรามิ วิภชามิ อุตฺตานีกโรมิ ปกาเสมีติ อชิต ตํ วทามิ เต.}

\prose{\textbf{ยตฺถ นามญฺจ รูปญฺจ, อเสสํ อุปรุชฺฌตี}ติ \textbf{นามนฺ}ติ จตฺตาโร อรูปิโน ขนฺธา. \textbf{รูปนฺ}ติ จตฺตาโร จ มหาภูตา จตุนฺนญฺจ มหาภูตานํ อุปาทายรูปํ. \textbf{อเสสนฺ}ติ สพฺเพน สพฺพํ สพฺพถา สพฺพํ อเสสํ นิสฺเสสํ ปริยาทิย\-น\-วจนเมตํ\csromansharedfootnote{ad7f6938b357dcc4}{ปริยาทายวจนเมตํ (สฺยา, ก)} อเสสนฺติ. \textbf{อุปรุชฺฌตี}ติ นิรุชฺฌติ วูปสมฺมติ อตฺถํ คจฺฉติ ปฏิปฺปสฺสมฺภตีติ ยตฺถ นามญฺจ รูปญฺจ, อเสสํ อุปรุชฺฌติ.}

\prose{\textbf{วิญฺญาณสฺส นิโรเธน, เอตฺเถตํ อุปรุชฺฌตี}ติ โสตาปตฺติ\-มคฺคญาเณน อภิสงฺขาร\-วิญฺญาณสฺส นิโรเธน สตฺต ภเว ฐเปตฺวา อนมตคฺเค สํสาเร เย อุปฺปชฺเชยฺยุํ นามญฺจ รูปญฺจ, เอตฺเถเต นิรุชฺฌนฺติ วูปสมฺมนฺติ อตฺถํ คจฺฉนฺติ ปฏิปฺปสฺสมฺภนฺติ. สกทาคามิ\-มคฺคญาเณน อภิสงฺขาร\-วิญฺญาณสฺส นิโรเธน เทฺว ภเว ฐเปตฺวา ปญฺจสุ ภเวสุ เย อุปฺปชฺเชยฺยุํ นามญฺจ รูปญฺจ. เอตฺเถเต นิรุชฺฌนฺติ วูปสมฺมนฺติ อตฺถํ คจฺฉนฺติ ปฏิปฺปสฺสมฺภนฺติ. อนาคามิ\-มคฺคญาเณน อภิสงฺขาร\-วิญฺญาณสฺส นิโรเธน เอกํ ภวํ ฐเปตฺวา รูปธาตุยา วา อรูปธาตุยา วา เย อุปฺปชฺเชยฺยุํ นามญฺจ รูปญฺจ, เอตฺเถเต นิรุชฺฌนฺติ วูปสมฺมนฺติ อตฺถํ คจฺฉนฺติ ปฏิปฺปสฺสมฺภนฺติ. อรหตฺตมคฺค\-ญาเณน อภิสงฺขาร\-วิญฺญาณสฺส นิโรเธน เย อุปฺปชฺเชยฺยุํ นามญฺจ รูปญฺจ, เอตฺเถเต นิรุชฺฌนฺติ วูปสมฺมนฺติ อตฺถํ คจฺฉนฺติ \csromanfolio{34}ปฏิปฺปสฺสมฺภนฺติ. อรหโต อนุปาทิเสสาย นิพฺพานธาตุยา ปรินิพฺพายนฺตสฺส จริม\-วิญฺญาณสฺส นิโรเธน ปญฺญา จ สติ จ นามญฺจ รูปญฺจ, เอตฺเถเต นิรุชฺฌนฺติ วูปสมฺมนฺติ อตฺถํ คจฺฉนฺติ ปฏิปฺปสฺสมฺภนฺตีติ วิญฺญาณสฺส นิโรเธน, เอตฺเถตํ อุปรุชฺฌติ. เตนาห ภควา\csromandash{}}

\setlength{\csromangathapagewidth}{0pt}
\setlength{\csromangathawakwidth}{0pt}
\csromangathameasuregroup{ยเมตํ ปญฺหํ อปุจฺฉิ, \\ ยตฺถ นามญฺจ รูปญฺจ, \\ วิญฺญาณสฺส นิโรเธน, \\ \textbf{เย จ สงฺขาต}\textbf{ธมฺ}มาเส, \\ \textbf{เตสํ เม นิปโก อิรฺ}อิยํ,}{\csromangathabat{ยเมตํ ปญฺหํ อปุจฺฉิ,}{อชิต ตํ วทามิ เต.} \\ \csromangathabat{ยตฺถ นามญฺจ รูปญฺจ,}{อเสสํ อุปรุชฺฌติ.} \\ \csromangathabat{วิญฺญาณสฺส นิโรเธน,}{เอตฺเถตํ อุปรุชฺฌตีติ.} \\ \csromangathabat{\textbf{เย จ สงฺขาต}\textbf{ธมฺ}มาเส,}{เย จ \textbf{เสขา}\textsuperscript{1} \textbf{ปุถู อิธ.}} \\ \csromangathabat{\textbf{เตสํ เม นิปโก อิรฺ}อิยํ,}{ปุฏฺโฐ ปพฺรูหิ มาริส.}}
\csromangathasetleft{ยเมตํ ปญฺหํ อปุจฺฉิ, \\ ยตฺถ นามญฺจ รูปญฺจ, \\ วิญฺญาณสฺส นิโรเธน, \\ \textbf{เย จ สงฺขาต}\textbf{ธมฺ}มาเส, \\ \textbf{เตสํ เม นิปโก อิรฺ}อิยํ,}
\csromangathagroup{\csromangathastanza{\csromangathabat{ยเมตํ ปญฺหํ อปุจฺฉิ,}{อชิต ตํ วทามิ เต.} \\ \csromangathabat{ยตฺถ นามญฺจ รูปญฺจ,}{อเสสํ อุปรุชฺฌติ.} \\ \csromangathabat{วิญฺญาณสฺส นิโรเธน,}{เอตฺเถตํ อุปรุชฺฌตีติ.}}\csromangathastanza{\csromangathaitembat{6}{\textbf{เย จ สงฺขาต}\-\textbf{ธมฺ}มาเส,}{เย จ \textbf{เสขา}\csromansharedfootnote{22352c351fad0591}{เสกฺขา (สฺยา, ก)} \textbf{ปุถู อิธ.}} \\ \csromangathacontbat{\textbf{เตสํ เม นิปโก อิรฺ}อิยํ,}{ปุฏฺโฐ ปพฺรูหิ มาริส.}}}

\prose{\textbf{เย จ สงฺขาตธมฺมา}\-\textbf{เส}ติ สงฺขาตธมฺมา วุจฺจนฺติ อรหนฺโต ขีณาสวา, กิํ การณา สงฺขาตธมฺมา วุจฺจนฺติ อรหนฺโต ขีณาสวา. เต สงฺขาตธมฺมา ญาตธมฺมา ตุลิตธมฺมา ตีริตธมฺมา วิภูตธมฺมา วิภาวิต\-ธมฺมา. “สพฺเพ สงฺขารา อนิจฺจา”ติ สงฺขาตธมฺมา ญาตธมฺมา ตุลิตธมฺมา ตีริตธมฺมา วิภูตธมฺมา วิภาวิต\-ธมฺมา. “สพฺเพ สงฺขารา ทุกฺขา”ติ สงฺขาตธมฺมา ฯเปฯ. “สพฺเพ ธมฺมา อนตฺตา”ติ สงฺขาตธมฺมา. “อวิชฺชาปจฺจยา สงฺขารา”ติ สงฺขาตธมฺมา ฯเปฯ. “ยํ กิญฺจิ สมุทย\-ธมฺมํ, สพฺพํ ตํ นิโรธธมฺมนฺ”ติ สงฺขาตธมฺมา ญาตธมฺมา ตุลิตธมฺมา ตีริตธมฺมา วิภูตธมฺมา วิภาวิต\-ธมฺมา. อถ วา เตสํ ขนฺธา สงฺขาตา ธาตุโย สงฺขาตา อายตนานิ สงฺขาตา คติโย สงฺขาตา อุปปตฺติโย สงฺขาตา ปฏิสนฺธิ สงฺขาตา ภวา สงฺขาตา สํสารา สงฺขาตา วฏฺฏา สงฺขาตา. อถ วา เต ขนฺธ\-ปริยนฺเต ฐิตา ธาตุปริยนฺเต ฐิตา อายตน\-ปริยนฺเต ฐิตา คติปริยนฺเต ฐิตา อุปปตฺติ\-ปริยนฺเต ฐิตา ปฏิสนฺธิ\-ปริยนฺเต ฐิตา ภวปริยนฺเต ฐิตา สํสาร\-ปริยนฺเต ฐิตา วฏฺฏปริยนฺเต ฐิตา อนฺติเม ภเว ฐิตา อนฺติเม สมุสฺสเย ฐิตา อนฺติม\-เทห\-ธรา อรหนฺโต.}

\setlength{\csromangathapagewidth}{0pt}
\setlength{\csromangathawakwidth}{0pt}
\csromangathameasuregroup{เตสํ จายํ\textsuperscript{1} ปจฺฉิมโก, \\ ชาติมรณสํสาโร,}{\csromangathabat{เตสํ จายํ\textsuperscript{1} ปจฺฉิมโก,}{จริโมยํ สมุสฺสโย.} \\ \csromangathabat{ชาติมรณสํสาโร,}{นตฺถิ เนสํ ปุนพฺภโวติ.}}
\csromangathasetleft{เตสํ จายํ\textsuperscript{1} ปจฺฉิมโก, \\ ชาติมรณสํสาโร,}
\csromangathagroup{\csromangathastanza{\csromangathabat{เตสํ จายํ\csromansharedfootnote{0d41cd4ebd3a51e4}{ยายํ (ก)} ปจฺฉิมโก,}{จริโมยํ สมุสฺสโย.} \\ \csromangathabat{ชาติ\-มรณ\-สํสาโร,}{นตฺถิ เนสํ ปุนพฺภโวติ.}}}

\prose{ตํการณา สงฺขาตธมฺมา วุจฺจนฺติ อรหนฺโต ขีณาสวาติ. \textbf{เย} จ \textbf{สงฺขาต}\-\textbf{ธมฺมาเส, เย จ เสขา ปุถู อิธา}ติ \textbf{เสขา}ติ กิํ การณา \csromanfolio{35}\textbf{วุจฺจนฺติ เสขา, สิกฺขนฺตีติ เสขา, กิญฺจ สิกฺขนฺติ, อธิสีลมฺปิ สิกฺขนฺติ.} \textbf{อธิจิตฺตมฺปิ สิกฺขนฺติ. อธิปญฺญมฺปิ สิกฺขนฺติ. กตมา อธิสีลสิกฺขา.} อิธ ภิกฺขุ สีลวา โหติ ปาติโมกฺข\-สํวร\-สํวุโต วิหรติ อาจารโคจร\-สมฺปนฺโน อณุมตฺเตสุ วชฺเชสุ ภยทสฺสาวี, สมาทาย สิกฺขติ สิกฺขาปเทสุ ขุทฺทโก สีลกฺขนฺโธ มหนฺโต สีลกฺขนฺโธ สีลํ ปติฏฺฐา อาทิ จรณํ สํยโม สํวโร มุขํ ปมุขํ กุสลานํ ธมฺมานํ สมาปตฺติยา. อยํ อธิสีลสิกฺขา.}

\proseitem{7}{กตมา อธิจิตฺต\-สิกฺขา. อิธ ภิกฺขุ วิวิจฺเจว กาเมหิ ฯเปฯ ปฐมํ ฌานํ. ทุติยํ ฌานํ. ตติยํ ฌานํ. จตุตฺถํ ฌานํ อุปสมฺปชฺช วิหรติ. อยํ อธิจิตฺต\-สิกฺขา.}

\prose{กตมา อธิปญฺญา\-สิกฺขา. อิธ ภิกฺขุ ปญฺญวา โหติ อุทย\-ตฺถคามินิยา ปญฺญาย สมนฺนาคโต อริยาย นิพฺเพธิกาย สมฺมา ทุกฺข\-กฺขย\-คามินิยา. โส “อิทํ ทุกฺขนฺ”ติ ยถาภูตํ ปชานาติ. “อยํ ทุกฺขสมุทโย”ติ. “อยํ ทุกฺขนิโรโธ”ติ. “อยํ ทุกฺขนิโรธ\-คามินี ปฏิปทา”ติ ยถาภูตํ ปชานาติ. “อิเม อาสวา”ติ. “อยํ อาสวสมุทโย”ติ. “อยํ อาสวนิโรโธ”ติ. “อยํ อาสว\-นิโรธ\-คามินี ปฏิปทา”ติ ยถาภูตํ ปชานาติ. อยํ อธิปญฺญา\-สิกฺขา. อิมา ติสฺโส สิกฺขาโย อาวชฺชนฺตา สิกฺขนฺติ. ชานนฺตา สิกฺขนฺติ. ปสฺสนฺตา สิกฺขนฺติ. จิตฺตํ อธิฏฺฐหนฺตา สิกฺขนฺติ. สทฺธาย อธิมุจฺจนฺตา สิกฺขนฺติ. วีริยํ\csromansharedfootnote{4c6eec0257bfc5c7}{วิริยํ (สฺยา)} ปคฺคณฺหนฺตา สิกฺขนฺติ. สติํ อุปฏฺฐเปนฺตา สิกฺขนฺติ. จิตฺตํ สมาทหนฺตา สิกฺขนฺติ. ปญฺญาย ปชานนฺตา สิกฺขนฺติ. อภิญฺเญยฺยํ อภิชานนฺตา สิกฺขนฺติ. ปริญฺเญยฺยํ ปริชานนฺตา สิกฺขนฺติ. ปหาตพฺพํ ปชหนฺตา สิกฺขนฺติ. ภาเวตพฺพํ ภาเวนฺตา สิกฺขนฺติ. สจฺฉิกาตพฺพํ สจฺฉิกโรนฺตา สิกฺขนฺติ อาจรนฺติ สมาจรนฺติ สมาทาย วตฺตนฺติ. ตํการณา วุจฺจนฺติ เสขา. \textbf{ปุถู}ติ พหุกา. เอเต เสขา โสตาปนฺนา จ ปฏิปนฺนา จ สกทาคามิโน จ ปฏิปนฺนา จ อนาคามิโน จ ปฏิปนฺนา จ อรหนฺโต จ ปฏิปนฺนา จ. \textbf{อิธา}ติ อิมิสฺสา ทิฏฺฐิยา อิมิสฺสา ขนฺติยา อิมิสฺสา รุจิยา อิมสฺมิํ อาทาเย อิมสฺมิํ ธมฺเม อิมสฺมิํ วินเย อิมสฺมิํ ธมฺมวินเย อิมสฺมิํ ปาวจเน อิมสฺมิํ พฺรหฺมจริเย อิมสฺมิํ สตฺถุสาสเน อิมสฺมิํ อตฺตภาเว อิมสฺมิํ มนุสฺสโลเกติ เย จ เสขา ปุถู อิธ.}

\proseitem{1}{\csromanfolio{36}\textbf{เตสํ เม นิปโก อิริยํ, ปุฏฺโฐ ปพฺรูหิ มาริสา}ติ ตฺวมฺปิ นิปโก ปณฺฑิโต ปญฺญวา พุทฺธิมา ญาณี เมธาวี. เตสํ สงฺขาตธมฺมานญฺจ เสกฺขานญฺจ อิริยํ จริยํ วุตฺติ ปวตฺติ อาจรํ โคจรํ วิหารํ ปฏิปทํ. \textbf{ปุฏฺโฐ}ติ ปุจฺฉิโต ยาจิโต อชฺเฌสิโต ปสาทิโต. \textbf{ปพฺรูหี}ติ พฺรูหิ อาจิกฺขาหิ เทเสหิ ปญฺญเปหิ ปฏฺฐเปหิ วิวราหิ วิภชาหิ อุตฺตานีกโรหิ ปกาเสหิ. \textbf{มาริสา}ติ ปิยวจนํ ครุวจนํ สคารวสปฺปติสฺสาธิวจนเมตํ มาริสาติ เตสํ เม นิปโก อิริยํ, ปุฏฺโฐ ปพฺรูหิ มาริส. เตนาห โส พฺราหฺมโณ\csromandash{}}

\setlength{\csromangathapagewidth}{0pt}
\setlength{\csromangathawakwidth}{0pt}
\csromangathameasuregroup{เย จ สงฺขาตธมฺมาเส,}{\csromangathabat{เย จ สงฺขาตธมฺมาเส,}{เย จ เสขา ปุถู อิธ.}}
\csromangathasetleft{เย จ สงฺขาตธมฺมาเส,}
\csromangathagroup{\csromangathastanza{\csromangathabat{เย จ สงฺขาต\-ธมฺมาเส,}{เย จ เสขา ปุถู อิธ.}}}

\prose{เตสํ เม นิปโก อิริยํ, ปุฏฺโฐ ปพฺรูหิ มาริสาติ.}

\setlength{\csromangathapagewidth}{0pt}
\setlength{\csromangathawakwidth}{0pt}
\csromangathameasuregroup{\textbf{กาเมสุ นาภิคิชฺชฺ}เหยฺย, \\ \textbf{กุสโล สพฺพธมฺมา}นํ,}{\csromangathabat{\textbf{กาเมสุ นาภิคิชฺชฺ}เหยฺย,}{มนสา’นาวิโล สิยา.} \\ \csromangathabat{\textbf{กุสโล สพฺพธมฺมา}นํ,}{สโต ภิกฺขุ ปริพฺพเช.}}
\csromangathasetleft{\textbf{กาเมสุ นาภิคิชฺชฺ}เหยฺย, \\ \textbf{กุสโล สพฺพธมฺมา}นํ,}
\csromangathagroup{\csromangathastanza{\csromangathaitembat{8}{\textbf{กาเมสุ นาภิคิชฺชฺ}เหยฺย,}{มนสา’นาวิโล สิยา.} \\ \csromangathacontbat{\textbf{กุสโล สพฺพธมฺมา}นํ,}{สโต ภิกฺขุ ปริพฺพเช.}}}

\prose{\textbf{กาเมสุ นาภิคิชฺเฌยฺยา}ติ กามาติ อุทฺทานโต เทฺว \textbf{กามา} วตฺถุกามา จ กิเลสกามา จ, กตเม วตฺถุกามา. มนาปิกา รูปา มนาปิกา สทฺทา มนาปิกา คนฺธา มนาปิกา รสา มนาปิกา โผฏฺฐพฺพา, อตฺถรณา ปาวุรณา\csromansharedfootnote{0706b3990bfdc13a}{ปาปุรณา (สฺยา) 2. ราชฐานิโย (ก) 3. นตฺถิ สฺยาโปตฺถํเก, ขุ 7. 2 ปิฏฺเฐปิ.} ทาสิทาสา อเชฬกา กุกฺกุฏสูกรา หตฺถิ\-ควา\-สฺส\-วฬวา เขตฺตํ วตฺถุ หิรญฺญํ สุวณฺณํ คามนิคม\-ราชธานิโย\csromansharedfootnote{5acf76565d82d73e}{ราชฐานิโย (ก)} รฏฺฐญฺจ ชนปโท จ โกโส จ โกฏฺฐาคารญฺจ, ยํ กิญฺจิ รชนียวตฺถุ วตฺถุกามา. อปิ จ อตีตา กามา อนาคตา กามา ปจฺจุปฺปนฺนา กามา, อชฺฌตฺตา กามา พหิทฺธา กามา อชฺฌตฺต\-พหิทฺธา กามา, หีนา กามา มชฺฌิมา กามา ปณีตา กามา, อาปายิกา กามา มานุสิกา กามา ทิพฺพา กามา, ปจฺจุปฏฺฐิตา กามา, นิมฺมิตา กามา ปรนิมฺมิตา กามา, ปริคฺคหิตา กามา อปริคฺคหิตา กามา มมายิตา กามา อมมายิตา กามา, สพฺเพปิ กามาวจรา ธมฺมา, สพฺเพปิ รูปาวจรา ธมฺมา, สพฺเพปิ อรูปาวจรา ธมฺมา, ตณฺหาวตฺถุกา ตณฺหารมฺมณา, กามนียฏฺเฐน รชนียฏฺเฐน มทนียฏฺเฐน รมณียฏฺเฐน\csromansharedfootnote{157be1c205410153}{นตฺถิ สฺยาโปตฺถํเก, ขุ 7. 2 ปิฏฺเฐปิ.} กามา. อิเม วุจฺจนฺติ วตฺถุกามา.}

\prose{กตเม กิเลสกามา. ฉนฺโท กาโม ราโค กาโม ฉนฺทราโค กาโม สงฺกปฺโป กาโม ราโค กาโม สงฺกปฺปราโค กาโม. โย กาเมสุ กามจฺฉนฺโท กามราโค กามนนฺที กามตณฺหา \csromanfolio{37}กามสิเนโห กามปิปาสา กามปริฬาโห กามเคโธ กามมุจฺฉา \textbf{กามชฺโฌสานํ กาโมโฆ กามโยโค กามุปาทานํ} กามจฺฉนฺท\-นีวรณํ.}

\setlength{\csromangathapagewidth}{0pt}
\setlength{\csromangathawakwidth}{0pt}
\csromangathameasuregroup{อทฺทสํ กาม เต มูลํ, \\ น ตํ สงฺกปฺปยิสฺสามิ,}{\csromangathabat{อทฺทสํ กาม เต มูลํ,}{สงฺกปฺปา กาม ชายสิ.} \\ \csromangathabat{น ตํ สงฺกปฺปยิสฺสามิ,}{เอวํ กาม น เหหิสีติ.}}
\csromangathasetleft{อทฺทสํ กาม เต มูลํ, \\ น ตํ สงฺกปฺปยิสฺสามิ,}
\csromangathagroup{\csromangathastanza{\csromangathabat{อทฺทสํ กาม เต มูลํ,}{สงฺกปฺปา กาม ชายสิ.} \\ \csromangathabat{น ตํ สงฺกปฺป\-ยิสฺสามิ,}{เอวํ กาม น เหหิสีติ.}}}

\prose{อิเม วุจฺจนฺติ กิเลสกามา. \textbf{เคโธ} วุจฺจติ ตณฺหา, โย ราโค สาราโค ฯเปฯ อภิชฺฌา โลโภ อกุสลมูลํ. \textbf{กาเมสุ นาภิคิชฺเฌยฺยา}ติ กิเลสกาเมน วตฺถุกาเมสุ นาภิคิชฺเฌยฺย น ปลิคิชฺเฌยฺย น ปลิพุนฺเธยฺย\csromansharedfootnote{30d16982da064cd2}{ปลิพุชฺเฌยฺย (สฺยา)} อคิทฺโธ อสฺส อคธิโต อมุจฺฉิโต อนชฺฌาปนฺโน\csromansharedfootnote{e611116e8ad840aa}{อนชฺโฌปนฺโน (สฺยา)} วีตเคโธ วิคตเคโธ จตฺตเคโธ วนฺตเคโธ มุตฺตเคโธ ปหีนเคโธ ปฏินิสฺสฏฺฐ\-เคโธ วีตราโค วิคตราโค จตฺตราโค วนฺตราโค มุตฺตราโค ปหีนราโค ปฏินิสฺสฏฺฐ\-ราโค นิจฺฉาโต นิพฺพุโต สีติภูโต สุขปฺปฏิสํเวที พฺรหฺมภูเตน อตฺตนา วิหเรยฺยาติ กาเมสุ นาภิคิชฺเฌยฺย.}

\prose{\textbf{มนสา’นาวิโล สิยา}ติ \textbf{มโน}ติ ยํ จิตฺตํ มโน มานสํ หทยํ ปณฺฑรํ มโน มนายตนํ มนินฺทฺริยํ วิญฺญาณํ วิญฺญาณ\-กฺขนฺโธ ตชฺชา มโนวิญฺญาณ\-ธาตุ. กาย\-ทุจฺจริเตน จิตฺตํ อาวิลํ โหติ ลุฬิตํ เอริตํ ฆฏฺฏิตํ จลิตํ ภนฺตํ อวูปสนฺตํ. วจีทุจฺจริเตน. มโน\-ทุจฺจริเตน. ราเคน. โทเสน. โมเหน. โกเธน. อุปนาเหน. มกฺเขน. ปฬาเสน. อิสฺสาย. มจฺฉริเยน. มายาย. สาเฐยฺเยน. ถมฺเภน. สารมฺเภน. มาเนน. อติมาเนน. มเทน. ปมาเทน. สพฺพกิเลเสหิ. สพฺพ\-ทุจฺจริเตหิ. สพฺพฑาเหหิ. สพฺพ\-ปริฬาเหหิ. สพฺพสนฺตาเปหิ. สพฺพากุสลา\-ภิสงฺขาเรหิ จิตฺตํ อาวิลํ โหติ ลุฬิตํ เอริตํ ฆฏฺฏิตํ จลิตํ ภนฺตํ อวูปสนฺตํ. \textbf{มนสา’นาวิโล สิยา}ติ จิตฺเตน อนาวิโล สิยา อลุฬิโต อเนริโต อฆฏฺฏิโต อจลิโต อภนฺโต วูปสนฺโต อาวิลกเร กิเลเส ชเหยฺย ปชเหยฺย วิโนเทยฺย พฺยนฺตีกเรยฺย\csromansharedfootnote{e838edc83087b9e2}{พฺยนฺติํ กเรยฺย (ก)} อนภาวํ คเมยฺย, อาวิลกเรหิ กิเลเสหิ จ อารโต\csromansharedfootnote{da0de921abd08b72}{อารโต อสฺส (ก) ขุ 7. 54 ปิฏฺเฐ ปสฺส.} วิรโต ปฏิวิรโต นิกฺขนฺโต นิสฺสโฏ วิปฺปมุตฺโต วิสญฺญุตฺโต วิมริยาทิกเตน เจตสา วิหเรยฺยาติ มนสฺสา’นาวิโล สิยา.}

\proseitem{1}{\csromanfolio{38}\textbf{กุสโล สพฺพ}\-\textbf{ธมฺมานนฺ}ติ “สพฺเพ สงฺขารา อนิจฺจา”ติ กุสโล สพฺพธมฺมานํ, “สพฺเพ สงฺขารา ทุกฺขา”ติ กุสโล สพฺพธมฺมานํ, “สพฺเพ ธมฺมา อนตฺตา”ติ กุสโล สพฺพธมฺมานํ ฯเปฯ “อวิชฺชาปจฺจยา สงฺขารา”ติ กุสโล สพฺพธมฺมานํ. “ยํ กิญฺจิ สมุทย\-ธมฺมํ, สพฺพํ ตํ นิโรธธมฺมนฺ”ติ กุสโล สพฺพธมฺมานํ. เอวมฺปิ กุสโล สพฺพธมฺมานํ.}

\prose{อถ วา อนิจฺจโต กุสโล สพฺพธมฺมานํ. ทุกฺขโต. โรคโต. คณฺฑโต. สลฺลโต. อฆโต. อาพาธโต. ปรโต. ปโลกโต. ¢ติโต. อุปทฺทวโต. ภยโต. อุปสคฺคโต. จลโต. ปภงฺคุโต. อทฺธุวโต\csromansharedfootnote{14053ce509a6afcb}{อธุวโต (ก) ขุ 7. 40 ปิฏฺเฐปิ.}. อตาณโต. อเลณโต. อสรณโต. อสรณีภูตโต. ริตฺตโต. ตุจฺฉโต. สุญฺญโต. อนตฺตโต. อาทีนวโต. วิปริณาม\-ธมฺมโต. อสารกโต. อฆมูลโต. วธกโต. วิภวโต. สาสวโต. สงฺขตโต. มารามิสโต. ชาติธมฺมโต. ชราธมฺมโต. พฺยาธิธมฺมโต. มรณธมฺมโต. โสก\-ปริเทว\-ทุกฺข\-โทมนสฺสุ\-ปายาส\-ธมฺมโต. สํกิเลสิก\-ธมฺมโต. สมุทยโต. อตฺถงฺคมโต. อสฺสาทโต. อาทีนวโต. นิสฺสรณโต กุสโล สพฺพธมฺมานํ. เอวมฺปิ กุสโล สพฺพธมฺมานํ.}

\prose{อถ วา ขนฺธกุสโล ธาตุกุสโล อายตนกุสโล ปฏิจฺจสมุปฺปาท\-กุสโล สติปฏฺฐาน\-กุสโล สมฺมปฺปธาน\-กุสโล อิทฺธิปาท\-กุสโล อินฺทฺริยกุสโล พลกุสโล โพชฺฌงฺค\-กุสโล มคฺคกุสโล ผลกุสโล นิพฺพานกุสโล. เอวมฺปิ กุสโล สพฺพธมฺมานํ.}

\prose{อถ วา \textbf{สพฺพธมฺมา} วุจฺจนฺติ ทฺวาทสา\-ยตนานิ. จกฺขุ เจว\csromansharedfootnote{6473e149c4db9951}{จกฺขุญฺเจว (ก)} รูปา จ, โสตญฺจ สทฺทา จ, ฆานญฺจ คนฺธา จ, ชิวฺหา จ รสา จ, กาโย จ โผฏฺฐพฺพา จ, มโน จ ธมฺมา จ. ยโต จ อชฺฌตฺติก\-พาหิเรสุ อายตเนสุ ฉนฺทราโค ปหีโน โหติ อุจฺฉินฺนมูโล ตาลาวตฺถุกโต อนภาวํกโต\csromansharedfootnote{001aaeb2b57ed59e}{อนภาวงฺคโต (สฺยา)} อายติํ อนุปฺปาทธมฺโม, เอตฺตาวตาปิ กุสโล สพฺพ\-ธมฺมานนฺติ กุสโล สพฺพธมฺมานํ.}

\prose{\textbf{สโต ภิกฺขุ ปริพฺพเช}ติ \textbf{สโต}ติ จตูหิ การเณหิ สโต, กาเย กายานุปสฺสนา\-สติปฏฺฐานํ ภาเวนฺโต สโต, เวทนาสุ \csromanfolio{39}\textbf{เวทนานุปสฺสนา}\-\textbf{สติปฏฺฐานํ ภเวนฺโต สโต, จิตฺเต} \textbf{จิตฺตานุปสฺสนา}\-\textbf{สติปฏฺฐานํ ภาเวนฺโต สโต, ธมฺเมสุ} ธมฺมานุปสฺสนา\-สติปฏฺฐานํ ภาเวนฺโต สโต.}

\prose{อปเรหิปิ จตูหิ การเณหิ สโต, อสติ\-ปริวชฺชนาย สโต, สติกรณียานํ ธมฺมานํ กตตฺตา สโต, สติปริพนฺธานํ\csromansharedfootnote{b1a36877d01cb359}{สติปฏิปกฺขานํ (สฺยา) ขุ 7. 7 ปิฏฺเฐ.} ธมฺมานํ หตตฺตา สโต, สตินิมิตฺตานํ ธมฺมานํ อสมฺมุฏฺฐตฺตา\csromansharedfootnote{417198b5491f26ca}{อปฺปมุฏฺฐตฺตา (สฺยา)} สโต.}

\prose{อปเรหิปิ จตูหิ การเณหิ สโต, สติยา สมนฺนาคตตฺตา สโต, สติยา วสิตตฺตา สโต, สติยา ปาคุญฺเญน สมนฺนาคตตฺตา สโต, สติยา อปจฺโจโรหณตาย สโต.}

\prose{อปเรหิปิ จตูหิ การเณหิ สโต, สติยา สมนฺนาคตตฺตา สโต, สนฺตตฺตา สโต, สมิตตฺตา สโต, สนฺต\-ธมฺม\-สมนฺนาคตตฺตา สโต. พุทฺธา\-นุสฺสติยา สโต. ธมฺมา\-นุสฺสติยา สโต. สํฆานุสฺสติยา สโต. สีลานุสฺสติยา สโต. จาคานุสฺสติยา สโต. เทวตา\-นุสฺสติยา สโต. อานาปานสฺสติยา สโต. มรณสฺสติยา สโต. กายคตาสติยา สโต. อุปสมา\-นุสฺสติยา สโต. ยา สติ อนุสฺสติ ฯเปฯ สมฺมาสติ สติสมฺโพชฺฌงฺโค เอกายนมคฺโค. อยํ วุจฺจติ สติ. อิมาย สติยา อุเปโต โหติ สมุเปโต อุปาคโต สมุปาคโต อุปปนฺโน สมุปปนฺโน\csromansharedfootnote{2235c5316b755fa2}{สมฺปนฺโน (ก)} สมนฺนาคโต, โส วุจฺจติ สโต. \textbf{ภิกฺขู}ติ สตฺตนฺนํ ธมฺมานํ ภินฺนตฺตา ภิกฺขุ. สกฺกายทิฏฺฐิ ภินฺนา โหติ. วิจิกิจฺฉา ภินฺนา โหติ. สีลพฺพต\-ปรามาโส ภินฺโน โหติ. ราโค ภินฺโน โหติ. โทโส ภินฺโน โหติ. โมโห ภินฺโน โหติ. มาโน ภินฺโน โหติ. ภินฺนา โหนฺติ ปาปกา อกุสลา ธมฺมา สํกิเลสิกา โปโนภวิกา\csromansharedfootnote{86d675b178881c1d}{โปโนพฺภวิกา (สฺยา, ก)} สทรา ทุกฺขวิปากา อายติํ ชาติ\-ชรา\-มรณิยา.}

\prose{ปชฺเชน กเตน\csromansharedfootnote{99012859bc6894f4}{ปชฺโชตกเตน (ก) ขุ 1. 357 ปิฏฺเฐ.} อตฺตนา, (สภิยาติ ภควา,)}

\prose{ปรินิพฺพาน\-คโต วิติณฺณกงฺโข.}

\setlength{\csromangathapagewidth}{0pt}
\setlength{\csromangathawakwidth}{0pt}
\csromangathameasuregroup{}{วิภวญฺจ ภวญฺจ วิปฺปหาย, \\ วุสิตวา ขีณปุนพฺภโว ส ภิกฺขูติ.}
\csromangathameasurewak{วิภวญฺจ ภวญฺจ วิปฺปหาย, \\ วุสิตวา ขีณปุนพฺภโว ส ภิกฺขูติ.}
\setlength{\csromangathaleftcol}{0pt}
\csromangathagroup{\csromangathastanza{\csromangathawak{วิภวญฺจ ภวญฺจ วิปฺปหาย, \\ วุสิตวา ขีณปุนพฺภโว ส ภิกฺขูติ.}}}

\csromanmatikaanchor{mtk.23}
\csromantocmarkref{subsection}{2. ติสฺสเมตฺเตยฺยมาณวปุจฺฉานิทฺเทส}{mtk.23}
\bookmark[dest={mtk.23},level=2]{2. ติสฺสเมตฺเตยฺยมาณวปุจฺฉานิทฺเทส}
\prose{\csromanfolio{40}\textbf{สโต ภิกฺขุ ปริพฺพเช}ติ สโต ภิกฺขุ ปริพฺพเช, สโต คจฺเฉยฺย. สโต ติฏฺเฐยฺย. สโต นิสีเทยฺย. สโต เสยฺยํ กปฺเปยฺย. สโต อภิกฺกเมยฺย. สโต ปฏิกฺกเมยฺย. สโต อาโลเกยฺย. สโต วิโลเกยฺย. สโต สมิญฺเชยฺย. สโต ปสาเรยฺย. สโต สงฺฆาฏิ\-ปตฺต\-จีวรํ ธาเรยฺย. สโต จเรยฺย วิหเรยฺย อิริเยยฺย วตฺเตยฺย ปาเลยฺย ยเปยฺย ยาเปยฺยาติ สโต ภิกฺขุ ปริพฺพเช. เตนาห ภควา\csromandash{}}

\setlength{\csromangathapagewidth}{0pt}
\setlength{\csromangathawakwidth}{0pt}
\csromangathameasuregroup{กาเมสุ นาภิคิชฺเฌยฺย, \\ กุสโล สพฺพธมฺมานํ,}{\csromangathabat{กาเมสุ นาภิคิชฺเฌยฺย,}{มนสา’นาวิโล สิยา.} \\ \csromangathabat{กุสโล สพฺพธมฺมานํ,}{สโต ภิกฺขุ ปริพฺพเชติ.}}
\csromangathasetleft{กาเมสุ นาภิคิชฺเฌยฺย, \\ กุสโล สพฺพธมฺมานํ,}
\csromangathagroup{\csromangathastanza{\csromangathabat{กาเมสุ นาภิคิชฺเฌยฺย,}{มนสา’นาวิโล สิยา.} \\ \csromangathabat{กุสโล สพฺพธมฺมานํ,}{สโต ภิกฺขุ ปริพฺพเชติ.}}}

\prose{สห คาถา\-ปริโยสานา เย เต พฺราหฺมเณน สทฺธิํ เอกจฺฉนฺทา เอกปโยคา เอกาธิปฺปายา เอก\-วาสน\-วาสิตา. เตสํ อเนก\-ปาณสหสฺสานํ วิรชํ วีตมลํ ธมฺมจกฺขุํ อุทปาทิ “ยํ กิญฺจิ สมุทย\-ธมฺมํ, สพฺพํ ตํ นิโรธธมฺมนฺ”ติ. ตสฺส พฺราหฺมณสฺส อนุปาทาย อาสเวหิ จิตฺตํ วิมุจฺจิ. สห อรหตฺตปฺปตฺตา อชิน\-ชฏา\-วากจีร\-ติทณฺฑ\-กมณฺฑลุ\-เกสา จ มสฺสู จ อนฺตรหิตา, ภณฺฑุ\-กาสายวตฺถวสโน สงฺฆาฏิปตฺตจีวร\-ธโร อนฺวตฺถ\-ปฏิปตฺติยา ปญฺชลิโก ภควนฺตํ นมสฺสมาโน นิสินฺโน โหติ “สตฺถา เม ภนฺเต ภควา, สาวโกหมสฺมี”ติ.}

\vspace{\tipitakaparskip}
\csromancenter{อชิตมาณวปุจฺฉา\-นิทฺเทโส ปฐโม.}
\proseitem{9}{\textbf{โกธ สนฺตุสิโต โลเก, (อิจฺจายสฺมา ติสฺสเมตฺเตยฺโย,)}}

\prose{\textbf{กสฺส โน สนฺติ อิญฺชิตา.}}

\setlength{\csromangathapagewidth}{0pt}
\setlength{\csromangathawakwidth}{0pt}
\csromangathameasuregroup{}{\textbf{โก อุภนฺตม}\textbf{ภิญฺญาย,} \\ \textbf{มชฺเฌ มนฺตา น ลิปฺปติ.} \\ \textbf{กํ พฺรูสิ มหาปุริโสติ,} \\ \textbf{โก อิธ สิพฺพินิมจฺจคา.}}
\csromangathameasurewak{\textbf{โก อุภนฺตม}\textbf{ภิญฺญาย,} \\ \textbf{มชฺเฌ มนฺตา น ลิปฺปติ.} \\ \textbf{กํ พฺรูสิ มหาปุริโสติ,} \\ \textbf{โก อิธ สิพฺพินิมจฺจคา.}}
\setlength{\csromangathaleftcol}{0pt}
\csromangathagroup{\csromangathastanza{\csromangathawak{\textbf{โก อุภนฺตม}\-\textbf{ภิญฺญาย,} \\ \textbf{มชฺเฌ มนฺตา น ลิปฺปติ.} \\ \textbf{กํ พฺรูสิ มหาปุริโสติ,} \\ \textbf{โก อิธ สิพฺพินิมจฺจคา.}}}}

\prose{\textbf{โกธ สนฺตุสิโต โลเก}ติ โก โลเก ตุฏฺโฐ สนฺตุฏฺโฐ อตฺตมโน ปริปุณฺณ\-สงฺกปฺโปติ โกธ สนฺตุสิโต โลเก.}

\prose{\csromanfolio{41}\textbf{อิจฺจายสฺมา ติสฺส}\-\textbf{เมตฺเตยฺโย}ติ \textbf{อิจฺจา}ติ ปทสนฺธิ ปทสํสคฺโค ปทปาริปูรี อกฺขร\-สมวาโย พฺยญฺชน\-สิลิฏฺฐตา ปทา\-นุปุพฺพตา\-เปตํ อิจฺจาติ. \textbf{อายสฺมา}ติ ปิยจนํ ครุวจนํ สคารวสปฺปติสฺสาธิวจนเมตํ อายสฺมาติ. \textbf{ติสฺโส}ติ ตสฺส พฺราหฺมณสฺส นามํ สงฺขา สมญฺญา ปญฺญตฺติ โวหาโร นามํ นามกมฺมํ นามเธยฺยํ นิรุตฺติ พฺยญฺชนํ อภิลาโป. \textbf{เมตฺเตยฺโย}ติ ตสฺส พฺราหฺมณสฺส โคตฺตํ สงฺขา สมญฺญา ปญฺญตฺติ โวหาโรติ อิจฺจายสฺมา ติสฺสเมตฺเตยฺโย.}

\prose{\textbf{กสฺส โน สนฺติ อิญฺชิตา}ติ ตณฺหิญฺชิตํ ทิฏฺฐิญฺชิตํ มานิญฺชิตํ กิเลสิญฺชิตํ กามิญฺชิตํ. กสฺสิเม อิญฺชิตา นตฺถิ น สนฺติ น สํวิชฺชนฺติ นุปลพฺภนฺติ, ปหีนา สมุจฺฉินฺนา วูปสนฺตา ปฏิปสฺสทฺธา อภพฺพุ\-ปฺปตฺติกา ญาณคฺคินา ทฑฺฒาติ กสฺส โน สนฺติ อิญฺชิตา.}

\prose{\textbf{โก อุภนฺตม}\-\textbf{ภิญฺญายา}ติ โก อุโภ อนฺเต อภิญฺญาย ชานิตฺวา ตุลยิตฺวา ตีรยิตฺวา วิภาวยิตฺวา วิภูตํ กตฺวาติ โก อุภนฺตม\-ภิญฺญาย.}

\prose{\textbf{มชฺเฌ มนฺตา น ลิปฺปตี}ติ มชฺเฌ มนฺตาย น ลิปฺปติ, อลิตฺโต อนุปลิตฺโต นิกฺขนฺโต นิสฺสโฏ วิปฺปมุตฺโต วิสญฺญุตฺโต วิมริยาทิกเตน เจตสา วิหรตีติ มชฺเฌ มนฺตา น ลิปฺปติ.}

\prose{\textbf{กํ พฺรูสิ มหาปุริโส}ติ มหาปุริโส อคฺคปุริโส เสฏฺฐปุริโส วิเสฏฺฐปุริโส ปาโมกฺขปุริโส อุตฺตมปุริโส ปธานปุริโส ปวรปุริโสติ กํ พฺรูสิ กํ กเถสิ กํ มญฺญสิ กํ ภณสิ กํ ปสฺสสิ กํ โวหรสีติ กํ พฺรูสิ มหาปุริโสติ.}

\prose{\textbf{โก อิธ สิพฺพินิม}\-\textbf{จฺจคา}ติ โก อิธ สิพฺพินิํ ตณฺหํ อชฺฌคา อุปจฺจคา อติกฺกนฺโต สมติกฺกนฺโต วีติวตฺโตติ โก อิธ สิพฺพินิมจฺจคา. เตนาห โส พฺราหฺมโณ\csromandash{}}

\vspace{\tipitakaparskip}
\csromancenter{โกธ สนฺตุสิโต โลเก, (อิจฺจายสฺมา ติสฺสเมตฺเตยฺโย,)}
\prose{กสฺส โน สนฺติ อิญฺชิตา.}

\setlength{\csromangathapagewidth}{0pt}
\setlength{\csromangathawakwidth}{0pt}
\csromangathameasuregroup{}{โก อุภนฺตมภิญฺญาย, \\ มชฺเฌ มนฺตา น ลิปฺปติ. \\ กํ พฺรูสิ มหาปุริโสติ, \\ โก อิธ สิพฺพินิมจฺจคาติ.}
\csromangathameasurewak{โก อุภนฺตมภิญฺญาย, \\ มชฺเฌ มนฺตา น ลิปฺปติ. \\ กํ พฺรูสิ มหาปุริโสติ, \\ โก อิธ สิพฺพินิมจฺจคาติ.}
\setlength{\csromangathaleftcol}{0pt}
\csromangathagroup{\csromangathastanza{\csromangathawak{โก อุภนฺตม\-ภิญฺญาย, \\ มชฺเฌ มนฺตา น ลิปฺปติ. \\ กํ พฺรูสิ มหาปุริโสติ, \\ โก อิธ สิพฺพินิม\-จฺจคาติ.}}}

\csromanhangingbegin
\hangingheaditem{10}{\csromanfolio{42}\textbf{กาเมสุ พฺรหฺมจริยวา, (เมตฺเตยฺยาติ ภควา,)}}
\hangingbody{\textbf{วีตตโณฺห สทา สโต.}}
\hangingbody{\textbf{สงฺขาย นิพฺพุโต ภิกฺขุ,}}
\hangingbody{\textbf{ตสฺส โน สนฺติ อิญฺชิตา.}}
\csromanhangingend

\prose{\textbf{กาเมสุ พฺรหฺมจริยวา}ติ กามาติ อุทฺทานโต เทฺว \textbf{กามา} วตฺถุกามา จ กิเลสกามา จ ฯเปฯ. อิเม วุจฺจนฺติ วตฺถุกามา ฯเปฯ. อิเม วุจฺจนฺติ กิเลสกามา. \textbf{พฺรหฺมจริยํ} วุจฺจติ อสทฺธมฺม\-สมาปตฺติยา อารติ วิรติ ปฏิวิรติ เวรมณี อกิริยา อกรณํ อนชฺฌาปตฺติ เวลาอนติกฺกโม. อปิ จ นิปฺปริยาเยน พฺรหฺมจริยํ วุจฺจติ อริโย อฏฺฐงฺคิโก มคฺโค. เสยฺยถิทํ, สมฺมาทิฏฺฐิ สมฺมาสงฺกปฺโป สมฺมาวาจา สมฺมากมฺมนฺโต สมฺมาอาชีโว สมฺมาวายาโม สมฺมาสติ สมฺมาสมาธิ. โย อิมินา อริเยน อฏฺฐงฺคิเกน มคฺเคน อุเปโต สมุเปโต อุปาคโต สมุปาคโต อุปปนฺโน สมุปปนฺโน สมนฺนาคโต. โส วุจฺจติ พฺรหฺมจริยวา, ยถา จ ธเนน ธนวาติ วุจฺจติ. โภเคน โภควาติ วุจฺจติ. ยเสน ยสวาติ วุจฺจติ. สิปฺเปน สิปฺปวาติ วุจฺจติ. สีเลน สีลวาติ วุจฺจติ. วีริเยน วีริยวาติ วุจฺจติ. ปญฺญาย ปญฺญวาติ วุจฺจติ. วิชฺชาย วิชฺชวาติ วุจฺจติ. เอวเมว โย อิมินา อริเยน อฏฺฐงฺคิเกน มคฺเคน อุเปโต สมุเปโต อุปาคโต สมุปาคโต อุปปนฺโน สมุปปนฺโน สมนฺนาคโต. โส วุจฺจติ พฺรหฺมจริยวาติ กาเมสุ พฺรหฺมจริยวา.}

\prose{\textbf{เมตฺเตยฺยา}ติ ภควา ตํ พฺราหฺมณํ โคตฺเตน อาลปติ. \textbf{ภควา}ติ คารวาธิวจนเมตํ ฯเปฯ สจฺฉิกา ปญฺญตฺติ, ยทิทํ ภควาติ เมตฺเตยฺยาติ ภควา.}

\prose{\textbf{วีตตโณฺห สทา สโต}ติ ตณฺหาติ รูป\textbf{ตณฺหา} ฯเปฯ ธมฺมตณฺหา, ยสฺเสสา ตณฺหา ปหีนา สมุจฺฉินฺนา วูปสนฺตา ปฏิปสฺสทฺธา อภพฺพุ\-ปฺปตฺติกา ญาณคฺคินา ทฑฺฒา. โส วุจฺจติ วีตตโณฺห จตฺตตโณฺห วนฺตตโณฺห มุตฺตตโณฺห ปหีนตโณฺห ปฏินิสฺสฏฺฐ\-ตโณฺห วีตราโค จตฺตราโค วนฺตราโค มุตฺตราโค ปหีนราโค ปฏินิสฺสฏฺฐ\-ราโค นิจฺฉาโต นิพฺพุโต สีติภูโต สุขปฺปฏิสํเวที พฺรหฺมภูเตน อตฺตนา วิหรติ. \textbf{สทา}ติ สทา สพฺพทา สพฺพกาลํ นิจฺจกาลํ ธุวกาลํ สตตํ สมิตํ อพฺโพกิณฺณํ \csromanfolio{43}โปงฺขานุโปงฺขํ\csromansharedfootnote{2f5ab1c0c98f6912}{โปขานุโปขํ (สฺยา)} อุทกูมิกชาตํ อวีจิ สนฺตติ สหิตํ\csromansharedfootnote{70457654c86d17ad}{อวีจิ สมงฺคิสหิตํ (สฺยา) ขุ 7. 14 ปิฏฺเฐปิ.} ผสฺสิตํ\csromansharedfootnote{7b8e2d6a3ad9b7c7}{ผุสิตํ (สฺยา)} ปุเรภตฺตํ ปจฺฉาภตฺตํ ปุริมยามํ มชฺฌิมยามํ ปจฺฉิมยามํ กาเฬ ชุเณฺห วสฺเส เหมนฺเต คิเมฺห ปุริเม วโยขนฺเธ มชฺฌิเม วโยขนฺเธ ปจฺฉิเม วโยขนฺเธ. \textbf{สโต}ติ จตูหิ การเณหิ สโต, กาเย กายานุปสฺสนา\-สติปฏฺฐานํ ภาเวนฺโต สโต, เวทนาสุ เวทนานุปสฺสนา\-สติปฏฺฐานํ ภาเวนฺโต สโต, จิตฺเต จิตฺตานุปสฺสนา\-สติปฏฺฐานํ ภาเวนฺโต สโต, ธมฺเมสุ ธมฺมานุปสฺสนา\-สติปฏฺฐานํ ภาเวนฺโต สโต ฯเปฯ. โส วุจฺจติ สโตติ วีตตโณฺห สทา สโต.}

\prose{\textbf{สงฺขาย นิพฺพุโต ภิกฺขู}ติ สงฺขา วุจฺจติ ญาณํ. ยา ปญฺญา ปชานนา วิจโย ปวิจโย ฯเปฯ อโมโห ธมฺมวิจโย สมฺมาทิฏฺฐิ. \textbf{สงฺขายา}ติ สงฺขาย ชานิตฺวา ตุลยิตฺวา ตีรยิตฺวา วิภาวยิตฺวา วิภูตํ กตฺวา, “สพฺเพ สงฺขารา อนิจฺจา”ติ สงฺขาย ชานิตฺวา ตุ ลยิตฺวา ตีรยิตฺวา วิภาวยิตฺวา วิภูตํ กตฺวา, “สพฺเพ สงฺขารา ทุกฺขา”ติ. “สพฺเพ ธมฺมา อนตฺตา”ติ ฯเปฯ “อวิชฺชาปจฺจยา สงฺขารา”ติ. “ยํ กิญฺจิ สมุทย\-ธมฺมํ, สพฺพํ ตํ นิโรธธมฺมนฺ”ติ สงฺขาย ชานิตฺวา ตุลยิตฺวา ตีรยิตฺวา วิภาวยิตฺวา วิภูตํ กตฺวา.}

\prose{อถ วา อนิจฺจโต สงฺขาย ชานิตฺวา. ทุกฺขโต. โรคโต. คณฺฑโต. สลฺลโต ฯเปฯ. นิสฺสรณโต สงฺขาย ชานิตฺวา ตุลยิตฺวา ตีรยิตฺวา วิภาวยิตฺวา วิภูตํ กตฺวา. \textbf{นิพฺพุโต}ติ ราคสฺส นิพฺพาปิตตฺตา นิพฺพุโต. โทสสฺส นิพฺพาปิตตฺตา นิพฺพุโต. โมหสฺส นิพฺพาปิตตฺตา นิพฺพุโต. โกธสฺส. อุปนาหสฺส. มกฺขสฺส. ปฬาสสฺส. อิสฺสาย. มจฺฉริยสฺส. มายาย. สาเฐยฺยสฺส. ถมฺภสฺส. สารมฺภสฺส. มานสฺส. อติมานสฺส. มทสฺส. ปมาทสฺส. สพฺพกิเลสานํ. สพฺพ\-ทุจฺจริตานํ. สพฺพ\-ทรถานํ. สพฺพ\-ปริฬาหานํ. สพฺพ\-สนฺตาปานํ. สพฺพา\-กุสลา\-ภิสงฺขารานํ นิพฺพาปิตตฺตา นิพฺพุโต. \textbf{ภิกฺขู}ติ สตฺตนฺนํ ธมฺมานํ ภินฺนตฺตา ภิกฺขุ ฯเปฯ วุสิตวา ขีณปุนพฺภโว ส ภิกฺขูติ สงฺขาย นิพฺพุโต ภิกฺขุ.}

\prose{\textbf{ตสฺส โน สนฺติ อิญฺชิตา}ติ \textbf{ตสฺสา}ติ อรหโต ขีณาสวสฺส. \textbf{อิญฺชิตา}ติ ตณฺหิญฺชิตํ ทิฏฺฐิญฺชิตํ มานิญฺชิตฺตํ กิเลสิญฺชิตํ กามิญฺชิตํ. ตสฺสิเม \csromanfolio{44}อิญฺชิตา นตฺถิ น สนฺติ น สํวิชฺชนฺติ นุปลพฺภนฺติ, ปหีนา สมุจฺฉินฺนา วูปสนฺตา ปฏิปสฺสทฺธา อภพฺพุ\-ปฺปตฺติกา ญาณคฺคินา ทฑฺฒาติ ตสฺส โน สนฺติ อิญฺชิตา. เตนาห ภควา\csromandash{}}

\proseitem{2}{กาเมสุ พฺรหฺมจริยวา, (เมตฺเตยฺยาติ ภควา,)}

\prose{วีตตโณฺห สทา สโต.}

\setlength{\csromangathapagewidth}{0pt}
\setlength{\csromangathawakwidth}{0pt}
\csromangathameasuregroup{\textbf{โส อุภนฺตม}\textbf{ภิญฺญา}ย, \\ \textbf{ตํ พฺรูมิ มหาปุริ}โสติ,}{สงฺขาย นิพฺพุโต ภิกฺขุ, \\ ตสฺส โน สนฺติ อิญฺชิตาติ. \\ \csromangathabat{\textbf{โส อุภนฺตม}\textbf{ภิญฺญา}ย,}{มชฺเฌ มฺ\textbf{อนฺตา} น ลิปฺปติ.} \\ \csromangathabat{\textbf{ตํ พฺรูมิ มหาปุริ}โสติ,}{โส อิธ สิพฺพินิมจฺจคา.}}
\csromangathameasurewak{สงฺขาย นิพฺพุโต ภิกฺขุ, \\ ตสฺส โน สนฺติ อิญฺชิตาติ.}
\csromangathasetleft{\textbf{โส อุภนฺตม}\textbf{ภิญฺญา}ย, \\ \textbf{ตํ พฺรูมิ มหาปุริ}โสติ,}
\csromangathagroup{\csromangathastanza{\csromangathawak{สงฺขาย นิพฺพุโต ภิกฺขุ, \\ ตสฺส โน สนฺติ อิญฺชิตาติ.}}\csromangathastanza{\csromangathaitembat{2}{\textbf{โส อุภนฺตม}\-\textbf{ภิญฺญา}ย,}{มชฺเฌ มฺ\textbf{อนฺตา} น ลิปฺปติ.} \\ \csromangathacontbat{\textbf{ตํ พฺรูมิ มหาปุริ}โสติ,}{โส อิธ สิพฺพินิมจฺจคา.}}}

\prose{\textbf{โส อุภนฺตม}\-\textbf{ภิญฺญาย, มชฺเฌ มนฺตา น ลิปฺปตี}ติ \textbf{อนฺตา}ติ ผสฺโส เอโก อนฺโต, ผสฺสสมุทโย ทุติโย อนฺโต, ผสฺสนิโรโธ มชฺเฌ. อตีตํ เอโก อนฺโต, อนาคตํ ทุติโย อนฺโต, ปจฺจุปฺปนฺนํ มชฺเฌ. สุขา เวทนา เอโก อนฺโต, ทุกฺขา เวทนา ทุติโย อนฺโต, อทุกฺขมสุขา เวทนา มชฺเฌ. นามํ เอโก อนฺโต, รูปํ ทุติโย อนฺโต, วิญฺญาณํ มชฺเฌ. ฉ อชฺฌตฺติกานิ อายตนานิ เอโก อนฺโต, ฉ พาหิรานิ อายตนานิ ทุติโย อนฺโต, วิญฺญาณํ มชฺเฌ. สกฺกาโย เอโก อนฺโต, สกฺกาย\-สมุทโย ทุติโย อนฺโต, สกฺกายนิโรโธ มชฺเฌ. \textbf{มนฺตา} วุจฺจติ ปญฺญา, ยา ปญฺญา ปชานนา ฯเปฯ อโมโห ธมฺมวิจโย สมฺมาทิฏฺฐิ.}

\prose{\textbf{เลปา}ติ เทฺว เลปา ตณฺหาเลโป จ ทิฏฺฐิเลโป จ. กตโม ตณฺหาเลโป. ยาวตา ตณฺหา\-สงฺขาเตน สีมกตํ โอธิกตํ\csromansharedfootnote{79c135138a8a1502}{มริยาทิกตํ โอธิกตํ (สฺยา)} ปริยนฺตกตํ ปริคฺคหิตํ มมายิตํ อิทํ มม, เอตํ มม, เอตฺตกํ มม, เอตฺตาวตา มม รูปา สทฺทา คนฺธา รสา โผฏฺฐพฺพา อตฺถรณา ปาวุรณา ทาสิทาสา อเชฬกา กุกฺกุฏสูกรา หตฺถิ\-ควา\-สฺส\-วฬวา เขตฺตํ วตฺถุ หิรญฺญํ สุวณฺณํ คามนิคม\-ราชธานิโย รฏฺฐญฺจ ชนปโท จ โกโส จ โกฏฺฐาคารญฺจ, เกวลมฺปิ มหาปถวิํ ตณฺหาวเสน มมายติ, ยาวตา อฏฺฐสต\-ตณฺหาวิจริตํ. อยํ ตณฺหาเลโป.}

\prose{\csromanfolio{45}กตโม ทิฏฺฐิเลโป. วีสติวตฺถุกา สกฺกายทิฏฺฐิ, ทสวตฺถุกา มิจฺฉาทิฏฺฐิ, ทสวตฺถุกา อนฺตคฺคาหิกา ทิฏฺฐิ. ยา เอวรูปา ทิฏฺฐิ ทิฏฺฐิคตํ ทิฏฺฐิคหนํ ทิฏฺฐิกนฺตาโร ทิฏฺฐิ\-วิสูกายิกํ ทิฏฺฐิ\-วิปฺผนฺทิตํ ทิฏฺฐิ\-สํโยชนํ คาโห ปฏิคฺคาโห อภินิเวโส ปรามาโส กุมฺมคฺโค มิจฺฉาปโถ มิจฺฉตฺตํ ติตฺถายตนํ วิปริยาส\-คฺคาโห\csromansharedfootnote{e959f2e8835affd3}{วิปริเยสคฺคาโห (พหูสุ)} วิปรีตคฺคาโห วิปลฺลาสคฺคาโห มิจฺฉาคาโห, อยาถาวกสฺมิํ ยาถาวกนฺติ คาโห, ยาวตา ทฺวาสฏฺฐิ ทิฏฺฐิคตานิ. อยํ ทิฏฺฐิเลโป.}

\prose{\textbf{โส อุภนฺตม}\-\textbf{ภิญฺญาย, มชฺเฌ มนฺตา น ลิปฺปตี}ติ โส อุโภ จ อนฺเต มชฺฌญฺจ มนฺตาย อภิญฺญาย ชานิตฺวา ตุลยิตฺวา ตีรยิตฺวา วิภาวยิตฺวา วิภูตํ กตฺวา น ลิปฺปติ น ปลิปฺปติ น อุปลิปฺปติ, อลิตฺโต อสํลิตฺโต อนุปลิตฺโต นิกฺขนฺโต นิสฺสโฏ วิปฺปมุตฺโต วิสญฺญุตฺโต วิมริยาทิกเตน เจตสา วิหรตีติ โส อุภนฺตม\-ภิญฺญาย, มชฺเฌ มนฺตา น ลิปฺปติ.}

\prose{\textbf{ตํ พฺรูมิ มหาปุริโสติ} มหาปุริโส อคฺคปุริโส เสฏฺฐปุริโส วิเสฏฺฐปุริโส ปาโมกฺขปุริโส อุตฺตมปุริโส ปวรปุริโส ตํ พฺรูมิ ตํ กเถมิ ตํ ภณามิ ตํ ทีเปมิ ตํ โวหรามิ.}

\prose{อายสฺมา สาริปุตฺโต\csromansharedfootnote{89cc2fbe69d8d7fe}{ปสฺส สํ 3. 137 ปิฏฺเฐ.} ภควนฺตํ เอตทโวจ “มหาปุริโส มหาปุริโส”ติ ภนฺเต วุจฺจติ. กิตฺตาวตา นุ โข ภนฺเต มหาปุริโส โหตีติ. วิมุตฺต\-จิตฺตตฺตา ขฺวาหํ สาริปุตฺต “มหาปุริโส”ติ วทามิ, อวิมุตฺต\-จิตฺตตฺตา โน “มหาปุริโส”ติ วทามิ.}

\prose{กถญฺจ สาริปุตฺต วิมุตฺตจิตฺโต โหติ. อิธ สาริปุตฺต ภิกฺขุ อชฺฌตฺตํ กาเย กายานุปสฺสี วิหรติ อาตาปี สมฺปชาโน สติมา วิเนยฺย โลเก อภิชฺฌา\-โทมนสฺสํ. ตสฺส กาเย กายานุปสฺสิโน วิหรโต จิตฺตํ วิรชฺชติ วิมุจฺจติ อนุปาทาย อาสเวหิ. เวทนาสุ. จิตฺเต. ธมฺเมสุ ธมฺมานุปสฺสี วิหรติ อาตาปี สมฺปชาโน สติมา วิเนยฺย โลเก อภิชฺฌา\-โทมนสฺสํ. ตสฺส ธมฺเมสุ ธมฺมา\-นุปสฺสิโน วิหรโต จิตฺตํ วิรชฺชติ วิมุจฺจติ อนุปาทาย อาสเวหิ. เอวํ โข สาริปุตฺต ภิกฺขุ วิมุตฺตจิตฺโต โหติ, วิมุตฺต\-จิตฺตตฺตา ขฺวาหํ สาริปุตฺต}

\csromanmatikaanchor{mtk.24}
\csromantocmarkref{subsection}{3. ปุณฺณกมาณวปุจฺฉานิทฺเทส}{mtk.24}
\bookmark[dest={mtk.24},level=2]{3. ปุณฺณกมาณวปุจฺฉานิทฺเทส}
\prose{\csromanfolio{46}“มหาปุริโส”ติ วทามิ. อวิมุตฺต\-จิตฺตตฺตา โน “มหาปุริโส”ติ วทามีติ ตํ พฺรูมิ มหาปุริโสติ.}

\prose{\textbf{โส อิธ สิพฺพินิม}\-\textbf{จฺจคา}ติ สิพฺพินี วุจฺจติ ตณฺหา, โย ราโค สาราโค ฯเปฯ อภิชฺฌา โลโภ อกุสลมูลํ. ยสฺเสสา สิพฺพินี ตณฺหา ปหีนา สมุจฺฉินฺนา วูปสนฺตา ปฏิปสฺสทฺธา อภพฺพุ\-ปฺปตฺติกา ญาณคฺคินา ทฑฺฒา. โส สิพฺพินิํ ตณฺหํ อจฺจคา อุปจฺจคา อติกฺกนฺโต สมติกฺกนฺโต วีติวตฺโตติ โส อิธ สิพฺพินิมจฺจคา. เตนาห ภควา\csromandash{}}

\setlength{\csromangathapagewidth}{0pt}
\setlength{\csromangathawakwidth}{0pt}
\csromangathameasuregroup{โส อุภนฺตมภิญฺญาย, \\ ตํ พฺรูมิ มหาปุริโสติ,}{\csromangathabat{โส อุภนฺตมภิญฺญาย,}{มชฺเฌ มนฺตา น ลิปฺปติ.} \\ \csromangathabat{ตํ พฺรูมิ มหาปุริโสติ,}{โส อิธ สิพฺพินิมจฺจคาติ.}}
\csromangathasetleft{โส อุภนฺตมภิญฺญาย, \\ ตํ พฺรูมิ มหาปุริโสติ,}
\csromangathagroup{\csromangathastanza{\csromangathabat{โส อุภนฺตม\-ภิญฺญาย,}{มชฺเฌ มนฺตา น ลิปฺปติ.} \\ \csromangathabat{ตํ พฺรูมิ มหาปุริโสติ,}{โส อิธ สิพฺพินิม\-จฺจคาติ.}}}

\prose{สห คาถา\-ปริโยสานา เย เต พฺราหฺมเณน สทฺธิํ เอกจฺฉนฺทา เอกปโยคา เอกาธิปฺปายา เอก\-วาสน\-วาสิตา, เตสํ อเนก\-ปาณสหสฺสานํ วิรชํ วีตมลํ ธมฺมจกฺขุํ อุทปาทิ “ยํ กิญฺจิ สมุทย\-ธมฺมํ, สพฺพํ ตํ นิโรธธมฺมนฺ”ติ. ตสฺส พฺราหฺมณสฺส อนุปาทาย อาสเวหิ จิตฺตํ วิมุจฺจิ. สห อรหตฺตปฺปตฺตา อชิน\-ชฏา\-วากจีร\-ติทณฺฑ\-กมณฺฑลุ\-เกสา จ มสฺสู จ อนฺตรหิตา. ภณฺฑุ\-กาสายวตฺถวสโน สงฺฆาฏิปตฺตจีวร\-ธโร อนฺวตฺถ\-ปฏิปตฺติยา ปญฺชลิโก ภควนฺตํ นมสฺสมาโน นิสินฺโน โหติ “สตฺถา เม ภนฺเต ภควา, สาวโกหมสฺมี”ติ.}

\vspace{\tipitakaparskip}
\csromancenter{ติสฺสเมตฺเตยฺยมาณวปุจฺฉา\-นิทฺเทโส ทุติโย.}
\proseitem{12}{\textbf{อเนชํ มูลทสฺสาวิํ, (อิจฺจายสฺมา ปุณฺณโก,)}}

\prose{\textbf{อตฺถิ ปเญฺหน อาคมํ.}}

\setlength{\csromangathapagewidth}{0pt}
\setlength{\csromangathawakwidth}{0pt}
\csromangathameasuregroup{}{\textbf{กิํ นิสฺสิตา อิสโย มนุชา,} \\ \textbf{ยญฺญม}\textbf{กปฺปยิํสุ ปุถุ’ธ โลเก,} \\ \textbf{ปุจฺฉามิ ตํ ภควา พฺรูหิ เมตํ.}}
\csromangathameasurewak{\textbf{กิํ นิสฺสิตา อิสโย มนุชา,} \\ \textbf{ยญฺญม}\textbf{กปฺปยิํสุ ปุถุ’ธ โลเก,} \\ \textbf{ปุจฺฉามิ ตํ ภควา พฺรูหิ เมตํ.}}
\setlength{\csromangathaleftcol}{0pt}
\csromangathagroup{\csromangathastanza{\csromangathawak{\textbf{กิํ นิสฺสิตา อิสโย มนุชา,} \\ \textbf{ยญฺญม}\-\textbf{กปฺปยิํสุ ปุถุ’ธ โลเก,} \\ \textbf{ปุจฺฉามิ ตํ ภควา พฺรูหิ เมตํ.}}}}

\prose{\textbf{อเนชํ มูลทสฺสาวินฺ}ติ \textbf{เอชา} วุจฺจติ ตณฺหา, โย ราโค สาราโค ฯเปฯ อภิชฺฌา โลโภ อกุสลมูลํ. สา เอชา ตณฺหา พุทฺธสฺส ภควโต \csromanfolio{47}ปหีนา อุจฺฉินฺนมูลา ตาลาวตฺถุกตา อนภาวํกตา อายติํ อนุปฺปาทธมฺมา. ตสฺมา พุทฺโธ อเนโช, เอชาย ปหีนตฺตา อเนโช. ภควา ลาเภปิ น อิญฺชติ, อลาเภปิ น อิญฺชติ, ยเสปิ น อิญฺชติ, อยเสปิ น อิญฺชติ, ปสํสายปิ น อิญฺชติ, นินฺทายปิ น อิญฺชติ, สุเขปิ น อิญฺชติ, ทุกฺเขปิ น อิญฺชติ น จลติ น เวธติ นปฺปเวธตีติ อเนชํ. \textbf{มูลทสฺสาวินฺ}ติ ภควา มูลทสฺสาวี เหตุทสฺสาวี นิทานทสฺสาวี สมฺภว\-ทสฺสาวี ปภวทสฺสาวี สมุฏฺฐาน\-ทสฺสาวี อาหารทสฺสาวี อารมฺมณ\-ทสฺสาวี ปจฺจยทสฺสาวี สมุทย\-ทสฺสาวี.}

\prose{ตีณิ อกุสลมูลานิ, โลโภ อกุสลมูลํ, โทโส อกุสลมูลํ, โมโห อกุสลมูลํ.}

\prose{วุตฺตํ เหตํ ภควตา\csromandash{}\csromansharedfootnote{e945739d162f58fb}{ปสฺส อํ 1. 266 ปิฏฺเฐ.} ตีณิมานิ ภิกฺขเว นิทานานิ กมฺมานํ สมุทยาย. กตมานิ ตีณิ, โลโภ นิทานํ กมฺมานํ สมุทยาย, โทโส นิทานํ กมฺมานํ สมุทยาย, โมโห นิทานํ กมฺมานํ สมุทยาย. น ภิกฺขเว โลภเชน กมฺเมน โทสเชน กมฺเมน โมหเชน กมฺเมน เทวา ปญฺญายนฺติ, มนุสฺสา ปญฺญายนฺติ, ยา วา ปนญฺญาปิ กาจิ สุคติโย. อถ โข ภิกฺขเว โลภเชน กมฺเมน โทสเชน กมฺเมน โมหเชน กมฺเมน นิรโย ปญฺญายติ, ติรจฺฉานโยนิ ปญฺญายติ, เปตฺติวิสโย ปญฺญายติ. ยา วา ปนญฺญาปิ กาจิ ทุคฺคติโย นิรเย ติรจฺฉาน\-โยนิยา เปตฺติวิสเย อตฺตภาวา\-ภินิพฺพตฺติยา. อิมานิ ตีณิ อกุสลมูลานีติ ภควา ชานาติ ปสฺสติ. เอวมฺปิ ภควา มูลทสฺสาวี ฯเปฯ สมุทย\-ทสฺสาวี. ตีณิ กุสลมูลานิ, อโลโภ กุสลมูลํ, อโทโส กุสลมูลํ, อโมโห กุสลมูลํ.}

\prose{วุตฺตํ เหตํ ภควตา\csromandash{}ตีณิมานิ ฯเปฯ. น ภิกฺขเว อโลภเชน กมฺเมน อโทสเชน กมฺเมน อโมหเชน กมฺเมน นิรโย ปญฺญายติ, ติรจฺฉานโยนิ ปญฺญายติ, เปตฺติวิสโย ปญฺญายติ, ยา วา ปนญฺญาปิ กาจิ ทุคฺคติโย. อถ โข ภิกฺขเว อโลภเชน กมฺเมน อโทสเชน กมฺเมน อโมหเชน กมฺเมน เทวา ปญฺญายนฺติ, มนุสฺสา ปญฺญายนฺติ, ยา วา ปนญฺญาปิ กาจิ สุคติโย เทเว จ มนุสฺเส จ อตฺตภาวา\-ภินิพฺพตฺติยา. อิมานิ ตีณิ กุสลมูลานีติ ภควา ชานาติ ปสฺสติ. เอวมฺปิ ภควา มูลทสฺสาวี ฯเปฯ สมุทย\-ทสฺสาวี.}

\proseitem{3}{\csromanfolio{48}วุตฺตํ เหตํ ภควตา\csromandash{}เย เกจิ ภิกฺขเว ธมฺมา อกุสลา อกุสลภาคิยา อกุสล\-ปกฺขิกา, สพฺเพ เต อวิชฺชามูลกา อวิชฺชา\-สโมสรณา. อวิชฺชา สมุคฺฆาตา สพฺเพ เต สมุคฺฆาตํ คจฺฉนฺตีติ ภควา ชานาติ ปสฺสติ. เอวมฺปิ ภควา มูลทสฺสาวี ฯเปฯ สมุทย\-ทสฺสาวี.}

\prose{วุตฺตํ เหตํ ภควตา\csromandash{}เย เกจิ ภิกฺขเว ธมฺมา กุสลา กุสลภาคิยา กุสลปกฺขิกา, สพฺเพ เต อปฺปมาทมูลกา อปฺปมาท\-สโมสรณา. อปฺปมาโท เตสํ ธมฺมานํ อคฺคมกฺขายตีติ ภควา ชานาติ ปสฺสติ. เอวมฺปิ ภควา มูลทสฺสาวี ฯเปฯ สมุทย\-ทสฺสาวี.}

\prose{อถ วา ภควา ชานาติ ปสฺสติ “อวิชฺชา มูลํ สงฺขารานํ. สงฺขารา มูลํ วิญฺญาณสฺส. วิญฺญาณํ มูลํ นามรูปสฺส. นามรูปํ มูลํ สฬายตนสฺส. สฬายตนํ มูลํ ผสฺสสฺส. ผสฺโส มูลํ เวทนาย. เวทนา มูลํ ตณฺหาย. ตณฺหา มูลํ อุปาทานสฺส. อุปาทานํ มูลํ ภวสฺส. ภโว มูลํ ชาติยา. ชาติ มูลํ ชรามรณสฺสา”ติ ภควา ชานาติ ปสฺสติ. เอวมฺปิ ภควา มูลทสฺสาวี ฯเปฯ สมุทย\-ทสฺสาวี.}

\prose{อถ วา ภควา ชานาติ ปสฺสติ “จกฺขุ มูลํ จกฺขุโรคานํ. โสตํ มูลํ โสตโรคานํ. ฆานํ มูลํ ฆานโรคานํ. ชิวฺหา มูลํ ชิวฺหาโรคานํ. กาโย มูลํ กายโรคานํ. มโน มูลํ เจตสิกานํ ทุกฺขานนฺ”ติ ภควา ชานาติ ปสฺสติ. เอวมฺปิ ภควา มูลทสฺสาวี เหตุทสฺสาวี นิทานทสฺสาวี สมฺภว\-ทสฺสาวี ปภวทสฺสาวี สมุฏฺฐาน\-ทสฺสาวี อาหารทสฺสาวี อารมฺมณ\-ทสฺสาวี ปจฺจยทสฺสาวี สมุทย\-ทสฺสาวีติ อเนชํ มูลทสฺสาวี.}

\prose{\textbf{อิจฺจายสฺมา ปุณฺณโก}ติ \textbf{อิจฺจา}ติ ปทสนฺธิ \textbf{ฯเปฯ} อายสฺมา ปุณฺณโก.}

\prose{\textbf{อตฺถิ ปเญฺหน อาคมนฺ}ติ ปเญฺหน อตฺถิโก อาคโตมฺหิ\csromansharedfootnote{975d2e293ea57dd5}{ปญฺหตฺถิกามฺห อาคตา (พหูสุ) ปสฺส ขุ 7. 369 ปิฏฺเฐ.}, ปญฺหํ ปุจฺฉิตุกาโม อาคโตมฺหิ, ปญฺหํ โสตุกาโม อาคโตมฺหีติ เอวมฺปิ อตฺถิ ปเญฺหน อาคมํ. อถ วา ปญฺหตฺถิกานํ ปญฺหํ ปุจฺฉิตุ\-กามานํ ปญฺหํ โสตุกามานํ อาคมนํ อภิกฺกมนํ อุปสงฺกมนํ ปยิรุปาสนํ อตฺถีติ เอวมฺปิ อตฺถิ ปเญฺหน อาคมํ. อถ วา ปญฺหาคโม ตุยฺหํ อตฺถิ, ตฺวมฺปิ ปหุ, วิสวี อลมตฺโต, มยา ปุจฺฉิตํ กเถตุํ วิสชฺเชตุํ, วหสฺเสตํ ภารนฺติ\csromansharedfootnote{645fb03078f4fb64}{วิสชฺเชตุํ สนฺทสฺเสตุํ ภณิตุนฺติ (สฺยา) วหสฺสุ + เอตํ.} เอวมฺปิ อตฺถิ ปเญฺหน อาคมํ.}

\prose{\csromanfolio{49}\textbf{กิํ นิสฺสิตา อิสโย มนุชา}ติ กิํ นิสฺสิตา อาสิตา อลฺลีนา อุปคตา อชฺโฌสิตา อธิมุตฺตา. \textbf{อิสโย}ติ อิสินามกา เย เกจิ อิสิปพฺพชฺชํ ปพฺพชิตา อาชีวกา นิคณฺฐา ชฏิลา ตาปสา. \textbf{มนุชา}ติ มนุสฺสา วุจฺจนฺตีติ กิํ นิสฺสิตา อิสโย มนุชา.}

\prose{\textbf{ขตฺติยา พฺราหฺมณา เทวตานนฺ}ติ \textbf{ขตฺติยา}ติ เย เกจิ ขตฺติยชาติกา. \textbf{พฺราหฺมณา}ติ เย เกจิ โภวาทิกา. \textbf{เทวตานนฺ}ติ อาชีวก\-สาวกานํ อาชีวกา เทวตา. นิคณฺฐ\-สาวกานํ นิคณฺฐา เทวตา. ชฏิล\-สาวกานํ ชฏิลา เทวตา. ปริพฺพาชก\-สาวกานํ ปริพฺพาชกา เทวตา. อวิรุทฺธ\-กสา\-วกานํ อวิรุทฺธกา\csromansharedfootnote{d0d60f1408142ca9}{อวรุทฺธกสาวกานํ อวรุทฺธกา (สฺยา)} เทวตา. หตฺถิ\-วติกานํ หตฺถี เทวตา. อสฺสวติกานํ อสฺสา เทวตา. โควติกานํ คาโว เทวตา. กุกฺกุร\-วติกานํ กุกฺกุรา เทวตา. กากวติกานํ กากา เทวตา. วาสุเทว\-วติกานํ วาสุเทโว เทวตา. พลเทว\-วติกานํ พลเทโว เทวตา. ปุณฺณ\-ภทฺทวติกานํ ปุณฺณภทฺโท เทวตา. มณิภทฺทวติกานํ มณิภทฺโท เทวตา. อคฺคิวติกานํ อคฺคิ เทวตา. นาควติกานํ นาคา เทวตา. สุปณฺณ\-วติกานํ สุปณฺณา เทวตา. ยกฺข\-วติกานํ ยกฺขา เทวตา. อสุรวติกานํ อสุรา เทวตา. คนฺธพฺพ\-วติกานํ คนฺธพฺพา เทวตา. มหาราช\-วติกานํ มหาราชาโน เทวตา. จนฺทวติกานํ จนฺโท เทวตา. สูริย\-วติกานํ สูริโย เทวตา. อินฺทวติกานํ อินฺโท เทวตา. พฺรหฺม\-วติกานํ พฺรหฺมา เทวตา. เทววติกานํ เทโว เทวตา. ทิสาวติกานํ ทิสา เทวตา. เย เยสํ ทกฺขิเณยฺยา เต เตสํ เทวตาติ ขตฺติย\-พฺราหฺมณา เทวตานํ.}

\prose{\textbf{ยญฺญม}\-\textbf{กปฺปยิํสุ ปุถูธ โลเก}ติ \textbf{ยญฺญํ} วุจฺจติ เทยฺยธมฺโม จีวร\-ปิณฺฑปาต\-เสนาสน\-คิลานปจฺจย\-เภสชฺช\-ปริกฺขารํ อนฺนํ ปานํ วตฺถํ ยานํ มาลาคนฺธาวิเลปนํ\csromansharedfootnote{96b3c2424a152f25}{มาลาคนฺธํ วิเลปนํ (สฺยา) ขุ 1. 240 ปิฏฺเฐ.} เสยฺยา\-วสถ\-ปทีเปยฺยํ. \textbf{ยญฺญมกปฺปยิํสู}ติ เยปิ ยญฺญํ เอสนฺติ คเวสนฺติ ปริเยสนฺติ จีวร\-ปิณฺฑปาต\-เสนาสน\-คิลานปจฺจย\-เภสชฺช\-ปริกฺขารํ อนฺนํ ปานํ วตฺถํ ยานํ มาลา\-คนฺธ\-วิเลปนํ เสยฺยา\-วสถ\-ปทีเปยฺยํ. เตปิ ยญฺญํ กปฺเปนฺติ. เยปิ ยญฺญํ อภิสงฺขาโรนฺติ จีวร\-ปิณฺฑปาต\-เสนาสน\-คิลานปจฺจย\-เภสชฺช\-ปริกฺขารํ อนฺนํ ปานํ ฯเปฯ เสยฺยา\-วสถ\-ปทีเปยฺยํ. เตปิ ยญฺญํ กปฺเปนฺติ. เยปิ ยญฺญํ เทนฺติ ยชนฺติ ปริจฺจชนฺติ จีวร ปิณฺฑปาต เสนาสน คิลาน ปจฺจย เภสชฺช ปริกฺขารํ อนฺนํ ปานํ ฯเปฯ}

\prose{\csromanfolio{50}เสยฺยา\-วสถ\-ปทีเปยฺยํ. เตปิ ยญฺญํ กปฺเปนฺติ. \textbf{ปุถู}ติ ยญฺญา วา เอเต ปุถู. ยญฺญยาชกา\csromansharedfootnote{fb12630412daeaf7}{ยญฺญยชกา (สฺยา)} วา เอเต ปุถู. ทกฺขิเณยฺยา วา เอเต ปุถู. กถํ ยญฺญา วา เอเต ปุถู. พหุกานํ เอเต ยญฺญา จีวร\-ปิณฺฑปาต\-เสนาสน\-คิลานปจฺจย\-เภสชฺช\-ปริกฺขารา อนฺนํ ปานํ วตฺถํ ยานํ มาลํ คนฺธํ วิเลปนํ เสยฺยา\-วสถ\-ปทีเปยฺยํ. เอวํ ยญฺญา วา เอเต ปุถู.}

\prose{กถํ ยญฺญยาชกา วา เอเต ปุถู. พหุกา เอเต ยญฺญยาชกา ขตฺติยา จ พฺราหฺมณา จ เวสฺสา จ สุทฺทา จ คหฏฺฐา จ ปพฺพชิตา จ เทวา จ มนุสฺสา จ. เอวํ ยญฺญยาชกา วา เอเต ปุถู.}

\prose{กถํ ทกฺขิเณยฺยา วา เอเต ปุถู. พหุกา เอเต ทกฺขิเณยฺยา ปุถู สมณพฺราหฺมณา กปณทฺธิก\-วนิพฺพกยาจกา\csromansharedfootnote{c490750efcd163b8}{...วณิพฺพกสาวกา (สฺยา) ขุ 1. 240 ปิฏฺเฐ.}. เอวํ ทกฺขิเณยฺยา วา เอเต ปุถู.  อิธ โลเกติ มนุสฺสโลเกติ ยญฺญม\-กปฺปยิํสุ ปุถูธ โลเก.}

\prose{\textbf{ปุจฺฉามิ ตํ ภควา พฺรูหิ เมตนฺ}ติ ปุจฺฉาติ ติสฺโส \textbf{ปุจฺฉา} อทิฏฺฐโชตนา ปุจฺฉา ทิฏฺฐ\-สํสนฺทนา ปุจฺฉา วิมติจฺเฉทนา ปุจฺฉา. กตมา อทิฏฺฐโชตนา ปุจฺฉา. ปกติยา ลกฺขณํ อญฺญาตํ โหติ อทิฏฺฐํ อตุลิตํ อตีริตํ อวิภูตํ อวิภาวิตํ, ตสฺส ญาณาย ทสฺสนาย ตุลนาย ตีรณาย วิภูตตฺถาย วิภาวน\-ตฺถาย ปญฺหํ ปุจฺฉติ. อยํ อทิฏฺฐโชตนา ปุจฺฉา.}

\prose{กตมา ทิฏฺฐ\-สํสนฺทนา ปุจฺฉา. ปกติยา ลกฺขณํ ญาตํ โหติ ทิฏฺฐํ ตุลิตํ ตีริตํ วิภูตํ วิภาวิตํ. อญฺเญหิ ปณฺฑิเตหิ สทฺธิํ สํสนฺทน\-ตฺถาย ปญฺหํ ปุจฺฉติ. อยํ ทิฏฺฐ\-สํสนฺทนา ปุจฺฉา.}

\prose{กตมา วิมติจฺเฉทนา ปุจฺฉา. ปกติยา สํสย\-ปกฺขนฺโท\csromansharedfootnote{5904e042bb970f82}{สํสยปกฺขนฺโน (สฺยา)} โหติ วิมติ\-ปกฺขนฺโท เทฺวฬฺหกชาโต “เอวํ นุ โข, น นุ โข, กิํ นุ โข, กถํ นุ โข”ติ. โส วิมติจฺเฉทน\-ตฺถาย ปญฺหํ ปุจฺฉติ. อยํ วิมติจฺเฉทนา ปุจฺฉา. อิมา ติสฺโส ปุจฺฉา.}

\prose{อปราปิ ติสฺโส ปุจฺฉา มนุสฺสปุจฺฉา อมนุสฺสปุจฺฉา นิมฺมิตปุจฺฉา. กตมา มนุสฺสปุจฺฉา. มนุสฺสา พุทฺธํ ภควนฺตํ อุปสงฺกมิตฺวา ปุจฺฉนฺติ, ภิกฺขู ปุจฺฉนฺติ, ภิกฺขุนิโย ปุจฺฉนฺติ, อุปาสกา ปุจฺฉนฺติ, อุปาสิกาโย ปุจฺฉนฺติ, \csromanfolio{51}\textbf{ราชาโน ปุจฺฉนฺติ, ขตฺติยา ปุจฺฉนฺติ, พฺราหฺมณา ปุจฺฉนฺติ, เวสฺสา ปุจฺฉนฺติ,} \textbf{สุทฺทา ปุจฺฉนฺติ, คหฏฺฐา ปุจฺฉนฺติ, ปพฺพชิตา ปุจฺฉนฺติ, อยํ} \textbf{มนุสฺสปุจฺฉา.}}

\proseitem{3}{กตมา อมนุสฺสปุจฺฉา. อมนุสฺสา พุทฺธํ ภควนฺตํ อุปสงฺกมิตฺวา ปญฺหํ ปุจฺฉนฺติ, นาคา ปุจฺฉนฺติ, สุปณฺณา ปุจฺฉนฺติ, ยกฺขา ปุจฺฉนฺติ, อสุรา ปุจฺฉนฺติ, คนฺธพฺพา ปุจฺฉนฺติ, มหาราชาโน ปุจฺฉนฺติ, อินฺทา ปุจฺฉนฺติ, พฺรหฺมาโน ปุจฺฉนฺติ, เทวตาโย ปุจฺฉนฺติ. อยํ อมนุสฺสปุจฺฉา.}

\prose{กตมา นิมฺมิตปุจฺฉา. ยํ ภควา รูปํ อภินิมฺมินาติ มโนมยํ สพฺพ\-งฺค\-ปจฺจงฺคํ อหีนินฺทฺริยํ. โส นิมฺมิโต พุทฺธํ ภควนฺตํ อุปสงฺกมิตฺวา ปญฺหํ ปุจฺฉติ. ภควา วิสชฺเชติ\csromansharedfootnote{e01478662d1a339a}{วิสฺสชฺเชติ (ก)}. อยํ นิมฺมิตปุจฺฉา. อิมา ติสฺโส ปุจฺฉา.}

\prose{อปราปิ ติสฺโส ปุจฺฉา อตฺตตฺถปุจฺฉา ปรตฺถปุจฺฉา อุภยตฺถ\-ปุจฺฉา. อปราปิ ติสฺโส ปุจฺฉา ทิฏฺฐธมฺมิก\-ตฺถ\-ปุจฺฉา สมฺปรายิกตฺถ\-ปุจฺฉา ปรมตฺถ\-ปุจฺฉา. อปราปิ ติสฺโส ปุจฺฉา อนวชฺช\-ตฺถ\-ปุจฺฉา นิกฺกิเลส\-ตฺถ\-ปุจฺฉา โวทาน\-ตฺถ\-ปุจฺฉา. อปราปิ ติสฺโส ปุจฺฉา อตีตปุจฺฉา อนาคตปุจฺฉา ปจฺจุปฺปนฺน\-ปุจฺฉา. อปราปิ ติสฺโส ปุจฺฉา อชฺฌตฺตปุจฺฉา พหิทฺธาปุจฺฉา อชฺฌตฺต\-พหิทฺธาปุจฺฉา. อปราปิ ติสฺโส ปุจฺฉา กุสลปุจฺฉา อกุสลปุจฺฉา อพฺยากตปุจฺฉา. อปราปิ ติสฺโส ปุจฺฉา ขนฺธปุจฺฉา ธาตุปุจฺฉา อายตนปุจฺฉา. อปราปิ ติสฺโส ปุจฺฉา สติปฏฺฐาน\-ปุจฺฉา สมฺมปฺปธาน\-ปุจฺฉา อิทฺธิปาท\-ปุจฺฉา. อปราปิ ติสฺโส ปุจฺฉา อินฺทฺริยปุจฺฉา พลปุจฺฉา โพชฺฌงฺค\-ปุจฺฉา. อปราปิ ติสฺโส ปุจฺฉา มคฺคปุจฺฉา ผลปุจฺฉา นิพฺพานปุจฺฉา. \textbf{ปุจฺฉามิ ต}นฺติ ปุจฺฉามิ ตํ ยาจามิ ตํ อชฺเฌสามิ ตํ ปสาเทมิ ตํ “กถยสฺสุ เม”ติ ปุจฺฉามิ ตํ. \textbf{ภควา}ติ คารวาธิวจนเมตํ ฯเปฯ สจฺฉิกา ปญฺญตฺติ ยทิทํ ภควาติ. \textbf{พฺรูหิเมตนฺ}ติ พฺรูหิ อาจิกฺขาหิ เทเสหิ ปญฺญเปหิ ปฏฺฐเปหิ วิวราหิ วิภชาหิ อุตฺตานีกโรหิ ปกาเสหีติ ปุจฺฉามิ ตํ ภควา พฺรูหิ เมตํ. เตนาห โส พฺราหฺมโณ\csromandash{}}

\vspace{\tipitakaparskip}
\csromancenter{อเนชํ มูลทสฺสาวิํ, (อิจฺจายสฺมา ปุณฺณโก,)}
\prose{อตฺถิ ปเญฺหน อาคมํ.}

\setlength{\csromangathapagewidth}{0pt}
\setlength{\csromangathawakwidth}{0pt}
\csromangathameasuregroup{}{กิํ นิสฺสิตา อิสโย มนุชา, \\ ยญฺญมกปฺปยิํสุ ปุถูธ โลเก, \\ ปุจฺฉามิ ตํ ภควา พฺรูหิ เมตนฺติ.}
\csromangathameasurewak{กิํ นิสฺสิตา อิสโย มนุชา, \\ ยญฺญมกปฺปยิํสุ ปุถูธ โลเก, \\ ปุจฺฉามิ ตํ ภควา พฺรูหิ เมตนฺติ.}
\setlength{\csromangathaleftcol}{0pt}
\csromangathagroup{\csromangathastanza{\csromangathawak{กิํ นิสฺสิตา อิสโย มนุชา, \\ ยญฺญม\-กปฺปยิํสุ ปุถูธ โลเก, \\ ปุจฺฉามิ ตํ ภควา พฺรูหิ เมตนฺติ.}}}

\proseitem{13}{\csromanfolio{52}\textbf{เย เกจิเม อิสโย มนุชา, (ปุณฺณกาติ ภควา,)}}

\setlength{\csromangathapagewidth}{0pt}
\setlength{\csromangathawakwidth}{0pt}
\csromangathameasuregroup{}{\textbf{ยญฺญม}\textbf{กปฺปยิํสุ ปุถูธ โลเก,} \\ \textbf{อาสีสมานา ปุณฺณก อิตฺถตฺตํ}\textsuperscript{1}\textbf{.} \\ \textbf{ชรํ สิตา ยญฺญม}\textbf{กปฺปยิํสุ.}}
\csromangathameasurewak{\textbf{ยญฺญม}\textbf{กปฺปยิํสุ ปุถูธ โลเก,} \\ \textbf{อาสีสมานา ปุณฺณก อิตฺถตฺตํ}\textsuperscript{1}\textbf{.} \\ \textbf{ชรํ สิตา ยญฺญม}\textbf{กปฺปยิํสุ.}}
\setlength{\csromangathaleftcol}{0pt}
\csromangathagroup{\csromangathastanza{\csromangathawak{\textbf{ยญฺญม}\-\textbf{กปฺปยิํสุ ปุถูธ โลเก,} \\ \textbf{อาสีสมานา ปุณฺณก อิตฺถตฺตํ}\csromansharedfootnote{f5300eedb1e3c944}{อิตฺถตํ (สฺยา), อิตฺถภาวํ (ก)}\textbf{.} \\ \textbf{ชรํ สิตา ยญฺญม}\-\textbf{กปฺปยิํสุ.}}}}

\prose{\textbf{เย เกจิเม อิสโย มนุชา}ติ \textbf{เย เกจี}ติ สพฺเพน สพฺพํ สพฺพถา สพฺพํ อเสสํ นิสฺเสสํ ปริยาทิย\-น\-วจนเมตํ เย เกจีติ. \textbf{อิสโย}ติ อิสินามกา เย เกจิ อิสิปพฺพชฺชํ ปพฺพชิตา อาชีวกา นิคณฺฐา ชฏิลา ตาปสา. \textbf{มนุชา}ติ มนุสฺสา วุจฺจนฺตีติ เย เกจิเม อิสโย มนุชา, ปุณฺณกาติ ภควา.}

\prose{\textbf{ขตฺติยา พฺราหฺมณา เทวตานนฺ}ติ \textbf{ขตฺติยา}ติ เย เกจิ ขตฺติยชาติกา. \textbf{พฺราหฺมณา}ติ เย เกจิ โภวาทิกา. \textbf{เทวตานนฺ}ติ อาชีวก\-สาวกานํ อาชีวกา เทวตา ฯเปฯ. ทิสาวติกานํ ทิสา เทวตา. เย เยสํ ทกฺขิเณยฺยา, เต เตสํ เทวตาติ ขตฺติยา พฺราหฺมณา เทวตานํ.}

\prose{\textbf{ยญฺญม}\-\textbf{กปฺปยิํสุ ปุถูธ โลเก}ติ \textbf{ยญฺญํ} วุจฺจติ เทยฺยธมฺโม, จีวร\-ปิณฺฑปาต\-เสนาสน\-คิลานปจฺจย\-เภสชฺช\-ปริกฺขารํ อนฺนํ ปานํ ฯเปฯ เสยฺยา\-วสถ\-ปทีเปยฺยํ. \textbf{ยญฺญมกปฺปยิํสู}ติ เยปิ ยญฺญํ เอสนฺติ คเวสนฺติ ปริเยสนฺติ ฯเปฯ เสยฺยา\-วสถ\-ปทีเปยฺยํ, เตปิ ยญฺญํ กปฺเปนฺติ. \textbf{ปุถู}ติ ยญฺญา วา เอเต ปุถู, ยญฺญยาชกา วา เอเต ปุถู, ทกฺขิเณยฺยา วา เอเต ปุถู ฯเปฯ. เอวํ ทกฺขิเณยฺยา วา เอเต ปุถู. \textbf{อิธ โลเก}ติ มนุสฺสโลเกติ ยญฺญม\-กปฺปยิํสุ ปุถูธ โลเก.}

\prose{\textbf{อาสีสมานา ปุณฺณก อิตฺถตฺตนฺ}ติ อาสีสมานาติ รูปปฏิลาภํ \textbf{อาสีสมานา} สทฺท\-ปฏิลาภํ อาสีสมานา คนฺธ\-ปฏิลาภํ อาสีสมานา รสปฏิลาภํ อาสีสมานา โผฏฺฐพฺพ\-ปฏิลาภํ อาสีสมานา ปุตฺต\-ปฏิลาภํ อาสีสมานา ทารปฏิลาภํ อาสีสมานา ธนปฏิลาภํ อาสีสมานา ยสปฏิลาภํ อาสีสมานา อิสฺสริย\-ปฏิลาภํ อาสีสมานา ขตฺติย\-มหาสาล\-กุเล อตฺตภาว\-ปฏิลาภํ อาสีสมานา \csromanfolio{53}\textbf{พฺราหฺมณ}\-\textbf{มหาสาล}\-\textbf{กุเล อตฺตภาว}\-\textbf{ปฏิลาภํ อาสีสมานา} \textbf{คหปติ}\-\textbf{มหาสาล}\-\textbf{กุเล อตฺตภาว}\-\textbf{ปฏิลาภํ อาสีสมานา จาตุมหาราชิเกสุ}\csromansharedfootnote{f9f1754494c64cd4}{จาตุมฺมหาราชิเกสุ (สฺยา)} \textbf{เทเวสุ อตฺตภาว}\-\textbf{ปฏิลาภํ อาสีสมานา ตาวติํเสสุ เทเวสุ. ยาเมสุ} \textbf{เทเวสุ. ตุสิเตสุ เทเวสุ. นิมฺมานรตีสุ เทเวสุ. ปรนิมฺมิต}\-\textbf{วส}\-\textbf{วตฺตีสุ} \textbf{เทเวสุ. พฺรหฺมกายิเกสุ เทเวสุ อตฺตภาว}\-\textbf{ปฏิลาภํ อาสีสมานา อิจฺฉมานา} \textbf{สาทิยมานา ปตฺถยมานา ปิหยมานา อภิชปฺป}\-\textbf{มานาติ อาสีสมานา.}}

\prose{\textbf{ปุณฺณก} \textbf{อิตฺถตฺตนฺ}ติ เอตฺถ อตฺตภาวา\-ภินิพฺพตฺติํ อาสีสมานา เอตฺถ ขตฺติย\-มหาสาล\-กุเล อตฺตภาวา\-ภินิพฺพตฺติํ อาสีสมานา ฯเปฯ. เอตฺถ พฺรหฺมกายิเกสุ เทเวสุ อตฺตภาวา\-ภินิพฺพตฺติํ อาสีสมานา อิจฺฉยมานา สาทิยมานา ปตฺถยมานา ปิหยมานา อภิชปฺป\-มานาติ อาสีสมานา ปุณฺณก อิตฺถตฺตํ.}

\prose{\textbf{ชรํ สิตา ยญฺญมกปฺปยิํสู}ติ ชรานิสฺสิตา พฺยาธินิสฺสิตา มรณนิสฺสิตา โสกปริเทวทุกฺขโทมนสฺสุปายาส\-นิสฺสิตา. ยเทว เต ชาตินิสฺสิตา, ตเทว เต ชรานิสฺสิตา. ยเทว เต ชรานิสฺสิตา, ตเทว เต พฺยาธินิสฺสิตา. ยเทว เต พฺยาธินิสฺสิตา, ตเทว เต มรณนิสฺสิตา. ยเทว เต มรณนิสฺสิตา, ตเทว เต โสกปริเทวทุกฺขโทมนสฺสุปายาส\-นิสฺสิตา. ยเทว เต โสกปริเทวทุกฺขโทมนสฺสุปายาส\-นิสฺสิตา, ตเทว เต คตินิสฺสิตา. ยเทว เต คตินิสฺสิตา, ตเทว เต อุปปตฺติ\-นิสฺสิตา. ยเทว เต อุปปตฺติ\-นิสฺสิตา, ตเทว เต ปฏิสนฺธิ\-นิสฺสิตา. ยเทว เต ปฏิสนฺธิ\-นิสฺสิตา, ตเทว เต ภวนิสฺสิตา. ยเทว เต ภวนิสฺสิตา, ตเทว เต สํสารนิสฺสิตา. ยเทว เต สํสารนิสฺสิตา, ตเทว เต วฏฺฏนิสฺสิตา อลฺลีนา อุปคตา อชฺโฌสิตา อธิมุตฺตาติ ชรํ สิตา ยญฺญม\-กปฺปยิํสุ. เตนาห ภควา\csromandash{}}

\vspace{\tipitakaparskip}
\csromancenter{เย เกจิเม อิสโย มนุชา, (ปุณฺณกาติ ภควา,)}
\setlength{\csromangathapagewidth}{0pt}
\setlength{\csromangathawakwidth}{0pt}
\csromangathameasuregroup{}{ยญฺญมกปฺปยิํสุ ปุถูธ โลเก, \\ อาสีสมานา ปุณฺณก อิตฺถตฺตํ.}
\csromangathameasurewak{ยญฺญมกปฺปยิํสุ ปุถูธ โลเก, \\ อาสีสมานา ปุณฺณก อิตฺถตฺตํ.}
\setlength{\csromangathaleftcol}{0pt}
\csromangathagroup{\csromangathastanza{\csromangathawak{ยญฺญม\-กปฺปยิํสุ ปุถูธ โลเก, \\ อาสีสมานา ปุณฺณก อิตฺถตฺตํ.}}}

\prose{ชรํ สิตา ยญฺญมกปฺปยิํสูติ.}

\csromanhangingbegin
\hangingheaditem{14}{\csromanfolio{54}\textbf{เย เกจิเม อิสโย มนุชา, (อิจฺจายสฺมา ปุณฺณาโก,)}}
\hangingbody{\textbf{ขตฺติยา พฺรหฺมณา เทวตานํ,}}
\hangingbody{\textbf{ยญฺญมกปฺปยิํสุ ปุถูธ โลเก,}}
\hangingbody{\textbf{กจฺจิสุ เต ภควา ยญฺญปเถ อปฺปมตฺตา.}}
\hangingbody{\textbf{อตารุ1} \textbf{ชาติญฺจ ชรญฺจ มาริส,}}
\hangingbody{\textbf{ปุจฺฉามิ ตํ ภควา พฺรูหิ เมตํ.}}
\csromanhangingend

\prose{\textbf{เย เกจิเม อิสโย มนุชา}ติ \textbf{เย เกจี}ติ ฯเปฯ.}

\prose{\textbf{กจฺจิสุ เต ภควา ยญฺญปเถ อปฺปมตฺตา}ติ \textbf{กจฺจิสู}ติ สํสยปุจฺฉา วิมติปุจฺฉา เทฺวฬฺหกปุจฺฉา อเนกํสปุจฺฉา “เอวํ นุ โข, น นุ โข กิํ นุ โข, กถํ นุ โข”ติ กจฺจิสุ. \textbf{เต}ติ ยญฺญยาชกา วุจฺจนฺติ. \textbf{ภควา}ติ คารวา\-ธิวจนํ ฯเปฯ สจฺฉิกา ปญฺญตฺติ, ยทิทํ ภควาติ กจฺจิสุ เต ภควา. ยญฺญปเถ อปฺปมตฺตาติ ยญฺโญเยว วุจฺจติ ยญฺญปโถ. ยถา อริยมคฺโค อริยปโถ เทวมคฺโค เทวปโถ พฺรหฺมมคฺโค พฺรหฺมปโถ. เอวเมว ยญฺโญเยว วุจฺจติ ยญฺญปโถ. อปฺปมตฺตาติ ยญฺญปเถ อปฺปมตฺตา สกฺกจฺจการิโน สาตจฺจการิโน อฏฺฐิตการิโน อโนลีนวุตฺติโน อนิกฺขิตฺตจฺฉนฺทา อนิกฺขิตฺตธุรา ตจฺจริตา ตพฺพหุลา ตคฺครุกา ตนฺนินฺนา ตปฺโปณา ตปฺปพฺภารา ตทธิมุตฺตา ตทธิปเตยฺยาติ เตปิ ยญฺญปเถ อปฺปมตฺตา. เยปิ ยญฺญํ เอสนฺติ คเวสนฺติ ปริเยสนฺติ จีวร\-ปิณฺฑปาต\-เสนาสน\-คิลานปจฺจย\-เภสชฺช\-ปริกฺขารํ อนฺนํ ปานํ ฯเปฯ เสยฺยา\-วสถ\-ปทีเปยฺยํ สกฺกจฺจการิโน ฯเปฯ ตทธิปเตยฺยา. เตปิ ยญฺญปเถ อปฺปมตฺตา. เยปิ ยญฺญํ อภิส\-งฺขโรนฺติ จีวร\-ปิณฺฑปาต\-เสนาสน\-คิลานปจฺจย\-เภสชฺช\-ปริกฺขารํ อนฺนํ ปานํ ฯเปฯ เสยฺยา\-วสถ\-ปทีเปยฺยํ สกฺกจฺจการิโน ฯเปฯ ตทธิปเตยฺยา. เตปิ ยญฺญปเถ อปฺปมตฺตา. เยปิ ยญฺญํ เทนฺติ ยชนฺติ ปริจฺจชนฺติ จีวร\-ปิณฺฑปาต\-เสนาสน\-คิลานปจฺจย\-เภสชฺช\-ปริกฺขารํ อนฺนํ ปานํ ฯเปฯ เสยฺยา\-วสถ\-ปทีเปยฺยํ สกฺกจฺจการิโน ฯเปฯ ตทธิปเตยฺยา. เตปิ ยญฺญปเถ อปฺปมตฺตาติ กจฺจิสุ เต ภควา ยญฺญปเถ อปฺปมตฺตา.}

\prose{\textbf{อตารุ ชาติญฺจ ชรญฺจ มาริสา}ติ ชรามรณํ อตริํสุ อุตฺตริํสุ ปตริํสุ สมติกฺกมิํสุ วีติวตฺติํสุ. \textbf{มาริสา}ติ ปิยวจนํ ครุวจนํ สคารวสปฺปติสฺสาธิวจนเมตํ มาริสาติ อตารุ ชาติญฺจ ชรญฺจ มาริส.}

\prose{\csromanfolio{55}\textbf{ปุจฺฉามิ ตํ ภควา พฺรูหิ เมตนฺ}ติ ปุจฺฉามิ \textbf{ตนฺ}ติ \textbf{ปุจฺฉามิ} ตํ ยาจามิ ตํ อชฺเฌสามิ ตํ ปสาเทมิ ตํ กถยสฺสุ เมติ ปุจฺฉามิ ตํ. \textbf{ภควา}ติ คารวา\-ธิวจนํ ฯเปฯ สจฺฉิกา ปญฺญตฺติ, ยทิทํ ภควาติ. \textbf{พฺรูหิ เมตนฺ}ติ พฺรูหิ อาจิกฺขาหิ เทเสหิ ปญฺญเปหิ ปฏฺฐเปหิ วิวราหิ วิภชาหิ อุตฺตานีกโรหิ ปกาเสหีติ ปุจฺฉามิ ตํ ภควา พฺรูหิ เมตํ. เตนาห โส พฺราหฺมโณ\csromandash{}}

\prose{เย เกจิเม อิสโย มนุชา, (อิจฺจายสฺมา ปุณฺณโก,)}

\prose{ขตฺติยา พฺรหฺมณา เทวตานํ.}

\setlength{\csromangathapagewidth}{0pt}
\setlength{\csromangathawakwidth}{0pt}
\csromangathameasuregroup{}{ยญฺญมกปฺปยิํสุ ปุถูธ โลเก, \\ กจฺจิสุ เต ภควา ยญฺญปเถ อปฺปมตฺตา. \\ อตารุ ชาติญฺจ ชรญฺจ มาริส,}
\csromangathameasurewak{ยญฺญมกปฺปยิํสุ ปุถูธ โลเก, \\ กจฺจิสุ เต ภควา ยญฺญปเถ อปฺปมตฺตา. \\ อตารุ ชาติญฺจ ชรญฺจ มาริส,}
\setlength{\csromangathaleftcol}{0pt}
\csromangathagroup{\csromangathastanza{\csromangathawak{ยญฺญม\-กปฺปยิํสุ ปุถูธ โลเก, \\ กจฺจิสุ เต ภควา ยญฺญปเถ อปฺปมตฺตา. \\ อตารุ ชาติญฺจ ชรญฺจ มาริส,}}}

\prose{ปุจฺฉามิ ตํ ภควา พฺรูหิ เม ตนฺติ.}

\csromanhangingbegin
\hangingheaditem{15}{อาสีสนฺติ\csromansharedfootnote{2ef2705b45a15d00}{อาสิํสนฺติ (สฺยา)} \textbf{โถมยนฺติ อภิชปฺปนฺติ ชุหนฺติ, (ปุณฺณกาติ ภควา,)}}
\hangingbody{\textbf{กามาภิชปฺปนฺติ ปฏิจฺจ ลาภํ.}}
\hangingbody{\textbf{เต ยาชโยคา ภวราครตฺตา,}}
\hangingbody{\textbf{นาตริํสุ ชาติชรนฺติ พฺรูมิ.}}
\csromanhangingend

\prose{\textbf{อาสีสนฺติ โถมยนฺติ อภิชปฺปนฺติ ชุหนฺติ}ติ \textbf{อาสีสนฺตี}ติ รูปปฏิลาภํ อาสีสนฺติ. สทฺท\-ปฏิลาภํ อาสีสนฺติ. คนฺธ\-ปฏิลาภํ อาสีสนฺติ. รสปฏิลาภํ อาสีสนฺติ. โผฏฺฐพฺพ\-ปฏิลาภํ อาสีสนฺติ. ปุตฺต\-ปฏิลาภํ อาสีสนฺติ. ทารปฏิลาภํ อาสีสนฺติ. ธนปฏิลาภํ อาสีสนฺติ. ยสปฏิลาภํ อาสีสนฺติ. อิสฺสริย\-ปฏิลาภํ อาสีสนฺติ ขตฺติย\-มหาสาล\-กุเล อตฺตภาว\-ปฏิลาภํ อาสีสนฺติ. พฺราหฺมณ\-มหาสาล\-กุเล. คหปติ\-มหาสาล\-กุเล อตฺตภาว\-ปฏิลาภํ อาสีสนฺติ. จาตุมหาราชิเกสุ เทเวสุ ฯเปฯ. พฺรหฺมกายิเกสุ เทเวสุ อตฺตภาว\-ปฏิลาภํ อาสีสนฺติ อิจฺฉนฺติ สาทิยนฺติ ปตฺถยนฺติ ปิหยนฺตีติ อาสีสนฺติ.}

\prose{\textbf{โถมยนฺตี}ติ ยญฺญํ วา โถเมนฺติ ผลํ วา โถเมนฺติ ทกฺขิเณยฺเย วา โถเมนฺติ. กถํ ยญฺญํ โถเมนฺติ. สุจิํ ทินฺนํ\csromansharedfootnote{8e887095aef2f427}{ปิยํ ทินฺนํ (สฺยา)} มนาปํ ทินฺนํ ปณีตํ ทินฺนํ \csromanfolio{56}กาเลน ทินฺนํ กปฺปิยํ ทินฺนํ วิเจยฺย ทินฺนํ อนวชฺชํ ทินฺนํ อภิณฺหํ ทินฺนํ ททํ จิตฺตํ ปสาทิตนฺติ โถเมนฺติ กิตฺเตนฺติ วณฺเณนฺติ ปสํสนฺติ. เอวํ ยญฺญํ โถเมนฺติ.}

\prose{กถํ ผลํ โถเมนฺติ. อิโต นิทานํ รูปปฏิลาโภ ภวิสฺสติ ฯเปฯ. พฺรหฺมกายิเกสุ เทเวสุ อตฺตภาว\-ปฏิลาโภ ภวิสฺสตีติ โถเมนฺติ กิตฺเตนฺติ วณฺเณนฺติ ปสํสนฺติ เอวํ ผลํ โถเมนฺติ.}

\prose{กถํ ทกฺขิเณยฺเย โถเมนฺติ. ทกฺขิเณยฺยา ชาติสมฺปนฺนา โคตฺตสมฺปนฺนา อชฺฌายกา มนฺตธรา ติณฺณํ เวทานํ ปารคู สนิฆณฺฑุ\-เกฏุภานํ สากฺขร\-ปฺปเภทานํ อิติหาส\-ปญฺจมานํ ปทกา เวยฺยากรณา โลกายต\-มหา\-ปุริส\-ลกฺขเณสุ อนวยาติ, วีตราคา วา ราควินยาย วา ปฏิปนฺนา, วีตโทสา วา โทสวินยาย วา ปฏิปนฺนา, วีตโมหา วา โมหวินยาย วา ปฏิปนฺนา, สทฺธาสมฺปนฺนา สีลสมฺปนฺนา สมาธิ\-สมฺปนฺนา ปญฺญาสมฺปนฺนา วิมุตฺติ\-สมฺปนฺนา วิมุตฺติ\-ญาณ\-ทสฺสน\-สมฺปนฺนาติ โถเมนฺติ กิตฺเตนฺติ วณฺเณนฺติ ปสํสนฺติ, เอวํ ทกฺขิเณยฺเย โถเมนฺตีติ อาสีสนฺติ โถมยนฺติ.}

\prose{\textbf{อภิชปฺปนฺตี}ติ รูปปฏิลาภํ อภิชปฺปนฺติ, สทฺท\-ปฏิลาภํ อภิชปฺปนฺติ. คนฺธ\-ปฏิลาภํ อภิชปฺปนฺติ. รสปฏิลาภํ อภิชปฺปนฺติ ฯเปฯ. พฺรหฺมกายิเกสุ เทเวสุ อตฺตภาว\-ปฏิลาภํ อภิชปฺปนฺตีติ อาสีสนฺติ โถมยนฺติ อภิชปฺปนฺติ. \textbf{ชุหนฺตี}ติ ชุหนฺติ เทนฺติ ยชนฺติ ปริจฺจชนฺติ จีวร\-ปิณฺฑปาต\-เสนาสน\-คิลานปจฺจย\-เภสชฺช\-ปริกฺขารํ อนฺนํ ปานํ วตฺถํ ยานํ มาลา\-คนฺธ\-วิเลปนํ เสยฺยา\-วสถ\-ปทีเปยฺยนฺติ อาสีสนฺติ โถมยนฺติ อภิชปฺปนฺติ ชุหนฺติ. ปุณฺณกาติ ภควา.}

\prose{\textbf{กามา}\-\textbf{ภิชปฺปนฺติ ปฏิจฺจ ลาภนฺ}ติ รูปปฏิลาภํ ปฏิจฺจ กาเม อภิชปฺปนฺติ สทฺท\-ปฏิลาภํ ปฏิจฺจ กาเม อภิชปฺปนฺติ ฯเปฯ. พฺรหฺมกายิเกสุ เทเวสุ อตฺตภาว\-ปฏิลาภํ ปฏิจฺจ กาเม อภิชปฺปนฺติ ปชปฺปนฺตีติ กามา\-ภิชปฺปนฺติ ปฏิจฺจ ลาภํ.}

\prose{\textbf{เต ยาชโยคา ภวราครตฺตา, นาตริํสุ ชาติชรนฺติ พฺรูมี}ติ \textbf{เต}ติ ยญฺญยาชกา วุจฺจนฺติ. \textbf{ยาชโยคา}ติ ยาชโยเคสุ ยุตฺตา ปยุตฺตา อายุตฺตา สมายุตฺตา ตจฺจริตา ตพฺพหุลา ตคฺครุกา ตนฺนินฺนา ตปฺโปณา ตปฺปพฺภารา ตทธิมุตฺตา ตทธิปเตยฺยาติ เต ยาชโยคา. \textbf{ภวราครตฺตา}ติ ภวราโค วุจฺจติ โย ภเวสุ ภวจฺฉนฺโท ภวราโค \csromanfolio{57}\textbf{ภวนนฺที ภวตณฺหา ภวสิเนโห ภวปริฬาโห ภวมุจฺฉา} ภวชฺโฌสานํ. ภวราเคน ภเวสุ รตฺตา คิทฺธา คธิตา มุจฺฉิตา อชฺโฌสนฺนา ลคฺคา ลคฺคิตา ปลิพุทฺธาติ เต ยาชโยคา ภวราครตฺตา.}

\prose{\textbf{นาตริํสุ ชาติชรนฺติ พฺรูมี}ติ เต ยาชโยคา ภวราครตฺตา ชาติ\-ชรา\-มรณํ นาตริํสุ น อุตฺตริํสุ น ปตริํสุ น สมติกฺกมิํสุ น วีติวตฺติํสุ ชาติชรามรณา อนิกฺขนฺตา อนิสฺสฏา อนติกฺกนฺตา อสมติกฺกนฺตา อวีติวตฺตา อนฺโต\-ชาติชรามรเณ ปริวตฺตนฺติ. อนฺโต\-สํสาร\-ปเถ ปริวตฺตนฺติ. ชาติยา อนุคตา ชราย อนุสฏา พฺยาธินา อภิภูตา มรเณน อพฺภาหตา อตาณา อเลณา อสรณา อสรณีภูตาติ พฺรูมิ อาจิกฺขามิ เทเสมิ ปญฺญเปมิ ปฏฺฐเปมิ วิวรามิ วิภชามิ อุตฺตานีกโรมิ ปกาเสมีติ เต ยาชโยคา ภวราครตฺตา, นาตริํสุ ชาติชรนฺติ พฺรูมิ. เตนาห ภควา\csromandash{}}

\prose{อาสีสนฺติ โถมยนฺติ อภิชปฺปนฺติ ชุหนฺติ, (ปุณฺณกาติ ภควา,)}

\prose{กามา’ภิชปฺปนฺติ ปฏิจฺจ ลาภํ.}

\prose{เต ยาชโยคา ภวราครตฺตา,}

\prose{นาตริํสุ ชาติชรนฺติ พฺรูมีติ.}

\csromanhangingbegin
\hangingheaditem{15}{\textbf{เต เจ นาตริํสุ ยาชโยคา, (อิจฺจายสฺมา ปุณฺณโก,)}}
\hangingbody{\textbf{ยญฺเญหิ ชาติญฺจ ชรญฺจ มาริส.}}
\hangingbody{\textbf{อถ โก จรหิ เทวมนุสฺสโลเก,}}
\hangingbody{\textbf{อตาริ ชาติญฺจ ชรญฺจ มาริส.}}
\hangingbody{\textbf{ปุจฺฉามิ ตํ ภควา พฺรูหิ เมตํ.}}
\csromanhangingend

\prose{\textbf{เต เจ นาตริํสุ ยาชโยคา}ติ เต ยญฺญยาชกา ยาชโยคา ภวราครตฺตา ชาติ\-ชรา\-มรณํ นาตริํสุ น อุตฺตริํสุ น ปตริํสุ น สมติกฺกมิํสุ น วีติวตฺติํสุ, ชาติชรามรณา อนิกฺขนฺตา อนิสฺสฏา อนติกฺกนฺตา อสมติกฺกนฺตา อวีติวตฺตา อนฺโต\-ชาติชรามรเณ ปริวตฺตนฺติ. อนฺโต\-สํสาร\-ปเถ ปริวตฺตนฺติ. ชาติยา อนุคตา ชราย อนุสฏา พฺยาธินา อภิภูตา มรเณน อพฺภาหตา อตาณา อเลณา อสรณา อสรณีภูตาติ เต เจ นาตริํสุ ยาชโยคา.}

\prose{\textbf{อิจฺจายสฺมา ปุณฺณโก}ติ อิจฺจาติ ปทสนฺธิ ฯเปฯ อายสฺมา ปุณฺณโก.}

\proseitem{16}{\csromanfolio{58}\textbf{ยญฺเญหิ ชาติญฺจ ชรญฺจ มาริสา}ติ \textbf{ยญฺเญหี}ติ ยญฺเญหิ ปหูเตหิ ยญฺเญหิ วิวิเธหิ ยญฺเญหิ ปุถูหิ. \textbf{มาริสา}ติ ปิยวจนํ ครุวจนํ สคารวสปฺปติสฺสาธิวจนเมตํ มาริสาติ ยญฺเญหิ ชาติญฺจ ชรญฺจ มาริส.}

\prose{\textbf{อถ โก จรหิ เทว}\-\textbf{มนุสฺส}\-\textbf{โลเก, อตาริ ชาติญฺจ ชรญฺจ มาริสา}ติ อถ โก เอโส สเทวเก โลเก สมารเก สพฺรหฺมเก สสฺสมณ\-พฺราหฺมณิยา ปชาย สเทว\-มนุสฺสาย ชาติ\-ชรา\-มรณํ อตริ อุตฺตริ ปตริ สมติกฺกมิ วีติวตฺตยิ\csromansharedfootnote{09559e842b5e6dd1}{วีติวตฺติ (ก)}. \textbf{มาริสา}ติ ปิยวจนํ ครุวจนํ สคารวสปฺปติสฺสาธิวจนเมตํ มาริสาติ อถ โก จรหิ เทว\-มนุสฺส\-โลเก, อตาริ ชาติญฺจ ชรญฺจ มาริส.}

\prose{\textbf{ปุจฺฉามิ ตํ ภควา พฺรูหิ เมตนฺ}ติ \textbf{ปุจฺฉามิ ตนฺ}ติ ปุจฺฉามิ ตํ ยาจามิ ตํ อชฺเฌสามิ ตํ ปสาเทมิ ตํ, กถยสฺสุ เมตนฺติ ปุจฺฉามิ ตํ. \textbf{ภควา}ติ คารวา\-ธิวจนํ ฯเปฯ สจฺฉิกา ปญฺญตฺติ, ยทิทํ ภควาติ. \textbf{พฺรูหิ เมตนฺ}ติ พฺรูหิ อาจิกฺขาหิ เทเสหิ ปญฺญเปหิ ปฏฺฐเปหิ วิวราหิ วิภชาหิ อุตฺตานีกโรหิ ปกาเสหีติ ปุจฺฉามิ ตํ ภควา พฺรูหิ เมตํ. เตนาห โส พฺราหฺมโณ\csromandash{}}

\prose{เต เจ นาตริํสุ ยาชโยคา, (อิจฺจายสฺมา ปุณฺณโก,)}

\prose{ยญฺเญหิ ชาติญฺจ ชรญฺจ มาริส.}

\prose{\textbf{อถ โก จรหิ เทว}\-\textbf{มนุสฺส}\-\textbf{โล}เก,}

\prose{อตาริ ชาติญฺจ ชรญฺจ มาริส.}

\prose{ปุจฺฉามิ ตํ ภควา พฺรูหิ เมตนฺติ.}

\csromanhangingbegin
\hangingheaditem{17}{สงฺขาย โลกสฺมิ\csromansharedfootnote{3cfa20bfd1c2b167}{โลกสฺมิํ (สฺยา, ก)} \textbf{ปโรปรานิ, (ปุณฺณกาติ ภควา,)}}
\hangingbody{\textbf{ยสฺสิญฺชิตํ นตฺถิ กุหิญฺจิ โลเก.}}
\hangingbody{\textbf{สนฺโต วิธูโม อนีโฆ นิราโส,}}
\hangingbody{\textbf{อตาริ โส ชาติชรนฺติ พฺรูมิ.}}
\csromanhangingend

\prose{\textbf{สงฺขาย โลกสฺมิ ปโรปรานี}ติ สงฺขา วุจฺจติ ญาณํ. ยา ปญฺญา ปชานนา ฯเปฯ อโมโห ธมฺมวิจโย สมฺมาทิฏฺฐิ. \textbf{ปโรปรานี}ติ \textbf{โอรํ} วุจฺจติ สกตฺตภาโว. \textbf{ปรํ} วุจฺจติ ปรตฺตภาโว. โอรํ วุจฺจติ สกรู\-ปเวทนา\-สญฺญา\-สงฺขาร\-วิญฺญาณํ. ปรํ วุจฺจติ ปรรูปเวทนาสญฺญาสงฺขาร- \csromanfolio{59}วิญฺญาณํ. โอรํ วุจฺจติ ฉ อชฺฌตฺติกานิ อายตนานิ. ปรํ วุจฺจติ ฉ พาหิรานิ อายตนานิ. โอรํ วุจฺจติ มนุสฺสโลโก. ปรํ วุจฺจติ เทวโลโก. โอรํ วุจฺจติ กามธาตุ. ปรํ วุจฺจติ รูปธาตุ อรูปธาตุ. โอรํ วุจฺจติ กามธาตุ รูปธาตุ. ปรํ วุจฺจติ อรูปธาตุ. สงฺขาย โลกสฺมิ ปโรปรานีติ ปโรปรานิ อนิจฺจโต สงฺขาย, ทุกฺขโต. โรคโต. คณฺฑโต ฯเปฯ. นิสฺสรณโต สงฺขาย ชานิตฺวา ตุลยิตฺวา ตีรยิตฺวา วิภาวยิตฺวา วิภูตํ กตฺวาติ สงฺขาย โลกสฺมิ ปโรปรานิ. \textbf{ปุณฺณกาติ ภควา}ติ \textbf{ปุณฺณกา}ติ ภควา ตํ พฺราหฺมณํ นาเมน อาลปติ. \textbf{ภควา}ติ คารวาธิวจนเมตํ ฯเปฯ ยทิทํ ภควาติ ปุณฺณกาติ ภควา.}

\prose{\textbf{ยสฺสิญฺชิตํ นตฺถิ กุหิญฺจิ โลเก}ติ \textbf{ยสฺสา}ติ อรหโต ขีณาสวสฺส. \textbf{อิญฺชิตนฺ}ติ ตณฺหิญฺชิตํ ทิฏฺฐิญฺชิตํ มานิญฺชิตํ กิเลสิญฺชิตํ กามิญฺชิตํ. ยสฺสิเม อิญฺชิตา นตฺถิ น สนฺติ น สํวิชฺชนฺติ นุปลพฺภนฺติ, ปหีนา สมุจฺฉินฺนา วูปสนฺตา ปฏิปสฺสทฺธา อภพฺพุ\-ปฺปตฺติกา ญาณคฺคินา ทฑฺฒา. \textbf{กุหิญฺจี}ติ กุหิญฺจิ กิสฺมิญฺจิ กตฺถจิ อชฺฌตฺตํ วา พหิทฺธา วา อชฺฌตฺต\-พหิทฺธา วา. \textbf{โลเก}ติ อปายโลเก ฯเปฯ อายตนโลเกติ ยสฺสิญฺชิตํ นตฺถิ กุหิญฺจิ โลเก.}

\prose{\textbf{สนฺโต วิธูโม อนีโฆ นิราโส, อตาริ โส ชาติชรนฺติ พฺรูมี}ติ \textbf{สนฺโต}ติ ราคสฺส สนฺตตฺตา สนฺโต. โทสสฺส. โมหสฺส. โกธสฺส. อุปนาหสฺส. มกฺขสฺส ฯเปฯ. สพฺพา\-กุสลา\-ภิสงฺขารานํ สนฺตตฺตา สมิตตฺตา วูปสมิตตฺตา วิชฺฌาตตฺตา\csromansharedfootnote{f0510903136555b6}{นิชฺฌาตตฺตา (ก) ขุ 7. 53 ปิฏฺเฐ.} นิพฺพุตตฺตา วิคตตฺตา ปฏิปสฺสทฺธ\-ตฺตา สนฺโต อุปสนฺโต วูปสนฺโต นิพฺพุโต ปฏิปสฺสทฺโธติ สนฺโต. \textbf{วิธูโม}ติ กายทุจฺจริตํ วิธูมิตํ วิธมิตํ โสสิตํ วิโสสิตํ พฺยนฺตีกตํ\csromansharedfootnote{e8f0ee77d888e9f9}{พฺยนฺติกตํ (ก)}. วจีทุจฺจริตํ. มโนทุจฺจริตํ วิธูมิตํ วิธมิตํ โสสิตํ วิโสสิตํ พฺยนฺตีกตํ ราโค, โทโส, โมโห วิธูมิโต วิธมิโต โสสิโต วิโสสิโต พฺยนฺตีกโต. โกโธ. อุปนาโห. มกฺโข. ปฬาโส. อิสฺสา. มจฺฉริยํ. มายา. สาเฐยฺยํ. ถมฺโภ. สารมฺโภ. มาโน. อติมาโน. มโท. ปมาโท. สพฺเพ กิเลสา. สพฺเพ ทุจฺจริตา. สพฺเพ ทรถา. สพฺเพ ปริฬาหา. สพฺเพ สนฺตาปา. สพฺพา\-กุสลา\-ภิสงฺขารา วิธูมิตา วิธมิตา โสสิตา วิโสสิตา พฺยนฺตีกตา. อถ วา โกโธ วุจฺจติ ธูโม\textbf{.}}

\setlength{\csromangathapagewidth}{0pt}
\setlength{\csromangathawakwidth}{0pt}
\csromangathameasuregroup{}{มาโน หิ เต พฺราหฺมณ ขาริภาโร, \\ โกโธ ธูโม ภสฺมนิ\textsuperscript{1} โมสวชฺชํ. \\ ชิวฺหา สุชา หทยํ\textsuperscript{1} โชติฐานํ, \\ อตฺตา สุทนฺโต ปุริสสฺส โชติ.}
\csromangathameasurewak{มาโน หิ เต พฺราหฺมณ ขาริภาโร, \\ โกโธ ธูโม ภสฺมนิ\textsuperscript{1} โมสวชฺชํ. \\ ชิวฺหา สุชา หทยํ\textsuperscript{1} โชติฐานํ, \\ อตฺตา สุทนฺโต ปุริสสฺส โชติ.}
\setlength{\csromangathaleftcol}{0pt}
\csromangathagroup{\csromangathastanza{\csromangathawak{\csromanfolio{60}มาโน หิ เต พฺราหฺมณ ขาริภาโร, \\ โกโธ ธูโม ภสฺมนิ\csromansharedfootnote{36f1b6c672dd4861}{คมฺมนิ (สฺยา)} โมสวชฺชํ. \\ ชิวฺหา สุชา หทยํ\csromansharedfootnote{a208bd24dbbc05d7}{ตปฺปรสฺส (สฺยา)} โชติฐานํ, \\ อตฺตา สุทนฺโต ปุริสสฺส โชติ.}}}

\proseitem{3}{อปิ จ ทสหากาเรหิ โกโธ ชายติ, “อนตฺถํ เม อจรี”ติ โกโธ ชายติ, “อนตฺถํ เม จรตี”ติ โกโธ ชายติ, “อนตฺถํ เม จริสฺสตี”ติ โกโธ ชายติ. “ปิยสฺส เม มนาปสฺส อนตฺถํ อจริ. อนตฺถํ จรติ. อนตฺถํ จริสฺสตี”ติ โกโธ ชายติ. “อปฺปิยสฺส เม อมนาปสฺส อตฺถํ อจริ. อตฺถํ จรติ. อตฺถํ จริสฺสตี”ติ โกโธ ชายติ. อฏฺฐาเน วา ปน โกโธ ชายติ. โย เอวรูโป จิตฺตสฺส อาฆาโต ปฏิฆาโต ปฏิฆํ ปฏิวิโรโธ โกโป ปโกโป สมฺปโกโป โทโส ปโทโส สมฺปโทโส จิตฺตสฺส พฺยาปตฺติ มโนปโทโส โกโธ กุชฺฌนา กุชฺฌิตตฺตํ โทโส ทุสฺสนา ทุสฺสิตตฺตํ พฺยาปตฺติ พฺยาปชฺชนา พฺยาปชฺชิตตฺตํ วิโรโธ ปฏิวิโรโธ จณฺฑิกฺกํ อสุโรโป\csromansharedfootnote{a1be0447fa64ec56}{อสฺสุโรโป (สฺยา)} อนตฺตมนตา จิตฺตสฺส. อยํ วุจฺจติ โกโธ.}

\prose{อปิ จ โกธสฺส อธิมตฺต\-ปริตฺตตา เวทิตพฺพา, อตฺถิ กญฺจิ\csromansharedfootnote{e1f0f314dba37ff9}{กิญฺจิ (ก) ขุ 7. 165 ปิฏฺเฐ.} กาลํ โกโธ จิตฺตา\-วิล\-กรณ\-มตฺโต โหติ, น จ ตาว มุขกุลานวิกุลาโน โหติ. อตฺถิ กญฺจิ กาลํ โกโธ มุขกุลานวิกุลานมตฺโต โหติ. น จ ตาว หนุสญฺโจปโน โหติ. อตฺถิ กญฺจิ กาลํ โกโธ หนุสญฺโจปน\-มตฺโต โหติ, น จ ตาว ผรุสวาจํ นิจฺฉารโณ\csromansharedfootnote{420480fc471869bc}{ผรุสวาจนิจฺฉารโณ (สฺยา)} โหติ. อตฺถิ กญฺจิ กาลํ โกโธ ผรุสวาจํ นิจฺฉารณ\-มตฺโต โหติ, น จ ตาว ทิสาวิทิสานุวิโลกโน โหติ. อตฺถิ กญฺจิ กาลํ โกโธ ทิสาวิทิสา\-นุ\-วิโลกน\-มตฺโต โหติ, น จ ตาว ทณฺฑ\-สตฺถ\-ปรา\-มสโน โหติ. อตฺถิ กญฺจิ กาลํ โกโธ ทณฺฑ\-สตฺถ\-ปรามสน\-มตฺโต โหติ, น จ ตาว ทณฺฑ\-สตฺถ\-อพฺภุ\-กฺกิรโณ โหติ. อตฺถิ กญฺจิ กาลํ โกโธ ทณฺฑ\-สตฺถ\-อพฺภุกฺกิรณ\-มตฺโต โหติ, น จ ตาว ทณฺฑสตฺถอภินิปาตโน โหติ. อตฺถิ กญฺจิ กาลํ โกโธ ทณฺฑสตฺถ- อภินิปาตนมตฺโต โหติ, \csromanfolio{61}น จ ตาว สมฺภญฺชน\-ปลิ\-ภ\-ญฺชโน โหติ. อตฺถิ กญฺจิ กาลํ โกโธ สมฺภญฺชน\-ปลิ\-ภ\-ญฺชน\-มตฺโต โหติ, น จ ตาว องฺคมงฺคอปกฑฺฒโน โหติ. อตฺถิ กญฺจิ กาลํ โกโธ องฺคมงฺค\-อปกฑฺฒน\-มตฺโต โหติ, น จ ตาว ชีวิตาโวโรปโน\csromansharedfootnote{c1986c6167e180ca}{ชีวิตปนาสโน (สฺยา) ขุ 7. 166 ปิฏฺเฐ.} โหติ. อตฺถิ กญฺจิ กาลํ โกโธ ชีวิตา\-โวโรปน\-มตฺโต โหติ, น จ ตาว สพฺพ\-จาค\-ปริจฺจาคาย สณฺฐิโต โหติ. ยโต โกโธ ปรํ ปุคฺคลํ ฆาเตตฺวา อตฺตานํ ฆาเตติ. เอตฺตาวตา โกโธ ปรมุ\-สฺสท\-คโต ปรม\-เวปุลฺล\-ปตฺโต โหติ. ยสฺส โส โหติ โกโธ ปหีโน สมุจฺฉินฺโน วูปสนฺโต ปฏิปสฺสทฺโธ อภพฺพุ\-ปฺปตฺติโก ญาณคฺคินา ทฑฺโฒ. โส วุจฺจติ วิธูโม.}

\prose{โกธสฺส ปหีนตฺตา วิธูโม โกธวตฺถุสฺส ปริญฺญาตตฺตา วิธูโม โกธเหตุสฺส ปริญฺญาตตฺตา วิธูโม โกธเหตุสฺส อุปจฺฉินฺนตฺตา วิธูโม. \textbf{อนีโฆ}ติ ราโค นีโฆ, โทโส นีโฆ, โมโห นีโฆ, โกโธ นีโฆ, อุปนาโห นีโฆ ฯเปฯ. สพฺพา\-กุสลา\-ภิสงฺขารา นีฆา, ยสฺเสเต นีฆา ปหีนา สมุจฺฉินฺนา วูปสนฺตา ปฏิปสฺสทฺธา อภพฺพุ\-ปฺปตฺติกา ญาณคฺคินา ทฑฺฒา, โส วุจฺจติ อนีโฆ.}

\prose{\textbf{นิราโส}ติ \textbf{อาสา} วุจฺจติ ตณฺหา. โย ราโค สาราโค ฯเปฯ อภิชฺฌา โลโภ อกุสลมูลํ, ยสฺเสสา อาสา ตณฺหา ปหีนา สมุจฺฉินฺนา วูปสนฺตา ปฏิปสฺสทฺธา อภพฺพุ\-ปฺปตฺติกา ญาณคฺคินา ทฑฺฒา, โส วุจฺจติ นิราโส. \textbf{ชาตี}ติ ยา เตสํ เตสํ สตฺตานํ ตมฺหิ ตมฺหิ สตฺตนิกาเย ชาติ สญฺชาติ โอกฺกนฺติ อภินิพฺพตฺติ ขนฺธานํ ปาตุภาโว อายตนานํ ปฏิลาโภ. \textbf{ชรา}ติ ยา เตสํ เตสํ สตฺตานํ ตมฺหิ ตมฺหิ สตฺตนิกาเย ชรา ชีรณตา ขณฺฑิจฺจํ ปาลิจฺจํ วลิตฺตจตา อายุโน สํหานิ อินฺทฺริยานํ ปริปาโก. \textbf{สนฺโต วิธูโม อนีโฆ นิราโส, อตาริ โส ชาติชรนฺติ พฺรูมี}ติ โย สนฺโต จ วิธูโม จ อนีโฆ จ นิราโส จ, โส ชาติ\-ชรา\-มรณํ อตริ อุตฺตริ ปตริ สมติกฺกมิ วีติวตฺตยีติ พฺรูมิ อาจิกฺขามิ เทเสมิ ปญฺญเปมิ ปฏฺฐเปมิ วิวรามิ วิภชามิ อุตฺตานีกโรมิ ปกาเสมีติ สนฺโต วิธูโม อนีโฆ นิราโส, อตาริ โส ชาติชรนฺติ พฺรูมิ. เตนาห ภควา\csromandash{}}

\prose{\csromanfolio{62}สงฺขาย โลกสฺมิ ปโรปรานิ, (ปุณฺณกาติ ภควา,)}

\prose{ยสฺสิญฺชิตํ นตฺถิ กุหิญฺจิ โลเก.}

\setlength{\csromangathapagewidth}{0pt}
\setlength{\csromangathawakwidth}{0pt}
\csromangathameasuregroup{}{สนฺโต วิธูโม อนีโฆ นิราโส, \\ อตาริ โส ชาติชรนฺติ พฺรูมีติ.}
\csromangathameasurewak{สนฺโต วิธูโม อนีโฆ นิราโส, \\ อตาริ โส ชาติชรนฺติ พฺรูมีติ.}
\setlength{\csromangathaleftcol}{0pt}
\csromangathagroup{\csromangathastanza{\csromangathawak{สนฺโต วิธูโม อนีโฆ นิราโส, \\ อตาริ โส ชาติชรนฺติ พฺรูมีติ.}}}

\prose{สหคาถา\-ปริโยสานา ฯเปฯ ปญฺชลิโก ภควนฺตํ นมสฺสมาโน นิสินฺโน โหติ “สตฺถา เม ภนฺเต ภควา, สาวโกหมสฺมี”ติ.}

\vspace{\tipitakaparskip}
\csromancenter{ปุณฺณก\-มาณว\-ปุจฺฉานิทฺเทโส ตติโย.}
\csromanmatikaanchor{mtk.25}
\csromantocmarkref{subsection}{4. เมตฺตคูมาณวปุจฺฉานิทฺเทส}{mtk.25}
\bookmark[dest={mtk.25},level=2]{4. เมตฺตคูมาณวปุจฺฉานิทฺเทส}
\csromanheader{2}{4. \textbf{เมตฺตคูมาณวปุจฺฉา}\-\textbf{นิทฺเทส}}
\csromanhangingbegin
\hangingheaditem{18}{\textbf{ปุจฺฉามิ ตํ ภควา พฺรูหิ เม’ตํ, (อิจฺจายสฺมา เมตฺตคู,)}}
\hangingbody{\textbf{มญฺญามิ ตํ เวทคู ภาวิตตฺตํ.}}
\hangingbody{\textbf{กุโต นุ ทุกฺขา สมุทาคตา อิเม,}}
\hangingbody{\textbf{เย เกจิ โลกสฺมิ’มเนกรูปา.}}
\csromanhangingend

\prose{\textbf{ปุจฺฉามิตํ ภควา พฺรูหิ เม’ตนฺ}ติ \textbf{ปุจฺฉามี}ติ ติสฺโส ปุจฺฉา อทิฏฺฐโชตนา ปุจฺฉา ทิฏฺฐ\-สํสนฺทนา ปุจฺฉา วิมติจฺเฉทนา ปุจฺฉา. กตมา อทิฏฺฐโชตนา ปุจฺฉา, ปกติยา ลกฺขณํ อญฺญาตํ โหติ อทิฏฺฐํ อตุลิตํ อตีริตํ อวิภูตํ อวิภาวิตํ, ตสฺส ญาณาย ทสฺสนาย ตุลนาย ตีรณาย วิภูตตฺถาย วิภาวน\-ตฺถาย ปญฺหํ ปุจฺฉติ, อยํ อทิฏฺฐโชตนา ปุจฺฉา.}

\prose{กตมา ทิฏฺฐ\-สํสนฺทนา ปุจฺฉา, ปกติยา ลกฺขณํ ญาตํ โหติ ทิฏฺฐํ ตุลิตํ ตีริตํ วิภูตํ วิภาวิตํ, อญฺเญหิ ปณฺฑิเตหิ สทฺธิํ สํสนฺทน\-ตฺถาย ปญฺหํ ปุจฺฉติ, อยํ ทิฏฺฐ\-สํสนฺทนา ปุจฺฉา.}

\prose{กตมา วิมติจฺเฉทนา ปุจฺฉา, ปกติยา สํสย\-ปกฺขนฺโท โหติ วิมติ\-ปกฺขนฺโท เทฺวฬฺหกชาโต “เอวํ นุ โข, น นุ โข กิํ นุ โข กถํ นุ โข”ติ, โส วิมติจฺเฉทน\-ตฺถาย ปญฺหํ ปุจฺฉติ. อยํ วิมติจฺเฉทนา ปุจฺฉา. อิมา ติสฺโส ปุจฺฉา.}

\prose{อปราปิ ติสฺโส ปุจฺฉา มนุสฺสปุจฺฉา อมนุสฺสปุจฺฉา นิมฺมิตปุจฺฉา. กตมา มนุสฺส ปุจฺฉา, มนุสฺสา พุทฺธํ ภควนฺตํ อุปสงฺกมิตฺวา ปญฺหํ ปุจฺฉนฺติ, ภิกฺขู ปุจฺฉนฺติ, ภิกฺขุนิโย ปุจฺฉนฺติ, อุปาสกา ปุจฺฉนฺติ, อุปาสิกาโย ปุจฺฉนฺติ, ราชาโน \csromanfolio{63}ปุจฺฉนฺติ, ขตฺติยา ปุจฺฉนฺติ, พฺราหฺมณา ปุจฺฉนฺติ, เวสฺสา ปุจฺฉนฺติ, สุทฺทา ปุจฺฉนฺติ, คหฏฺฐา ปุจฺฉนฺติ, ปพฺพชิตา ปุจฺฉนฺติ, อยํ มนุสฺสปุจฺฉา.}

\prose{กตมา อมนุสฺสปุจฺฉา, อมนุสฺสา พุทฺธํ ภควนฺตํ อุปสงฺกมิตฺวา ปญฺหํ ปุจฺฉนฺติ, นาคา ปุจฺฉนฺติ, สุปณฺณา ปุจฺฉนฺติ, ยกฺขา ปุจฺฉนฺติ, อสุรา ปุจฺฉนฺติ, คนฺธพฺพา ปุจฺฉนฺติ, มหาราชาโน ปุจฺฉนฺติ, อินฺทา ปุจฺฉนฺติ, พฺรหฺมา ปุจฺฉนฺติ, เทวา ปุจฺฉนฺติ, อยํ อมนุสฺสปุจฺฉา.}

\prose{กตมา นิมฺมิตปุจฺฉา, ภควา รูปํ อภินิมฺมินาติ มโนมยํ สพฺพ\-งฺค\-ปจฺจงฺคํ อหีนินฺทฺริยํ, โส นิมฺมิโต พุทฺธํ ภควนฺตํ อุปสงฺกมิตฺวา ปญฺหํ ปุจฺฉติ, ภควา วิสชฺเชติ, อยํ นิมฺมิตปุจฺฉา. อิมา ติสฺโส ปุจฺฉา.}

\prose{อปราปิ ติสฺโส ปุจฺฉา อตฺตตฺถปุจฺฉา ปรตฺถปุจฺฉา อุภยตฺถ\-ปุจฺฉา.  อปราปิ ติสฺโส ปุจฺฉา ทิฏฺฐธมฺมิก\-ตฺถ\-ปุจฺฉา สมฺปรายิกตฺถ\-ปุจฺฉา ปรมตฺถ\-ปุจฺฉา.  อปราปิ ติสฺโส ปุจฺฉา อนวชฺช\-ตฺถ\-ปุจฺฉา นิกฺกิเลส\-ตฺถ\-ปุจฺฉา โวทาน\-ตฺถ\-ปุจฺฉา.  อปราปิ ติสฺโส ปุจฺฉา อตีตปุจฺฉา อนาคตปุจฺฉา ปจฺจุปฺปนฺน\-ปุจฺฉา.  อปราปิ ติสฺโส ปุจฺฉา อชฺฌตฺตปุจฺฉา พหิทฺธาปุจฺฉา อชฺฌตฺต\-พหิทฺธาปุจฺฉา.  อปราปิ ติสฺโส ปุจฺฉา กุสลปุจฺฉา อกุสลปุจฺฉา อพฺยากตปุจฺฉา.  อราปิ ติสฺโส ปุจฺฉา ขนฺธปุจฺฉา ธาตุปุจฺฉา อายตนปุจฺฉา.  อปราปิ ติสฺโส ปุจฺฉา สติปฏฺฐาน\-ปุจฺฉา สมฺมปฺปธาน\-ปุจฺฉา อิทฺธิปาท\-ปุจฺฉา.  อปราปิ ติสฺโส ปุจฺฉา อินฺทฺริยปุจฺฉา พลปุจฺฉา โพชฺฌงฺค\-ปุจฺฉา.  อปราปิ ติสฺโส ปุจฺฉา มคฺคปุจฺฉา ผลปุจฺฉา นิพฺพานปุจฺฉา.}

\prose{\textbf{ปุจฺฉามิ ตนฺ}ติ ปุจฺฉามิ ตํ ยาจามิ ตํ อชฺเฌสามิ ตํ ปสาเทมิ ตํ “กถยสฺสุ เม”ติ ปุจฺฉามิ ตํ. \textbf{ภควา}ติ คารวาธิวจนเมตํ ฯเปฯ สจฺฉิกา ปญฺญตฺติ, ยทิทํ ภควาติ. \textbf{พฺรูหิ เม ตนฺ}ติ พฺรูหิ อาจิกฺขาหิ เทเสหิ ปญฺญเปหิ ปฏฺฐเปหิ วิวราหิ วิภชาหิ อุตฺตานีกโรหิ ปกาเสหีติ ปุจฺฉามิ ตํ ภควา พฺรูหิ เมตํ.}

\prose{\textbf{อิจฺจายสฺมา เมตฺตคู}ติ อิจฺจาติ ปทสนฺธิ \textbf{ฯเปฯ} อิจฺจายสฺมา เมตฺตคู.}

\prose{\textbf{มญฺญามิ ตํ เวทคู ภาวิตตฺตนฺ}ติ \textbf{เวทคู}ติ ตํ มญฺญามิ, ภาวิตตฺโตติ ตํ มญฺญามิ เอวํ ชานามิ เอวํ อาชานามิ เอวํ ปฏิชานามิ เอวํ ปฏิวิชฺฌามิ. \textbf{เวทคู ภาวิตตฺโต}ติ กถญฺจ ภควา เวทคู, \textbf{เวทา} วุจฺจนฺติ จตูสุ มคฺเคสุ ญาณํ ปญฺญา ปญฺญินฺทฺริยํ ปญฺญาพลํ ฯเปฯ ธมฺมวิจย\-สมฺโพชฺฌงฺโค \csromanfolio{64}วีมสา วิปสฺสนา สมฺมาทิฏฺฐิ. ภควา เตหิ เวเทหิ ชาติ\-ชรา\-มรณสฺส อนฺตคโต อนฺตปฺปตฺโต โกฏิคโต โกฏิปฺปตฺโต ปริยนฺตคโต ปริยนฺต\-ปฺปตฺโต โวสานคโต โวสานปฺปตฺโต ตาณคโต ตาณปฺปตฺโต เลณคโต เลณปฺปตฺโต สรณคโต สรณปฺปตฺโต อภยคโต อภยปฺปตฺโต อจฺจุตคโต อจฺจุตปฺปตฺโต อมตคโต อมตปฺปตฺโต นิพฺพานคโต นิพฺพานปฺปตฺโต. เวทานํ วา อนฺตคโตติ เวทคู. เวเทหิ วา อนฺตคโตติ เวทคู. สตฺตนฺนํ วา ธมฺมานํ วิทิตตฺตา เวทคู. สกฺกายทิฏฺฐิ วิทิตา โหติ. วิจิกิจฺฉา วิทิตา โหติ. สีลพฺพต\-ปรามาโส วิทิโต โหติ. ราโค โทโส โมโห มาโน วิทิโต โหติ. วิทิตาสฺส โหนฺติ ปาปกา อกุสลา ธมฺมา สํกิเลสิกา โปโนภวิกา สทรา ทุกฺขวิปากา อายติํ ชาติ\-ชรา\-มรณิยา.}

\vspace{\tipitakaparskip}
\csromancenter{เวทานิ วิเจยฺย เกวลานิ, (สภิยาติ ภควา,)}
\prose{สมณานํ ยานีธตฺถิ\csromansharedfootnote{26105e250f174363}{ยานิ ปตฺถิ (สฺยา), ยานิ อตฺถิ (ก) ขุ 1. 360 ปิฏฺเฐสุ.} พฺราหฺมณานํ.}

\setlength{\csromangathapagewidth}{0pt}
\setlength{\csromangathawakwidth}{0pt}
\csromangathameasuregroup{}{สพฺพเวทนาสุ วีตราโค, \\ สพฺพํ เวทมติจฺจ เวทคู โสติ.}
\csromangathameasurewak{สพฺพเวทนาสุ วีตราโค, \\ สพฺพํ เวทมติจฺจ เวทคู โสติ.}
\setlength{\csromangathaleftcol}{0pt}
\csromangathagroup{\csromangathastanza{\csromangathawak{สพฺพเวทนาสุ วีตราโค, \\ สพฺพํ เวทมติจฺจ เวทคู โสติ.}}}

\prose{เอวํ ภควา เวทคู.}

\prose{กถํ ภควา ภาวิตตฺโต. ภควา ภาวิตกาโย ภาวิตสีโล ภาวิตจิตฺโต ภาวิตปญฺโญ ภาวิต\-สติปฏฺฐาโน ภาวิต\-สมฺมปฺปธาโน ภาวิตอิทฺธิปาโท ภาวิตอินฺทฺริโย ภาวิตพโล ภาวิต\-โพชฺฌงฺโค ภาวิตมคฺโค ปหีนกิเลโส ปฏิวิทฺธา\-กุปฺโป สจฺฉิกต\-นิโรโธ. ทุกฺขํ ตสฺส ปริญฺญาตํ, สมุทโย ปหีโน, มคฺโค ภาวิโต, นิโรโธ สจฺฉิกโต. อภิญฺเญยฺยํ อภิญฺญาตํ, ปริญฺเญยฺยํ ปริญฺญาตํ, ปหาตพฺพํ ปหีนํ, ภาเวตพฺพํ ภาวิตํ. สจฺฉิกาตพฺพํ สจฺฉิกตํ. อปริตฺโต มหนฺโต คมฺภีโร อปฺปเมยฺโย ทุปฺปริโยคาโฬฺห พหุรตโน สาครูปโม\csromansharedfootnote{be53090be1310263}{สาครสโม (ก)} ฉฬงฺคุ\-เปกฺขาย สมนฺนาคโต โหติ.}

\prose{จกฺขุนา รูปํ ทิสฺวา เนว สุมโน โหติ น ทุมฺมโน. อุเปกฺขโก วิหรติ สโต สมฺปชาโน. โสเตน สทฺทํ สุตฺวา. ฆาเนน คนฺธํ ฆายิตฺวา. ชิวฺหาย รสํ สายิตฺวา. กาเยน โผฏฺฐพฺพํ ผุสิตฺวา. มนสา \csromanfolio{65}ธมฺมํ วิญฺญาย เนว สุมโน โหติ น ทุมฺมโน. อุเปกฺขโก วิหรติ สโต สมฺปชาโน.}

\prose{จกฺขุนา รูปํ ทิสฺวา มนาปํ รูปํ นาภิคิชฺฌติ นาภิหํสติ\csromansharedfootnote{57b1e63a7e92aac1}{นาภิปิหยติ (สฺยา) ขุ 7. 186 ปิฏฺเฐ.} น ราคํ ชเนติ. ตสฺส ฐิโตว กาโย โหติ, ฐิตํ จิตฺตํ อชฺฌตฺตํ สุสณฺฐิตํ สุวิมุตฺตํ. จกฺขุนา โข ปเนว รูปํ ทิสฺวา อมนาปํ น มงฺกุ โหติ อปฺปติฏฺฐิต\-จิตฺโต\csromansharedfootnote{25bd7e83a64c33ef}{อปฺปติฏฺฐีนจิตฺโต (สฺยา)} อลีนมนโส\csromansharedfootnote{b4ad4bcedf097077}{อาทินมนโส (สฺยา) ขุ 7. 186 ปิฏฺเฐปิ.} อพฺยาปนฺน\-เจตโส. ตสฺส ฐิโตว กาโย โหติ, ฐิตํ จิตฺตํ อชฺฌตฺตํ สุสณฺฐิตํ สุวิมุตฺตํ. โสเตน สทฺทํ สุตฺวา. ฆาเนน คนฺธํ ฆายิตฺวา. ชิวฺหาย รสํ สายิตฺวา. กาเยน โผฏฺฐพฺพํ ผุสิตฺวา. มนสา ธมฺมํ วิญฺญาย มนาปํ นาภิคิชฺฌติ นาภิหํสติ น ราคํ ชเนติ. ตสฺส ฐิโตว กาโย โหติ, ฐิตํ จิตฺตํ อชฺฌตฺตํ สุสณฺฐิตํ สุวิมุตฺตํ. มนสา เยว โข ปน ธมฺมํ วิญฺญาย อมนาปํ น มงฺกุ โหติ อปฺปติฏฺฐิต\-จิตฺโต อลีนมนโส อพฺยาปนฺน\-เจตโส. ตสฺส ฐิโตว กาโย โหติ, ฐิตํ จิตฺตํ อชฺฌตฺตํ สุสณฺฐิตํ สุวิมุตฺตํ.}

\prose{จกฺขุนา รูปํ ทิสฺวา มนาปามนาเปสุ รูเปสุ ฐิโตว กาโย โหติ, ฐิตํ จิตฺตํ อชฺฌตฺตํ สุสณฺฐิตํ สุวิมุตฺตํ. โสเตน สทฺทํ สุตฺวา ฯเปฯ. มนสา ธมฺมํ วิญฺญาย มนาปามนาเปสุ ธมฺเมสุ ฐิโตว กาโย โหติ ฐิตํ จิตฺตํ อชฺฌตฺตํ สุสณฺฐิตํ สุวิมุตฺตํ.}

\prose{จกฺขุนา รูปํ ทิสฺวา รชนีเย น รชฺชติ. ทุสฺสนีเย\csromansharedfootnote{c59273e7f1dddc14}{โทสนีเย (สฺยา, ก) ขุ 7. 186 ปิฏฺเฐปิ.} น ทุสฺสติ. โมหนีเย น มุยฺหติ. โกปนีเย น กุปฺปติ. มทนีเย น มชฺชติ. กิเลสนีเย น กิลิสฺสติ. โสเตน สทฺทํ สุตฺวา ฯเปฯ. มนสา ธมฺมํ วิญฺญาย รชนีเย น รชฺชติ. ทุสฺสนีเย น ทุสฺสติ. โมหนีเย น มุยฺหติ. โกปนีเย น กุปฺปติ. มทนีเย น มชฺชติ. กิเลสนีเย น กิลิสฺสติ.}

\prose{ทิฏฺเฐ ทิฏฺฐมตฺโต, สุเต สุตมตฺโต, มุเต มุตมตฺโต, วิญฺญาเต วิญฺญาตมตฺโต, ทิฏฺเฐ น ลิมฺปติ, สุเต น ลิมฺปติ, มุเต น ลิมฺปติ, วิญฺญาเต น ลิมฺปติ, ทิฏฺเฐ อนูปโย\csromansharedfootnote{26751dcdfca87e85}{อนุปโย (สฺยา), อนุสโย (ก) ขุ 7. 187 ปิฏฺเฐ.} อนปาโย อนิสฺสิโต \csromanfolio{66}อปฺปฏิพทฺโธ วิปฺปมุตฺโต วิสญฺญุตฺโต วิมริยาทิกเตน เจตสา วิหรติ. สุเต. มุเต. วิญฺญาเต อนูปโย\csromansharedfootnote{26751dcdfca87e85}{อนุปโย (สฺยา), อนุสโย (ก) ขุ 7. 187 ปิฏฺเฐ.} อนปาโย อนิสฺสิโต อปฺปฏิพทฺโธ วิปฺปมุตฺโต วิสญฺญุตฺโต วิมริยาทิกเตน เจตสา วิหรติ.}

\prose{สํวิชฺชติ ภควโต จกฺขุ, ปสฺสติ ภควา จกฺขุนา รูปํ, ฉนฺทราโค ภควโต นตฺถิ, สุวิมุตฺตจิตฺโต ภควา. สํวิชฺชติ ภควโต โสตํ, สุณาติ ภควา โสเตน สทฺทํ, ฉนฺทราโค ภควโต นตฺถิ, สุวิมุตฺตจิตฺโต ภควา. สํวิชฺชติ ภควโต ฆานํ, ฆายติ ภควา ฆาเนน คนฺธํ, ฉนฺทราโค ภควโต นตฺถิ, สุวิมุตฺตจิตฺโต ภควา. สํวิชฺชติ ภควโต ชิวฺหา, สายติ ภควา ชิวฺหาย รสํ, ฉนฺทราโค ภควโต นตฺถิ, สุวิมุตฺตจิตฺโต ภควา. สํวิชฺชติ ภควโต กาโย, ผุสติ ภควา กาเยน โผฏฺฐพฺพํ, ฉนฺทราโค ภควโต นตฺถิ, สุวิมุตฺตจิตฺโต ภควา. สํวิชฺชติ ภควโต มโน, วิชานาติ ภควา มนสา ธมฺมํ, ฉนฺทราโค ภควโต นตฺถิ, สุวิมุตฺตจิตฺโต ภควา.}

\prose{จกฺขุ รูปารามํ รูปรตํ รูปสมฺมุทิตํ, ตํ ภควโต\csromansharedfootnote{e8091bccd4d00383}{ภควตา (สฺยา) ขุ 7. 187 ปิฏฺเฐ.} ทนฺตํ คุตฺตํ รกฺขิตํ สํวุตํ, ตสฺส จ สํวราย ธมฺมํ เทเสติ. โสตํ สทฺทารามํ สทฺทรตํ. ฆานํ คนฺธารามํ คนฺธรตํ. ชิวฺหา รสารามา รสรตา รสสมฺมุทิตา, สา ภควโต ทนฺตา คุตฺตา รกฺขิตา สํวุตา, ตสฺส จ สํวราย ธมฺมํ เทเสติ. กาโย โผฏฺฐพฺพาราโม โผฏฺฐพฺพรโต โผฏฺฐพฺพ\-สมฺมุทิโต. มโน ธมฺมาราโม ธมฺมรโต ธมฺม\-สมฺมุทิโต, โส ภควโต ทนฺโต คุตฺโต รกฺขิโต สํวุโต, ตสฺส จ สํวราย ธมฺมํ เทเสติ.}

\setlength{\csromangathapagewidth}{0pt}
\setlength{\csromangathawakwidth}{0pt}
\csromangathameasuregroup{“ทนฺตํ นยนฺติ สมิติํ, \\ ทนฺโต เสฏฺโฐ มนุสฺเสสุ, \\ วรมสฺสตรา ทนฺตา, \\ กุญฺชรา จ\textsuperscript{1} มหานาคา, \\ น หิ เอเตหิ ยาเนหิ, \\ ยถา’ตฺตนา สุทนฺเตน, \\ วิธาสุ น วิกมฺปนฺติ, \\ ทนฺตภูมิํ อนุปฺปตฺตา,}{\csromangathabat{“ทนฺตํ นยนฺติ สมิติํ,}{ทนฺตํ ราชาภิรูหติ.} \\ \csromangathabat{ทนฺโต เสฏฺโฐ มนุสฺเสสุ,}{โย’ติวากฺยํ ติติกฺขติ.} \\ \csromangathabat{วรมสฺสตรา ทนฺตา,}{อาชานียา จ\textsuperscript{1} สินฺธวา.} \\ \csromangathabat{กุญฺชรา จ\textsuperscript{1} มหานาคา,}{อตฺตทนฺโต ตโต วรํ.} \\ \csromangathabat{น หิ เอเตหิ ยาเนหิ,}{คจฺเฉยฺย อคตํ ทิสํ.} \\ \csromangathabat{ยถา’ตฺตนา สุทนฺเตน,}{ทนฺโต ทนฺเตน คจฺฉติ.} \\ \csromangathabat{วิธาสุ น วิกมฺปนฺติ,}{วิปฺปมุตฺตา ปุนพฺภวา.} \\ \csromangathabat{ทนฺตภูมิํ อนุปฺปตฺตา,}{เต โลเก วิชิตาวิโน.} \\ ยสฺสินฺทฺริยานิ ภาวิตานิ, \\ อชฺฌตฺตญฺจ พหิทฺธา จ สพฺพโลเก. \\ นิพฺพิชฺฌ อิมํ ปรญฺจ โลกํ, \\ กาลํ กงฺขติ ภาวิโต ส ทนฺโต”ติ\textsuperscript{1}.}
\csromangathameasurewak{ยสฺสินฺทฺริยานิ ภาวิตานิ, \\ อชฺฌตฺตญฺจ พหิทฺธา จ สพฺพโลเก. \\ นิพฺพิชฺฌ อิมํ ปรญฺจ โลกํ, \\ กาลํ กงฺขติ ภาวิโต ส ทนฺโต”ติ\textsuperscript{1}.}
\csromangathasetleft{“ทนฺตํ นยนฺติ สมิติํ, \\ ทนฺโต เสฏฺโฐ มนุสฺเสสุ, \\ วรมสฺสตรา ทนฺตา, \\ กุญฺชรา จ\textsuperscript{1} มหานาคา, \\ น หิ เอเตหิ ยาเนหิ, \\ ยถา’ตฺตนา สุทนฺเตน, \\ วิธาสุ น วิกมฺปนฺติ, \\ ทนฺตภูมิํ อนุปฺปตฺตา,}
\csromangathagroup{\csromangathastanza{\csromangathabat{“ทนฺตํ นยนฺติ สมิติํ,}{ทนฺตํ ราชาภิรูหติ.} \\ \csromangathabat{ทนฺโต เสฏฺโฐ มนุสฺเสสุ,}{โย’ติวากฺยํ ติติกฺขติ.}}\csromangathastanza{\csromangathabat{วรมสฺสตรา ทนฺตา,}{อาชานียา จ\csromansharedfootnote{fd20cbf6dd42968f}{อาชานิยาว (สฺยา) ขุ 1. 59 ปิฏฺเฐ.} สินฺธวา.} \\ \csromangathabat{กุญฺชรา จ\csromansharedfootnote{fd20cbf6dd42968f}{อาชานิยาว (สฺยา) ขุ 1. 59 ปิฏฺเฐ.} มหานาคา,}{อตฺตทนฺโต ตโต วรํ.}}\csromangathastanza{\csromangathabat{น หิ เอเตหิ ยาเนหิ,}{คจฺเฉยฺย อคตํ ทิสํ.} \\ \csromangathabat{ยถา’ตฺตนา สุทนฺเตน,}{ทนฺโต ทนฺเตน คจฺฉติ.}}\csromangathastanza{\csromanfolio{67}\csromangathabat{วิธาสุ น วิกมฺปนฺติ,}{วิปฺปมุตฺตา ปุนพฺภวา.} \\ \csromangathabat{ทนฺตภูมิํ อนุปฺปตฺตา,}{เต โลเก วิชิตาวิโน.}}\csromangathastanza{\csromangathawak{ยสฺสินฺทฺริยานิ ภาวิตานิ, \\ อชฺฌตฺตญฺจ พหิทฺธา จ สพฺพโลเก. \\ นิพฺพิชฺฌ อิมํ ปรญฺจ โลกํ, \\ กาลํ กงฺขติ ภาวิโต ส ทนฺโต”ติ\csromansharedfootnote{b0865415e4fbb8f6}{สุทนฺโตติ (สฺยา) ขุ 1. 358 ปิฏฺเฐ; ขุ 7. 188 ปิฏฺเฐปิ.}.}}}

\prose{เอวํ ภควา ภาวิตตฺโตติ.}

\prose{\textbf{มญฺญามิ ตํ เวทคู ภาวิตตฺตํ.} กุโต นุ ทุกฺขา สมุทาคตา อิเมติ \textbf{กุโต นู}ติ สํสยปุจฺฉา วิมติปุจฺฉา เทฺวฬฺหกปุจฺฉา อเนกํสปุจฺฉา “เอวํ นุ โข, นนุ โข, กิํ นุ โข, กถํ นุ โข”ติ กุโต นุ. \textbf{ทุกฺขา}ติ ชาติทุกฺขํ ชราทุกฺขํ พฺยาธิทุกฺขํ มรณทุกฺขํ โสกปริเทวทุกฺขโทมนสฺสุปายาส\-ทุกฺขํ,พฺยสนํ ทุกฺขํ เนรยิกํ ทุกฺขํ ติรจฺฉาน\-โยนิกํ ทุกฺขํ เปตฺติวิสยิกํ ทุกฺขํ, มานุสิกํ ทุกฺขํ คพฺโภ\-กฺกนฺติ\-มูลกํ ทุกฺขํ คพฺภ\-ฏฺฐิติ\-มูลกํ ทุกฺขํ คพฺภวุฏฺฐาน\-มูลกํ ทุกฺขํ ชาตสฺสู\-ปนิพนฺธกํ ทุกฺขํ ชาตสฺส ปราเธยฺยกํ ทุกฺขํ, อตฺตูปกฺกมํ ทุกฺขํ ปรูปกฺกมํ ทุกฺขํ ทุกฺขทุกฺขํ สงฺขาร\-ทุกฺขํ วิปริณาม\-ทุกฺขํ, จกฺขุโรโค โสตโรโค ฆานโรโค ชิวฺหาโรโค กายโรโค สีสโรโค กณฺณโรโค มุขโรโค ทนฺตโรโค กาโส สาโส ปินาโส ฑาโห ชโร กุจฺฉิโรโค มุจฺฉา ปกฺขนฺทิกา สูลา วิสูจิกา กุฏฺฐํ คณฺโฑ กิลาโส โสโส อปมาโร ททฺทุ กณฺฑุ กจฺฉุ รขสา วิตจฺฉิกา โลหิตปิตฺตํ มธุเมโห อํสา ปิฬกา ภคนฺทลา ปิตฺต\-สมุฏฺฐานา อาพาธา เสมฺห\-สมุฏฺฐานา อาพาธา วาตสมุฏฺฐานา อาพาธา สนฺนิปาติกา อาพาธา อุตุปริณามชา อาพาธา วิสม\-ปริหารชา อาพาธา โอปกฺกมิกา อาพาธา กมฺมวิปากชา อาพาธา สีตํ อุณฺหํ ชิฆจฺฉา ปิปาสา อุจฺจาโร ปสฺสาโว ฑํสมกส\-วาตาตป\-สรีสป- สมฺผสฺสํ ทุกฺขํ, มาตุมรณํ ทุกฺขํ ปิตุมรณํ ทุกฺขํ ภาตุมรณํ ทุกฺขํ ภคินิมรณํ ทุกฺขํ ปุตฺตมรณํ ทุกฺขํ ธีตุมรณํ ทุกฺขํ ญาติพฺยสนํ ทุกฺขํ โรคพฺยสนํ ทุกฺขํ โภคพฺยสนํ ทุกฺขํ สีลพฺยสนํ ทุกฺขํ ทิฏฺฐิพฺยสนํ ทุกฺขํ, เยสํ ธมฺมานํ อาทิโต สมุทาคมนํ ปญฺญายติ, อตฺถงฺคมโต นิโรโธ ปญฺญายติ, กมฺม\-สนฺนิสฺสิโต วิปาโก, วิปาก\-สนฺนิสฺสิตํ กมฺมํ, นาม\-สนฺนิสฺสิตํ รูปํ,}

\prose{\csromanfolio{68}รูปสนฺนิสฺสิตํ นามํ, ชาติยา อนุคตํ, ชราย อนุสฏํ, พฺยาธินา อภิภูตํ, มรเณน อพฺภาหตํ, ทุกฺเข ปติฏฺฐิตํ, อตาณํ อเลณํ อสรณํ อสรณีภูตํ, อิเม วุจฺจนฺติ ทุกฺขา. อิเม ทุกฺขา กุโต สมุทาคตา กุโต ชาตา กุโต สญฺชาตา กุโต นิพฺพตฺตา กุโต อภินิพฺพตฺตา กุโต ปาตุภูตา, กิํนิทานา กิํสมุทยา กิํชาติกา กิํปภวาติ อิเมสํ ทุกฺขานํ มูลํ ปุจฺฉติ, เหตุํ ปุจฺฉติ, นิทานํ ปุจฺฉติ, สมฺภวํ ปุจฺฉติ, ปภวํ ปุจฺฉติ, สมุฏฺฐานํ ปุจฺฉติ, อาหารํ ปุจฺฉติ, อารมฺมณํ ปุจฺฉติ, ปจฺจยํ ปุจฺฉติ, สมุทยํ ปุจฺฉติ ปปุจฺฉติ ยาจติ อชฺเฌสติ ปสาเทตีติ กุโต นุ ทุกฺขา สมุทาคตา อิเม.}

\prose{\textbf{เย เกจิ โลกสฺมิ’มเนกรูปา}ติ \textbf{เย เกจี}ติ สพฺเพน สพฺพํ สพฺพถา สพฺพํ อเสสํ นิสฺเสสํ ปริยาทิย\-น\-วจนเมตํ เย เกจีติ. \textbf{โลกสฺมินฺ}ติ อปายโลเก มนุสฺสโลเก เทวโลเก ขนฺธโลเก ธาตุโลเก อายตนโลเก. \textbf{อเนกรูปา}ติ อเนกวิธา นานาปฺปการา ทุกฺขาติ เย เกจิ โลกสฺมิ’มเนกรูปา. เตนาห โส พฺราหฺมโณ\csromandash{}}

\prose{ปุจฺฉามิ ตํ ภควา พฺรูหิ เม’ตํ, (อิจฺจายสฺมา เมตฺตคู,)}

\prose{มญฺญามิ ตํ เวทคู ภาวิตตฺตํ.}

\prose{กุโต นุ ทุกฺขา สมุทาคตา อิเม,}

\prose{เย เกจิ โลกสฺมิ’มเนกรูปาติ.}

\csromanhangingbegin
\hangingheaditem{19}{\textbf{ทุกฺขสฺส เว มํ ปภวํ อปุจฺฉสิ, (เมตฺตคูติ ภควา,)}}
\hangingbody{\textbf{ตํ เต ปวกฺขามิ ยถา ปชานํ.}}
\hangingbody{\textbf{อุปธินิทานา ปภวนฺติ ทุกฺขา,}}
\hangingbody{\textbf{เย เกจิ โลกสฺมิ’มเนกรูปา.}}
\csromanhangingend

\prose{\textbf{ทุกฺขสฺส เว มํ ปภวํ อปุจฺฉสี}ติ \textbf{ทุกฺขสฺสา}ติ ชาติทุกฺขสฺส ชราทุกฺขสฺส พฺยาธิ\-ทุกฺขสฺส มรณ\-ทุกฺขสฺส โสกปริเทวทุกฺขโทมนสฺสุปายาส\-ทุกฺขสฺส. \textbf{ปภวํ อปุจฺฉสี}ติ ทุกฺขสฺส มูลํ ปุจฺฉสิ, เหตุํ ปุจฺฉสิ, นิทานํ ปุจฺฉสิ, สมฺภวํ ปุจฺฉสิ, ปภวํ ปุจฺฉสิ, สมุฏฺฐานํ ปุจฺฉสิ, อาหารํ ปุจฺฉสิ, อารมฺมณํ ปุจฺฉสิ, ปจฺจยํ ปุจฺฉสิ, สมุทยํ ปุจฺฉสิ ยาจสิ อชฺเฌสสิ ปสาเทสีติ ทุกฺขสฺส เว มํ ปภวํ อปุจฺฉสิ. \textbf{เมตฺตคู}ติ ภควา ตํ พฺราหฺมณํ นาเมน อาลปติ. \csromanfolio{69}\textbf{ภควา}ติ คารวาธิวจนเมตํ ฯเปฯ สจฺฉิกา ปญฺญตฺติ, ยทิทํ ภควาติ เมตฺตคูติ ภควา.}

\prose{\textbf{ตํ เต ปวกฺขามิ ยถา ปชานนฺ}ติ \textbf{ตนฺ}ติ ทุกฺขสฺส มูลํ ปวกฺขามิ, เหตุํ ปวกฺขามิ, นิทานํ ปวกฺขามิ, สมฺภวํ ปวกฺขามิ, ปภวํ ปวกฺขามิ, สมุฏฺฐานํ ปวกฺขามิ, อาหารํ ปวกฺขามิ, อารมฺมณํ ปวกฺขามิ, ปจฺจยํ ปวกฺขามิ, สมุทยํ ปวกฺขามิ อาจิกฺขิสฺสามิ เทเสสฺสามิ ปญฺญเปสฺสามิ ปฏฺฐเปสฺสามิ วิวริสฺสามิ วิภชิสฺสามิ อุตฺตานีกริสฺสามิ ปกาเสสฺสามีติ ตํ เต ปวกฺขามิ. \textbf{ยถา ปชานนฺ}ติ ยถา ปชานนฺโต อาชานนฺโต วิชานนฺโต ปฏิวิชานนฺโต ปฏิวิชฺฌนฺโต. น อิติหีติหํ น อิติกิราย น ปรมฺปราย น ปิฏก\-สมฺปทาย\csromansharedfootnote{a50a0df1ce234e58}{น ปิฏกสมฺปทาเนน (ก) ขุ 7. 280 ปิฏฺเฐ.} น ตกฺกเหตุ น นยเหตุ น อาการ\-ปริวิตกฺเกน น ทิฏฺฐิ\-นิชฺฌาน\-กฺขนฺติยา สามํ สยม\-ภิญฺญาตํ อตฺต\-ปจฺจกฺขธมฺมํ, ตํ กถยิสฺสามีติ ตํ เต ปวกฺขามิ ยถา ปชานํ.}

\prose{\textbf{อุปธินิทานา ปภวนฺติ ทุกฺขา}ติ อุปธีติ ทส \textbf{อุปธี}, ตณฺหูปธิ, ทิฏฺฐูปธิ, กิเลสูปธิ, กมฺมูปธิ, ทุจฺจริตูปธิ, อาหารูปธิ, ปฏิฆูปธิ, จตสฺโส อุปาทินฺน\-ธาตุโย อุปธี, ฉ อชฺฌตฺติกานิ อายตนานิ อุปธี, ฉ วิญฺญาณกายา อุปธี, สพฺพมฺปิ ทุกฺขํ ทุกฺขม\-นฏฺเฐน\csromansharedfootnote{aeb6c367c3a78114}{ทุกฺขฏฺเฐน (สฺยา)} อุปธิ, อิเม วุจฺจนฺติ ทส อุปธี. \textbf{ทุกฺขา}ติ ชาติทุกฺขํ ชราทุกฺขํ พฺยาธิทุกฺขํ มรณทุกฺขํ โสกปริเทวทุกฺขโทมนสฺสุปายาส- ทุกฺขํ เนรยิกํ ทุกฺขํ ฯเปฯ ทิฏฺฐิพฺยสนํ ทุกฺขํ. เยสํ ธมฺมานํ อาทิโต สมุทาคมนํ ปญฺญายติ, อตฺถงฺคมโต นิโรโธ ปญฺญายติ, กมฺม\-สนฺนิสฺสิโต วิปาโก, วิปาก\-สนฺนิสฺสิตํ กมฺมํ, นาม\-สนฺนิสฺสิตํ รูปํ. รูปสนฺนิสฺสิตํ นามํ. ชาติยา อนุคตํ ชราย อนุสฏํ, พฺยาธินา อภิภูตํ, มรเณน อพฺภาหตํ, ทุกฺเข ปติฏฺฐิตํ อตาณํ อเลณํ อสรณํ อสรณีภูตํ. อิเม วุจฺจนฺติ ทุกฺขา. อิเม ทุกฺขา อุปธินิทานา อุปธิเหตุกา อุปธิปจฺจยา อุปธิการณา โหนฺติ ปภวนฺติ สมฺภวนฺติ ชายนฺติ สญฺชายนฺติ นิพฺพตฺตนฺติ ปาตุภวนฺตีติ อุปธินิทานา ปภวนฺติ ทุกฺขา.}

\prose{\textbf{เย เกจิ โลกสฺมิ’มเนกรูปา}ติ \textbf{เย เกจี}ติ สพฺเพน สพฺพํ สพฺพถา สพฺพํ อเสสํ นิสฺเสสํ ปริยาทิย\-น\-วจนเมตํ เย เกจีติ. \textbf{โลกสฺมินฺ}ติ อปายโลเก มนุสฺสโลเก เทวโลเก ขนฺธโลเก \csromanfolio{70}ธาตุโลเก อายตนโลเก. \textbf{อเนกรูปา}ติ อเนกวิธา นานปฺปการา ทุกฺขาติ เย เกจิ โลกสฺมิ’มเนกรูปา. เตนาห ภควา\csromandash{}}

\proseitem{4}{ทุกฺขสฺส เว มํ ปภวํ อปุจฺฉสิ, (เมตฺตคูติ ภควา,)}

\prose{ตํ เต ปวกฺขามิ ยถา ปชานํ.}

\prose{อุปธินิทานา ปภวนฺติ ทุกฺขา.}

\prose{เย เกจิ โลกสฺมิ’มเนกรูปาติ.}

\setlength{\csromangathapagewidth}{0pt}
\setlength{\csromangathawakwidth}{0pt}
\csromangathameasuregroup{}{\textbf{โย เว อวิทฺวา อุปธิํ กโรติ,} \\ \textbf{ปุนปฺปุนํ ทุกฺข’มุเปติ มนฺโท.} \\ \textbf{ตสฺมา ปชานํ อุปธิํ น กยิรา,} \\ \textbf{ทุกฺขสฺส ชาติปฺปภวา}\textbf{นุปสฺสี.}}
\csromangathameasurewak{\textbf{โย เว อวิทฺวา อุปธิํ กโรติ,} \\ \textbf{ปุนปฺปุนํ ทุกฺข’มุเปติ มนฺโท.} \\ \textbf{ตสฺมา ปชานํ อุปธิํ น กยิรา,} \\ \textbf{ทุกฺขสฺส ชาติปฺปภวา}\textbf{นุปสฺสี.}}
\setlength{\csromangathaleftcol}{0pt}
\csromangathagroup{\csromangathastanza{\csromangathawak{\csromangathaitemline{20}{\textbf{โย เว อวิทฺวา อุปธิํ กโรติ,}} \\ \csromangathacontline{\textbf{ปุนปฺปุนํ ทุกฺข’มุเปติ มนฺโท.}} \\ \csromangathacontline{\textbf{ตสฺมา ปชานํ อุปธิํ น กยิรา,}} \\ \csromangathacontline{\textbf{ทุกฺขสฺส ชาติปฺปภวา}\-\textbf{นุปสฺสี.}}}}}

\prose{\textbf{โย เว อวิทฺวา อุปธิํ กโรตี}ติ \textbf{โย}ติ โย ยาทิโส ยถายุตฺโต ยถาวิหิโต ยถาปกาโร ยํฐานปฺปตฺโต ยํธมฺมสมนฺนาคโต ขตฺติโย วา พฺราหฺมโณ วา เวสฺโส วา สุทฺโท วา คหฏฺโฐ วา ปพฺพชิโต วา เทโว วา มนุสฺโส วา. \textbf{อวิทฺวา}ติ อวิชฺชาคโต อญฺญาณี อวิภาวี ทุปฺปญฺโญ. \textbf{อุปธิํ กโรตี}ติ ตณฺหูปธิํ กโรติ, ทิฏฺฐูปธิํ กโรติ, กิเลสูปธิํ กโรติ, กมฺมูปธิํ กโรติ, ทุจฺจริตูปธิํ กโรติ, อาหารูปธิํ กโรติ, ปฏิฆูปธิํ กโรติ. จตสฺโส อุปาทินฺน\-ธาตุโย อุปธี กโรติ, ฉ อชฺฌตฺติกานิ อายตนานิ อุปธี กโรติ, ฉ วิญฺญาณกาเย อุปธี กโรติ ชเนติ สญฺชเนติ นิพฺพตฺเตติ อภินิพฺพตฺเตตีติ อวิทฺวา อุปธิํ กโรติ.}

\prose{\textbf{ปุนปฺปุนํ ทุกฺข’มุเปติ มนฺโท}ติ ปุนปฺปุนํ ชาติทุกฺขํ ชราทุกฺขํ พฺยาธิทุกฺขํ มรณทุกฺขํ โสกปริเทว\-ทุกฺข- โทมนสฺสุ\-ปายา\-สทุกฺขํ เอติ สมุเปติ อุปคจฺฉติ คณฺหาติ ปรามสติ อภินิวิสตีติ ปุนปฺปุนํ ทุกฺข’มุเปติ. \textbf{มนฺโท}ติ มนฺโท โมมุโห อวิทฺวา อวิชฺชาคโต อญฺญาณี อวิภาวี ทุปฺปญฺโญติ ปุนปฺปุนํ ทุกฺข’มุเปติ มนฺโท.}

\prose{\textbf{ตสฺมา ปชานํ อุปธิํ น กยิรา}ติ \textbf{ตสฺมา}ติ ตํการณา ตํเหตุ ตปฺปจฺจยา ตํนิทานา เอตํ อาทีนวํสมฺปสฺสมาโน อุปธีสูติ ตสฺมา. \textbf{ปชานนฺ}ติ ปชานนฺโต อาชานนฺโต วิชานนฺโต ปฏิวิชานนฺโต ปฏิวิชฺฌนฺโต, “สพฺเพ สงฺขารา อนิจฺจา”ติ ปชานนฺโต อาชานนฺโต วิชานนฺโต ปฏิวิชานนฺโต \csromanfolio{71}ปฏิวิชฺฌนฺโต, “สพฺเพ สงฺขารา ทุกฺขา”ติ ฯเปฯ “สพฺเพ ธมฺมา อนตฺตา”ติ ฯเปฯ. “ยํ กิญฺจิ สมุทย\-ธมฺมํ, สพฺพํ ตํ นิโรธธมฺมนฺ”ติ ปชานนฺโต อาชานนฺโต วิชานนฺโต ปฏิวิชานนฺโต ปฏิวิชฺฌนฺโต. \textbf{อุปธิํ น กยิรา}ติ \textbf{ตณฺหูปธิํ น กเรยฺย, ทิฏฺฐูปธิํ น กเรยฺย, กิเลสูปธิํ น} กเรยฺย, ทุจฺจริตูปธิํ น กเรยฺย, อาหารูปธิํ น กเรยฺย, ปฏิฆูปธิํ น กเรยฺย, จตสฺโส อุปาทินฺน\-ธาตุโย อุปธี น กเรยฺย, ฉ อชฺฌตฺติกานิ อายตนานิ อุปธี น กเรยฺย, ฉ วิญฺญาณกาเย อุปธี น \textbf{กเรยฺย น ชเนยฺย น สญฺชเนยฺย น นิพฺพตฺเตยฺย นาภินิพฺพตฺเตยฺยาติ ตสฺมา} \textbf{ปชานํ อุปธิํ น กยิรา.}}

\prose{\textbf{ทุกฺขสฺสา}ติ ชาติทุกฺขสฺส ชราทุกฺขสฺส พฺยาธิ\-ทุกฺขสฺส มรณ\-ทุกฺขสฺส โสกปริเทวทุกฺขโทมนสฺสุปายาส\-ทุกฺขสฺส. \textbf{ปภวา}\-\textbf{นุปสฺสี}ติ ทุกฺขสฺส มูลานุปสฺสี เหตานุปสฺสี นิทานานุปสฺสี สมฺภวา\-นุปสฺสี ปภวานุปสฺสี สมุฏฺฐา\-นานุปสฺสี อาหารานุปสฺสี อารมฺมณา\-นุปสฺสี ปจฺจยานุปสฺสี สมุทยา\-นุปสฺสี. \textbf{อนุปสฺสนา} วุจฺจติ ญาณํ. ยา ปญฺญา ปชานนา ฯเปฯ อโมโห ธมฺมวิจโย สมฺมาทิฏฺฐิ. อิมาย อนุปสฺสนาย ปญฺญาย อุเปโต โหติ สมุเปโต อุปาคโต สมุปาคโต อุปปนฺโน สมุปปนฺโน สมนฺนาคโต, โส วุจฺจติ อนุปสฺสีติ ทุกฺขสฺส ชาติปฺปภวา\-นุปสฺสี. เตนาห ภควา\csromandash{}}

\setlength{\csromangathapagewidth}{0pt}
\setlength{\csromangathawakwidth}{0pt}
\csromangathameasuregroup{}{โย เว อวิทฺวา อุปธิํ กโรติ, \\ ปุนปฺปุนํ ทุกฺข’มุเปติ มนฺโท. \\ ตสฺมา ปชานํ อุปธิํ น กยิรา, \\ ทุกฺขสฺส ชาติปฺปภวานุปสฺสีติ.}
\csromangathameasurewak{โย เว อวิทฺวา อุปธิํ กโรติ, \\ ปุนปฺปุนํ ทุกฺข’มุเปติ มนฺโท. \\ ตสฺมา ปชานํ อุปธิํ น กยิรา, \\ ทุกฺขสฺส ชาติปฺปภวานุปสฺสีติ.}
\setlength{\csromangathaleftcol}{0pt}
\csromangathagroup{\csromangathastanza{\csromangathawak{โย เว อวิทฺวา อุปธิํ กโรติ, \\ ปุนปฺปุนํ ทุกฺข’มุเปติ มนฺโท. \\ ตสฺมา ปชานํ อุปธิํ น กยิรา, \\ ทุกฺขสฺส ชาติปฺปภวานุปสฺสีติ.}}}

\proseitem{21}{\textbf{ยํ ตํ อปุจฺฉิมฺห อกิตฺตยี โน,}}

\prose{\textbf{อญฺญํ ตํ ปุจฺฉาม ตทิงฺฆ พฺรูหิ.}}

\setlength{\csromangathapagewidth}{0pt}
\setlength{\csromangathawakwidth}{0pt}
\csromangathameasuregroup{}{\textbf{กถํ นุ ธีรา วิตรนฺติ โอฆํ,} \\ \textbf{ชาติํ ชรํ โสกปริทฺทวญฺจ.} \\ \textbf{ตํ เม มุนี สาธุ วิยากโรหิ,}}
\csromangathameasurewak{\textbf{กถํ นุ ธีรา วิตรนฺติ โอฆํ,} \\ \textbf{ชาติํ ชรํ โสกปริทฺทวญฺจ.} \\ \textbf{ตํ เม มุนี สาธุ วิยากโรหิ,}}
\setlength{\csromangathaleftcol}{0pt}
\csromangathagroup{\csromangathastanza{\csromangathawak{\textbf{กถํ นุ ธีรา วิตรนฺติ โอฆํ,} \\ \textbf{ชาติํ ชรํ โสกปริทฺทวญฺจ.} \\ \textbf{ตํ เม มุนี สาธุ วิยากโรหิ,}}}}

\prose{\textbf{ยํ ตํ อปุจฺฉิมฺห อกิตฺตยี โน}ติ ยํ ตํ อปุจฺฉิมฺห อยาจิมฺห อชฺเฌสิมฺห ปสาทิมฺห. \textbf{อกิตฺตยี โน}ติ กิตฺติตํ\csromansharedfootnote{089016ce156e0143}{อกิตฺติ ตํ (สฺยา) เอวมีทิเสสุ ปเทสุ อตีตวิภตฺติวเสน. ขุ 7. 217 ปิฏฺเฐ.} ปกิตฺติตํ อาจิกฺขิตํ เทสิตํ \csromanfolio{72}ปญฺญปิตํ\csromansharedfootnote{a989c87cb4360486}{ปญฺญาปิตํ (ก)} ปฏฺฐปิตํ วิวริตํ วิภตฺตํ อุตฺตานีกตํ ปกาสิตนฺติ ยํ ตํ อปุจฺฉิมฺห อกิตฺตยี โน.}

\prose{\textbf{อญฺญํ ตํ ปุจฺฉาม ตทิงฺฆ พฺรูหี}ติ อญฺญํ ตํ ปุจฺฉาม, อญฺญํ ตํ ยาจาม, อญฺญํ ตํ อชฺเฌสาม, อญฺญํ ตํ ปสาเทม, อุตฺตริ ตํ ปุจฺฉาม. \textbf{ตทิงฺฆ พฺรูหี}ติ อิงฺฆ พฺรูหิ อาจิกฺขาหิ เทเสหิ ปญฺญเปหิ ปฏฺฐเปหิ วิวราหิ วิภชาหิ อุตฺตานีกโรหิ ปกาเสหีติ อญฺญํ ตํ ปุจฺฉาม ตทิงฺฆ พฺรูหิ.}

\prose{\textbf{กถํ นุ ธีรา วิตรนฺติ โอฆํ, ชาติํ ชรํ โสกปริทฺทว}\-\textbf{ญฺจา}ติ \textbf{กถํ นู}ติ สํสยปุจฺฉา วิมติปุจฺฉา เทฺวฬฺหกปุจฺฉา อเนกํสปุจฺฉา, “เอวํ นุ โข, นนุ โข, กิํ นุ โข, กถํ นุ โข”ติ กถํ นุ. \textbf{ธีรา}ติ ธีรา ปณฺฑิตา ปญฺญวนฺโต พุทฺธิมนฺโต ญาณิโน วิภาวิโน เมธาวิโน. \textbf{โอฆนฺ}ติ กาโมฆํ ภโวฆํ ทิฏฺโฐฆํ อวิชฺโชฆํ. \textbf{ชาตี}ติ ยา เตสํ เตสํ สตฺตานํ ตมฺหิ ตมฺหิ สตฺตนิกาเย ชาติ สญฺชาติ โอกฺกนฺติ นิพฺพตฺติ อภินิพฺพตฺติ ขนฺธานํ ปาตุภาโว อายตนานํ ปฏิลาโภ. \textbf{ชรา}ติ ยา เตสํ เตสํ สตฺตานํ ตมฺหิ ตมฺหิ สตฺตนิกาเย ชรา ชีรณตา ขณฺฑิจฺจํ ปาลิจฺจํ วลิตฺตจตา อายุโน สํหานิ อินฺทฺริยานํ ปริปาโก. \textbf{โสโก}ติ ญาติพฺยสเนน วา ผุฏฺฐสฺส โภคพฺยสเนน วา ผุฏฺฐสฺส โรคพฺยสเนน วา ผุฏฺฐสฺส สีลพฺยสเนน วา ผุฏฺฐสฺส ทิฏฺฐิ\-พฺยสเนน วา ผุฏฺฐสฺส อญฺญตร\-ญฺญตเรน พฺยสเนน วา สมนฺนาคตสฺส อญฺญตร\-ญฺญตเรน ทุกฺขธมฺเมน วา ผุฏฺฐสฺส โสโก โสจนา โสจิตตฺตํ อนฺโตโสโก อนฺโต ปริโสโก อนฺโตฑาโห อนฺโตปริฑาโห เจตโส ปริชฺฌายนา โทมนสฺสํ โสกสลฺลํ. \textbf{ปริเทโว}ติ ญาติพฺยสเนน วา ผุฏฺฐสฺส โภคพฺยสเนน วา ผุฏฺฐสฺส โรคพฺยสเนน วา ผุฏฺฐสฺส สีลพฺยสเนน วา ผุฏฺฐสฺส ทิฏฺฐิ\-พฺยสเนน วา ผุฏฺฐสฺส อญฺญตร\-ญฺญตเรน พฺยสเนน วา สมนฺนาคตสฺส อญฺญตร\-ญฺญตเรน ทุกฺขธมฺเมน วา ผุฏฺฐสฺส อาเทโว ปริเทโว อาเทวนา ปริเทวนา อาเทวิตตฺตํ ปริเทวิตตฺตํ วาจา ปลาโป\csromansharedfootnote{5dee2baf03761ab6}{ลาโป ปลาโป (สฺยา) อภิ 1. 105 ปิฏฺเฐปิ.} วิปฺปลาโป ลาลปฺโป ลาลปฺปนา ลาลปฺปิตตฺตํ\csromansharedfootnote{05f183904a46676b}{ลาลปฺปายนา ลาลปฺปายิตตฺตํ (พหูสุ) ชราสุตฺตนิทฺเทสฏฺฐกถา โอโลเกตพฺพา.}.}

\prose{\textbf{กถํ นุ ธีรา วิตรนฺติ โอฆํ, ชาติํ ชรํ โสกปริทฺทว}\-\textbf{ญฺจา}ติ ธีรา กถํ โอฆญฺจ ชาติญฺจ ชรญฺจ โสกญฺจ ปริเทวญฺจ ตรนฺติ อุตฺตรนฺติ ปตรนฺติ \csromanfolio{73}สมติกฺกมนฺติ วีติวตฺตนฺตีติ กถํ นุ ธีรา วิตรนฺติ โอฆํ, ชาติํ ชรํ โสกปริทฺทวญฺจ.}

\prose{\textbf{ตํ เม มุนี สาธุ วิยากโรหี}ติ \textbf{ตนฺ}ติ ยํ ปุจฺฉามิ, ยํ ยาจามิ, ยํ อชฺเฌสามิ, ยํ ปสาเทมิ. \textbf{มุนี}ติ โมนํ วุจฺจติ ญาณํ. ยา ปญฺญา ปชานนา ฯเปฯ อโมโห ธมฺมวิจโย สมฺมาทิฏฺฐิ. ภควา เตน ญาเณน สมนฺนาคโต มุนิ โมนปฺปตฺโต. ตีณิ โมเนยฺยานิ กายโมเนยฺยํ วจีโมเนยฺยํ มโนโมเนยฺยํ.}

\prose{กตมํ กายโมเนยฺยํ, ติวิธานํ กาย\-ทุจฺจริตานํ ปหานํ กายโมเนยฺยํ, ติวิธํ กายสุจริตํ กายโมเนยฺยํ, กายารมฺมเณ ญาณํ กายโมเนยฺยํ, กายปริญฺญา กายโมเนยฺยํ, ปริญฺญา\-สหคโต มคฺโค กายโมเนยฺยํ, กาเย ฉนฺทราคสฺส ปหานํ กายโมเนยฺยํ, กาย\-สงฺขาร\-นิโรโธ จตุตฺถชฺฌาน\-สมาปตฺติ กายโมเนยฺยํ. อิทํ กายโมเนยฺยํ.}

\prose{กตมํ วจีโมเนยฺยํ, จตุพฺพิธานํ วจีทุจฺจริตานํ ปหานํ วจีโมเนยฺยํ, จตุพฺพิธํ วจีสุจริตํ วจีโมเนยฺยํ, วาจารมฺมเณ ญาณํ วจีโมเนยฺยํ, วาจาปริญฺญา วจีโมเนยฺยํ, ปริญฺญา\-สหคโต มคฺโค วจีโมเนยฺยํ, วาจาย ฉนฺทราคสฺส ปหานํ วจีโมเนยฺยํ, วจีสงฺขาร\-นิโรโธ ทุติย\-ชฺฌาน\-สมาปตฺติ วจีโมเนยฺยํ. อิทํ วจีโมเนยฺยํ.}

\prose{กตมํ มโนโมเนยฺยํ, ติวิธานํ มโน\-ทุจฺจริตานํ ปหานํ มโนโมเนยฺยํ, ติวิธํ มโนสุจริตํ มโนโมเนยฺยํ, จิตฺตารมฺมเณ ญาณํ มโนโมเนยฺยํ, จิตฺตปริญฺญา มโนโมเนยฺยํ, ปริญฺญา\-สหคโต มคฺโค มโนโมเนยฺยํ, จิตฺเต ฉนฺทราคสฺส ปหานํ มโนโมเนยฺยํ, จิตฺต\-สงฺขาร\-นิโรโธ สญฺญา\-เวทยิต\-นิโรธ\-สมาปตฺติ มโนโมเนยฺยํ. อิทํ มโนโมเนยฺยํ.}

\setlength{\csromangathapagewidth}{0pt}
\setlength{\csromangathawakwidth}{0pt}
\csromangathameasuregroup{กายมุนิํ วจีมุนิํ\textsuperscript{1}, \\ มุนิํ โมเนยฺยสมฺปนฺนํ, \\ กายมุนิํ วจีมุนิํ, \\ มุนิํ โมเนยฺยสมฺปนฺนํ,}{\csromangathabat{กายมุนิํ วจีมุนิํ\textsuperscript{1},}{มโนมุนิมนาสวํ.} \\ \csromangathabat{มุนิํ โมเนยฺยสมฺปนฺนํ,}{อาหุ สพฺพปฺปหายินํ.} \\ \csromangathabat{กายมุนิํ วจีมุนิํ,}{มโนมุนิมนาสวํ.} \\ \csromangathabat{มุนิํ โมเนยฺยสมฺปนฺนํ,}{อาหุ นินฺหาตปาปกนฺติ.}}
\csromangathasetleft{กายมุนิํ วจีมุนิํ\textsuperscript{1}, \\ มุนิํ โมเนยฺยสมฺปนฺนํ, \\ กายมุนิํ วจีมุนิํ, \\ มุนิํ โมเนยฺยสมฺปนฺนํ,}
\csromangathagroup{\csromangathastanza{\csromangathabat{กายมุนิํ วจีมุนิํ\csromansharedfootnote{260f2f555c7ea6be}{วาจามุนิํ (พหูสุ) ขุ 1. 234 ปิฏฺเฐ.},}{มโนมุนิม\-นาสวํ.} \\ \csromangathabat{มุนิํ โมเนยฺย\-สมฺปนฺนํ,}{อาหุ สพฺพ\-ปฺปหายินํ.}}\csromangathastanza{\csromangathabat{กายมุนิํ วจีมุนิํ,}{มโนมุนิม\-นาสวํ.} \\ \csromangathabat{มุนิํ โมเนยฺย\-สมฺปนฺนํ,}{อาหุ นินฺหาต\-ปาปกนฺติ.}}}

\proseitem{4}{\csromanfolio{74}อิเมหิ ตีหิ โมเนยฺเยหิ ธมฺเมหิ สมนฺนาคตา, ฉ มุนิโน\csromansharedfootnote{422a37c497cc58e3}{มุนโย (สฺยา) ขุ 7. 44 ปิฏฺเฐปิ.} อคารมุนิโน อนคารมุนิโน เสขมุนิโน\csromansharedfootnote{87cbadba4d7604b1}{เสกฺขมุนิโน (สฺยา, ก)} อเสขมุนิโน ปจฺเจกมุนิโน มุนิมุนิโนติ. กตเม อคารมุนิโน, เย เต อคาริกา ทิฏฺฐปทา วิญฺญาตสาสนา, อิเม อคารมุนิโน.  กตเม อนคารมุนิโน, เย เต ปพฺพชิตา ทิฏฺฐปทา วิญฺญาตสาสนา, อิเม อนคารมุนิโน. สตฺต เสขา เสขมุนิโน. อรหนฺโต อเสขมุนิโน. ปจฺเจก\-สมฺพุทฺธา ปจฺเจกมุนิโน. ตถาคตา อรหนฺโต สมฺมาสมฺพุทฺธา มุนิมุนิโน.}

\setlength{\csromangathapagewidth}{0pt}
\setlength{\csromangathawakwidth}{0pt}
\csromangathameasuregroup{น โมเนน มุนี\textsuperscript{1} โหติ, \\ โย จ ตุลํว ปคฺคยฺห, \\ ปาปานิ ปริวชฺเชติ, \\ โย มุนาติ อุโภ โลเก,}{\csromangathabat{น โมเนน มุนี\textsuperscript{1} โหติ,}{มูฬฺหรูโป อวิทฺทสุ.} \\ \csromangathabat{โย จ ตุลํว ปคฺคยฺห,}{วรมาทาย ปณฺฑิโต.} \\ \csromangathabat{ปาปานิ ปริวชฺเชติ,}{ส มุนี เตน โส มุนิ.} \\ \csromangathabat{โย มุนาติ อุโภ โลเก,}{“มุนิ” เตน ปวุจฺจติ.} \\ อสตญฺจ สตญฺจ ญตฺวา ธมฺมํ, \\ อชฺฌตฺตํ พหิทฺธา จ สพฺพโลเก. \\ เทวมนุสฺเสหิ ปูชนีโย\textsuperscript{1}, \\ สงฺคชาลมติจฺจ\textsuperscript{1} โส มุนีติ.}
\csromangathameasurewak{อสตญฺจ สตญฺจ ญตฺวา ธมฺมํ, \\ อชฺฌตฺตํ พหิทฺธา จ สพฺพโลเก. \\ เทวมนุสฺเสหิ ปูชนีโย\textsuperscript{1}, \\ สงฺคชาลมติจฺจ\textsuperscript{1} โส มุนีติ.}
\csromangathasetleft{น โมเนน มุนี\textsuperscript{1} โหติ, \\ โย จ ตุลํว ปคฺคยฺห, \\ ปาปานิ ปริวชฺเชติ, \\ โย มุนาติ อุโภ โลเก,}
\csromangathagroup{\csromangathastanza{\csromangathabat{น โมเนน มุนี\csromansharedfootnote{0dbb9afe422efc37}{มุนิ (สฺยา, ก) ขุ 1. 51 ปิฏฺเฐ.} โหติ,}{มูฬฺหรูโป อวิทฺทสุ.} \\ \csromangathabat{โย จ ตุลํว ปคฺคยฺห,}{วรมาทาย ปณฺฑิโต.}}\csromangathastanza{\csromangathabat{ปาปานิ ปริวชฺเชติ,}{ส มุนี เตน โส มุนิ.} \\ \csromangathabat{โย มุนาติ อุโภ โลเก,}{“มุนิ” เตน ปวุจฺจติ.}}\csromangathastanza{\csromangathawak{อสตญฺจ สตญฺจ ญตฺวา ธมฺมํ, \\ อชฺฌตฺตํ พหิทฺธา จ สพฺพโลเก. \\ เทวมนุสฺเสหิ ปูชนีโย\csromansharedfootnote{0d17715fb8f4359e}{ปูชิโต (สฺยา, ก) ขุ 7. 45 ปิฏฺเฐปิ.}, \\ สงฺค\-ชาลม\-ติจฺจ\csromansharedfootnote{24d03025906a3874}{สงฺคํ ชาลมติจฺจ, ขุ 1. 360 ปิฏฺเฐ.} โส มุนีติ.}}}

\prose{\textbf{สาธุ วิยากโรหี}ติ ตํ สาธุ อาจิกฺขาหิ เทเสหิ ปญฺญเปหิ ปฏฺฐเปหิ วิวราหิ วิภชาหิ อุตฺตานีกโรหิ ปกาเสหีติ ตํ เม มุนี สาธุ วิยากโรหิ. \textbf{ตถา หิ เต วิทิโต เอส ธมฺโม}ติ ตถา หิ เต วิทิโต ตุลิโต ตีริโต วิภูโต วิภาวิโต เอส ธมฺโมติ ตถา หิ เต วิทิโต เอส ธมฺโม. เตนาห โส พฺราหฺมโณ\csromandash{}}

\setlength{\csromangathapagewidth}{0pt}
\setlength{\csromangathawakwidth}{0pt}
\csromangathameasuregroup{}{ยํ ตํ อปุจฺฉิมฺห อกิตฺตยี โน, \\ อญฺญํ ตํ ปุจฺฉาม ตทิงฺฆ พฺรูหิ. \\ กถํ นุ ธีรา วิตรนฺติ โอฆํ, \\ ชาติํ ชรํ โสกปริทฺทวญฺจ. \\ ตํ เม มุนี สาธุ วิยากโรหิ,}
\csromangathameasurewak{ยํ ตํ อปุจฺฉิมฺห อกิตฺตยี โน, \\ อญฺญํ ตํ ปุจฺฉาม ตทิงฺฆ พฺรูหิ. \\ กถํ นุ ธีรา วิตรนฺติ โอฆํ, \\ ชาติํ ชรํ โสกปริทฺทวญฺจ. \\ ตํ เม มุนี สาธุ วิยากโรหิ,}
\setlength{\csromangathaleftcol}{0pt}
\csromangathagroup{\csromangathastanza{\csromangathawak{ยํ ตํ อปุจฺฉิมฺห อกิตฺตยี โน, \\ อญฺญํ ตํ ปุจฺฉาม ตทิงฺฆ พฺรูหิ. \\ กถํ นุ ธีรา วิตรนฺติ โอฆํ, \\ ชาติํ ชรํ โสกปริทฺทวญฺจ.}}\csromangathastanza{\csromangathawak{ตํ เม มุนี สาธุ วิยากโรหิ,}}}

\prosecont{ตถา หิ เต วิทิโต เอส ธมฺโมติ.}

\csromanhangingbegin
\hangingheaditem{22}{\csromanfolio{75}\textbf{กิตฺตยิสฺสามิ เต ธมฺมํ, (เมตฺตคูติ ภควา,)}}
\hangingbody{\textbf{ทิฏฺเฐ ธมฺเม อนีติหํ.}}
\hangingbody{\textbf{ยํ วิทิตฺวา สโต จรํ, ตเร โลเก วิสตฺติกํ.}}
\csromanhangingend

\prose{\textbf{กิตฺตยิสฺสามิ เต ธมฺมนฺ}ติ ธมฺมนฺติ อาทิกลฺยาณํ มชฺเฌกลฺยาณํ ปริโยสาน\-กลฺยาณํ สาตฺถํ สพฺยญฺชนํ เกวล\-ปริปุณฺณํ ปริสุทฺธํ พฺรหฺมจริยํ จตฺตาโร สติปฏฺฐาเน จตฺตาโร สมฺมปฺปธาเน จตฺตาโร อิทฺธิปาเท ปญฺจินฺทฺริยานิ ปญฺจ พลานิ สตฺต โพชฺฌงฺเค อริยํ อฏฺฐงฺคิกํ มคฺคํ นิพฺพานญฺจ นิพฺพานคามินิญฺจ ปฏิปทํ กิตฺตยิสฺสามิ อาจิกฺขิสฺสามิ เทเสสฺสามิ ปญฺญเปสฺสามิ ปฏฺฐเปสฺสามิ วิวริสฺสามิ วิภชิสฺสามิ อุตฺตานีกริสฺสามิ ปกาสิสฺสามีติ กิตฺตยิสฺสามิ เต ธมฺมํ. \textbf{เมตฺตคู}ติ ภควา ตํ พฺราหฺมณํ นาเมน อาลปติ.}

\prose{\textbf{ทิฏฺเฐ ธมฺเม อนีติหนฺ}ติ ทิฏฺเฐ ธมฺเมติ \textbf{ทิฏฺเฐ ธมฺเม} ญาเต ธมฺเม ตุลิเต ธมฺเม ตีริเต ธมฺเม วิภูเต ธมฺเม วิภาวิเต ธมฺเม “สพฺเพ สงฺขารา อนิจฺจา”ติ ฯเปฯ. “ยํ กิญฺจิ สมุทย\-ธมฺมํ, สพฺพํ ตํ นิโรธธมฺมนฺ”ติ ทิฏฺเฐ ธมฺเม ญาเต ธมฺเม ตุลิเต ธมฺเม ตีริเต ธมฺเม วิภูเต ธมฺเม วิภาวิเต ธมฺเมติ เอวมฺปิ ทิฏฺเฐ ธมฺเม กถยิสฺสามิ. อถ วา ทุกฺเข ทิฏฺเฐ ทุกฺขํ กถยิสฺสามิ, สมุทเย ทิฏฺเฐ สมุทยํ กถยิสฺสามิ, มคฺเค ทิฏฺเฐ มคฺคํ กถยิสฺสามิ, นิโรเธ ทิฏฺเฐ นิโรธํ กถยิสฺสามีติ เอวมฺปิ ทิฏฺเฐ ธมฺเม กถยิสฺสามิ. อถ วา ทิฏฺเฐ ธมฺเม สนฺทิฏฺฐิกํ อกาลิกํ เอหิปสฺสิกํ โอปเนยฺยิกํ ปจฺจตฺตํ เวทิตพฺพํ วิญฺญูหีติ เอวมฺปิ ทิฏฺเฐ ธมฺเม กถยิสฺสามีติ ทิฏฺเฐ ธมฺเม. อนีติหนฺติ น อิติหีติหํ น อิติกิราย น ปรมฺปราย น ปิฏก\-สมฺปทาย น ตกฺกเหตุ น นยเหตุ น อาการ\-ปริวิตกฺเกน น ทิฏฺฐิ\-นิชฺฌาน\-กฺขนฺติยา สามํ สยม\-ภิญฺญาตํ อตฺต\-ปจฺจกฺขธมฺมํ, ตํ กถยิสฺสามีติ ทิฏฺเฐ ธมฺเม อนีติหํ.}

\prose{\textbf{ยํ วิทิตฺวา สโต จรนฺ}ติ ยํ วิทิตํ กตฺวา ตุลยิตฺวา ตีรยิตฺวา วิภาวยิตฺวา วิภูตํ กตฺวา, “สพฺเพ สงฺขารา อนิจฺจา”ติ วิทิตํ กตฺวา ตุลยิตฺวา ตีรยิตฺวา วิภาวยิตฺวา วิภูตํ กตฺวา, “สพฺเพ สงฺขารา ทุกฺขา”ติ ฯเปฯ. “สพฺเพ ธมฺมา อนตฺตา”ติ. “ยํ กิญฺจิ สมุทย\-ธมฺมํ, สพฺพํ ตํ นิโรธธมฺมนฺ”ติ วิทิตํ กตฺวา ตุลยิตฺวา ตีรยิตฺวา วิภาวยิตฺวา วิภูตํ กตฺวา. \textbf{สโต}ติ จตูหิ การเณหิ สโต, กาเย กายานุปสฺสนา\-สติปฏฺฐานํ ภาเวนฺโต สโต ฯเปฯ โส วุจฺจติ สโต. \csromanfolio{76}\textbf{จรนฺ}ติ จรนฺโต วิหรนฺโต อิริยนฺโต วตฺเตนฺโต ปาเลนฺโต ยเปนฺโต ยาเปนฺโตติ ยํ \textbf{วิทิตฺวา สโต จรํ.}}

\prose{\textbf{ตเร โลเก วิสตฺติกนฺ}ติ วิสตฺติกา วุจฺจติ ตณฺหา. โย ราโค สาราโค ฯเปฯ อภิชฺฌา โลโภ อกุสลมูลํ. \textbf{วิสตฺติกา}ติ เกนฏฺเฐน วิสตฺติกา, วิสตาติ วิสตฺติกา, วิสาลาติ วิสตฺติกา, วิสฏาติ วิสตฺติกา, วิสมาติ วิสตฺติกา, วิสกฺกตีติ วิสตฺติกา, วิสํหรตีติ วิสตฺติกา, วิสํวาทิกาติ วิสตฺติกา, วิสมูลาติ วิสตฺติกา, วิสผลาติ วิสตฺติกา, วิสปริโภคาติ วิสตฺติกา, วิสาลา วา ปน สา ตณฺหา รูเป สทฺเท คนฺเธ รเส โผฏฺฐพฺเพ กุเล คเณ อาวาเส ลาเภ ยเส ปสํสาย สุเข จีวเร ปิณฺฑปาเต เสนาสเน คิลานปจฺจย\-เภสชฺชปริกฺขาเร กามธาตุยา รูปธาตุยา อรูปธาตุยา กามภเว รูปภเว อรูปภเว สญฺญาภเว อสญฺญาภเว เนว\-สญฺญา\-นา\-สญฺญาภเว เอกโวการภเว จตุโวการภเว ปญฺจโวการภเว อตีเต อนาคเต ปจฺจุปฺปนฺเน ทิฏฺฐ\-สุต\-มุต\-วิญฺญาตพฺเพสุ ธมฺเมสุ วิสฏา วิตฺถตาติ วิสตฺติกา. \textbf{โลเก}ติ อปายโลเก มนุสฺสโลเก เทวโลเก ขนฺธโลเก ธาตุโลเก อายตนโลเก. \textbf{ตเร โลเก วิสตฺติกนฺ}ติ โลเก เวสา วิสตฺติกา\csromansharedfootnote{08aed5d9896c7774}{ยา สา โลเก วิสตฺติกา (สฺยา) กามสุตฺตนิทฺเทสฏฺฐกถา โอโลเกตพฺพา.}, โลเก เวตํ วิสตฺติกํ สโต ตเรยฺย อุตฺตเรยฺย ปตเรยฺย สมติกฺกเมยฺย วีติวตฺเตยฺยาติ ตเร โลเก วิสตฺติกํ. เตนาห ภควา\csromandash{}}

\prose{กิตฺตยิสฺสามิ เต ธมฺมํ, (เมตฺตคูติ ภควา,)}

\prose{ทิฏฺเฐ ธมฺเม อนีติหํ.}

\setlength{\csromangathapagewidth}{0pt}
\setlength{\csromangathawakwidth}{0pt}
\csromangathameasuregroup{ยํ วิทิตฺวา สโต จรํ, \\ \textbf{ตญฺจาหํ อภินนฺทา}มิ, \\ ยํ วิทิตฺวา สโต จรํ,}{\csromangathabat{ยํ วิทิตฺวา สโต จรํ,}{ตเร โลเก วิสตฺติกนฺติ.} \\ \csromangathabat{\textbf{ตญฺจาหํ อภินนฺทา}มิ,}{มเหสิ ธมฺม’มุตฺตมํ.} \\ \csromangathabat{ยํ วิทิตฺวา สโต จรํ,}{ตเร โลเก วิสตฺติกํ.}}
\csromangathasetleft{ยํ วิทิตฺวา สโต จรํ, \\ \textbf{ตญฺจาหํ อภินนฺทา}มิ, \\ ยํ วิทิตฺวา สโต จรํ,}
\csromangathagroup{\csromangathastanza{\csromangathabat{ยํ วิทิตฺวา สโต จรํ,}{ตเร โลเก วิสตฺติกนฺติ.}}\csromangathastanza{\csromangathaitembat{23}{\textbf{ตญฺจาหํ อภินนฺทา}มิ,}{มเหสิ ธมฺม’มุตฺตมํ.} \\ \csromangathacontbat{ยํ วิทิตฺวา สโต จรํ,}{ตเร โลเก วิสตฺติกํ.}}}

\prose{\textbf{ตญฺจาหํ อภินนฺทามี}ติ \textbf{ตนฺ}ติ ตุยฺหํ วจนํ พฺยปฺปถํ\csromansharedfootnote{d2749849e4790f58}{พฺยปถํ (สฺยา, ก)} เทสนํ อนุสาสนํ อนุสิฏฺฐํ\csromansharedfootnote{ccb4b5b8faf67e3a}{เทสนํ อนุสนฺธิํ (สฺยา)}. \textbf{นนฺทามี}ติ อภินนฺทามิ โมทามิ อนุโมทามิ อิจฺฉามิ สาทิยามิ ยาจามิ ปตฺถยามิ ปิหยามิ อภิชปฺปามีติ ตญฺจาหํ อภินนฺทามิ.}

\prose{\csromanfolio{77}\textbf{มเหสิ ธมฺม’มุตฺตมนฺ}ติ \textbf{มเหสี}ติ กิํ มเหสิ ภควา, มหนฺตํ สีลกฺขนฺธํ เอสี คเวสี\csromansharedfootnote{0407874f7301592c}{เอสิ คเวสิ (สฺยา) ขุ 7. 266 ปิฏฺเฐ.} ปริเยสีติ มเหสิ, มหนฺตํ สมาธิ\-กฺขนฺธํ. มหนฺตํ ปญฺญากฺขนฺธํ. มหนฺตํ วิมุตฺติ\-กฺขนฺธํ. มหนฺตํ วิมุตฺติ\-ญาณ\-ทสฺสนกฺขนฺธํ เอสี คเวสี ปริเยสีติ มเหสิ. มหโต ตโมกายสฺส ปทาลนํ เอสี คเวสี ปริเยสีติ มเหสิ, มหโต วิปลฺลาสสฺส ปเภทนํ เอสี คเวสี ปริเยสีติ มเหสิ, มหโต ตณฺหาสลฺลสฺส อพฺพหนํ\csromansharedfootnote{05035d751417b05d}{อพฺพูหนํ (พหูสุ) อพฺพูหํ (สีฏฺฐ)} เอสี คเวสี ปริเยสีติ มเหสิ, มหโต ทิฏฺฐิสํฆาฏสฺส วินิเวฐนํ เอสี คเวสี ปริเยสีติ มเหสิ, มหโต มานธชสฺส ปาตนํ เอสี คเวสี ปริเยสีติ มเหสิ, มหโต อภิสงฺขารสฺส วูปสมํ เอสี คเวสี ปริเยสีติ มเหสิ, มหโต โอฆสฺส นิตฺถรณํ เอสี คเวสี ปริเยสีติ มเหสิ, มหโต ภารสฺส นิกฺเขปนํ เอสี คเวสี ปริเยสีติ มเหสิ, มหโต สํสาร\-วฏฺฏสฺส อุปจฺเฉทํ เอสี คเวสี ปริเยสีติ มเหสิ, มหโต สนฺตาปสฺส นิพฺพาปนํ เอสี คเวสี ปริเยสีติ มเหสิ, มหโต ปริฬาหสฺส ปฏิปฺปสฺสทฺธํ เอสี คเวสี ปริเยสีติ มเหสิ, มหโต ธมฺมธชสฺส อุสฺสาปนํ เอสี คเวสี ปริเยสีติ มเหสิ, มหนฺเต สติปฏฺฐาเน. มหนฺเต สมฺมปฺปธาเน. มหนฺเต อิทฺธิปาเท. มหนฺตานิ อินฺทฺริยานิ. มหนฺตานิ พลานิ. มหนฺเต โพชฺฌงฺเค. มหนฺตํ อริยํ อฏฺฐงฺคิกํ มคฺคํ. มหนฺตํ ปรมตฺถํ อมตํ นิพฺพานํ เอสี คเวสี ปริเยสีติ มเหสิ, มเหสกฺเขหิ สตฺเตหิ เอสิโต คเวสิโต ปริเยสิโต “กหํ พุทฺโธ, กหํ ภควา, กหํ เทวเทโว, กหํ นราสโภ”ติ มเหสิ. \textbf{ธมฺม’มุตฺตมนฺ}ติ ธมฺมมุตฺตมํ วุจฺจติ อมตํ นิพฺพานํ. โย โส สพฺพ\-สงฺขาร\-สมโถ สพฺพู\-ปธิ\-ปฏินิสฺสคฺโค ตณฺหกฺขโย วิราโค นิโรโธ นิพฺพานํ. \textbf{อุตฺตมนฺ}ติ อคฺคํ เสฏฺฐํ วิเสฏฺฐํ ปาโมกฺขํ อุตฺตมํ ปวรํ ธมฺมนฺติ มเหสิ ธมฺม’มุตฺตมํ.}

\prose{\textbf{ยํ วิทิตฺวา สโต จรนฺ}ติ วิทิตํ กตฺวา ตุลยิตฺวา ตีรยิตฺวา วิภาวยิตฺวา วิภูตํ กตฺวา, “สพฺเพ สงฺขารา อนิจฺจา”ติ วิทิตํ กตฺวา ตุลยิตฺวา ตีรยิตฺวา วิภาวยิตฺวา วิภูตํ กตฺวา, “สพฺเพ สงฺขารา ทุกฺขา”ติ ฯเปฯ. “สพฺเพ ธมฺมา อนตฺตา”ติ. “ยํ กิญฺจิ สมุทย\-ธมฺมํ, สพฺพํ ตํ นิโรธธมฺมนฺ”ติ วิทิตํ กตฺวา ตุลยิตฺวา ตีรยิตฺวา วิภาวยิตฺวา วิภูตํ กตฺวา. \textbf{สโต}ติ จตูหิ การเณหิ สโต, กาเย}

\prose{\csromanfolio{78}กายานุปสฺสนา\-สติปฏฺฐานํ ภาเวนฺโต สโต. เวทนาสุ. จิตฺเต. ธมฺเมสุ ธมฺมานุปสฺสนา\-สติปฏฺฐานํ ภาเวนฺโต สโต ฯเปฯ โส วุจฺจติ สโต. \textbf{จรนฺ}ติ จรนฺโต วิหรนฺโต อิริยนฺโต วตฺเตนฺโต ปาเลนฺโต ยเปนฺโต ยาเปนฺโตติ ยํ วิทิตฺวา สโต จรํ.}

\prose{\textbf{ตเร โลเก วิสตฺติกนฺ}ติ วิสตฺติกา วุจฺจติ ตณฺหา. โย ราโค สาราโค ฯเปฯ อภิชฺฌา โลโภ อกุสลมูลํ. \textbf{วิสตฺติกา}ติ เกนฏฺเฐน วิสตฺติกา ฯเปฯ วิสฏา วิตฺถตาติ วิสตฺติกา. \textbf{โลเก}ติ อปายโลเก ฯเปฯ อายตนโลเก. \textbf{ตเร โลเก วิสตฺติกนฺ}ติ โลเก เวสา วิสตฺติกา, โลเก เวตํ วิสตฺติกํ สโต ตเรยฺย อุตฺตเรยฺย ปตเรยฺย สมติกฺกเมยฺย วีติวตฺเตยฺยาติ ตเร โลเก วิสตฺติกํ. เตนาห โส พฺราหฺมโณ\csromandash{}}

\setlength{\csromangathapagewidth}{0pt}
\setlength{\csromangathawakwidth}{0pt}
\csromangathameasuregroup{ตญฺจาหํ อภินนฺทามิ, \\ ยํ วิทิตฺวา สโต จรํ,}{\csromangathabat{ตญฺจาหํ อภินนฺทามิ,}{มเหสิ ธมฺม’มุตฺตมํ.} \\ \csromangathabat{ยํ วิทิตฺวา สโต จรํ,}{ตเร โลเก วิสตฺติกนฺติ.}}
\csromangathasetleft{ตญฺจาหํ อภินนฺทามิ, \\ ยํ วิทิตฺวา สโต จรํ,}
\csromangathagroup{\csromangathastanza{\csromangathabat{ตญฺจาหํ อภินนฺทามิ,}{มเหสิ ธมฺม’มุตฺตมํ.} \\ \csromangathabat{ยํ วิทิตฺวา สโต จรํ,}{ตเร โลเก วิสตฺติกนฺติ.}}}

\csromanhangingbegin
\hangingheaditem{24}{\textbf{ยํ กิญฺจิ สมฺปชานาสิ, (เมตฺตคูติ ภควา,)}}
\hangingbody{\textbf{อุทฺธํ อโธ ติริยญฺจาปิ มชฺเฌ.}}
\hangingbody{\textbf{เอเตสุ นนฺทิญฺจ นิเวสนญฺจ,}}
\hangingbody{\textbf{ปนุชฺช วิญฺญาณํ ภเว น ติฏฺเฐ.}}
\csromanhangingend

\prose{\textbf{ยํ กิญฺจิ สมฺปชานาสี}ติ ยํ กิญฺจิ ปชานาสิ อาชานาสิ วิชานาสิ ปฏิวิชานาสิ ปฏิวิชฺฌสีติ ยํ กิญฺจิ สมฺปชานาสิ. \textbf{เมตฺตคู}ติ ภควา ตํ พฺราหฺมณํ นาเมน อาลปติ. \textbf{ภควา}ติ คารวาธิวจนเมตํ ฯเปฯ สจฺฉิกา ปญฺญตฺติ, ยทิทํ ภควาติ เมตฺตคูติ ภควา.}

\prose{\textbf{อุทฺธํ อโธ ติริยญฺจาปิ มชฺเฌ}ติ \textbf{อุทฺธนฺ}ติ อนาคตํ\csromansharedfootnote{96b5e8e3c4349bee}{อุทฺธํ วุจฺจติ อนาคตํ (สฺยา, ก)}. \textbf{อโธ}ติ อตีตํ. \textbf{ติริยญฺจาปิ มชฺเฌ}ติ ปจฺจุปฺปนฺนํ. \textbf{อุทฺธนฺ}ติ เทวโลโก. \textbf{อโธ}ติ นิรยโลโก. \textbf{ติริยญฺจาปิ มชฺเฌ}ติ มนุสฺสโลโก. อถ วา \textbf{อุทฺธนฺ}ติ กุสลา ธมฺมา. \textbf{อโธ}ติ อกุสลา ธมฺมา. \textbf{ติริยญฺจาปิ มชฺเฌ}ติ อพฺยากตา ธมฺมา. \textbf{อุทฺธนฺ}ติ อรูปธาตุ. \textbf{อโธ}ติ กามธาตุ. \textbf{ติริยญฺจาปิ มชฺเฌ}ติ รูปธาตุ. \textbf{อุทฺธนฺ}ติ สุขา เวทนา. \textbf{อโธ}ติ ทุกฺขา เวทนา. \textbf{ติริยญฺจาปิ มชฺเฌ}ติ อทุกฺขมสุขา เวทนา. \textbf{อุทฺธนฺ}ติ อุทฺธํ ปาทตลา. \csromanfolio{79}\textbf{อโธ}ติ อโธ เกสมตฺถกา. \textbf{ติริยญฺจาปิ มชฺเฌ}ติ เวมชฺเฌติ อุทฺธํ อโธ \textbf{ติริยญฺจาปิ มชฺเฌ.}}

\prose{\textbf{เอเตสุ นนฺทิญฺจ นิเวสนญฺจ, ปนุชฺช วิญฺญาณํ ภเว น ติฏฺเฐ}ติ \textbf{เอเตสู}ติ อาจิกฺขิเตสุ เทสิเตสุ ปญฺญปิเตสุ ปฏฺฐปิเตสุ วิวริเตสุ วิภชิเตสุ อุตฺตานีกเตสุ ปกาสิเตสุ. \textbf{นนฺที} วุจฺจติ ตณฺหา. โย ราโค สาราโค ฯเปฯ อภิชฺฌา โลโภ อกุสลมูลํ. \textbf{นิเวสนนฺ}ติ เทฺว นิเวสนา, ตณฺหานิเวสนา จ ทิฏฺฐินิเวสนา จ. กตมา ตณฺหานิเวสนา, ยาวตา ตณฺหา\-สงฺขาเตน ฯเปฯ อยํ ตณฺหานิเวสนา. กตมา ทิฏฺฐินิเวสนา, วีสติวตฺถุกา สกฺกายทิฏฺฐิ ฯเปฯ อยํ ทิฏฺฐินิเวสนา.}

\prose{\textbf{ปนุชฺช วิญฺญาณนฺ}ติ ปุญฺญาภิสงฺขาร\-สหคตํ วิญฺญาณํ อปุญฺญาภิสงฺขาร\-สหคตํ วิญฺญาณํ อาเนญฺชาภิสงฺขาร\-สหคตํ วิญฺญาณํ, เอเตสุ นนฺทิญฺจ นิเวสนญฺจ อภิสงฺขาร\-สหคตญฺจ วิญฺญาณํ นุชฺช ปนุชฺช นุท ปนุท ชห ปชห วิโนเทหิ พฺยนฺตีกโรหิ อนภาวํ คเมหีติ เอเตสุ นนฺทิญฺจ นิเวสนญฺจ, ปนุชฺช วิญฺญาณํ.}

\prose{\textbf{ภเว น ติฏฺเฐ}ติ \textbf{ภวา}ติ เทฺว ภวา กมฺมภโว จ ปฏิสนฺธิโก จ ปุนพฺภโว. กตโม กมฺมภโว, ปุญฺญา\-ภิสงฺขาโร อปุญฺญา\-ภิสงฺขาโร อาเนญฺชา\-ภิสงฺขาโร, อยํ กมฺมภโว. กตโม ปฏิสนฺธิโก ปุนพฺภโว, ปฏิสนฺธิกา รูปํ เวทนา สญฺญา สงฺขารา วิญฺญาณํ, อยํ ปฏิสนฺธิโก ปุนพฺภโว. \textbf{ภเว น ติฏฺเฐ}ติ นนฺทิญฺจ นิเวสนญฺจ อภิสงฺขาร\-สหคตํ วิญฺญาณญฺจ กมฺมภวญฺจ ปฏิสนฺธิ\-กญฺจ ปุนพฺภวํ ปชหนฺโต วิโนเทนฺโต พฺยนฺตีกโรนฺโต อนภาวํ คเมนฺโต กมฺมภเว น ติฏฺเฐยฺย, ปฏิสนฺธิเก ปุนพฺภเว น ติฏฺเฐยฺย น สนฺติฏฺเฐยฺยาติ ปนุชฺช วิญฺญาณํ ภเว น ติฏฺเฐ. เตนาห ภควา\csromandash{}}

\prose{ยํ กิญฺจิ สมฺปชานาสิ, (เมตฺตคูติ ภควา,)}

\prose{อุทฺธํ อโธ ติริยญฺจาปิ มชฺเฌ.}

\prose{\textbf{เอเตสุ นนฺทิญฺจ นิเวสนญฺจ,}}

\prose{ปนุชฺช วิญฺญาณํ ภเว น ติฏฺเฐติ.}

\setlength{\csromangathapagewidth}{0pt}
\setlength{\csromangathawakwidth}{0pt}
\csromangathameasuregroup{}{\textbf{เอวํวิหารี สโต อปฺปมตฺโต,} \\ \textbf{ภิกฺขุ จรํ หิตฺวา มมายิตานิ.} \\ \textbf{ชาติํ ชรํ โสกปริทฺทวญฺจ,} \\ \textbf{อิเธว วิทฺวา ปชเหยฺย ทุกฺขํ.}}
\csromangathameasurewak{\textbf{เอวํวิหารี สโต อปฺปมตฺโต,} \\ \textbf{ภิกฺขุ จรํ หิตฺวา มมายิตานิ.} \\ \textbf{ชาติํ ชรํ โสกปริทฺทวญฺจ,} \\ \textbf{อิเธว วิทฺวา ปชเหยฺย ทุกฺขํ.}}
\setlength{\csromangathaleftcol}{0pt}
\csromangathagroup{\csromangathastanza{\csromangathawak{\csromangathaitemline{25}{\textbf{เอวํวิหารี สโต อปฺปมตฺโต,}} \\ \csromangathacontline{\textbf{ภิกฺขุ จรํ หิตฺวา มมายิตานิ.}} \\ \csromangathacontline{\textbf{ชาติํ ชรํ โสกปริทฺทวญฺจ,}} \\ \csromangathacontline{\textbf{อิเธว วิทฺวา ปชเหยฺย ทุกฺขํ.}}}}}

\prose{\csromanfolio{80}\textbf{เอวํวิหารี สโต อปฺปมตฺโต}ติ \textbf{เอวํวิหารี}ติ นนฺทิญฺจ นิเวสนญฺจ อภิสงฺขาร\-สหคต\-วิญฺญาณญฺจ กมฺมภวญฺจ ปฏิสนฺธิ\-กญฺจ ปุนพฺภวํ ปชหนฺโต วิโนเทนฺโต พฺยนฺตีกโรนฺโต อนภาวํ คเมนฺโตติ เอวํวิหารี. \textbf{สโต}ติ จตูหิ การเณหิ สโต, กาเย กายานุปสฺสนา\-สติปฏฺฐานํ ภาเวนฺโต ฯเปฯ โส วุจฺจติ สโต. \textbf{อปฺปมตฺโต}ติ สกฺกจฺจการี สาตจฺจการี อฏฺฐิตการี อโนลีนวุตฺตีโก\csromansharedfootnote{94e5e2cba323d139}{อโนลีนวุตฺติโก (ก) ขุ 7. 45 ปิฏฺเฐ.} อนิกฺขิตฺตจฺฉนฺโท อนิกฺขิตฺตธุโร อปฺปมตฺโต กุสเลสุ ธมฺเมสุ “กถาหํ\csromansharedfootnote{1e79fde666862bce}{กทาหํ (สฺยา)} อปริปูรํ วา สีลกฺขนฺธํ ปริปูเรยฺยํ, ปริปูรํ วา สีลกฺขนฺธํ ตตฺถ ตตฺถ ปญฺญาย อนุคฺคเณฺหยฺยนฺ”ติ โย ตตฺถ ฉนฺโท จ วายาโม จ อุสฺสาโห จ อุสฺโสฬฺหี จ อปฺปฏิวานี\csromansharedfootnote{09ccf9b2e4367479}{อปฺปฏิวานิ (ก) ขุ 7. 45 ปิฏฺเฐ.} จ สติ จ สมฺปชญฺญญฺจ อาตปฺปํ ปธานํ อธิฏฺฐานํ อนุโยโค อปฺปมตฺโต อปฺปมาโท กุสเลสุ ธมฺเมสุ. “กถาหํ อปริปูรํ วา สมาธิ\-กฺขนฺธํ. ปญฺญากฺขนฺธํ. วิมุตฺติ\-กฺขนฺธํ. วิมุตฺติ\-ญาณ\-ทสฺสนกฺขนฺธํ ปริปูเรยฺยํ, ปริปูรํ วา วิมุตฺติ\-ญาณ\-ทสฺสนกฺขนฺธํ ตตฺถ ตตฺถ ปญฺญาย อนุคฺคเณฺหยฺยนฺ”ติ โย ตตฺถ ฉนฺโท จ วายาโม จ อุสฺสาโห จ อุสฺโสฬฺหี จ อปฺปฏิวานี จ สติ จ สมฺปชญฺญญฺจ อาตปฺปํ ปธานํ อธิฏฺฐานํ อนุโยโค อปฺปมตฺโต อปฺปมาโท กุสเลสุ ธมฺเมสุ. “กถาหํ อปริญฺญาตํ วา ทุกฺขํ ปริชาเนยฺยํ, อปฺปหีเน วา กิเลเส ปชเหยฺยํ, อภาวิตํ วา มคฺคํ ภาเวยฺยํ, อสจฺฉิกตํ วา นิโรธํ สจฺฉิกเรยฺยนฺ”ติ โย ตตฺถ ฉนฺโท จ วายาโม จ อุสฺสาโห จ อุสฺโสฬฺหี จ อปฺปฏิวานี จ สติ จ สมฺปชญฺญญฺจ อาตปฺปํ ปธานํ อธิฏฺฐานํ อนุโยโค อปฺปมตฺโต อปฺปมาโท กุสเลสุ ธมฺเมสูติ เอวํวิหารี สโต อปฺปมตฺโต.}

\prose{\textbf{ภิกฺขุ จรํ หิตฺวา มมายิตานี}ติ \textbf{ภิกฺขู}ติ ปุถุชฺชน\-กลฺยาณโก\csromansharedfootnote{a90ef8afbff1ac5f}{กลฺยาณปุถุชฺชโน (สฺยา), เอวมีทิเสสุ ฐาเนสุ.} วา ภิกฺขุ เสกฺโข วา ภิกฺขุ. \textbf{จรนฺ}ติ จรนฺโต วิหรนฺโต อิริยนฺโต วตฺเตนฺโต ปาเลนฺโต ยเปนฺโต ยาเปนฺโต. \textbf{มมตฺตา}ติ เทฺว มมตฺตา ตณฺหา\-มมตฺตญฺจ ทิฏฺฐิ\-มมตฺตญฺจ ฯเปฯ อิทํ ตณฺหามมตฺตํ ฯเปฯ อิทํ ทิฏฺฐิมมตฺตํ. ตณฺหามมตฺตํ ปหาย ทิฏฺฐิมมตฺตํ ปฏินิสฺสชฺชิตฺวา มมตฺเต ชหิตฺวา จชิตฺวา ปชหิตฺวา วิโนเทตฺวา พฺยนฺตีกริตฺวา อนภาวํ คเมตฺวาติ ภิกฺขุ จรํ หิตฺวา มมายิตานิ.}

\prose{\csromanfolio{81}\textbf{ชาติํ ชรํ โสกปริทฺทวญฺจ, อิเธว วิทฺวา ปชเหยฺย ทุกฺขนฺ}ติ \textbf{ชาตี}ติ ยา เตสํ เตสํ สตฺตานํ ฯเปฯ. \textbf{ชรนฺ}ติ ยา เตสํ เตสํ สตฺตานํ ฯเปฯ. \textbf{โสโก}ติ ญาติพฺยสเนน วา ผุฏฺฐสฺส ฯเปฯ. \textbf{ปริเทโว}ติ ญาติพฺยสเนน วา ผุฏฺฐสฺส ฯเปฯ. \textbf{อิธา}ติ อิมิสฺสา ทิฏฺฐิยา ฯเปฯ. อิมสฺมิํ มนุสฺสโลเก. \textbf{วิทฺวา}ติ วิชฺชาคโต ญาณี วิภาวี เมธาวี. \textbf{ทุกฺขนฺ}ติ ชาติทุกฺขํ ฯเปฯ โทมนสฺสุ\-ปายา\-สทุกฺขํ. ชาติํ ชรํ โสกปริทฺทวญฺจ, อิเธว วิทฺวา ปชเหยฺย ทุกฺขนฺติ วิชฺชาคโต ญาณี วิภาวี เมธาวี อิเธว ชาติญฺจ ชรญฺจ โสกปริทฺทวญฺจ ทุกฺขญฺจ ปชเหยฺย วิโนเทยฺย พฺยนฺตีกเรยฺย อนภาวํ คเมยฺยาติ ชาติํ ชรํ โสกปริทฺทวญฺจ, อิเธว วิทฺวา ปชเหยฺย ทุกฺขํ. เตนาห ภควา\csromandash{}}

\setlength{\csromangathapagewidth}{0pt}
\setlength{\csromangathawakwidth}{0pt}
\csromangathameasuregroup{}{เอวํวิหารี สโต อปฺปมตฺโต, \\ ภิกฺขุ จรํ หิตฺวา มมายิตานิ.}
\csromangathameasurewak{เอวํวิหารี สโต อปฺปมตฺโต, \\ ภิกฺขุ จรํ หิตฺวา มมายิตานิ.}
\setlength{\csromangathaleftcol}{0pt}
\csromangathagroup{\csromangathastanza{\csromangathawak{เอวํวิหารี สโต อปฺปมตฺโต, \\ ภิกฺขุ จรํ หิตฺวา มมายิตานิ.}}}

\prose{\textbf{ชาติํ ชรํ โสกปริทฺทวญฺจฺ}อ,}

\prose{อิเธว วิทฺวา ปชเหยฺย ทุกฺขนฺติ.}

\proseitem{26}{\textbf{เอตา’ภินนฺทามิ วโจ มเหสิโน,}}

\prose{\textbf{สุกิตฺติตํ โคตม’นูปธีกํ.}}

\prose{\textbf{อทฺธา หิ ภควา ปหาสิ ทุกฺขํ,}}

\prose{(9)}

\prose{\textbf{เอตา’ภินนฺทามิ วโจ มเหสิโน}ติ \textbf{เอตนฺ}ติ ตุยฺหํ วจนํ พฺยปฺปถํ เทสนํ อนุสาสนํ อนุสิฏฺฐํ นนฺทามิ อภินนฺทามิ โมทามิ อนุโมทามิ อิจฺฉามิ สาทิยามิ ปตฺถยามิ ปิหยามิ อภิชปฺปามิ. \textbf{มเหสิโน}ติ กิํ มเหสิ ภควา, มหนฺตํ สีลกฺขนฺธํ เอสี คเวสี ปริเยสีติ มเหสิ ฯเปฯ “กหํ นราสโภ”ติ มเหสีติ เอตาภินนฺทามิ วโจ มเหสิโน.}

\prose{\textbf{สุกิตฺติตํ โคตมนูปธีกนฺ}ติ \textbf{สุกิตฺติตนฺ}ติ สุกิตฺติตํ สุอาจิกฺขิตํ สุเทสิตํ สุปญฺญปิตํ สุปฏฺฐปิตํ สุวิวริตํ สุวิภชิตํ สุ- อุตฺตานีกตํ สุปกาสิตนฺติ สุกิตฺติตํ. \textbf{โคตมนูปธีกนฺ}ติ อุปธี วุจฺจนฺติ กิเลสา จ ขนฺธา จ อภิสงฺขารา จ. อุปธิปฺปหานํ อุปธิ\-วูปสมํ อุปธิ\-ปฏินิสฺสคฺคํ อุปธิ\-ปฏิปสฺสทฺธํ อมตํ นิพฺพานนฺติ สุกิตฺติตํ โคตม\-นูปธีกํ.}

\prose{\csromanfolio{82}\textbf{อทฺธา หิ ภควา ปหาสิ ทุกฺขนฺ}ติ \textbf{อทฺธา}ติ เอกํสวจนํ นิสฺสํสยวจนํ นิกฺกงฺขา\-วจนํ อเทฺวชฺฌวจนํ อเทฺวฬฺหก\-วจนํ นิโรธวจนํ อปณฺณก\-วจนํ อวตฺถาปน\-วจนเมตํ อทฺธาติ. \textbf{ภควา}ติ คารวาธิวจนเมตํ ฯเปฯ สจฺฉิกา ปญฺญตฺติ, ยทิทํ ภควาติ. \textbf{ปหาสิ ทุกฺขนฺ}ติ ชาติทุกฺขํ ชราทุกฺขํ พฺยาธิทุกฺขํ มรณทุกฺขํ โสกปริเทวทุกฺขโทมนสฺสุปายาส\-ทุกฺขํ ปหาสิ ปชหิ วิโนเทสิ พฺยนฺตีกโรสิ อนภาวํ คเมสีติ อทฺธา หิ ภควา ปหาสิ ทุกฺขํ.}

\prose{\textbf{ตถา หิ เต วิทิโต เอส ธมฺโม}ติ ตถา หิ เต วิทิโต ตุลิโต ตีริโต วิภูโต วิภาวิโต เอส ธมฺโมติ ตถา หิ เต วิทิโต เอส ธมฺโม. เตนาห โส พฺราหฺมโณ\csromandash{}}

\setlength{\csromangathapagewidth}{0pt}
\setlength{\csromangathawakwidth}{0pt}
\csromangathameasuregroup{}{เอตา’ภินนฺทามิ วโจ มเหสิโน, \\ สุกิตฺติตํ โคตม’นูปธีกํ.}
\csromangathameasurewak{เอตา’ภินนฺทามิ วโจ มเหสิโน, \\ สุกิตฺติตํ โคตม’นูปธีกํ.}
\setlength{\csromangathaleftcol}{0pt}
\csromangathagroup{\csromangathastanza{\csromangathawak{เอตา’ภินนฺทามิ วโจ มเหสิโน, \\ สุกิตฺติตํ โคตม’นูปธีกํ.}}}

\prose{อทฺธา หิ ภควา ปหาสิ ทุกฺขํ,}

\prose{ตถา หิ เต วิทิโต เอส ธมฺโมติ.}

\proseitem{27}{\textbf{เต จาปิ นูนปฺปชเหยฺยุ ทุกฺขํ,}}

\prose{\textbf{เย ตฺวํ มุนี อฏฺฐิตํ โอวเทยฺย.}}

\prose{\textbf{ตํ ตํ นมสฺสามิ สเมจฺจ นาคํ,}}

\prose{\textbf{อปฺเปว มํ ภควา อฏฺฐิตํ โอวเทยฺย.} (10)}

\prose{\textbf{เต จาปิ นูนปฺปชเหยฺยุ ทุกฺขนฺ}ติ \textbf{เต จาปี}ติ ขตฺติยา จ พฺราหฺมณา จ เวสฺสา จ สุทฺทา จ คหฏฺฐา จ ปพฺพชิตา จ เทวา จ มนุสฺสา จ. \textbf{ปชเหยฺยุ ทุกฺขนฺ}ติ ชาติทุกฺขํ ชราทุกฺขํ พฺยาธิทุกฺขํ มรณทุกฺขํ โสกปริเทว ทุกฺขโทมนสฺสุ\-ปายา\-สทุกฺขํ ปชเหยฺยุํ วิโนเทยฺยุํ พฺยนฺตีกเรยฺยุํ อนภาวํ คเมยฺยุนฺติ เต จาปิ นูนปฺปชเหยฺยุ ทุกฺขํ.}

\prose{\textbf{เย ตฺวํ มุนี อฏฺฐิตํ โอวเทยฺยา}ติ เยติ ขตฺติ\textbf{เย} จ พฺราหฺมเณ จ เวสฺเส จ สุทฺเท จ คหฏฺเฐ จ ปพฺพชิเต จ เทเว จ มนุสฺเส จ. \textbf{ตฺวนฺ}ติ ภควนฺตํ ภณติ. \textbf{มุนี}ติ โมนํ วุจฺจติ ญาณํ ฯเปฯ สงฺค\-ชาลม\-ติจฺจ โส มุนิ. \textbf{อฏฺฐิตํ โอวเทยฺยา}ติ อฏฺฐิตํ โอวเทยฺย, สกฺกจฺจํ โอวเทยฺย, อภิณฺหํ โอวเทยฺย, ปุนปฺปุนํ โอวเทยฺย อนุสาเสยฺยาติ เย ตฺวํ มุนี อฏฺฐิตํ โอวเทยฺย.}

\prose{\textbf{ตํ ตํ นมสฺสามิ สเมจฺจ นาคนฺ}ติ \textbf{ตนฺ}ติ ภควนฺตํ ภณติ. \textbf{นมสฺสามี}ติ กาเยน วา นมสฺสามิ, วาจาย วา นมสฺสามิ, จิตฺเตน วา นมสฺสามิ, \csromanfolio{83}อนฺวตฺถ\-ปฏิปตฺติยา วา นมสฺสามิ, ธมฺมานุธมฺม\-ปฏิปตฺติยา วา นมสฺสามิ สกฺกโรมิ ครุํ กโรมิ\csromansharedfootnote{ef6821a7d2194587}{ครุกโรมิ (สฺยา)} มาเนมิ ปูเชมิ. \textbf{สเมจฺจา}ติ สเมจฺจ อภิสเมจฺจ สมาคนฺตฺวา อภิสมาคนฺตฺวา สมฺมุขา ตํ นมสฺสามิ. \textbf{นาคนฺ}ติ นาโค จ ภควา, อาคุํ น กโรตีติ นาโค, น คจฺฉตีติ นาโค, น อาคจฺฉตีติ นาโค, กถํ ภควา อาคุํ น กโรตีติ นาโค. อาคุ วุจฺจติ \textbf{ปาปกา อกุสลา ธมฺมา สํกิเลสิกา โปโนภวิกา สทรา ทุกฺขวิปากา} \textbf{อายติํ ชาติ}\-\textbf{ชรา}\-\textbf{มรณิยา.}}

\vspace{\tipitakaparskip}
\csromancenter{อาคุํ น กโรติ กิญฺจิ โลเก, (สภิยาติ ภควา,)}
\prose{สพฺพสํโยเค\csromansharedfootnote{38f592f752342f9f}{สพฺพโยเค (ก) ขุ 1. 359 ปิฏฺเฐ.} วิสชฺช พนฺธนานิ.}

\setlength{\csromangathapagewidth}{0pt}
\setlength{\csromangathawakwidth}{0pt}
\csromangathameasuregroup{}{สพฺพตฺถ น สชฺชตี วิมุตฺโต, \\ นาโค ตาทิ ปวุจฺจเต ตถตฺตาติ.}
\csromangathameasurewak{สพฺพตฺถ น สชฺชตี วิมุตฺโต, \\ นาโค ตาทิ ปวุจฺจเต ตถตฺตาติ.}
\setlength{\csromangathaleftcol}{0pt}
\csromangathagroup{\csromangathastanza{\csromangathawak{สพฺพตฺถ น สชฺชตี วิมุตฺโต, \\ นาโค ตาทิ ปวุจฺจเต ตถตฺตาติ.}}}

\prose{เอวํ ภควา อาคุํ น กโรตีติ นาโค.}

\prose{กถํ ภควา น คจฺฉตีติ นาโค, ภควา น ฉนฺทาคติํ คจฺฉติ, น โทสาคติํ คจฺฉติ, น โมหาคติํ คจฺฉติ, น ภยาคติํ คจฺฉติ, น ราควเสน คจฺฉติ, น โทสวเสน คจฺฉติ, น โมหวเสน คจฺฉติ, น มานวเสน คจฺฉติ, น ทิฏฺฐิวเสน คจฺฉติ, น อุทฺธจฺจวเสน คจฺฉติ, น วิจิกิจฺฉา\-วเสน คจฺฉติ, น อนุสยวเสน คจฺฉติ, น วคฺเคหิ ธมฺเมหิ ยายติ นียติ\csromansharedfootnote{ee843794401f93eb}{นิยฺยติ (สฺยา, ก)} วุยฺหติ สํหรียติ. เอวํ ภควา น คจฺฉตีติ นาโค.}

\prose{กถํ ภควา น อาคจฺฉตีติ นาโค, โสตาปตฺติ\-มคฺเคน เย กิเลสา ปหีนา, เต กิเลเส น ปุเนติ น ปจฺเจติ น ปจฺจาคจฺฉติ. สกทาคามิ\-มคฺเคน. อนาคามิมคฺเคน. อรหตฺต\-มคฺเคน เย กิเลสา ปหีนา, เต กิเลเส น ปุเนติ น ปจฺเจติ น ปจฺจาคจฺฉติ. เอวํ ภควา น อาคจฺฉตีติ นาโคติ ตํ ตํ นมสฺสามิ สเมจฺจ นาค.}

\prose{\textbf{อปฺเปว มํ ภควา อฏฺฐิตํ โอวเทยฺยา}ติ อปฺเปว มํ ภควา อฏฺฐิตํ โอวเทยฺย, สกฺกจฺจํ โอวเทยฺย, อภิณฺหํ โอวเทยฺย, ปุนปฺปุนํ โอวเทยฺย, อนุสาเสยฺยาติ อปฺเปว มํ ภควา อฏฺฐิตํ โอวเทยฺย. เตนาห โส พฺราหฺมโณ\csromandash{}}

\proseitem{4}{\csromanfolio{84}เต จาปิ นูนปฺปชเหยฺยุ ทุกฺขํ.}

\prose{เย ตฺวํ มุนี อฏฺฐิตํ โอวเทยฺย.}

\setlength{\csromangathapagewidth}{0pt}
\setlength{\csromangathawakwidth}{0pt}
\csromangathameasuregroup{}{ตํ ตํ นมสฺสามิ สเมจฺจ นาคํ, \\ อปฺเปว มํ ภควา อฏฺฐิตํ โอวเทยฺยาติ.}
\csromangathameasurewak{ตํ ตํ นมสฺสามิ สเมจฺจ นาคํ, \\ อปฺเปว มํ ภควา อฏฺฐิตํ โอวเทยฺยาติ.}
\setlength{\csromangathaleftcol}{0pt}
\csromangathagroup{\csromangathastanza{\csromangathawak{ตํ ตํ นมสฺสามิ สเมจฺจ นาคํ, \\ อปฺเปว มํ ภควา อฏฺฐิตํ โอวเทยฺยาติ.}}}

\proseitem{28}{\textbf{ยํ พฺราหฺมณํ เวทคุ’มาภิชญฺญา,}}

\prose{\textbf{อกิญฺจนํ กามภเว อสตฺตํ.}}

\prose{\textbf{อทฺธา หิ โส โอฆมิมํ อตาริ,}}

\prose{\textbf{ติณฺโณ จ ปารํ อขิโล อกงฺโข.} (11)}

\prose{\textbf{ยํ พฺราหฺมณํ เวทคุมา}\-\textbf{ภิชญฺญา}ติ \textbf{พฺราหฺมโณ}ติ สตฺตนฺนํ ธมฺมานํ พาหิตตฺตา พฺราหฺมโณ, สกฺกายทิฏฺฐิ พาหิตา โหติ, วิจิกิจฺฉา พาหิตา โหติ, สีลพฺพต\-ปรามาโส พาหิโต โหติ, ราโค พาหิโต โหติ, โทโส พาหิโต โหติ, โมโห พาหิโต โหติ, มาโน พาหิโต โหติ, พาหิตาสฺส โหนฺติ ปาปกา อกุสลา ธมฺมา สํกิเลสิกา โปโนภวิกา สทรา ทุกฺขวิปากา อายติํ ชาติ\-ชรา\-มรณิยา.}

\prose{พาหิตฺวา สพฺพปาปกานิ, (สภิยาติ ภควา,)}

\prose{วิมโล สาธุสมาหิโต ฐิตตฺโต.}

\setlength{\csromangathapagewidth}{0pt}
\setlength{\csromangathawakwidth}{0pt}
\csromangathameasuregroup{}{สํสารมติจฺจ เกวลี โส, \\ อสิโต\textsuperscript{1} ตาทิ ปวุจฺจเต ส พฺรหฺมา.}
\csromangathameasurewak{สํสารมติจฺจ เกวลี โส, \\ อสิโต\textsuperscript{1} ตาทิ ปวุจฺจเต ส พฺรหฺมา.}
\setlength{\csromangathaleftcol}{0pt}
\csromangathagroup{\csromangathastanza{\csromangathawak{สํสารมติจฺจ เกวลี โส, \\ อสิโต\csromansharedfootnote{f51e78b05f9f2cc7}{อนิสฺสิโต (สฺยา) ขุ 1. 358 ปิฏฺเฐ} ตาทิ ปวุจฺจเต ส พฺรหฺมา.}}}

\prose{\textbf{เวทคู}ติ เวโท วุจฺจติ จตูสุ มคฺเคสุ ญาณํ ฯเปฯ สพฺพํ เวทมติจฺจ เวทคู โสติ. \textbf{อภิชญฺญา}ติ อภิชาเนยฺย อาชาเนยฺย วิชาเนยฺย ปฏิวิชาเนยฺย ปฏิวิชฺเฌยฺยาติ ยํ พฺราหฺมณํ เวทคุมา\-ภิชญฺญา.}

\prose{\textbf{อกิญฺจนํ กามภเว อสตฺตนฺ}ติ \textbf{อกิญฺจนนฺ}ติ ราคกิญฺจนํ โทสกิญฺจนํ โมหกิญฺจนํ มานกิญฺจนํ ทิฏฺฐิกิญฺจนํ กิเลสกิญฺจนํ ทุจฺจริต\-กิญฺจนํ. ยสฺเสเต กิญฺจนา ปหีนา สมุจฺฉินฺนา วูปสนฺตา ปฏิปสฺสทฺธา อภพฺพุ\-ปฺปตฺติกา ญาณคฺคินา ทฑฺฒา, โส วุจฺจติ อกิญฺจโน. \textbf{กามา}ติ อุทฺทานโต เทฺว กามา วตฺถุกามา จ กิเลสกามา จ ฯเปฯ. อิเม วุจฺจนฺติ วตฺถุกามา ฯเปฯ. อิเม วุจฺจนฺติ กิเลสกามา. \textbf{ภวา}ติ เทฺว ภวา กมฺมภโว จ ปฏิสนฺธิโก จ \csromanfolio{85}ปุนพฺภโว ฯเปฯ อยํ ปฏิสนฺธิโก ปุนพฺภโว. \textbf{อกิญฺจนํ กามภเว อสตฺตนฺ}ติ อกิญฺจนํ ปุคฺคลํ กามภเว จ อสตฺตํ อลคฺคํ อลคฺคิตํ อปลิพุทฺธํ นิกฺขนฺตํ นิสฺสฏํ วิปฺปมุตฺตํ วิสญฺญุตฺตํ วิมริยาทิกเตน เจตสา วิหรนฺตนฺติ อกิญฺจนํ กามภเว อสตฺตํ.}

\proseitem{28}{\textbf{อทฺธา หิ โส โอฆมิมํ อตารี}ติ \textbf{อทฺธา}ติ เอกํสวจนํ ฯเปฯ อวตฺถาปน\-วจนเมตํ อทฺธาติ. \textbf{โอฆนฺ}ติ กาโมฆํ ภโวฆํ ทิฏฺโฐฆํ อวิชฺโชฆํ. \textbf{อตารี}ติ อุตฺตริ ปตริ สมติกฺกมิ วีติวตฺตยีติ อทฺธา หิ โส โอฆมิมํ อตาริ.}

\prose{\textbf{ติณฺโณ จ ปารํ อขิโล อกงฺโข}ติ ติณฺโณติ กาโมฆํ \textbf{ติณฺโณ}, ภโวฆํ ติณฺโณ, ทิฏฺโฐฆํ ติณฺโณ, อวิชฺโชฆํ ติณฺโณ, สํสารปถํ ติณฺโณ อุตฺติณฺโณ นิตฺถิณฺโณ\csromansharedfootnote{e57ed0227787ac85}{นิตฺติณฺโณ (สฺยา)} อติกฺกนฺโต สมติกฺกนฺโต วีติวตฺโต. โส วุตฺถวาโส\csromansharedfootnote{1636ad2b04a62d9b}{วุฏฺฐวาโส (สฺยา) ขุ 7. 16 ปิฏฺเฐปิ.} จิณฺณจรโณ คตทฺโธ คตทิโส คตโกฏิโก ปาลิต\-พฺรหฺมจริโย อุตฺตม\-ทิฏฺฐิปฺปตฺโต ภาวิตมคฺโค, ปหีนกิเลโส ปฏิวิทฺธา\-กุปฺโป สจฺฉิกต\-นิโรโธ, ทุกฺขํ ตสฺส ปริญฺญาตํ, สมุทโย ปหีโน, มคฺโค ภาวิโต, นิโรโธ สจฺฉิกโต, อภิญฺเญยฺยํ อภิญฺญาตํ, ปริญฺเญยฺยํ ปริญฺญาตํ, ปหาตพฺพํ ปหีนํ, ภาเวตพฺพํ ภาวิตํ, สจฺฉิกาตพฺพํ สจฺฉิกตํ. โส อุกฺขิตฺต\-ปลิโฆ สํกิณฺณ\-ปริกฺโข อพฺพุเฬฺหสิโก นิรคฺคโฬ อริโย ปนฺนทฺธโช ปนฺนภาโร วิสญฺญุตฺโต ปญฺจ\-งฺค\-วิปฺปหีโน ฉฬงฺค\-สมนฺนาคโต เอการกฺโข จตุราปสฺเสโน ปนุณฺณ\-ปจฺเจกสจฺโจ\csromansharedfootnote{01e8869c4ddf6fc6}{ปณุนฺนปจฺเจกสจฺโจ (ก)} สมวย\-สฏฺเฐ\-สโน อนาวิล\-สงฺกปฺโป ปสฺสทฺธ\-กายสงฺขาโร สุวิมุตฺตจิตฺโต สุวิมุตฺตปญฺโญ เกวลี วุสิตวา อุตฺตมปุริโส ปรมปุริโส ปรม\-ปตฺติ\-ปฺปตฺโต, โส เนว อาจินาติ น อปจินาติ อปจินิตฺวา ฐิโต, เนว ปชหติ น อุปาทิยติ ปชหิตฺวา ฐิโต, เนว วิสิเนติ น อุสฺสิเนติ วิสิเนตฺวา ฐิโต, เนว วิธูเปติ น สนฺธูเปติ วิธูเปตฺวา ฐิโต, อเสกฺเขน สีลกฺขนฺเธน สมนฺนาคตตฺตา ฐิโต, อเสกฺเขน สมาธิ\-กฺขนฺเธน. ปญฺญา\-กฺขนฺเธน. วิมุตฺติ\-กฺขนฺเธน. วิมุตฺติ\-ญาณ\-ทสฺสนกฺขนฺเธน สมนฺนาคตตฺตา ฐิโต, สจฺจํ สมฺปฏิปาทยิตฺวา\csromansharedfootnote{8fdbdd3423333248}{ปฏิปาทยิตฺวา (สฺยา)} ฐิโต, เอชํ สมติกฺกมิตฺวา ฐิโต, กิเลสคฺคิํ ปริยาทิยิตฺวา ฐิโต, อปริคมนตาย ฐิโต, กถํ\csromansharedfootnote{095d699da8d61db6}{กฏํ (สฺยา) กามสุตฺตนิทฺเทสฏฺฐกถา โอโลเกตพฺพา.} สมาทาย ฐิโต, วิมุตฺติ\-ปฏิเสว\-นตาย \csromanfolio{86}ฐิโต, เมตฺตาย ปาริสุทฺธิยา ฐิโต, กรุณาย. มุทิตาย. อุเปกฺขาย ปาริสุทฺธิยา ฐิโต, อจฺจนฺต\-ปาริสุทฺธิยา ฐิโต, อตมฺมยตาย\csromansharedfootnote{58131d55d44549f9}{อกมฺมญฺญตาย (สฺยา)} ปาริสุทฺธิยา ฐิโต, วิมุตฺตตฺตา ฐิโต, สนฺตุสฺสิตตฺตา ฐิโต, ขนฺธ\-ปริยนฺเต ฐิโต, ธาตุปริยนฺเต ฐิโต, อายตน\-ปริยนฺเต ฐิโต, คติปริยนฺเต ฐิโต, อุปปตฺติ\-ปริยนฺเต ฐิโต, ปฏิสนฺธิ\-ปริยนฺเต ฐิโต, ภวปริยนฺเต ฐิโต, สํสาร\-ปริยนฺเต ฐิโต, วฏฺฏปริยนฺเต ฐิโต, อนฺติมภเว ฐิโต, อนฺติเม \textbf{สมุสฺสเย ฐิโต, อนฺติม}\-\textbf{เทห}\-\textbf{ธโร อรหา.}}

\setlength{\csromangathapagewidth}{0pt}
\setlength{\csromangathawakwidth}{0pt}
\csromangathameasuregroup{ตสฺสายํ ปจฺฉิมโก ภโว, \\ ชาติมรณสํสาโร,}{\csromangathabat{ตสฺสายํ ปจฺฉิมโก ภโว,}{จริโมยํ สมุสฺสโย.} \\ \csromangathabat{ชาติมรณสํสาโร,}{นตฺถิ ตสฺส ปุนพฺภโวติ.}}
\csromangathasetleft{ตสฺสายํ ปจฺฉิมโก ภโว, \\ ชาติมรณสํสาโร,}
\csromangathagroup{\csromangathastanza{\csromangathabat{ตสฺสายํ ปจฺฉิมโก ภโว,}{จริโมยํ สมุสฺสโย.} \\ \csromangathabat{ชาติ\-มรณ\-สํสาโร,}{นตฺถิ ตสฺส ปุนพฺภโวติ.}}}

\proseitem{4}{\textbf{ติณฺโณ จ ปารนฺ}ติ ปารํ วุจฺจติ อมตํ นิพฺพานํ. โย โส สพฺพ\-สงฺขาร\-สมโถ สพฺพู\-ปธิ\-ปฏินิสฺสคฺโค ตณฺหกฺขโย วิราโค นิโรโธ นิพฺพานํ, โส ปารคโต ปารปฺปตฺโต อนฺตคโต อนฺตปฺปตฺโต โกฏิคโต โกฏิปฺปตฺโต ปริยนฺตคโต ปริยนฺต\-ปฺปตฺโต โวสานคโต โวสานปฺปตฺโต ตาณคโต ตาณปฺปตฺโต เลณคโต เลณปฺปตฺโต สรณคโต สรณปฺปตฺโต อภยคโต อภยปฺปตฺโต อจฺจุตคโต อจฺจุตปฺปตฺโต อมตคโต อมตปฺปตฺโต นิพฺพานคโต นิพฺพานปฺปตฺโต. โส วุตฺถวาโส จิณฺณจรโณ ฯเปฯ ชาติ\-มรณ\-สํสาโร นตฺถิ ตสฺส ปุนพฺภโวติ ติณฺโณ จ ปารํ.}

\prose{\textbf{อขิโล}ติ ราโค ขิโล, โทโส ขิโล, โมโห ขิโล, โกโธ ขิโล, อุปนาโห ขิโล ฯเปฯ สพฺพา\-กุสลา\-ภิสงฺขารา ขิลา. ยสฺเสเต ขิลา ปหีนา สมุจฺฉินฺนา วูปสนฺตา ปฏิปสฺสทฺธา อภพฺพุ\-ปฺปตฺติกา ญาณคฺคินา ทฑฺฒา, โส วุจฺจติ อขิโล. \textbf{อกงฺโข}ติ ทุกฺเข กงฺขา, ทุกฺขสมุทเย กงฺขา, ทุกฺขนิโรเธ กงฺขา, ทุกฺขนิโรธ\-คามินิยา ปฏิปทาย กงฺขา, ปุพฺพนฺเต กงฺขา, อปรนฺเต กงฺขา, ปุพฺพนฺตา\-ปรนฺเต กงฺขา, อิทปฺปจฺจยตา\-ปฏิจฺจ\-สมุปฺปนฺเนสุ ธมฺเมสุ กงฺขา, ยา เอวรูปา กงฺขา กงฺขายนา กงฺขายิตตฺตํ วิมติ วิจิกิจฺฉา เทฺวฬฺหกํ เทฺวธาปโถ สํสโย อเนกํสคฺคาโห อาสปฺปนา ปริสปฺปนา อปริโยคาหนา ฉมฺภิตตฺตํ จิตฺตสฺส มโนวิเลโข, ยสฺเสเต กงฺขา ปหีนา สมุจฺฉินฺนา วูปสนฺตา ปฏิปสฺสทฺธา อภพฺพุ\-ปฺปตฺติกา ญาณคฺคินา ทฑฺฒา, \csromanfolio{87}โส วุจฺจติ อกงฺโขติ ติณฺโณ จ ปารํ อขิโล อกงฺโข. เตนาห ภควา\csromandash{}}

\setlength{\csromangathapagewidth}{0pt}
\setlength{\csromangathawakwidth}{0pt}
\csromangathameasuregroup{}{ยํ พฺราหฺมณํ เวทคุมาภิชญฺญา, \\ อกิญฺจนํ กามภเว อสตฺตํ. \\ อทฺธา หิ โส โอฆมิมํ อตาริ, \\ ติณฺโณ จ ปารํ อขิโล อกงฺโขติ.}
\csromangathameasurewak{ยํ พฺราหฺมณํ เวทคุมาภิชญฺญา, \\ อกิญฺจนํ กามภเว อสตฺตํ. \\ อทฺธา หิ โส โอฆมิมํ อตาริ, \\ ติณฺโณ จ ปารํ อขิโล อกงฺโขติ.}
\setlength{\csromangathaleftcol}{0pt}
\csromangathagroup{\csromangathastanza{\csromangathawak{ยํ พฺราหฺมณํ เวทคุมา\-ภิชญฺญา, \\ อกิญฺจนํ กามภเว อสตฺตํ. \\ อทฺธา หิ โส โอฆมิมํ อตาริ, \\ ติณฺโณ จ ปารํ อขิโล อกงฺโขติ.}}}

\proseitem{29}{\textbf{วิทฺวา จ โย เวทคู นโร อิธ,}}

\prose{\textbf{ภวาภเว สงฺคมิมํ วิสชฺช.}}

\prose{\textbf{โส วีตตโณฺห อนีโฆ นิราโส,}}

\prose{\textbf{อตาริ โส ชาติชรนฺติ พฺรูมิ. (12)}}

\prose{\textbf{วิทฺวา จ โย เวทคู นโร อิธา}ติ \textbf{วิทฺวา}ติ วิชฺชาคโต ญาณี วิภาวี เมธาวี. \textbf{โย}ติ โย ยาทิโส ฯเปฯ มนุสฺโส วา. \textbf{เวทคู}ติ เวโท วุจฺจติ จตูสุ มคฺเคสุ ญาณํ. ปญฺญา ปญฺญินฺทฺริยํ ปญฺญาพลํ ธมฺมวิจย\-สมฺโพชฺฌงฺโค วีมํสา วิปสฺสนา สมฺมาทิฏฺฐิ\csromansharedfootnote{28924fcb9cbae8f8}{ญาณํ ฯเปฯ สพฺพเวทมติจฺจ เวทคู โสติ (สฺยา) ปสฺส ขุ 7. 158 ปิฏฺเฐ.}. เตหิ เวเทหิ ชาติ\-ชรา\-มรณสฺส อนฺตคโต อนฺตปฺปตฺโต โกฏิคโต โกฏิปฺปตฺโต ปริยนฺตคโต ปริยนฺต\-ปฺปตฺโต โวสานคโต โวสานปฺปตฺโต ตาณคโต ตาณปฺปตฺโต เลณคโต เลณปฺปตฺโต สรณคโต สรณปฺปตฺโต อภยคโต อภยปฺปตฺโต อจฺจุตคโต อจฺจุตปฺปตฺโต อมตคโต อมตปฺปตฺโต นิพฺพานคโต นิพฺพานปฺปตฺโต. เวทานํ วา อนฺตคโตติ เวทคู, เวเทหิ วา อนฺตคโตติ เวทคู, สตฺตนฺนํ วา ธมฺมานํ วิทิตตฺตา เวทคู, สกฺกายทิฏฺฐิ วิทิตา โหติ. วิจิกิจฺฉา. สีลพฺพต\-ปรามาโส. ราโค. โทโส. โมโห. มาโน วิทิโต โหติ, วิทิตาสฺส โหนฺติ ปาปกา อกุสลา ธมฺมา สํกิเลสิกา โปโนภวิกา สทรา ทุกฺขวิปากา อายติํ ชาติ\-ชรา\-มรณิยา.}

\prose{เวทานิ วิเจยฺย เกวลานิ, (สภิยาติ ภควา,)}

\prose{สมณานํ ยานีธตฺถิ พฺราหฺมณานํ.}

\setlength{\csromangathapagewidth}{0pt}
\setlength{\csromangathawakwidth}{0pt}
\csromangathameasuregroup{}{สพฺพเวทนาสุ วีตราโค, \\ สพฺพํ เวทมติจฺจ เวทคู โส.}
\csromangathameasurewak{สพฺพเวทนาสุ วีตราโค, \\ สพฺพํ เวทมติจฺจ เวทคู โส.}
\setlength{\csromangathaleftcol}{0pt}
\csromangathagroup{\csromangathastanza{\csromangathawak{สพฺพเวทนาสุ วีตราโค, \\ สพฺพํ เวทมติจฺจ เวทคู โส.}}}

\prose{\csromanfolio{88}\textbf{นโร}ติ สตฺโต นโร มานโว โปโส ปุคฺคโล ชีโว ชาคุ\csromansharedfootnote{202e5f9b49a6a41b}{ชาตุ (สฺยา)} ชนฺตุ อินฺทคุ\csromansharedfootnote{cfd6db94f21f8023}{อินฺทคู (สฺยา)} มนุโช. \textbf{อิธา}ติ อิมิสฺสา ทิฏฺฐิยา ฯเปฯ อิมสฺมิํ มนุสฺสโลเกติ วิทฺวา จ โย เวทคู นโร อิธ.}

\prose{\textbf{ภวาภเว สงฺคมิมํ วิสชฺชา}ติ ภวาภเวติ \textbf{ภวาภเว} กมฺมภเว ปุนพฺภเว กามภเว, กมฺมภเว กามภเว ปุนพฺภเว รูปภเว, กมฺมภเว รูปภเว ปุนพฺภเว อรูปภเว, กมฺมภเว อรูปภเว ปุนพฺภเว ปุนปฺปุนภเว, ปุนปฺปุน\-คติยา ปุนปฺปุนฺ\-เอาปปตฺติยา ปุนปฺปุน\-ปฏิสนฺธิยา ปุนปฺปุน- อตฺตภาวา\-ภินิพฺพตฺติยา. สงฺคาติ สตฺต สงฺคา ราคสงฺโค โทสสงฺโค โมหสงฺโค มานสงฺโค ทิฏฺฐิสงฺโค กิเลสสงฺโค ทุจฺจริตสงฺโค. วิสชฺชาติ สงฺเค โวสชฺเชตฺวา วา วิสชฺช. อถ วา สงฺเค พนฺเธ วิพนฺเธ อาพนฺเธ ลคฺเค ลคฺคิเต ปลิพุทฺเธ พนฺธเน โผฏยิตฺวา\csromansharedfootnote{b0d0dc2021a097af}{โมจยิตฺวา (สฺยา)} วา วิสชฺช. ยถา ยานํ วา วยฺหํ วา รถํ วา สกฏํ วา สนฺทมานิกํ วา สชฺชํ วิสชฺชํ กโรนฺติ วิโกเปนฺติ, เอวเมว เต สงฺเค โวสชฺเชตฺวา วา วิสชฺช. อถ วา สงฺเค พนฺเธ วิพนฺเธ อาพนฺเธ ลคฺเค ลคฺคิเต ปลิพุทฺเธ พนฺธเน โผฏยิตฺวา วา วิสชฺชาติ ภวาภเว สงฺคมิมํ วิสชฺช.}

\csromanmatikaanchor{mtk.26}
\csromantocmarkref{subsection}{5. โธตกมาณวปุจฺฉานิทฺเทส}{mtk.26}
\bookmark[dest={mtk.26},level=2]{5. โธตกมาณวปุจฺฉานิทฺเทส}
\prose{\textbf{โส วีตตโณฺห อนีโฆ นิราโส, อตาริ โส ชาติชรนฺติ พฺรูมี}ติ ตณฺหาติ รูป\textbf{ตณฺหา} ฯเปฯ ธมฺมตณฺหา. ยสฺเสสา ตณฺหา ปหีนา สมุจฺฉินฺนา วูปสนฺตา ปฏิปสฺสทฺธา อภพฺพุ\-ปฺปตฺติกา ญาณคฺคินา ทฑฺฒา, โส วุจฺจติ วีตตโณฺห วิคตตโณฺห จตฺตตโณฺห วนฺตตโณฺห มุตฺตตโณฺห ปหีนตโณฺห ปฏินิสฺสฏฺฐ\-ตโณฺห, วีตราโค จตฺตราโค ปหีนราโค ปฏินิสฺสฏฺฐ\-ราโค นิจฺฉาโต นิพฺพุโต สีติภูโต สุขปฺปฏิสํเวที พฺรหฺมภูเตน อตฺตนา วิหรตีติ โส วีตตโณฺห. \textbf{อนีโฆ}ติ ราโค นีโฆ, โทโส นีโฆ, โมโห นีโฆ, โกโธ นีโฆ, อุปนาโห นีโฆ ฯเปฯ สพฺพา\-กุสลา\-ภิสงฺขารา นีฆา. ยสฺเสเต นีฆา ปหีนา สมุจฺฉินฺนา วูปสนฺตา ปฏิปสฺสทฺธา อภพฺพุ\-ปฺปตฺติกา ญาณคฺคินา ทฑฺฒา, โส วุจฺจติ อนีโฆ. \textbf{นิราโส}ติ อาสา วุจฺจติ ตณฺหา. โย ราโค สาราโค ฯเปฯ อภิชฺฌา โลโภ อกุสลมูลํ. ยสฺเสสา อาสา ตณฺหา ปหีนา สมุจฺฉินฺนา วูปสนฺตา ปฏิปสฺสทฺธา อภพฺพุ\-ปฺปตฺติกา ญาณคฺคินา ทฑฺฒา, โส วุจฺจติ นิราโส. \textbf{ชาตี}ติ ยา เตสํ เตสํ สตฺตานํ ฯเปฯ อายตนานํ ปฏิลาโภ. \textbf{ชรา}ติ ยา เตสํ เตสํ \csromanfolio{89}สตฺตานํ ฯเปฯ อินฺทฺริยานํ ปริปาโก. อยํ วุจฺจติ ชรา. \textbf{โส วีตตโณฺห อนีโฆ นิราโส, อตาริ โส ชาติชรนฺติ พฺรูมี}ติ โย โส วีตตโณฺห อนีโฆ จ นิราโส จ, โส โข ชาติ\-ชรา\-มรณํ อตริ อุตฺตริ ปตริ สมติกฺกมิ วีติวตฺตยีติ พฺรูมิ อาจิกฺขามิ เทเสมิ ปญฺญเปมิ ปฏฺฐเปมิ วิวรามิ วิภชามิ อุตฺตานีกโรมิ ปกาเสมีติ โส วีตตโณฺห อนีโฆ นิราโส, อตาริ โส ชาติชรนฺติ พฺรูมิ. เตนาห ภควา\csromandash{}}

\setlength{\csromangathapagewidth}{0pt}
\setlength{\csromangathawakwidth}{0pt}
\csromangathameasuregroup{}{วิทฺวา จ โย เวทคู นโร อิธ, \\ ภวาภเว สงฺคมิมํ วิสชฺช. \\ โส วีตตโณฺห อนีโฆ นิราโส, \\ อตาริ โส ชาติชรนฺติ พฺรูมีติ.}
\csromangathameasurewak{วิทฺวา จ โย เวทคู นโร อิธ, \\ ภวาภเว สงฺคมิมํ วิสชฺช. \\ โส วีตตโณฺห อนีโฆ นิราโส, \\ อตาริ โส ชาติชรนฺติ พฺรูมีติ.}
\setlength{\csromangathaleftcol}{0pt}
\csromangathagroup{\csromangathastanza{\csromangathawak{วิทฺวา จ โย เวทคู นโร อิธ, \\ ภวาภเว สงฺคมิมํ วิสชฺช. \\ โส วีตตโณฺห อนีโฆ นิราโส, \\ อตาริ โส ชาติชรนฺติ พฺรูมีติ.}}}

\proseitem{29}{สห คาถา\-ปริโยสานา ฯเปฯ “สตฺถา เม ภนฺเต ภควา, สาวโกหมสฺมี”ติ.}

\vspace{\tipitakaparskip}
\csromancenter{เมตฺตคูมาณวปุจฺฉา\-นิทฺเทโส จตุตฺโถ.}
\csromanhangingbegin
\hangingheaditem{30}{\textbf{ปุจฺฉามิ ตํ ภควา พฺรูหิ เม’ตํ, (อิจฺจายสฺมา โธตโก,)}}
\hangingbody{\textbf{วาจาภิกงฺขามิ มเหสิ ตุยฺหํ.}}
\hangingbody{\textbf{ตว สุตฺวาน นิคฺโฆสํ, สิกฺเข นิพฺพานมตฺตโน.}}
\csromanhangingend

\prose{\textbf{ปุจฺฉามิ ตํ ภควา พฺรูหิ เม’ตนฺ}ติ \textbf{ปุจฺฉามี}ติ ติสฺโส ปุจฺฉา อทิฏฺฐโชตนา ปุจฺฉา ทิฏฺฐ\-สํสนฺทนา ปุจฺฉา วิมติจฺเฉทนา ปุจฺฉา ฯเปฯ. อิมา ติสฺโส ปุจฺฉา ฯเปฯ นิพฺพานปุจฺฉา. \textbf{ปุจฺฉามิ ตนฺ}ติ ปุจฺฉามิ ตํ, ยาจามิ ตํ, อชฺเฌสามิ ตํ, ปสาเทมิ ตํ, “กถยสฺสุ เม”ติ ปุจฺฉามิ ตํ. \textbf{ภควา}ติ คารวาธิวจนเมตํ ฯเปฯ สจฺฉิกา ปญฺญตฺติ, ยทิทํ ภควาติ. \textbf{พฺรูหิ เมตนฺ}ติ พฺรูหิ อาจิกฺขาหิ เทเสหิ ปญฺญเปหิ ปฏฺฐเปหิ วิวราหิ วิภชาหิ อุตฺตานีกโรหิ ปกาเสหีติ ปุจฺฉามิ ตํ ภควา พฺรูหิ \textbf{เม’ตํ.}}

\prose{\textbf{อิจฺจายสฺมา โธตโก}ติ \textbf{อิจฺจา}ติ ปทสนฺธิ ฯเปฯ. \textbf{อายสฺมา}ติ ปิยวจนํ ครุวจนํ สคารวสปฺปติสฺสาธิวจนเมตํ อายสฺมาติ. \textbf{โธตโก}ติ \csromanfolio{90}\textbf{ตสฺส พฺราหฺมณสฺส นามํ สงฺขา สมญฺญา ปญฺญตฺติ โวหาโร นามํ} \textbf{นามกมฺมํ นามเธยฺยํ นิรุตฺติ พฺยญฺชนํ อภิลาโปติ อิจฺจายสฺมา} \textbf{โธตโก.}}

\prose{\textbf{วาจาภิกงฺขามิ มเหสิ ตุยฺหนฺ}ติ ตุยฺหํ วจนํ พฺยปฺปถํ เทสนํ อนุสาสนํ อนุสิฏฺฐํ กงฺขามิ อภิกงฺขามิ อิจฺฉามิ สาทิยามิ ปตฺถยามิ ปิหยามิ อภิชปฺปามิ. \textbf{มเหสี}ติ กิํ มเหสิ ภควา, มหนฺตํ สีลกฺขนฺธํ เอสี คเวสี ปริเยสีติ มเหสิ ฯเปฯ. “กหํ นราสโภ”ติ มเหสีติ วาจาภิกงฺขามิ มเหสิ ตุยฺหํ.}

\prose{\textbf{ตว สุตฺวาน นิคฺโฆสนฺ}ติ ตุยฺหํ วจนํ พฺยปฺปถํ เทสนํ อนุสาสนํ อนุสิฏฺฐํ สุตฺวา สุณิตฺวา อุคฺคเหตฺวา อุปธารยิตฺวา อุปลกฺขยิตฺวาติ ตว สุตฺวาน นิคฺโฆสํ.}

\prose{\textbf{สิกฺเข นิพฺพานมตฺตโน}ติ สิกฺขาติ ติสฺโส \textbf{สิกฺขา} อธิสีลสิกฺขา อธิจิตฺต\-สิกฺขา อธิปญฺญา\-สิกฺขา ฯเปฯ อยํ อธิปญฺญา\-สิกฺขา. นิพฺพาน\textbf{มตฺตโน}ติ อตฺตโน ราคสฺส นิพฺพาปนาย, โทสสฺส นิพฺพาปนาย, โมหสฺส นิพฺพาปนาย, โกธสฺส นิพฺพาปนาย, อุปนาหสฺส นิพฺพาปนาย ฯเปฯ สพฺพา\-กุสลา\-ภิสงฺขารานํ สมาย อุปสมาย วูปสมาย นิพฺพาปนาย ปฏินิสฺสคฺคาย ปฏิปสฺสทฺธิยา อธิสีลมฺปิ สิกฺเขยฺย, อธิจิตฺตมฺปิ สิกฺเขยฺย, อธิปญฺญมฺปิ สิกฺเขยฺย, อิมา ติสฺโส สิกฺขาโย อาวชฺชนฺโต สิกฺเขยฺย, ชานนฺโต สิกฺเขยฺย, ปสฺสนฺโต สิกฺเขยฺย, ปจฺจเวกฺขนฺโต สิกฺเขยฺย, จิตฺตํ ปทหนฺโต สิกฺเขยฺย, สทฺธาย อธิมุจฺจนฺโต สิกฺเขยฺย, วีริยํ ปคฺคณฺหนฺโต สิกฺเขยฺย, สติํ อุปฏฺฐเปนฺโต สิกฺเขยฺย, จิตฺตํ สมาทหนฺโต สิกฺเขยฺย, ปญฺญาย ปชานนฺโต สิกฺเขยฺย, อภิญฺเญยฺยํ อภิชานนฺโต สิกฺเขยฺย, ปริญฺเญยฺยํ ปริชานนฺโต สิกฺเขยฺย, ปหาตพฺพํ ปชหนฺโต สิกฺเขยฺย, ภาเวตพฺพํ ภาเวนฺโต สิกฺเขยฺย, สจฺฉิกาตพฺพํ สจฺฉิกโรนฺโต สิกฺเขยฺย อาจเรยฺย สมาจเรยฺย สมาทาย วตฺเตยฺยาติ สิกฺเข นิพฺพานมตฺตโน. เตนาห โส พฺราหฺมโณ\csromandash{}}

\prose{ปุจฺฉามิ ตํ ภควา พฺรูหิ เมตํ, (อิจฺจายสฺมา โธตโก,)}

\prose{วาจาภิกงฺขามิ มเหสิ ตุยฺหํ.}

\setlength{\csromangathapagewidth}{0pt}
\setlength{\csromangathawakwidth}{0pt}
\csromangathameasuregroup{ตว สุตฺวาน นิคฺโฆสํ,}{\csromangathabat{ตว สุตฺวาน นิคฺโฆสํ,}{สิกฺเข นิพฺพาน\textbf{มตฺตโน}ติ.}}
\csromangathasetleft{ตว สุตฺวาน นิคฺโฆสํ,}
\csromangathagroup{\csromangathastanza{\csromangathabat{ตว สุตฺวาน นิคฺโฆสํ,}{สิกฺเข นิพฺพาน\textbf{มตฺตโน}ติ.}}}

\csromanhangingbegin
\hangingheaditem{31}{\csromanfolio{91}\textbf{เตนหาตปฺปํ กโรหิ, (โธตกาติ ภควา,)}}
\hangingbody{\textbf{อิเธว นิปโก สโต.}}
\hangingbody{\textbf{อิโต สุตฺวาน นิคฺโฆสํ, สิกฺเข นิพฺพานมตฺตโน.}}
\csromanhangingend

\prose{\textbf{เตนหาตปฺปํ กโรหี}ติ อาตปฺปํ กโรหิ, อุสฺสาหํ กโรหิ, อุสฺโสฬฺหิํ กโรหิ, ถามํ กโรหิ, ธิติํ กโรหิ, วีริยํ กโรหิ, ฉนฺทํ ชเนหิ สญฺชเนหิ อุปฏฺฐเปหิ สมุฏฺฐเปหิ นิพฺพตฺเตหิ อภินิพฺพตฺเตหีติ เตนหาตปฺปํ กโรหิ.}

\prose{\textbf{โธตกา}ติ ภควา ตํ พฺราหฺมณํ นาเมน อาลปติ. \textbf{ภควา}ติ คารวาธิวจนเมตํ ฯเปฯ สจฺฉิกา ปญฺญตฺติ, ยทิทํ ภควาติ โธตกาติ ภควา.}

\prose{\textbf{อิเธว นิปโก สโต}ติ \textbf{อิธา}ติ อิมิสฺสา ทิฏฺฐิยา อิมิสฺสา ขนฺติยา อิมิสฺสา รุจิยา อิมสฺมิํ อาทาเย อิมสฺมิํ ธมฺเม อิมสฺมิํ วินเย อิมสฺมิํ ธมฺมวินเย อิมสฺมิํ ปาวจเน อิมสฺมิํ พฺรหฺมจริเย อิมสฺมิํ สตฺถุสาสเน อิมสฺมิํ อตฺตภาเว อิมสฺมิํ มนุสฺสโลเก. \textbf{นิปโก}ติ นิปโก ปณฺฑิโต ปญฺญวา พุทฺธิมา ญาณี วิภาวี เมธาวี. \textbf{สโต}ติ จตูหิ การเณหิ สโต, กาเย กายานุปสฺสนา\-สติปฏฺฐานํ ภาเวนฺโต สโต ฯเปฯ โส วุจฺจติ สโตติ อิเธว นิปโก สโต.}

\prose{\textbf{อิโต สุตฺวาน นิคฺโฆสนฺ}ติ อิโต มยฺหํ วจนํ พฺยปฺปถํ เทสนํ อนุสาสนํ อนุสิฏฺฐํ สุตฺวา สุณิตฺวา อุคฺคณฺหิตฺวา อุปธารยิตฺวา อุปลกฺขยิตฺวาติ อิโต สุตฺวาน นิคฺโฆสํ.}

\prose{\textbf{สิกฺเข นิพฺพานมตฺตโน}ติ สิกฺขาติ ติสฺโส \textbf{สิกฺขา} อธิสีลสิกฺขา อธิจิตฺต\-สิกฺขา อธิปญฺญา\-สิกฺขา ฯเปฯ อยํ อธิปญฺญา\-สิกฺขา. นิพฺพาน\textbf{มตฺตโน}ติ อตฺตโน ราคสฺส นิพฺพาปนาย, โทสสฺส นิพฺพาปนาย, โมหสฺส นิพฺพาปนาย, โกธสฺส นิพฺพาปนาย, อุปนาหสฺส นิพฺพาปนาย ฯเปฯ สพฺพา\-กุสลา\-ภิสงฺขารานํ สมาย อุปสมาย วูปสมาย นิพฺพาปนาย ปฏินิสฺสคฺคาย ปฏิปสฺสทฺธิยา อธิสีลมฺปิ สิกฺเขยฺย, อธิจิตฺตมฺปิ สิกฺเขยฺย, อธิปญฺญมฺปิ สิกฺเขยฺย, อิมา ติสฺโส สิกฺขาโย อาวชฺชนฺโต สิกฺเขยฺย, ชานนฺโต สิกฺเขยฺย ฯเปฯ สจฺฉิกาตพฺพํ สจฺฉิกโรนฺโต สิกฺเขยฺย \csromanfolio{92}อาจเรยฺย สมาจเรยฺย, สมาทาย วตฺเตยฺยาติ สิกฺเข นิพฺพานมตฺตโน. เตนาห ภควา\csromandash{}}

\prose{เตนหาตปฺปํ กโรหิ, (โธตกาติ ภควา,)}

\prose{อิเธว นิปโก สโต.}

\setlength{\csromangathapagewidth}{0pt}
\setlength{\csromangathawakwidth}{0pt}
\csromangathameasuregroup{อิโต สุตฺวาน นิคฺโฆสํ,}{\csromangathabat{อิโต สุตฺวาน นิคฺโฆสํ,}{สิกฺเข นิพฺพานมตฺตโนติ.} \\ \textbf{ปสฺสามหํ เทว}\textbf{มนุสฺส}\textbf{โลเก,} \\ \textbf{อกิญฺจนํ พฺราหฺมณมิ}\textbf{ริยมานํ.} \\ \textbf{ตํ ตํ นมสฺสามิ สมนฺตจกฺขุ,} \\ \textbf{ปมุญฺจ มํ สกฺก กถํกถาหิ.}}
\csromangathameasurewak{\textbf{ปสฺสามหํ เทว}\textbf{มนุสฺส}\textbf{โลเก,} \\ \textbf{อกิญฺจนํ พฺราหฺมณมิ}\textbf{ริยมานํ.} \\ \textbf{ตํ ตํ นมสฺสามิ สมนฺตจกฺขุ,} \\ \textbf{ปมุญฺจ มํ สกฺก กถํกถาหิ.}}
\csromangathasetleft{อิโต สุตฺวาน นิคฺโฆสํ,}
\csromangathagroup{\csromangathastanza{\csromangathabat{อิโต สุตฺวาน นิคฺโฆสํ,}{สิกฺเข นิพฺพานมตฺตโนติ.}}\csromangathastanza{\csromangathawak{\csromangathaitemline{32}{\textbf{ปสฺสามหํ เทว}\-\textbf{มนุสฺส}\-\textbf{โลเก,}} \\ \csromangathacontline{\textbf{อกิญฺจนํ พฺราหฺมณมิ}\-\textbf{ริยมานํ.}} \\ \csromangathacontline{\textbf{ตํ ตํ นมสฺสามิ สมนฺตจกฺขุ,}} \\ \csromangathacontline{\textbf{ปมุญฺจ มํ สกฺก กถํกถาหิ.}}}}}

\prose{\textbf{ปสฺสามหํ เทว}\-\textbf{มนุสฺส}\-\textbf{โลเก}ติ เทวาติ ตโย \textbf{เทวา} สมฺมุติเทวา อุปปตฺติเทวา วิสุทฺธิเทวา. กตเม สมฺมุติเทวา. สมฺมุติเทวา วุจฺจนฺติ ราชาโน จ ราชกุมารา จ เทวิโย จ, อิเม วุจฺจนฺติ สมฺมุติเทวา. กตเม อุปปตฺติเทวา, อุปปตฺติเทวา วุจฺจนฺติ จาตุมหาราชิกา เทวา ตาวติํสา เทวา ยามา เทวา ตุสิตา เทวา นิมฺมานรตี เทวา ปรนิมฺมิต\-วส\-วตฺตี เทวา พฺรหฺมกายิกา เทวา, เย จ เทวา ตทุตฺตริ\csromansharedfootnote{69eb4ed415087a85}{ตตฺรุปริ (สฺยา)}, อิเม วุจฺจนฺติ อุปปตฺติเทวา. กตเม วิสุทฺธิเทวา, วิสุทฺธิเทวา วุจฺจนฺติ ตถาคต\-สาวกา อรหนฺโต ขีณาสวา, เย จ ปจฺเจกพุทฺธา, อิเม วุจฺจนฺติ วิสุทฺธิเทวา. ภควา สมฺมุติเทวานญฺจ อุปปตฺติเทวานญฺจ วิสุทฺธิเทวานญฺจ เทโว จ อติเทโว จ เทวาติเทโว จ สีหสีโห นาคนาโค คณิคณี มุนิมุนี ราชราชา. \textbf{ปสฺสามหํ เทว}\-\textbf{มนุสฺส}\-\textbf{โลเก}ติ มนุสฺสโลเก เทวํ ปสฺสามิ, อติเทวํ ปสฺสามิ, เทวาติเทวํ ปสฺสามิ ทกฺขามิ โอโลเกมิ นิชฺฌายามิ อุปปริกฺขามีติ ปสฺสามหํ เทว\-มนุสฺส\-โลเก.}

\prose{\textbf{อกิญฺจนํ พฺราหฺมณมิริยมานนฺ}ติ \textbf{อกิญฺจนนฺ}ติ ราคกิญฺจนํ โทสกิญฺจนํ โมหกิญฺจนํ มานกิญฺจนํ ทิฏฺฐิกิญฺจนํ กิเลสกิญฺจนํ ทุจฺจริต\-กิญฺจนํ, เต กิญฺจนา พุทฺธสฺส ภควโต ปหีนา อุจฺฉินฺนมูลา ตาลาวตฺถุกตา อนภาวํกตา อายติํ อนุปฺปาทธมฺมา, ตสฺมา พุทฺโธ อกิญฺจโน. \textbf{พฺราหฺมโณ}ติ ภควา สตฺตนฺนํ ธมฺมานํ พาหิตตฺตา พฺราหฺมโณ, สกฺกายทิฏฺฐิ พาหิตา โหติ, วิจิกิจฺฉา พาหิตา โหติ, สีลพฺพต\-ปรามาโส พาหิโต โหติ, ราโค พาหิโต โหติ, โทโส \csromanfolio{93}พาหิโต โหติ, โมโห พาหิโต โหติ, มาโน พาหิโต โหติ, พาหิตาสฺส โหนฺติ ปาปกา อกุสลา ธมฺมา สํกิเลสิกา โปโนภวิกา สทรา ทุกฺขวิปากา อายติํ ชาติ\-ชรา\-มรณิยา.}

\vspace{\tipitakaparskip}
\csromancenter{พาหิตฺวา สพฺพปาปกานิ, (สภิยาติ ภควา,)}
\prose{วิมโล สาธุสมาหิโต ฐิตตฺโต.}

\setlength{\csromangathapagewidth}{0pt}
\setlength{\csromangathawakwidth}{0pt}
\csromangathameasuregroup{}{สํสารมติจฺจ เกวลี โส, \\ อสิโต ตาทิ ปวุจฺจเต ส พฺรหฺมาติ.}
\csromangathameasurewak{สํสารมติจฺจ เกวลี โส, \\ อสิโต ตาทิ ปวุจฺจเต ส พฺรหฺมาติ.}
\setlength{\csromangathaleftcol}{0pt}
\csromangathagroup{\csromangathastanza{\csromangathawak{สํสารมติจฺจ เกวลี โส, \\ อสิโต ตาทิ ปวุจฺจเต ส พฺรหฺมาติ.}}}

\prose{\textbf{อิริยมานนฺ}ติ จรนฺตํ วิหรนฺตํ อิริยนฺตํ วตฺเตนฺตํ ปาเลนฺตํ ยเปนฺตํ ยาเปนฺตนฺติ อกิญฺจนํ พฺราหฺมณมิ\-ริยมานํ.}

\prose{\textbf{ตํ ตํ นมสฺสามิ สมนฺต}\-\textbf{จกฺขู}ติ \textbf{ตนฺ}ติ ภควนฺตํ ภณติ. \textbf{นมสฺสามี}ติ กาเยน วา นมสฺสามิ, วาจาย วา นมสฺสามิ, จิตฺเตน วา นมสฺสามิ, อนฺวตฺถ\-ปฏิปตฺติยา วา นมสฺสามิ, ธมฺมานุธมฺม\-ปฏิปตฺติยา วา นมสฺสามิ สกฺกโรมิ ครุํ กโรมิ มาเนมิ ปูเชมิ. \textbf{สมนฺต}\-\textbf{จกฺขู}ติ สมนฺตจกฺขุ วุจฺจติ สพฺพญฺญุต\-ญาณํ. ภควา สพฺพญฺญุต\-ญาเณน อุเปโต สมุเปโต อุปาคโต สมุปาคโต อุปปนฺโน สมุปปนฺโน สมนฺนาคโต.}

\setlength{\csromangathapagewidth}{0pt}
\setlength{\csromangathawakwidth}{0pt}
\csromangathameasuregroup{}{“น ตสฺส อทฺทิฏฺฐมิธตฺถิ\textsuperscript{1} กิญฺจิ, \\ อโถ อวิญฺญาตมชานิตพฺพํ. \\ สพฺพํ อภิญฺญาสิ ยทตฺถิ เนยฺยํ, \\ ตถาคโต เตน สมนฺตจกฺขู”ติ.}
\csromangathameasurewak{“น ตสฺส อทฺทิฏฺฐมิธตฺถิ\textsuperscript{1} กิญฺจิ, \\ อโถ อวิญฺญาตมชานิตพฺพํ. \\ สพฺพํ อภิญฺญาสิ ยทตฺถิ เนยฺยํ, \\ ตถาคโต เตน สมนฺตจกฺขู”ติ.}
\setlength{\csromangathaleftcol}{0pt}
\csromangathagroup{\csromangathastanza{\csromangathawak{“น ตสฺส อทฺทิฏฺฐมิธตฺถิ\csromansharedfootnote{50b27029c4e9d3bb}{อทิฏฺฐมิธตฺถิ (สฺยา, ก) ขุ 7. 280 ปิฏฺเฐปิ.} กิญฺจิ, \\ อโถ อวิญฺญาตมชานิตพฺพํ. \\ สพฺพํ อภิญฺญาสิ ยทตฺถิ เนยฺยํ, \\ ตถาคโต เตน สมนฺตจกฺขู”ติ.}}}

\prose{ตํ ตํ นมสฺสามิ สมนฺตจกฺขุ.}

\prose{\textbf{ปมุญฺจ มํ สกฺก กถํกถาหี}ติ \textbf{สกฺกา}ติ สกฺโก, ภควา สกฺยกุลา ปพฺพชิโตติปิ สกฺโก. อถ วา อฑฺโฒ\csromansharedfootnote{5551507e310a15f6}{อทฺโธ (สฺยา, ก)} มหทฺธโน ธนวาติปิ สกฺโก. ตสฺสิมานิ ธนานิ. เสยฺยถิทํ, สทฺธาธนํ สีลธนํ หิริธนํ โอตฺตปฺปธนํ สุตธนํ จาคธนํ ปญฺญาธนํ สติปฏฺฐานธนํ สมฺมปฺปธานธนํ อิทฺธิปาทธนํ อินฺทฺริยธนํ พลธนํ โพชฺฌงฺคธนํ มคฺคธนํ ผลธนํ นิพฺพานธนํ, อิเมหิ อเนกวิเธหิ ธนรตเนหิ อฑฺโฒ มหทฺธโน ธนวาติปิ สกฺโก. อถ วา สกฺโก ปหุ วิสวี อลมตฺโต สูโร วีโร วิกฺกนฺโต อภีรู อจฺฉมฺภี อนุตฺราสี อปลายี ปหีน\-ภยเภรโว วิคต\-โลมหํโสติปิ \csromanfolio{94}สกฺโก. กถํกถา วุจฺจติ วิจิกิจฺฉา, ทุกฺเข กงฺขา ทุกฺขสมุทเย กงฺขา ทุกฺขนิโรเธ กงฺขา ทุกฺขนิโรธ\-คามินิยา ปฏิปทาย กงฺขา ปุพฺพนฺเต กงฺขา อปรนฺเต กงฺขา ปุพฺพนฺตา\-ปรนฺเต กงฺขา อิทปฺปจฺจยตา\-ปฏิจฺจ\-สมุปฺปนฺเนสุ ธมฺเมสุ กงฺขา, ยา เอวรูปา กงฺขา กงฺขายนา กงฺขายิตตฺตํ วิมติ วิจิกิจฺฉา เทฺวฬฺหกํ เทฺวธาปโถ สํสโย อเนกํสคฺคาโห อาสปฺปนา ปริสปฺปนา อปริโยคาหนา ฉมฺภิตตฺตํ จิตฺตสฺส มโนวิเลโข. \textbf{ปมุญฺจ มํ สกฺก กถํกถาหี}ติ มุญฺจ มํ ปมุญฺจ มํ โมเจหิ มํ ปโมเจหิ มํ อุทฺธร มํ สมุทฺธร มํ วุฏฺฐาเปหิ มํ กถํกถา\-สลฺลโตติ ปมุญฺจ มํ สกฺก กถํกถาหิ. เตนาห โส พฺราหฺมโณ\csromandash{}}

\setlength{\csromangathapagewidth}{0pt}
\setlength{\csromangathawakwidth}{0pt}
\csromangathameasuregroup{}{ปสฺสามหํ เทวมนุสฺสโลเก, \\ อกิญฺจนํ พฺราหฺมณมิริยมานํ. \\ ตํ ตํ นมสฺสามิ สมนฺตจกฺขุ, \\ ปมุญฺจ มํ สกฺก กถํกถาหีติ. \\ \textbf{นาหํ สหิสฺสามิ ปโมจนาย,} \\ \textbf{กถํกถิํ โธตก กญฺจิ โลเก.} \\ \textbf{ธมฺมญฺจ เสฏฺฐํ อาชานมาโน,} \\ \textbf{เอวํ ตุวํ โอฆมิมํ ตเรสิ.}}
\csromangathameasurewak{ปสฺสามหํ เทวมนุสฺสโลเก, \\ อกิญฺจนํ พฺราหฺมณมิริยมานํ. \\ ตํ ตํ นมสฺสามิ สมนฺตจกฺขุ, \\ ปมุญฺจ มํ สกฺก กถํกถาหีติ. \\ \textbf{นาหํ สหิสฺสามิ ปโมจนาย,} \\ \textbf{กถํกถิํ โธตก กญฺจิ โลเก.} \\ \textbf{ธมฺมญฺจ เสฏฺฐํ อาชานมาโน,} \\ \textbf{เอวํ ตุวํ โอฆมิมํ ตเรสิ.}}
\setlength{\csromangathaleftcol}{0pt}
\csromangathagroup{\csromangathastanza{\csromangathawak{ปสฺสามหํ เทว\-มนุสฺส\-โลเก, \\ อกิญฺจนํ พฺราหฺมณมิ\-ริยมานํ. \\ ตํ ตํ นมสฺสามิ สมนฺตจกฺขุ, \\ ปมุญฺจ มํ สกฺก กถํกถาหีติ.}}\csromangathastanza{\csromangathawak{\csromangathaitemline{33}{\textbf{นาหํ สหิสฺสามิ ปโมจนาย,}} \\ \csromangathacontline{\textbf{กถํกถิํ โธตก กญฺจิ โลเก.}} \\ \csromangathacontline{\textbf{ธมฺมญฺจ เสฏฺฐํ อาชานมาโน,}} \\ \csromangathacontline{\textbf{เอวํ ตุวํ โอฆมิมํ ตเรสิ.}}}}}

\prose{\textbf{นาหํ สหิสฺสามิ}\csromansharedfootnote{7507f9de272e9532}{สมีหามิ (ก)} \textbf{ปโมจนายา}ติ นาหํ ตํ สกฺโกมิ มุญฺจิตุํ ปมุญฺจิตุํ โมเจตุํ ปโมเจตุํ อุทฺธริตุํ สมุทฺธริตุํ อุฏฺฐาเปตุํ สมุฏฺฐาเปตุํ กถํกถา\-สลฺลโตติ, เอวมฺปิ นาหํ สหิสฺสามิ ปโมจนาย. อถ วา น อีหามิ น สมีหามิ น อุสฺสหามิ น วายมามิ น อุสฺสาหํ กโรมิ น อุสฺโสฬฺหิํ กโรมิ น ถามํ กโรมิ น ธิติํ กโรมิ น วีริยํ กโรมิ น ฉนฺทํ ชเนมิ น สญฺชเนมิ น นิพฺพตฺเตมิ น อภินิพฺพตฺเตมิ อสฺสทฺเธ ปุคฺคเล อจฺฉนฺทิเก กุสีเต หีนวีริเย อปฺปฏิปชฺชมาเน ธมฺมเทสนายาติ, เอวมฺปิ นาหํ สหิสฺสามิ ปโมจนาย. อถ วา นตฺถญฺโญ โกจิ โมเจตา, เต ยทิ โมเจยฺยุํ, สเกน ถาเมน สเกน พเลน สเกน วีริเยน สเกน ปรกฺกเมน สเกน ปุริสถาเมน สเกน ปุริสพเลน สเกน ปุริสวีริเยน สเกน ปุริส\-ปรกฺกเมน อตฺตนา \csromanfolio{95}\textbf{สมฺมาปฏิปทํ อนุโลม}\-\textbf{ปฏิปทํ อปจฺจนีกปฏิปทํ} อนฺวตฺถ\-ปฏิปทํ ธมฺมานุธมฺม\-ปฏิปทํ ปฏิปชฺชมานา โมเจยฺยุนฺติ, เอวมฺปิ นาหํ สหิสฺสามิ ปโมจนาย.}

\prose{วุตฺตเญฺหตํ ภควตา\csromandash{}โส วต จุนฺท อตฺตนา ปลิป\-ปลิปนฺโน ปรํ ปลิป\-ปลิปนฺนํ อุทฺธริสฺสตีติ เนตํ ฐานํ วิชฺชติ, โส วต จุนฺท อตฺตนา อทนฺโต อวินีโต อปรินิพฺพุโต ปรํ ทเมสฺสติ วิเนสฺสติ ปรินิพฺพาเปสฺสตีติ เนตํ ฐานํ วิชฺชตีติ, เอวมฺปิ นาหํ สหิสฺสามิ ปโมจนาย. วุตฺตเญฺหตํ ภควตา\csromandash{}}

\setlength{\csromangathapagewidth}{0pt}
\setlength{\csromangathawakwidth}{0pt}
\csromangathameasuregroup{“อตฺตนา หิ\textsuperscript{1} กตํ ปาปํ, \\ อตฺตนา อกตํ ปาปํ, \\ สุทฺธี อสุทฺธิ ปจฺจตฺตํ,}{\csromangathabat{“อตฺตนา หิ\textsuperscript{1} กตํ ปาปํ,}{อตฺตนา สํกิลิสฺสติ.} \\ \csromangathabat{อตฺตนา อกตํ ปาปํ,}{อตฺตนาว วิสุชฺฌติ.} \\ \csromangathabat{สุทฺธี อสุทฺธิ ปจฺจตฺตํ,}{นาญฺโญ อญฺญํ วิโสธเย”ติ.}}
\csromangathasetleft{“อตฺตนา หิ\textsuperscript{1} กตํ ปาปํ, \\ อตฺตนา อกตํ ปาปํ, \\ สุทฺธี อสุทฺธิ ปจฺจตฺตํ,}
\csromangathagroup{\csromangathastanza{\csromangathabat{“อตฺตนา หิ\csromansharedfootnote{38e271ab9b04dc41}{อตฺตนาว (พหูสุ) ขุ 1. 38 ปิฏฺเฐ.} กตํ ปาปํ,}{อตฺตนา สํกิลิสฺสติ.} \\ \csromangathabat{อตฺตนา อกตํ ปาปํ,}{อตฺตนาว วิสุชฺฌติ.} \\ \csromangathabat{สุทฺธี อสุทฺธิ ปจฺจตฺตํ,}{นาญฺโญ อญฺญํ วิโสธเย”ติ.}}}

\prose{เอวมฺปิ นาหํ สหิสฺสามิ ปโมจนาย.}

\prose{วุตฺตเญฺหตํ ภควตา\csromandash{}เอวเมว โข พฺราหฺมณ ติฏฺฐเตว นิพฺพานํ, ติฏฺฐติ นิพฺพานคามิ\-มคฺโค, ติฏฺฐามหํ สมาทเปตา, อถ จ ปน มม สาวกา มยา เอวํ โอวทิยมานา เอวํ อนุสาสิยมานา อปฺเปกจฺเจ อจฺจนฺตนิฏฺฐํ นิพฺพานํ อาราเธนฺติ, เอกจฺเจ นาราเธนฺตีติ, เอตฺถ กฺยาหํ พฺราหฺมณ กโรมิ, มคฺคกฺขายี พฺราหฺมณ ตถาคโต มคฺคํ พุทฺโธ อาจิกฺขติ อตฺตนา ปฏิปชฺชมานา มุจฺเจยฺยุนฺติ\csromansharedfootnote{4663139aa6784f1c}{มุญฺเจยฺยุนฺติ (สฺยา)}, เอวมฺปิ นาหํ สหิสฺสามิ ปโมจนาย.}

\prose{\textbf{กถํกถิํ โธตก กญฺจิ โลเก}ติ กถํกถิํ ปุคฺคลํ สกงฺขํ สขิลํ สเทฺวฬฺหกํ สวิจิกิจฺฉํ. \textbf{กญฺจี}ติ กญฺจิ ขตฺติยํ วา พฺราหฺมณํ วา เวสฺสํ วา สุทฺทํ วา คหฏฺฐํ วา ปพฺพชิตํ วา เทวํ วา มนุสฺสํ วา. \textbf{โลเก}ติ อปายโลเก ฯเปฯ อายตนโลเกติ กถํกถิํ โธตก กญฺจิ โลเก.}

\prose{\textbf{ธมฺมญฺจ เสฏฺฐํ อาชานมาโน}ติ ธมฺมํ เสฏฺฐํ วุจฺจติ อมตํ นิพฺพานํ. โย โส สพฺพ\-สงฺขาร\-สมโถ สพฺพู\-ปธิ\-ปฏินิสฺสคฺโค ตณฺหกฺขโย วิราโค นิโรโธ นิพฺพานํ. \textbf{เสฏฺฐนฺ}ติ อคฺคํ เสฏฺฐํ วิเสฏฺฐํ ปาโมกฺขํ อุตฺตมํ ปวรํ ธมฺมํ อาชานมาโน วิชานมาโน ปฏิวิชานมาโน ปฏิวิชฺฌมาโนติ ธมฺมญฺจ เสฏฺฐํ อาชานมาโน.}

\proseitem{5}{\csromanfolio{96}\textbf{เอวํ ตุวํ โอฆมิมํ ตเรสี}ติ เอวํ กาโมฆํ ภโวฆํ ทิฏฺโฐฆํ อวิชฺโชฆํ ตเรยฺยาสิ อุตฺตเรยฺยาสิ ปตเรยฺยาสิ สมติกฺกเมยฺยาสิ วีติวตฺเตยฺยาสีติ เอวํ ตุวํ โอฆมิมํ ตเรสิ. เตนาห ภควา\csromandash{}}

\setlength{\csromangathapagewidth}{0pt}
\setlength{\csromangathawakwidth}{0pt}
\csromangathameasuregroup{}{นาหํ สหิสฺสามิ ปโมจนาย, \\ กถํกถิํ โธตก กญฺจิ โลเก. \\ ธมฺมญฺจ เสฏฺฐํ อาชานมาโน,}
\csromangathameasurewak{นาหํ สหิสฺสามิ ปโมจนาย, \\ กถํกถิํ โธตก กญฺจิ โลเก. \\ ธมฺมญฺจ เสฏฺฐํ อาชานมาโน,}
\setlength{\csromangathaleftcol}{0pt}
\csromangathagroup{\csromangathastanza{\csromangathawak{นาหํ สหิสฺสามิ ปโมจนาย, \\ กถํกถิํ โธตก กญฺจิ โลเก. \\ ธมฺมญฺจ เสฏฺฐํ อาชานมาโน,}}}

\prose{เอวํ ตุวํ โอฆมิมํ ตเรสีติ.}

\setlength{\csromangathapagewidth}{0pt}
\setlength{\csromangathawakwidth}{0pt}
\csromangathameasuregroup{}{\textbf{อนุสาส พฺรหฺเม กรุณายมาโน,} \\ \textbf{วิเวกธมฺมํ ยมหํ วิชญฺญํ.} \\ \textbf{ยถาหํ อากาโสว}\textsuperscript{1}\textbf{ อพฺยาปชฺชมาโน}\textsuperscript{1}\textbf{,} \\ \textbf{อิเธว สนฺโต อสิโต จเรยฺยํ.}}
\csromangathameasurewak{\textbf{อนุสาส พฺรหฺเม กรุณายมาโน,} \\ \textbf{วิเวกธมฺมํ ยมหํ วิชญฺญํ.} \\ \textbf{ยถาหํ อากาโสว}\textsuperscript{1}\textbf{ อพฺยาปชฺชมาโน}\textsuperscript{1}\textbf{,} \\ \textbf{อิเธว สนฺโต อสิโต จเรยฺยํ.}}
\setlength{\csromangathaleftcol}{0pt}
\csromangathagroup{\csromangathastanza{\csromangathawak{\csromangathaitemline{34}{\textbf{อนุสาส พฺรหฺเม กรุณายมาโน,}} \\ \csromangathacontline{\textbf{วิเวกธมฺมํ ยมหํ วิชญฺญํ.}} \\ \csromangathacontline{\textbf{ยถาหํ อากาโสว}\csromansharedfootnote{b98e5441dace0ba8}{อากาโส จ (สฺยา)}\textbf{ อพฺยาปชฺชมาโน}\csromansharedfootnote{07487fc68f51d4e3}{อพฺยาปชฺฌมาโน (สฺยา)}\textbf{,}} \\ \csromangathacontline{\textbf{อิเธว สนฺโต อสิโต จเรยฺยํ.}}}}}

\prose{\textbf{อนุสาส พฺรหฺเม กรุณายมาโน}ติ อนุสาส พฺรหฺเม อนุคฺคณฺห พฺรหฺเม อนุกมฺป พฺรหฺเมติ อนุสาส พฺรหฺเม. \textbf{กรุณายมาโน}ติ กรุณายมาโน อนุทยมาโน\csromansharedfootnote{ec5c57869577e46d}{อนุทฺทยมาโน (พหูสุ) 4. นตฺถิ ... สฺยา...โปตฺถเก.} อนุรกฺขมาโน อนุคฺคณฺหมาโน อนุกมฺปมาโนติ อนุสาส พฺรหฺเม กรูณายมาโน.}

\prose{\textbf{วิเวกธมฺมํ ยมหํ วิชญฺญนฺ}ติ วิเวกธมฺมํ วุจฺจติ อมตํ นิพฺพานํ. โย โส สพฺพ\-สงฺขาร\-สมโถ สพฺพู\-ปธิ\-ปฏินิสฺสคฺโค ตณฺหกฺขโย วิราโค นิโรโธ นิพฺพานํ. \textbf{ยมหํ วิชญฺญนฺ}ติ ยมหํ ชาเนยฺยํ อาชาเนยฺยํ วิชาเนยฺยํ ปฏิวิชาเนยฺยํ ปฏิวิชฺเฌยฺยํ อธิคจฺเฉยฺยํ ผสฺเสยฺยํ สจฺฉิ\-กเรยฺยนฺติ วิเวกธมฺมํ ยมหํ วิชญฺญํ.}

\prose{\textbf{ยถาหํ อากาโสว อพฺยาปชฺชมาโน}ติ ยถา อากาโส อ ปชฺชติ น คณฺหติ\csromansharedfootnote{3b66742ab76f2481}{หลิทฺเทน (สฺยา)} น พชฺฌติ น ปลิพชฺฌติ, เอวํ อปชฺชมาโน อคณฺหมาโน\csromansharedfootnote{76726718967ad02d}{นีเลน (สฺยา)} อพชฺฌมาโน อปลิพชฺฌมาโนติ เอวมฺปิ อากาโสว อพฺยาปชฺชมาโน. ยถา อากาโส น รชฺชติ ลาขาย วา หลิทฺทิยา\csromansharedfootnote{b1b89630d68a2d17}{อกิลิยมาโน (สฺยา)} วา นีลิยา\csromansharedfootnote{b2a5f4d4e3b3fa11}{นตฺถิ ... สฺยา...โปตฺถเก.} วา มญฺเชฏฺฐาย วา, เอวํ อรชฺชมาโน อทุสฺสมาโน อมุยฺหมาโน อกิลิสฺสมาโนติ\csromansharedfootnote{b2a5f4d4e3b3fa11}{นตฺถิ ... สฺยา...โปตฺถเก.} เอวมฺปิ อากาโสว อพฺยาปชฺชมาโน. ยถา อากาโส \csromanfolio{97}น กุปฺปติ น พฺยาปชฺชติ น ปติลียติ\csromansharedfootnote{f6b2122c4a5c9e96}{ปติฏฺฐิยติ (ก)} น ปฏิหญฺญติ, เอวํ อกุปฺปมาโน อพฺยาปชฺชมาโน อปฺปติลียมาโน อปฺปฏิหญฺญมาโน อปฺปฏิหตมาโนติ เอวมฺปิ อากาโสว อพฺยาปชฺชมาโน.}

\prose{\textbf{อิเธว สนฺโต อสิโต จเรยฺยนฺ}ติ อิเธว สนฺโตติ \textbf{อิเธว สนฺโต} อิเธว สมาโน อิเธว นิสินฺโน สมาโน อิมสฺมิํเยว อาสเน นิสินฺโน สมาโน อิมิสฺสาเยว ปริสาย นิสินฺโน สมาโนติ เอวมฺปิ อิเธว สนฺโต. อถ วา อิเธว สนฺโต อุปสนฺโต วูปสนฺโต นิพฺพุโต ปฏิปฺปสฺสทฺโธติ เอวมฺปิ อิเธว สนฺโต. \textbf{อสิโต}ติ เทฺว นิสฺสยา ตณฺหานิสฺสโย จ ทิฏฺฐินิสฺสโย จ ฯเปฯ อยํ ตณฺหานิสฺสโย ฯเปฯ อยํ ทิฏฺฐินิสฺสโย. ตณฺหานิสฺสยํ ปหาย, ทิฏฺฐินิสฺสยํ ปฏินิสฺสชฺชิตฺวา, จกฺขุํ อนิสฺสิโต, โสตํ อนิสฺสิโต, ฆานํ อนิสฺสิโต, ชิวฺหํ อนิสฺสิโต, กายํ อนิสฺสิโต, มนํ อนิสฺสิโต, รูเป. สทฺเท. คนฺเธ. รเส. โผฏฺฐพฺเพ. ธมฺเม. กุลํ. คณํ. อาวาสํ. ลาภํ. ยสํ. ปสํสํ. สุขํ. จีวรํ. ปิณฺฑปาตํ. เสนาสนํ. คิลานปจฺจย\-เภสชฺช\-ปริกฺขารํ. กามธาตุํ. รูปธาตุํ. อรูปธาตุํ. กามภวํ. รูปภวํ. อรูปภวํ. สญฺญาภวํ. อสญฺญาภวํ. เนว\-สญฺญา\-นา\-สญฺญาภวํ. เอกโวการภวํ. จตุโวการภวํ. ปญฺจโวการภวํ. อตีตํ. อนาคตํ. ปจฺจุปฺปนฺนํ. ทิฏฺฐ\-สุต\-มุต\-วิญฺญาตพฺเพ\csromansharedfootnote{bca1cd786bd42418}{ทิฏฺฐํ, สุตํ, มุตํ , วิญฺญาตํ, สพฺเพ. ขุ 7. 103 ปิฏฺเฐ, 189 ปิฏฺเฐปิ ปสฺสิตพฺพํ.} ธมฺเม อสิโต อนิสฺสิโต อนลฺลีโน อนุปคโต อนชฺโฌสิโต อนธิมุตฺโต นิกฺขนฺโต นิสฺสโฏ\csromansharedfootnote{054145bc8436a5be}{นิสฺสฏฺโฐ (สฺยา)} วิปฺปมุตฺโต วิสํยุตฺโต วิมริยาทิกเตน เจตสา. \textbf{จเรยฺยนฺ}ติ จเรยฺยํ วิหเรยฺยํ อิริเยยฺยํ วตฺเตยฺยํ ยเปยฺยํ ยาเปยฺยนฺติ อิเธว สนฺโต อสิโต จเรยฺยํ. เตนาห โส พฺราหฺมโณ\csromandash{}}

\setlength{\csromangathapagewidth}{0pt}
\setlength{\csromangathawakwidth}{0pt}
\csromangathameasuregroup{}{อนุสาส พฺรหฺเม กรุณายมาโน, \\ วิเวกธมฺมํ ยมหํ วิชญฺญํ. \\ ยถาหํ อากาโสว อพฺยาปชฺชมาโน,}
\csromangathameasurewak{อนุสาส พฺรหฺเม กรุณายมาโน, \\ วิเวกธมฺมํ ยมหํ วิชญฺญํ. \\ ยถาหํ อากาโสว อพฺยาปชฺชมาโน,}
\setlength{\csromangathaleftcol}{0pt}
\csromangathagroup{\csromangathastanza{\csromangathawak{อนุสาส พฺรหฺเม กรุณายมาโน, \\ วิเวกธมฺมํ ยมหํ วิชญฺญํ. \\ ยถาหํ อากาโสว อพฺยาปชฺชมาโน,}}}

\prose{อิเธว สนฺโต อสิโต จเรยฺยนฺติ.}

\csromanhangingbegin
\hangingheaditem{35}{\textbf{กิตฺตยิสฺสามิ เต สนฺติํ, (โธตกาติ ภควา,)}}
\hangingbody{\textbf{ทิฏฺเฐ ธมฺเม อนีติหํ.}}
\hangingbody{\textbf{ยํ วิทิตฺวา สโต จรํ, ตเร โลเก วิสตฺติกํ.}}
\csromanhangingend

\prose{\csromanfolio{98}\textbf{กิตฺตยิสฺสามิ เต สนฺตินฺ}ติ ราคสฺส สนฺติํ, โทสสฺส สนฺติํ, โมหสฺส สนฺติํ, โกธสฺส สนฺติํ, อุปนาหสฺส. มกฺขสฺส. ปฬาสสฺส. อิสฺสาย. มจฺฉริยสฺส. มายาย. สาเฐยฺยสฺส. ถมฺภสฺส. สารมฺภสฺส. มานสฺส. อติมานสฺส. มทสฺส. ปมาทสฺส. สพฺพกิเลสานํ. สพฺพ\-ทุจฺจริตานํ. สพฺพ\-ทรถานํ. สพฺพ\-ปริฬาหานํ. สพฺพ\-สนฺตาปานํ. สพฺพา\-กุสลา\-ภิสงฺขารานํ สนฺติํ อุปสนฺติํ วูปสนฺติํ นิพฺพุติํ ปฏิปฺปสฺสทฺธิํ กิตฺตยิสฺสามิ ปกิตฺตยิสฺสามิ อาจิกฺขิสฺสามิ เทเสสฺสามิ ปญฺญเปสฺสามิ ปฏฺฐเปสฺสามิ วิวริสฺสามิ วิภชิสฺสามิ อุตฺตานีกริสฺสามิ ปกาสิสฺสามีติ กิตฺตยิสฺสามิ เต สนฺติํ.}

\prose{\textbf{โธตกาติ ภควา}ติ โธตกาติ ภควา ตํ พฺราหฺมณํ นาเมน อาลปติ. ภควฺ\textbf{อา}ติ คารวาธิวจนเมตํ ฯเปฯ สจฺฉิกา ปญฺญตฺติ, ยทิทํ ภควาติ โธตกาติ ภควา.}

\prose{\textbf{ทิฏฺเฐ ธมฺเม อนีติหนฺ}ติ ทิฏฺเฐ ธมฺเมติ \textbf{ทิฏฺเฐ ธมฺเม} ญาเต ธมฺเม ตุลิเต ธมฺเม ตีริเต ธมฺเม วิภูเต ธมฺเม วิภาวิเต ธมฺเม “สพฺเพ สงฺขารา อนิจฺจา”ติ ฯเปฯ. “ยํ กิญฺจิ สมุทย\-ธมฺมํ, สพฺพํ ตํ นิโรธธมฺมนฺ”ติ ทิฏฺเฐ ธมฺเม ญาเต ธมฺเม ตุลิเต ธมฺเม ตีริเต ธมฺเม วิภาวิเต ธมฺเม วิภูเต ธมฺเมติ เอวมฺปิ ทิฏฺเฐ ธมฺเม. อถ วา ทุกฺเข ทิฏฺเฐ ทุกฺขํ กถยิสฺสามิ, สมุทเย ทิฏฺเฐ สมุทยํ กถยิสฺสามิ, มคฺเค ทิฏฺเฐ มคฺคํ กถยิสฺสามิ, นิโรเธ ทิฏฺเฐ นิโรธํ กถยิสฺสามีติ เอวมฺปิ ทิฏฺเฐ ธมฺเม. อถ วา สนฺทิฏฺฐิกํ อกาลิกํ เอหิปสฺสิกํ โอปเนยฺยิกํ\csromansharedfootnote{00ca7034c8e8547f}{โอปนยิกํ (สฺยา, ก)} ปจฺจตฺตํ เวทิตพฺพํ วิญฺญูหีติ เอวมฺปิ ทิฏฺเฐ ธมฺเม. \textbf{อนีติหนฺ}ติ น อิติหีติหํ น อิติกิราย น ปรมฺปราย น ปิฏก\-สมฺปทาย น ตกฺกเหตุ น นยเหตุ น อาการ\-ปริวิตกฺเกน น ทิฏฺฐิ\-นิชฺฌาน\-กฺขนฺติยา สามํ สยม\-ภิญฺญาตํ อตฺต\-ปจฺจกฺขธมฺมํ, ตํ กถยิสฺสามีติ ทิฏฺเฐ ธมฺเม อนีติหํ.}

\prose{\textbf{ยํ วิทิตฺวา สโต จรนฺ}ติ ยํ วิทิตํ กตฺวา ตุลยิตฺวา ตีรยิตฺวา วิภาวยิตฺวา วิภูตํ กตฺวา, “สพฺเพ สงฺขารา อนิจฺจา”ติ วิทิตํ กตฺวา ตุลยิตฺวา ตีรยิตฺวา วิภาวยิตฺวา วิภูตํ กตฺวา, “สพฺเพ สงฺขารา ทุกฺขา”ติ ฯเปฯ. “สพฺเพ ธมฺมา อนตฺตา”ติ ฯเปฯ. “ยํ กิญฺจิ สมุทย\-ธมฺมํ, สพฺพํ ตํ นิโรธธมฺมนฺ”ติ วิทิตํ กตฺวา ตุลยิตฺวา ตีรยิตฺวา วิภาวยิตฺวา วิภูตํ \csromanfolio{99}กตฺวา. \textbf{สโต}ติ จตูหิ การเณหิ สโต, กาเย กายานุปสฺสนา\-สติปฏฺฐานํ ภาเวนฺโต สโต ฯเปฯ โส วุจฺจติ สโต. \textbf{จรนฺ}ติ จรนฺโต วิหรนฺโต อิริยนฺโต วตฺเตนฺโต ปาเลนฺโต ยเปนฺโต ยาเปนฺโตติ ยํ วิทิตฺวา สโต จรํ.}

\prose{\textbf{ตเร โลเก วิสตฺติกนฺ}ติ วิสตฺติกา วุจฺจติ ตณฺหา. โย ราโค สาราโค ฯเปฯ อภิชฺฌา โลโภ อกุสลมูลํ. \textbf{วิสตฺติกา}ติ เกนฏฺเฐน วิสตฺติกา ฯเปฯ วิสฏา วิตฺถตาติ วิสตฺติกา. \textbf{โลเก}ติ อปายโลเก ฯเปฯ อายตนโลเก. \textbf{ตเร โลเก วิสตฺติกนฺ}ติ โลเก เวสา วิสตฺติกา, โลเก เวตํ วิสตฺติกํ สโต ตเรยฺย อุตฺตเรยฺย ปตเรยฺย สมติกฺกเมยฺย วีติวตฺเตยฺยาติ ตเร โลเก วิสตฺติกํ. เตนาห ภควา\csromandash{}}

\prose{กิตฺตยิสฺสามิ เต สนฺติํ, (โธตกาติ ภควา,)}

\prose{ทิฏฺเฐ ธมฺเม อนีติหํ.}

\setlength{\csromangathapagewidth}{0pt}
\setlength{\csromangathawakwidth}{0pt}
\csromangathameasuregroup{ยํ วิทิตฺวา สโต จรํ, \\ \textbf{ตญฺจาหํ อภินนฺทฺ}อามิ, \\ ยํ วิทิตฺวา สโต จรํ,}{\csromangathabat{ยํ วิทิตฺวา สโต จรํ,}{ตเร โลเก วิสตฺติกนฺติ.} \\ \csromangathabat{\textbf{ตญฺจาหํ อภินนฺทฺ}อามิ,}{มเหสิ สนฺติ’มุตฺตมํ.} \\ \csromangathabat{ยํ วิทิตฺวา สโต จรํ,}{ตเร โลเก วิสตฺติกํ.}}
\csromangathasetleft{ยํ วิทิตฺวา สโต จรํ, \\ \textbf{ตญฺจาหํ อภินนฺทฺ}อามิ, \\ ยํ วิทิตฺวา สโต จรํ,}
\csromangathagroup{\csromangathastanza{\csromangathabat{ยํ วิทิตฺวา สโต จรํ,}{ตเร โลเก วิสตฺติกนฺติ.}}\csromangathastanza{\csromangathaitembat{36}{\textbf{ตญฺจาหํ อภินนฺทฺ}อามิ,}{มเหสิ สนฺติ’มุตฺตมํ.} \\ \csromangathacontbat{ยํ วิทิตฺวา สโต จรํ,}{ตเร โลเก วิสตฺติกํ.}}}

\prose{\textbf{ตญฺจาหํ อภินนฺทามี}ติ \textbf{ตนฺ}ติ ตุยฺหํ วจนํ พฺยปฺปถํ เทสนํ อนุสาสนํ อนุสิฏฺฐํ นนฺทามิ อภินนฺทามิ โมทามิ อนุโมทามิ อิจฺฉามิ สาทิยามิ ปตฺถยามิ ปิหยามิ อภิชปฺปามีติ ตญฺจาหํ อภินนฺทามิ.}

\prose{\textbf{มเหสิ สนฺติ’มุตฺตมนฺ}ติ \textbf{มเหสี}ติ กิํ มเหสิ ภควา, มหนฺตํ สีลกฺขนฺธํ เอสี คเวสี ปริเยสีติ มเหสิ, มหนฺตํ สมาธิ\-กฺขนฺธํ ฯเปฯ “กหํ นราสโภ”ติ มเหสิ. \textbf{สนฺติ’มุตฺตมนฺ}ติ สนฺติ วุจฺจติ อมตํ นิพฺพานํ. โย โส สพฺพ\-สงฺขาร\-สมโถ สพฺพู\-ปธิ\-ปฏินิสฺสคฺโค ตณฺหกฺขโย วิราโค นิโรโธ นิพฺพานํ. \textbf{อุตฺตมนฺ}ติ อคฺคํ เสฏฺฐํ วิเสฏฺฐํ ปาโมกฺขํ อุตฺตมํ ปวรนฺติ มเหสิ สนฺติ’มุตฺตมํ.}

\prose{\textbf{ยํ วิทิตฺวา สโต จรนฺ}ติ ยํ วิทิตํ กตฺวา ฯเปฯ “สพฺเพ สงฺขารา อนิจฺจา”ติ วิทิตํ กตฺวา ตุลยิตฺวา ตีรยิตฺวา วิภาวยิตฺวา วิภูตํ กตฺวา, “สพฺเพ สงฺขารา ทุกฺขา”ติ. “สพฺเพ ธมฺมา อนตฺตา”ติ ฯเปฯ. “ยํ กิญฺจิ สมุทย\-ธมฺมํ, สพฺพํ ตํ นิโรธธมฺมนฺ”ติ วิทิตํ กตฺวา ตุลยิตฺวา ตีรยิตฺวา \csromanfolio{100}วิภาวยิตฺวา วิภูตํ กตฺวา. \textbf{สโต}ติ จตูหิ การเณหิ สโต, กาเย กายานุปสฺสนา\-สติปฏฺฐานํ ภเวนฺโต สโต ฯเปฯ โส วุจฺจติ สโต. \textbf{จรนฺ}ติ \textbf{จรนฺโต ฯเปฯ ยาเปนฺโตติ ยํ วิทิตฺวา สโต จรํ.}}

\prose{\textbf{ตเร โลเก วิสตฺติกนฺ}ติ วิสตฺติกา วุจฺจติ ตณฺหา. โย ราโค สาราโค ฯเปฯ อภิชฺฌา โลโภ อกุสลมูลํ. \textbf{วิสตฺติกา}ติ เกนฏฺเฐน วิสตฺติกา ฯเปฯ วิสฏา วิตฺถตาติ วิสตฺติกา. \textbf{โลเก}ติ อปายโลเก ฯเปฯ อายตนโลเก. \textbf{ตเร โลเก วิสตฺติกนฺ}ติ โลเก เวสา วิสตฺติกา, โลเก เวตํ วิสตฺติกํ สโต ตเรยฺยํ อุตฺตเรยฺยํ ฯเปฯ วีติวตฺเตยฺยนฺติ ตเร โลเก วิสตฺติกํ. เตนาห โส พฺราหฺมโณ\csromandash{}}

\setlength{\csromangathapagewidth}{0pt}
\setlength{\csromangathawakwidth}{0pt}
\csromangathameasuregroup{ตญฺจาหํ อภินนฺทามิ, \\ ยํ วิทิตฺวา สโต จรํ,}{\csromangathabat{ตญฺจาหํ อภินนฺทามิ,}{มเหสิ สนฺติ’มุตฺตมํ.} \\ \csromangathabat{ยํ วิทิตฺวา สโต จรํ,}{ตเร โลเก วิสตฺติกนฺติ.}}
\csromangathasetleft{ตญฺจาหํ อภินนฺทามิ, \\ ยํ วิทิตฺวา สโต จรํ,}
\csromangathagroup{\csromangathastanza{\csromangathabat{ตญฺจาหํ อภินนฺทามิ,}{มเหสิ สนฺติ’มุตฺตมํ.} \\ \csromangathabat{ยํ วิทิตฺวา สโต จรํ,}{ตเร โลเก วิสตฺติกนฺติ.}}}

\csromanhangingbegin
\hangingheaditem{37}{\textbf{ยํ กิญฺจิ สมฺปชานาสิ, (โธตกาติ ภควา,)}}
\hangingbody{\textbf{อุทฺธํ อโธ ติริยญฺจาปิ มชฺเฌ.}}
\hangingbody{\textbf{เอตํ วิทิตฺวา สงฺโคติ โลเก,}}
\hangingbody{\textbf{ภวาภวาย มากาสิ ตณฺหํ.}}
\csromanhangingend

\prose{\textbf{ยํ กิญฺจิ สมฺปชานาสี}ติ ยํ กิญฺจิ สมฺปชานาสิ อาชานาสิ ปฏิวิชานาสิ ปฏิวิชฺฌสีติ ยํ กิญฺจิ สมฺปชานาสิ. \textbf{โธตกาติ ภควา}ติ \textbf{โธตกา}ติ ภควา ตํ พฺราหฺมณํ นาเมน อาลปติ. \textbf{ภควา}ติ คารวาธิวจนเมตํ ฯเปฯ สจฺฉิกา ปญฺญตฺติ, ยทิทํ ภควาติ โธตกาติ ภควา.}

\csromanmatikaanchor{mtk.27}
\csromantocmarkref{subsection}{6. อุปสีวมาณวปุจฺฉานิทฺเทส}{mtk.27}
\bookmark[dest={mtk.27},level=2]{6. อุปสีวมาณวปุจฺฉานิทฺเทส}
\prose{\textbf{อุทฺธํ อโธ ติริยญฺจาปิ มชฺเฌ}ติ \textbf{อุทฺธนฺ}ติ อนาคตํ. \textbf{อโธ}ติ อตีตํ. \textbf{ติริยญฺจาปิ มชฺเฌ}ติ ปจฺจุปฺปนฺนํ. \textbf{อุทฺธนฺ}ติ เทวโลโก. \textbf{อโธ}ติ อปายโลโก. \textbf{ติริยญฺจาปิ มชฺเฌ}ติ มนุสฺสโลโก. อถ วา \textbf{อุทฺธนฺ}ติ กุสลา ธมฺมา. \textbf{อโธ}ติ อกุสลา ธมฺมา. \textbf{ติริยญฺจาปิ มชฺเฌ}ติ อพฺยากตา ธมฺมา. \textbf{อุทฺธนฺ}ติ อรูปธาตุ. \textbf{อโธ}ติ กามธาตุ. \textbf{ติริยญฺจาปิ มชฺเฌ}ติ รูปธาตุ. \textbf{อุทฺธนฺ}ติ สุขา เวทนา. \textbf{อโธ}ติ ทุกฺขา เวทนา. \textbf{ติริยญฺจาปิ มชฺเฌ}ติ อทุกฺขมสุขา เวทนา. \textbf{อุทฺธนฺ}ติ อุทฺธํ ปาทตลา. \csromanfolio{101}\textbf{อโธติ อโธ เกสมตฺถกา.}\textbf{ ติริยญฺจาปิ มชฺเฌติ เวมชฺเฌติ อุทฺธํ อโธ} \textbf{ติริยญฺจาปิ มชฺเฌ.}}

\prose{\textbf{เอตํ วิทิตฺวา สงฺโคติ โลเก}ติ สงฺโค เอโส, ลคฺคนํ เอตํ, พนฺธนํ เอตํ, ปลิโพโธ เอโสติ ญ ตฺวา ชานิตฺวา ตุลยิตฺวา ตีรยิตฺวา วิภาวยิตฺวา วิภูตํ กตฺวาติ เอตํ วิทิตฺวา สงฺโคติ โลเก.}

\prose{\textbf{ภวาภวาย มากาสิ ตณฺหนฺ}ติ ตณฺหาติ รูป\textbf{ตณฺหา} สทฺทตณฺหา ฯเปฯ ธมฺมตณฺหา. \textbf{ภวาภวายา}ติ ภวาภวาย กมฺมภวาย ปุนพฺภวาย กามภวาย, กมฺมภวาย กามภวาย ปุนพฺภวาย รูปภวาย, กมฺมภวาย รูปภวาย ปุนพฺภวาย อรูปภวาย, กมฺมภวาย อรูปภวาย ปุนพฺภวาย ปุนปฺปุนพฺภวาย, ปุนปฺปุน\-คติยา ปุนปฺปุนฺ\-เอาปปตฺติยา ปุนปฺปุน\-ปฏิสนฺธิยา ปุนปฺปุน- อตฺตภาวา\-ภินิพฺพตฺติยา ตณฺหํ มากาสิ มา ชเนสิ มา สญฺชเนสิ มา นิพฺพตฺเตสิ มาภินิพฺพตฺเตสิ, ปชห วิโนเทหิ พฺยนฺตีกโรหิ อนภาวํ คเมหีติ ภวาภวาย มากาสิ ตณฺหนฺติ. เตนาห ภควา\csromandash{}}

\prose{ยํ กิญฺจิ สมฺปชานาสิ, (โธตกาติ ภควา,)}

\prose{อุทฺธํ อโธ ติริยญฺจาปิ มชฺเฌ.}

\prose{เอตํ วิทิตฺวา สงฺโคติ โลเก,}

\prose{ภวาภวาย มากาสิ ตณฺหนฺติ.}

\prose{สห คาถา\-ปริโยสานา ฯเปฯ “สตฺถา เม ภนฺเต ภควา, สาวโกหมสฺมี”ติ.}

\vspace{\tipitakaparskip}
\csromancenter{โธตกมาณวปุจฺฉา\-นิทฺเทโส ปญฺจโม.}
\csromanhangingbegin
\hangingheaditem{38}{\textbf{เอโก อหํ สกฺก มหนฺตโมฆํ, (อิจฺจายสฺมา อุปสีโว,)}}
\hangingbody{\textbf{อนิสฺสิโต โน วิสหามิ ตาริตุํ.}}
\hangingbody{\textbf{อารมฺมณํ1} \textbf{พฺรูหิ สมนฺตจกฺขุ,}}
\hangingbody{\textbf{ยํ นิสฺสิโต โอฆมิมํ ตเรยฺยํ.}}
\csromanhangingend

\prose{\csromanfolio{102}\textbf{เอโก อหํ สกฺก มหนฺตโมฆนฺ}ติ \textbf{เอโก}ติ ปุคฺคโล วา เม ทุติโย นตฺถิ, ธมฺโม วา เม ทุติโย นตฺถิ, ยํ วา ปุคฺคลํ นิสฺสาย, ธมฺมํ วา นิสฺสาย มหนฺตํ กาโมฆํ ภโวฆํ ทิฏฺโฐฆํ อวิชฺโชฆํ ตเรยฺยํ อุตฺตเรยฺยํ ปตเรยฺยํ สมติกฺกเมยฺยํ วีติวตฺเตยฺยนฺติ. \textbf{สกฺกา}ติ สกฺโก, ภควา สกฺยกุลา ปพฺพชิโตติปิ สกฺโก. อถ วา อฑฺโฒ มหทฺธโน ธนวาติปิ สกฺโก. ตสฺสิมานิ ธนานิ. เสยฺยถิทํ, สทฺธาธนํ สีลธนํ หิริธนํ โอตฺตปฺปธนํ สุตธนํ จาคธนํ ปญฺญาธนํ สติปฏฺฐานธนํ ฯเปฯ นิพฺพานธนํ. อิเมหิ อเนเกหิ ธนรตเนหิ อฑฺโฒ มหทฺธโน ธนวาติปิ สกฺโก. อถ วา สกฺโก ปหุ วิสวี อลมตฺโต สูโร วีโร วิกฺกนฺโต อภีรู อฉมฺภี อนุตฺราสี อปลายี, ปหีน\-ภยเภรโว วิคต\-โลมหํโสติปิ สกฺโกติ เอโก อหํ สกฺก มหนฺตโมฆํ.}

\prose{\textbf{อิจฺจายสฺมา อุปสีโว}ติ \textbf{อิจฺจา}ติ ปทสนฺธิ ฯเปฯ. \textbf{อายสฺมา}ติ ปิยวจนํ ฯเปฯ. \textbf{อุปสีโว}ติ ตสฺส พฺราหฺมณสฺสนามํ ฯเปฯ อภิลาโปติ อิจฺจายสฺมา อุปสีโว.}

\prose{\textbf{อนิสฺสิโต โน วิสหามิ ตาริตุนฺ}ติ อนิสฺสิโตติ ปุคฺคลํ วา \textbf{อนิสฺสิโต} ธมฺมํ วา อนิสฺสิโต โน วิสหามิ น อุสฺสหามิ น สกฺโกมิ น ปฏิพโล. มหนฺตํ กาโมฆํ ภโวฆํ ทิฏฺโฐฆํ อวิชฺโชฆํ ตริตุํ อุตฺตริตุํ ปตริตุํ สมติกฺกมิตุํ วีติวตฺติตุนฺติ อนิสฺสิโต โน วิสหามิ ตาริตุํ.}

\prose{\textbf{อารมฺมณํ พฺรูหิ สมนฺต}\-\textbf{จกฺขู}ติ อารมฺมณํ อาลมฺพณํ นิสฺสยํ อุปนิสฺสยํ พฺรูหิ อาจิกฺขาหิ เทเสหิ ปญฺญเปหิ ปฏฺฐเปหิ วิวราหิ วิภชาหิ อุตฺตานีกโรหิ ปกาเสหิ. \textbf{สมนฺต}\-\textbf{จกฺขู}ติ สมนฺตจกฺขุ วุจฺจติ สพฺพญฺญุต\-ญาณํ. ภควา เตน สพฺพญฺญุต\-ญาเณน อุเปโต สมุเปโต อุปาคโต สมุปาคโต อุปปนฺโน สมุปปนฺโน สมนฺนาคโต.}

\setlength{\csromangathapagewidth}{0pt}
\setlength{\csromangathawakwidth}{0pt}
\csromangathameasuregroup{}{น ตสฺส อทิฏฺฐมิธตฺถิ กิญฺจิ, \\ อโถ อวิญฺญาตมชานิตพฺพํ. \\ สพฺพํ อภิญฺญาสิ ยทตฺถิ เนยฺยํ, \\ ตถาคโต เตน สมนฺตจกฺขูติ.}
\csromangathameasurewak{น ตสฺส อทิฏฺฐมิธตฺถิ กิญฺจิ, \\ อโถ อวิญฺญาตมชานิตพฺพํ. \\ สพฺพํ อภิญฺญาสิ ยทตฺถิ เนยฺยํ, \\ ตถาคโต เตน สมนฺตจกฺขูติ.}
\setlength{\csromangathaleftcol}{0pt}
\csromangathagroup{\csromangathastanza{\csromangathawak{น ตสฺส อทิฏฺฐมิ\-ธตฺถิ กิญฺจิ, \\ อโถ อวิญฺญาตมชานิตพฺพํ. \\ สพฺพํ อภิญฺญาสิ ยทตฺถิ เนยฺยํ, \\ ตถาคโต เตน สมนฺต\-จกฺขูติ.}}}

\prose{อารมฺมณํ พฺรูหิ สมนฺตจกฺขุ.}

\prose{\csromanfolio{103}\textbf{ยํ นิสฺสิโต โอฆมิมํ ตเรยฺยนฺ}ติ \textbf{ยํ นิสฺสิโต}ติ ยํ ปุคฺคลํ วา นิสฺสิโต ธมฺมํ วา นิสฺสิโต มหนฺตํ กาโมฆํ ภโวฆํ ทิฏฺโฐฆํ อวิชฺโชฆํ ตเรยฺยํ อุตฺตเรยฺยํ ปตเรยฺยํ สมติกฺกเมยฺยํ วีติวตฺเตยฺยนฺติ ยํ นิสฺสิโต โอฆมิมํ ตเรยฺยํ. เตนาห โส พฺราหฺมโณ\csromandash{}}

\prose{เอโก อหํ สกฺก มหนฺตโมฆํ, (อิจฺจายสฺมา อุปสีโว,)}

\prose{อนิสฺสิโต โน วิสหามิ ตาริตุํ.}

\prose{อารมฺมณํ พฺรูหิ สมนฺตจกฺขุ,}

\prose{ยํ นิสฺสิโต โอฆมิมํ ตเรยฺยนฺติ.}

\csromanhangingbegin
\hangingheaditem{39}{\textbf{อากิญฺจญฺญํ เปกฺขมาโน สติมา, (อุปสีวาติ ภควา,)}}
\hangingbody{\textbf{นตฺถีติ นิสฺสาย ตรสฺสุ โอฆํ.}}
\hangingbody{\textbf{กาเม ปหาย วิรโต กถาหิ,}}
\hangingbody{\textbf{ตณฺหกฺขยํ นตฺตมหา}’ภิปสฺส1.}
\csromanhangingend

\prose{\textbf{อากิญฺจญฺญํ เปกฺขมาโน สติมา}ติ โส พฺราหฺมโณ ปกติยา อากิญฺจญฺญายตน\-สมาปตฺติํ ลาภีเยว นิสฺสยํ น ชานาติ “อยํ เม นิสฺสโย”ติ. ตสฺส ภควา นิสฺสยญฺจ อาจิกฺขติ อุตฺตริญฺจ นิยฺยานปถํ. อากิญฺจญฺญายตน\-สมาปตฺติํ สโต สมาปชฺชิตฺวา ตโต วุฏฺฐหิตฺวา ตตฺถ ชาเต จิตฺตเจตสิเก ธมฺเม อนิจฺจโต เปกฺขมาโน, ทุกฺขโต. โรคโต. คณฺฑโต. สลฺลโต. อฆโต. อาพาธโต. ปรโต. ปโลกโต. ¢ติโต. อุปทฺทวโต. ภยโต. อุปสคฺคโต. จลโต. ปภงฺคุโต. อทฺธุวโต. อตาณโต. อเลณโต. อสรณโต. อสรณีภูตโต. ริตฺตโต. ตุจฺฉโต. สุญฺญโต. อนตฺตโต. อาทีนวโต. วิปริณาม\-ธมฺมโต. อสารกโต. อฆมูลโต. ภวโต. วิภวโต. สาสวโต. สงฺขตโต. มารามิสโต. ชาติธมฺมโต. ชราธมฺมโต. พฺยาธิธมฺมโต. มรณธมฺมโต. โสก\-ปริเทว\-ทุกฺข\-โทมนสฺสุ\-ปายาส\-ธมฺมโต. สมุทย\-ธมฺมโต. อตฺถงฺคมโต. อสฺสาทโต. อาทีนวโต. นิสฺสรณโต เปกฺขมาโน ทกฺขมาโน โอโลกยมาโน นิชฺฌายมาโน อุปปริกฺขมาโน.}

\prose{\csromanfolio{104}\textbf{สติมา}ติ ยา สติ อนุสฺสติ ปฏิสฺสติ ฯเปฯ สมฺมาสติ, อยํ วุจฺจติ สติ. อิมาย สติยา อุเปโต โหติ ฯเปฯ สมนฺนาคโต. โส วุจฺจติ สติมาติ อากิญฺจญฺญํ เปกฺขมาโน สติมา.}

\prose{\textbf{อุปสีวาติ ภควา}ติ อุปสีวาติ ภควา ตํ พฺราหฺมณํ นาเมน อาลปติ. \textbf{ภควา}ติ คารวาธิวจนเมตํ ฯเปฯ สจฺฉิกา ปญฺญตฺติ, ยทิทํ ภควาติ อุปสีวาติ ภควา.}

\prose{\textbf{นตฺถีติ นิสฺสาย ตรสฺสุ โอฆนฺ}ติ “นตฺถิ กิญฺจี”ติ อากิญฺจญฺญายตน\-สมาปตฺติ. กิํการณา “นตฺถิ กิญฺจี”ติ อากิญฺจญฺญายตน\-สมาปตฺติ, วิญฺญาณญฺจา\-ยตนสมาปตฺติํ สโต สมาปชฺชิตฺวา ตโต วุฏฺฐหิตฺวา ตญฺเญว วิญฺญาณํ อภาเวติ วิภาเวติ อนฺตรธาเปติ “นตฺถิ กิญฺจี”ติ ปสฺสติ, ตํการณา “นตฺถิ กิญฺจี”ติ อากิญฺจญฺญายตน\-สมาปตฺติํ นิสฺสาย อุปนิสฺสาย อาลมฺพณํ กริตฺวา กาโมฆํ ภโวฆํ ทิฏฺโฐฆํ อวิชฺโชฆํ ตรสฺสุ อุตฺตรสฺสุ ปตรสฺสุ สมติกฺกมสฺสุ วีติวตฺตสฺสูติ นตฺถีติ นิสฺสาย ตรสฺสุ โอฆํ.}

\prose{\textbf{กาเม ปหาย วิรโต กถาหี}ติ กามาติ อุทฺทานโต เทฺว \textbf{กามา} วตฺถุกามา จ กิเลสกามา จ ฯเปฯ. อิเม วุจฺจนฺติ วตฺถุกามา ฯเปฯ. อิเม วุจฺจนฺติ กิเลสกามา. \textbf{กาเม ปหายา}ติ วตฺถุกาเม ปริชานิตฺวา กิเลสกาเม ปหาย ปชหิตฺวา วิโนเทตฺวา พฺยนฺตีกริตฺวา อนภาวํ คเมตฺวาติ กาเม ปหาย. \textbf{วิรโต กถาหี}ติ กถํกถา วุจฺจติ วิจิกิจฺฉา. ทุกฺเข กงฺขา ฯเปฯ ฉมฺภิตตฺตํ จิตฺตสฺส มโนวิเลโข, กถํกถาย อารโต วิรโต ปฏิวิรโต นิกฺขนฺโต นิสฺสโฏ วิปฺปมุตฺโต วิสญฺญุตฺโต วิมริยาทิกเตน เจตสา วิหรตีติ เอวมฺปิ วิรโต กถาหิ. อถ วา ทฺวตฺติํสาย ติรจฺฉาน\-กถาย อารโต วิรโต ปฏิวิรโต นิกฺขนฺโต นิสฺสโฏ วิปฺปมุตฺโต วิสญฺญุตฺโต วิมริยาทิกเตน เจตสา วิหรตีติ เอวมฺปิ วิรโต กถาหีติ กาเม ปหาย วิรโต กถาหิ.}

\prose{\textbf{ตณฺหกฺขยํ นตฺตมหา’ภิปสฺสา}ติ ตณฺหาติ รูป\textbf{ตณฺหา} ฯเปฯ ธมฺมตณฺหา. \textbf{นตฺตํ} วุจฺจติ รตฺติ. \textbf{อโห}ติ ทิวโส. รตฺติญฺจ ทิวา จ ตณฺหกฺขยํ ราคกฺขยํ โทสกฺขยํ โมหกฺขยํ คติกฺขยํ อุปปตฺติ\-กฺขยํ ปฏิสนฺธิ\-กฺขยํ ภวกฺขยํ \csromanfolio{105}สํสารกฺขยํ วฏฺฏกฺขยํ ปสฺส อภิปสฺส ทกฺข โอโลกย นิชฺฌาย อุปปริกฺขาติ ตณฺหกฺขยํ นตฺตมหา’ภิปสฺส. เตนาห ภควา\csromandash{}}

\prose{อากิญฺจญฺญํ เปกฺขมาโน สติมา, (อุปสีวาติ ภควา,)}

\prose{นตฺถีติ นิสฺสาย ตรสฺสุ โอฆํ.}

\setlength{\csromangathapagewidth}{0pt}
\setlength{\csromangathawakwidth}{0pt}
\csromangathameasuregroup{}{กาเม ปหาย วิรโต กถาหิ, \\ ตณฺหกฺขยํ นตฺตมหา’ภิปสฺสาติ.}
\csromangathameasurewak{กาเม ปหาย วิรโต กถาหิ, \\ ตณฺหกฺขยํ นตฺตมหา’ภิปสฺสาติ.}
\setlength{\csromangathaleftcol}{0pt}
\csromangathagroup{\csromangathastanza{\csromangathawak{กาเม ปหาย วิรโต กถาหิ, \\ ตณฺหกฺขยํ นตฺตมหา’ภิปสฺสาติ.}}}

\proseitem{40}{\textbf{สพฺเพสุ กาเมสุ โย วีตราโค, (อิจฺจายสฺมา อุปสีโว,)}}

\prose{\textbf{อากิญฺจญฺญํ นิสฺสิโต หิตฺวา มญฺญํ.}}

\prose{\textbf{สญฺญาวิโมกฺเข ปรเมธิมุตฺโต,}}

\prose{\textbf{ติฏฺเฐ นุ โส ตตฺถ อนานุยายี.(3)}}

\prose{\textbf{สพฺเพสุ กาเมสุ โย วีตราโค}ติ \textbf{สพฺเพสู}ติ สพฺเพน สพฺพํ สพฺพถา สพฺพํ อเสสํ นิสฺเสสํ ปริยาทิย\-น\-วจนเมตํ สพฺเพสูติ. \textbf{กาเมสู}ติ \textbf{กามา}ติ อุทฺทานโต เทฺว กามา วตฺถุกามา จ กิเลสกามา จ ฯเปฯ. อิเม วุจฺจนฺติ วตฺถุกามา ฯเปฯ. อิเม วุจฺจนฺติ กิเลสกามา. สพฺเพสุ กาเมสุ โย วีตราโคติ สพฺเพสุ กาเมสุ โย วีตราโค วิคตราโค จตฺตราโค วนฺตราโค มุตฺตราโค ปหีนราโค ปฏินิสฺสฏฺฐ\-ราโค วิกฺขมฺภนโตติ สพฺเพสุ กาเมสุ โย วีตราโค.}

\prose{\textbf{อิจฺจายสฺมา อุปสีโว}ติ \textbf{อิจฺจา}ติ ปทสนฺธิ ฯเปฯ. \textbf{อายสฺมา}ติ ปิยวจนํ ฯเปฯ. \textbf{อุปสีโว}ติ ตสฺส พฺราหฺมณสฺส นามํ ฯเปฯ อภิลาโปติ อิจฺจายสฺมา อุปสีโว.}

\prose{\textbf{อากิญฺจญฺญํ นิสฺสิโต หิตฺวา มญฺญนฺ}ติ เหฏฺฐิมา ฉ สมาปตฺติโย หิตฺวา จชิตฺวา ปริจฺจชิตฺวา อติกฺกมิตฺวา สมติกฺกมิตฺวา วีติวตฺติตฺวา อากิญฺจญฺญายตน\-สมาปตฺติํ นิสฺสิโต อลฺลีโน อุปคโต สมุปคโต อชฺโฌสิโต อธิมุตฺโตติ อากิญฺจญฺญํ นิสฺสิโต หิตฺวา มญฺญํ.}

\prose{\textbf{สญฺญาวิโมกฺเข ปรเมธิมุตฺโต}ติ สญฺญาวิโมกฺขา วุจฺจนฺติ สตฺต สญฺญา\-สมาปตฺติโย. ตาสํ สญฺญา\-สมาปตฺตีนํ อากิญฺจญฺญายตน\-สมาปตฺติวิโมกฺโข\csromansharedfootnote{f2a352e3b6fdbb28}{วิโมกฺขา (ก) เอวมญฺเญสุ ปเทสุ พหุวจเนน.} อคฺโค จ เสฏฺโฐ จ วิเสฏฺโฐ จ ปาโมกฺโข จ อุตฺตโม จ ปวโร จ, ปรเม อคฺเค เสฏฺเฐ วิเสฏฺเฐ ปาโมกฺเข อุตฺตเม}

\prose{\csromanfolio{106}ปวเร อธิมุตฺติ\-วิโมกฺเขน อธิมุตฺโต ตตฺราธิมุตฺโต ตทธิมุตฺโต ตจฺจริโต ตพฺพหุโล ตคฺครุโก ตนฺนินฺโน ตปฺโปโณ ตปฺปพฺภาโร ตทธิมุตฺโต ตทธิปเตยฺโยติ สญฺญาวิโมกฺเข ปรเมธิมุตฺโต.}

\prose{\textbf{ติฏฺเฐ นุ โส ตตฺถ อนานุยายี}ติ \textbf{ติฏฺเฐ นู}ติ สํสยปุจฺฉา วิมติปุจฺฉา เทฺวฬฺหกปุจฺฉา อเนกํสปุจฺฉา “เอวํ นุ โข, นนุ โข, กิํ นุ โข, กถํ นุ โข”ติ ติฏฺเฐ นุ. \textbf{ตตฺถา}ติ อากิญฺจญฺญา\-ยตเน. \textbf{อนานุยายี}ติ อนานุยายี อวิจฺจมาโน\csromansharedfootnote{e1765460d896a90a}{อเวธมาโน (สฺยา)} อวิคจฺฉมาโน อนนฺตรธายมาโน อปริหายมาโน. อถ วา อรชฺชมาโน อทุสฺสมาโน อมุยฺหมาโน อกิลิสฺสมาโนติ ติฏฺเฐ นุ โส ตตฺถ อนานุยายี. เตนาห โส พฺราหฺมโณ\csromandash{}}

\prose{สพฺเพสุ กาเมสุ โย วีตราโค, (อิจฺจายสฺมา อุปสีโว,)}

\prose{อากิญฺจญฺญํ นิสฺสิโต หิตฺวา มญฺญํ.}

\prose{สญฺญาวิโมกฺเข ปรเมธิมุตฺโต,}

\prose{ติฏฺเฐ นุ โส ตตฺถ อนานุยายีติ.}

\csromanhangingbegin
\hangingheaditem{41}{\textbf{สพฺเพสุ กาเมสุ โย วีตราโค, (อุปสีวาติ ภควา,)}}
\hangingbody{อากิญฺจญฺญํ นิสฺสิโต หิตฺวา มญฺญํ.}
\hangingbody{สญฺญาวิโมกฺเข ปรเมธิมุตฺโต,}
\hangingbody{\textbf{ติฏฺเฐยฺย โส ตตฺถ อนานุยายี.}}
\csromanhangingend

\prose{\textbf{สพฺเพสุ กาเมสุ โย วีตราโค}ติ \textbf{สพฺเพสู}ติ สพฺเพน สพฺพํ สพฺพถา สพฺพํ อเสสํ นิสฺเสสํ ปริยาทิย\-น\-วจนเมตํ สพฺเพสูติ. \textbf{กาเมสู}ติ \textbf{กามา}ติ อุทฺทานโต เทฺว กามา วตฺถุกามา จ กิเลสกามา จ ฯเปฯ. อิเม วุจฺจนฺติ วตฺถุกามา ฯเปฯ. อิเม วุจฺจนฺติ กิเลสกามา. \textbf{สพฺเพสุ กาเมสุ โย วีตราโค}ติ สพฺเพสุ กาเมสุ โย วีตราโค ฯเปฯ ปฏินิสฺสฏฺฐ\-ราโค วิกฺขมฺภนโตติ สพฺเพสุ กาเมสุ โย วีตราโค.}

\prose{\textbf{อุปสีวาติ ภควา}ติ อุปสีวาติ ภควา ตํ พฺราหฺมณํ นาเมน อาลปติ. \textbf{ภควา}ติ คารวาธิวจนเมตํ ฯเปฯ สจฺฉิกา ปญฺญตฺติ, ยทิทํ ภควาติ อุปสีวาติ ภควา.}

\prose{\csromanfolio{107}\textbf{อากิญฺจญฺญํ นิสฺสิโต หิตฺวา มญฺญนฺ}ติ เหฏฺฐิมา ฉ สมาปตฺติโย หิตฺวา จชิตฺวา ปริจฺจชิตฺวา อติกฺกมิตฺวา สมติกฺกมิตฺวา วีติวตฺติตฺวา อากิญฺจญฺญายตน\-สมาปตฺติํ นิสฺสิโต อลฺลีโน อุปคโต สมุปคโต อชฺโฌสิโต อธิมุตฺโตติ อากิญฺจญฺญํ นิสฺสิโต หิตฺวา มญฺญํ.}

\prose{\textbf{สญฺญาวิโมกฺเข ปรเมธิมุตฺโต}ติ สญฺญาวิโมกฺขา วุจฺจนฺติ สตฺต สญฺญา\-สมาปตฺติโย. ตาสํ สญฺญา\-สมาปตฺตีนํ อากิญฺจญฺญายตน\-สมาปตฺติวิโมกฺโข อคฺโค จ เสฏฺโฐ จ วิเสฏฺโฐ จ ปาโมกฺโข จ อุตฺตโม จ ปวโร จ, ปรเม อคฺเค เสฏฺเฐ วิเสฏฺเฐ ปาโมกฺเข อุตฺตเม ปวเร อธิมุตฺติ\-วิโมกฺเขน อธิมุตฺโต ตตฺราธิมุตฺโต ตทธิมุตฺโต ฯเปฯ ตทธิปเตยฺโยติ สญฺญาวิโมกฺเข ปรเมธิมุตฺโต.}

\prose{\textbf{ติฏฺฐยฺย โส ตตฺถ อนานุยายี}ติ \textbf{ติฏฺเฐยฺยา}ติ ติฏฺเฐยฺย สฏฺฐิ\-กปฺปสหสฺสานิ. \textbf{ตตฺถา}ติ อากิญฺจญฺญา\-ยตเน. \textbf{อนานุยายี}ติ อนานุยายี อวิจฺจมาโน อวิคจฺฉมาโน อนนฺตรธายมาโน อปริหายมาโน. อถ วา อรชฺชมาโน อทุสฺสมาโน อมุยฺหมาโน อกิลิสฺสมาโนติ ติฏฺเฐยฺย โส ตตฺถ อนานุยายี. เตนาห ภควา\csromandash{}}

\prose{สพฺเพสุ กาเมสุ โย วีตราโค, (อุปสีวาติ ภควา,)}

\prose{อากิญฺจญฺญํ นิสฺสิโต หิตฺวา มญฺญํ.}

\prose{สญฺญาวิโมกฺเข ปรเมธิมุตฺโต,}

\prose{ติฏฺเฐยฺย โส ตตฺถ อนานุยายีติ.}

\setlength{\csromangathapagewidth}{0pt}
\setlength{\csromangathawakwidth}{0pt}
\csromangathameasuregroup{}{\textbf{ติฏฺเฐ เจ โส ตตฺถ อนานุยายี,} \\ \textbf{ปูคมฺปิ วสฺสานิ}\textsuperscript{1}\textbf{ สมนฺตจกฺขุ.} \\ \textbf{ตตฺเถว โส สีติสิยา วิมุตฺโต,} \\ \textbf{จเวถ วิญฺญาณํ ตถาวิธสฺส.}}
\csromangathameasurewak{\textbf{ติฏฺเฐ เจ โส ตตฺถ อนานุยายี,} \\ \textbf{ปูคมฺปิ วสฺสานิ}\textsuperscript{1}\textbf{ สมนฺตจกฺขุ.} \\ \textbf{ตตฺเถว โส สีติสิยา วิมุตฺโต,} \\ \textbf{จเวถ วิญฺญาณํ ตถาวิธสฺส.}}
\setlength{\csromangathaleftcol}{0pt}
\csromangathagroup{\csromangathastanza{\csromangathawak{\csromangathaitemline{42}{\textbf{ติฏฺเฐ เจ โส ตตฺถ อนานุยายี,}} \\ \csromangathacontline{\textbf{ปูคมฺปิ วสฺสานิ}\csromansharedfootnote{85b1df472a261fcd}{วสฺสานํ (สฺยา, ก)}\textbf{ สมนฺตจกฺขุ.}} \\ \csromangathacontline{\textbf{ตตฺเถว โส สีติสิยา วิมุตฺโต,}} \\ \csromangathacontline{\textbf{จเวถ วิญฺญาณํ ตถาวิธสฺส.}}}}}

\prose{\textbf{ติฏฺเฐ เจ โส ตตฺถ อนานุยายี}ติ สเจ โส ติฏฺเฐยฺย สฏฺฐิ\-กปฺปสหสฺสานิ. \textbf{ตตฺถา}ติ อากิญฺจญฺญา\-ยตเน. \textbf{อนานุยายี}ติ อนานุยายี อวิจฺจมาโน อวิคจฺฉมาโน อนนฺตรธายมาโน อปริหายมาโน. อถ วา อรชฺชมาโน อทุสฺสมาโน อมุยฺหมาโน อกิลิสฺสมาโนติ ติฏฺเฐ เจ โส ตตฺถ อนานุยายี.}

\prose{\csromanfolio{108}\textbf{ปูคมฺปิ วสฺสานิ สมนฺต}\-\textbf{จกฺขู}ติ \textbf{ปูคมฺปิ วสฺสานี}ติ ปูคมฺปิ วสฺสานิ พหูนิ วสฺสานิ\csromansharedfootnote{addb1fdfdbadcc8a}{พหุนฺนํ วสฺสานํ (สฺยา)} พหูนิ วสฺสสตานิ พหูนิ วสฺสสหสฺสานิ พหูนิ วสฺส\-สตสหสฺสานิ พหูนิ กปฺปานิ พหูนิ กปฺปสตานิ พหูนิ กปฺปสหสฺสานิ พหูนิ กปฺป\-สตสหสฺสานิ. \textbf{สมนฺต}\-\textbf{จกฺขู}ติ สมนฺตจกฺขุ วุจฺจติ สพฺพญฺญุต\-ญาณํ ฯเปฯ ตถาคโต เตน สมนฺต\-จกฺขูติ ปูคมฺปิ วสฺสานิ สมนฺตจกฺขุ.}

\prose{\textbf{ตตฺเถว โส สีติสิยา วิมุตฺโต, จเวถ วิญฺญาณํ ตถาวิธสฺสา}ติ ตตฺเถว โส สีติภาวม\-นุปฺปตฺโต นิจฺโจ ธุโว สสฺสโต อวิปริณาม\-ธมฺโม สสฺสติสมํ ตเถว ติฏฺเฐยฺย. อถ วา ตสฺส วิญฺญาณํ จเวยฺย อุจฺฉิชฺเชยฺย นสฺเสยฺย วินสฺเสยฺย น ภเวยฺยาติ ปุนพฺภว\-ปฏิสนฺธิ\-วิญฺญาณํ นิพฺพตฺเตยฺย กามธาตุยา วา รูปธาตุยา วา อรูปธาตุยา วาติ อากิญฺจญฺญา\-ยตนํ สมาปนฺนสฺส สสฺสตญฺจ อุจฺเฉทญฺจ ปุจฺฉติ. อุทาหุ ตตฺเถว อนุปาทิเสสาย นิพฺพานธาตุยา ปรินิพฺพาเยยฺย. อถ วา ตสฺส วิญฺญาณํ จเวยฺย ปุน ปฏิสนฺธิ\-วิญฺญาณํ นิพฺพตฺเตยฺย กามธาตุยา วา รูปธาตุยา วา อรูปธาตุยา วาติ อากิญฺจญฺญา\-ยตนํ อุปปนฺนสฺส ปรินิพฺพานญฺจ ปฏิสนฺธิญฺจ ปุจฺฉติ. \textbf{ตถาวิธสฺสา}ติ ตถาวิธสฺส ตาทิสสฺส ตสฺสณฺฐิตสฺส ตปฺปการสฺส ตปฺปฏิภาคสฺส อากิญฺจญฺญา\-ยตนํ อุปปนฺนสฺสาติ ตตฺเถว โส สีติสิยา วิมุตฺโต, จเวถ วิญฺญาณํ ตถาวิธสฺส. เตนาห โส พฺราหฺมโณ\csromandash{}}

\prose{ติฏฺเฐ เจ โส ตตฺถ อนานุยายี,}

\prose{ปูคมฺปิ วสฺสานิ สมนฺตจกฺขุ.}

\prose{\textbf{ตตฺเถว โส สีติสิยา วิมุตฺโต},}

\prose{จเวถ วิญฺญาณํ ตถาวิธสฺสาติ.}

\csromanhangingbegin
\hangingheaditem{43}{\textbf{อจฺจิ ยถา วาตเวเคน ขิตฺตา, (อุปสีวาติ ภควา,)}}
\hangingbody{\textbf{อตฺถํ ปเลติ น อุเปติ สงฺขํ.}}
\hangingbody{\textbf{เอวํ มุนี นามกายา วิมุตฺโต,}}
\hangingbody{\textbf{อตฺถํ ปเลติ น อุเปติ สงฺขํ.}}
\csromanhangingend

\prose{\textbf{อจฺจิ ยถา วาตเวเคน ขิตฺตา}ติ อจฺจิ วุจฺจติ ชาลสิขา. \textbf{วาตา}ติ ปุรตฺถิมา วาตา ปจฺฉิมา วาตา อุตฺตรา วาตา ทกฺขิณา วาตา สรชา วาตา อรชา วาตา สีตา วาตา อุณฺหา วาตา ปริตฺตา วาตา \csromanfolio{109}อธิมตฺตา วาตา\csromansharedfootnote{38b8c97df5c7a3d6}{กาฬวาตา (ก)} เวรมฺภวาตา ปกฺขวาตา สุปณฺณวาตา ตาลปณฺณวาตา วิธูปนวาตา. \textbf{วาตเวเคน ขิตฺตา}ติ วาตเวเคน ขิตฺตา\csromansharedfootnote{3732a37ef4e49561}{ขิตฺตํ (สฺยา) เอวมญฺเญสุ ปเทสุ นิคฺคหีตนฺตวเสน.} อุกฺขิตฺตา นุนฺนา ปนุนฺนา ขมฺภิตา วิกฺขมฺภิตาติ อจฺจิ ยถา วาตเวเคน ขิตฺตา. อุปสีวาติ ภควาติ \textbf{อุปสีวาติ ภควา} ตํ พฺราหฺมณํ นาเมน อาลปติ. \textbf{ภควา}ติ คารวาธิวจนเมตํ ฯเปฯ สจฺฉิกา ปญฺญตฺติ, ยทิทํ ภควาติ อุปสีวาติ ภควา.}

\prose{\textbf{อตฺถํ ปเลติ น อุเปติ สงฺขนฺ}ติ \textbf{อตฺถํ ปเลตี}ติ อตฺถํ ปเลติ, อตฺถํ คเมติ, อตฺถํ คจฺฉติ นิรุชฺฌติ วูปสมติ ปฏิปฺปสฺสมฺภติ. \textbf{น อุเปติ สงฺขนฺ}ติ สงฺขํ\csromansharedfootnote{9e7df0d4974abf8e}{อมุกํ นาม ทิสํ คโตติ สงฺขํ (สฺยา)} น อุเปติ, อุทฺเทสํ น อุเปติ, คณนํ น อุเปติ, ปณฺณตฺติํ น อุเปติ, “ปุรตฺถิมํ วา ทิสํ คตา, ปจฺฉิมํ วา ทิสํ คตา, อุตฺตรํ วา ทิสํ คตา, ทกฺขิณํ วา ทิสํ คตา, อุทฺธํ วา คตา, อโธ วา คตา, ติริยํ วา คตา, วิทิสํ วา คตา”ติ, โส เหตุ นตฺถิ, ปจฺจโย นตฺถิ, การณํ นตฺถิ, เยน สงฺขํ คจฺเฉยฺยาติ อตฺถํ ปเลติ น อุเปติ สงฺขํ.}

\prose{\textbf{เอวํ มุนี นามกายา วิมุตฺโต}ติ \textbf{เอวนฺ}ติ โอปมฺม\-สมฺปฏิปาทนํ. \textbf{มุนี}ติ โมนํ วุจฺจติ ญาณํ ฯเปฯ สงฺค\-ชาลม\-ติจฺจ โส มุนิ. \textbf{นามกายา วิมุตฺโต}ติ โส มุนิ ปกติยา ปุพฺเพว รูปกายา วิมุตฺโต. ตทงฺคํ สมติกฺกมา\csromansharedfootnote{b6e93ad63d1cb950}{ตทงฺคํ สมติกฺกมฺม (ก)} วิกฺขมฺภน\-ปฺปหาเนน ปหีโน. ตสฺส มุนิโน ภวนฺตํ อาคมฺม จตฺตาโร อริยมคฺคา ปฏิลทฺธา โหนฺติ, จตุนฺนํ อริยมคฺคานํ ปฏิลทฺธตฺตา นามกาโย จ รูปกาโย จ ปริญฺญาตา โหนฺติ. นามกายสฺส จ รูปกายสฺส จ ปริญฺญาตตฺตา นามกายา จ รูปกายา จ มุตฺโต วิมุตฺโต สุวิมุตฺโต อจฺจนฺต\-อนุปาทาวิโมกฺเขนาติ เอวํ มุนี นามกายา วิมุตฺโต.}

\prose{\textbf{อตฺถํ ปเลติ น อุเปติ สงฺขนฺ}ติ \textbf{อตฺถํ ปเลตี}ติ อนุปาทิเสสาย นิพฺพานธาตุยา ปรินิพฺพายติ. \textbf{น อุเปติ สงฺขนฺ}ติ อนุปาทิเสสาย นิพฺพานธาตุยา ปรินิพฺพุโต สงฺขํ น อุเปติ, อุทฺเทสํ น อุเปติ, คณนํ น อุเปติ, ปณฺณตฺติํ น อุเปติ “ขตฺติโย”ติ วา “พฺราหฺมโณ”ติ วา “เวสฺโส”ติ วา “สุทฺโท”ติ วา “คหฏฺโฐ”ติ วา “ปพฺพชิโต”ติ วา “เทโว”ติ วา “มนุสฺโส”ติ วา “รูปี”ติ วา “อรูปี”ติ วา “สญฺญี”ติ เวตฺ “อสญฺญี”ติ วา “เนว\-สญฺญี\-นา\-สญฺญี”ติ วา, โส เหตุ นตฺถิ, ปจฺจโย นตฺถิ, การณํ \csromanfolio{110}\textbf{นตฺถิ, เยน สงฺขํ คจฺเฉยฺยาติ อตฺถํ ปเลติ น อุเปติ สงฺขํ.}\textbf{ เตนาห} \textbf{ภควา\csromandash{}}}

\proseitem{6}{อจฺจิ ยถา วาตเวเคน ขิตฺตา, (อุปสีวาติ ภควา,)}

\prose{อตฺถํ ปเลติ น อุเปติ สงฺขํ.}

\setlength{\csromangathapagewidth}{0pt}
\setlength{\csromangathawakwidth}{0pt}
\csromangathameasuregroup{}{เอวํ มุนี นามกายา วิมุตฺโต, \\ อตฺถํ ปเลติ น อุเปติ สงฺขนฺติ.}
\csromangathameasurewak{เอวํ มุนี นามกายา วิมุตฺโต, \\ อตฺถํ ปเลติ น อุเปติ สงฺขนฺติ.}
\setlength{\csromangathaleftcol}{0pt}
\csromangathagroup{\csromangathastanza{\csromangathawak{เอวํ มุนี นามกายา วิมุตฺโต, \\ อตฺถํ ปเลติ น อุเปติ สงฺขนฺติ.}}}

\proseitem{44}{\textbf{อตฺถงฺคโต โส อุท วา โส นตฺถิ,}}

\prose{\textbf{อุทาหุ เว สสฺสติยา อโรโค.}}

\prose{\textbf{ตํ เม มุนี สาธุ วิยากโรหิ,}}

\prose{\textbf{อตฺถงฺคโต โส อุท วา โส นตฺถี}ติ โส อตฺถงฺคโต อุทาหุ นตฺถิ โส นิรุทฺโธ อุจฺฉินฺโน วินฏฺโฐติ อตฺถงฺคโต โส อุท วา โส นตฺถิ.}

\prose{\textbf{อุทาหุ เว สสฺสติยา อโรโค}ติ อุทาหุ นิจฺโจ ธุโว สสฺสโต อวิปริณาม\-ธมฺโม สสฺสติสมํ ตเถว ติฏฺเฐยฺยาติ อุทาหุ เว สสฺสติยา อโรโค.}

\prose{\textbf{ตํ เม มุนี สาธุ วิยากโรหี}ติ \textbf{ตนฺ}ติ ยํ ปุจฺฉามิ, ยํ ยาจามิ, ยํ อชฺเฌสามิ, ยํ ปสาเทมิ. \textbf{มุนี}ติ โมนํ วุจฺจติ ญาณํ ฯเปฯ สงฺค\-ชาลม\-ติจฺจ โส มุนิ. \textbf{สาธุ วิยากโรหี}ติ สาธุ อาจิกฺขาหิ เทเสหิ ปญฺญเปหิ ปฏฺฐเปหิ วิวราหิ วิภชาหิ อุตฺตานีกโรหิ ปกาเสหีติ ตํ เม มุนี สาธุ วิยากโรหิ.}

\prose{\textbf{ตถา หิ เต วิทิโต เอส ธมฺโม}ติ ตถา หิ เต วิทิโต ตุลิโต ตีริโต วิภูโต วิภาวิโต เอส ธมฺโมติ, ตถา หิ เต วิทิโต เอส ธมฺโม. เตนาห โส พฺราหฺมโณ\csromandash{}}

\prose{\textbf{อตฺถงฺคโต โส อุท วา โส นตฺถิ,}}

\prose{\textbf{อุทาหุ เว สสฺสติยา อโรโค.}}

\setlength{\csromangathapagewidth}{0pt}
\setlength{\csromangathawakwidth}{0pt}
\csromangathameasuregroup{}{\textbf{ตํ เม มุนี สาธุ วิยากโรหิ,} \\ ตถา หิ เต วิทิโต เอส ธมฺโมติ.}
\csromangathameasurewak{\textbf{ตํ เม มุนี สาธุ วิยากโรหิ,} \\ ตถา หิ เต วิทิโต เอส ธมฺโมติ.}
\setlength{\csromangathaleftcol}{0pt}
\csromangathagroup{\csromangathastanza{\csromangathawak{\textbf{ตํ เม มุนี สาธุ วิยากโรหิ,} \\ ตถา หิ เต วิทิโต เอส ธมฺโมติ.}}}

\csromanhangingbegin
\hangingheaditem{45}{\csromanfolio{111}\textbf{อตฺถงฺคตสฺส น ปมาณ’มตฺถิ, (อุปสีวาติ ภควา,)}}
\hangingbody{\textbf{เยน นํ วชฺชุํ ตํ ตสฺส นตฺถิ.}}
\hangingbody{\textbf{สพฺเพสุ ธมฺเมสุ สมูหเตสุ,}}
\hangingbody{\textbf{สมูหตา วาทปถาปิ สพฺเพ.}}
\csromanhangingend

\prose{\textbf{อตฺถงฺคตสฺส น ปมาณ’มตฺถี}ติ อตฺถงฺคตสฺส อนุปาทิเสสาย นิพฺพานธาตุยา ปรินิพฺพุตสฺส รูปปมาณํ นตฺถิ. เวทนาปมาณํ นตฺถิ. สญฺญาปมาณํ นตฺถิ. สงฺขาร\-ปมาณํ นตฺถิ. วิญฺญาณปมาณํ นตฺถิ น สติ น สํวิชฺชติ นุปลพฺภติ ปหีนํ สมุจฺฉินฺนํ วูปสนฺตํ ปฏิปฺปสฺสทฺธํ อภพฺพุ\-ปฺปตฺติกํ ญาณคฺคินา ทฑฺฒนฺติ อตฺถงฺคตสฺส น ปมาณ’มตฺถิ. อุปสีวาติ ภควาติ \textbf{อุปสีวาติ ภควา} ตํ พฺราหฺมณํ นาเมน อาลปติ. \textbf{ภควา}ติ คารวาธิวจนเมตํ ฯเปฯ สจฺฉิกา ปญฺญตฺติ, ยทิทํ ภควาติ อุปสีวาติ ภควา.}

\prose{\textbf{เยน นํ วชฺชุํ ตํ ตสฺส นตฺถี}ติ เยน ตํ ราเคน\csromansharedfootnote{8ae42ebb935675b3}{เยน ราเคน (สฺยา, ก) ขุ 7. 192 ปิฏฺเฐ.} วเทยฺยุํ. เยน โทเสน วเทยฺยุํ. เยน โมเหน วเทยฺยุํ. เยน มาเนน วเทยฺยุํ. ยาย ทิฏฺฐิยา วเทยฺยุํ. เยน อุทฺธจฺเจน วเทยฺยุํ. ยาย วิจิกิจฺฉาย วเทยฺยุํ. เยหิ อนุสเยหิ วเทยฺยุํ “รตฺโต”ติ วา “ทุฏฺโฐ”ติ วา “มูโฬฺห”ติ วา “วินิพทฺโธ”ติ วา “ปรามฏฺโฐ”ติ วา “วิกฺเขปคโต”ติ วา “อนิฏฺฐงฺคโต”ติ\csromansharedfootnote{b9ae1d9f1e81d2c1}{อนิฏฺฐา คโตติ (ก)} วา “ถามคโต”ติ วา, เต อภิสงฺขารา ปหีนา, อภิสงฺขารานํ ปหีนตฺตา คติยา เยน ตํ วเทยฺยุํ “เนรยิโก”ติ วา “ติรจฺฉาน\-โยนิโก”ติ วา “เปตฺติวิสยิโก”ติ วา “มนุสฺโส”ติ วา “เทโว”ติ วา “รูปี”ติ วา “อรูปี”ติ วา “สญฺญี”ติ วา “อสญฺญี”ติ วา “เนว\-สญฺญี\-นา\-สญฺญี”ติ วา, โส เหตุ นตฺถิ, ปจฺจโย นตฺถิ, การณํ นตฺถิ. เยน วเทยฺยุํ กเถยฺยุํ ภเณยฺยุํ ทีเปยฺยุํ โวหเรยฺยุนฺติ เยน นํ วชฺชุํ ตํ ตสฺส นตฺถิ.}

\prose{\textbf{สพฺเพสุ ธมฺเมสุ สมูหเตสู}ติ สพฺเพสุ ธมฺเมสุ สพฺเพสุ ขนฺเธสุ สพฺเพสุ อายตเนสุ สพฺพาสุ ธาตูสุ สพฺพาสุ คตีสุ สพฺพาสุ อุปปตฺตีสุ สพฺพาสุ ปฏิสนฺธีสุ สพฺเพสุ ภเวสุ สพฺเพสุ สํสาเรสุ สพฺเพสุ วฏฺเฏสุ อูหเตสุ สมูหเตสุ อุทฺธเตสุ สมุทฺธเตสุ อุปฺปาฏิเตสุ สมุปฺปาฏิเตสุ ปหีเนสุ สมุจฺฉินฺเนสุ วูปสนฺเตสุ ปฏิปฺปสฺสทฺเธสุ อภพฺพุ\-ปฺปตฺติเกสุ ญาณคฺคินา ทฑฺเฒสูติ สพฺเพสุ ธมฺเมสุ สมูหเตสุ.}

\csromanmatikaanchor{mtk.28}
\csromantocmarkref{subsection}{7. นนฺทมาณวปุจฺฉานิทฺเทส}{mtk.28}
\bookmark[dest={mtk.28},level=2]{7. นนฺทมาณวปุจฺฉานิทฺเทส}
\prose{\csromanfolio{112}\textbf{สมูหตา วาทปถาปิ สพฺเพ}ติ วาทปถา วุจฺจนฺติ กิเลสา จ ขนฺธา จ อภิสงฺขารา จ. ตสฺส วาทา จ วาทปถา จ อธิวจนานิ จ อธิวจน\-ปถา จ นิรุตฺติ จ นิรุตฺติปถา จ ปญฺญตฺติ จ ปญฺญตฺติปถา จ อูหตา สมูหตา อุทฺธตา สมุทฺธตา อุปฺปาฏิตา สมุปฺปาฏิตา ปหีนา สมุจฺฉินฺนา วูปสนฺตา ปฏิปฺปสฺสทฺธา อภพฺพุ\-ปฺปตฺติกา ญาณคฺคินา ทฑฺฑาติ สมูหตา วาทปถาปิ สพฺเพ. เตนาห ภควา\csromandash{}}

\prose{อตฺถงฺคตสฺส น ปมาณมตฺถิ, (อุปสีวาติ ภควา,)}

\prose{เยน นํ วชฺชุํ ตํ ตสฺส นตฺถิ.}

\prose{สพฺเพสุ ธมฺเมสุ สมูหเตสุ,}

\prose{สมูหตา วาทปถาปิ สพฺเพติ.}

\prose{สห คาถา\-ปริโยสานา เย เต พฺราหฺมเณน สทฺธิํ ฯเปฯ ปญฺชลิโก นมสฺสมาโน นิสินฺโน โหติ “สตฺถา เม ภนฺเต ภควา, สาวโกหมสฺมี”ติ.}

\vspace{\tipitakaparskip}
\csromancenter{อุปสีวมาณวปุจฺฉา\-นิทฺเทโส ฉฏฺโฐ.}
\csromanhangingbegin
\hangingheaditem{46}{“\textbf{สนฺติ โลเก มุนโย”, (อิจฺจายสฺมา นนฺโท,)}}
\hangingbody{\textbf{ชนา วทนฺติ ตยิทํ กถํสุ.}}
\hangingbody{\textbf{ญาณูปปนฺนํ มุนิ โน วทนฺติ,}}
\hangingbody{\textbf{อุทาหุ เว ชีวิเตนูปปฺ}อนฺนํ1.}
\csromanhangingend

\prose{\textbf{สนฺติ โลเก มุนโย}ติ \textbf{สนฺตี}ติ สนฺติ สํวิชฺชนฺติ อตฺถิ อุปลพฺภนฺติ. \textbf{โลเก}ติ อปายโลเก ฯเปฯ อายตนโลเก. \textbf{มุนโย}ติ มุนินามกา อาชีวกา นิคณฺฐา ชฏิลา ตาปสา (เทวา โลเก มุนโยติ สญฺชานนฺติ, น จ เต มุนโย)\csromansharedfootnote{d2e359541fead2a7}{( ) เอตฺถนฺตเร ปาฐา นตฺถิ สฺยาโปตฺถเก.}ติ สนฺติ โลเก มุนโย. \textbf{อิจฺจายสฺมา นนฺโท}ติ \textbf{อิจฺจา}ติ ปทสนฺธิ ฯเปฯ. \textbf{อายสฺมา}ติ ปิยวจนํ ฯเปฯ. \textbf{นนฺโท}ติ ตสฺส พฺราหฺมณสฺส นามํ ฯเปฯ อภิลาโปติ อิจฺจายสฺมา นนฺโท.}

\prose{\csromanfolio{113}\textbf{ชนา วทนฺติ ตยิทํ กถํสู}ติ \textbf{ชนา}ติ ขตฺติยา จ พฺราหฺมณา จ เวสฺสา จ สุทฺทา จ คหฏฺฐา จ ปพฺพชิตา จ เทวา จ มนุสฺสา จ. \textbf{วทนฺตี}ติ กเถนฺติ ภณนฺติ ทีปยนฺติ โวหรนฺติ. \textbf{ตยิทํ กถํสู}ติ สํสยปุจฺฉา วิมติปุจฺฉา เทฺวฬฺหกปุจฺฉา อเนกํสปุจฺฉา “เอวํ นุ โข, น นุ โข, กิํ นุ โข, กถํ นุ โข”ติ ชนา วทนฺติ ตยิทํ กถํสุ.}

\prose{\textbf{ญาณูปปนฺนํ มุนิ โน วทนฺตี}ติ อฏฺฐ สมาปตฺติ\-ญาเณน วา ปญฺจาภิญฺญา\-ญาเณน วา อุเปตํ สมุเปตํ อุปาคตํ สมุปาคตํ อุปปนฺนํ สมุปปนฺนํ สมนฺนาคตํ มุนิํ วทนฺติ กเถนฺติ ภณนฺติ ทีปยนฺติ โวหรนฺตีติ ญาณูปปนฺนํ มุนิ โน วทนฺติ.}

\prose{\textbf{อุทาหุ เว ชีวิเตนู}\-\textbf{ปปนฺนนฺ}ติ อุทาหุ อเนกวิวิธอติปรมทุกฺกร- การิกลูขชีวิตานุโยเคน อุเปตํ สมุเปตํ อุปาคตํ สมุปาคตํ อุปปนฺนํ สมุปปนฺนํ สมนฺนาคตํ มุนิํ วทนฺติ กเถนฺติ ภณนฺติ ทีปยนฺติ โวหรนฺตีติ อุทาหุ เว ชีวิเตนู\-ปปนฺนํ. เตนาห โส พฺราหฺมโณ\csromandash{}}

\prose{สนฺติ โลเก มุนโย, (อิจฺจายสฺมา นนฺโท,)}

\prose{ชนา วทนฺติ ตยิทํ กถํสุ.}

\setlength{\csromangathapagewidth}{0pt}
\setlength{\csromangathawakwidth}{0pt}
\csromangathameasuregroup{}{ญาณูปปนฺนํ มุนิ โน วทนฺติ, \\ อุทาหุ เว ชีวิเตนุปปนฺนนฺติ. \\ \textbf{น ทิฏฺฐิยา น สุติยา น ญาเณน,} \\ \textbf{มุนีธ นนฺท กุสลา วทนฺติ.} \\ \textbf{วิเสนิกตฺวา}\textsuperscript{1}\textbf{ อนีฆา นิราสา,} \\ \textbf{จรนฺติ เย เต มุนโยติ พฺรูมิ.}}
\csromangathameasurewak{ญาณูปปนฺนํ มุนิ โน วทนฺติ, \\ อุทาหุ เว ชีวิเตนุปปนฺนนฺติ. \\ \textbf{น ทิฏฺฐิยา น สุติยา น ญาเณน,} \\ \textbf{มุนีธ นนฺท กุสลา วทนฺติ.} \\ \textbf{วิเสนิกตฺวา}\textsuperscript{1}\textbf{ อนีฆา นิราสา,} \\ \textbf{จรนฺติ เย เต มุนโยติ พฺรูมิ.}}
\setlength{\csromangathaleftcol}{0pt}
\csromangathagroup{\csromangathastanza{\csromangathawak{ญาณูปปนฺนํ มุนิ โน วทนฺติ, \\ อุทาหุ เว ชีวิเตนุปปนฺนนฺติ.}}\csromangathastanza{\csromangathawak{\csromangathaitemline{47}{\textbf{น ทิฏฺฐิยา น สุติยา น ญาเณน,}} \\ \csromangathacontline{\textbf{มุนีธ นนฺท กุสลา วทนฺติ.}} \\ \csromangathacontline{\textbf{วิเสนิกตฺวา}\csromansharedfootnote{6072f35d9db9de99}{วิเสนิํกตฺวา (ก) ขุ 7. 135 ปิฏฺเฐ.}\textbf{ อนีฆา นิราสา,}} \\ \csromangathacontline{\textbf{จรนฺติ เย เต มุนโยติ พฺรูมิ.}}}}}

\prose{\textbf{น ทิฏฺฐิยา น สุติยา น ญาเณนา}ติ \textbf{น ทิฏฺฐิยา}ติ น ทิฏฺฐสุทฺธิยา. \textbf{น สุติยา}ติ น สุตสุทฺธิยา. \textbf{น ญาเณนา}ติ นปิ อฏฺฐ\-สมาปตฺติ\-ญาเณน นปิ ปญฺจาภิญฺญา\-ญาเณน นปิ มิจฺฉาญาเณนาติ น ทิฏฺฐิยา น สุติยา น ญาเณน.}

\prose{\textbf{มุนีธ นนฺท กุสลา วทนฺตี}ติ \textbf{กุสลา}ติ เย เต ขนฺธกุสลา ธาตุกุสลา อายตนกุสลา ปฏิจฺจสมุปฺปาท\-กุสลา สติปฏฺฐาน\-กุสลา สมฺมปฺปธาน\-กุสลา อิทฺธิปาท\-กุสลา อินฺทฺริยกุสลา พลกุสลา \csromanfolio{114}โพชฺฌงฺค\-กุสลา มคฺคกุสลา ผลกุสลา นิพฺพานกุสลา ทิฏฺฐสุทฺธิยา วา สุตสุทฺธิยา วา อฏฺฐ\-สมาปตฺติ\-ญาเณน วา ปญฺจาภิญฺญา\-ญาเณน วา มิจฺฉาญาเณน วา ทิฏฺเฐน วา สุเตน วา อุเปตํ สมุเปตํ อุปาคตํ \textbf{สมุปาคตํ อุปปนฺนํ สมุปปนฺนํ สมนฺนาคตํ มุนิํ น วทนฺติ} \textbf{น กเถนฺติ น ภณนฺติ น ทีปยนฺติ น โวหรนฺตีติ มุนีธ นนฺท กุสลา} วทนฺติ.}

\prose{\textbf{วิเสนิกตฺวา อนีฆา นิราสา, จรนฺติ เย เต มุนโยติ พฺรูมี}ติ \textbf{เสนา} วุจฺจติ มารเสนา, กายทุจฺจริตํ มารเสนา, วจีทุจฺจริตํ มารเสนา, มโนทุจฺจริตํ มารเสนา, ราโค มารเสนา, โทโส มารเสนา, โมโห มารเสนา, โกโธ. อุปนาโห. มกฺโข. ปฬาโส. อิสฺสา. มจฺฉริยํ. มายา. สาเฐยฺยํ. ถมฺโภ. สารมฺโภ. มาโน. อติมาโน. มโท. ปมาโท. สพฺเพ กิเลสา. สพฺเพ ทุจฺจริตา. สพฺเพ ทรถา. สพฺเพ ปริฬาหา. สพฺเพ สนฺตาปา. สพฺพา\-กุสลา\-ภิสงฺขารา มารเสนา. วุตฺตเญฺหตํ ภควตา\csromandash{}}

\setlength{\csromangathapagewidth}{0pt}
\setlength{\csromangathawakwidth}{0pt}
\csromangathameasuregroup{กามา เต ปฐมา เสนา, \\ ตติยา ขุปฺปิปาสา เต, \\ ปญฺจมํ ถินมิทฺธํ เต, \\ สตฺตมี วิจิกิจฺฉา เต, \\ เอสา นมุจิ เต เสนา\textsuperscript{1}, \\ น นํ อสูโร ชินาติ,}{\csromangathabat{กามา เต ปฐมา เสนา,}{ทุติยา อรติ วุจฺจติ.} \\ \csromangathabat{ตติยา ขุปฺปิปาสา เต,}{จตุตฺถี ตณฺหา ปวุจฺจติ.} \\ \csromangathabat{ปญฺจมํ ถินมิทฺธํ เต,}{ฉฏฺฐา ภีรู ปวุจฺจติ.} \\ \csromangathabat{สตฺตมี วิจิกิจฺฉา เต,}{มกฺโข ถมฺโภ เต อฏฺฐโม.} \\ ลาโภ สิโลโก สกฺกาโร, \\ มิจฺฉาลทฺโธ จ โย ยโส. \\ โย จตฺตานํ สมุกฺกํเส, \\ ปเร จ อวชานาติ. \\ \csromangathabat{เอสา นมุจิ เต เสนา\textsuperscript{1},}{กณฺหสฺสาภิปฺปหารินี.} \\ \csromangathabat{น นํ อสูโร ชินาติ,}{เชตฺวา จ ลภเต สุขนฺติ.}}
\csromangathameasurewak{ลาโภ สิโลโก สกฺกาโร, \\ มิจฺฉาลทฺโธ จ โย ยโส. \\ โย จตฺตานํ สมุกฺกํเส, \\ ปเร จ อวชานาติ.}
\csromangathasetleft{กามา เต ปฐมา เสนา, \\ ตติยา ขุปฺปิปาสา เต, \\ ปญฺจมํ ถินมิทฺธํ เต, \\ สตฺตมี วิจิกิจฺฉา เต, \\ เอสา นมุจิ เต เสนา\textsuperscript{1}, \\ น นํ อสูโร ชินาติ,}
\csromangathagroup{\csromangathastanza{\csromangathabat{กามา เต ปฐมา เสนา,}{ทุติยา อรติ วุจฺจติ.} \\ \csromangathabat{ตติยา ขุปฺปิปาสา เต,}{จตุตฺถี ตณฺหา ปวุจฺจติ.}}\csromangathastanza{\csromangathabat{ปญฺจมํ ถินมิทฺธํ เต,}{ฉฏฺฐา ภีรู ปวุจฺจติ.} \\ \csromangathabat{สตฺตมี วิจิกิจฺฉา เต,}{มกฺโข ถมฺโภ เต อฏฺฐโม.}}\csromangathastanza{\csromangathawak{ลาโภ สิโลโก สกฺกาโร, \\ มิจฺฉาลทฺโธ จ โย ยโส. \\ โย จตฺตานํ สมุกฺกํเส, \\ ปเร จ อวชานาติ.}}\csromangathastanza{\csromangathabat{เอสา นมุจิ เต เสนา\csromansharedfootnote{db8ceaac26287e3c}{เอสา เต นมุจิ เสนา (สฺยา, ก) ขุ 1. 343 ปิฏฺเฐ.},}{กณฺหสฺสา\-ภิปฺปหารินี.} \\ \csromangathabat{น นํ อสูโร ชินาติ,}{เชตฺวา จ ลภเต สุขนฺติ.}}}

\prose{ยโต จตูหิ อริยมคฺเคหิ สพฺพา จ มารเสนา สพฺเพ จ ปฏิเสนิกรา\csromansharedfootnote{6072f35d9db9de99}{วิเสนิํกตฺวา (ก) ขุ 7. 135 ปิฏฺเฐ.} กิเลสา ชิตา จ ปราชิตา จ ภคฺคา วิปฺปลุคฺคา\csromansharedfootnote{78861019ce557543}{วิปฺปลุคฺคตา (สฺยา), ปสฺส ขุ 7. 74 ปิฏฺเฐ.} ปรมฺมุขา, เตน วุจฺจนฺติ วิเสนิกตฺวา. \textbf{อนีฆา}ติ ราโค นีโฆ. โทโส นีโฆ. โมโห นีโฆ. โกโธ นีโฆ. อุปนาโห นีโฆ ฯเปฯ สพฺพา\-กุสลา\-ภิสงฺขารา นีฆา, เยสํ เอเต นีฆา ปหีนา สมุจฺฉินฺนา วูปสนฺตา ปฏิปฺปสฺสทฺธา อภพฺพุ\-ปฺปตฺติกา ญาณคฺคินา ทฑฺฒา. เต วุจฺจนฺติ อนีฆา. \csromanfolio{115}\textbf{นิราสา}ติ อาสา วุจฺจติ ตณฺหา, โย ราโค สาราโค ฯเปฯ อวิชฺชา โลโภ อกุสลมูลํ. เยสํ เอสา อาสา ตณฺหา ปหีนา สมุจฺฉินฺนา วูปสนฺตา ปฏิปฺปสฺสทฺธา อภพฺพุ\-ปฺปตฺติกา ญาณคฺคินา ทฑฺฒา. เต วุจฺจนฺติ นิราสา อรหนฺโต ขีณาสวา. \textbf{วิเสนิกตฺวา อนีฆา นิราสา, จรนฺติ เย เต มุนโยติ พฺรูมี}ติ เย เต วิเสนิกตฺวาว อนีฆา จ นิราสา จ จรนฺติ วิหรนฺติ อิริยนฺติ วตฺเตนฺติ ปาเลนฺติ ยเปนฺติ ยาเปนฺติ. เต โลเก มุนโยติ พฺรูมิ อาจิกฺขามิ เทเสมิ ปญฺญเปมิ ปฏฺฐเปมิ วิวรามิ วิภชามิ อุตฺตานีกโรมิ ปกาเสมีติ วิเสนิกตฺวา อนีฆา นิราสา, จรนฺติ เย เต มุนโยติ พฺรูมิ. เตนาห ภควา\csromandash{}}

\setlength{\csromangathapagewidth}{0pt}
\setlength{\csromangathawakwidth}{0pt}
\csromangathameasuregroup{}{น ทิฏฺฐิยา น สุติยา น ญาเณน, \\ มุนีธ นนฺท กุสลา วทนฺติ.}
\csromangathameasurewak{น ทิฏฺฐิยา น สุติยา น ญาเณน, \\ มุนีธ นนฺท กุสลา วทนฺติ.}
\setlength{\csromangathaleftcol}{0pt}
\csromangathagroup{\csromangathastanza{\csromangathawak{น ทิฏฺฐิยา น สุติยา น ญาเณน, \\ มุนีธ นนฺท กุสลา วทนฺติ.}}}

\prosecont{\textbf{วิเสนิกตฺวา อนีฆา นิราสา},}

\proseitem{7}{จรนฺติ เย เต มุนโยติ พฺรูมีติ.}

\csromanhangingbegin
\hangingheaditem{48}{\textbf{เย เกจิเม สมณ}\-\textbf{พฺราหฺมณาเส, (อิจฺจายสฺมา นนฺโท,)}}
\hangingbody{\textbf{ทิฏฺถสฺสุเตนาปิ วทนฺติ สุทฺธิํ.}}
\hangingbody{\textbf{สีลพฺพเตนาปิ วทนฺติ สุทฺธิํ,}}
\hangingbody{\textbf{อเนกรูเปน วทนฺติ สุทฺธิํ.}}
\hangingbody{\textbf{กจฺจิสฺสุ เต ภควา ตตฺถ ยตา จรนฺตา,}}
\hangingbody{\textbf{อตารุ ชาติญฺจ ชรญฺจ มาริส.}}
\hangingbody{\textbf{ปุจฺฉามิ ตํ ภควา พฺรูหิ เมตํ.}}
\csromanhangingend

\prose{\textbf{เย เกจิเม สมณ}\-\textbf{พฺราหฺมณา}\-\textbf{เส}ติ \textbf{เย เกจี}ติ สพฺเพน สพฺพํ สพฺพถา สพฺพํ อเสสํ นิสฺเสสํ ปริยาทิย\-น\-วจนเมตํ เย เกจีติ. \textbf{สมณา}ติ เย เกจิ อิโต พหิทฺธา ปพฺพชฺชูปคตา ปริพฺพาชก\-สมาปนฺนา. \textbf{พฺราหฺมณา}ติ เย เกจิ โภวาทิกาติ เย เกจิเม สมณ\-พฺราหฺมณาเส. \textbf{อิจฺจายสฺมา นนฺโท}ติ \textbf{อิจฺจา}ติ ปทสนฺธิ ฯเปฯ. \textbf{อายสฺมา}ติ ปิยวจนํ ฯเปฯ. \textbf{นนฺโท}ติ ตสฺส พฺราหฺมณสฺส นามํ ฯเปฯ อภิลาโปติ อิจฺจายสฺมา นนฺโท.}

\prose{\textbf{ทิฏฺฐสฺสุเตนาปิ วทนฺติ สุทฺธินฺ}ติ ทิฏฺเฐนปิ สุทฺธิํ วิสุทฺธิํ ปริสุทฺธิํ, มุตฺติํ วิมุตฺติํ ปริมุตฺติํ วทนฺติ กเถนฺติ ภณนฺติ ทีปยนฺติ โวหรนฺติ. สุเตนปิ สุทฺธิํ วิสุทฺธิํ ปริสุทฺธิํ, มุตฺติํ วิมุตฺติํ ปริมุตฺติํ วทนฺติ กเถนฺติ ภณนฺติ \csromanfolio{116}ทีปยนฺติ โวหรนฺติ. ทิฏฺฐสฺสุเตนปิ สุทฺธิํ วิสุทฺธิํ ปริสุทฺธิํ, มุตฺติํ วิมุตฺติํ ปริมุตฺติํ วทนฺติ กเถนฺติ ภณนฺติ ทีปยนฺติ โวหรนฺตีติ ทิฏฺฐสฺสุเตนาปิ วทนฺติ สุทฺธิํ.}

\prose{\textbf{สีลพฺพเตนาปิ วทนฺติ สุทฺธินฺ}ติ สีเลนปิ สุทฺธิํ วิสุทฺธิํ ปริสุทฺธิํ, มุตฺติํ วิมุตฺติํ ปริมุตฺติํ วทนฺติ กเถนฺติ ภณนฺติ ทีปยนฺติ โวหรนฺติ. วเตนปิ สุทฺธิํ วิสุทฺธิํ ปริสุทฺธิํ, มุตฺติํ วิมุตฺติํ ปริมุตฺติํ วทนฺติ กเถนฺติ ภณนฺติ ทีปยนฺติ โวหรนฺติ. สีลพฺพเตนาปิ สุทฺธิํ วิสุทฺธิํ ปริสุทฺธิํ, มุตฺติํ วิมุตฺติํ ปริมุตฺติํ วทนฺติ กเถนฺติ ภณนฺติ ทีปยนฺติ โวหรนฺตีติ สีลพฺพเตนาปิ วทนฺติ สุทฺธิํ.}

\prose{\textbf{อเนกรูเปน วทนฺติ สุทฺธินฺ}ติ อเนกวิธ\-โกตูหล\-มงฺคเลน\csromansharedfootnote{d7b1922281d80040}{อเนกวิธวตฺต กุตูหลมงฺคเลน (สฺยา)} สุทฺธิํ วิสุทฺธิํ ปริสุทฺธิํ, มุตฺติํ วิมุตฺติํ ปริมุตฺติํ วทนฺติ กเถนฺติ ภณนฺติ ทีปยนฺติ โวหรนฺตีติ อเนกรูเปน วทนฺติ สุทฺธิํ.}

\prose{\textbf{กจฺจิสฺสุ เต ภควา ตตฺถ ยตา จรนฺตา}ติ \textbf{กจฺจิสฺสู}ติ สํสยปุจฺฉา วิมติปุจฺฉา เทฺวฬฺหกปุจฺฉา อเนกํสปุจฺฉา “เอวํ นุ โข, น นุ โข, กิํ นุ โข, กถํ นุ โข”ติ กจฺจิสฺสุ. \textbf{เต}ติ ทิฏฺฐิคติกา. \textbf{ภควา}ติ คารวาธิวจนเมตํ ฯเปฯ สจฺฉิกา ปญฺญตฺติ, ยทิทํ ภควาติ กจฺจิสฺสุ เต ภควา. \textbf{ตตฺถ ยตา จรนฺตา}ติ \textbf{ตตฺถา}ติ สกาย ทิฏฺฐิยา สกาย ขนฺติยา สกาย รุจิยา สกาย ลทฺธิยา. \textbf{ยตา}ติ ยตฺตา ปฏิยตฺตา\csromansharedfootnote{1652958e93dfa58e}{ยตา ปฏิยตา (สฺยา)} คุตฺตา โคปิตา รกฺขิตา สํวุตา. \textbf{จรนฺตา}ติ จรนฺตา วิหรนฺตา อิริยนฺตา วตฺเตนฺตา ปาเลนฺตา ยเปนฺตา ยาเปนฺตาติ กจฺจิสฺสุ เต ภควา ตตฺถ ยตา จรนฺตา.}

\prose{\textbf{อตารุ ชาติญฺจ ชรญฺจ มาริสา}ติ ชาติ\-ชรา\-มรณํ อตริํสุ อุตฺตริํสุ ปตริํสุ สมติกฺกมิํสุ วีติวตฺติํสุ. \textbf{มาริสา}ติ ปิยวจนํ ครุวจนํ สคารวสปฺปติสฺสาธิวจนเมตํ มาริสาติ อตารุ ชาติญฺจ ชรญฺจ มาริส.}

\prose{\textbf{ปุจฺฉามิ ตํ ภควา พฺรูหิ เมตนฺ}ติ \textbf{ปุจฺฉามิ ตนฺ}ติ ปุจฺฉามิ ตํ, ยาจามิ ตํ, อชฺเฌสามิ ตํ, “กถยสฺสุ เม”ติ ปุจฺฉามิ ตํ. \textbf{ภควา}ติ ฯเปฯ สจฺฉิกา ปญฺญตฺติ, ยทิทํ ภควาติ. \textbf{พฺรูหิ เมตนฺ}ติ พฺรูหิ อาจิกฺขาหิ เทเสหิ \csromanfolio{117}ปญฺญเปหิ ปฏฺฐเปหิ วิวราหิ วิภชาหิ อุตฺตานีกโรหิ ปกาเสหีติ ปุจฺฉามิ ตํ ภควา พฺรูหิ เมตํ. เตนาห โส พฺราหฺมโณ\csromandash{}}

\prose{เย เกจิเม สมณ\-พฺราหฺมณาเส, (อิจฺจายสฺมา นนฺโท,)}

\prose{ทิฏฺฐสฺสุเตนาปิ วทนฺติ สุทฺธิํ.}

\setlength{\csromangathapagewidth}{0pt}
\setlength{\csromangathawakwidth}{0pt}
\csromangathameasuregroup{}{สีลพฺพเตนาปิ วทนฺติ สุทฺธิํ, \\ อเนกรูเปน วทนฺติ สุทฺธิํ. \\ กจฺจิสฺสุ เต ภควา ตตฺถ ยตา จรนฺตา, \\ อตารุ ชาติญฺจ ชรญฺจ มาริส. \\ ปุจฺจามิ ตํ ภควา พฺรูหิ เมตนฺติ.}
\csromangathameasurewak{สีลพฺพเตนาปิ วทนฺติ สุทฺธิํ, \\ อเนกรูเปน วทนฺติ สุทฺธิํ. \\ กจฺจิสฺสุ เต ภควา ตตฺถ ยตา จรนฺตา, \\ อตารุ ชาติญฺจ ชรญฺจ มาริส.}
\setlength{\csromangathaleftcol}{0pt}
\csromangathagroup{\csromangathastanza{\csromangathawak{สีลพฺพเตนาปิ วทนฺติ สุทฺธิํ, \\ อเนกรูเปน วทนฺติ สุทฺธิํ. \\ กจฺจิสฺสุ เต ภควา ตตฺถ ยตา จรนฺตา, \\ อตารุ ชาติญฺจ ชรญฺจ มาริส.}}\csromangathastanza{ปุจฺจามิ ตํ ภควา พฺรูหิ เมตนฺติ.}}

\csromanhangingbegin
\hangingheaditem{49}{\textbf{เย เกจิเม สมณ}\-\textbf{พฺราหฺมณาเส, (นนฺทาติ ภควา,)}}
\hangingbody{ทิฏฺฐสฺสุเตนาปิ วทนฺติ สุทฺธิํ.}
\hangingbody{สีลพฺพเตนาปิ วทนฺติ สุทฺธิํ,}
\hangingbody{อเนกรูเปน วทนฺติ สุทฺธิํ.}
\hangingbody{\textbf{กิญฺจาปิ เต ตตฺถ ยตา จรนฺติ,}}
\hangingbody{\textbf{นาตริํสุ ชาติชรนฺติ พฺรูมิ.}}
\csromanhangingend

\prose{\textbf{เย เกจิเม สมณ}\-\textbf{พฺราหฺมณา}\-\textbf{เส}ติ \textbf{เย เกจี}ติ สพฺเพน สพฺพํ สพฺพถา สพฺพํ อเสสํ นิสฺเสสํ ปริยาทิย\-น\-วจนเมตํ เย เกจีติ. \textbf{สมณา}ติ เย เกจิ อิโต พหิทฺธา ปพฺพชฺชูปคตา ปริพฺพาชก\-สมาปนฺนา. \textbf{พฺราหฺมณา}ติ เย เกจิ โภวาทิกาติ เย เกจิเม สมณ\-พฺราหฺมณาเส. \textbf{นนฺทาติ ภควา}ติ \textbf{นนฺทา}ติ ภควา ตํ พฺราหฺมณํ นาเมน อาลปติ. \textbf{ภควา}ติ คารวาธิวจนเมตํ ฯเปฯ สจฺฉิกา ปญฺญตฺติ, ยทิทํ ภควาติ นนฺทาติ ภควา.}

\prose{\textbf{ทิฏฺฐสฺสุเตนาปิ วทนฺติ สุทฺธินฺ}ติ ทิฏฺเฐนปิ สุทฺธิํ วิสุทฺธิํ ปริสุทฺธิํ, มุตฺติํ วิมุตฺติํ ปริมุตฺติํ วทนฺติ กเถนฺติ ภณนฺติ ทีปยนฺติ โวหรนฺติ. สุเตนปิ สุทฺธิํ ฯเปฯ ทิฏฺฐสฺสุเตนปิ สุทฺธิํ วิสุทฺธิํ ปริสุทฺธิํ, มุตฺติํ วิมุตฺติํ ปริมุตฺติํ วทนฺติ กเถนฺติ ภณนฺติ ทีปยนฺติ โวหรนฺตีติ ทิฏฺฐสฺสุเตนาปิ วทนฺติ สุทฺธิํ.}

\prose{\textbf{สีลพฺพเตนาปิ วทนฺติ สุทฺธินฺ}ติ สีเลนปิ สุทฺธิํ วิสุทฺธิํ ปริสุทฺธิํ, มุตฺติํ วิมุตฺติํ ปริมุตฺติํ วทนฺติ กเถนฺติ ภณนฺติ ทีปยนฺติ โวหรนฺติ. วเตนปิ สุทฺธิํ ฯเปฯ โวหรนฺติ. สีลพฺพเตนาปิ สุทฺธิํ วิสุทฺธิํ ปริสุทฺธิํ, มุตฺติํ วิมุตฺติํ \csromanfolio{118}\textbf{ปริมุตฺติํ วทนฺติ กเถนฺติ ภณนฺติ ทีปยนฺติ โวหรนฺตีติ สีลพฺพเตนาปิ} \textbf{วทนฺติ สุทฺธิํ.}}

\prose{\textbf{อเนกรูเปน วทนฺติ สุทฺธินฺ}ติ อเนกวิธ\-โกตูหล\-มงฺคเลน สุทฺธิํ วิสุทฺธิํ ปริสุทฺธิํ, มุตฺติํ วิมุตฺติํ ปริมุตฺติํ วทนฺติ กเถนฺติ ภณนฺติ ทีปยนฺติ โวหรนฺตีติ อเนกรูเปน วทนฺติ สุทฺธิํ.}

\prose{\textbf{กิญฺจาปิ เต ตตฺถ ยตา จรนฺตี}ติ \textbf{กิญฺจาปี}ติ ปทสนฺธิ ปทสํสคฺโค ปทปาริปูรี อกฺขร\-สมวาโย พฺยญฺชน\-สิลิฏฺฐตา ปทา\-นุปุพฺพตา\-เปตํ กิญฺจาปีติ. \textbf{เต}ติ ทิฏฺฐิคติกา. \textbf{ตตฺถา}ติ สกาย ทิฏฺฐิยา สกาย ขนฺติยา สกาย รุจิยา สกาย ลทฺธิยา. \textbf{ยตา}ติ ยตฺตา ปฏิยตฺตา คุตฺตา โคปิตา รกฺขิตา สํวุตา. \textbf{จรนฺตี}ติ จรนฺติ วิหรนฺติ อิริยนฺติ วตฺเตนฺติ ปาเลนฺติ ยเปนฺติ ยาเปนฺตีติ กิญฺจาปิ เต ตตฺถ ยตา จรนฺติ.}

\prose{\textbf{นาตริํสุ ชาติชรนฺติ พฺรูมี}ติ ชาติ\-ชรา\-มรณํ น ตริํสุ น อุตฺตริํสุ น ปตริํสุ น สมติกฺกมิํสุ น วีติวตฺติํสุ, ชาติชรามรณา อนิกฺขนฺตา อนิสฺสฏา อนติกฺกนฺตา อสมติกฺกนฺตา อวีติวตฺตา, อนฺโต\-ชาติชรามรเณ ปริวตฺเตนฺติ, อนฺโต\-สํสาร\-ปเถ ปริวตฺเตนฺติ, ชาติยา อนุคตา, ชราย อนุสฏา, พฺยาธินา อภิภูตา, มรเณน อพฺภาหตา อตาณา อเลณา อสรณา อสรณีภูตาติ พฺรูมิ อาจิกฺขามิ เทเสมิ ปญฺญเปมิ ปฏฺฐเปมิ วิวรามิ วิภชามิ อุตฺตานีกโรมิ ปกาเสมีติ นาตริํสุ ชาติชรนฺติ พฺรูมิ. เตนาห ภควา\csromandash{}}

\prose{เย เกจิเม สมณ\-พฺราหฺมณาเส, (นนฺทาติ ภควา,)}

\prose{ทิฏฺฐสฺสุเตนาปิ วทนฺติ สุทฺธิํ.}

\setlength{\csromangathapagewidth}{0pt}
\setlength{\csromangathawakwidth}{0pt}
\csromangathameasuregroup{}{สีลพฺพเตนาปิ วทนฺติ สุทฺธิํ, \\ อเนกรูเปน วทนฺติ สุทฺธิํ.}
\csromangathameasurewak{สีลพฺพเตนาปิ วทนฺติ สุทฺธิํ, \\ อเนกรูเปน วทนฺติ สุทฺธิํ.}
\setlength{\csromangathaleftcol}{0pt}
\csromangathagroup{\csromangathastanza{\csromangathawak{สีลพฺพเตนาปิ วทนฺติ สุทฺธิํ, \\ อเนกรูเปน วทนฺติ สุทฺธิํ.}}}

\prose{กิญฺจาปิ เต ตตฺถ ยตา จรนฺติ,}

\prose{นาตริํสุ ชาติชรนฺติ พฺรูมีติ.}

\csromanhangingbegin
\hangingheaditem{50}{\textbf{เย เกจิเม สมณ}\-\textbf{พฺราหฺมณาเส, (อิจฺจายสฺมา นนฺโท,)}}
\hangingbody{ทิฏฺฐสฺสุเตนาปิ วทนฺติ สุทฺธิํ.}
\hangingbody{สีลพฺพเตนาปิ วทนฺติ สุทฺธิํ,}
\hangingbody{อเนกรูเปน วทนฺติ สุทฺธิํ.}
\csromanhangingend

\setlength{\csromangathapagewidth}{0pt}
\setlength{\csromangathawakwidth}{0pt}
\csromangathameasuregroup{}{\textbf{เต เจ มุนี พฺรูสิ อโนฆติณฺเณ,} \\ \textbf{อถ โก จรหิ เทว}\textbf{มนุสฺส}\textbf{โลเก.} \\ \textbf{อตาริ ชาติญฺจ ชรญฺจ มาริส,} \\ \textbf{ปุจฺฉามิ ตํ ภควา พฺรูหิ เมตํ.}}
\csromangathameasurewak{\textbf{เต เจ มุนี พฺรูสิ อโนฆติณฺเณ,} \\ \textbf{อถ โก จรหิ เทว}\textbf{มนุสฺส}\textbf{โลเก.} \\ \textbf{อตาริ ชาติญฺจ ชรญฺจ มาริส,} \\ \textbf{ปุจฺฉามิ ตํ ภควา พฺรูหิ เมตํ.}}
\setlength{\csromangathaleftcol}{0pt}
\csromangathagroup{\csromangathastanza{\csromangathawak{\csromanfolio{119}\textbf{เต เจ มุนี พฺรูสิ อโนฆติณฺเณ,} \\ \textbf{อถ โก จรหิ เทว}\-\textbf{มนุสฺส}\-\textbf{โลเก.} \\ \textbf{อตาริ ชาติญฺจ ชรญฺจ มาริส,} \\ \textbf{ปุจฺฉามิ ตํ ภควา พฺรูหิ เมตํ.}}}}

\prose{\textbf{เย เกจิเม สมณ}\-\textbf{พฺราหฺมณา}\-\textbf{เส}ติ \textbf{เย เกจี}ติ สพฺเพน สพฺพํ สพฺพถา สพฺพํ อเสสํ นิสฺเสสํ ปริยาทิย\-น\-วจนเมตํ เย เกจีติ. \textbf{สมณา}ติ เย เกจิ อิโต พหิทฺธา ปพฺพชฺชูปคตา ปริพฺพาชก\-สมาปนฺนา. \textbf{พฺราหฺมณา}ติ เย เกจิ โภวาทิกาติ เย เกจิเม สมณ\-พฺราหฺมณาเส. \textbf{อิจฺจายสฺมา นนฺโท}ติ \textbf{อิจฺจา}ติ ปทสนฺธิ ฯเปฯ อิจฺจายสฺมา นนฺโท.}

\prose{\textbf{ทิฏฺฐสฺสุเตนาปิ วทนฺติ สุทฺธินฺ}ติ ทิฏฺเฐนปิ สุทฺธิํ วิสุทฺธิํ ปริสุทฺธิํ, มุตฺติํ วิมุตฺติํ ปริมุตฺติํ วทนฺติ กเถนฺติ ภณนฺติ ทีปยนฺติ โวหรนฺติ. สุเตนาปิ สุทฺธิํ ฯเปฯ ทิฏฺฐสฺสุเตนาปิ สุทฺธิํ วิสุทฺธิํ ปริสุทฺธิํ, มุตฺติํ วิมุตฺติํ ปริมุตฺติํ วทนฺติ กเถนฺติ ภณนฺติ ทีปยนฺติ โวหรนฺตีติ ทิฏฺฐสฺสุเตนาปิ วทนฺติ สุทฺธิํ.}

\prose{\textbf{สีลพฺพเตนาปิ วทนฺติ สุทฺธินฺ}ติ สีเลนาปิ สุทฺธิํ วิสุทฺธิํ ปริสุทฺธิํ, มุตฺติํ วิมุตฺติํ ปริมุตฺติํ วทนฺติ กเถนฺติ ภณนฺติ ทีปยนฺติ โวหรนฺติ. วเตนาปิ สุทฺธิํ ฯเปฯ โวหรนฺติ. สีลพฺพเตนาปิ สุทฺธิํ วิสุทฺธิํ ปริสุทฺธิํ, มุตฺติํ วิมุตฺติํ ปริมุตฺติํ วทนฺติ กเถนฺติ ภณนฺติ ทีปยนฺติ โวหรนฺตีติ สีลพฺพเตนาปิ วทนฺติ สุทฺธิํ.}

\prose{\textbf{อเนกรูเปน วทนฺติ สุทฺธินฺ}ติ อเนกวิธ\-โกตูหล\-มงฺคเลน สุทฺธิํ วิสุทฺธิํ ปริสุทฺธิํ, มุตฺติํ วิมุตฺติํ ปริมุตฺติํ วทนฺติ กเถนฺติ ภณนฺติ ทีปยนฺติ โวหรนฺตีติ อเนกรูเปน วทนฺติ สุทฺธิํ.}

\prose{\textbf{เต เจ มุนี พฺรูสิ อโนฆติณฺเณ}ติ \textbf{เต เจ}ติ ทิฏฺฐิคติเก. \textbf{มุนี}ติ โมนํ วุจฺจติ ญาณํ ฯเปฯ สงฺค\-ชาลม\-ติจฺจ โส มุนิ. \textbf{พฺรูสิ อโนฆติณฺเณ}ติ กาโมฆํ ภโวฆํ ทิฏฺโฐฆํ อวิชฺโชฆํ อติณฺเณ อนติกฺกนฺเต อสมติกฺกนฺเต อวีติวตฺเต, อนฺโต\-ชาติชรามรเณ ปริวตฺเตนฺเต, อนฺโต\-สํสาร\-ปเถ ปริวตฺเตนฺเต, ชาติยา อนุคเต, ชราย อนุสเฏ, พฺยาธินา อภิภูเต, มรเณน อพฺภาหเต อตาเณ อเลเณ อสรเณ อสรณีภูเต. \textbf{พฺรูสี}ติ พฺรูสิ อาจิกฺขสิ เทเสสิ ปญฺญเปสิ ปฏฺฐเปสิ}

\prose{\csromanfolio{120}วิวรสิ วิภชสิ อุตฺตานีกโรสิ ปกาเสสีติ เต เจ มุนี พฺรูสิ อโนฆติณฺเณ.}

\prose{\textbf{อถ โก จรหิ เทว}\-\textbf{มนุสฺส}\-\textbf{โลเก, อตาริ ชาติญฺจ ชรญฺจ มาริสา}ติ อถ โก เอโส สเทวเก โลเก สมารเก สพฺรหฺมเก สสฺสมณ\-พฺราหฺมณิยา ปชาย สเทว\-มนุสฺสาย ชาติ\-ชรา\-มรณํ อตริ อุตฺตริ ปตริ สมติกฺกมิ วีติวตฺตยิ. \textbf{มาริสา}ติ ปิยวจนํ ครุวจนํ สคารวสปฺปติสฺสาธิวจนเมตํ มาริสาติ อถ โก จรหิ เทว\-มนุสฺส\-โลเก, อตาริ ชาติญฺจ ชรญฺจ มาริส.}

\prose{\textbf{ปุจฺฉามิ ตํ ภควา พฺรูหิ เมตนฺ}ติ \textbf{ปุจฺฉามิ ตนฺ}ติ ปุจฺฉามิ ตํ, ยาจามิ ตํ, อชฺเฌสามิ ตํ, ปสาเทมิ ตํ. \textbf{ภควา}ติ คารวาธิวจนเมตํ ฯเปฯ สจฺฉิกา ปญฺญตฺติ, ยทิทํ ภควาติ. \textbf{พฺรูหิ เมตนฺ}ติ พฺรูหิ อาจิกฺขาหิ เทเสหิ ปญฺญเปหิ ปฏฺฐเปหิ วิวราหิ วิภชาหิ อุตฺตานีกโรหิ ปกาเสหีติ ปุจฺฉามิ ตํ ภควา พฺรูหิ เมตํ. เตนาห โส พฺราหฺมโณ\csromandash{}}

\prose{เย เกจิเม สมณ\-พฺราหฺมณาเส, (อิจฺจายสฺมา นนฺโท,)}

\prose{ทิฏฺฐสฺสุเตนาปิ วทนฺติ สุทฺธิํ.}

\setlength{\csromangathapagewidth}{0pt}
\setlength{\csromangathawakwidth}{0pt}
\csromangathameasuregroup{}{สีลพฺพเตนาปิ วทนฺติ สุทฺธิํ, \\ อเนกรูเปน วทนฺติ สุทฺธิํ. \\ เต เจ มุนี พฺรูสิ อโนฆติณฺเณ,}
\csromangathameasurewak{สีลพฺพเตนาปิ วทนฺติ สุทฺธิํ, \\ อเนกรูเปน วทนฺติ สุทฺธิํ. \\ เต เจ มุนี พฺรูสิ อโนฆติณฺเณ,}
\setlength{\csromangathaleftcol}{0pt}
\csromangathagroup{\csromangathastanza{\csromangathawak{สีลพฺพเตนาปิ วทนฺติ สุทฺธิํ, \\ อเนกรูเปน วทนฺติ สุทฺธิํ. \\ เต เจ มุนี พฺรูสิ อโนฆติณฺเณ,}}}

\prose{อถ โก จรหิ เทว\-มนุสฺส\-โลเก.}

\prose{อตาริ ชาติญฺจ ชรญฺจ มาริส,}

\prose{ปุจฺฉามิ ตํ ภควา พฺรูหิ เมตนฺติ.}

\csromanhangingbegin
\hangingheaditem{51}{\textbf{นาหํ สพฺเพ สมณ}\-\textbf{พฺราหฺมณาเส, (นนฺทาติ ภควา,)}}
\hangingbody{\textbf{ชาติชราย นิวุตาติ พฺรูมิ.}}
\hangingbody{\textbf{เย สีธ ทิฏฺฐํ ว สุตํ มุตํ วา,}}
\hangingbody{\textbf{สีลพฺพตํ วาปิ ปหาย สพฺพํ.}}
\hangingbody{\textbf{อเนกรูปมฺปิ ปหาย สพฺพํ,}}
\hangingbody{\textbf{ตณฺหํ ปริญฺญาย อนา}สวาเส1.}
\hangingbody{\textbf{เต เว นรา โอฆติณฺณาติ พฺรูมิ.}}
\csromanhangingend

\prose{\csromanfolio{121}นาหํ สพฺเพ สมณ\-พฺราหฺมณาเส, นนฺทาติ ภควา. \textbf{ ชาติชราย นิวุตาติ พฺรูมี}ติ นาหํ นนฺท สพฺเพ สมณพฺราหฺมณา ชาติชราย อาวุตา นิวุตา โอวุตา ปิหิตา ปฏิจฺฉนฺนา ปฏิกุชฺชิตาติ วทามิ. อตฺถิ เต สมณพฺราหฺมณา. เยสํ ชาติ จ ชรามรณญฺจ ปหีนา อุจฺฉินฺนมูลา ตาลาวตฺถุกตา อนภาวํกตา อายติํ อนุปฺปาทธมฺมาติ พฺรูมิ อาจิกฺขามิ เทเสมิ ปญฺญเปมิ ปฏฺฐเปมิ วิวรามิ วิภชามิ อุตฺตานีกโรมิ ปกาเสมีติ นาหํ สพฺเพ สมณ\-พฺราหฺมณาเส. นนฺทาติ ภควา ชาติชราย นิวุตาติ พฺรูมิ.}

\prose{\textbf{เย สีธ ทิฏฺฐํ ว สุตํ มุตํ วา, สีลพฺพตํ วาปิ ปหาย สพฺพนฺ}ติ เย สพฺพา ทิฏฺฐสุทฺธิโย ปหาย ชหิตฺวา ปชหิตฺวา วิโนเทตฺวา พฺยนฺตีกริตฺวา อนภาวํ คเมตฺวา. เย สพฺพา สุตสุทฺธิโย ปหาย ฯเปฯ. เย สพฺพา มุตสุทฺธิโย ปหาย. เย สพฺพา ทิฏฺฐสุต\-มุต\-สุทฺธิโย ปหาย. เย สพฺพา สีลสุทฺธิโย ปหาย. เย สพฺพา วตสุทฺธิโย ปหาย. เย สพฺพา สีลพฺพต\-สุทฺธิโย ปหาย ชหิตฺวา ปชหิตฺวา วิโนเทตฺวา พฺยนฺตีกริตฺวา อนภาวํ คเมตฺวาติ เย สีธ ทิฏฺฐํ ว สุตํ มุตํ วา, สีลพฺพตํ วาปิ ปหาย สพฺพํ.}

\prose{\textbf{อเนกรูปมฺปิ ปหาย สพฺพนฺ}ติ อเนกวิธ\-โกตูหล\-มงฺคเลน สุทฺธิํ วิสุทฺธิํ ปริสุทฺธิํ, มุตฺติํ วิมุตฺติํ ปริมุตฺติํ ปหาย ชหิตฺวา ปชหิตฺวา วิโนเทตฺวา พฺยนฺตีกริตฺวา อนภาวํ คเมตฺวาติ อเนกรูปมฺปิ ปหาย สพฺพํ.}

\prose{\textbf{ตณฺหํ ปริญฺญาย อนาสวา เส, เต เว นรา โอฆติณฺณาติ พฺรูมี}ติ \textbf{ตณฺหา}ติ รูปตณฺหา สทฺทตณฺหา คนฺธตณฺหา รสตณฺหา โผฏฺฐพฺพ\-ตณฺหา ธมฺมตณฺหา. \textbf{ตณฺหํ ปริญฺญายา}ติ ตณฺหํ ตีหิ ปริญฺญาหิ ปริชานิตฺวา ญาตปริญฺญาย ตีรณ\-ปริญฺญาย ปหาน\-ปริญฺญาย. กตมา ญาตปริญฺญา, ตณฺหํ ชานาติ\csromansharedfootnote{c4ae649136875770}{ปชานาติ (สฺยา) ปริชานาติ (ก) ขุ 7. 40 ปิฏฺเฐปิ.} “อยํ รูปตณฺหา, อยํ สทฺทตณฺหา, อยํ คนฺธตณฺหา, อยํ รสตณฺหา, อยํ โผฏฺฐพฺพ\-ตณฺหา, อยํ ธมฺมตณฺหา”ติ ชานาติ ปสฺสติ. อยํ ญาตปริญฺญา.}

\prose{กตมา ตีรณปริญฺญา, เอวํ ญาตํ กตฺวา ตณฺหํ ตีเรติ อนิจฺจโต. ทุกฺขโต. โรคโต. คณฺฑโต ฯเปฯ นิสฺสรณโต ตีเรติ. อยํ ตีรณปริญฺญา.}

\prose{\csromanfolio{122}กตมา ปหานปริญฺญา, เอวํ ตีรยิตฺวา ตณฺหํ ปชหติ วิโนเทติ พฺยนฺตีกโรติ อนภาวํ คเมติ. วุตฺตเญฺหตํ ภควตา\csromandash{}โย ภิกฺขเว ตณฺหาย ฉนฺทราโค, ตํ ปชหถ. เอวํ สา ตณฺหา ปหีนา ภวิสฺสติ อุจฺฉินฺนมูลา ตาลาวตฺถุกตา อนภาวํกตา อายติํ อนุปฺปาทธมฺมา. อยํ ปหานปริญฺญา.  ตณฺหํ ปริญฺญายาติ ตณฺหํ อิมาหิ ตีหิ ปริญฺญาหิ ปริชานิตฺวา. อนาสวาติ จตฺตาโร อาสวา กามาสโว ภวาสโว ทิฏฺฐาสโว อวิชฺชาสโว. เยสํ อิเม อาสวา ปหีนา อุจฺฉินฺนมูลา ตาลาวตฺถุกตา อนภาวํกตา อายติํ อนุปฺปาทธมฺมา. เต วุจฺจนฺติ อนาสวา อรหนฺโต ขีณาสวา. ตณฺหํ ปริญฺญาย อนาสวา.}

\prose{\textbf{เต เว นรา โอฆติณฺณาติ พฺรูมี}ติ เย ตณฺหํ ปริญฺญาย อนาสวา. เต กาโมฆํ ติณฺณา, ภโวฆํ ติณฺณา, ทิฏฺโฐฆํ ติณฺณา, อวิชฺโชฆํ ติณฺณา, สพฺพ\-สํสารปถํ ติณฺณา อุตฺติณฺณา นิตฺถิณฺณา อติกฺกนฺตา สมติกฺกนฺตา วีติวตฺตาติ พฺรูมิ อาจิกฺขามิ เทเสมิ ปญฺญเปมิ ปฏฺฐเปมิ วิวรามิ วิภชามิ อุตฺตานีกโรมิ ปกาเสมีติ ตณฺหํ ปริญฺญาย อนาสวาเส, เต เว นรา โอฆติณฺณาติ พฺรูมิ. เตนาห ภควา\csromandash{}}

\prose{นาหํ สพฺเพ สมณ\-พฺราหฺมณาเส, (นนฺทาติ ภควา,)}

\prose{ชาติชราย นิวุตาติ พฺรูมิ.}

\setlength{\csromangathapagewidth}{0pt}
\setlength{\csromangathawakwidth}{0pt}
\csromangathameasuregroup{}{เย สีธ ทิฏฺฐํ ว สุตํ มุตํ วา, \\ สีลพฺพตํ วาปิ ปหาย สพฺพํ. \\ อเนกรูปมฺปิ ปหาย สพฺพํ, \\ ตณฺหํ ปริญฺญาย อนาสวาเส.}
\csromangathameasurewak{เย สีธ ทิฏฺฐํ ว สุตํ มุตํ วา, \\ สีลพฺพตํ วาปิ ปหาย สพฺพํ. \\ อเนกรูปมฺปิ ปหาย สพฺพํ, \\ ตณฺหํ ปริญฺญาย อนาสวาเส.}
\setlength{\csromangathaleftcol}{0pt}
\csromangathagroup{\csromangathastanza{\csromangathawak{เย สีธ ทิฏฺฐํ ว สุตํ มุตํ วา, \\ สีลพฺพตํ วาปิ ปหาย สพฺพํ. \\ อเนกรูปมฺปิ ปหาย สพฺพํ, \\ ตณฺหํ ปริญฺญาย อนาสวาเส.}}}

\prose{เต เว นรา โอฆติณฺณาติ พฺรูมีติ.}

\setlength{\csromangathapagewidth}{0pt}
\setlength{\csromangathawakwidth}{0pt}
\csromangathameasuregroup{}{\textbf{เอตา’ภินนฺทามิ วโจ มเหสิโน,} \\ \textbf{สุกิตฺติตํ โคตม’นูปธีกํ.} \\ เย สีธ ทิฏฺฐํ ว สุตํ มุตํ วา, \\ สีลพฺพตํ วาปิ ปหาย สพฺพํ. \\ อเนกรูปมฺปิ ปหาย สพฺพํ, \\ ตณฺหํ ปริญฺญาย อนาสวาเส. \\ \textbf{อหมฺปิ เต โอฆติณฺณาติ พฺรูมิ.}}
\csromangathameasurewak{\textbf{เอตา’ภินนฺทามิ วโจ มเหสิโน,} \\ \textbf{สุกิตฺติตํ โคตม’นูปธีกํ.} \\ เย สีธ ทิฏฺฐํ ว สุตํ มุตํ วา, \\ สีลพฺพตํ วาปิ ปหาย สพฺพํ. \\ อเนกรูปมฺปิ ปหาย สพฺพํ, \\ ตณฺหํ ปริญฺญาย อนาสวาเส. \\ \textbf{อหมฺปิ เต โอฆติณฺณาติ พฺรูมิ.}}
\setlength{\csromangathaleftcol}{0pt}
\csromangathagroup{\csromangathastanza{\csromangathawak{\csromangathaitemline{52}{\textbf{เอตา’ภินนฺทามิ วโจ มเหสิโน,}} \\ \csromangathacontline{\textbf{สุกิตฺติตํ โคตม’นูปธีกํ.}} \\ \csromangathacontline{เย สีธ ทิฏฺฐํ ว สุตํ มุตํ วา,} \\ \csromangathacontline{สีลพฺพตํ วาปิ ปหาย สพฺพํ.} \\ \csromangathacontline{อเนกรูปมฺปิ ปหาย สพฺพํ,} \\ \csromangathacontline{ตณฺหํ ปริญฺญาย อนาสวาเส.} \\ \csromangathacontline{\textbf{อหมฺปิ เต โอฆติณฺณาติ พฺรูมิ.}}}}}

\prose{\csromanfolio{123}\textbf{เอตา’ภินนฺทามิ วโจ มเหสิโน}ติ \textbf{เอตนฺ}ติ ตุยฺหํ วจนํ พฺยปฺปถํ เทสนํ อนุสาสนํ อนุสิฏฺฐํ นนฺทามิ อภินนฺทามิ โมทามิ อนุโมทามิ อิจฺฉามิ สาทิยามิ ปตฺถยามิ ปิหยามิ อภิชปฺปามิ. \textbf{มเหสิโน}ติ กิํ มเหสิ ภควา. มหนฺตํ สีลกฺขนฺธํ เอสี คเวสี ปริเยสีติ มเหสิ ฯเปฯ “กหํ นราสโภ”ติ มเหสีติ เอตาภินนฺทามิ วโจ มเหสิโน.}

\prose{\textbf{สุกิตฺติตํ โคตม’นูปธีกนฺ}ติ \textbf{สุกิตฺติตนฺ}ติ สุกิตฺติตํ สุอาจิกฺขิตํ\csromansharedfootnote{13475f1732aa1070}{สฺวาจิกฺขิตํ (ก)} สุเทสิตํ สุปญฺญปิตํ สุปฏฺฐปิตํ สุวิวฏํ สุวิภตฺตํ สุ- อุตฺตานีกตํ สุปกาสิตํ. \textbf{โคตม’นูปธีกนฺ}ติ อุปธี วุจฺจนฺติ กิเลสา จ ขนฺธา จ อภิสงฺขารา จ. อุปธิปฺปหานํ อุปธิ\-วูปสมํ อุปธิ\-นิสฺสคฺคํ อุปธิ\-ปฏิปฺปสฺสทฺธํ อมตํ นิพฺพานนฺติ สุกิตฺติตํ โคตม’นูปธีกํ.}

\prose{\textbf{เย สีธ ทิฏฺฐํ ว สุตํ มุตํ วา, สีลพฺพตํ วาปิ ปหาย สพฺพนฺ}ติ เย สพฺพา ทิฏฺฐสุทฺธิโย ปหาย ชหิตฺวา ปชหิตฺวา วิโนเทตฺวา พฺยนฺตีกริตฺวา อนภาวํ คเมตฺวา. เย สพฺพา สุตสุทฺธิโย ฯเปฯ. เย สพฺพา มุตสุทฺธิโย. เย สพฺพา ทิฏฺฐสุต\-มุต\-สุทฺธิโย. เย สพฺพา สีลสุทฺธิโย. เย สพฺพา วตสุทฺธิโย. เย สพฺพา สีลพฺพต\-สุทฺธิโย ปหาย ชหิตฺวา ปชหิตฺวา วิโนเทตฺวา พฺยนฺตีกริตฺวา อนภาวํ คเมตฺวาติ เย สีธ ทิฏฺฐํ ว สุตํ มุตํ วา, สีลพฺพตํ วาปิ ปหาย สพฺพํ.}

\prose{\textbf{อเนกรูปมฺปิ ปหาย สพฺพนฺ}ติ อเนกวิธ\-โกตูหล\-มงฺคเลน สุทฺธิํ วิสุทฺธิํ ปริสุทฺธิํ, มุตฺติํ วิมุตฺติํ ปริมุตฺติํ ปหาย ชหิตฺวา ปชหิตฺวา วิโนเทตฺวา พฺยนฺตีกริตฺวา อนภาวํ คเมตฺวาติ อเนกรูปมฺปิ ปหาย สพฺพํ.}

\prose{\textbf{ตณฺหํ ปริญฺญาย อนาสวาเส, อหมฺปิ เต โอฆติณฺณาติ พฺรูมี}ติ \textbf{ตณฺหา}ติ รูปตณฺหา, สทฺทตณฺหา, คนฺธตณฺหา, รสตณฺหา, โผฏฺฐพฺพ\-ตณฺหา, ธมฺมตณฺหา. \textbf{ตณฺหํ ปริญฺญายา}ติ ตณฺหํ ตีหิ ปริญฺญาหิ ปริชานิตฺวา ญาตปริญฺญาย ตีรณ\-ปริญฺญาย\csromansharedfootnote{1652342efb19f03d}{ติรณปริญฺญาย (สฺยา)} ปหาน\-ปริญฺญาย. กตมา ญาตปริญฺญา, ตณฺหํ ชานาติ “อยํ รูปตณฺหา, อยํ สทฺทตณฺหา, อยํ คนฺธตณฺหา, อยํ รสตณฺหา, อยํ โผฏฺฐพฺพ\-ตณฺหา, อยํ ธมฺมตณฺหา”ติ ชานาติ ปสฺสติ. อยํ ญาตปริญฺญา.}

\prose{กตมา ตีรณปริญฺญา, เอวํ ญาตํ กตฺวา ตณฺหํ ตีเรติ\csromansharedfootnote{e167d2cb38d4d382}{ติเรติ (สฺยา)} อนิจฺจโต ทุกฺขโต โรคโต คณฺฑโต สลฺลโต อฆโต อาพาธโต}

\prose{\csromanfolio{124}ปรโต ปโลกโต อีติโต อุปทฺทวโต ภยโต อุปสคฺคโต จลโต ปภงฺคุโต อทฺธุวโต อตาณโต อเลณโต อสรณโต อสรณีภูตโต ริตฺตโต ตุจฺฉโต สุญฺญโต อนตฺตโต อาทีนวโต วิปริณาม\-ธมฺมโต อสารกโต อฆมูลโต วธกโต วิภวโต สาสวโต สงฺขตโต มารามิสโต ชาติธมฺมโต ชราธมฺมโต พฺยาธิธมฺมโต มรณธมฺมโต โสกปริเทวทุกฺขโทมนสฺสุปายาส- ธมฺมโต สํกิเลส\-ธมฺมโต สมุทยโต อตฺถงฺคมโต อสฺสาทโต อาทีนวโต นิสฺสรณโต ตีเรติ. อยํ ตีรณปริญฺญา.}

\prose{กตมา ปหานปริญฺญา, เอวํ ตีรยิตฺวา ตณฺหํ ปชหติ วิโนเทติ พฺยนฺตีกโรติ อนภาวํ คเมติ. อยํ ปหานปริญฺญา.}

\prose{\textbf{ตณฺหํ ปริญฺญายา}ติ ตณฺหํ อิมาหิ ตีหิ ปริญฺญาหิ ปริชานิตฺวา. อ\textbf{นาสวา}ติ จตฺตาโร อาสวา กามาสโว ภวาสโว ทิฏฺฐาสโว อวิชฺชาสโว. เยสํ อิเม อาสวา ปหีนา อุจฺฉินฺนมูลา ตาลาวตฺถุกตา อนภาวํกตา อายติํ อนุปฺปาทธมฺมา. เต วุจฺจนฺติ อนาสวา. อรหนฺโต ขีณาสวา. \textbf{ตณฺหํ ปริญฺญาย อนาสวาเส, อหมฺปิ เต โอฆติณฺณาติ พฺรูมี}ติ เย ตณฺหํ ปริญฺญาย อนาสวา, อหมฺปิ เต กาโมฆํ ติณฺณา, ภโวฆํ ติณฺณา, ทิฏฺโฐฆํ ติณฺณา, อวิชฺโชฆํ ติณฺณา, สพฺพ\-สํสารปถํ ติณฺณา อุตฺติณฺณา นิตฺถิณฺณา อติกฺกนฺตา สมติกฺกนฺตา วีติวตฺตาติ พฺรูมิ วทามีติ ตณฺหํ ปริญฺญาย อนาสวาเส, อหมฺปิ เต โอฆติณฺณาติ พฺรูมิ. เตนาห โส พฺราหฺมโณ\csromandash{}}

\setlength{\csromangathapagewidth}{0pt}
\setlength{\csromangathawakwidth}{0pt}
\csromangathameasuregroup{}{เอตา’ภินนฺทามิ วโจ มเหสิโน, \\ สุกิตฺติตํ โคตม’นูปธีกํ. \\ เย สีธ ทิฏฺฐํ ว สุตํ มุตํ วา, \\ สีลพฺพตํ วาปิ ปหาย สพฺพํ. \\ อเนกรูปมฺปิ ปหาย สพฺพํ. \\ ตณฺหํ ปริญฺญาย อนาสวาเส. \\ อหมฺปิ เต โอฆติณฺณาติ พฺรูมีติ.}
\csromangathameasurewak{เอตา’ภินนฺทามิ วโจ มเหสิโน, \\ สุกิตฺติตํ โคตม’นูปธีกํ. \\ เย สีธ ทิฏฺฐํ ว สุตํ มุตํ วา, \\ สีลพฺพตํ วาปิ ปหาย สพฺพํ. \\ อเนกรูปมฺปิ ปหาย สพฺพํ. \\ ตณฺหํ ปริญฺญาย อนาสวาเส. \\ อหมฺปิ เต โอฆติณฺณาติ พฺรูมีติ.}
\setlength{\csromangathaleftcol}{0pt}
\csromangathagroup{\csromangathastanza{\csromangathawak{เอตา’ภินนฺทามิ วโจ มเหสิโน, \\ สุกิตฺติตํ โคตม’นูปธีกํ. \\ เย สีธ ทิฏฺฐํ ว สุตํ มุตํ วา, \\ สีลพฺพตํ วาปิ ปหาย สพฺพํ.}}\csromangathastanza{\csromangathawak{อเนกรูปมฺปิ ปหาย สพฺพํ. \\ ตณฺหํ ปริญฺญาย อนาสวาเส. \\ อหมฺปิ เต โอฆติณฺณาติ พฺรูมีติ.}}}

\vspace{\tipitakaparskip}
\csromancenter{นนฺทมาณวปุจฺฉา\-นิทฺเทโส สตฺตโม.}
\csromanmatikaanchor{mtk.29}
\csromantocmarkref{subsection}{8. เหมกมาณวปุจฺฉานิทฺเทส}{mtk.29}
\bookmark[dest={mtk.29},level=2]{8. เหมกมาณวปุจฺฉานิทฺเทส}
\csromanheader{2}{8. เหมกมาณวปุจฺฉา\-นิทฺเทส \textbf{8. เหมกมาณวปุจฺฉา}\-\textbf{นิทฺเทส}}
\csromanhangingbegin
\hangingheaditem{53}{\csromanfolio{125}\textbf{เย เม ปุพฺเพ วิยากํสุ, (อิจฺจายสฺมา เหมโก,)}}
\hangingbody{\textbf{หุรํ โคตมสาสนา.}}
\hangingbody{\textbf{อิจฺจาสิ อิติ ภวิสฺสติ, สพฺพํ ตํ อิติหีติหํ.}}
\hangingbody{\textbf{สพฺพํ ตํ ตกฺกวฑฺฒนํ, นาหํ ตตฺถ อภิรมิํ. (1)}}
\csromanhangingend

\prose{\textbf{เย เม ปุพฺเพ วิยากํสู}ติ โย จ พาวรี พฺราหฺมโณ เย จญฺเญ ตสฺส อาจริยา, เต สกํ ทิฏฺฐิํ สกํ ขนฺติํ สกํ รุจิํ สกํ ลทฺธิํ สกํ อชฺฌาสยํ สกํ อธิปฺปายํ พฺยากํสุ อาจิกฺขิํสุ เทสยิํสุ ปญฺญปิํสุ ปฏฺฐปิํสุ วิวริํสุ วิภชิํสุ อุตฺตานีอกํสุ ปกาเสสุนฺติ เย เม ปุพฺเพ วิยากํสุ. \textbf{อิจฺจายสฺมา เหมโก}ติ \textbf{อิจฺจา}ติ ปทสนฺธิ ฯเปฯ ปทา\-นุปุพฺพตา\-เปตํ อิจฺจาติ. \textbf{อายสฺมา}ติ ปิยวจนํ ฯเปฯ. \textbf{เหมโก}ติ ตสฺส พฺราหฺมณสฺส นามํ ฯเปฯ อภิลาโปติ อิจฺจายสฺมา เหมโก.}

\prose{\textbf{หุรํ โคตมสาสนา}ติ หุรํ โคตมสาสนา ปรํ โคตมสาสนา ปุเร โคตมสาสนา ปฐมตรํ โคตมสาสนา พุทฺธสาสนา ชินสาสนา ตถาคต\-สาสนา\csromansharedfootnote{eff388456dc9a221}{ตถาคตสาสนา เทวสาสนา (ก)} อรหนฺต\-สาสนาติ หุรํ โคตมสาสนา.}

\prose{\textbf{อิจฺจาสิ อิติ ภวิสฺสตี}ติ เอวํ กิร อาสิ, เอวํ กิร ภวิสฺสตีติ อิจฺจาสิ อิติ ภวิสฺสติ.}

\prose{\textbf{สพฺพํ ตํ อิติหีติหนฺ}ติ สพฺพํ ตํ อิติหีติหํ อิติกิราย ปรํปราย ปิฏก\-สมฺปทาย ตกฺกเหตุ นยเหตุ อาการ\-ปริวิตกฺเกน ทิฏฺฐิ\-นิชฺฌาน\-กฺขนฺติยา น สามํ สยม\-ภิญฺญาตํ น อตฺต\-ปจฺจกฺขธมฺมํ กถยิํสูติ สพฺพํ ตํ อิติหีติหํ.}

\prose{\textbf{สพฺพํ ตํ ตกฺก}\-\textbf{วฑฺฒนนฺ}ติ สพฺพํ ตํ ตกฺกวฑฺฒนํ วิตกฺก\-วฑฺฒนํ สงฺกปฺป\-วฑฺฒนํ กามวิตกฺกวฑฺฑฺธนํ พฺยาปาท\-วิตกฺก\-วฑฺฒนํ วิหิํสา\-วิตกฺก\-วฑฺฒนํ ญาติ\-วิตกฺก\-วฑฺฒนํ ชนปท\-วิตกฺก\-วฑฺฒนํ อมรา\-วิตกฺก\-วฑฺฒนํ ปรานุทยตาปฏิสํยุตฺต\-วิตกฺก\-วฑฺฒนํ ลาภ\-สกฺการ\-สิโลก- ปฏิสํยุตฺต\-วิตกฺก\-วฑฺฒนํ อนวญฺญตฺติปฏิสํยุตฺตวิตกฺก\-วฑฺฒนนฺติ สพฺพํ ตํ ตกฺกวฑฺฒนํ.}

\prose{\csromanfolio{126}\textbf{นาหํ ตตฺถ อภิรมินฺ}ติ นาหํ ตตฺถ อภิรมิํ น วินฺทิํ นาธิคจฺฉิํ น ปฏิลภินฺติ นาหํ ตตฺถ อภิรมิํ. เตนาห โส พฺราหฺมโณ\csromandash{}}

\prose{เย เม ปุพฺเพ วิยากํสุ, (อิจฺจายสฺมา เหมโก,)}

\prose{หุรํ โคตมสาสนา.}

\setlength{\csromangathapagewidth}{0pt}
\setlength{\csromangathawakwidth}{0pt}
\csromangathameasuregroup{อิจฺจาสิ อิติ ภวิสฺสติ, \\ สพฺพํ ตํ ตกฺกวฑฺฒนํ, \\ \textbf{ตฺวญฺจ เม ธมฺมมกฺ}ขาหิ, \\ \textbf{ยํ วิทิตฺวา สโต จฺ}อรํ,}{\csromangathabat{อิจฺจาสิ อิติ ภวิสฺสติ,}{สพฺพํ ตํ อิติหีติหํ.} \\ \csromangathabat{สพฺพํ ตํ ตกฺกวฑฺฒนํ,}{นาหํ ตตฺถ อภิรมินฺติ.} \\ \csromangathabat{\textbf{ตฺวญฺจ เม ธมฺมมกฺ}ขาหิ,}{\textbf{ตณฺหา}นิคฺฆาตนํ มุนิ.} \\ \csromangathabat{\textbf{ยํ วิทิตฺวา สโต จฺ}อรํ,}{ตเร โลเก วิสตฺติกํ.}}
\csromangathasetleft{อิจฺจาสิ อิติ ภวิสฺสติ, \\ สพฺพํ ตํ ตกฺกวฑฺฒนํ, \\ \textbf{ตฺวญฺจ เม ธมฺมมกฺ}ขาหิ, \\ \textbf{ยํ วิทิตฺวา สโต จฺ}อรํ,}
\csromangathagroup{\csromangathastanza{\csromangathabat{อิจฺจาสิ อิติ ภวิสฺสติ,}{สพฺพํ ตํ อิติหีติหํ.} \\ \csromangathabat{สพฺพํ ตํ ตกฺกวฑฺฒนํ,}{นาหํ ตตฺถ อภิรมินฺติ.}}\csromangathastanza{\csromangathaitembat{54}{\textbf{ตฺวญฺจ เม ธมฺมมกฺ}ขาหิ,}{\textbf{ตณฺหา}\-นิคฺฆาตนํ มุนิ.} \\ \csromangathacontbat{\textbf{ยํ วิทิตฺวา สโต จฺ}อรํ,}{ตเร โลเก วิสตฺติกํ.}}}

\prose{\textbf{ตฺวญฺจ เม ธมฺม}\-\textbf{มกฺขาหี}ติ \textbf{ตฺวนฺ}ติ ภควนฺตํ ภณติ. \textbf{ธมฺม}\-\textbf{มกฺขาหี}ติ \textbf{ธมฺมนฺ}ติ อาทิกลฺยาณํ มชฺเฌกลฺยาณํ ปริโยสาน\-กลฺยาณํ สาตฺถํ สพฺยญฺชนํ เกวล\-ปริปุณฺณํ ปริสุทฺธํ พฺรหฺมจริยํ จตฺตาโร สติปฏฺฐาเน จตฺตาโร สมฺมปฺปธาเน จตฺตาโร อิทฺธิปาเท ปญฺจินฺทฺริยานิ ปญฺจ พลานิ สตฺต โพชฺฌงฺเค อริยํ อฏฺฐงฺคิกํ มคฺคํ นิพฺพานญฺจ นิพฺพานคามินิญฺจ ปฏิปทํ อกฺขาหิ อาจิกฺขาหิ เทเสหิ ปญฺญเปหิ ปฏฺฐเปหิ วิวราหิ วิภชาหิ อุตฺตานีกโรหิ ปกาเสหีติ ตฺวญฺจ เม ธมฺมมกฺขาหิ.}

\prose{\textbf{ตณฺหา}\-\textbf{นิคฺฆาตนํ มุนี}ติ ตณฺหาติ รูป\textbf{ตณฺหา} ฯเปฯ ธมฺมตณฺหา. ตณฺหา\-นิคฺฆาตนํ ตณฺหาปหานํ ตณฺหาวูปสมํ ตณฺหา\-ปฏินิสฺสคฺคํ ตณฺหา\-ปฏิปฺปสฺสทฺธิํ อมตํ นิพฺพานํ. \textbf{มุนี}ติ โมนํ วุจฺจติ ญาณํ ฯเปฯ สงฺค\-ชาลม\-ติจฺจ โส มุนีติ ตณฺหา\-นิคฺฆาตนํ มุนิ.}

\prose{\textbf{ยํ วิทิตฺวา สโต จรนฺ}ติ ยํ วิทิตํ กตฺวา ตุลยิตฺวา ตีรยิตฺวา วิภาวยิตฺวา วิภูตํ กตฺวา. “สพฺเพ สงฺขารา อนิจฺจา”ติ วิทิตํ กตฺวา ตุลยิตฺวา ตีรยิตฺวา วิภาวยิตฺวา วิภูตํ กตฺวา. “สพฺเพ สงฺขารา ทุกฺขา”ติ ฯเปฯ. “สพฺเพ ธมฺมา อนตฺตา”ติ. “ยํ กิญฺจิ สมุทย\-ธมฺมํ, สพฺพํ ตํ นิโรธธมฺมนฺ”ติ วิทิตํ กตฺวา ตุลยิตฺวา ตีรยิตฺวา วิภาวยิตฺวา วิภูตํ กตฺวา. \textbf{สโต}ติ จตูหิ การเณหิ สโต, กาเย กายานุปสฺสนา\-สติปฏฺฐานํ ภาเวนฺโต สโต ฯเปฯ โส วุจฺจติ สโต. \textbf{จรนฺ}ติ จรนฺโต วิหรนฺโต อิริยนฺโต วตฺเตนฺโต ปาเลนฺโต ยเปนฺโต ยาเปนฺโตติ ยํ วิทิตฺวา สโต จรํ.}

\prose{\textbf{ตเร โลเก วิสตฺติกนฺ}ติ วิสตฺติกา วุจฺจติ ตณฺหา. โย ราโค สาราโค ฯเปฯ อภิชฺฌา โลโภ อกุสลมูลํ, \textbf{วิสตฺติกา}ติ \csromanfolio{127}เกนฏฺเฐน วิสตฺติกา ฯเปฯ วิสฏา วิตฺถตาติ วิสตฺติกา. \textbf{โลเก}ติ อปายโลเก มนุสฺสโลเก เทวโลเก, ขนฺธโลเก ธาตุโลเก อายตนโลเก. \textbf{ตเร โลเก วิสตฺติกนฺ}ติ โลเก เวสา วิสตฺติกา, โลเก เวตํ วิสตฺติกํ สโต ตเรยฺยํ อุตฺตเรยฺยํ ปตเรยฺยํ สมติกฺกเมยฺยํ วีติวตฺเตยฺยนฺติ ตเร โลเก วิสตฺติกํ. เตนาห โส พฺราหฺมโณ\csromandash{}}

\setlength{\csromangathapagewidth}{0pt}
\setlength{\csromangathawakwidth}{0pt}
\csromangathameasuregroup{ตฺวญฺจ เม ธมฺมมกฺขาหิ, \\ ยํ วิทิตฺวา สโต จรํ, \\ \textbf{อิธ ทิฏฺฐ}\textbf{สุต}\textbf{มุต}วิญฺญาเตสุ, \\ \textbf{ฉนฺทราค}\textbf{วิโนท}นํ,}{\csromangathabat{ตฺวญฺจ เม ธมฺมมกฺขาหิ,}{ตณฺหานิคฺฆาตนํ มุนิ.} \\ \csromangathabat{ยํ วิทิตฺวา สโต จรํ,}{ตเร โลเก วิสตฺติกนฺติ.} \\ \csromangathabat{\textbf{อิธ ทิฏฺฐ}\textbf{สุต}\textbf{มุต}วิญฺญาเตสุ,}{ปิยรูเปสุ เหมก.} \\ \csromangathabat{\textbf{ฉนฺทราค}\textbf{วิโนท}นํ,}{นิพฺพานปทมจฺจุตํ.}}
\csromangathasetleft{ตฺวญฺจ เม ธมฺมมกฺขาหิ, \\ ยํ วิทิตฺวา สโต จรํ, \\ \textbf{อิธ ทิฏฺฐ}\textbf{สุต}\textbf{มุต}วิญฺญาเตสุ, \\ \textbf{ฉนฺทราค}\textbf{วิโนท}นํ,}
\csromangathagroup{\csromangathastanza{\csromangathabat{ตฺวญฺจ เม ธมฺมมกฺขาหิ,}{ตณฺหา\-นิคฺฆาตนํ มุนิ.} \\ \csromangathabat{ยํ วิทิตฺวา สโต จรํ,}{ตเร โลเก วิสตฺติกนฺติ.}}\csromangathastanza{\csromangathaitembat{55}{\textbf{อิธ ทิฏฺฐ}\-\textbf{สุต}\-\textbf{มุต}\-วิญฺญาเตสุ,}{ปิยรูเปสุ เหมก.} \\ \csromangathacontbat{\textbf{ฉนฺทราค}\-\textbf{วิโนท}นํ,}{นิพฺพานปทม\-จฺจุตํ.}}}

\prose{อิธ ทิฏฺฐ\-สุต\-มุต\-วิญฺญาเตสูติ ทิฏฺฐนฺ\textbf{ติ} \textbf{จกฺขุนา ทิฏฺฐํ.} สุตนฺ\textbf{ติ} \textbf{โสเตน สุตํ.} มุตนฺติ \textbf{ฆาเนน ฆายิตํ, ชิวฺหาย สายิตํ, กาเยน ผุฏฺฐํ.} วิญฺญาตนฺ\textbf{ติ} \textbf{มนสา วิญฺญาตนฺติ อิธ ทิฏฺฐ}\-\textbf{สุต}\-\textbf{มุต}\-\textbf{วิญฺญาเตสุ.}}

\prose{\textbf{ปิยรูเปสุ เหมกา}ติ กิญฺจ โลเก ปิยรูปํ สาตรูปํ, จกฺขุ\csromansharedfootnote{b15203dbe0312afb}{จกฺขุํ (สฺยา, ก)} โลเก ปิยรูปํ สาตรูปํ. โสตํ โลเก. ฆานํ โลเก. ชิวฺหา โลเก. กาโย โลเก. มโน โลเก ปิยรูปํ สาตรูปํ. รูปา โลเก ปิยรูปํ สาตรูปํ. สทฺทา โลเก. คนฺธา โลเก. รสา โลเก. โผฏฺฐพฺพา โลเก. ธมฺมา โลเก ปิยรูปํ สาตรูปํ. จกฺขุวิญฺญาณํ โลเก ปิยรูปํ สาตรูปํ. โสตวิญฺญาณํ โลเก ปิยรูปํ สาตรูปํ. ฆานวิญฺญาณํ โลเก. ชิวฺหาวิญฺญาณํ โลเก. กายวิญฺญาณํ โลเก. มโนวิญฺญาณํ โลเก ปิยรูปํ สาตรูปํ. จกฺขุ\-สมฺผสฺโส โลเก. โสตสมฺผสฺโส โลเก. ฆานสมฺผสฺโส โลเก. ชิวฺหาสมฺผสฺโส โลเก. กายสมฺผสฺโส โลเก. มโนสมฺผสฺโส โลเก ปิยรูปํ สาตรูปํ. จกฺขุ\-สมฺผสฺสชา เวทนา โลเก ปิยรูปํ สาตรูปํ. โสต\-สมฺผสฺสชา เวทนา. ฆาน\-สมฺผสฺสชา เวทนา. ชิวฺหา\-สมฺผสฺสชา เวทนา. กาย\-สมฺผสฺสชา เวทนา. มโน\-สมฺผสฺสชา เวทนา โลเก ปิยรูปํ สาตรูปํ. รูปสญฺญา โลเก. สทฺทสญฺญา โลเก. คนฺธสญฺญา โลเก. รสสญฺญา โลเก. โผฏฺฐพฺพ\-สญฺญา โลเก. ธมฺมสญฺญา โลเก ปิยรูปํ สาตรูปํ. \csromanfolio{128}รูปสญฺเจตนา โลเก. สทฺทสญฺเจตนา โลเก. คนฺธ\-สญฺเจตนา โลเก. รสสญฺเจตนา โลเก. โผฏฺฐพฺพ\-สญฺเจตนา โลเก. ธมฺม\-สญฺเจตนา โลเก ปิยรูปํ สาตรูปํ. รูปตณฺหา โลเก. สทฺทตณฺหา โลเก. คนฺธตณฺหา โลเก. รสตณฺหา โลเก. โผฏฺฐพฺพ\-ตณฺหา โลเก. ธมฺมตณฺหา โลเก ปิยรูปํ สาตรูปํ. รูปวิตกฺโก โลเก. สทฺทวิตกฺโก โลเก. คนฺธวิตกฺโก โลเก. รสวิตกฺโก โลเก. โผฏฺฐพฺพ\-วิตกฺโก โลเก. ธมฺมวิตกฺโก โลเก ปิยรูปํ สาตรูปํ. รูปวิจาโร โลเก ปิยรูปํ สาตรูปํ. สทฺทวิจาโร โลเก. คนฺธวิจาโร โลเก. รสวิจาโร โลเก. โผฏฺฐพฺพ\-วิจาโร โลเก. ธมฺมวิจาโร โลเก ปิยรูปํ สาตรูปนฺติ ปิยรูเปสุ เหมก.}

\prose{\textbf{ฉนฺทราค}\-\textbf{วิโนทนนฺ}ติ \textbf{ฉนฺทราโค}ติ โย กาเมสุ กามจฺฉนฺโท กามราโค กามนนฺที กามตณฺหา กามสิเนโห กามปริฬาโห กามมุจฺฉา กามชฺโฌสานํ กาโมโฆ กามโยโค กามุปาทานํ กามจฺฉนฺท\-นีวรณํ. \textbf{ฉนฺทราค}\-\textbf{วิโนทนนฺ}ติ ฉนฺท\-ราค\-ปฺปหานํ ฉนฺทราค\-วูปสมํ ฉนฺทราค\-ปฏินิสฺสคฺคํ ฉนฺทราค\-ปฏิปฺปสฺสทฺธํ อมตํ นิพฺพานนฺติ ฉนฺทราค\-วิโนทนํ.}

\prose{\textbf{นิพฺพานปทมจฺจุตนฺ}ติ นิพฺพานปทํ ตาณปทํ เลณปทํ สรณปทํ อภยปทํ. \textbf{อจฺจุตนฺ}ติ นิจฺจํ ธุวํ สสฺสตํ อวิปริณาม\-ธมฺมนฺติ นิพฺพานปทม\-จฺจุตํ. เตนาห ภควา\csromandash{}}

\setlength{\csromangathapagewidth}{0pt}
\setlength{\csromangathawakwidth}{0pt}
\csromangathameasuregroup{อิธ ทิฏฺฐสุตมุตวิญฺญาเตสุ, \\ ฉนฺทราควิโนทนํ, \\ \textbf{เอตทญฺญาย เย สตฺ}อา, \\ \textbf{อุปสนฺตา จ เต ส}ทา,}{\csromangathabat{อิธ ทิฏฺฐสุตมุตวิญฺญาเตสุ,}{ปิยรูเปสุ เหมก.} \\ \csromangathabat{ฉนฺทราควิโนทนํ,}{นิพฺพานปทมจฺจุตนฺติ.} \\ \csromangathabat{\textbf{เอตทญฺญาย เย สตฺ}อา,}{ทิฏฺฐธมฺมาภินิพฺพุตา.} \\ \csromangathabat{\textbf{อุปสนฺตา จ เต ส}ทา,}{ติณฺณา โลเก วิสตฺติกํ.}}
\csromangathasetleft{อิธ ทิฏฺฐสุตมุตวิญฺญาเตสุ, \\ ฉนฺทราควิโนทนํ, \\ \textbf{เอตทญฺญาย เย สตฺ}อา, \\ \textbf{อุปสนฺตา จ เต ส}ทา,}
\csromangathagroup{\csromangathastanza{\csromangathabat{อิธ ทิฏฺฐ\-สุต\-มุต\-วิญฺญาเตสุ,}{ปิยรูเปสุ เหมก.} \\ \csromangathabat{ฉนฺทราค\-วิโนทนํ,}{นิพฺพานปทมจฺจุตนฺติ.}}\csromangathastanza{\csromangathaitembat{56}{\textbf{เอตทญฺญาย เย สตฺ}อา,}{ทิฏฺฐธมฺมา\-ภินิพฺพุตา.} \\ \csromangathacontbat{\textbf{อุปสนฺตา จ เต ส}ทา,}{ติณฺณา โลเก วิสตฺติกํ.}}}

\prose{\textbf{เอตทญฺญาย เย สตา}ติ \textbf{เอตนฺ}ติ อมตํ นิพฺพานํ. โย โส สพฺพ\-สงฺขาร\-สมโถ สพฺพู\-ปธิ\-ปฏินิสฺสคฺโค ตณฺหกฺขโย วิราโค นิโรโธ นิพฺพานํ. \textbf{อญฺญายา}ติ อญฺญาย ชานิตฺวา ตุลยิตฺวา ตีรยิตฺวา วิภาวยิตฺวา วิภูตํ กตฺวา. “สพฺเพ สงฺขารา อนิจฺจา”ติ อญฺญาย ชานิตฺวา ตุลยิตฺวา ตีรยิตฺวา วิภาวยิตฺวา วิภูตํ กตฺวา. “สพฺเพ สงฺขารา ทุกฺขา”ติ. “สพฺเพ ธมฺมา อนตฺตา”ติ ฯเปฯ. “ยํ กิญฺจิ สมุทย\-ธมฺมํ, สพฺพํ ตํ นิโรธธมฺมนฺ”ติ อญฺญาย ชานิตฺวา ตุลยิตฺวา ตีรยิตฺวา \csromanfolio{129}วิภาวยิตฺวา วิภูตํ กตฺวา. \textbf{เย}ติ อรหนฺโต ขีณาสวา. \textbf{สตา}ติ จตูหิ การเณหิ สตา. กาเย กายานุปสฺสนา\-สติปฏฺฐานํ ภาวิตตฺตา สตา ฯเปฯ เต วุจฺจนฺติ สตาติ เอตทญฺญาย เย สตา.}

\prose{\textbf{ทิฏฺฐธมฺมา}\-\textbf{ภินิพฺพุตา}ติ \textbf{ทิฏฺฐธมฺมา}ติ ทิฏฺฐธมฺมา ญาตธมฺมา ตุลิตธมฺมา ตีริตธมฺมา วิภูตธมฺมา วิภาวิต\-ธมฺมา. “สพฺเพ สงฺขารา อนิจฺจา”ติ ทิฏฺฐธมฺมา ฯเปฯ. “ยํ กิญฺจิ สมุทย\-ธมฺมํ, สพฺพํ ตํ นิโรธธมฺมนฺ”ติ ทิฏฺฐธมฺมา ญาตธมฺมา ตุลิตธมฺมา ตีริตธมฺมา วิภูตธมฺมา วิภาวิต\-ธมฺมา. \textbf{อภินิพฺพุตา}ติ ราคสฺส นิพฺพาปิตตฺตา นิพฺพุตา. โทสสฺส นิพฺพาปิตตฺตา นิพฺพุตา. โมหสฺส นิพฺพาปิตตฺตา นิพฺพุตา. โกธสฺส. อุปานาหสฺส ฯเปฯ สพฺพา\-กุสลา\-ภิสงฺขารานํ สนฺตตฺตา สมิตตฺตา วูปสมิตตฺตา นิชฺฌาตตฺตา นิพฺพุตตฺตา วิคตตฺตา ปฏิปฺปสฺสทฺธ\-ตฺตา สนฺตา อุปสนฺตา วูปสนฺตา นิพฺพุตา ปฏิปฺปสฺสทฺธาติ ทิฏฺฐธมฺมา\-ภินิพฺพุตา.}

\prose{\textbf{อุปสนฺตา จ เต สทา}ติ \textbf{อุปสนฺตา}ติ ราคสฺส อุปสมิตตฺตา นิพฺพาปิตตฺตา อุปสนฺตา. โทสสฺส. โมหสฺส. โกธสฺส. อุปนาหสฺส ฯเปฯ สพฺพา\-กุสลา\-ภิสงฺขารานํ สนฺตตฺตา สมิตตฺตา วูปสมิตตฺตา นิชฺฌาตตฺตา นิพฺพุตตฺตา วิคตตฺตา ปฏิปฺปสฺสทฺธ\-ตฺตา สนฺตา อุปสนฺตา วูปสนฺตา นิพฺพุตา ปฏิปฺปสฺสทฺธาติ อุปสนฺตา. \textbf{เต}ติ อรหนฺโต กฺกฺ. \textbf{สทา}ติ สทา สพฺพกาลํ นิจฺจกาลํ ธุวกาลํ สตตํ สมิตํ อพฺโพกิณฺณํ โปงฺขานุโปงฺขํ อุทกูมิกชาตํ อวีจิ สนฺตติ สหิตํ ผสฺสิตํ ปุเรภตฺตํ ปจฺฉาภตฺตํ ปุริมยามํ มชฺฌิมยามํ ปจฺฉิมยามํ กาเฬ ชุเณฺห วสฺเส เหมนฺเต คิเมฺห ปุริเม วโยขนฺเธ มชฺฌิเม วโยขนฺเธ ปจฺฉิเม วโยขนฺเธติ อุปสนฺตา จ เต สทา.}

\prose{\textbf{ติณฺณา โลเก วิสตฺติกนฺ}ติ วิสตฺติกา วุจฺจติ ตณฺหา. โย ราโค สาราโค ฯเปฯ อภิชฺฌา โลโภ อกุสลมูลํ. \textbf{วิสตฺติกา}ติ เกนฏฺเฐน วิสตฺติกา ฯเปฯ วิสฏา วิตฺถตาติ วิสตฺติกา. \textbf{โลเก}ติ อปายโลเก ฯเปฯ อายตนโลเก. \textbf{ติณฺณา โลเก วิสตฺติกนฺ}ติ โลเก เวสา วิสตฺติกา, โลเก เวตํ วิสตฺติกํ ติณฺณา อุตฺติณฺณา นิตฺถิณฺณา อติกฺกนฺตา สมติกฺกนฺตา วีติวตฺตาติ ติณฺณา โลเก วิสตฺติกํ. เตนาห ภควา\csromandash{}}

\setlength{\csromangathapagewidth}{0pt}
\setlength{\csromangathawakwidth}{0pt}
\csromangathameasuregroup{เอตทญฺญาย เย สตา, \\ อุปสนฺตา จ เต สทา,}{\csromangathabat{เอตทญฺญาย เย สตา,}{ทิฏฺฐธมฺมาภินิพฺพุตา.} \\ \csromangathabat{อุปสนฺตา จ เต สทา,}{ติณฺณา โลเก วิสตฺติกนฺติ.}}
\csromangathasetleft{เอตทญฺญาย เย สตา, \\ อุปสนฺตา จ เต สทา,}
\csromangathagroup{\csromangathastanza{\csromanfolio{130}\csromangathabat{เอตทญฺญาย เย สตา,}{ทิฏฺฐธมฺมา\-ภินิพฺพุตา.} \\ \csromangathabat{อุปสนฺตา จ เต สทา,}{ติณฺณา โลเก วิสตฺติกนฺติ.}}}

\proseitem{8}{สห คาถา\-ปริโยสานา ฯเปฯ “สตฺถา เม ภนฺเต ภควา, สาวโกหมสฺมี”ติ.}

\vspace{\tipitakaparskip}
\csromancenter{เหมกมาณวปุจฺฉา\-นิทฺเทโส อฏฺฐโม.}
\csromanmatikaanchor{mtk.30}
\csromantocmarkref{subsection}{9. โตเทยฺยมาณวปุจฺฉานิทฺเทส}{mtk.30}
\bookmark[dest={mtk.30},level=2]{9. โตเทยฺยมาณวปุจฺฉานิทฺเทส}
\csromanheader{2}{9. \textbf{โตเทยฺยมาณวปุจฺฉา}\-\textbf{นิทฺเทส}}
\csromanhangingbegin
\hangingheaditem{57}{\textbf{ยสฺมิํ กามา น วสนฺติ, (อิจฺจายสฺมา โตเทยฺโย,)}}
\hangingbody{\textbf{ตณฺหา ยสฺส น วิชฺชติ.}}
\hangingbody{\textbf{กถํกถา จ โย ติณฺโณ,}}
\hangingbody{\textbf{วิโมกฺโข ตสฺส กีทิโส.}}
\csromanhangingend

\prose{\textbf{ยสฺมิํ กามา น วสนฺตี}ติ ยสฺมิํ กามา น วสนฺติ น สํวสนฺติ น อาวสนฺติ น ปริวสนฺตีติ ยสฺมิํ กามา น วสนฺติ. \textbf{อิจฺจายสฺมา โตเทยฺโย}ติ \textbf{อิจฺจา}ติ’ปทสนฺธิ ฯเปฯ ปทา\-นุปุพฺพตา\-เปตํ อิจฺจาติ. \textbf{อายสฺมา}ติ ปิยวจนํ ฯเปฯ. \textbf{โตเทยฺโย}ติ ตสฺส พฺราหฺมณสฺส นามํ ฯเปฯ อภิลาโปติ อิจฺจายสฺมา โตเทยฺโย.}

\prose{\textbf{ตณฺหา ยสฺส น วิชฺชตี}ติ ตณฺหา ยสฺส นตฺถิ น สติ น สํวิชฺชติ นุปลพฺภติ ญาณคฺคินา ทฑฺฒาติ ตณฺหา ยสฺส น วิชฺชติ.}

\prose{\textbf{กถํกถา จ โย ติณฺโณ}ติ กถํกถา จ โย ติณฺโณ อุตฺติณฺโณ นิตฺถิณฺโณ อติกฺกนฺโต สมติกฺกนฺโต วีติวตฺโตติ กถํกถา จ โย ติณฺโณ.}

\prose{\textbf{วิโมกฺโข ตสฺส กีทิโส}ติ วิโมกฺโข ตสฺส กีทิโส กิํสณฺฐิโต กิํปกาโร กิํปฏิภาโค อิจฺฉิตพฺโพติ วิโมกฺขํ ปุจฺฉตีติ วิโมกฺโข ตสฺส กีทิโส. เตนาห โส พฺราหฺมโณ\csromandash{}}

\prose{\textbf{ยสฺมิํ กามา น วสนฺติ, (อิจฺจายสฺมา โตเทยฺโย,)}}

\prose{\textbf{ตณฺหา ยสฺส น วิชฺชติ.}}

\setlength{\csromangathapagewidth}{0pt}
\setlength{\csromangathawakwidth}{0pt}
\csromangathameasuregroup{}{\textbf{กถํกถา จ โย ติณฺโณ,} \\ วิโมกฺโข ตสฺส กีทิโสติ.}
\csromangathameasurewak{\textbf{กถํกถา จ โย ติณฺโณ,} \\ วิโมกฺโข ตสฺส กีทิโสติ.}
\setlength{\csromangathaleftcol}{0pt}
\csromangathagroup{\csromangathastanza{\csromangathawak{\textbf{กถํกถา จ โย ติณฺโณ,} \\ วิโมกฺโข ตสฺส กีทิโสติ.}}}

\csromanheader{2}{9. โตเทยฺยมาณวปุจฺฉา\-นิทฺเทส}
\csromanhangingbegin
\hangingheaditem{58}{\csromanfolio{131}\textbf{ยสฺมิํ กามา น วสนฺติ, (โตเทยฺยาติ ภควา,)}}
\hangingbody{\textbf{ตณฺหา ยสฺส น วิชฺชติ.}}
\hangingbody{\textbf{กถํกถา จ โย ติณฺโณ,}}
\hangingbody{\textbf{วิโมกฺโข ตสฺส นา’ปโร.}}
\csromanhangingend

\prose{\textbf{ยสฺมิํ กามา น วสนฺตี}ติ \textbf{ยสฺมินฺ}ติ ยสฺมิํ ปุคฺคเล อรหนฺเต ขีณาสเว. \textbf{กามา}ติ อุทฺทานโต เทฺว กามา วตฺถุกามา จ กิเลสกามา จ ฯเปฯ. อิเม วุจฺจนฺติ วตฺถุกามา ฯเปฯ. อิเม วุจฺจนฺติ กิเลสกามา. \textbf{ยสฺมิํ กามา น วสนฺตี}ติ ยสฺมิํ กามา น วสนฺติ น สํวสนฺติ น อาวสนฺติ น ปริวสนฺตีติ ยสฺมิํ กามา น วสนฺติ.}

\prose{\textbf{โตเทยฺยาติ ภควา}ติ โตเทยฺยาติ ภควา ตํ พฺราหฺมณํ นาเมน อาลปติ. \textbf{ภควา}ติ คารวาธิวจนเมตํ ฯเปฯ สจฺฉิกา ปญฺญตฺติ, ยทิทํ ภควาติ โตเทยฺยาติ ภควา.}

\prose{\textbf{ตณฺหา ยสฺส น วิชฺชตี}ติ ตณฺหาติ รูป\textbf{ตณฺหา} สทฺทตณฺหา คนฺธตณฺหา รสตณฺหา โผฏฺฐพฺพ\-ตณฺหา ธมฺมตณฺหา. \textbf{ยสฺสา}ติ อรหโต ขีณาสวสฺส. \textbf{ตณฺหา ยสฺส น วิชฺชตี}ติ ตณฺหา ยสฺส นตฺถิ น อตฺถิ น สํวิชฺชติ นุปลพฺภติ, ปหีนา สมุจฺฉินฺนา วูปสนฺตา ปฏิปฺปสฺสทฺธา อภพฺพุ\-ปฺปตฺติกา ญาณคฺคินา ทฑฺฒาติ ตณฺหา ยสฺส น วิชฺชติ.}

\prose{\textbf{กถํกถา จ โย ติณฺโณ}ติ กถํกถา วุจฺจติ วิจิกิจฺฉา. ทุกฺเข กงฺขา ฯเปฯ ฉมฺภิตตฺตํ จิตฺตสฺส มโนวิเลโข. \textbf{โย}ติ โย โส อรหํ ขีณาสโว. \textbf{กถํกถา จ โย ติณฺโณ}ติ กถํกถา จ โย ติณฺโณ อุตฺติณฺโณ นิตฺถิณฺโณ อติกฺกนฺโต สมติกฺกนฺโต วีติวตฺโตติ กถํกถา จ โย ติณฺโณ.}

\prose{\textbf{วิโมกฺโข ตสฺส นา’ปโร}ติ นตฺถิ ตสฺส อปโร วิโมกฺโข. เยน วิโมกฺเขน วิมุจฺเจยฺย วิมุตฺโต โส. กตํ ตสฺส วิโมกฺเขน กรณียนฺติ วิโมกฺโข ตสฺส นา’ปโร. เตนาห ภควา\csromandash{}}

\prose{\textbf{ยสฺมิํ กามา น วสนฺติ, (โตเทยฺยาติ ภควา,)}}

\prose{\textbf{ตณฺหา ยสฺส น วิชฺชติ.}}

\setlength{\csromangathapagewidth}{0pt}
\setlength{\csromangathawakwidth}{0pt}
\csromangathameasuregroup{}{\textbf{กถํกถา จ โย ติณฺโณ,} \\ วิโมกฺโข ตสฺส นา’ปโรติ. \\ \textbf{นิราสโส โส อุท อาสสาโน,} \\ \textbf{ปญฺญาณวา โส อุท ปญฺญกปฺปี.} \\ \textbf{มุนิํ อหํ สกฺก ยถา วิชญฺญํ,} \\ \textbf{ตํ เม วิยาจิกฺข สมนฺตจกฺขุ.}}
\csromangathameasurewak{\textbf{กถํกถา จ โย ติณฺโณ,} \\ วิโมกฺโข ตสฺส นา’ปโรติ. \\ \textbf{นิราสโส โส อุท อาสสาโน,} \\ \textbf{ปญฺญาณวา โส อุท ปญฺญกปฺปี.} \\ \textbf{มุนิํ อหํ สกฺก ยถา วิชญฺญํ,} \\ \textbf{ตํ เม วิยาจิกฺข สมนฺตจกฺขุ.}}
\setlength{\csromangathaleftcol}{0pt}
\csromangathagroup{\csromangathastanza{\csromangathawak{\textbf{กถํกถา จ โย ติณฺโณ,} \\ วิโมกฺโข ตสฺส นา’ปโรติ.}}\csromangathastanza{\csromangathawak{\csromanfolio{132}\csromangathaitemline{59}{\textbf{นิราสโส โส อุท อาสสาโน,}} \\ \csromangathacontline{\textbf{ปญฺญาณวา โส อุท ปญฺญกปฺปี.}} \\ \csromangathacontline{\textbf{มุนิํ อหํ สกฺก ยถา วิชญฺญํ,}} \\ \csromangathacontline{\textbf{ตํ เม วิยาจิกฺข สมนฺตจกฺขุ.}}}}}

\prose{\textbf{นิราสโส โส อุท อาสสาโน}ติ นิตฺตโณฺห โส, อุทาหุ สตโณฺห รูเป อาสีสติ\csromansharedfootnote{c0edbe5f1b32117b}{อาสิํสติ (สฺยา)}. สทฺเท. คนฺเธ. รเส. โผฏฺฐพฺเพ. กุลํ. คณํ. อาวาสํ. ลาภํ. ยสํ. ปสํสํ. สุขํ. จีวรํ. ปิณฺฑปาตํ. เสนาสนํ. คิลานปจฺจย\-เภสชฺช\-ปริกฺขารํ. กามธาตุํ. รูปธาตุํ. อรูปธาตุํ. กามภวํ. รูปภวํ. อรูปภวํ. สญฺญาภวํ. อสญฺญาภวํ. เนว\-สญฺญา\-นา\-สญฺญาภวํ. เอกโวการภวํ. จตุโวการภวํ. ปญฺจโวการภวํ. อตีตํ. อนาคตํ. ปจฺจุปฺปนฺนํ. ทิฏฺฐ\-สุต\-มุต\-วิญฺญาตพฺเพ ธมฺเม อาสีสติ สาทิยติ ปตฺเถติ ปิเหติ อภิชปฺปตีติ นิราสโส โส อุท อาสสาโน.}

\prose{\textbf{ปญฺญาณวา โส อุท ปญฺญกปฺปี}ติ \textbf{ปญฺญาณวา โส}ติ ปณฺฑิโต ปญฺญวา พุทฺธิมา ญาณี วิภาวี เมธาวี. \textbf{อุท ปญฺญกปฺปี}ติ อุทาหุ อฏฺฐ\-สมาปตฺติ\-ญาเณน วา ปญฺจาภิญฺญา\-ญาเณน วา มิจฺฉาญาเณน วา ตณฺหากปฺปํ วา ทิฏฺฐิกปฺปํ วา กปฺเปติ ชเนติ สญฺชเนติ นิพฺพตฺเตติ อภินิพฺพตฺเตตีติ ปญฺญณวา โส อุท ปญฺญกปฺปี.}

\prose{\textbf{มุนิํ อหํ สกฺก ยถา วิชญฺญนฺ}ติ \textbf{สกฺกา}ติ สกฺโก, ภควา สกฺยกุลา ปพฺพชิโตติปิ สกฺโก. อถ วา อฑฺโฒ มหทฺธโน ธนวาติปิ สกฺโก. ตสฺสิมานิ ธนานิ. เสยฺยถิทํ, สทฺธาธนํ สีลธนํ หิริธนํ โอตฺตปฺปธนํ สุตธนํ จาคธนํ ปญฺญาธนํ สติปฏฺฐานธนํ สมฺมปฺปธานธนํ อิทฺธิปาทธนํ อินฺทฺริยธนํ พลธนํ โพชฺฌงฺคธนํ มคฺคธนํ ผลธนํ นิพฺพานธนนฺติ, เตหิ อเนกวิเธหิ ธนรตเนหิ อฑฺโฒ มหทฺธโน ธนวาติปิ สกฺโก. อถ วา ปหุ วิสวี อลมตฺโต สูโร วีโร วิกฺกนฺโต อภีรู อจฺฉมฺภี อนุตฺราสี อปลายี ปหีน\-ภยเภรโว วิคต\-โลมหํโสติปิ สกฺโก. \textbf{มุนิํ อหํ สกฺก ยถา วิชญฺญนฺ}ติ สกฺก ยถาหํ มุนิํ ชาเนยฺยํ อาชาเนยฺยํ วิชาเนยฺยํ ปฏิวิชาเนยฺยํ ปฏิวิชฺเฌยฺยนฺติ มุนิํ อหํ สกฺก ยถา วิชญฺญํ.}

\csromanheader{2}{9. โตเทยฺยมาณวปุจฺฉา\-นิทฺเทส}
\prose{\csromanfolio{133}\textbf{ตํ เม วิยาจิกฺข สมนฺต}\-\textbf{จกฺขู}ติ \textbf{ตนฺ}ติ ยํ ปุจฺฉามิ, ยํ ยาจามิ, ยํ อชฺเฌสามิ, ยํ ปสาเทมิ. \textbf{วิยาจิกฺขา}ติ อาจิกฺขาหิ เทเสหิ ปญฺญเปหิ ปฏฺฐเปหิ วิวราหิ วิภชาหิ อุตฺตานีกโรหิ ปกาเสหิ. \textbf{สมนฺต}\-\textbf{จกฺขู}ติ สมนฺตจกฺขุ วุจฺจติ สพฺพญฺญุต\-ญาณํ ฯเปฯ ตถาคโต เตน สมนฺต\-จกฺขูติ ตํ เม วิยาจิกฺข สมนฺตจกฺขุ. เตนาห โส พฺราหฺมโณ\csromandash{}}

\setlength{\csromangathapagewidth}{0pt}
\setlength{\csromangathawakwidth}{0pt}
\csromangathameasuregroup{}{นิราสโส โส อุท อาสสาโน, \\ ปญฺญาณวา โส อุท ปญฺญกปฺปี. \\ มุนิํ อหํ สกฺก ยถา วิชญฺญํ,}
\csromangathameasurewak{นิราสโส โส อุท อาสสาโน, \\ ปญฺญาณวา โส อุท ปญฺญกปฺปี. \\ มุนิํ อหํ สกฺก ยถา วิชญฺญํ,}
\setlength{\csromangathaleftcol}{0pt}
\csromangathagroup{\csromangathastanza{\csromangathawak{นิราสโส โส อุท อาสสาโน, \\ ปญฺญาณวา โส อุท ปญฺญกปฺปี. \\ มุนิํ อหํ สกฺก ยถา วิชญฺญํ,}}}

\prose{ตํ เม วิยาจิกฺข สมนฺต\-จกฺขูติ.}

\setlength{\csromangathapagewidth}{0pt}
\setlength{\csromangathawakwidth}{0pt}
\csromangathameasuregroup{}{\textbf{นิราสโส โส น จ อาสสาโน,} \\ \textbf{ปญฺญาณวา โส น จ ปญฺญกปฺปี.} \\ \textbf{เอวมฺปิ โตเทยฺย มุนิํ วิชาน,} \\ \textbf{อกิญฺจนํ กามภเว อสตฺตํ.}}
\csromangathameasurewak{\textbf{นิราสโส โส น จ อาสสาโน,} \\ \textbf{ปญฺญาณวา โส น จ ปญฺญกปฺปี.} \\ \textbf{เอวมฺปิ โตเทยฺย มุนิํ วิชาน,} \\ \textbf{อกิญฺจนํ กามภเว อสตฺตํ.}}
\setlength{\csromangathaleftcol}{0pt}
\csromangathagroup{\csromangathastanza{\csromangathawak{\csromangathaitemline{60}{\textbf{นิราสโส โส น จ อาสสาโน,}} \\ \csromangathacontline{\textbf{ปญฺญาณวา โส น จ ปญฺญกปฺปี.}} \\ \csromangathacontline{\textbf{เอวมฺปิ โตเทยฺย มุนิํ วิชาน,}} \\ \csromangathacontline{\textbf{อกิญฺจนํ กามภเว อสตฺตํ.}}}}}

\prose{\textbf{นิราสโส โส น จ อาสสาโน}ติ นิตฺตโณฺห โส, น โส สตโณฺห รูเป นาสีสติ. สทฺเท. คนฺเธ ฯเปฯ ทิฏฺฐ\-สุต\-มุต\-วิญฺญาตพฺเพ ธมฺเม นาสีสติ น อิจฺฉติ น สาทิยติ น ปตฺเถติ น ปิเหติ นาภิชปฺปตีติ นิราสโส โส น จ อาสสาโน.}

\prose{\textbf{ปญฺญาณวา โส น จ ปญฺญกปฺปี}ติ \textbf{ปญฺญาณวา}ติ ปณฺฑิโต ปญฺญวา พุทฺธิมา ญาณี วิภาวี เมธาวี. \textbf{น จ ปญฺญกปฺปี}ติ อฏฺฐ\-สมาปตฺติ\-ญาเณน วา ปญฺจาภิญฺญา\-ญาเณน วา มิจฺฉาญาเณน วา ตณฺหากปฺปํ วา น กปฺเปติ ทิฏฺฐิกปฺปํ วา น กปฺเปติ น ชเนติ น สญฺชเนติ น นิพฺพตฺเตติ นาภินิพฺพตฺเตตีติ ปญฺญาณวา โส น จ ปญฺญกปฺปี.}

\prose{\textbf{เอวมฺปิ โตเทยฺย มุนิํ วิชานา}ติ \textbf{มุนี}ติ โมนํ วุจฺจติ ญาณํ ฯเปฯ สงฺค\-ชาลม\-ติจฺจ โส มุนิ. \textbf{เอวมฺปิ โตเทยฺย มุนิํ วิชานา}ติ โตเทยฺย เอวํ มุนิํ ชาน ปฏิชาน ปฏิวิชาน ปฏิวิชฺฌาติ เอวมฺปิ โตเทยฺย มุนิํ วิชาน.}

\csromanmatikaanchor{mtk.31}
\csromantocmarkref{subsection}{10. กปฺปมาณวปุจฺฉานิทฺเทส}{mtk.31}
\bookmark[dest={mtk.31},level=2]{10. กปฺปมาณวปุจฺฉานิทฺเทส}
\prose{\textbf{อกิญฺจนํ กามภเว อสตฺตนฺ}ติ \textbf{อกิญฺจนนฺ}ติ ราคกิญฺจนํ โทสกิญฺจนํ โมหกิญฺจนํ มานกิญฺจนํ ทิฏฺฐิกิญฺจนํ กิเลสกิญฺจนํ ทุจฺจริต\-กิญฺจนํ. ยสฺเสตานิ\csromansharedfootnote{7786cdbd2f930042}{ยสฺเสเต (สฺยา)} \csromanfolio{134}กิญฺจนานิ\csromansharedfootnote{da7845f8e7c1d28b}{กิญฺจนา (สฺยา)} ปหีนานิ สมุจฺฉินฺนานิ วูปสนฺตานิ ปฏิปฺปสฺสทฺธานิ อภพฺพุ\-ปฺปตฺติกานิ ญาณคฺคินา ทฑฺฒานิ. โส วุจฺจติ อกิญฺจโน. \textbf{กามา}ติ อุทฺทานโต เทฺว กามา วตฺถุกามา จ กิเลสกามา จ ฯเปฯ. อิเม วุจฺจนฺติ วตฺถุกามา ฯเปฯ. อิเม วุจฺจนฺติ กิเลสกามา. \textbf{ภวา}ติ เทฺว ภวา กมฺมภโว จ ปฏิสนฺธิโก จ ปุนพฺภโว ฯเปฯ อยํ ปฏิสนฺธิโก ปุนพฺภโว.}

\prose{\textbf{อกิญฺจนํ กามภเว อสตฺตนฺ}ติ อกิญฺจนํ ปุคฺคลํ กาเม จ ภเว จ อสตฺตํ อลคฺคํ อลคฺคิตํ อปลิพุทฺธํ นิกฺขนฺตํ นิสฺสฏํ วิปฺปมุตฺตํ วิสญฺญุตฺตํ วิมริยาทิกเตน เจตสา วิหรนฺตนฺติ อกิญฺจนํ กามภเว อสตฺตํ. เตนาห ภควา\csromandash{}}

\setlength{\csromangathapagewidth}{0pt}
\setlength{\csromangathawakwidth}{0pt}
\csromangathameasuregroup{}{นิราสโส โส นา จ อาสสาโน, \\ ปญฺญาณวา โส น จ ปญฺญกปฺปี. \\ เอวมฺปิ โตเทยฺย มุนิํ วิชาน,}
\csromangathameasurewak{นิราสโส โส นา จ อาสสาโน, \\ ปญฺญาณวา โส น จ ปญฺญกปฺปี. \\ เอวมฺปิ โตเทยฺย มุนิํ วิชาน,}
\setlength{\csromangathaleftcol}{0pt}
\csromangathagroup{\csromangathastanza{\csromangathawak{นิราสโส โส นา จ อาสสาโน, \\ ปญฺญาณวา โส น จ ปญฺญกปฺปี. \\ เอวมฺปิ โตเทยฺย มุนิํ วิชาน,}}}

\prose{อกิญฺจนํ กามภเว อสตฺตนฺติ.}

\prose{สห คาถา\-ปริโยสานา ฯเปฯ “สตฺถา เม ภนฺเต ภควา, สาวโก หมสฺมี”ติ.}

\vspace{\tipitakaparskip}
\csromancenter{โตเทยฺยมาณวปุจฺฉา\-นิทฺเทโส นวโม.}
\csromanhangingbegin
\hangingheaditem{61}{\textbf{มชฺเฌ สรสฺมิํ ติฏฺฐตํ, (อิจฺจายสฺมา กปฺโป,)}}
\hangingbody{\textbf{โอเฆ ชาเต มหพฺภเย.}}
\hangingbody{\textbf{ชรามจฺจุปเรตานํ, ทีปํ ปพฺรูหิ มาริส.}}
\hangingbody{\textbf{ตฺวญฺจ เม ทีปมกฺขาหิ, ยถายิทํ นาปรํ สิยา.}}
\csromanhangingend

\prose{\textbf{มชฺเฌ สรสฺมิํ ติฏฺฐตนฺ}ติ สโร วุจฺจติ สํสาโร อาคมนํ คมนํ คมนาคมนํ กาลํ คติ ภวาภโว จุติ จ อุปปตฺติ จ นิพฺพตฺติ จ เภโท จ ชาติ จ ชรา จ มรณญฺจ. สํสารสฺส ปุริมาปิ โกฏิ น ปญฺญายติ, ปจฺฉิมาปิ โกฏิ น ปญฺญายติ. มชฺเฌว สํสาเร สตฺตา ฐิตา ปติฏฺฐิตา อลฺลีนา อุปคตา อชฺโฌสิตา อธิมุตฺตา.}

\prose{\csromanfolio{135}กถํ สํสารสฺส ปุริมา โกฏิ น ปญฺญายติ, เอตฺตกา ชาติโย วฏฺฏํ วตฺติ, ตโต ปรํ น วตฺตตีติ เหวํ นตฺถิ, เอวมฺปิ สํสารสฺส ปุริมา โกฏิ น ปญฺญายติ. เอตฺตกานิ ชาติสตานิ วฏฺฏํ วตฺติ, ตโต ปรํ น วตฺตตีติ เหวํ นตฺถิ, เอวมฺปิ สํสารสฺส ปุริมา โกฏิ น ปญฺญายติ. เอตฺตกานิ ชาติสหสฺสานิ วฏฺฏํ วตฺติ, ตโต ปรํ น วตฺตตีติ เหวํ นตฺถิ, เอวมฺปิ สํสารสฺส ปุริมา โกฏิ น ปญฺญายติ. เอตฺตกานิ ชาติ\-สต\-สหสฺสานิ วฏฺฏํ วตฺติ, ตโต ปรํ น วตฺตตีติ เหวํ นตฺถิ, เอวมฺปิ สํสารสฺส ปุริมา โกติ น ปญฺญายติ. เอตฺตกา ชาติโกฏิโย วฏฺฏํ วตฺติ, ตโต ปรํ น วตฺตตีติ เหวํ นตฺถิ, เอวมฺปิ สํสารสฺส ปุริมา โกฏิ น ปญฺญายติ. เอตฺตกานิ ชาติโกฏิสตานิ วฏฺฏํ วตฺติ, ตโต ปรํ น วตฺตตีติ เหวํ นตฺถิ, เอวมฺปิ สํสารสฺส ปุริมา โกฏิ น ปญฺญายติ. เอตฺตกานิ ชาติ\-โกฏิสหสฺสานิ วฏฺฏํ วตฺติ, ตโต ปรํ น วตฺตตีติ เหวํ นตฺถิ, เอวมฺปิ สํสารสฺส ปุริมา โกฏิ น ปญฺญายติ. เอตฺตกานิ ชาติ\-โกฏิ\-สต\-สหสฺสานิ วฏฺฏํ วตฺติ, ตโต ปรํ น วตฺตตีติ เหวํ นตฺถิ, เอวมฺปิ สํสารสฺส ปุริมา โกฏิ น ปญฺญายติ.}

\prose{เอตฺตกานิ วสฺสานิ วฏฺฏํ วตฺติ, ตโต ปรํ น วตฺตตีติ เหวํ นตฺถิ, เอวมฺปิ สํสารสฺส ปุริมา โกฏิ น ปญฺญายติ. เอตฺตกานิ วสฺสสตานิ วฏฺฏํ วตฺติ, ตโต ปรํ น วตฺตตีติ เหวํ นตฺถิ, เอวมฺปิ สํสารสฺส ปุริมา โกฏิ น ปญฺญายติ. เอตฺตกานิ วสฺสสหสฺสานิ วฏฺฏํ วตฺติ, ตโต ปรํ น วตฺตตีติ เหวํ นตฺถิ, เอวมฺปิ สํสารสฺส ปุริมา โกฏิ น ปญฺญายติ. เอตฺตกานิ วสฺส\-สตสหสฺสานิ วฏฺฏํ วตฺติ, ตโต ปรํ น วตฺตตีติ เหวํ นตฺถิ, เอวมฺปิ สํสารสฺส ปุริมา โกฏิ น ปญฺญายติ. เอตฺตกา วสฺสโกฏิโย วฏฺฏํ วตฺติ, ตโต ปรํ น วตฺตตีติ เหวํ นตฺถิ, เอวมฺปิ สํสารสฺส ปุริมา โกฏิ น ปญฺญายติ. เอตฺตกานิ วสฺส\-โกฏิสตานิ วฏฺฏํ วตฺติ, ตโต ปรํ น วตฺตตีติ เหวํ นตฺถิ, เอวมฺปิ สํสารสฺส ปุริมา โกฏิ น ปญฺญายติ. เอตฺตกานิ วสฺส\-โกฏิสหสฺสานิ วฏฺฏํ วตฺติ, ตโต ปรํ น วตฺตตีติ เหวํ นตฺถิ, เอวมฺปิ สํสารสฺส ปุริมา โกฏิ น ปญฺญายติ. เอตฺตกานิ วสฺส\-โกฏิ\-สต\-สหสฺสานิ วฏฺฏํ วตฺติ, ตโต ปรํ น วตฺตตีติ เหวํ นตฺถิ, เอวมฺปิ สํสารสฺส ปุริมา โกฏิ น ปญฺญายติ.}

\prose{เอตฺตกานิ กปฺปานิ วฏฺฏํ วตฺติ, ตโต ปรํ น วตฺตตีติ เหวํ นตฺถิ, เอวมฺปิ สํสารสฺส ปุริมา โกฏิ น ปญฺญายติ. เอตฺตกานิ กปฺปสตานิ}

\prose{\csromanfolio{136}วฏฺฏํ วตฺติ, ตโต ปรํ น วตฺตตีติ เหวํ นตฺถิ, เอวมฺปิ สํสารสฺส ปุริมา โกฏิ น ปญฺญายติ. เอตฺตกานิ กปฺปสหสฺสานิ วฏฺฏํ วตฺติ, ตโต ปรํ น วตฺตตีติ เหวํ นตฺถิ, เอวมฺปิ สํสารสฺส ปุริมา โกฏิ น ปญฺญายติ. เอตฺตกานิ กปฺป\-สตสหสฺสานิ วฏฺฏํ วตฺติ, ตโต ปรํ น วตฺตตีติ เหวํ นตฺถิ, เอวมฺปิ สํสารสฺส ปุริมา โกฏิ น ปญฺญายติ. เอตฺตกา กปฺปโกฏิโย วฏฺฏํ วตฺติ, ตโต ปรํ น วตฺตตีติ เหวํ นตฺถิ, เอวมฺปิ สํสารสฺส ปุริมา โกฏิ น ปญฺญายติ. เอตฺตกานิ กปฺป\-โกฏิสตานิ วฏฺฏํ วตฺติ, ตโต ปรํ น วตฺตตีติ เหวํ นตฺถิ, เอวมฺปิ สํสารสฺส ปุริมา โกฏิ น ปญฺญายติ. เอตฺตกานิ กปฺป\-โกฏิสหสฺสานิ วฏฺฏํ วตฺติ, ตโต ปรํ น วตฺตตีติ เหวํ นตฺถิ. เอวมฺปิ สํสารสฺส ปุริมา โกฏิ น ปญฺญายติ. เอตฺตกานิ กปฺป\-โกฏิ\-สต\-สหสฺสานิ วฏฺฏํ วตฺติ, ตโต ปรํ น วตฺตตีติ เหวํ นตฺถิ, เอวมฺปิ สํสารสฺส ปุริมา โกฏิ น ปญฺญายติ.}

\prose{วุตฺตํ เหตํ ภควตา\csromandash{}อนมตคฺโคยํ ภิกฺขเว สํสาโร ปุพฺพา โกฏิ น ปญฺญายติ อวิชฺชา\-นีวรณานํ สตฺตานํ ตณฺหา\-สํโยชนานํ สนฺธาวตํ สํสรตํ. เอวํ ทีฆรตฺตํ โข ภิกฺขเว ทุกฺขํ ปจฺจนุภูตํ ติพฺพํ ปจฺจนุภูตํ พฺยสนํ ปจฺจนุภูตํ กฏสี วฑฺฒิตา\csromansharedfootnote{ce15c0a25adc2079}{กฏสีววฑฺฒิตํ (สฺยา) ปสฺส สํ 1. 387 ปิฏฺเฐ.}, ยาวญฺจิทํ ภิกฺขเว อลเมว สพฺพ\-สงฺขาเรสุ นิพฺพินฺทิตุํ, อลํ วิรชฺชิตุํ, อลํ วิมุจฺจิตุนฺติ เอวมฺปิ สํสารสฺส ปุริมา โกฏิ น ปญฺญายติ.}

\prose{กถํ สํสารสฺส ปจฺฉิมา โกฏิ น ปญฺญายติ. เอตฺตกา ชาติโย วฏฺฏํ วตฺติสฺสติ, ตโต ปรํ น วตฺติสฺสตีติ เหวํ นตฺถิ, เอวมฺปิ สํสารสฺส ปจฺฉิมา โกฏิ น ปญฺญายติ. เอตฺตกานิ ชาติสตานิ. เอตฺตกานิ ชาติสหสฺสานิ. เอตฺตกานิ ชาติ\-สต\-สหสฺสานิ ฯเปฯ. เอตฺตกา ชาติโกฏิโย. เอตฺตกานิ ชาติโกฏิสตานิ. เอตฺตกานิ ชาติ\-โกฏิสหสฺสานิ. เอตฺตกานิ ชาติ\-โกฏิ\-สต\-สหสฺสานิ. เอตฺตกานิ วสฺสานิ. เอตฺตกานิ วสฺสาสตานิ. เอตฺตกานิ วสฺสสหสฺสานิ. เอตฺตกานิ วสฺส\-สตสหสฺสานิ. เอตฺตกา วสฺสโกฏิโย. เอตฺตกานิ วสฺส\-โกฏิสตานิ. เอตฺตกานิ วสฺส\-โกฏิสหสฺสานิ. เอตฺตกานิ วสฺส\-โกฏิ\-สต\-สหสฺสานิ. เอตฺตกานิ กปฺปานิ. เอตฺตกานิ กปฺปสตานิ. เอตฺตกานิ กปฺปสหสฺสานิ. เอตฺตกานิ กปฺป\-สตสหสฺสานิ. เอตฺตกา กปฺปโกฏิโย. เอตฺตกานิ \csromanfolio{137}กปฺป\-โกฏิสตานิ. เอตฺตกานิ กปฺป\-โกฏิสหสฺสานิ. เอตฺตกานิ กปฺป\-โกฏิ\-สต\-สหสฺสานิ วฏฺฏํ วตฺติสฺสติ. ตโต ปรํ น วตฺติสฺสตีติ เหวํ นตฺถิ, เอวมฺปิ สํสารสฺส ปจฺฉิมา โกฏิ น ปญฺญายติ. เอวมฺปิ สํสารสฺส ปุริมาปิ โกฏิ น ปญฺญายติ, ปจฺฉิมาปิ โกฏิ น ปญฺญายติ. มชฺเฌว สํสาเร สตฺตา ฐิตา ปติฏฺฐิตา อลฺลีนา อุปคตา อชฺโฌสิตา อธิมุตฺตาติ มชฺเฌ สรสฺมิํ ติฏฺฐตํ. \textbf{อิจฺจายสฺมา กปฺโป}ติ \textbf{อิจฺจา}ติ ปทสนฺธิ ฯเปฯ. \textbf{อายสฺมา}ติ ปิยวจนํ ฯเปฯ. \textbf{กปฺโป}ติ ตสฺส พฺราหฺมณสฺส นามํ ฯเปฯ อภิลาโปติ อิจฺจายสฺมา กปฺโป.}

\proseitem{61}{\textbf{โอเฆ ชาเต มหพฺภเย}ติ กาโมเฆ ภโวเฆ ทิฏฺโฐเฆ อวิชฺโชเฆ ชาเต สญฺชาเต นิพฺพตฺเต อภินิพฺพตฺเต ปาตุภูเต. \textbf{มหพฺภเย}ติ ชาติภเย ชราภเย พฺยาธิภเย มรณภเยติ โอเฆ ชาเต มหพฺภเย.}

\prose{\textbf{ชรามจฺจุ}\-\textbf{ปเรตานนฺ}ติ ชราย ผุฏฺฐานํ ปเรตานํ สโมหิตานํ สมนฺนาคตานํ. มจฺจุนา ผุฏฺฐานํ ปเรตานํ สโมหิตานํ สมนฺนาคตานํ, ชาติยา อนุคตานํ ชราย อนุสฏานํ พฺยาธินา อภิภูตานํ มรเณน อพฺภาหตานํ อตาณานํ อเลณานํ อสรณานํ อสรณีภูตานนฺติ ชรามจฺจุ\-ปเรตานํ.}

\prose{\textbf{ทีปํ ปพฺรูหิ มาริสา}ติ ทีปํ ตาณํ เลณํ สรณํ คติํ ปรายนํ\csromansharedfootnote{9aa6c70ca921ee33}{คติปรายนํ (สฺยา) เอวมุปริปิ.} พฺรูหิ อาจิกฺขาหิ เทเสหิ ปญฺญเปหิ ปฏฺฐเปหิ วิวราหิ วิภชาหิ อุตฺตานีกโรหิ ปกาเสหิ. \textbf{มาริสา}ติ ปิยวจนํ ครุวจนํ สคารวสปฺปติสฺสาธิวจนเมตํ มาริสาติ ทีปํ ปพฺรูหิ มาริส.}

\prose{\textbf{ตฺวญฺจ เม ทีปมกฺขาหี}ติ \textbf{ตฺวนฺ}ติ ภควนฺตํ ภณติ. \textbf{ทีปมกฺขาหี}ติ ทีปํ ตาณํ เลณํ สรณํ คติํ ปรายนํ อกฺขาหิ อาจิกฺขาหิ เทเสหิ ปญฺญเปหิ ปฏฺฐเปหิ วิวราหิ วิภชาหิ อุตฺตานีกโรหิ ปกาเสหีติ ตฺวญฺจ เม ทีปมกฺขาหิ.}

\prose{\textbf{ยถายิทํ นาปรํ สิยา}ติ ยถยิทํ ทุกฺขํ อิเธว นิรุชฺเฌยฺย วูปสเมยฺย อตฺถํ คจฺเฉยฺย ปฏิปฺปสฺสมฺเภยฺย, ปุนปฏิสนฺธิกํ ทุกฺขํ น นิพฺพตฺเตยฺย, กามธาตุยา วา รูปธาตุยา วา อรูปธาตุยา วา กามภเว วา รูปภเว \csromanfolio{138}วา อรูปภเว วา สญฺญาภเว วา อสญฺญาภเว วา เนว\-สญฺญา\-นา\-สญฺญาภเว วา เอกโวการภเว วา จตุโวการภเว วา ปญฺจโวการภเว วา ปุนคติยา วา อุปปตฺติยา วา ปฏิสนฺธิยา วา ภเว วา สํสาเร วา วฏฺเฏ วา น ชเนยฺย น สญฺชเนยฺย น นิพฺพตฺเตยฺย นาภินิพฺพตฺเตยฺย อิเธว นิรุชฺเฌยฺย วูปสเมยฺย อตฺถํ คจฺเฉยฺย ปฏิปฺปสฺสมฺเภยฺยาติ ยถายิทํ นาปรํ \textbf{สิยา. เตนาห โส พฺราหฺมโณ\csromandash{}}}

\proseitem{10}{มชฺเฌ สรสฺมิํ ติฏฺฐตํ, (อิจฺจายสฺมา กปฺโป,)}

\prose{โอเฆ ชาเต มหพฺภเย.}

\setlength{\csromangathapagewidth}{0pt}
\setlength{\csromangathawakwidth}{0pt}
\csromangathameasuregroup{ชรามจฺจุปเรตานํ, \\ ตฺวญฺจ เม ทีปมกฺขาหิ,}{\csromangathabat{ชรามจฺจุปเรตานํ,}{ทีปํ ปพฺรูหิ มาริส.} \\ \csromangathabat{ตฺวญฺจ เม ทีปมกฺขาหิ,}{ยถายิทํ นาปรํ สิยาติ.}}
\csromangathasetleft{ชรามจฺจุปเรตานํ, \\ ตฺวญฺจ เม ทีปมกฺขาหิ,}
\csromangathagroup{\csromangathastanza{\csromangathabat{ชรามจฺจุ\-ปเรตานํ,}{ทีปํ ปพฺรูหิ มาริส.} \\ \csromangathabat{ตฺวญฺจ เม ทีปมกฺขาหิ,}{ยถายิทํ นาปรํ สิยาติ.}}}

\csromanhangingbegin
\hangingheaditem{62}{\textbf{มชฺเฌ สรสฺมิํ ติฏฺฐตํ, (กปฺปาติ ภควา,)}}
\hangingbody{โอเฆ ชาเต มหพฺภเย.}
\hangingbody{\textbf{ชรามจฺจุปเรตานํ, ทีปํ ปพฺรูมิ กปฺป เต.}}
\csromanhangingend

\prose{\textbf{มชฺเฌ สรสฺมิํ ติฏฺฐตนฺ}ติ สโร วุจฺจติ สํสาโร อาคมนํ คมนํ คมนาคมนํ กาลํ คติ ภวาภโว จุติ จ อุปปตฺติ จ นิพฺพตฺติ จ เภโท จ ชาติ จ ชรา จ มรณญฺจ. สํสารสฺส ปุริมาปิ โกฏิ น ปญฺญายติ, ปจฺฉิมาปิ โกฏิ น ปญฺญายติ. มชฺเฌว สํสาเร สตฺตา ฐิตา ปติฏฺฐิตา อลฺลีนา อุปคตา อชฺโฌสิตา อธิมุตฺตา.}

\prose{กถํ สํสารสฺส ปุริมา โกฏิ น ปญฺญายติ ฯเปฯ. เอวํ สํสารสฺส ปุริมา โกฏิ น ปญฺญายติ. กถํ สํสารสฺส ปจฺฉิมา โกฏิ น ปญฺญายติ ฯเปฯ. เอวํ สํสารสฺส ปจฺฉิมา โกฏิ น ปญฺญายติ. เอวํ สํสารสฺส ปุริมาปิ โกฏิ น ปญฺญายติ, ปจฺฉิมาปิ โกฏิ น ปญฺญายติ. มชฺเฌว สํสาเร สตฺตา ฐิตา ปติฏฺฐิตา อลฺลีนา อุปคตา อชฺโฌสิตา อธิมุตฺตาติ มชฺเฌ สรสฺมิํ ติฏฺฐตํ. \textbf{กปฺปาติ ภควา}ติ \textbf{กปฺปา}ติ ภควา ตํ พฺราหฺมณํ นาเมน อาลปติ. \textbf{ภควา}ติ คารวาธิวจนเมตํ ฯเปฯ สจฺฉิกา ปญฺญตฺติ, ยทิทํ ภควาติ กปฺปาติ ภควา.}

\prose{\textbf{โอเฆ ชาเต มหพฺภเย}ติ กาโมเฆ ภโวเฆ ทิฏฺโฐเฆ อวิชฺโชเฆ ชาเต สญฺชาเต นิพฺพตฺเต อภินิพฺพตฺเต ปาตุภูเต. \textbf{มหพฺภเย}ติ ชาติภเย ชราภเย พฺยาธิภเย มรณภเยติ โอเฆ ชาเต มหพฺภเย.}

\prose{\csromanfolio{139}\textbf{ชรามจฺจุ}\-\textbf{ปเรตานนฺ}ติ ชราย ผุฏฺฐานํ ปเรตานํ สโมหิตานํ สมนฺนาคตานํ. มจฺจุนา ผุฏฺฐานํ ปเรตานํ สโมหิตานํ สมนฺนาคตานํ ชาติยา อนุคตานํ ชราย อนุสฏานํ พฺยาธินา อภิภูตานํ มรเณน อพฺภาหตานํ อตาณานํ อเลณานํ อสรณานํ อสรณีภูตานนฺติ ชรามจฺจุ\-ปเรตานํ.}

\prose{\textbf{ทีปํ ปพฺรูมิ กปฺป เต}ติ ทีปํ ตาณํ เลณํ สรณํ คติํ ปรายนํ พฺรูมิ อาจิกฺขามิ เทเสมิ ปญฺญเปมิ ปฏฺฐเปมิ วิวรามิ วิภชามิ อุตฺตานีกโรมิ ปกาเสมีติ ทีปํ ปพฺรูมิ กปฺป เต. เตนาห ภควา\csromandash{}}

\prose{มชฺเฌ สรสฺมิํ ติฏฺฐตํ, (กปฺปาติ ภควา,)}

\prose{โอเฆ ชาเต มหพฺภเย.}

\setlength{\csromangathapagewidth}{0pt}
\setlength{\csromangathawakwidth}{0pt}
\csromangathameasuregroup{ชรามจฺจุปเรตานํ, \\ \textbf{อกิญฺจนํ อนาทานฺ}อํ, \\ \textbf{นิพฺพานํ อิติ นํ พฺ}รูมิ,}{\csromangathabat{ชรามจฺจุปเรตานํ,}{ทีปํ ปพฺรูมิ กปฺป เตติ.} \\ \csromangathabat{\textbf{อกิญฺจนํ อนาทานฺ}อํ,}{เอตํ ทีปํ อนาปรํ.} \\ \csromangathabat{\textbf{นิพฺพานํ อิติ นํ พฺ}รูมิ,}{ชรามจฺจุปริกฺขยํ.}}
\csromangathasetleft{ชรามจฺจุปเรตานํ, \\ \textbf{อกิญฺจนํ อนาทานฺ}อํ, \\ \textbf{นิพฺพานํ อิติ นํ พฺ}รูมิ,}
\csromangathagroup{\csromangathastanza{\csromangathabat{ชรามจฺจุ\-ปเรตานํ,}{ทีปํ ปพฺรูมิ กปฺป เตติ.}}\csromangathastanza{\csromangathaitembat{63}{\textbf{อกิญฺจนํ อนาทานฺ}อํ,}{เอตํ ทีปํ อนาปรํ.} \\ \csromangathacontbat{\textbf{นิพฺพานํ อิติ นํ พฺ}รูมิ,}{ชรามจฺจุ\-ปริกฺขยํ.}}}

\prose{\textbf{อกิญฺจนํ อนาทานนฺ}ติ \textbf{กิญฺจนนฺ}ติ ราคกิญฺจนํ โทสกิญฺจนํ โมหกิญฺจนํ มานกิญฺจนํ ทิฏฺฐิกิญฺจนํ กิเลสกิญฺจนํ ทุจฺจริต\-กิญฺจนํ, กิญฺจน\-ปฺปหานํ กิญฺจน\-วูปสมํ\csromansharedfootnote{90c6ff24d5106dc6}{กิญฺจนวูปสโม (สฺยา) เอวมีทิเสสุ ฐาเนสุ.} กิญฺจน\-ปฏินิสฺสคฺคํ\csromansharedfootnote{bde69735c292366c}{กิญฺจนปฏินิสฺสคฺโค (สฺยา)} กิญฺจน\-ปฏิปฺปสฺสทฺธิํ อมตํ นิพฺพานนฺติ อกิญฺจนํ. \textbf{อนาทานนฺ}ติ \textbf{อาทานํ} วุจฺจติ ตณฺหา. โย ราโค สาราโค ฯเปฯ อภิชฺฌา โลโภ อกุสลมูลํ, อาทานปฺปหานํ อาทานวูปสมํ อาทาน\-ปฏินิสฺสคฺคํ อาทาน\-ปฏิปฺปสฺสทฺธิํ อมตํ นิพฺพานนฺติ อกิญฺจนํ อนาทานํ.}

\prose{\textbf{เอตํ ทีปํ อนาปรนฺ}ติ เอตํ ทีปํ ตาณํ เลณํ สรณํ คติ ปรายนํ. \textbf{อนาปรนฺ}ติ ตมฺหา ปโร อญฺโญ ทีโป นตฺถิ, อถ โข โส เอวํ ทีโป อคฺโค จ เสฏฺโฐ จ วิเสฏฺโฐ จ ปาโมกฺโข จ อุตฺตโม จ ปวโร จาติ เอตํ ทีปํ อนาปรํ.}

\prose{\textbf{นิพฺพานํ อิติ นํ พฺรูมี}ติ \textbf{วานํ} วุจฺจติ ตณฺหา. โย ราโค สาราโค ฯเปฯ อภิชฺฌา โลโภ อกุสลมูลํ, วานปฺปหานํ วานวูปสมํ วาน\-ปฏินิสฺสคฺคํ วาน\-ปฏิปฺปสฺสทฺธิํ อมตํ นิพฺพานํ. \textbf{อิตี}ติ ปทสนฺธิ ปทสํสคฺโค ปทปาริปูรี อกฺขร\-สมวาโย พฺยญฺชน\-สิลิฏฺฐตา ปทา\-นุปุพฺพตา\-เปตํ อิตีติ. \csromanfolio{140}\textbf{พฺรูมี}ติ พฺรูมิ อาจิกฺขามิ เทเสมิ ปญฺญเปมิ ปฏฺฐเปมิ วิวรามิ วิภชามิ \textbf{อุตฺตานีกโรมิ ปกาเสมีติ นิพฺพานํ อิติ นํ พฺรูมิ.}}

\prose{\textbf{ชรามจฺจุ}\-\textbf{ปริกฺขยนฺ}ติ ชรามรณสฺส ปหานํ วูปสมํ ปฏินิสฺสคฺคํ ปฏิปฺปสฺสทฺธิํ อมตํ นิพฺพานนฺติ ชรามจฺจุ\-ปริกฺขยํ. เตนาห ภควา\csromandash{}}

\setlength{\csromangathapagewidth}{0pt}
\setlength{\csromangathawakwidth}{0pt}
\csromangathameasuregroup{\textbf{เอตทญฺญาย เย สตฺ}อา, \\ \textbf{น เต มารวสานุคฺ}อา,}{อกิญฺจนํ อนาทานํ, เอตํ ทีปํ อนาปรํ, \\ นิพฺพานํ อิติ นํ พฺรูมิ, \\ ชรามจฺจุปริกฺขยนฺติ. \\ \csromangathabat{\textbf{เอตทญฺญาย เย สตฺ}อา,}{\textbf{ทิฏฺฐธมฺมา}ภินิพฺพุตา.} \\ \csromangathabat{\textbf{น เต มารวสานุคฺ}อา,}{น เต มารสฺส ปทฺธคู\textsuperscript{1}.}}
\csromangathameasurewak{อกิญฺจนํ อนาทานํ, เอตํ ทีปํ อนาปรํ, \\ นิพฺพานํ อิติ นํ พฺรูมิ, \\ ชรามจฺจุปริกฺขยนฺติ.}
\csromangathasetleft{\textbf{เอตทญฺญาย เย สตฺ}อา, \\ \textbf{น เต มารวสานุคฺ}อา,}
\csromangathagroup{\csromangathastanza{\csromangathawak{อกิญฺจนํ อนาทานํ, เอตํ ทีปํ อนาปรํ, \\ นิพฺพานํ อิติ นํ พฺรูมิ, \\ ชรามจฺจุ\-ปริกฺขยนฺติ.}}\csromangathastanza{\csromangathaitembat{64}{\textbf{เอตทญฺญาย เย สตฺ}อา,}{\textbf{ทิฏฺฐธมฺมา}\-ภินิพฺพุตา.} \\ \csromangathacontbat{\textbf{น เต มารวสานุคฺ}อา,}{น เต มารสฺส ปทฺธคู\csromansharedfootnote{32b0771cb3b81575}{ปฏฺฐคู (สฺยา, ก)}.}}}

\prose{\textbf{เอตทญฺญาย เย สตา}ติ \textbf{เอตนฺ}ติ อมตํ นิพฺพานํ. โย โส สพฺพ\-สงฺขาร\-สมโถ สพฺพู\-ปธิ\-ปฏินิสฺสคฺโค ตณฺหกฺขโย วิราโค นิโรโธ นิพฺพานํ. \textbf{อญฺญายา}ติ อญฺญาย ชานิตฺวา ตุลยิตฺวา ตีรยิตฺวา วิภาวยิตฺวา วิภูตํ กตฺวา, “สพฺเพ สงฺขารา อนิจฺจา”ติ ฯเปฯ. “ยํ กิญฺจิ สมุทย\-ธมฺมํ, สพฺพํ ตํ นิโรธธมฺมนฺ”ติ อญฺญาย ชานิตฺวา ตุลยิตฺวา ตีรยิตฺวา วิภาวยิตฺวา วิภูตํ กตฺวา. \textbf{เย}ติ อรหนฺโต ขีณาสวา. \textbf{สตา}ติ จตูหิ การเณหิ สตา, กาเย กายานุปสฺสานาสติปฏฺฐานํ ภาเวนฺตา\csromansharedfootnote{821631925b479f78}{ภาวิตตฺตา (ก)} สตา ฯเปฯ เต วุจฺจนฺติ สตาติ เอตทญฺญาย เย สตา.}

\prose{\textbf{ทิฏฺฐธมฺมา}\-\textbf{ภินิพฺพุตา}ติ \textbf{ทิฏฺฐธมฺมา}ติ ทิฏฺฐธมฺมา ญาตธมฺมา ตุลิตธมฺมา ตีริตธมฺมา วิภูตธมฺมา วิภาวิต\-ธมฺมา. \textbf{อภินิพฺพุตา}ติ ราคสฺส นิพฺพาปิตตฺตา นิพฺพุตา, โทสสฺส ฯเปฯ สพฺพา\-กุสลา\-ภิสงฺขารานํ สนฺตตฺตา สมิตตฺตา วูปสมิตตฺตา นิชฺฌาตตฺตา นิพฺพุตตฺตา ปฏิปฺปสฺสทฺธ\-ตฺตา สนฺตา อุปสนฺตา วูปสนฺตา นิพฺพุตา ปฏิปฺปสฺสทฺธาติ ทิฏฺฐธมฺมา\-ภินิพฺพุตา.}

\prose{\textbf{น เต มารวสานุคา}ติ \textbf{มาโร}ติ โย โส มาโร กโณฺห อธิปติ อนฺตคู นมุจิ ปมตฺตพนฺธุ. \textbf{น เต มารวสานุคา}ติ น เต มารสฺส วเส วตฺตนฺติ, นาปิ มาโร เตสุ วสํ วตฺเตติ, เต มารญฺจ มารปกฺขญฺจ มารปาสญฺจ มารพฬิสญฺจ\csromansharedfootnote{19a88e009660bbf3}{มารพลิสญฺจ (ก)} มารามิสญฺจ มารวิสยญฺจ มารนิวาสญฺจ มารโคจรญฺจ มารพนฺธนญฺจ อภิภุยฺย อภิภวิตฺวา อชฺโฌตฺถริตฺวา ปริยาทิยิตฺวา มทฺทิตฺวา จรนฺติ วิหรนฺติ อิริยนฺติ วตฺเตนฺติ ปาเลนฺติ ยเปนฺติ ยาเปนฺตีติ น เต มารวสานุคา.}

\csromanmatikaanchor{mtk.32}
\csromantocmarkref{subsection}{11. ชตุกณฺณิมาณวปุจฺฉานิทฺเทส}{mtk.32}
\bookmark[dest={mtk.32},level=2]{11. ชตุกณฺณิมาณวปุจฺฉานิทฺเทส}
\prose{\csromanfolio{141}\textbf{น เต มารสฺส ปทฺธคู}ติ น เต มารสฺส ปทฺธา ปทฺธจรา\csromansharedfootnote{9e3a70abc383f5da}{ปฏฺฐา ปฏฺฐจรา (สฺยา, ก)} ปริจาริกา สิยา, พุทฺธสฺส เต ภควโต ปทฺธา ปทฺธจรา ปริจาริกา สิยาติ น เต มารสฺส ปทฺธคู. เตนาห ภควา\csromandash{}}

\setlength{\csromangathapagewidth}{0pt}
\setlength{\csromangathawakwidth}{0pt}
\csromangathameasuregroup{เอตทญฺญาย เย สตา, \\ น เต มารวสานุคา,}{\csromangathabat{เอตทญฺญาย เย สตา,}{ทิฏฺฐธมฺมาภินิพฺพุตา.} \\ \csromangathabat{น เต มารวสานุคา,}{น เต มารสฺส ปทฺธคูติ.}}
\csromangathasetleft{เอตทญฺญาย เย สตา, \\ น เต มารวสานุคา,}
\csromangathagroup{\csromangathastanza{\csromangathabat{เอตทญฺญาย เย สตา,}{ทิฏฺฐธมฺมา\-ภินิพฺพุตา.} \\ \csromangathabat{น เต มารวสานุคา,}{น เต มารสฺส ปทฺธคูติ.}}}

\prose{สห คาถา\-ปริโยสานา ฯเปฯ “สตฺถา เม ภนฺเต ภควา, สาวโกหมสฺมี”ติ.}

\vspace{\tipitakaparskip}
\csromancenter{กปฺปมาณวปุจฺฉา\-นิทฺเทโส ทสโม.}
\csromanhangingbegin
\hangingheaditem{65}{\textbf{สุตฺวานหํ วีร อกามกามิํ, (อิจฺจายสฺมา ชตุกณฺณิ,)}}
\hangingbody{\textbf{โอฆาติคํ ปุฏฺฐุมกามมาคมํ.}}
\hangingbody{\textbf{สนฺติปทํ พฺรูหิ สหชเนตฺต,}}
\hangingbody{\textbf{ยถาตจฺฉํ ภควา พฺรูหิ เม’ตํ.}}
\csromanhangingend

\prose{\textbf{สุตฺวานหํ วีร อกามกามินฺ}ติ สุตฺวา สุณิตฺวา อุคฺคเหตฺวา อุปธาเรตฺวา อุปลกฺขยิตฺวา “อิติปิ โส ภควา อรหํ ฯเปฯ พุทฺโธ ภควา”ติ สุตฺวานหํ. \textbf{วีรนฺ}ติ วีโร, ภควา วีริยวาติ วีโร, ปหูติ วีโร, วิสวีติ วีโร, อลมตฺโตติ วีโร, สูโรติ วีโร, วิกฺกนฺโต อภีรู อจฺฉมฺภี อนุตฺราสี อปลายี ปหีน\-ภยเภรโว วิคต\-โลมหํโสติ วีโร.}

\setlength{\csromangathapagewidth}{0pt}
\setlength{\csromangathawakwidth}{0pt}
\csromangathameasuregroup{}{วิรโต อิธ สพฺพปาปเกหิ, \\ นิรยทุกฺขํ อติจฺจ วีริยวาโส\textsuperscript{1}. \\ โส วีริยวา ปธานวา, \\ วีโร ตาทิ ปวุจฺจเต ตถตฺตาติ.}
\csromangathameasurewak{วิรโต อิธ สพฺพปาปเกหิ, \\ นิรยทุกฺขํ อติจฺจ วีริยวาโส\textsuperscript{1}. \\ โส วีริยวา ปธานวา, \\ วีโร ตาทิ ปวุจฺจเต ตถตฺตาติ.}
\setlength{\csromangathaleftcol}{0pt}
\csromangathagroup{\csromangathastanza{\csromangathawak{วิรโต อิธ สพฺพปาปเกหิ, \\ นิรยทุกฺขํ อติจฺจ วีริยวาโส\csromansharedfootnote{2470a4c8af1b6536}{วิริยวา (สฺยา) ขุ 1. 360 ปิฏฺเฐปิ.}. \\ โส วีริยวา ปธานวา, \\ วีโร ตาทิ ปวุจฺจเต ตถตฺตาติ.}}}

\prose{สุตฺวานหํ วีร. \textbf{อกามกามินฺ}ติ กามาติ อุทฺทานโต เทฺว กามา วตฺถุกามา จ กิเลสกามา จ ฯเปฯ. อิเม วุจฺจนฺติ วตฺถุกามา ฯเปฯ. อิเม วุจฺจนฺติ กิเลสกามา. พุทฺธสฺส ภควโต วตฺถุกามา ปริญฺญาตา, กิเลสกามา ปหีนา, วตฺถุกามานํ ปริญฺญาตตฺตา กิเลสกามานํ \csromanfolio{142}ปหีนตฺตา ภควา น กาเม กาเมติ, น กาเม ปตฺเถติ, น กาเม ปิเหติ, น กาเม อภิชปฺปติ. เย กาเม กาเมนฺติ, กาเม ปตฺเถนฺติ, กาเม ปิเหนฺติ, กาเม อภิชปฺปนฺติ, เต กามกามิโน ราคราคิโน สญฺญาสญฺญิโน, ภควา น กาเม กาเมติ, น กาเม ปตฺเถติ, น กาเม ปิเหติ, น กาเม อภิชปฺปติ, ตสฺมา พุทฺโธ อกาโม นิกฺกาโม จตฺตกาโม วนฺตกาโม มุตฺตกาโม ปหีนกาโม ปฏินิสฺสฏฺฐ\-กาโม วีตราโค วิคตราโค จตฺตราโค วนฺตราโค มุตฺตราโค ปหีนราโค ปฏินิสฺสฏฺฐ\-ราโค นิจฺฉาโต นิพฺพุโต สีติภูโต สุขปฺปฏิสํเวที \textbf{พฺรหฺมภูเตน อตฺตนา วิหรตีติ สุตฺวานหํ วีร อกามกามิํ.}}

\prose{\textbf{อิจฺจายสฺมา ชตุกณฺณี}ติ \textbf{อิจฺจา}ติ ปทสนฺธิ ฯเปฯ ปทา\-นุปุพฺพตา\-เปตํ อิจฺจาติ. \textbf{อายสฺมา}ติ ปิยวจนํ สคารวสปฺปติสฺสาธิวจนเมตํ อายสฺมาติ. \textbf{ชตุกณฺณี}ติ ตสฺส พฺราหฺมณสฺส โคตฺตํ สงฺขา สมญฺญา ปญฺญตฺติ โวหาโรติ อิจฺจายสฺมา ชตุกณฺณิ.}

\prose{\textbf{โอฆาติคํ ปุฏฺฐุม}\-\textbf{กามมา}\-\textbf{คมนฺ}ติ \textbf{โอฆาติคนฺ}ติ โอฆาติคํ โอฆํ อติกฺกนฺตํ สมติกฺกนฺตํ วีติวตฺตนฺติ โอฆาติคํ. \textbf{ปุฏฺฐุ}นฺติ ปุฏฺฐุํ ปุจฺฉิตุํ ยาจิตุํ อชฺเฌสิตุํ ปสาเทตุํ. \textbf{อกามมาคมนฺ}ติ อกามํ ปุฏฺฐุํ นิกฺกามํ จตฺตกามํ วนฺตกามํ มุตฺตกามํ ปหีนกามํ ปฏินิสฺสฏฺฐ\-กามํ วีตราคํ วิคตราคํ จตฺตราคํ วนฺตราคํ มุตฺตราคํ ปหีนราคํ ปฏินิสฺสฏฺฐ\-ราคํ อาคมฺหา อาคตมฺหา อุปาคตมฺหา สมฺปตฺตมฺหา ตยา สทฺธิํ สมาคตมฺหาติ โอฆาติคํ ปุฏฺฐุม\-กามมาคมํ.}

\prose{\textbf{สนฺติปทํ พฺรูหิ สหชเนตฺตา}ติ \textbf{สนฺตี}ติ เอเกน อากาเรน สนฺติปิ สนฺติปทมฺปิ\csromansharedfootnote{72be4fd91ea73178}{สนฺติปทนฺติ (ก)} ตํเยว อมตํ นิพฺพานํ. โย โส สพฺพ\-สงฺขาร\-สมโถ สพฺพู\-ปธิ\-ปฏินิสฺสคฺโค ตณฺหกฺขโย วิราโค นิโรโธ นิพฺพานํ. วุตฺตํ เหตํ ภควตา\csromandash{}สนฺตเมตํ ปทํ, ปณีตเมตํ ปทํ, ยทิทํ สพฺพ\-สงฺขาร\-สมโถ สพฺพู\-ปธิ\-ปฏินิสฺสคฺโค ตณฺหกฺขโย วิราโค นิโรโธ นิพฺพานนฺติ. อถาปเรนา\-กาเรน เย ธมฺมา สนฺตาธิคมาย สนฺติผุสนาย สนฺติส\-จฺฉิ\-กิริยาย สํวตฺตนฺติ. เสยฺยถิทํ, จตฺตาโร สติปฏฺฐานา จตฺตาโร สมฺมปฺปธานา จตฺตาโร อิทฺธิปาทา ปญฺจินฺทฺริยานิ ปญฺจ พลานิ สตฺต โพชฺฌงฺคา \csromanfolio{143}อริโย อฏฺฐงฺคิโก มคฺโค, อิเม วุจฺจนฺติ สนฺติปทา. สนฺติปทํ ตาณปทํ เลณปทํ สรณปทํ อภยปทํ อจฺจุตปทํ อมตปทํ นิพฺพานปทํ พฺรูหิ อาจิกฺขาหิ เทเสหิ ปญฺญเปหิ ปฏฺฐเปหิ วิวราหิ วิภชาหิ อุตฺตานีกโรหิ ปกาเสหิ. \textbf{สหชเนตฺตา}ติ \textbf{เนตฺตํ} วุจฺจติ สพฺพญฺญุต\-ญฺญาณํ. พุทฺธสฺส ภควโต เนตฺตญฺจ ชินภาโว จ โพธิยา มูเล อปุพฺพํ อจริมํ เอกสฺมิํ ขเณ อุปฺปนฺโน ตสฺมา พุทฺโธ สหชเนตฺโตติ สนฺติปทํ พฺรูหิ สหชเนตฺต.}

\prose{\textbf{ยถาตจฺฉํ ภควา พฺรูหิ เม’ตนฺ}ติ ยถาตจฺฉํ วุจฺจติ อมตํ นิพฺพานํ ฯเปฯ นิโรโธ นิพฺพานํ. \textbf{ภควา}ติ คารวา\-ธิวจนํ ฯเปฯ สจฺฉิกา ปญฺญตฺติ, ยทิทํ ภควาติ. \textbf{พฺรูหิ เมตนฺ}ติ พฺรูหิ อาจิกฺขาหิ ฯเปฯ ปกาเสหีติ ยถาตจฺฉํ ภควา พฺรูหิ เม’ตํ. เตนาห โส พฺราหฺมโณ\csromandash{}}

\prose{สุตฺวานหํ วีร อกามกามิํ, (อิจฺจายสฺมา ชตุกณฺณิ)}

\prose{โอฆาติคํ ปุฏฺฐุม\-กามมาคมํ.}

\prose{สนฺติปทํ พฺรูหิ สหชเนตฺต,}

\prose{ยถาตจฺฉํ ภควา พฺรูหิ เม’ตนฺติ.}

\setlength{\csromangathapagewidth}{0pt}
\setlength{\csromangathawakwidth}{0pt}
\csromangathameasuregroup{}{\textbf{ภควา หิ กาเม อภิภุยฺย อิริยติ,} \\ \textbf{อาทิจฺโจว ปถวิํ เตชี เตชสา.} \\ \textbf{ปริตฺต}\textbf{ปญฺญสฺส เม ภูริปญฺโญ,} \\ \textbf{อาจิกฺข ธมฺมํ ยมหํ วิชญฺญํ.} \\ \textbf{ชาติชราย อิธ วิปฺปหานํ.}}
\csromangathameasurewak{\textbf{ภควา หิ กาเม อภิภุยฺย อิริยติ,} \\ \textbf{อาทิจฺโจว ปถวิํ เตชี เตชสา.} \\ \textbf{ปริตฺต}\textbf{ปญฺญสฺส เม ภูริปญฺโญ,} \\ \textbf{อาจิกฺข ธมฺมํ ยมหํ วิชญฺญํ.} \\ \textbf{ชาติชราย อิธ วิปฺปหานํ.}}
\setlength{\csromangathaleftcol}{0pt}
\csromangathagroup{\csromangathastanza{\csromangathawak{\csromangathaitemline{66}{\textbf{ภควา หิ กาเม อภิภุยฺย อิริยติ,}} \\ \csromangathacontline{\textbf{อาทิจฺโจว ปถวิํ เตชี เตชสา.}} \\ \csromangathacontline{\textbf{ปริตฺต}\-\textbf{ปญฺญสฺส เม ภูริปญฺโญ,}} \\ \csromangathacontline{\textbf{อาจิกฺข ธมฺมํ ยมหํ วิชญฺญํ.}} \\ \csromangathacontline{\textbf{ชาติชราย อิธ วิปฺปหานํ.}}}}}

\prose{\textbf{ภควา หิ กาเม อภิภุยฺย อิริยตี}ติ \textbf{ภควา}ติ คารวา\-ธิวจนํ ฯเปฯ สจฺฉิกา ปญฺญตฺติ, ยทิทํ ภควาติ. \textbf{กามา}ติ อุทฺทานโต เทฺว กามา วตฺถุกามา จ กิเลสกามา จ ฯเปฯ. อิเม วุจฺจนฺติ วตฺถุกามา ฯเปฯ. อิเม วุจฺจนฺติ กิเลสกามา. ภควา วตฺถุกาเม ปริชานิตฺวา กิเลสกาเม ปหาย อภิภุยฺย อภิภวิตฺวา อชฺโฌตฺถริตฺวา ปริยาทิยิตฺวา จรติ วิหรติ อิริยติ วตฺเตติ ปาเลติ ยเปติ ยาเปตีติ ภควา หิ กาเม อภิภุยฺย อิริยติ.}

\prose{\textbf{อาทิจฺโจว ปถวิํ เตชี เตชสา}ติ อาทิจฺโจ วุจฺจติ สูริโย\csromansharedfootnote{9314c6f2c359aad9}{สุริโย (สฺยา)}. ปถวี วุจฺจติ ชคตี\csromansharedfootnote{c1a933a11aa86389}{ชรา (สฺยา)}. ยถา สูริโย เตชี เตเชน สมนฺนาคโต \csromanfolio{144}\textbf{ปถวิํ อภิภุยฺย อภิภวิตฺวา อชฺโฌตฺถริตฺวา ปริยาทิยิตฺวา สนฺตาปยิตฺวา} \textbf{สพฺพํ อากาสคตํ ตมคตํ อภิวิหจฺจ อนฺธการํ วิธมิตฺวา อาโลกํ} \textbf{ทสฺสยิตฺวา อากาเส อนฺตลิกฺเข คคนปเถ}\csromansharedfootnote{8102752c744898ba}{คมนปเถ (สฺยา) อฏฺฐกถา โอโลเกตพฺพา.}\textbf{ คจฺฉติ. เอวเมว ภควา} \textbf{ญาณเตชี ญาณเตเชน สมนฺนาคโต สพฺพํ อภิสงฺขาร}\-\textbf{สมุทยํ ฯเปฯ} กิเลสตมํ อวิชฺช\-นฺธการํ วิธมิตฺวา ญาณาโลกํ ทสฺเสตฺวา วตฺถุกาเม ปริชานิตฺวา กิเลสกาเม ปหาย อภิภุยฺย อภิภวิตฺวา อชฺโฌตฺถริตฺวา ปริยาทิยิตฺวา มทฺทิตฺวา จรติ วิหรติ อิริยติ วตฺเตติ ปาเลติ ยเปติ ยาเปตีติ อาทิจฺโจว ปถวิํ เตชี เตชสา.}

\prose{\textbf{ปริตฺต}\-\textbf{ปญฺญสฺส เม ภูริปญฺโญ}ติ อหมสฺมิ ปริตฺตปญฺโญ โอมกปญฺโญ ลามกปญฺโญ ฉตุกฺกปญฺโญ. ตฺวมฺปิ มหาปญฺโญ ปุถุปญฺโญ หาสปญฺโญ ชวนปญฺโญ ติกฺขปญฺโญ นิพฺเพธิก\-ปญฺโญ. ภูริ วุจฺจติ ปถวี. ภควา ตาย ปถวิสมาย ปญฺญาย วิปุลาย วิตฺถตาย สมนฺนาคโตติ ปริตฺต\-ปญฺญสฺส เม ภูริปญฺโญ.}

\prose{\textbf{อาจิกฺข ธมฺมํ ยมหํ วิชญฺญนฺ}ติ \textbf{ธมฺมนฺ}ติ อาทิกลฺยาณํ มชฺเฌกลฺยาณํ ปริโยสาน\-กลฺยาณํ สาตฺถํ สพฺยญฺชนํ เกวล\-ปริปุณฺณํ ปริสุทฺธํ พฺรหฺมจริยํ จตฺตาโร สติปฏฺฐาเน ฯเปฯ นิพฺพานญฺจ นิพฺพานคามินิญฺจ ปฏิปทํ อาจิกฺขาหิ เทเสหิ ปญฺญเปหิ ปฏฺฐเปหิ วิวราหิ วิภชาหิ อุตฺตานีกโรหิ ปกาเสหิ. \textbf{ยมหํ วิชญฺญนฺ}ติ ยมหํ ชาเนยฺยํ อาชาเนยฺยํ วิชาเนยฺยํ ปฏิชาเนยฺยํ ปฏิวิชฺเฌยฺยํ อธิคจฺเฉยฺยํ ผสฺเสยฺยํ สจฺฉิ\-กเรยฺยนฺติ อาจิกฺข ธมฺมํ ยมหํ วิชญฺญํ.}

\prose{\textbf{ชาติชราย อิธ วิปฺปหานนฺ}ติ อิเธว ชาติชราย มรณสฺส ปหานํ วูปสมํ ปฏินิสฺสคฺคํ ปฏิปฺปสฺสทฺธิํ อมตํ นิพฺพานนฺติ ชาติชราย อิธ วิปฺปหานํ. เตนาห โส พฺราหฺมณฺ\csromandash{}}

\setlength{\csromangathapagewidth}{0pt}
\setlength{\csromangathawakwidth}{0pt}
\csromangathameasuregroup{}{ภควา หิ กาเม อภิภุยฺย อิริยติ, \\ อาทิจฺโจว ปถวิํ เตชี เตชสา.}
\csromangathameasurewak{ภควา หิ กาเม อภิภุยฺย อิริยติ, \\ อาทิจฺโจว ปถวิํ เตชี เตชสา.}
\setlength{\csromangathaleftcol}{0pt}
\csromangathagroup{\csromangathastanza{\csromangathawak{ภควา หิ กาเม อภิภุยฺย อิริยติ, \\ อาทิจฺโจว ปถวิํ เตชี เตชสา.}}}

\prose{ปริตฺต\-ปญฺญสฺส เม ภูริปญฺโญ,}

\vspace{\tipitakaparskip}
\csromancenter{อาจิกฺข ธมฺมํ ยมหํ วิชญฺญํ.}
\csromancenter{ชาติชราย อิธ วิปฺปหานนฺติ.}
\csromanhangingbegin
\hangingheaditem{67}{\csromanfolio{145}\textbf{กาเมสุ วินย เคธํ (ชตุกณฺณีติ ภควา,)}}
\hangingbody{\textbf{เนกฺขมฺมํ ทฏฺฐุ เขมโต.}}
\hangingbody{\textbf{อุคฺคหิตํ นิรตฺตํ วา, มา เต วิชฺชิตฺถ กิญฺจนํ.}}
\csromanhangingend

\prose{\textbf{กาเมสุ วินย เคธนฺ}ติ กามาติ อุทฺทานโต เทฺว \textbf{กามา} วตฺถุกามา จ กิเลสกามา จ ฯเปฯ. อิเม วุจฺจนฺติ วตฺถุกามา ฯเปฯ. อิเม วุจฺจนฺติ กิเลสกามา. \textbf{เคธนฺ}ติ เคโธ วุจฺจติ ตณฺหา. โย ราโค สาราโค ฯเปฯ อภิชฺฌา โลโภ อกุสลมูลํ. \textbf{กาเมสุ วินย เคธนฺ}ติ กาเมสุ เคธํ วินย ปฏิวินย ปชห วิโนเทหิ พฺยนฺตีกโรหิ อนภาวํ คเมหีติ กาเมสุ วินย เคธํ. \textbf{ชตุกณฺณี}ติ ภควา ตํ พฺราหฺมณํ โคตฺเตน อาลปติ. \textbf{ภควา}ติ คารวาธิวจนเมตํ ฯเปฯ สจฺฉิกา ปญฺญตฺติ, ยทิทํ ภควาติ ชตุกณฺณีติ ภควา.}

\prose{\textbf{เนกฺขมฺมํ ทฏฺฐุ เขมโต}ติ \textbf{เนกฺขมฺมนฺ}ติ สมฺมาปฏิปทํ อนุโลม\-ปฏิปทํ อปจฺจนีกปฏิปทํ อนฺวตฺถ\-ปฏิปทํ ธมฺมานุธมฺม\-ปฏิปทํ สีเลสุ ปริปูร\-การิตํ อินฺทฺริเยสุ คุตฺตทฺวารตํ โภชเน มตฺตญฺญุตํ ชาคริยานุโยคํ สติสมฺปชญฺญํ จตฺตาโร สติปฏฺฐาเน จตฺตาโร สมฺมปฺปธาเน จตฺตาโร อิทฺธิปาเท ปญฺจินฺทฺริยานิ ปญฺจ พลานิ สตฺต โพชฺฌงฺเค อริยํ อฏฺฐงฺคิกํ มคฺคํ นิพฺพานญฺจ นิพฺพานคามินิญฺจ ปฏิปทํ เขมโต ตาณโต เลณโต สรณโต สรณีภูตโต อภยโต อจฺจุตโต อมตโต นิพฺพานโต ทฏฺฐุํ ปสฺสิตฺวา ตุลยิตฺวา ตีรยิตฺวา วิภาวยิตฺวา วิภูตํ กตฺวาติ เนกฺขมฺมํ ทฏฺฐุ เขมโต.}

\prose{\textbf{อุคฺคหิตํ นิรตฺตํ วา}ติ \textbf{อุคฺคหิตนฺ}ติ ตณฺหาวเสน ทิฏฺฐิวเสน คหิตํ ปรามฏฺฐํ อภินิวิฏฺฐํ อชฺโฌสิตํ อธิมุตฺตํ. \textbf{นิรตฺตํ วา}ติ นิรตฺตํ วา มุญฺจิตพฺพํ วิชหิตพฺพํ วิโนทิตพฺพํ พฺยนฺตีกาตพฺพํ อนภาวํ คเมตพฺพนฺติ อุคฺคหิตํ นิรตฺตํ วา.}

\prose{\textbf{มา เต วิชฺชิตฺถ กิญฺจนนฺ}ติ ราคกิญฺจนํ โทสกิญฺจนํ โมหกิญฺจนํ มานกิญฺจนํ ทิฏฺฐิกิญฺจนํ กิเลสกิญฺจน ทุจฺจริต\-กิญฺจนํ. อิทํ กิญฺจนํ\csromansharedfootnote{1a34b3cfb5288023}{อิเม กิญฺจนา (ก)} ตุยฺหํ มา วิชฺชิตฺถ มา ปวิชฺชิตฺถ มา สํวิชฺชิตฺถ ปชห วิโนเทหิ พฺยนฺตีกโรหิ อนภาวํ คเมหีติ มา เต วิชฺชิตฺถ กิญฺจนํ. เตนาห ภควา\csromandash{}}

\prose{\csromanfolio{146}กาเมสุ วินย เคธํ, (ชตุกณฺณีติ ภควา,)}

\prose{เนกฺขมฺมํ ทฏฺฐุ เขมโต.}

\setlength{\csromangathapagewidth}{0pt}
\setlength{\csromangathawakwidth}{0pt}
\csromangathameasuregroup{อุคฺคหิตํ นิรตฺตํ วา, \\ \textbf{ยํ ปุพฺเพ ตํ วิโสสฺ}เอหิ, \\ \textbf{มชฺเฌ เจ โน คเห}สฺสสิ,}{\csromangathabat{อุคฺคหิตํ นิรตฺตํ วา,}{มา เต วิชฺชิตฺถ กิญฺจนนฺติ.} \\ \csromangathabat{\textbf{ยํ ปุพฺเพ ตํ วิโสสฺ}เอหิ,}{ปจฺฉา เต มาหุ กิญฺจนํ.} \\ \csromangathabat{\textbf{มชฺเฌ เจ โน คเห}สฺสสิ,}{อุปสนฺโต จริสฺสสิ.}}
\csromangathasetleft{อุคฺคหิตํ นิรตฺตํ วา, \\ \textbf{ยํ ปุพฺเพ ตํ วิโสสฺ}เอหิ, \\ \textbf{มชฺเฌ เจ โน คเห}สฺสสิ,}
\csromangathagroup{\csromangathastanza{\csromangathabat{อุคฺคหิตํ นิรตฺตํ วา,}{มา เต วิชฺชิตฺถ กิญฺจนนฺติ.}}\csromangathastanza{\csromangathaitembat{68}{\textbf{ยํ ปุพฺเพ ตํ วิโสสฺ}เอหิ,}{ปจฺฉา เต มาหุ กิญฺจนํ.} \\ \csromangathacontbat{\textbf{มชฺเฌ เจ โน คเห}สฺสสิ,}{อุปสนฺโต จริสฺสสิ.}}}

\prose{\textbf{ยํ ปุพฺเพ ตํ วิโสเสหี}ติ อตีเต สงฺขาเร อารพฺภ เย กิเลสา อุปฺปชฺเชยฺยุํ, เต กิเลเส โสเสหิ วิโสเสหิ สุกฺขาเปหิ วิสุกฺขาเปหิ อพีชํ กโรหิ ปชห วิโนเทหิ พฺยนฺตีกโรหิ อนภาวํ คเมหีติ เอวมฺปิ ยํ ปุพฺเพ ตํ วิโสเสหิ. อถ วา เย อตีตา กมฺมา\-ภิสงฺขารา อวิปกฺกวิปากา, เต กมฺมา\-ภิสงฺขาเร โสเสหิ วิโสเสหิ สุกฺขาเปหิ วิสุกฺขาเปหิ อพีชํ\csromansharedfootnote{8ae096419eee472e}{อวีชํ (สฺยา)} กโรหิ ปชห วิโนเทหิ พฺยนฺตีกโรหิ อนภาวํ คเมหีติ เอวมฺปิ ยํ ปุพฺเพ ตํ วิโสเสหิ.}

\prose{\textbf{ปจฺฉา เต มาหุ กิญฺจนนฺ}ติ ปจฺฉา วุจฺจติ อนาคเต สงฺขาเร อารพฺภ ราคกิญฺจนํ โทสกิญฺจนํ โมหกิญฺจนํ มานกิญฺจนํ ทิฏฺฐิกิญฺจนํ กิเลสกิญฺจนํ ทุจฺจริต\-กิญฺจนํ, อิทํ กิญฺจนํ ตุยฺหํ มา อหุ มา อโหสิ มา ชเนสิ\csromansharedfootnote{0605d1ee939d0fb5}{มา ชเนหิ (สฺยา) ตถาวเสเสสุ ทฺวีสุ ปเทสุปิ.} มา สญฺชเนสิ มาภินิพฺพตฺเตสิ ปชห วิโนเทหิ พฺยนฺตีกโรหิ อนภาวํ คเมหีติ ปจฺฉา เต มาหุ กิญฺจนํ.}

\prose{\textbf{มชฺเฌ เจ โน คเหสฺสสี}ติ มชฺเฌ วุจฺจติ ปจฺจุปฺปนฺนํ รูปํ เวทนา สญฺญา สงฺขารา วิญฺญาณํ, ปจฺจุปฺปนฺเน สงฺขาเร ตณฺหาวเสน ทิฏฺฐิวเสน น คเหสฺสสิ น คณฺหิสฺสสิ น ปรามสิสฺสสิ น นนฺทิสฺสสิ นาภินนฺทิสฺสสิ น อชฺโฌสิสฺสสิ, อภินนฺทนํ อภิวทนํ อชฺโฌสานํ คาหํ ปรามาสํ อภินิเวสํ ปชหิสฺสสิ วิโนเทสฺสสิ พฺยนฺตีกริสฺสสิ อนภาวํ คเมสฺสสีติ มชฺเฌ เจ โน คเหสฺสสิ.}

\prose{\textbf{อุปสนฺโต จริสฺสสี}ติ ราคสฺส อุปสมิตตฺตา อุปสนฺโต จริสฺสสิ, โทสสฺส ฯเปฯ สพฺพา\-กุสลา\-ภิสงฺขารานํ สนฺตตฺตา สมิตตฺตา อุปสมิตตฺตา วูปสมิตตฺตา นิชฺฌาตตฺตา นิพฺพุตตฺตา วิคตตฺตา ปฏิปฺปสฺสทฺธ\-ตฺตา สนฺโต อุปสนฺโต วูปสนฺโต นิพฺพุโต ปฏิปฺปสฺสทฺโธ จริสฺสสิ วิหริสฺสสิ \csromanfolio{147}อิริยิสฺสสิ วตฺติสฺสสิ ปาเลสฺสสิ ยเปสฺสสิ ยาเปสฺสสีติ อุปสนฺโต จริสฺสสิ. เตนาห ภควา\csromandash{}}

\setlength{\csromangathapagewidth}{0pt}
\setlength{\csromangathawakwidth}{0pt}
\csromangathameasuregroup{ยํ ปุพฺเพ ตํ วิโสเสหิ, \\ มชฺเฌ เจ โน คเหสฺสสิ, \\ \textbf{สพฺพโส นามรูปสฺ}มิํ, \\ \textbf{อาสวาสฺส น วิชฺช}นฺติ,}{\csromangathabat{ยํ ปุพฺเพ ตํ วิโสเสหิ,}{ปจฺฉา เต มาหุ กิญฺจนํ.} \\ \csromangathabat{มชฺเฌ เจ โน คเหสฺสสิ,}{อุปสนฺโต จริสฺสสีติ.} \\ \csromangathabat{\textbf{สพฺพโส นามรูปสฺ}มิํ,}{วีตเคธสฺส พฺราหฺมณ.} \\ \csromangathabat{\textbf{อาสวาสฺส น วิชฺช}นฺติ,}{เยหิ มจฺจุวสํ วเช.}}
\csromangathasetleft{ยํ ปุพฺเพ ตํ วิโสเสหิ, \\ มชฺเฌ เจ โน คเหสฺสสิ, \\ \textbf{สพฺพโส นามรูปสฺ}มิํ, \\ \textbf{อาสวาสฺส น วิชฺช}นฺติ,}
\csromangathagroup{\csromangathastanza{\csromangathabat{ยํ ปุพฺเพ ตํ วิโสเสหิ,}{ปจฺฉา เต มาหุ กิญฺจนํ.} \\ \csromangathabat{มชฺเฌ เจ โน คเหสฺสสิ,}{อุปสนฺโต จริสฺสสีติ.}}\csromangathastanza{\csromangathaitembat{69}{\textbf{สพฺพโส นามรูปสฺ}มิํ,}{วีตเคธสฺส พฺราหฺมณ.} \\ \csromangathacontbat{\textbf{อาสวาสฺส น วิชฺช}นฺติ,}{เยหิ มจฺจุวสํ วเช.}}}

\prose{\textbf{สพฺพโส นามรูปสฺมิํ, วีตเคธสฺส พฺราหฺมณา}ติ \textbf{สพฺพโส}ติ สพฺเพน สพฺพํ สพฺพถา สพฺพํ อเสสํ นิสฺเสสํ ปริยาทิย\-น\-วจนเมตํ สพฺพโสติ. \textbf{นามนฺ}ติ จตฺตาโร อรูปิโน ขนฺธา. \textbf{รูปนฺ}ติ จตฺตาโร จ มหาภูตา จตุนฺนญฺจ มหาภูตานํ อุปาทาย รูปํ. เคโธ วุจฺจติ ตณฺหา. โย ราโค สาราโค ฯเปฯ อภิชฺฌา โลโภ อกุสลมูลํ. \textbf{สพฺพโส นามรูปสฺมิํ, วีตเคธสฺส พฺราหฺมณา}ติ สพฺพโส นามรูปสฺมิํ วีตเคธสฺส วิคตเคธสฺส จตฺตเคธสฺส วนฺตเคธสฺส มุตฺตเคธสฺส ปหีนเคธสฺส ปฏินิสฺสฏฺฐ\-เคธสฺส, วีตราคสฺส วิคตราคสฺส จตฺตราคสฺส วนฺตราคสฺส มุตฺตราคสฺส ปหีนราคสฺส ปฏินิสฺสฏฺฐ\-ราคสฺสาติ สพฺพโส นามรูปสฺมิํ, วีตเคธสฺส พฺราหฺมณ.}

\prose{\textbf{อาสวาสฺส น วิชฺชนฺตี}ติ \textbf{อาสวา}ติ จตฺตาโร อาสวา กามาสโว ภวาสโว ทิฏฺฐาสโว อวิชฺชาสโว. \textbf{อสฺสา}ติ อรหโต ขีณาสวสฺส. \textbf{น วิชฺชนฺตี}ติ อิเม อาสวา ตสฺส นตฺถิ น สนฺติ น สํวิชฺชนฺติ นุปลพฺภนฺติ, ปหีนา สมุจฺฉินฺนา วูปสนฺตา ปฏิปฺปสฺสทฺธา อภพฺพุ\-ปฺปตฺติกา ญาณคฺคินา ทฑฺฒาติ อาสวาสฺส น วิชฺชนฺติ.}

\prose{\textbf{เยหิ มจฺจุวสํ วเช}ติ เยหิ อาสเวหิ มจฺจุโน วา วสํ คจฺเฉยฺย, มรณสฺส วา วสํ คจฺเฉยฺย, มารปกฺขสฺส วา วสํ คจฺเฉยฺย, เต อาสวา ตสฺส นตฺถิ น สนฺติ น สํวิชฺชนฺติ นุปลพฺภนฺติ, ปหีนา สมุจฺฉินฺนา วูปสนฺตา ปฏิปฺปสฺสทฺธา อภพฺพุ\-ปฺปตฺติกา ญาณคฺคินา ทฑฺฒาติ เยหิ มจฺจุวสํ วเช. เตนาห ภควา\csromandash{}}

\setlength{\csromangathapagewidth}{0pt}
\setlength{\csromangathawakwidth}{0pt}
\csromangathameasuregroup{\textbf{สพฺพโส นามรูปสฺ}มิํ, \\ \textbf{อาสวาสฺส น วิชฺช}นฺติ,}{\csromangathabat{\textbf{สพฺพโส นามรูปสฺ}มิํ,}{วีตเคธสฺส พฺราหฺมณ.} \\ \csromangathabat{\textbf{อาสวาสฺส น วิชฺช}นฺติ,}{เยหิ มจฺจุวสํ วเชติ.}}
\csromangathasetleft{\textbf{สพฺพโส นามรูปสฺ}มิํ, \\ \textbf{อาสวาสฺส น วิชฺช}นฺติ,}
\csromangathagroup{\csromangathastanza{\csromangathaitembat{69}{\textbf{สพฺพโส นามรูปสฺ}มิํ,}{วีตเคธสฺส พฺราหฺมณ.} \\ \csromangathacontbat{\textbf{อาสวาสฺส น วิชฺช}นฺติ,}{เยหิ มจฺจุวสํ วเชติ.}}}

\csromanmatikaanchor{mtk.33}
\csromantocmarkref{subsection}{12. ภทฺราวุธมาณวปุจฺฉานิทฺเทส}{mtk.33}
\bookmark[dest={mtk.33},level=2]{12. ภทฺราวุธมาณวปุจฺฉานิทฺเทส}
\prose{\csromanfolio{148}สห คาถา\-ปริโยสานา ฯเปฯ “สตฺถา เม ภนฺเต ภควา, สาวโกหมสฺมี”ติ.}

\vspace{\tipitakaparskip}
\csromancenter{ชาตุกณฺณิมาณวปุจฺฉานิทฺเทโส เอกาทสโม.}
\csromanhangingbegin
\hangingheaditem{70}{\textbf{โอกญฺชหํ ตณฺหจฺฉิทํ อเนชํ, (อิจฺจายสฺมา ภทฺราวุโธ,)}}
\hangingbody{\textbf{นนฺทิญฺชหํ โอฆติณฺณํ วิมุตฺตํ.}}
\hangingbody{\textbf{กปฺปญฺชหํ อภิยาเจ สุเมธํ,}}
\hangingbody{\textbf{สุตฺวาน นาคสฺส อปนมิ}สฺสนฺติ1 อิโต.}
\csromanhangingend

\prose{\textbf{โอกญฺชหํ ตณฺหจฺฉิทํ อเนชนฺ}ติ \textbf{โอกญฺชหนฺ}ติ รูปธาตุยา โย ฉนฺโท โย ราโค ยา นนฺที ยา ตณฺหา เย อุปยุปาทานา\csromansharedfootnote{afa1eb32207e6040}{อุปายุปาทานา (ก)} เจตโส อธิฏฺฐานา\-ภินิเวสา\-นุสยา, เต พุทฺธสฺส ภควโต ปหีนา อุจฺฉินฺนมูลา ตาลาวตฺถุกตา อนภาวํกตา อายติํ อนุปฺปาทธมฺมา. ตสฺมา พุทฺโธ โอกญฺชโห. เวทนาธาตุยา. สญฺญาธาตุยา. สงฺขาร\-ธาตุยา. วิญฺญาณธาตุยา โย ฉนฺโท โย ราโค ยา นนฺที ยา ตณฺหา เย อุปยุปาทานา เจตโส อธิฏฺฐานา\-ภินิเวสา\-นุสยา, เต พุทฺธสฺส ภควโต ปหีนา อุจฺฉินฺนมูลา ตาลาวตฺถุกตา อนภาวํกตา อายติํ อนุปฺปาทธมฺมา, ตสฺมา พุทฺโธ โอกญฺชโห.}

\prose{\textbf{ตณฺหจฺฉิทนฺ}ติ ตณฺหาติ รูป\textbf{ตณฺหา} ฯเปฯ ธมฺมตณฺหา, สา ตณฺหา พุทฺธสฺส ภควโต ฉินฺนา อุจฺฉินฺนา สมุจฺฉินฺนา วูปสนฺตา ปฏิปฺปสฺสทฺธา อภพฺพุ\-ปฺปตฺติกา ญาณคฺคินา ทฑฺฒา, ตสฺมา พุทฺโธ ตณฺหจฺฉิโท. \textbf{อเนโช}ติ \textbf{เอชา} วุจฺจติ ตณฺหา, โย ราโค สาราโค ฯเปฯ อภิชฺฌา โลโภ อกุสลมูลํ. สา เอชา ตณฺหา พุทฺธสฺส ภควโต ปหีนา อุจฺฉินฺนมูลา ตาลาวตฺถุกตา อนภาวํกตา อายติํ อนุปฺปาทธมฺมา, ตสฺมา พุทฺโธ อเนโช. เอชาย ปหีนตฺตา อเนโช. ภควา ลาเภปิ น อิญฺชติ, อลาเภปิ น อิญฺชติ, ยเสปิ น อิญฺชติ, อยเสปิ น อิญฺชติ, ปสํสายปิ น อิญฺชติ, นินฺทายปิ น อิญฺชติ, สุเขปิ น อิญฺชติ, ทุกฺเขปิ น อิญฺชติ น จลติ น เวธติ น ปเวธติ น สมฺปเวธตีติ ตสฺมา พุทฺโธ \csromanfolio{149}อเนโชติ โอกญฺชหํ ตณฺหจฺฉิทํ อเนชํ. \textbf{อิจฺจายสฺมา ภทฺราวุโธ}ติ \textbf{อิจฺจา}ติ ปทสนฺธิ ฯเปฯ. \textbf{อายสฺมา}ติ ปิยวจนํ ฯเปฯ. \textbf{ภทฺราวุโธ}ติ ตสฺส พฺราหฺมณสฺส นามํ ฯเปฯ อภิลาโปติ อิจฺจายสฺมา ภทฺราวุโธ.}

\prose{\textbf{นนฺทิญฺชหํ โอฆติณฺณํ วิมุตฺตนฺ}ติ นนฺที วุจฺจติ ตณฺหา. โย ราโค สาราโค ฯเปฯ อภิชฺฌา โลโภ อกุสลมูลํ, สา นนฺที สา ตณฺหา พุทฺธสฺส ภควโต ปหีนา อุจฺฉินฺนมูลา ตาลาวตฺถุกตา อนภาวํกตา อายติํ อนุปฺปาทธมฺมา, ตสฺมา พุทฺโธ นนฺทิญฺชิโห. \textbf{โอฆติณฺณนฺ}ติ ภควา กาโมฆํ ติณฺโณ, ภโวฆํ ติณฺโณ, ทิฏฺโฐฆํ ติณฺโณ, อวิชฺโชฆํ ติณฺโณ, สพฺพ\-สํสารปถํ ติณฺโณ อุตฺติณฺโณ นิตฺถิณฺโณ อติกฺกนฺโต สมติกฺกนฺโต วีติวตฺโต, โส วุตฺถวาโส จิณฺณจรโณ ฯเปฯ ชาติ\-มรณ\-สํสาโร, นตฺถิ ตสฺส ปุนพฺภโวติ นนฺทิญฺชหํ โอฆติณฺณํ. \textbf{วิมุตฺตนฺ}ติ ภควโต ราคา จิตฺตํ มุตฺตํ วิมุตฺตํ สุวิมุตฺตํ. โทสา จิตฺตํ. โมหา จิตฺตํ ฯเปฯ สพฺพากุสลา\-ภิสงฺขาเรหิ จิตฺตํ มุตฺตํ วิมุตฺตํ สุวิมุตฺตนฺติ นนฺทิญฺชหํ โอฆติณฺณํ วิมุตฺตํ.}

\prose{\textbf{กปฺปญฺชหํ อภิยาเจ สุเมธนฺ}ติ กปฺปาติ เทฺว \textbf{กปฺปา} ตณฺหากปฺโป จ ทิฏฺฐิกปฺโป จ ฯเปฯ อยํ ตณฺหากปฺโป ฯเปฯ อยํ ทิฏฺฐิกปฺโป. พุทฺธสฺส ภควโต ตณฺหากปฺโป ปหีโน, ทิฏฺฐิกปฺโป ปฏินิสฺสฏฺโฐ, ตณฺหากปฺปสฺส ปหีนตฺตา, ทิฏฺฐิกปฺปสฺส ปฏินิสฺสฏฺฐ\-ตฺตา, ตสฺมา พุทฺโธ กปฺปญฺชโห. \textbf{อภิยาเจ}ติ ยาจามิ อภิยาจามิ อชฺเฌสามิ สาทิยามิ ปตฺถยามิ ปิหยามิ ชปฺปามิ อภิชปฺปามิ. สุเมธา วุจฺจติ ปญฺญา, ยา ปญฺญา ปชานนา ฯเปฯ อโมโห ธมฺมวิจโย สมฺมาทิฏฺฐิ. ภควา อิมาย เมธาย ปญฺญาย อุเปโต สมุเปโต อุปาคโต สมุปาคโต อุปปนฺโน สมุปปนฺโน สมนฺนาคโต, ตสฺมา พุทฺโธ สุเมโธติ กปฺปญฺชหํ อภิยาเจ สุเมธํ.}

\prose{\textbf{สุตฺวาน นาคสฺส อปนมิสฺสนฺติ อิโต}ติ \textbf{นาคสฺสา}ติ นาโค, ภควา อาคุํ น กโรตีติ นาโค, น คจฺฉตีติ นาโค, น อาคจฺฉตีติ นาโค ฯเปฯ. เอวํ ภควา น คจฺฉตีติ นาโค. \textbf{สุตฺวาน นาคสฺส อปนมิสฺสนฺติ อิโต}ติ ตุยฺหํ วจนํ พฺยปฺปถํ เทสนํ อนุสาสนํ อนุสิฏฺฐํ สุตฺวา สุณิตฺวา อุคฺคเหตฺวา อุปธารยิตฺวา อุปลกฺขยิตฺวา อิโต อปนมิสฺสนฺติ วชิสฺสนฺติ ปกฺกมิสฺสนฺติ ทิสาวิทิสํ คมิสฺสนฺตีติ สุตฺวาน นาคสฺส อปนมิสฺสนฺติ อิโต. เตนาห โส พฺราหฺมโณ\csromandash{}}

\proseitem{12}{\csromanfolio{150}โอกญฺชหํ ตณฺหจฺฉิทํ อเนชํ, (อิจฺจายสฺมา ภทฺราวุโธ,)}

\prose{นนฺทิญฺชหํ โอฆติณฺณํ วิมุตฺตํ.}

\setlength{\csromangathapagewidth}{0pt}
\setlength{\csromangathawakwidth}{0pt}
\csromangathameasuregroup{}{กปฺปญฺชหํ อภิยาเจ สุเมธํ, \\ สุตฺวาน นาคสฺส อปนมิสฺสนฺติ อิโตติ.}
\csromangathameasurewak{กปฺปญฺชหํ อภิยาเจ สุเมธํ, \\ สุตฺวาน นาคสฺส อปนมิสฺสนฺติ อิโตติ.}
\setlength{\csromangathaleftcol}{0pt}
\csromangathagroup{\csromangathastanza{\csromangathawak{กปฺปญฺชหํ อภิยาเจ สุเมธํ, \\ สุตฺวาน นาคสฺส อปนมิสฺสนฺติ อิโตติ.}}}

\proseitem{71}{\textbf{นานาชนา ชนปเทหิ สงฺคตา,}}

\prose{\textbf{ตว วีร วากฺยํ อภิกงฺขมานา.}}

\prose{\textbf{เตสํ ตุวํ สาธุ วิยากโรหิ,}}

\prose{\textbf{นานาชนา ชนปเทหิ สงฺคตา}ติ \textbf{นานาชนา}ติ ขตฺติยา จ พฺราหฺมณา จ เวสฺสา จ สุทฺทา จ คหฏฺฐา จ ปพฺพชิตา จ เทวา จ มนุสฺสา จ. \textbf{ชนปเทหิ สงฺคตา}ติ องฺคา จ มคธา จ กลิงฺคา จ กาสิยา จ โกสลา จ วชฺชิยา จ มลฺลา จ เจติยมฺหา จ\csromansharedfootnote{369ca5d37dbb65a8}{เจติยมฺหา จ สาครมฺหา จ (สฺยา)} วํสา จ กุรุมฺหา จ ปญฺจาลา จ มจฺฉา จ สุรเสนา จ อสฺสกา จ อวนฺติยา จ โยนา\csromansharedfootnote{47cbc1a7ed6b7835}{โยนกา (ก) ขุ 7. 120 ปิฏฺเฐปิ.} จ กมฺโพชา จ. \textbf{สงฺคตา}ติ สงฺคตา สมาคตา สโมหิตา สนฺนิปติตาติ นานาชนา ชนปเทหิ สงฺคตา.}

\setlength{\csromangathapagewidth}{0pt}
\setlength{\csromangathawakwidth}{0pt}
\csromangathameasuregroup{}{\textbf{ตว วีร วากฺยํ อภิกงฺข}\textbf{มานา}ติ \textbf{วีรา}ติ วีโร, \\ ภควา วีริยวาติ วีโร, ปหูติ วีโร, วิสวีติ วีโร, อลมตฺโตติ วีโร, วิคตโลมหํโสติปิ วีโร. \\ วิรโต อิธ สพฺพปาปเกหิ, \\ นิรยทุกฺขํ อติจฺจ วีริยวาโส. \\ โส วีริยวา ปธานวา, \\ วีโร ตาทิ ปวุจฺจเต ตถตฺตาติ.}
\csromangathameasurewak{\textbf{ตว วีร วากฺยํ อภิกงฺข}\textbf{มานา}ติ \textbf{วีรา}ติ วีโร, \\ ภควา วีริยวาติ วีโร, ปหูติ วีโร, วิสวีติ วีโร, อลมตฺโตติ วีโร, วิคตโลมหํโสติปิ วีโร. \\ วิรโต อิธ สพฺพปาปเกหิ, \\ นิรยทุกฺขํ อติจฺจ วีริยวาโส. \\ โส วีริยวา ปธานวา, \\ วีโร ตาทิ ปวุจฺจเต ตถตฺตาติ.}
\setlength{\csromangathaleftcol}{0pt}
\csromangathagroup{\csromangathastanza{\csromangathawak{\textbf{ตว วีร วากฺยํ อภิกงฺข}\-\textbf{มานา}ติ \textbf{วีรา}ติ วีโร, \\ ภควา วีริยวาติ วีโร, ปหูติ วีโร, วิสวีติ วีโร, อลมตฺโตติ วีโร, วิคต\-โลมหํโสติปิ วีโร. \\ วิรโต อิธ สพฺพปาปเกหิ, \\ นิรยทุกฺขํ อติจฺจ วีริยวาโส. \\ โส วีริยวา ปธานวา, \\ วีโร ตาทิ ปวุจฺจเต ตถตฺตาติ.}}}

\prose{\textbf{ตว วีร วากฺยํ อภิกงฺข}\-\textbf{มานา}ติ ตุยฺหํ วจนํ พฺยปฺปถํ เทสนํ อนุสาสนํ อนุสิฏฺฐํ. \textbf{อภิกงฺข}\-\textbf{มานา}ติ อภิกงฺขมานา อิจฺฉมานา สาทิยมานา ปตฺถยมานา ปิหยมานา อภิชปฺป\-มานาติ ตว วีร วากฺยํ อภิกงฺขมานา.}

\prose{\textbf{เตสํ ตุวํ สาธุ วิยากโรหี}ติ \textbf{เตสนฺ}ติ เตสํ ขตฺติยานํ พฺราหฺมณานํ เวสฺสานํ สุทฺทานํ คหฏฺฐานํ ปพฺพชิตานํ เทวานํ มนุสฺสานํ. \textbf{ตุวนฺ}ติ ภควนฺตํ ภณติ. \textbf{สาธุ วิยากโรหี}ติ สาธุ อาจิกฺขาหิ \csromanfolio{151}เทเสหิ ปญฺญเปหิ ปฏฺฐเปหิ วิวราหิ วิภชาหิ อุตฺตานีกโรหิ ปกาเสหีติ เตสํ ตุวํ สาธุ วิยากโรหิ.}

\prose{\textbf{ตถา หิ เต วิทิโต เอส ธมฺโม}ติ ตถา หิ เต วิทิโต ตุลิโต ตีริโต วิภูโต วิภาวิโต เอส ธมฺโมติ ตถา หิ เต วิทิโต เอส ธมฺโม. เตนาห โส พฺราหฺมโณ\csromandash{}}

\setlength{\csromangathapagewidth}{0pt}
\setlength{\csromangathawakwidth}{0pt}
\csromangathameasuregroup{}{นานาชนา ชนปเทหิ สงฺคตา, \\ ตว วีร วากฺยํ อภิกงฺขมานา. \\ เตสํ ตุวํ สาธุ วิยากโรหิ,}
\csromangathameasurewak{นานาชนา ชนปเทหิ สงฺคตา, \\ ตว วีร วากฺยํ อภิกงฺขมานา. \\ เตสํ ตุวํ สาธุ วิยากโรหิ,}
\setlength{\csromangathaleftcol}{0pt}
\csromangathagroup{\csromangathastanza{\csromangathawak{นานาชนา ชนปเทหิ สงฺคตา, \\ ตว วีร วากฺยํ อภิกงฺขมานา. \\ เตสํ ตุวํ สาธุ วิยากโรหิ,}}}

\prose{ตถา หิ เต วิทิโต เอส ธมฺโมติ.}

\csromanhangingbegin
\hangingheaditem{72}{\textbf{อาทานตณฺหํ วินเยถ สพฺพํ, (ภทฺราวุธาติ ภควา,)}}
\hangingbody{\textbf{อุทฺธํ อโธ ติริยญฺจาปิ มชฺเฌ.}}
\hangingbody{\textbf{ยํ ยญฺหิ โลกสฺมิมุปาทิยนฺติ,}}
\hangingbody{\textbf{เตเนว มาโร อเนฺวติ ชนฺตุํ.}}
\csromanhangingend

\prose{\textbf{อาทานตณฺหํ วินเยถ สพฺพนฺ}ติ อาทานตณฺหํ วุจฺจติ รูปตณฺหา ฯเปฯ อาทานตณฺหาติ กิํการณา วุจฺจติ อาทานตณฺหา, ตาย ตณฺหาย รูปํ อาทิยนฺติ อุปาทิยนฺติ คณฺหนฺติ ปรามสนฺติ อภินิวิสนฺติ. เวทนํ. สญฺญํ. สงฺขาเร. วิญฺญาณํ. คติํ. อุปปตฺติํ. ปฏิสนฺธิํ. ภวํ. สํสารํ. วฏฺฏํ อาทิยนฺติ อุปาทิยนฺติ คณฺหนฺติ ปรามสนฺติ อภินิวิสนฺติ, ตํการณา วุจฺจติ อาทานตณฺหา. \textbf{อาทานตณฺหํ วินเยถ สพฺพนฺ}ติ สพฺพํ อาทานตณฺหํ วินเยยฺย ปฏิวินเยยฺย ปชเหยฺย วิโนเทยฺย พฺยนฺตีกเรยฺย อนภาวํ คเมยฺยาติ อาทานตณฺหํ วินเยถ สพฺพํ, \textbf{ภทฺราวุธาติ ภควา}ติ \textbf{ภทฺราวุธา}ติ ภควา ตํ พฺราหฺมณํ นาเมน อาลปติ. \textbf{ภควา}ติ คารวาธิวจนเมตํ ฯเปฯ สจฺฉิกา ปญฺญตฺติ, ยทิทํ ภควาติ ภทฺราวุธาติ ภควา.}

\prose{\textbf{อุทฺธํ อโธ ติริยญฺจาปิ มชฺเฌ}ติ \textbf{อุทฺธนฺ}ติ อนาคตํ. \textbf{อโธ}ติ อตีตํ. \textbf{ติริยญฺจาปิ มชฺเฌ}ติ ปจฺจุปฺปนฺนํ. \textbf{อุทฺธนฺ}ติ เทวโลโก. \textbf{อโธ}ติ นิรยโลโก. \textbf{ติริยญฺจาปิ มชฺเฌ}ติ มนุสฺสโลโก. อถ วา \textbf{อุทฺธนฺ}ติ กุสลา ธมฺมา. \textbf{อโธ}ติ อกุสลา ธมฺมา. \textbf{ติริยญฺจาปิ มชฺเฌ}ติ อพฺยากตา ธมฺมา. \textbf{อุทฺธนฺ}ติ อรูปธาตุ. \textbf{อโธ}ติ กามธาตุ. \textbf{ติริยญฺจาปิ มชฺเฌ}ติ รูปธาตุ. \textbf{อุทฺธนฺ}ติ สุขา เวทนา. \textbf{อโธ}ติ ทุกฺขา เวทนา. \textbf{ติริยญฺจาปิ มชฺเฌ}ติ อทุกฺขมสุขา เวทนา. \textbf{อุทฺธนฺ}ติ อุทฺธํ ปาทตลา. \csromanfolio{152}\textbf{อโธติ อโธ เกสมตฺถกา.}\textbf{ ติริยญฺจาปิ มชฺเฌติ เวมชฺเฌติ อุทฺธํ อโธ} \textbf{ติริยญฺจาปิ มชฺเฌ.}}

\prose{\textbf{ยํ ยญฺหิ โลกสฺมิมุ}\-\textbf{ปาทิยนฺตี}ติ ยํ ยํ รูปคตํ เวทนาคตํ สญฺญาคตํ สงฺขารคตํ วิญฺญาณคตํ อาทิยนฺติ อุปาทิยนฺติ คณฺหนฺติ ปรามสนฺติ อภินิวิสนฺติ. \textbf{โลกสฺมินฺ}ติ อปายโลเก ฯเปฯ อายตนโลเกติ ยํ ยญฺหิ โลกสฺมิมุ\-ปาทิยนฺติ.}

\prose{\textbf{เตเนว มาโร อเนฺวติ ชนฺตุนฺ}ติ เตเนว กมฺมา\-ภิสงฺขาร\-วเสน ปฏิสนฺธิโก ขนฺธมาโร ธาตุมาโร อายตนมาโร คติมาโร อุปปตฺติมาโร ปฏิสนฺธิมาโร ภวมาโร สํสารมาโร วฏฺฏมาโร อเนฺวติ อนุคจฺฉติ อนฺวายิโก โหติ. \textbf{ชนฺตุนฺ}ติ สตฺตํ ชนํ นรํ มาณวํ\csromansharedfootnote{6141be559683362c}{มาณวํ (สฺยา)} โปสํ ปุคฺคลํ ชีวํ ชาคุํ ชนฺตุํ อินฺทคุํ มนุชนฺติ เตเนว มาโร อเนฺวติ ชนฺตุํ. เตนาห ภควา\csromandash{}}

\prose{อาทานตณฺหํ วินเยถ สพฺพํ, (ภทฺราวุธาติ ภควา,)}

\prose{อุทฺธํ อโธ ติริยญฺจาปิ มชฺเฌ.}

\setlength{\csromangathapagewidth}{0pt}
\setlength{\csromangathawakwidth}{0pt}
\csromangathameasuregroup{}{ยํ ยญฺหิ โลกสฺมิมุปาทิยนฺติ, \\ เตเนว มาโร อเนฺวติ ชนฺตุนฺติ. \\ \textbf{ตสฺมา ปชานํ น อุปาทิเยถ,} \\ \textbf{ภิกฺขุ สโต กิญฺจนํ สพฺพโลเก.} \\ \textbf{อาทานสตฺเต อิติ เปกฺขมาโน,} \\ \textbf{ปชํ อิมํ มจฺจุเธยฺเย วิสตฺตํ.}}
\csromangathameasurewak{ยํ ยญฺหิ โลกสฺมิมุปาทิยนฺติ, \\ เตเนว มาโร อเนฺวติ ชนฺตุนฺติ. \\ \textbf{ตสฺมา ปชานํ น อุปาทิเยถ,} \\ \textbf{ภิกฺขุ สโต กิญฺจนํ สพฺพโลเก.} \\ \textbf{อาทานสตฺเต อิติ เปกฺขมาโน,} \\ \textbf{ปชํ อิมํ มจฺจุเธยฺเย วิสตฺตํ.}}
\setlength{\csromangathaleftcol}{0pt}
\csromangathagroup{\csromangathastanza{\csromangathawak{ยํ ยญฺหิ โลกสฺมิมุ\-ปาทิยนฺติ, \\ เตเนว มาโร อเนฺวติ ชนฺตุนฺติ.}}\csromangathastanza{\csromangathawak{\csromangathaitemline{73}{\textbf{ตสฺมา ปชานํ น อุปาทิเยถ,}} \\ \csromangathacontline{\textbf{ภิกฺขุ สโต กิญฺจนํ สพฺพโลเก.}} \\ \csromangathacontline{\textbf{อาทานสตฺเต อิติ เปกฺขมาโน,}} \\ \csromangathacontline{\textbf{ปชํ อิมํ มจฺจุเธยฺเย วิสตฺตํ.}}}}}

\prose{\textbf{ตสฺมา ปชานํ น อุปาทิเยถา}ติ ตสฺมาติ \textbf{ตสฺมา} ตํการณา ตํเหตุ ตปฺปจฺจยา ตํนิทานา เอตํ อาทีนวํ สมฺปสฺสมาโน อาทานตณฺหายาติ ตสฺมา. \textbf{ปชานนฺ}ติ ชานนฺโต ปชานนฺโต อาชานนฺโต วิชานนฺโต ปฏิวิชานนฺโต ปฏิวิชฺฌนฺโต “สพฺเพ สงฺขารา อนิจฺจา”ติ ฯเปฯ. “ยํ กิญฺจิ สมุทย\-ธมฺมํ, สพฺพํ ตํ นิโรธธมฺมนฺ”ติ ชานนฺโต ปชานนฺโต อาชานนฺโต วิชานนฺโต ปฏิวิชานนฺโต ปฏิวิชฺฌนฺโต. \textbf{น อุปาทิเยถา}ติ รูปํ นาทิเยยฺย น อุปาทิเยยฺย น คเณฺหยฺย น ปรามเสยฺย นาภินิวิเสยฺย. เวทนํ. สญฺญํ. สงฺขาเร. วิญฺญาณํ, คติํ, อุปปตฺติํ, ปฏิสนฺธิํ, ภวํ. สํสารํ. วฏฺฏํ นาทิเยยฺย น อุปาทิเยยฺย น คเณฺหยฺย น ปรามเสยฺย นาภิ\-นิวิเสยฺยาติ ตสฺมา ปชานํ น อุปาทิเยถ.}

\prose{\csromanfolio{153}\textbf{ภิกฺขุ สโต กิญฺจนํ สพฺพโลเก}ติ \textbf{ภิกฺขู}ติ ปุถุชฺชน\-กลฺยาณโก วา ภิกฺขุ, เสกฺโข วา ภิกฺขุ. \textbf{สโต}ติ จตูหิ การเณหิ สโต, กาเย กายานุปสฺสนา\-สติปฏฺฐานํ ภาเวนฺโต สโต ฯเปฯ โส วุจฺจติ สโตติ ภิกฺขุ สโต. \textbf{กิญฺจนนฺ}ติ กิญฺจิ รูปคตํ เวทนาคตํ สญฺญาคตํ สงฺขารคตํ วิญฺญาณคตํ. \textbf{สพฺพโลเก}ติ สพฺพอปายโลเก สพฺพ\-มนุสฺสโลเก สพฺพเทวโลเก สพฺพ\-ขนฺธโลเก สพฺพธาตุโลเก สพฺพอายตนโลเกติ ภิกฺขุ สโต กิญฺจนํ สพฺพโลเก.}

\prose{\textbf{อาทานสตฺเต อิติ เปกฺขมาโน}ติ อาทานสตฺตา วุจฺจนฺติ เย รูปํ อาทิยนฺติ อุปาทิยนฺติ คณฺหนฺติ ปรามสนฺติ อภินิวิสนฺติ. เวทนํ. สญฺญํ. สงฺขาเร. วิญฺญาณํ. คติํ. อุปปตฺติํ. ปฏิสนฺธิํ. ภวํ. สํสารํ. วฏฺฏํ อาทิยนฺติ อุปาทิยนฺติ คณฺหนฺติ ปรามสนฺติ อภินิวิสนฺติ. \textbf{อิตี}ติ ปทสนฺธิ ฯเปฯ ปทา\-นุปุพฺพตา\-เปตํ อิตีติ. \textbf{เปกฺขมาโน}ติ เปกฺขมาโน ทกฺขมาโน ทิสฺสมาโน ปสฺสมาโน โอโลกยมาโน นิชฺฌายมาโน อุปปริกฺขมาโนติ อาทานสตฺเต อิติ เปกฺขมาโน.}

\prose{\textbf{ปชํ อิมํ มจฺจุเธยฺเย วิสตฺตนฺ}ติ \textbf{ปชา}ติ สตฺตา\-ธิวจนํ. มจฺจุเธยฺยา วุจฺจนฺติ กิเลสา จ ขนฺธา จ อภิสงฺขารา จ, ปชา มจฺจุเธยฺเย มารเธยฺเย มรณเธยฺเย สตฺตา วิสตฺตา อาสตฺตา ลคฺคา ลคฺคิตา ปลิพุทฺธา, ยถา ภิตฺติขิเล วา นาคทนฺเต วา ภณฺฑํ สตฺตํ วิสตฺตํ อาสตฺตํ ลคฺคํ ลคฺคิตํ ปลิพุทฺธํ, เอวเมว ปชา มจฺจุเธยฺเย มารเธยฺเย มรณเธยฺเย สตฺตา วิสตฺตา อาสตฺตา ลคฺคา ลคฺคิตา ปลิพุทฺธาติ ปชํ อิมํ มจฺจุเธยฺเย วิสตฺตํ. เตนาห ภควา\csromandash{}}

\prose{ตสฺมา ปชานํ น อุปาทิเยถ,}

\prose{ภิกฺขุ สโต กิญฺจนํ สพฺพโลเก.}

\prose{อาทานสตฺเต อิติ เปกฺขมาโน,}

\prose{ปชํ อิมํ มจฺจุเธยฺเย วิสตฺตนฺติ.}

\prose{สห คาถา\-ปริโยสานา ฯเปฯ “สตฺถา เม ภนฺเต ภควา, สาวโกหมสฺมี”ติ.}

\vspace{\tipitakaparskip}
\csromancenter{ภทฺราวุธมาณวปุจฺฉา\-นิทฺเทโส ทฺวาทสโม.}
\csromanmatikaanchor{mtk.34}
\csromantocmarkref{subsection}{13. อุทยมาณวปุจฺฉานิทฺเทส}{mtk.34}
\bookmark[dest={mtk.34},level=2]{13. อุทยมาณวปุจฺฉานิทฺเทส}
\csromanhangingbegin
\hangingheaditem{74}{\csromanfolio{154}\textbf{ฌายิํ วิรชมาสีนํ, (อิจฺจายสฺมา อุทโย,)}}
\hangingbody{\textbf{กตกิจฺจํ อนาสวํ.}}
\hangingbody{\textbf{ปารคุํ สพฺพธมฺมานํ, อตฺถิ ปเญฺหน อาคมํ.}}
\hangingbody{\textbf{อญฺญาวิโมกฺขํ ปพฺ}รูหิ1, อวิชฺชาย ปเภทนํ.}
\csromanhangingend

\prose{\textbf{ฌายิํ วิรชมาสีนนฺ}ติ \textbf{ฌายินฺ}ติ ฌายี, ภควา ปฐเมนปิ ฌาเนน ฌายี, ทุติเยนปิ ฌาเนน ฌายี, ตติเยนปิ ฌาเนน ฌายี, จตุตฺเถนปิ ฌาเนน ฌายี, สวิตกฺกสวิจาเรนปิ ฌาเนน ฌายี, อวิตกฺกวิจารมตฺเตนปิ ฌาเนน ฌายี, อวิตกฺกอวิจาเรนปิ ฌาเนน ฌายี, สปฺปีติเกนปิ ฌาเนน ฌายี, นิปฺปีติเกนปิ ฌาเนน ฌายี, สาตสหคเตนปิ ฌาเนน ฌายี, อุเปกฺขา\-สหคเตนปิ ฌาเนน ฌายี, สุญฺญเตนปิ ฌาเนน ฌายี, อนิมิตฺเตนปิ ฌาเนน ฌายี, อปฺปณิหิเตนปิ ฌาเนน ฌายี, โลกิเยนปิ ฌาเนน ฌายี, โลกุตฺตเรนปิ ฌาเนน ฌายี, ฌานรโต เอกตฺตม\-นุยุตฺโต สทตฺถครุโกติ ฌายิํ. วิรชนฺติ ราโค รโช, โทโส รโช, โมโห รโช, โกโธ รโช, อุปนาโห รโช ฯเปฯ สพฺพา\-กุสลา\-ภิสงฺขารา รชา, เต รชา พุทฺธสฺส ภควโต ปหีนา อุจฺฉินฺนมูลา ตาลาวตฺถุกตา อนภาวํกตา อายติํ อนุปฺปาทธมฺมา, ตสฺมา พุทฺโธ อรโช วิรโช นิรโช รชาปคโต รชวิปฺปหีโน รชวิปฺปยุตฺโต สพฺพ\-รช\-วีติวตฺโต.}

\setlength{\csromangathapagewidth}{0pt}
\setlength{\csromangathawakwidth}{0pt}
\csromangathameasuregroup{}{ราโค รโช น จ ปน เรณุ วุจฺจติ, \\ ราคสฺเสตํ อธิวจนํ รโชติ. \\ เอตํ รชํ วิปฺปชหิตฺวา\textsuperscript{1} จกฺขุมา, \\ ตสฺมา ชิโน วิคตรโชติ วุจฺจติ. \\ โทโส รโช น จ ปน เรณุ วุจฺจติ, \\ โทสสฺเสตํ อธิวจนํ รโชติ. \\ เอตํ รชํ วิปฺปชหิตฺวา จกฺขุมา, \\ ตสฺมา ชิโน วิคตรโชติ วุจฺจติ. \\ โมโห รโช น จ ปน เรณุ วุจฺจติ, \\ โมหสฺเสตํ อธิวจนํ รโชติ. \\ เอตํ รชํ วิปฺปชหิตฺวา จกฺขุมา, \\ ตสฺมา ชิโน วิคตรโชติ วุจฺจตีติ.}
\csromangathameasurewak{ราโค รโช น จ ปน เรณุ วุจฺจติ, \\ ราคสฺเสตํ อธิวจนํ รโชติ. \\ เอตํ รชํ วิปฺปชหิตฺวา\textsuperscript{1} จกฺขุมา, \\ ตสฺมา ชิโน วิคตรโชติ วุจฺจติ. \\ โทโส รโช น จ ปน เรณุ วุจฺจติ, \\ โทสสฺเสตํ อธิวจนํ รโชติ. \\ เอตํ รชํ วิปฺปชหิตฺวา จกฺขุมา, \\ ตสฺมา ชิโน วิคตรโชติ วุจฺจติ. \\ โมโห รโช น จ ปน เรณุ วุจฺจติ, \\ โมหสฺเสตํ อธิวจนํ รโชติ. \\ เอตํ รชํ วิปฺปชหิตฺวา จกฺขุมา, \\ ตสฺมา ชิโน วิคตรโชติ วุจฺจตีติ.}
\setlength{\csromangathaleftcol}{0pt}
\csromangathagroup{\csromangathastanza{\csromangathawak{ราโค รโช น จ ปน เรณุ วุจฺจติ, \\ ราคสฺเสตํ อธิวจนํ รโชติ. \\ เอตํ รชํ วิปฺปชหิตฺวา\csromansharedfootnote{6157fe761b8c9961}{ปฏิวิโนทิตฺวา (ก) ขุ 7. 406 ปิฏฺเฐปิ.} จกฺขุมา, \\ ตสฺมา ชิโน วิคตรโชติ วุจฺจติ.}}\csromangathastanza{\csromangathawak{โทโส รโช น จ ปน เรณุ วุจฺจติ, \\ โทสสฺเสตํ อธิวจนํ รโชติ. \\ เอตํ รชํ วิปฺปชหิตฺวา จกฺขุมา, \\ ตสฺมา ชิโน วิคตรโชติ วุจฺจติ.}}\csromangathastanza{\csromangathawak{\csromanfolio{155}โมโห รโช น จ ปน เรณุ วุจฺจติ, \\ โมหสฺเสตํ อธิวจนํ รโชติ. \\ เอตํ รชํ วิปฺปชหิตฺวา จกฺขุมา, \\ ตสฺมา ชิโน วิคตรโชติ วุจฺจตีติ.}}}

\prose{วิรชํ. \textbf{อาสีนนฺ}ติ นิสินฺโน ภควา ปาสาณเก เจติเยติ อาสีโน.}

\setlength{\csromangathapagewidth}{0pt}
\setlength{\csromangathawakwidth}{0pt}
\csromangathameasuregroup{นคสฺส\textsuperscript{1} ปสฺเส อาสีนํ,}{\csromangathabat{นคสฺส\textsuperscript{1} ปสฺเส อาสีนํ,}{มุนิํ ทุกฺขสฺส ปารคุํ.}}
\csromangathasetleft{นคสฺส\textsuperscript{1} ปสฺเส อาสีนํ,}
\csromangathagroup{\csromangathastanza{\csromangathabat{นคสฺส\csromansharedfootnote{747983ce23375e73}{นครสฺส (ก)} ปสฺเส อาสีนํ,}{มุนิํ ทุกฺขสฺส ปารคุํ.}}}

\vspace{\tipitakaparskip}
\csromancenter{สาวกา ปยิรุปาสนฺติ, เตวิชฺชา มจฺจุหายิโนติ\csromandash{}}
\prose{เอวมฺปิ ภควา อาสีโน.  อถ วา ภควา สพฺโพ\-สฺสุกฺก\-ปฏิปฺปสฺสทฺธ\-ตฺตา อาสีโน วุตฺถวาโส จิณฺณจรโณ ฯเปฯ ชาติ\-มรณ\-สํสาโร, นตฺถิ ตสฺส ปุนพฺภโวติ เอวมฺปิ ภควา อาสีโนติ ฌายิํ วิรชมาสีนํ.}

\prose{\textbf{อิจฺจายสฺมา อุทโย}ติ \textbf{อิจฺจา}ติ ปทสนฺธิ ฯเปฯ. \textbf{อายสฺมา}ติ ปิยวจนํ ฯเปฯ. \textbf{อุทโย}ติ ตสฺส พฺราหฺมณสฺส นามํ ฯเปฯ อภิลาโปติ อิจฺจายสฺมา อุทโย.}

\prose{\textbf{กตกิจฺจํ อนาสวนฺ}ติ พุทฺธสฺส ภควโต กิจฺจากิจฺจํ กรณียา\-กรณียํ ปหีนํ อุจฺฉินฺนมูลํ ตาลา\-วตฺถุกตํ อนภาวํกตํ อายติํ อนุปฺปาท\-ธมฺมํ, ตสฺมา พุทฺโธ กตกิจฺโจ.}

\setlength{\csromangathapagewidth}{0pt}
\setlength{\csromangathawakwidth}{0pt}
\csromangathameasuregroup{ยสฺส จ วิสตา\textsuperscript{1} นตฺถิ, \\ กิจฺจากิจฺจปฺปหีนสฺส,}{\csromangathabat{ยสฺส จ วิสตา\textsuperscript{1} นตฺถิ,}{ฉินฺนโสตสฺส ภิกฺขุโน.} \\ \csromangathabat{กิจฺจากิจฺจปฺปหีนสฺส,}{ปริฬาโห น วิชฺชตีติ.}}
\csromangathasetleft{ยสฺส จ วิสตา\textsuperscript{1} นตฺถิ, \\ กิจฺจากิจฺจปฺปหีนสฺส,}
\csromangathagroup{\csromangathastanza{\csromangathabat{ยสฺส จ วิสตา\csromansharedfootnote{25c8ba7be179707c}{ยสฺส ปริปตา (สฺยา) ปสฺส ขุ 1. 390 ปิฏฺเฐ.} นตฺถิ,}{ฉินฺนโสตสฺส ภิกฺขุโน.} \\ \csromangathabat{กิจฺจากิจฺจ\-ปฺปหีนสฺส,}{ปริฬาโห น วิชฺชตีติ.}}}

\prose{\textbf{กตกิจฺจํ อนาสวนฺ}ติ \textbf{อาสวา}ติ จตฺตาโร อาสวา กามาสโว ภวาสโว ทิฏฺฐาสโว อวิชฺชาสโว, เต อาสวา พุทฺธสฺส ภควโต ปหีนา อุจฺฉินฺนมูลา ตาลาวตฺถุกตา อนภาวํกตา อายติํ อนุปฺปาทธมฺมา, ตสฺมา พุทฺโธ อนาสโวติ กตกิจฺจํ อนาสวํ.}

\prose{\textbf{ปารคุํ สพฺพ}\-\textbf{ธมฺมานนฺ}ติ ภควา สพฺพธมฺมานํ อภิญฺญาปารคู ปริญฺญาปารคู ปหานปารคู ภาวนาปารคู สจฺฉิกิริยา\-ปารคู สมาปตฺติ\-ปารคู. อภิญฺญาปารคู สพฺพธมฺมานํ, ปริญฺญาปารคู สพฺพทุกฺขานํ, ปหานปารคู สพฺพกิเลสานํ, ภาวนาปารคู จตุนฺนํ มคฺคานํ, สจฺฉิกิริยา\-ปารคู นิโรธสฺส, สมาปตฺติ\-ปารคู สพฺพ\-สมาปตฺตีนํ. โส วสิปฺปตฺโต ปารมิปฺปตฺโต อริยสฺมิํ สีลสฺมิํ,}

\prose{\csromanfolio{156}วสิปฺปตฺโต ปารมิปฺปตฺโต อริยสฺมิํ สมาธิสฺมิํ, วสิปฺปตฺโต ปารมิปฺปตฺโต อริยาย ปญฺญาย, วสิปฺปตฺโต ปารมิปฺปตฺโต อริยาย วิมุตฺติยา. โส ปารคโต ปารปฺปตฺโต อนฺตคโต อนฺตปฺปตฺโต โกฏิคโต โกฏิปฺปตฺโต ปริยนฺตคโต ปริยนฺต\-ปฺปตฺโต โวสานคโต โวสานปฺปตฺโต ตาณคโต ตาณปฺปตฺโต เลณคโต เลณปฺปตฺโต สรณคโต สรณปฺปตฺโต อภยคโต อภยปฺปตฺโต อจฺจุตคโต อจฺจุตปฺปตฺโต อมตคโต อมตปฺปตฺโต นิพฺพานคโต นิพฺพานปฺปตฺโต, โส วุตฺถวาโส จิณฺณจรโณ ฯเปฯ ชาติ\-มรณ\-สํสาโร, นตฺถิ ตสฺส ปุนพฺภโวติ ปารคุํ สพฺพธมฺมานํ.}

\prose{\textbf{อตฺถิ ปเญฺหน อาคมนฺ}ติ ปเญฺหน อตฺถิโก อาคโตมฺหิ, ปญฺหํ ปุจฺฉิตุกาโม อาคโตมฺหิ, ปญฺหํ โสตุกาโม อาคโตมฺหีติ เอวมฺปิ อตฺถิ ปเญฺหน อาคมํ. อถ วา ปญฺหตฺถิกานํ ปญฺหํ ปุจฺฉิตุ\-กามานํ ปญฺหํ โสตุกามานํ อาคมนํ อภิกฺกมนํ อุปสงฺกมนํ ปยิรุปาสนํ อตฺถีติ เอวมฺปิ อตฺถิ ปเญฺหน อาคมํ. อถ วา ปญฺหาคโม ตุยฺหํ อตฺถิ, ตฺวมฺปิ ปหุ, วิสวี อลมตฺโต มยา ปุจฺฉิตํ กเถตุํ วิสชฺเชตุํ วหสฺเสตํ ภารนฺติ เอวมฺปิ อตฺถิ ปเญฺหน อาคมํ.}

\prose{\textbf{อญฺญาวิโมกฺขํ ปพฺรูหี}ติ อญฺญาวิโมกฺโข วุจฺจติ อรหตฺต\-วิโมกฺโข, อรหตฺต\-วิโมกฺขํ ปพฺรูหิ อาจิกฺขาหิ เทเสหิ ปญฺญเปหิ ปฏฺฐเปหิ วิวราหิ วิภชาหิ อุตฺตานีกโรหิ ปกาเสหีติ อญฺญาวิโมกฺขํ ปพฺรูหิ.}

\prose{\textbf{อวิชฺชาย ปเภทนนฺ}ติ อวิชฺชาย เภทนํ ปเภทนํ ปหานํ วูปสมํ ปฏินิสฺสคฺคํ ปฏิปฺปสฺสทฺธํ อมตํ นิพฺพานนฺติ อวิชฺชาย ปเภทนํ. เตนาห โส พฺราหฺมโณ\csromandash{}}

\prose{ฌายิํ วิรชมาสีนํ, (อิจฺจายสฺมา อุทโย,)}

\prose{กตกิจฺจํ อนาสวํ.}

\setlength{\csromangathapagewidth}{0pt}
\setlength{\csromangathawakwidth}{0pt}
\csromangathameasuregroup{ปารคุํ สพฺพธมฺมานํ, \\ อญฺญาวิโมกฺขํ ปพฺรูหิ,}{\csromangathabat{ปารคุํ สพฺพธมฺมานํ,}{อตฺถิ ปเญฺหน อาคมํ.} \\ \csromangathabat{อญฺญาวิโมกฺขํ ปพฺรูหิ,}{อวิชฺชาย ปเภทนนฺติ.}}
\csromangathasetleft{ปารคุํ สพฺพธมฺมานํ, \\ อญฺญาวิโมกฺขํ ปพฺรูหิ,}
\csromangathagroup{\csromangathastanza{\csromangathabat{ปารคุํ สพฺพธมฺมานํ,}{อตฺถิ ปเญฺหน อาคมํ.} \\ \csromangathabat{อญฺญาวิโมกฺขํ ปพฺรูหิ,}{อวิชฺชาย ปเภทนนฺติ.}}}

\csromanhangingbegin
\hangingheaditem{75}{\csromanfolio{157}\textbf{ปหานํ กามจฺฉนฺทานํ, (อุทยาติ ภควา,)}}
\hangingbody{\textbf{โทมนสฺสาน จูภยํ.}}
\hangingbody{\textbf{ถินสฺส1} จ ปนูทนํ, กุกฺกุจฺจานํ นิวารณํ.}
\csromanhangingend

\prose{\textbf{ปหานํ กาม}\-\textbf{จฺฉนฺทานนฺ}ติ ฉนฺโทติ โย กาเมสุ กามจฺ\textbf{ฉนฺโท} กามราโค กามนนฺที กามตณฺหา กามสิเนโห กามปิปาสา กามปริฬาโห กามมุจฺฉา กามชฺโฌสานํ กาโมโฆ กามโยโค กามุปาทานํ กามจฺฉนฺท\-นีวรณํ. \textbf{ปหานํ กาม}\-\textbf{จฺฉนฺทานนฺ}ติ กามจฺฉนฺทานํ ปหานํ วูปสมํ ปฏินิสฺสคฺคํ ปฏิปฺปสฺสทฺธิํ อมตํ นิพฺพานนฺติ ปหานํ กามจฺฉนฺทานํ. อุทยาติ ภควาติ \textbf{อุทยาติ ภควา} ตํ พฺราหฺมณํ นาเมน อาลปติ. \textbf{ภควา}ติ คารวาธิวจนเมตํ ฯเปฯ สจฺฉิกา ปญฺญตฺติ, ยทิทํ ภควาติ อุทยาติ ภควา.}

\prose{\textbf{โทมนสฺสาน จูภยนฺ}ติ \textbf{โทมนสฺสา}ติ ยํ เจตสิกํ อสาตํ, เจตสิกํ ทุกฺขํ เจโต\-สมฺผสฺสชํ อสาตํ ทุกฺขํ เวทยิตํ, เจโต\-สมฺผสฺสชา อสาตา ทุกฺขา เวทนา. \textbf{โทมนสฺสาน จูภยนฺ}ติ กามจฺฉนฺทสฺส จ โทมนสฺสสฺส จ อุภินฺนํ ปหานํ วูปสมํ ปฏินิสฺสคฺคํ ปฏิปฺปสฺสทฺธิํ อมตํ นิพฺพานนฺติ โทมนสฺสาน จูภยํ.}

\prose{\textbf{ถินสฺส จ ปนูทนนฺ}ติ \textbf{ถินนฺ}ติ ยา จิตฺตสฺส อกลฺยตา อกมฺมญฺญตา โอลียนา สลฺลียนา ลีนา ลียนา ลียิตตฺตํ ถินํ ถิยนา\csromansharedfootnote{343484cbc3da6448}{ถีนํ ถียนา (สฺยา)} ถิยิตตฺตํ จิตฺตสฺส. \textbf{ปนูทนนฺ}ติ ถินสฺส จ ปนูทนํ ปหานํ วูปสมํ ปฏินิสฺสคฺคํ ปฏิปฺปสฺสทฺธิํ อมตํ นิพฺพานนฺติ ถินสฺส จ ปนูทนํ.}

\prose{\textbf{กุกฺกุจฺจานํ นิวารณนฺ}ติ \textbf{กุกฺกุจฺจนฺ}ติ หตฺถ\-กุกฺกุจฺจมฺปิ กุกฺกุจฺจํ, ปาท\-กุกฺกุจฺจมฺปิ กุกฺกุจฺจํ, หตฺถปาท\-กุกฺกุจฺจมฺปิ กุกฺกุจฺจํ, อกปฺปิเย กปฺปิยสญฺญิตา, กปฺปิเย อกปฺปิย\-สญฺญิตา, อวชฺเช วชฺชสญฺญิตา, วชฺเช อวชฺชสญฺญิตา, ยํ เอวรูปํ กุกฺกุจฺจํ กุกฺกุจฺจายนา กุกฺกุจฺจายิ\-ตตฺตํ เจตโส วิปฺปฏิสาโร มโนวิเลโข, อิทํ วุจฺจติ กุกฺกุจฺจํ. อปิ จ ทฺวีหิ การเณหิ อุปฺปชฺชติ กุกฺกุจฺจํ เจตโส วิปฺปฏิสาโร มโนวิเลโข กตตฺตา จ อกตตฺตา จ. กถํ กตตฺตา จ อกตตฺตา จ อุปฺปชฺชติ กุกฺกุจฺจํ เจตโส วิปฺปฏิสาโร มโนวิเลโข, “กตํ เม กายทุจฺจริตํ, อกตํ เม กายสุจริตนฺ”ติ อุปฺปชฺชติ กุกฺกุจฺจํ เจตโส วิปฺปฏิสาโร \csromanfolio{158}\textbf{มโนวิเลโข, “กตํ เม วจีทุจฺจริตํ, อกตํ เม วจีสุจริตนฺ”ติ ฯเปฯ} “\textbf{กตํ เม มโนทุจฺจริตํ, อกตํ เม มโนสุจริตนฺ”ติ ฯเปฯ “กโต เม} ปาณาติปาโต, อกตา เม ปาณาติปาตา เวรมณี”ติ ฯเปฯ “กตํ เม อทินฺนาทานํ, อกตา เม อทินฺนาทานา เวรมณี”ติ ฯเปฯ “กโต เม กาเมสุ\-มิจฺฉาจาโร, อกตา เม กาเมสุ\-มิจฺฉาจารา เวรมณี”ติ ฯเปฯ “กโต เม มุสาวาโท, อกตา เม มุสาวาทา เวรมณี”ติ ฯเปฯ “กตา เม ปิสุณา วาจา\csromansharedfootnote{6a383bc90811e18d}{ปิสุณวาจา (ก)}, อกตา เม ปิสุณาย วาจาย เวรมณี”ติ ฯเปฯ “กตา เม ผรุสา วาจา, อกตา เม ผรุสาย วาจาย เวรมณี”ติ ฯเปฯ “กโต เม สมฺผปฺปลาโป, อกตา เม สมฺผปฺปลาปา เวรมณี”ติ ฯเปฯ “กตา เม อภิชฺฌา, อกตา เม อนภิชฺฌา”ติ ฯเปฯ “กโต เม พฺยาปาโท, อกโต เม อพฺยาปาโท”ติ ฯเปฯ “กตา เม มิจฺฉาทิฏฺฐิ, อกตา เม สมฺมาทิฏฺถิ”ติ อุปฺปชฺชติ กุกฺกุจฺจํ เจตโส วิปฺปฏิสาโร มโนวิเลโข, เอวํ กตตฺตา จ อกตตฺตา จ อุปฺปชฺชติ กุกฺกุจฺจํ เจตโส วิปฺปฏิสาโร มโนวิเลโข.}

\prose{อถ วา “สีเลสุ’มฺหิ อปริปูรการี”ติ อุปฺปชฺชติ กุกฺกุจฺจํ เจตโส วิปฺปฏิสาโร มโนวิเลโข, “อินฺทฺริเยสุ’มฺหิ อคุตฺตทฺวาโร”ติ. “โภชเน อมตฺตญฺญุมฺหี”ติ. “ชาคริยํ อนนุยุตฺโตมฺหี”ติ. “น สติสมฺปชญฺเญน สมนฺนาคโตมฺหี”ติ. “อภาวิตา เม จตฺตาโร สติปฏฺฐานา”ติ. “จตฺตาโร สมฺมปฺปธานา”ติ. “จตฺตาโร อิทฺธิปาทา”ติ. “ปญฺจินฺทฺริยานี”ติ. “ปญฺจ พลานี”ติ. “สตฺต โพชฺฌงฺคา”ติ. “อริโย อฏฺฐงฺคิโก มคฺโค”ติ. “ทุกฺขํ เม อปริญฺญาตํ. สมุทโย เม อปฺปหีโน. มคฺโค เม อภาวิโต. นิโรโธ เม อสจฺฉิกโต”ติ อุปฺปชฺชติ กุกฺกุจฺจํ เจตโส วิปฺปฏิสาโร มโนวิเลโข.}

\prose{\textbf{กุกฺกุจฺจานํ นิวารณนฺ}ติ กุกฺกุจฺจานํ อาวรณํ นีวรณํ ปหานํ อุปสมํ วูปสมํ ปฏินิสฺสคฺคํ ปฏิปฺปสฺสทฺธิํ อมตํ นิพฺพานนฺติ กุกฺกุจฺจานํ นิวารณํ. เตนาห ภควา\csromandash{}}

\vspace{\tipitakaparskip}
\csromancenter{ปหานํ กามจฺฉนฺทานํ, (อุทยาติ ภควา,)}
\prose{โทมนสฺสาน จูภยํ.}

\setlength{\csromangathapagewidth}{0pt}
\setlength{\csromangathawakwidth}{0pt}
\csromangathameasuregroup{ถินสฺส จ ปนูทนํ, \\ \textbf{อุเปกฺขา}\textbf{สติ}\textbf{สํสุทฺ}ธํ, \\ \textbf{อญฺญาวิโมกฺขํ ปพฺ}รูมิ,}{\csromangathabat{ถินสฺส จ ปนูทนํ,}{กุกฺกุจฺจานํ นิวารณนฺติ.} \\ \csromangathabat{\textbf{อุเปกฺขา}\textbf{สติ}\textbf{สํสุทฺ}ธํ,}{ธมฺมตกฺกปุเรชวํ.} \\ \csromangathabat{\textbf{อญฺญาวิโมกฺขํ ปพฺ}รูมิ,}{\textbf{อวิชฺชา}ย ปเภทนํ.}}
\csromangathasetleft{ถินสฺส จ ปนูทนํ, \\ \textbf{อุเปกฺขา}\textbf{สติ}\textbf{สํสุทฺ}ธํ, \\ \textbf{อญฺญาวิโมกฺขํ ปพฺ}รูมิ,}
\csromangathagroup{\csromangathastanza{\csromangathabat{ถินสฺส จ ปนูทนํ,}{กุกฺกุจฺจานํ นิวารณนฺติ.}}\csromangathastanza{\csromanfolio{159}\csromangathaitembat{76}{\textbf{อุเปกฺขา}\-\textbf{สติ}\-\textbf{สํสุทฺ}ธํ,}{ธมฺมตกฺก\-ปุเรชวํ.} \\ \csromangathacontbat{\textbf{อญฺญาวิโมกฺขํ ปพฺ}รูมิ,}{\textbf{อวิชฺชา}ย ปเภทนํ.}}}

\prose{\textbf{อุเปกฺขา}\-\textbf{สติ}\-\textbf{สํสุทฺธนฺ}ติ \textbf{อุเปกฺขา}ติ ยา จตุตฺเถ ฌาเน อุเปกฺขา อุเปกฺขนา อชฺฌุเปกฺขนา จิตฺตสมตา\csromansharedfootnote{7ed4acfd2c39efcd}{จิตฺตสมโถ (สฺยา) ขุ 7. 402 ปิฏฺเฐปิ.} จิตฺต\-ปฺปสฺสทฺธตา มชฺฌตฺตตา จิตฺตสฺส. \textbf{สตี}ติ ยา จตุตฺเถ ฌาเน อุเปกฺขํ อารพฺภสติ อนุสฺสติ ฯเปฯ สมฺมาสติ. \textbf{อุเปกฺขา}\-\textbf{สติ}\-\textbf{สํสุทฺธนฺ}ติ จตุตฺเถ ฌาเน อุเปกฺขา จ สติ จ สุทฺธา โหนฺติ วิสุทฺธา สํสุทฺธา ปริสุทฺธา ปริโยทาตา อนงฺคณา วิคตู\-ปกฺกิเลสา มุทุภูตา กมฺมนิยา ฐิตา อาเนญฺชปฺปตฺตาติ อุเปกฺขา\-สติ\-สํสุทฺธํ.}

\prose{\textbf{ธมฺมตกฺก}\-\textbf{ปุเรชวนฺ}ติ ธมฺมตกฺโก วุจฺจติ สมฺมาสงฺกปฺโป, โส อาทิโต โหติ, ปุรโต โหติ, ปุพฺพงฺคโม โหติ อญฺญาวิโมกฺขสฺสาติ เอวมฺปิ ธมฺมตกฺก\-ปุเรชวํ. อถ วา ธมฺมตกฺโก วุจฺจติ สมฺมาทิฏฺฐิ, สา อาทิโต โหติ, ปุรโต โหติ, ปุพฺพงฺคโม โหติ อญฺญาวิโมกฺขสฺสาติ เอวมฺปิ ธมฺมตกฺก\-ปุเรชวํ. อถ วา ธมฺมตกฺโก วุจฺจติ จตุนฺนํ มคฺคานํ ปุพฺพภาค\-วิปสฺสนา, สา อาทิโต โหติ, ปุรโต โหติ, ปุพฺพงฺคโม โหติ อญฺญาวิโมกฺขสฺสาติ เอวมฺปิ ธมฺมตกฺก\-ปุเรชวํ.}

\prose{\textbf{อญฺญาวิโมกฺขํ ปพฺรูมี}ติ อญฺญาวิโมกฺโข วุจฺจติ อรหตฺต\-วิโมกฺโข, อรหตฺต\-วิโมกฺขํ ปพฺรูมิ อาจิกฺขามิ เทเสมิ ปญฺญเปมิ ปฏฺฐเปมิ วิวรามิ วิภชามิ อุตฺตานีกโรมิ ปกาเสมีติ อญฺญาวิโมกฺขํ ปพฺรูมิ.}

\prose{\textbf{อวิชฺชาย ปเภทนนฺ}ติ \textbf{อวิชฺชา}ติ ทุกฺเข อญฺญาณํ ฯเปฯ อวิชฺชา โมโห อกุสลมูลํ. \textbf{ปเภทนนฺ}ติ อวิชฺชาย ปเภทนํ ปหานํ วูปสมํ ปฏินิสฺสคฺคํ ปฏิปฺปสฺสทฺธิํ อมตํ นิพฺพานนฺติ อวิชฺชาย ปเภทนํ. เตนาห ภควา\csromandash{}}

\setlength{\csromangathapagewidth}{0pt}
\setlength{\csromangathawakwidth}{0pt}
\csromangathameasuregroup{\textbf{อุเปกฺขา}\textbf{สติ}\textbf{สํสุทฺ}ธํ, \\ \textbf{อญฺญาวิโมกฺขํ ปพฺ}รูมิ, \\ \textbf{กิํสุ สํโยชโน ลฺ}โอโก, \\ \textbf{กิสฺสสฺส วิปฺปหานฺ}เอน,}{\csromangathabat{\textbf{อุเปกฺขา}\textbf{สติ}\textbf{สํสุทฺ}ธํ,}{ธมฺมตกฺกปุเรชวํ.} \\ \csromangathabat{\textbf{อญฺญาวิโมกฺขํ ปพฺ}รูมิ,}{\textbf{อวิชฺชา}ย ปเภทนนฺติ.} \\ \csromangathabat{\textbf{กิํสุ สํโยชโน ลฺ}โอโก,}{กิํสุ ตสฺส วิจารณํ.} \\ \csromangathabat{\textbf{กิสฺสสฺส วิปฺปหานฺ}เอน,}{นิพฺพานํ อิติ วุจฺจติ.}}
\csromangathasetleft{\textbf{อุเปกฺขา}\textbf{สติ}\textbf{สํสุทฺ}ธํ, \\ \textbf{อญฺญาวิโมกฺขํ ปพฺ}รูมิ, \\ \textbf{กิํสุ สํโยชโน ลฺ}โอโก, \\ \textbf{กิสฺสสฺส วิปฺปหานฺ}เอน,}
\csromangathagroup{\csromangathastanza{\csromangathaitembat{76}{\textbf{อุเปกฺขา}\-\textbf{สติ}\-\textbf{สํสุทฺ}ธํ,}{ธมฺมตกฺก\-ปุเรชวํ.} \\ \csromangathacontbat{\textbf{อญฺญาวิโมกฺขํ ปพฺ}รูมิ,}{\textbf{อวิชฺชา}ย ปเภทนนฺติ.}}\csromangathastanza{\csromangathaitembat{77}{\textbf{กิํสุ สํโยชโน ลฺ}โอโก,}{กิํสุ ตสฺส วิจารณํ.} \\ \csromangathacontbat{\textbf{กิสฺสสฺส วิปฺปหานฺ}เอน,}{นิพฺพานํ อิติ วุจฺจติ.}}}

\prose{\csromanfolio{160}\textbf{กิํสุ สํโยชโน โลโก}ติ โลกสฺส สํโยชนํ ลคฺคนํ พนฺธนํ อุปกฺกิเลโส, เกน โลโก ยุตฺโต ปยุตฺโต อายุตฺโต สมายุตฺโต ลคฺโค ลคฺคิโต ปลิพุทฺโธติ กิํสุ สํโยชโน โลโก.}

\prose{\textbf{กิํสุ ตสฺส วิจารณนฺ}ติ กิํสุ ตสฺส จารณํ วิจารณํ ปฏิวิจารณํ, เกน โลโก จรติ วิจรติ ปฏิวิจรตีติ กิํสุ ตสฺส วิจารณํ. \textbf{กิสฺสสฺส วิปฺปหาเนน, นิพฺพานํ อิติ วุจฺจตี}ติ กิสฺสสฺส วิปฺปหาเนน วูปสเมน ปฏินิสฺสคฺเคน ปฏิปฺปสฺสทฺธิยา นิพฺพานํ อิติ วุจฺจติ ปวุจฺจติ กถียติ ภณียติ ทีปียติ โวหรียตีติ กิสฺสสฺส วิปฺปหาเนน, นิพฺพานํ อิติ วุจฺจติ. เตนาห โส พฺราหฺมโณ\csromandash{}}

\setlength{\csromangathapagewidth}{0pt}
\setlength{\csromangathawakwidth}{0pt}
\csromangathameasuregroup{กิํสุ สํโยชโน โลโก, \\ กิสฺสสฺส วิปฺปหาเนน, \\ \textbf{นนฺทิสํโยชโน ลฺ}โอโก, \\ \textbf{ตณฺหาย วิปฺปหานฺ}เอน,}{\csromangathabat{กิํสุ สํโยชโน โลโก,}{กิํสุ ตสฺส วิจารณํ.} \\ \csromangathabat{กิสฺสสฺส วิปฺปหาเนน,}{นิพฺพานํ อิติ วุจฺจตีติ.} \\ \csromangathabat{\textbf{นนฺทิสํโยชโน ลฺ}โอโก,}{วิตกฺกสฺส วิจารณา.} \\ \csromangathabat{\textbf{ตณฺหาย วิปฺปหานฺ}เอน,}{นิพฺพานํ อิติ วุจฺจติ.}}
\csromangathasetleft{กิํสุ สํโยชโน โลโก, \\ กิสฺสสฺส วิปฺปหาเนน, \\ \textbf{นนฺทิสํโยชโน ลฺ}โอโก, \\ \textbf{ตณฺหาย วิปฺปหานฺ}เอน,}
\csromangathagroup{\csromangathastanza{\csromangathabat{กิํสุ สํโยชโน โลโก,}{กิํสุ ตสฺส วิจารณํ.} \\ \csromangathabat{กิสฺสสฺส วิปฺปหาเนน,}{นิพฺพานํ อิติ วุจฺจตีติ.}}\csromangathastanza{\csromangathaitembat{78}{\textbf{นนฺทิสํโยชโน ลฺ}โอโก,}{วิตกฺกสฺส วิจารณา.} \\ \csromangathacontbat{\textbf{ตณฺหาย วิปฺปหานฺ}เอน,}{นิพฺพานํ อิติ วุจฺจติ.}}}

\prose{\textbf{นนฺทิสํโยชโน โลโก}ติ นนฺที วุจฺจติ ตณฺหา. โย ราโค สาราโค ฯเปฯ อภิชฺฌา โลโภ อกุสลมูลํ, อยํ วุจฺจติ นนฺที. ยา นนฺที โลกสฺส สํโยชนํ ลคฺคนํ พนฺธนํ อุปกฺกิเลโส, อิมาย นนฺทิยา โลโก ยุตฺโต ปยุตฺโต อายุตฺโต สมายุตฺโต ลคฺโค ลคฺคิโต ปลิพุทฺโธติ นนฺทิสํโยชโน โลโก.}

\prose{\textbf{วิตกฺกสฺส วิจารณา}ติ \textbf{วิตกฺกา}ติ นว วิตกฺกา กามวิตกฺโก พฺยาปาทวิตกฺโก วิหิํสาวิตกฺโก ญาติวิตกฺโก ชนปท\-วิตกฺโก อมราวิตกฺโก ปรา\-นุทยตา\-ปฏิสํยุตฺโต วิตกฺโก ลาภ\-สกฺการ\-สิโลก\-ปฏิสํยุตฺโต วิตกฺโก อนวญฺญตฺติ\-ปฏิสํยุตฺโต วิตกฺโก, อิเม วุจฺจนฺติ นว วิตกฺกา. อิเม นว วิตกฺกา โลกสฺส จารณา วิจารณา ปฏิวิจารณา, อิเมหิ นวหิ วิตกฺเกหิ โลโก จรติ วิจรติ ปฏิวิจรตีติ วิตกฺกสฺส วิจารณา}

\prose{\textbf{ตณฺหาย วิปฺปหาเนน, นิพฺพานํ อิติ วุจฺจตี}ติ ตณฺหาติ รูป\textbf{ตณฺหา} ฯเปฯ ธมฺมตณฺหา. \textbf{ตณฺหาย วิปฺปหาเนน, นิพฺพานํ อิติ วุจฺจตี}ติ ตณฺหาย วิปฺปหาเนน วูปสเมน ปฏินิสฺสคฺเคน ปฏิปฺปสฺสทฺธิยา นิพฺพานํ อิติ วุจฺจติ ปวุจฺจติ กถียติ ภณียติ ทีปียติ โวหรียตีติ ตณฺหาย วิปฺปหาเนน, นิพฺพานํ อิติ วุจฺจติ. เตนาห ภควา\csromandash{}}

\setlength{\csromangathapagewidth}{0pt}
\setlength{\csromangathawakwidth}{0pt}
\csromangathameasuregroup{นนฺทิสํโยชโน โลโก, \\ ตณฺหาย วิปฺปหาเนน, \\ \textbf{กถํ สตสฺส จ}รโต, \\ \textbf{ภควนฺตํ ปุฏฺฐุ}มาคมา,}{\csromangathabat{นนฺทิสํโยชโน โลโก,}{วิ\textbf{ต}กฺกสฺส วิจารณา.} \\ \csromangathabat{ตณฺหาย วิปฺปหาเนน,}{นิพฺพานํ อิติ วุจฺจตีติ.} \\ \csromangathabat{\textbf{กถํ สตสฺส จ}รโต,}{วิญฺญาณํ อุปรุชฺฌติ.} \\ \csromangathabat{\textbf{ภควนฺตํ ปุฏฺฐุ}มาคมา,}{ตํ สุโณม วโจ ตว.}}
\csromangathasetleft{นนฺทิสํโยชโน โลโก, \\ ตณฺหาย วิปฺปหาเนน, \\ \textbf{กถํ สตสฺส จ}รโต, \\ \textbf{ภควนฺตํ ปุฏฺฐุ}มาคมา,}
\csromangathagroup{\csromangathastanza{\csromanfolio{161}\csromangathabat{นนฺทิสํโยชโน โลโก,}{วิ\textbf{ต}กฺกสฺส วิจารณา.} \\ \csromangathabat{ตณฺหาย วิปฺปหาเนน,}{นิพฺพานํ อิติ วุจฺจตีติ.}}\csromangathastanza{\csromangathaitembat{79}{\textbf{กถํ สตสฺส จ}รโต,}{วิญฺญาณํ อุปรุชฺฌติ.} \\ \csromangathacontbat{\textbf{ภควนฺตํ ปุฏฺฐุ}มาคมา,}{ตํ สุโณม วโจ ตว.}}}

\prose{\textbf{กถํ สตสฺส จรโต}ติ กถํ สตสฺส สมฺปชานสฺส จรโต วิหรโต อิริยโต วตฺตยโต ปาลยโต ยปยโต ยาปยโตติ กถํ สตสฺส จรโต.}

\prose{\textbf{วิญฺญาณํ อุปรุชฺฌตี}ติ วิญฺญาณํ นิรุชฺฌติ วูปสมฺมติ อตฺถํ คจฺฉติ ปฏิปฺปสฺสมฺภตีติ วิญฺญาณํ อุปรุชฺฌติ.}

\prose{\textbf{ภควนฺตํ ปุฏฺฐุมาคมา}ติ พุทฺธํ ภควนฺตํ ปุฏฺฐุํ ปุจฺฉิตุํ ยาจิตุํ อชฺเฌสิตุํ ปสาเทตุํ อาคมฺหา อาคตมฺหา อุปาคตมฺหา สมฺปตฺตมฺหา ตยา สทฺธิํ สมาคตมฺหาติ ภควนฺตํ ปุฏฺฐุมาคมา.}

\prose{\textbf{ตํ สุโณม วโจ ตวา}ติ ตนฺติ ตุยฺหํ วจนํ พฺยปฺปถํ เทสนํ อนุสาสนํ อนุสิฏฺฐํ สุโณม อุคฺคณฺหาม ธาเรม อุปธาเรม อุปลกฺเขมาติ ตํ สุโณม วโจ ตว. เตนาห โส พฺราหฺมโณ\csromandash{}}

\setlength{\csromangathapagewidth}{0pt}
\setlength{\csromangathawakwidth}{0pt}
\csromangathameasuregroup{\textbf{กถํ สตสฺส จ}รโต, \\ \textbf{ภควนฺตํ ปุฏฺฐุ}มาคมา, \\ \textbf{อชฺฌตฺตญฺจ พหิทฺทฺ}หา จ, \\ \textbf{เอวํ สตสฺส จรฺ}อโต,}{\csromangathabat{\textbf{กถํ สตสฺส จ}รโต,}{วิญฺญาณํ อุปรุชฺฌติ.} \\ \csromangathabat{\textbf{ภควนฺตํ ปุฏฺฐุ}มาคมา,}{ตํ สุโณม วโจ ตวาติ.} \\ \csromangathabat{\textbf{อชฺฌตฺตญฺจ พหิทฺทฺ}หา จ,}{เวทนํ นาภินนฺทโต.} \\ \csromangathabat{\textbf{เอวํ สตสฺส จรฺ}อโต,}{วิญฺญาณํ อุปรุชฺฌติ.}}
\csromangathasetleft{\textbf{กถํ สตสฺส จ}รโต, \\ \textbf{ภควนฺตํ ปุฏฺฐุ}มาคมา, \\ \textbf{อชฺฌตฺตญฺจ พหิทฺทฺ}หา จ, \\ \textbf{เอวํ สตสฺส จรฺ}อโต,}
\csromangathagroup{\csromangathastanza{\csromangathaitembat{79}{\textbf{กถํ สตสฺส จ}รโต,}{วิญฺญาณํ อุปรุชฺฌติ.} \\ \csromangathacontbat{\textbf{ภควนฺตํ ปุฏฺฐุ}มาคมา,}{ตํ สุโณม วโจ ตวาติ.}}\csromangathastanza{\csromangathaitembat{80}{\textbf{อชฺฌตฺตญฺจ พหิทฺทฺ}หา จ,}{เวทนํ นาภินนฺทโต.} \\ \csromangathacontbat{\textbf{เอวํ สตสฺส จรฺ}อโต,}{วิญฺญาณํ อุปรุชฺฌติ.}}}

\prose{\textbf{อชฺฌตฺตญฺจ พหิทฺธา จ, เวทนํ นาภินนฺทโต}ติ อชฺฌตฺตํ เวทนาสุ เวทนานุปสฺสี วิหรนฺโต เวทนํ นาภินนฺทติ นาภิวทติ น อชฺโฌเสติ\csromansharedfootnote{2ab06f3389135301}{น อชฺโฌสาย ติฏฺฐติ (สฺยา)}, อภินนฺทนํ อภิวทนํ อชฺโฌสานํ คาหํ ปรามาสํ อภินิเวสํ ปชหติ วิโนเทติ พฺยนฺตีกโรติ อนภาวํ คเมติ. พหิทฺธา เวทนาสุ เวทนานุปสฺสี วิหรนฺโต เวทนํ นาภินนฺทติ นาภิวทติ น อชฺโฌเสติ, อภินนฺทนํ อภิวทนํ อชฺโฌสานํ คาหํ ปรามาสํ อภินิเวสํ ปชหติ วิโนเทติ พฺยนฺตีกโรติ อนภาวํ คเมติ. อชฺฌตฺต\-พหิทฺธา เวทนาสุ เวทนานุปสฺสี วิหรนฺโต เวทนํ นาภินนฺทติ นาภิวทติ น อชฺโฌเสติ, อภินนฺทนํ \csromanfolio{162}อภิวทนํ อชฺโฌสานํ คาหํ ปรามาสํ อภินิเวสํ ปชหติ วิโนเทติ พฺยนฺตีกโรติ อนภาวํ คเมติ. อชฺฌตฺตํ สมุทยธมฺมา\-นุปสฺสี เวทนาสุ เวทนานุปสฺสี\csromansharedfootnote{ec6b59b1f1ff3a79}{อิทํ ปทํ นตฺถิ สฺยาโปตฺถเก.} วิหรนฺโต เวทนํ นาภินนฺทติ นาภิวทติ น อชฺโฌเสติ, อภินนฺทนํ อภิวทนํ อชฺโฌสานํ คาหํ ปรามฺสํ อภินิเวสํ ปชหติ วิโนเทติ พฺยนฺตีกโรติ อนภาวํ คเมติ. อชฺฌตฺตํ วยธมฺมา\-นุปสฺสี เวทนาสุ เวทนานุปสฺสี วิหรนฺโต ฯเปฯ. อชฺฌตฺตํ สมุทย\-วย\-ธมฺมานุปสฺสี เวทนาสุ เวทนานุปสฺสี วิหรนฺโต ฯเปฯ. พหิทฺธา สมุทยธมฺมา\-นุปสฺสี เวทนาสุ เวทนานุปสฺสี วิหรนฺโต เวทนํ นาภินนฺทติ นาภิวทติ น อชฺโฌเสติ, อภินนฺทนํ อภิวทนํ อชฺโฌสนํ คาหํ ปรามาสํ อภินิเวสํ ปชหติ วิโนเทติ พฺยนฺตีกโรติ อนภาวํ คเมติ. พหิทฺธา วยธมฺมา\-นุปสฺสี เวทนาสุ เวทนานุปสฺสี วิหรนฺโต ฯเปฯ. พหิทฺธา สมุทย\-วย\-ธมฺมานุปสฺสี เวทนาสุ เวทนานุปสฺสี วิหรนฺโต ฯเปฯ. อชฺฌตฺต\-พหิทฺธา สมุทยธมฺมา\-นุปสฺสี เวทนาสุ เวทนานุปสฺสี วิหรนฺโต ฯเปฯ. อชฺฌตฺต\-พหิทฺธา วยธมฺมา\-นุปสฺสี เวทนาสุ เวทนานุปสฺสี วิหรนฺโต ฯเปฯ. อชฺฌตฺต\-พหิทฺธา สมุทย\-วย\-ธมฺมานุปสฺสี เวทนาสุ เวทนานุปสฺสี วิหรนฺโต เวทนํ นาภินนฺทติ นาภิวทติ น อชฺโฌเสติ, อภินนฺทนํ อภิวทนํ อชฺโฌสานํ คาหํ ปรามาสํ อภินิเวสํ ปชหติ วิโนเทติ พฺยนฺตีกโรติ อนภาวํ คเมติ. อิเมหิ ทฺวาทสหิ อากาเรหิ เวทนาสุ เวทนานุปสฺสี วิหรนฺโต ฯเปฯ อนภาวํ คเมติ.}

\prose{อถ วา เวทนํ อนิจฺจโต ปสฺสนฺโต เวทนํ นาภินนฺทติ นาภิวทติ น อชฺโฌเสติ, อภินนฺทนํ อภิวทนํ อชฺโฌสานํ คาหํ ปรามาสํ อภินิเวสํ ปชหติ วิโนเทติ พฺยนฺตีกโรติ อนภาวํ คเมติ. เวทนํ ทุกฺขโต โรคโต คณฺฑโต สลฺลโต อฆโต อาพาธโต ฯเปฯ นิสฺสรณโต ปสฺสนฺโต เวทนํ นาภินนฺทติ นาภิวทติ น อชฺโฌเสติ, อภินนฺทนํ อภิวทนํ อชฺโฌสานํ คาหํ ปรามาสํ อภินิเวสํ ปชหติ วิโนเทติ พฺยนฺตีกโรติ อนภาวํ คเมติ. อิเมหิ จตฺตาลีสาย\csromansharedfootnote{63813da74b6de580}{ทฺวาจตฺตาฬีสาย (สฺยา)} อากาเรหิ เวทนาสุ เวทนานุปสฺสี วิหรนฺโต เวทนํ นาภินนฺทติ นาภิวทติ น อชฺโฌเสติ, อภินนฺทนํ อภิวทนํ อชฺโฌสานํ คาหํ ปรามาสํ อภินิเวสํ ปชหติ วิโนเทติ พฺยนฺตีกโรติ อนภาวํ คเมตีติ อชฺฌตฺตญฺจ พหิทฺธา จ, เวทนํ นาภินนฺทโต.}

\csromanmatikaanchor{mtk.35}
\csromantocmarkref{subsection}{14. โปสาลมาณวปุจฺฉานิทฺเทส}{mtk.35}
\bookmark[dest={mtk.35},level=2]{14. โปสาลมาณวปุจฺฉานิทฺเทส}
\prose{\csromanfolio{163}\textbf{เอวํ สตสฺส จรโต}ติ เอวํ สตสฺส สมฺปชานสฺส จรโต วิหรโต อิริยโต วตฺตยโต ปาลยโต ยปยโต ยาปยโตติ เอวํ สตสฺส จรโต.}

\prose{\textbf{วิญฺญาณํ อุปรุชฺฌตี}ติ ปุญฺญาภิสงฺขาร\-สหคตํ วิญฺญาณํ อปุญฺญาภิสงฺขาร\-สหคตํ วิญฺญาณํ อาเนญฺชาภิสงฺขาร\-สหคตํ วิญฺญาณํ นิรุชฺฌติ วูปสมฺมติ อตฺถํ คจฺฉติ ปฏิปฺปสฺสมฺภตีติ วิญฺญาณํ อุปรุชฺฌติ. เตนาห ภควา\csromandash{}}

\setlength{\csromangathapagewidth}{0pt}
\setlength{\csromangathawakwidth}{0pt}
\csromangathameasuregroup{อชฺฌตฺตญฺจ พหิทฺธา จ, \\ เอวํ สตสฺส จรโต,}{\csromangathabat{อชฺฌตฺตญฺจ พหิทฺธา จ,}{เวทนํ นาภินนฺทโต.} \\ \csromangathabat{เอวํ สตสฺส จรโต,}{วิญฺญาณํ อุปรุชฺฌตีติ.}}
\csromangathasetleft{อชฺฌตฺตญฺจ พหิทฺธา จ, \\ เอวํ สตสฺส จรโต,}
\csromangathagroup{\csromangathastanza{\csromangathabat{อชฺฌตฺตญฺจ พหิทฺธา จ,}{เวทนํ นาภินนฺทโต.} \\ \csromangathabat{เอวํ สตสฺส จรโต,}{วิญฺญาณํ อุปรุชฺฌตีติ.}}}

\prose{สห คาถา\-ปริโยสานา ฯเปฯ “สตฺถา เม ภนฺเต ภควา, สาวโกหมสฺมี”ติ.}

\vspace{\tipitakaparskip}
\csromancenter{อุทย\-มาณว\-ปุจฺฉานิทฺเทโส เตรสโม.}
\csromanhangingbegin
\hangingheaditem{81}{\textbf{โย อตีตํ อาทิสติ, (อิจฺจายสฺมา โปสาโล,)}}
\hangingbody{\textbf{อเนโช ฉินฺนสํสโย.}}
\hangingbody{\textbf{ปารคุํ1} สพฺพธมฺมานํ, อตฺถิ ปเญฺหน อาคมํ.}
\csromanhangingend

\prose{\textbf{โย อตีตํ อาทิสตี}ติ \textbf{โย}ติ โย โส ภควา สยมฺภู อนาจริยโก ปุพฺเพ อนนุสฺสุเตสุ ธมฺเมสุ สามํ สจฺจานิ อภิสมฺพุชฺฌิ, ตตฺถ จ สพฺพญฺญุตํ ปตฺโต, พเลสุ จ วสีภาวํ. \textbf{อตีตํ อาทิสตี}ติ ภควา อตฺตโน จ ปเรสญฺจ อตีตมฺปิ อาทิสติ, อนาคตมฺปิ อาทิสติ, ปจฺจุปฺปนฺนมฺปิ อาทิสติ.}

\prose{กถํ ภควา อตฺตโน อตีตํ อาทิสติ, ภควา อตฺตโน อตีตํ เอกมฺปิ ชาติํ อาทิสติ, เทฺวปิ ชาติโย อาทิสติ, ติสฺโสปิ ชาติโย อาทิสติ, จตสฺโสปิ ชาติโย อาทิสติ, ปญฺจปิ ชาติโย อาทิสติ, ทสปิ ชาติโย อาทิสติ, วีสมฺปิ ชาติโย อาทิสติ, ติํสมฺปิ ชาติโย อาทิสติ, จตฺตาลีสมฺปิ ชาติโย \csromanfolio{164}อาทิสติ, ปญฺญาสมฺปิ ชาติโย อาทิสติ, ชาติสตมฺปิ. ชาติสหสฺสมฺปิ. ชาติ\-สต\-สหสฺสมฺปิ. อเนเกปิ สํวฏฺฏกปฺเป. อเนเกปิ วิวฏฺฏกปฺเป. อเนเกปิ สํวฏฺฏ\-วิวฏฺฏ\-กปฺเป อาทิสติ, อมุตฺราสิํ เอวํนาโม เอวํโคตฺโต เอวํวณฺโณ เอวมาหาโร เอวํ\-สุขทุกฺข\-ปฺปฏิสํเวที เอวมา\-ยุปริยนฺโต โส ตโต จุโต อมุตฺร อุทปาทิํ, ตตฺราปาสิํ “เอวํนาโม เอวํโคตฺโต เอวํวณฺโณ เอวมาหาโร เอวํ\-สุขทุกฺข\-ปฺปฏิสํเวที เอวมา\-ยุปริยนฺโต, โส ตโต จุโต อิธูปปนฺโน”ติ อิติ สาการํ สอุทฺเทสํ อเนกวิหิตํ ปุพฺเพนิวาสํ อาทิสติ. เอวํ ภควา อตฺตโน อตีตํ อาทิสติ.}

\prose{กถํ ภควา ปเรสํ อตีตํ อาทิสติ, ภควา ปเรสํ อตีตํ เอกมฺปิ ชาติํ อาทิสติ, เทฺวปิ ชาติโย อาทิสติ ฯเปฯ. อเนเกปิ สํวฏฺฏ\-วิวฏฺฏ\-กปฺเป อาทิสติ, อมุตฺราสิ เอวํนาโม เอวํโคตฺโต เอวํวณฺโณ เอวมาหาโร เอวํ\-สุขทุกฺข\-ปฺปฏิสํเวที เอวมา\-ยุปริยนฺโต, โส ตโต จุโต อมุตฺร อุทปาทิ, ตตฺราปาสิ “เอวํนาโม เอวํโคตฺโต เอวํวณฺโณ เอวมาหาโร เอวํ\-สุขทุกฺข\-ปฺปฏิสํเวที เอวมา\-ยุปริยนฺโต, โส ตโต จุโต อิธูปปนฺโน”ติ อิติ สาการํ สอุทฺเทสํ อเนกวิหิตํ ปุพฺเพนิวาสํ อาทิสติ. เอวํ ภควา ปเรสํ อตีตํ อาทิสติ.}

\prose{ภควา ปญฺจ ชาตกสตานิ ภาสนฺโต อตฺตโน จ ปเรสญฺจ อตีตํ อาทิสติ, มหาปทา\-นิย\-สุตฺตนฺตํ\csromansharedfootnote{34bf7543d7ce1433}{มหาธนิยสุตฺตํ (สฺยา)} ภาสนฺโต อตฺตโน จ ปเรสญฺจ อตีตํ อาทิสติ, มหา\-สุทสฺส\-นิย\-สุตฺตนฺตํ ภาสนฺโต อตฺตโน จ ปเรสญฺจ อตีตํ อาทิสติ, มหา\-โควินฺทิย\-สุตฺตนฺตํ ภาสนฺโต อตฺตโน จ ปเรสญฺจ อตีตํ อาทิสติ, มาฆเทวิยสุตฺตนฺตํ ภาสนฺโต อตฺตโน จ ปเรสญฺจ อาทิสติ.}

\prose{วุตฺตํ เหตํ ภควตา\csromandash{}อตีตํ โข จุนฺท อทฺธานํ อารพฺภ ตถาคตสฺส สตานุสาริ\-ญาณํ\csromansharedfootnote{a1c38a75ce9a6211}{สตานุสฺสริยญาณํ (ก) ปสฺส ที 3. 111 ปิฏฺเฐ.} โหติ, โส ยาวตกํ อากงฺขติ, ตาวตกํ อนุสฺสรติ. อนาคตญฺจ โข จุนฺท. ปจฺจุปฺปนฺนญฺจ โข จุนฺท อทฺธานํ อารพฺภ ตถาคตสฺส โพธิชํ ญาณํ อุปฺปชฺชติ “อยมนฺติมา ชาติ, นตฺถิทานิ ปุนพฺภโว”ติ.}

\prose{\csromanfolio{165}อินฺทฺริยปโรปริยตฺต\-ญาณํ\csromansharedfootnote{230888062fcddfb9}{อินฺทฺริยปโรปริยตฺติญาณํ (ก) อฏฺฐกถา โอโลเกตพฺพา.} ตถาคตสฺส ตถาคตพลํ, สตฺตานํ อาสยา\-นุสย\-ญาณํ ตถาคตสฺส ตถาคตพลํ, ยมกปาฏิหีเร ญาณํ\csromansharedfootnote{327fa50734e25ba4}{ยมกปาฏิหิริยญาณํ (สฺยา)} ตถาคตสฺส ตถาคตพลํ, มหากรุณา\-สมาปตฺติยา ญาณํ ตถาคตสฺส ตถาคตพลํ, สพฺพญฺญุต\-ญาณํ ตถาคตสฺส ตถาคตพลํ, อนาวรณญาณํ ตถาคตสฺส ตถาคตพลํ, สพฺพตฺถ อสงฺคม\-ปฺปฏิหตม\-นาวรณญาณํ ตถาคตสฺส ตหาคตพลํ. เอวํ ภควา อตฺตโน จ ปเรสญฺจ อตีตมฺปิ อาทิสติ, อนาคตมฺปิ อาทิสติ, ปจฺจุปฺปนฺนมฺปิ อาทิสติ อาจิกฺขติ เทเสติ ปญฺญเปติ ปฏฺฐเปติ วิวรติ วิภชติ อุตฺตานีกโรติ ปกาเสตีติ โย อตีตํ อาทิสติ.}

\prose{\textbf{อิจฺจายสฺมา โปสาโล}ติ \textbf{อิจฺจา}ติ ปทสนฺธิ ฯเปฯ. อา\textbf{ยสฺมา}ติ ปิยวจนํ ฯเปฯ. \textbf{โปสาโล}ติ ตสฺส พฺราหฺมณสฺส นามํ ฯเปฯ อภิลาโปติ อิจฺจายสฺมา โปสาโล.}

\prose{\textbf{อเนโช ฉินฺน}\-\textbf{สํสโย}ติ เอชา วุจฺจติ ตณฺหา. โย ราโค สาราโค ฯเปฯ อภิชฺฌา โลโภ อกุสลมูลํ. สา เอชา ตณฺหา พุทฺธสฺส ภควโต ปหีนา อุจฺฉินฺนมูลา ตาลาวตฺถุกตา อนภาวํกตา อายติํ อนุปฺปาทธมฺมา, ตสฺมา พุทฺโธ อเนโช. เอชาย ปหีนตฺตา อเนโช, ภควา ลาเภปิ น อิญฺชติ ฯเปฯ ทุกฺเขปิ น อิญฺชติ น จลติ น เวธติ นปฺปเวธติ น สมฺปเวธตีติ อเนโช. \textbf{ฉินฺน}\-\textbf{สํสโย}ติ สํสโย วุจฺจติ วิจิกิจฺฉา. ทุกฺเข กงฺขา ฯเปฯ ฉมฺภิตตฺตํ จิตฺตสฺส มโนวิเลโข, โส สํสโย พุทฺธสฺส ภควโต ปหีโน ฉินฺโน อุจฺฉินฺโน สมุจฺฉินฺโน วูปสนฺโต ปฏินิสฺสคฺโค ปฏิปฺปสฺสทฺโธ อภพฺพุ\-ปฺปตฺติโก ญาณคฺคินา ทฑฺโฒ, ตสฺมา พุทฺโธ ฉินฺน\-สํสโยติ อเนโช ฉินฺนสํสโย.}

\prose{\textbf{ปารคุํ สพฺพ}\-\textbf{ธมฺมานนฺ}ติ ภควา สพฺพธมฺมานํ อภิญฺญาปารคู ปริญฺญาปารคู ปหานปารคู ภาวนาปารคู สจฺฉิกิริยา\-ปารคู สมาปตฺติ\-ปารคู. อภิญฺญาปารคู สพฺพธมฺมานํ ฯเปฯ ชาติ\-มรณ\-สํสาโร, นตฺถิ ตสฺส ปุนพฺภโวติ ปารคู สพฺพธมฺมานํ.}

\prose{\textbf{อตฺถิ ปเญฺหน อาคมนฺ}ติ ปเญฺหน อตฺถิโก อาคโตมฺหิ ฯเปฯ วหสฺเสตํ ภารนฺติ อตฺถิ ปเญฺหน อาคมํ. เตนาห โส พฺราหฺมโณ\csromandash{}}

\prose{\csromanfolio{166}โย อตีตํ อาทิสติ, (อิจฺจายสฺมา โปสาโล,)}

\prose{อเนโช ฉินฺนสํสโย.}

\setlength{\csromangathapagewidth}{0pt}
\setlength{\csromangathawakwidth}{0pt}
\csromangathameasuregroup{ปารคุํ สพฺพธมฺมานํ, \\ \textbf{วิภูต}\textbf{รูป}\textbf{สญฺญิ}สฺส, \\ \textbf{อชฺฌตฺตญฺจ พหิทฺธฺ}อา จ, \\ \textbf{ญาณํ สกฺกา’นุปุจฺฉฺ}อามิ,}{\csromangathabat{ปารคุํ สพฺพธมฺมานํ,}{อตฺถิ ปเญฺหน อาคมนฺติ.} \\ \csromangathabat{\textbf{วิภูต}\textbf{รูป}\textbf{สญฺญิ}สฺส,}{สพฺพกายปฺปหายิโน.} \\ \csromangathabat{\textbf{อชฺฌตฺตญฺจ พหิทฺธฺ}อา จ,}{นตฺถิ กิญฺจีติ ปสฺสโต.} \\ \csromangathabat{\textbf{ญาณํ สกฺกา’นุปุจฺฉฺ}อามิ,}{กถํ เนยฺโย ตถาวิโธ.}}
\csromangathasetleft{ปารคุํ สพฺพธมฺมานํ, \\ \textbf{วิภูต}\textbf{รูป}\textbf{สญฺญิ}สฺส, \\ \textbf{อชฺฌตฺตญฺจ พหิทฺธฺ}อา จ, \\ \textbf{ญาณํ สกฺกา’นุปุจฺฉฺ}อามิ,}
\csromangathagroup{\csromangathastanza{\csromangathabat{ปารคุํ สพฺพธมฺมานํ,}{อตฺถิ ปเญฺหน อาคมนฺติ.}}\csromangathastanza{\csromangathaitembat{82}{\textbf{วิภูต}\-\textbf{รูป}\-\textbf{สญฺญิ}สฺส,}{สพฺพกาย\-ปฺปหายิโน.} \\ \csromangathacontbat{\textbf{อชฺฌตฺตญฺจ พหิทฺธฺ}อา จ,}{นตฺถิ กิญฺจีติ ปสฺสโต.} \\ \csromangathacontbat{\textbf{ญาณํ สกฺกา’นุปุจฺฉฺ}อามิ,}{กถํ เนยฺโย ตถาวิโธ.}}}

\prose{\textbf{วิภูต}\-\textbf{รูป}\-\textbf{สญฺญิสฺสา}ติ กตมา รูปสญฺญา, รูปาวจร\-สมาปตฺติํ สมาปนฺนสฺส วา อุปปนฺนสฺส วา ทิฏฺฐธมฺม\-สุขวิหาริสฺส วา สญฺญา สญฺชานนา สญฺชานิตตฺตํ, อยํ รูปสญฺญา. \textbf{วิภูต}\-\textbf{รูป}\-\textbf{สญฺญิสฺสา}ติ จตสฺโส อรูป\-สมาปตฺติโย ปฏิลทฺธสฺส\csromansharedfootnote{635090a9f16ebf99}{ลาภิสฺส (สฺยา)} รูปสญฺญา วิภูตา โหนฺติ วิคตา อติกฺกนฺตา สมติกฺกนฺตา วีติวตฺตาติ วิภูต\-รูป\-สญฺญิสฺส.}

\prose{\textbf{สพฺพกาย}\-\textbf{ปฺปหายิโน}ติ สพฺโพ ตสฺส ปฏิสนฺธิโก รูปกาโย ปหีโน, ตทงฺค\-สมติกฺกมา วิกฺขมฺภน\-ปฺปหาเนน ปหีโน ตสฺส รูปกาโยติ สพฺพกาย\-ปฺปหายิโน.}

\prose{\textbf{อชฺฌตฺตญฺจ พหิทฺธา จ, นตฺถิ กิญฺจีติ ปสฺสโตติ นตฺถิ กิญฺจี}ติ อากิญฺจญฺญายตน\-สมาปตฺติ. กิํ การณา “นตฺถิ กิญฺจี”ติ อากิญฺจญฺญายตน\-สมาปตฺติ. ยํ วิญฺญาณญฺจา\-ยตนสมาปตฺติํ สโต สมาปชฺชิตฺวา ตโต วุฏฺฐหิตฺวา ตญฺเญว วิญฺญาณํ อภาเวติ วิภาเวติ อนฺตรธาเปติ, “นตฺถิ กิญฺจี”ติ ปสฺสติ, ตํการณา “นตฺถิ กิญฺจี”ติ อากิญฺจญฺญายตนสมาปตฺตีติ อชฺฌตฺตญฺจ พหิทฺธา จ, นตฺถิ กิญฺจีติ ปสฺสโต.}

\prose{\textbf{ญาณํ สกฺกา’นุปุจฺฉามี}ติ สกฺกาติ สกฺโก, ภควา สกฺยกุลา ปพฺพชิโตติปิ สกฺโก ฯเปฯ ปหีน\-ภยเภรโว วิคต\-โลมหํโสติปิ สกฺโก. \textbf{ญาณํ สกฺกา’นุปุจฺฉามี}ติ ตสฺส ญาณํ ปุจฺฉามิ, ปญฺญํ ปุจฺฉามิ, สมฺพุทฺธํ ปุจฺฉามิ. กีทิสํ กิํ สณฺฐิตํ กิํปการํ กิํปฏิภาคํ ญาณํ อิจฺฉิตพฺพนฺติ ญาณํ สกฺกา’นุปุจฺฉามิ.}

\prose{\textbf{กถํ เนยฺโย ตถาวิโธ}ติ กถํ โส เนตพฺโพ วิเนตพฺโพ อนุเนตพฺโพ ปญฺญาเปตพฺโพ นิชฺฌาเปตพฺโพ เปกฺเขตพฺโพ ปสาเทตพฺโพ, กถํ เตน\csromansharedfootnote{29c842e194aa0a63}{กถมสฺส (สฺยา)} อุตฺตริ ญาณํ อุปฺปาเทตพฺพํ. \textbf{ตถาวิโธ}ติ ตถาวิโธ \csromanfolio{167}ตาทิโส ตสฺสณฺฐิโต ตปฺปกาโร ตปฺปฏิภาโค, โย โส อากิญฺจญฺญายตนสมาปตฺติ\-ลาภีติ กถํ เนยฺโย ตถาวิโธ. เตนาห โส พฺราหฺมโณ\csromandash{}}

\setlength{\csromangathapagewidth}{0pt}
\setlength{\csromangathawakwidth}{0pt}
\csromangathameasuregroup{วิภูตรูปสญฺญิสฺส, \\ อชฺฌตฺตญฺจ พหิทฺธา จ, \\ ญาณํ สกฺกา’นุปุจฺฉามิ,}{\csromangathabat{วิภูตรูปสญฺญิสฺส,}{สพฺพกายปฺปหายิโน.} \\ \csromangathabat{อชฺฌตฺตญฺจ พหิทฺธา จ,}{นตฺถิ กิญฺจีติ ปสฺสโต.} \\ \csromangathabat{ญาณํ สกฺกา’นุปุจฺฉามิ,}{กถํ เนยฺโย ตถาวิโธติ.}}
\csromangathasetleft{วิภูตรูปสญฺญิสฺส, \\ อชฺฌตฺตญฺจ พหิทฺธา จ, \\ ญาณํ สกฺกา’นุปุจฺฉามิ,}
\csromangathagroup{\csromangathastanza{\csromangathabat{วิภูต\-รูป\-สญฺญิสฺส,}{สพฺพกาย\-ปฺปหายิโน.} \\ \csromangathabat{อชฺฌตฺตญฺจ พหิทฺธา จ,}{นตฺถิ กิญฺจีติ ปสฺสโต.} \\ \csromangathabat{ญาณํ สกฺกา’นุปุจฺฉามิ,}{กถํ เนยฺโย ตถาวิโธติ.}}}

\csromanhangingbegin
\hangingheaditem{83}{\textbf{วิญฺญาณ}\-\textbf{ฏฺฐิติโย สพฺพา, (โปสาลาติ ภควา,)}}
\hangingbody{\textbf{อภิชานํ ตถาคโต.}}
\hangingbody{\textbf{ติฏฺฐนฺตเมนํ ชานาติ, ธิมุตฺตํ ตปฺปรายณํ.}}
\csromanhangingend

\prose{\textbf{วิญฺญาณ}\-\textbf{ฏฺฐิติโย สพฺพา}ติ ภควา อภิสงฺขาร\-วเสน จตสฺโส วิญฺญาณ\-ฏฺฐิติโย ชานาติ, ปฏิสนฺธิ\-วเสน สตฺต วิญฺญาณ\-ฏฺฐิติโย ชานาติ. กถํ ภควา อภิสงฺขาร\-วเสน จตสฺโส วิญฺญาณ\-ฏฺฐิติโย ชานาติ, วุตฺตํ เหตํ ภควตา\csromandash{}รูปุปยํ\csromansharedfootnote{7fd7b9f620cd927c}{รูปูปายํ (สฺยา, ก) ปสฺส สํ 2. 43 ปิฏฺเฐ.} วา ภิกฺขเว วิญฺญาณํ ติฏฺฐมานํ ติฏฺเฐยฺย\csromansharedfootnote{d65e5455d229ed17}{ติฏฺฐติ (สฺยา, ก)}, รูปารมฺมณํ รูปปฺปติฏฺฐํ นนฺทูปเสจนํ วุทฺธิํ วิรูฬฺหิํ เวปุลฺลํ อาปชฺเชยฺย. เวทนุปยํ วา ภิกฺขเว ฯเปฯ. สญฺญุปยํ วา ภิกฺขเว ฯเปฯ. สงฺขารุปยํ วา ภิกฺขเว วิญฺญาณํ ติฏฺฐมานํ ติฏฺเฐยฺย, สงฺขารา\-รมฺมณํ สงฺขาร\-ปฺปติฏฺฐํ นนฺทูปเสจนํ วุทฺธิํ วิรูฬฺหิํ เวปุลฺลํ อาปชฺเชยฺยา”ติ, เอวํ ภควา อภิสงฺขาร\-วเสน จตสฺโส วิญฺญาณ\-ฏฺฐิติโย ชานาติ.}

\prose{กถํ ภควา ปฏิสนฺธิ\-วเสน สตฺต วิญฺญาณ\-ฏฺฐิติโย ชานาติ, วุตฺตํ เหตํ ภควตา\csromandash{}สนฺติ ภิกฺขเว สตฺตา นานตฺตกายา นานตฺตสญฺญิโน, เสยฺยถาปิ มนุสฺสา เอกจฺเจ จ เทวา เอกจฺเจ จ วินิปาติกา, อยํ ปฐมา วิญฺญาณฏฺฐิติ.}

\prose{สนฺติ ภิกฺขเว สตฺตา นานตฺตกายา เอกตฺตสญฺญิโน, เสยฺยถาปิ เทวา พฺรหฺมกายิกา ปฐมา\-ภินิพฺพตฺตา, อยํ ทุติยา วิญฺญาณฏฺฐิติ.}

\prose{สนฺติ ภิกฺขเว สตฺตา เอกตฺตกายา นานตฺตสญฺญิโน, เสยฺยถาปิ เทวา อาภสฺสรา, อยํ ตติยา วิญฺญาณฏฺฐิติ.}

\prose{\csromanfolio{168}สนฺติ ภิกฺขเว สตฺตา เอกตฺตกายา เอกตฺตสญฺญิโน, เสยฺยถาปิ เทวา สุภกิณฺหา, อยํ จตุตฺถี\csromansharedfootnote{cecc9302dfe959ea}{จตุตฺถา (สฺยา) ปสฺส อํ 2. 428 ปิฏฺเฐ.} วิญฺญาณฏฺฐิติ.}

\prose{สนฺติ ภิกฺขเว สตฺตา สพฺพโส รูปสญฺญานํ สมติกฺกมา ปฏิฆ\-สญฺญานํ อตฺถงฺคมา นานตฺต\-สญฺญานํ อมนสิการา “อนนฺโต อากาโส”ติ อากาสานญฺจายตนู\-ปคา, อยํ ปญฺจมี\csromansharedfootnote{b0019c4f847490ca}{ปญฺจมา (สฺยา)} วิญฺญาณฏฺฐิติ.}

\prose{สนฺติ ภิกฺขเว สตฺตา สพฺพโส อากาสา\-นญฺจา\-ยตนํ สมติกฺกมฺม “อนนฺตํ วิญฺญาณนฺ”ติ วิญฺญาณญฺจายตนู\-ปคา, อยํ ฉฏฺฐี\csromansharedfootnote{f762e42250cbe94a}{ฉฏฺฐา (สฺยา)} วิญฺญาณฏฺฐิติ.}

\prose{สนฺติ ภิกฺขเว สตฺตา สพฺพโส วิญฺญาณญฺจา\-ยตนํ สมติกฺกมฺม “นตฺถิ กิญฺจี”ติ อากิญฺจญฺญายตนู\-ปคา, อยํ สตฺตมี\csromansharedfootnote{8011e7e0903f3c6e}{สตฺตมา (สฺยา)} วิญฺญาณฏฺฐิติ. เอวํ ภควา ปฏิสนฺธิ\-วเสน สตฺต วิญฺญาณ\-ฏฺฐิติโย ชานาตีติ วิญฺญาณ\-ฏฺฐิติโย สพฺพา.}

\prose{\textbf{โปสาลาติ ภควา}ติ โปสาลาติ ภควา ตํ พฺราหฺมณํ นาเมน อาลปติ. \textbf{ภควา}ติ คารวาธิวจนเมตํ ฯเปฯ สจฺฉิกา ปญฺญตฺติ, ยทิทํ ภควาติ โปสาลาติ ภควา.}

\prose{\textbf{อภิชานํ ตถาคโต}ติ \textbf{อภิชานนฺ}ติ อภิชานนฺโต วิชานนฺโต ปฏิวิชานนฺโต ปฏิวิชฺฌนฺโต ตถาคโต. วุตฺตํ เหตํ ภควตา\csromansharedfootnote{88235f8e045116b4}{ปสฺส ที 3. 111 ปิฏฺเฐ.} อตีตญฺเจปิ โข จุนฺท โหติ อภูตํ อตจฺฉํ อนตฺถสญฺหิตํ, น ตํ ตถาคโต พฺยากโรติ. อตีตญฺเจปิ จุนฺท โหติ ภูตํ ตจฺฉํ อนตฺถสญฺหิตํ, ตมฺปิ ตถาคโต น พฺยากโรติ. อตีตญฺเจปิ โข จุนฺท โหติ ภูตํ ตจฺฉํ อตฺถสญฺหิตํ, ตตฺร กาลญฺญู ตถาคโต โหติ ตสฺเสว ปญฺหสฺส เวยฺยากรณาย. อนาคตญฺเจปิ จุนฺท โหติ ฯเปฯ. ปจฺจุปฺปนฺนญฺเจปิ จุนฺท โหติ อภูตํ อตจฺฉํ อนตฺถสญฺหิตํ, น ตํ ตถาคโต พฺยากโรติ. ปจฺจุปฺปนฺนํ เจปิ จุนฺท โหติ ภูตํ ตจฺฉํ อนตฺถสญฺหิตํ, ตมฺปิ ตถาคโต น พฺยากโรติ. ปจฺจุปฺปนฺนญฺเจปิ จุนฺท โหติ ภูตํ ตจฺฉํ อตฺถสญฺหิตํ, ตตฺร กาลญฺญู ตถาคโต โหติ ตสฺส ปญฺหสฺส เวยฺยากรณาย. อิติ โข จุนฺท อตีตา\-นาคต\-ปจฺจุปฺปนฺเนสุ ธมฺเมสุ ตถาคโต กาลวาที ภูตวาที อตฺถวาที ธมฺมวาที วินยวาที, ตสฺมา “ตถาคโต”ติ วุจฺจติ.}

\prose{ยํ โข จุนฺท สเทวกสฺส โลกสฺส สมารกสฺส สพฺรหฺมกสฺส สสฺสมณ\-พฺราหฺมณิยา ปชาย สเทว\-มนุสฺสาย ทิฏฺฐํ สุตํ มุตํ วิญฺญาตํ ปตฺตํ \csromanfolio{169}ปริเยสิตํ อนุวิจริตํ มนสา, สพฺพํ ตํ ตถาคเตน อภิสมฺพุทฺธํ, ตสฺมา “ตถาคโต”ติ วุจฺจติ. ยญฺจ จุนฺท รตฺติํ ตถาคโต อนุตฺตรํ สมฺมาสมฺโพธิํ อภิสมฺ\-พุชฺฌติ, ยญฺจ รตฺติํ อนุปาทิเสสาย นิพฺพานธาตุยา ปรินิพฺพายติ, ยํ เอตสฺมิํ อนฺตเร ภาสติ ลปติ นิทฺทิสติ. สพฺพํ ตํ ตเถว โหติ โน อญฺญถา, ตสฺมา “ตถาคโต”ติ วุจฺจติ. ยถาวาที จุนฺท ตถาคโต ตถาการี, ยถาการี ตถาวาที, อิติ ยถาวาที ตถาการี, ยถาการี ตถาวาที, ตสฺมา “ตถาคโต”ติ วุจฺจติ. สเทวเก จุนฺท โลเก สมารเก สพฺรหฺมเก สสฺสมณ\-พฺราหฺมณิยา ปชาย สเทว\-มนุสฺสาย ตถาคโต อภิภู อนภิภูโต อญฺญทตฺถุทโส วสวตฺถี, ตสฺมา “ตถาคโต”ติ วุจฺจตีติ อภิชานํ ตถาคโต.}

\prose{\textbf{ติฏฺฐนฺตเมนํ ชานาตี}ติ ภควา อิธตฺถญฺเญว\csromansharedfootnote{6e028f37216226a6}{อิธฏฺฐญฺเญว (สฺยา)} ชานาติ กมฺมา\-ภิสงฺขาร\-วเสน “อยํ ปุคฺคโล กายสฺส เภทา ปรํ มรณา อปายํ ทุคฺคติํ วินิปาตํ นิรยํ อุปปชฺชิสฺสตี”ติ. ภควา อิธตฺถญฺเญว ชานาติ กมฺมา\-ภิสงฺขาร\-วเสน “อยํ ปุคฺคโล กายสฺส เภทา ปรํ มรณา ติรจฺฉาน\-โยนิํ อุปปชฺชิสฺสตี”ติ. ภควา- อิธตฺถญฺเญว ชานาติ กมฺมา\-ภิสงฺขาร\-วเสน “อยํ ปุคฺคโล กายสฺส เภทา ปรํ มรณา เปตฺติวิสยํ อุปปชฺชิสฺสตี”ติ. ภควา อิธตฺถญฺเญว ชานาติ กมฺมา\-ภิสงฺขาร\-วเสน “อยํ ปุคฺคโล กายสฺส เภทา ปรํ มรณา มนุสฺเสสุ อุปฺปชฺชิสฺสตี”ติ. ภควา อิธตฺถญฺเญว ชานาติ กมฺมา\-ภิสงฺขาร\-วเสน “อยํ ปุคฺคโล สุปฺปฏิปนฺโน กายสฺส เภทา ปรํ มรณา สุคติํ สคฺคํ โลกํ อุปปชฺชิสฺสตี”ติ.}

\prose{วุตฺตํ เหตํ ภควตา\csromandash{}อิธ ปนาหํ สาริปุตฺต เอกจฺจํ ปุคฺคลํ เอวํ เจตสา เจโต ปริจฺจ ปชานามิ “ตถายํ ปุคฺคโล ปฏิปนฺโน, ตถา จ อิริยติ, ตญฺจ มคฺคํ สมารูโฬฺห, ยถา กายสฺส เภทา ปรํ มรณา อปายํ ทุคฺคติํ วินิปาตํ นิรยํ อุปปชฺชิสฺสตี”ติ.}

\prose{อิธ ปนาหํ สาริปุตฺต เอกจฺจํ ปุคฺคลํ เอวํ เจตสา เจโต ปริจฺจ ปชานามิ “ตถายํ ปุคฺคโล ปฏิปนฺโน, ตถา จ อิริยติ, ตญฺจ มคฺคํ สมารูโฬฺห, ยถา กายสฺส เภทา ปรํ มรณา ติรจฺฉาน\-โยนิํ อุปปชฺชิสฺสตี”ติ.}

\proseitem{14}{\csromanfolio{170}อิธ ปนาหํ สาริปุตฺต เอกจฺจํ ปุคฺคลํ เอวํ เจตสา เจโต ปริจฺจ ปชานามิ “ตถายํ ปุคฺคโล ปฏิปนฺโน, ตถา จ อิริยติ, ตญฺจ มคฺคํ สมารูโฬฺห, ยถา กายสฺส เภทา ปรํ มรณา เปตฺติวิสยํ อุปปชฺชิสฺสตี”ติ.}

\prose{อิธ ปนาหํ สาริปุตฺต เอกจฺจํ ปุคฺคลํ เอวํ เจตสา เจโต ปริจฺจ ปชานามิ “ตถายํ ปุคฺคโล ปฏิปนฺโน, ตถา จ อิริยติ, ตญฺจ มคฺคํ สมารูโฬฺห, ยถา กายสฺส เภทา ปรํ มรณา มนุสฺเสสุ อุปฺปชฺชิสฺสตี”ติ.}

\prose{อิธ ปนาหํ สาริปุตฺต เอกจฺจํ ปุคฺคลํ เอวํ เจตสา เจโต ปริจฺจ ปชานามิ “ตถายํ ปุคฺคโล ปฏิปนฺโน, ตถา จ อิริยติ, ตญฺจ มคฺคํ สมารูโฬฺห, ยถา กายสฺส เภทา ปรํ มรณา สุคติํ สคฺคํ โลกํ อุปปชฺชิสฺสตี”ติ.}

\prose{อิธ ปนาหํ สาริปุตฺต เอกจฺจํ ปุคฺคลํ เอวํ เจตสา เจโต ปริจฺจ ปชานามิ “ตถายํ ปุคฺคโล ปฏิปนฺโน, ตถา จ อิริยติ, ตญฺจ มคฺคํ สมารูโฬฺห, ยถา อาสวานํ ขยา อนาสวํ เจโตวิมุตฺติํ ปญฺญาวิมุตฺติํ ทิฏฺเฐว ธมฺเม สยํ อภิญฺญา สจฺฉิกตฺวา อุปสมฺปชฺช วิหรตี”ติ ติฏฺฐนฺตเมนํ ชานาติ.}

\prose{\textbf{ธิมุตฺตํ ตปฺปรายณนฺ}ติ ธิมุตฺตนฺติ อากิญฺจญฺญา\-ยตนํ ธิมุตฺตนฺติ วิโมกฺเขน ธิมุตฺตํ ตตฺราธิมุตฺตํ ตทธิมุตฺตํ ตทา\-ธิปเตยฺยํ. อถ วา ภควา ชานาติ “อยํ ปุคฺคโล รูปาธิมุตฺโต สทฺทาธิมุตฺโต คนฺธาธิมุตฺโต รสาธิมุตฺโต โผฏฺฐพฺพา\-ธิมุตฺโต กุลาธิมุตฺโต คณาธิมุตฺโต อาวาสาธิมุตฺโต ลาภาธิมุตฺโต ยสาธิมุตฺโต ปสํสา\-ธิมุตฺโต สุขาธิมุตฺโต จีวราธิมุตฺโต ปิณฺฑปาตา\-ธิมุตฺโต เสนาสนา\-ธิมุตฺโต คิลานปจฺจย\-เภสชฺชปริกฺขาราธิมุตฺโต สุตฺตนฺตา\-ธิมุตฺโต วินยาธิมุตฺโต อภิธมฺมา\-ธิมุตฺโต อารญฺญกงฺคา\-ธิมุตฺโต ปิณฺฑปาติกงฺคา\-ธิมุตฺโต ปํสุกูลิกงฺคา\-ธิมุตฺโต เตจีวริกงฺคา\-ธิมุตฺโต สปทานจาริกงฺคา\-ธิมุตฺโต ขลุปจฺฉาภตฺติกงฺคา\-ธิมุตฺโต เนสชฺชิกงฺคา\-ธิมุตฺโต ยถาสนฺถติกงฺคา\-ธิมุตฺโต ปฐมชฺฌานา\-ธิมุตฺโต ทุติยชฺฌานา\-ธิมุตฺโต ตติยชฺฌานา\-ธิมุตฺโต จตุตฺถชฺฌานา\-ธิมุตฺโต อากาสานญฺจายตนสมาปตฺตา\-ธิมุตฺโต วิญฺญาณญฺจายตนสมาปตฺตา\-ธิมุตฺโต อากิญฺจญฺญายตนสมาปตฺตา\-ธิมุตฺโต เนวสญฺญานาสญฺญายตนสมาปตฺตาธิมุตฺโต”ติ ธิมุตฺตํ.}

\prose{\csromanfolio{171}\textbf{ตปฺปรายณนฺ}ติ อากิญฺจญฺญายตนมยํ ตปฺปรายณํ กมฺมปรายณํ วิปาก\-ปรายณํ กมฺมครุกํ ปฏิสนฺธิ\-ครุกํ. อถ วา ภควา ชานาติ “อยํ ปุคฺคโล รูปปรายโณ ฯเปฯ เนว\-สญฺญา\-นา\-สญฺญา\-ยตนสมาปตฺติปรายโณ”ติ ธิมุตฺตํ ตปฺปรายณํ. เตนาห ภควา\csromandash{}}

\prose{วิญฺญาณ\-ฏฺฐิติโย สพฺพา, (โปสาลาติ ภควา,)}

\prose{อภิชานํ ตถาคโต.}

\setlength{\csromangathapagewidth}{0pt}
\setlength{\csromangathawakwidth}{0pt}
\csromangathameasuregroup{ติฏฺฐนฺตเมนํ ชานาติ, \\ \textbf{อากิญฺจญฺญา}\textbf{สมฺภวํ} ญตฺวา, \\ \textbf{เอวเมตํ อภิญฺญา}ย, \\ \textbf{เอตํ ญาณํ ตถํ} ตสฺส,}{\csromangathabat{ติฏฺฐนฺตเมนํ ชานาติ,}{ธิมุตฺตํ ตปฺปรายณนฺติ.} \\ \csromangathabat{\textbf{อากิญฺจญฺญา}\textbf{สมฺภวํ} ญตฺวา,}{นนฺทิสํโยชนํ อิติ.} \\ \csromangathabat{\textbf{เอวเมตํ อภิญฺญา}ย,}{ตโต ตตฺถ วิปสฺสติ.} \\ \csromangathabat{\textbf{เอตํ ญาณํ ตถํ} ตสฺส,}{พฺราหฺมณสฺส วุสีมโต.}}
\csromangathasetleft{ติฏฺฐนฺตเมนํ ชานาติ, \\ \textbf{อากิญฺจญฺญา}\textbf{สมฺภวํ} ญตฺวา, \\ \textbf{เอวเมตํ อภิญฺญา}ย, \\ \textbf{เอตํ ญาณํ ตถํ} ตสฺส,}
\csromangathagroup{\csromangathastanza{\csromangathabat{ติฏฺฐนฺตเมนํ ชานาติ,}{ธิมุตฺตํ ตปฺปรายณนฺติ.}}\csromangathastanza{\csromangathaitembat{84}{\textbf{อากิญฺจญฺญา}\-\textbf{สมฺภวํ} ญตฺวา,}{นนฺทิ\-สํโยชนํ อิติ.} \\ \csromangathacontbat{\textbf{เอวเมตํ อภิญฺญา}ย,}{ตโต ตตฺถ วิปสฺสติ.} \\ \csromangathacontbat{\textbf{เอตํ ญาณํ ตถํ} ตสฺส,}{พฺราหฺมณสฺส วุสีมโต.}}}

\prose{\textbf{อากิญฺจญฺญา}\-\textbf{สมฺภวํ ญตฺวา}ติ อากิญฺจญฺญาสมฺภโวติ วุจฺจติ อากิญฺจญฺญายตน\-สํวตฺตนิโก กมฺมา\-ภิสงฺขาโร, อากิญฺจญฺญายตน\-สํวตฺตนิกํ กมฺมา\-ภิสงฺขารํ อากิญฺจญฺญาสมฺภโวติ ญตฺวา, ลคฺคนนฺติ ญตฺวา, พนฺธนนฺติ ญตฺวา, ปลิโพโธติ ญตฺวา ชานิตฺวา ตุลยิตฺวา ตีรยิตฺวา วิภาวยิตฺวา วิภูตํ กตฺวาติ อากิญฺจญฺญา\-สมฺภวํ ญตฺวา.}

\prose{\textbf{นนฺทิ}\-\textbf{สํโยชนํ อิตี}ติ นนฺทิ\-สํโยชนํ วุจฺจติ อรูปราโค, อรูปราเคน ตํ กมฺมํ ลคฺคํ ลคฺคิตํ ปลิพุทฺธํ อรูปราคํ นนฺทิ\-สํโยชนนฺติ ญ ตฺวา, ลคฺคนนฺติ ญ ตฺวา, พนฺธนนฺติ ญ ตฺวา, ปลิโพโธติ ญ ตฺวา ชานิตฺวา ตุลยิตฺวา ตีรยิตฺวา วิภาวยิตฺวา วิภูตํ กตฺวา. \textbf{อิตี}ติ ปทสนฺธิ ปทสํสคฺโค ปทปาริปูรี อกฺขร\-สมวาโย พฺยญฺชน\-สิลิฏฺฐตา ปทา\-นุปุพฺพตา\-เปตํ อิตีติ นนฺทิ\-สํโยชนํ อิติ.}

\prose{\textbf{เอวเมตํ อภิญฺญายา}ติ เอวํ เอตํ อภิญฺญาย ชานิตฺวา ตุลยิตฺวา ตีรยิตฺวา วิภาวยิตฺวา วิภูตํ กตฺวาติ เอวเมตํ อภิญฺญาย.}

\prose{\textbf{ตโต ตตฺถ วิปสฺสตี}ติ \textbf{ตตฺถา}ติ อากิญฺจญฺญา\-ยตนํ สมาปชฺชิตฺวา ตโต วุฏฺฐหิตฺวา ตตฺถ ชาเต จิตฺตเจตสิเก ธมฺเม อนิจฺจโต วิปสฺสติ, ทุกฺขโต วิปสฺสติ, โรคโต ฯเปฯ นิสฺสรณโต วิปสฺสติ ทกฺขติ โอโลเกติ นิชฺฌายติ อุปปริ\-กฺขตีติ ตโต ตตฺถ วิปสฺสติ.}

\prose{\textbf{เอตํ ญาณํ ตถํ ตสฺสา}ติ เอตํ ญาณํ ตจฺฉํ ภูตํ ยาถาวํ อวิปรีตํ ตสฺสาติ เอตํ ญาณํ ตถํ ตสฺส.}

\csromanmatikaanchor{mtk.36}
\csromantocmarkref{subsection}{15. โมฆราชมาณวปุจฺฉานิทฺเทส}{mtk.36}
\bookmark[dest={mtk.36},level=2]{15. โมฆราชมาณวปุจฺฉานิทฺเทส}
\prose{\csromanfolio{172}\textbf{พฺราหฺมณสฺส วุสีมโต}ติ \textbf{พฺราหฺมโณ}ติ สตฺตนฺนํ ธมฺมานํ พาหิตตฺตา พฺราหฺมโณ ฯเปฯ อสิโต ตาทิ ปวุจฺจเต ส พฺรหฺมาติ. \textbf{พฺราหฺมณสฺส วุสีมโต}ติ ปุถุชฺชน\-กลฺยาณํ อุปาทาย สตฺต เสกฺขา อปฺปตฺตสฺส ปตฺติยา อนธิคตสฺส อธิคมาย อสจฺฉิกตสฺส สจฺฉิกิริยาย วสนฺติ สํวสนฺติ อาวสนฺติ ปริวสนฺติ, อรหา วุสิตวา กตกรณีโย โอหิตภาโร อนุปฺปตฺต\-สทตฺโถ ปริกฺขีณ\-ภวสํโยชโน สมฺมทญฺญา วิมุตฺโต, โส วุตฺถวาโส จิณฺณจรโณ ฯเปฯ ชาติ\-มรณ\-สํสาโร, นตฺถิ ตสฺส ปุนพฺภโวติ พฺราหฺมณสฺส วุสีมโต. เตนาห ภควา\csromandash{}}

\setlength{\csromangathapagewidth}{0pt}
\setlength{\csromangathawakwidth}{0pt}
\csromangathameasuregroup{อากิญฺจญฺญาสมฺภวํ ญ ตฺวา, \\ เอวเมตํ อภิญฺญาย, \\ เอตํ ญาณํ ตถํ ตสฺส,}{\csromangathabat{อากิญฺจญฺญาสมฺภวํ ญ ตฺวา,}{นนฺทิสํโยชนํ อิติ.} \\ \csromangathabat{เอวเมตํ อภิญฺญาย,}{ตโต ตตฺถ วิปสฺสติ.} \\ \csromangathabat{เอตํ ญาณํ ตถํ ตสฺส,}{พฺราหฺมณสฺส วุสีมโตติ.}}
\csromangathasetleft{อากิญฺจญฺญาสมฺภวํ ญ ตฺวา, \\ เอวเมตํ อภิญฺญาย, \\ เอตํ ญาณํ ตถํ ตสฺส,}
\csromangathagroup{\csromangathastanza{\csromangathabat{อากิญฺจญฺญา\-สมฺภวํ ญ ตฺวา,}{นนฺทิ\-สํโยชนํ อิติ.} \\ \csromangathabat{เอวเมตํ อภิญฺญาย,}{ตโต ตตฺถ วิปสฺสติ.} \\ \csromangathabat{เอตํ ญาณํ ตถํ ตสฺส,}{พฺราหฺมณสฺส วุสีมโตติ.}}}

\prose{สห คาถา\-ปริโยสานา ฯเปฯ “สตฺถา เม ภนฺเต ภควา, สาวโกหมสฺมี”ติ.}

\vspace{\tipitakaparskip}
\csromancenter{โปสาล\-มาณว\-ปุจฺฉา\-นิทฺเทโส จุทฺทสโม.}
\csromanhangingbegin
\hangingheaditem{85}{\textbf{ทฺวาหํ สกฺกํ อปุจฺฉิสฺสํ, (อิจฺจายสฺมา โมฆราชา,)}}
\hangingbody{\textbf{น เม พฺยากาสิ จกฺขุมา.}}
\hangingbody{\textbf{ยาวตติยญฺจ เทวีสิ1,}}
\hangingbody{\textbf{พฺยากโรตีติ เม สุตํ.}}
\csromanhangingend

\prose{\textbf{ทฺวาหํ สกฺกํ อปุจฺฉิสฺสนฺ}ติ โส พฺราหฺมโณ ทฺวิกฺขตฺตุํ พุทฺธํ ภควนฺตํ ปญฺหํ อปุจฺฉิ, ตสฺส ภควา ปญฺหํ ปุฏฺโฐ น พฺยากาสิ “ตทนฺตรา\csromansharedfootnote{23f094d72be552d6}{จกฺขุสมนนฺตรา (สฺยา)} อิมสฺส พฺราหฺมณสฺส อินฺทฺริย\-ปริปาโก ภวิสฺสตี”ติ. \textbf{สกฺกนฺ}ติ สกฺโก, ภควา สกฺยกุลา ปพฺพชิโตติปิ สกฺโก. อถ วา อฑฺโฒ มหทฺธโน ธนวาติปิ สกฺโก. ตสฺสิมานิ ธนานิ. เสยฺยถิทํ, สทฺธาธนํ สีลธนํ หิริธนํ โอตฺตปฺปธนํ สุตธนํ จาคธนํ ปญฺญาธนํ สติปฏฺฐานธนํ สมฺมปฺปธานธนํ อิทฺธิปาทธนํ อินฺทฺริยธนํ พลธนํ โพชฺฌงฺคธนํ มคฺคธนํ ผลธนํ นิพฺพานธนํ. อิเมหิ อเนกวิเธหิ ธนรตเนหิ อฑฺโฒ มหทฺธโน ธนวาติปิ สกฺโก. อถ วา \csromanfolio{173}สกฺโก ปหุ วิสวี อลมตฺโต สูโร วีโร วิกฺกนฺโต อภีรู อจฺฉมฺภี อนุตฺราสี อปลายี ปหีน\-ภยเภรโว วิคต\-โลมหํโสติปิ สกฺโก. \textbf{ทฺวาหํ สกฺกํ อปุจฺฉิสฺสนฺ}ติ ทฺวาหํ สกฺกํ อปุจฺฉิสฺสํ อยาจิสฺสํ อชฺเฌสิสฺสํ ปสาทยิสฺสนฺติ ทฺวาหํ สกฺกํ อปุจฺฉิสฺสํ.}

\prose{\textbf{อิจฺจายสฺมา โมฆราชา}ติ \textbf{อิจฺจา}ติ ปทสนฺธิ ฯเปฯ. \textbf{อายสฺมา}ติ ปิยวจนํ ฯเปฯ. \textbf{โมฆราชา}ติ ตสฺส พฺราหฺมณสฺส นามํ ฯเปฯ อภิลาโปติ อิจฺจายสฺมา โมฆราชา.}

\prose{\textbf{น เม พฺยากาสิ จกฺขุมา}ติ \textbf{น เม พฺยากาสี}ติ น เม พฺยากาสิ น อาจิกฺขิ น เทเสสิ น ปญฺญเปสิ น ปฏฺฐเปสิ น วิวริ น วิภชิ น อุตฺตานีอกาสิ น ปกาเสสิ. \textbf{จกฺขุมา}ติ ภควา ปญฺจหิ จกฺขูหิ จกฺขุมา, มํสจกฺขุนาปิ จกฺขุมา, ทิพฺพจกฺขุนาปิ\csromansharedfootnote{49299501b5f13b20}{ทิพฺเพน จกฺขุนาปิ (ก)} จกฺขุมา, ปญฺญาจกฺขุนาปิ จกฺขุมา, พุทฺธจกฺขุนาปิ จกฺขุมา, สมนฺตจกฺขุนาปิ จกฺขุมา.}

\prose{กถํ ภควา มํสจกฺขุนาปิ จกฺขุมา. มํสจกฺขุมฺหิ ภควโต ปญฺจ วณฺณา สํวิชฺชนฺติ นีโล จ วณฺโณ ปีตโก จ วณฺโณ โลหิตโก จ วณฺโณ กโณฺห จ วณฺโณ โอทาโต จ วณฺโณ. ยตฺถ จ อกฺขิโลมานิ ปติฏฺฐิตานิ, ตํ นีลํ โหติ สุนีลํ ปาสาทิกํ ทสฺสเนยฺยํ อุมาปุปฺผ\-สมานํ\csromansharedfootnote{1fe2a04ccef7a268}{อุมฺมารปุปฺผสมานํ (สฺยา) ขุ 7. 276 ปิฏฺเฐปิ. 3. อฬาริฏฺฐกสมานํ (สฺยา)}. ตสฺส ปรโต ปีตกํ โหติ สุปีตกํ สุวณฺณวณฺณํ ปาสาทิกํ ทสฺสเนยฺยํ กณิการปุปฺผ\-สมานํ, อุภโต จ อกฺขิกูฏานิ ภควโต โลหิตกานิ โหนฺติ สุโลหิตกานิ ปาสาทิกานิ ทสฺสเนยฺยานิ อินฺทโคปก\-สมานานิ. มชฺเฌ กณฺหํ โหติ สุกณฺหํ อลูขํ สินิทฺธํ ปาสาทิกํ ทสฺสเนยฺยํ อทฺทา\-ริฏฺฐ\-กสมานํ\csromansharedfootnote{380e37481cc0b694}{อตฺถงฺคมิโต (สฺยา, ก) 5. อกาลเมโฆ (สฺยา, ก) ปสฺส ขุ 7. 276 ปิฏฺเฐ.}. ตสฺส ปรโต โอทาตํ โหติ สุโอทาตํ เสตํ ปณฺฑรํ ปาสาทิกํ ทสฺสเนยฺยํ โอสธิตารก\-สมานํ. เตน ภควา ปากติเกน มํสจกฺขุนา อตฺตภาว\-ปริยาปนฺเนน ปุริม\-สุจริต\-กมฺมา\-ภินิพฺพตฺเตน สมนฺตา โยชนํ ปสฺสติ ทิวา เจว รตฺติํ จ. ยทา หิ จตุรงฺคสมนฺนาคโต อนฺธกาโร โหติ, สูริโย จ อตฺถงฺคโต\csromansharedfootnote{b5e198b9dc9569aa}{อฬาริฏฺฐกสมานํ (สฺยา)} โหติ, กาฬปกฺโข จ อุโปสโถ โหติ, ติพฺโพ จ วนสณฺโฑ โหติ, มหา จ กาฬเมโฆ\csromansharedfootnote{7a59e9ea9e1c32f7}{อกาลเมโฆ (สฺยา, ก) ปสฺส ขุ 7. 276 ปิฏฺเฐ.} อพฺภุฏฺฐิโต โหติ. เอวรูเป จตุรงฺคสมนฺนาคเต อนฺธกาเร \csromanfolio{174}สมนฺตา โยชนํ ปสฺสติ, นตฺถิ โส กุฏฺโฏ วา กวาโฏ วา ปากาโร วา ปพฺพโต วา คจฺโฉ วา ลตา วา อาวรณํ รูปานํ ทสฺสนาย. เอกํ เจ ติลผลํ นิมิตฺตํ กตฺวา ติลวาเห ปกฺขิเปยฺย, ตํเยว ติลผลํ อุทฺธเรยฺย. เอวํ ปริสุทฺธํ ภควโต ปากติกํ มํสจกฺขุ. เอวํ ภควา มํสจกฺขุนาปิ จกฺขุมา.}

\proseitem{15}{กถํ ภควา ทิพฺเพน จกฺขุนาปิ จกฺขุมา, ภควา ทิพฺเพน จกฺขุนา วิสุทฺเธน อติกฺกนฺต\-มานุสเกน สตฺเต ปสฺสติ จวมาเน อุปปชฺชมาเน หีเน ปณีเต สุวณฺเณ ทุพฺพณฺเณ สุคเต ทุคฺคเต, ยถากมฺมูปเค สตฺเต ปชานาติ “อิเม วต โภนฺโต สตฺตา กาย\-ทุจฺจริเตน สมนฺนาคตา วจีทุจฺจริเตน สมนฺนาคตา มโน\-ทุจฺจริเตน สมนฺนาคตา อริยานํ อุปวาทกา มิจฺฉาทิฏฺฐิกา มิจฺฉาทิฏฺฐิ\-กมฺม\-สมาทานา, เต กายสฺส เภทา ปรํ มรณา อปายํ ทุคฺคติํ วินิปาตํ นิรยํ อุปปนฺนา, อิเม วา ปน โภนฺโต สตฺตา กายสุจริเตน สมนฺนาคตา วจีสุจริเตน สมนฺนาคตา มโนสุจริเตน สมนฺนาคตา อริยานํ อนุปวาทกา สมฺมาทิฏฺฐิกา สมฺมา\-ทิฏฺฐิ\-กมฺม\-สมาทานา, เต กายสฺส เภทา ปรํ มรณา สุคติํ สคฺคํ โลกํ อุปปนฺนา”ติ อิติ ทิพฺเพน จกฺขุนา วิสุทฺเธน อติกฺกนฺต\-มานุสเกน สตฺเต ปสฺสติ จวมาเน อุปปชฺชมาเน หีเน ปณีเต สุวณฺเณ ทุพฺพณฺเณ สุคเต ทุคฺคเต, ยถากมฺมูปเค สตฺเต ปชานาติ. อากงฺขมาโน จ ภควา เอกมฺปิ โลกธาตุํ ปสฺเสยฺย, เทฺวปิ โลกธาตุโย ปสฺเสยฺย, ติสฺโสปิ โลกธาตุโย ปสฺเสยฺย, จตสฺโสปิ โลกธาตุโย ปสฺเสยฺย, ปญฺจปิ โลกธาตุโย ปสฺเสยฺย, ทสปิ โลกธาตุโย ปสฺเสยฺย, วีสมฺปิ โลกธาตุโย ปสฺเสยฺย, ติํสมฺปิ โลกธาตุโย ปสฺเสยฺย, จตฺตาลีสมฺปิ โลกธาตุโย ปสฺเสยฺย, ปญฺญาสมฺปิ โลกธาตุโย ปสฺเสยฺย, สตมฺปิ โลกธาตุโย ปสฺเสยฺย, สหสฺสิมฺปิ จูฬนิกํ โลกธาตุํ ปสฺเสยฺย, ทฺวิสหสฺสิมฺปิ มชฺฌิมิกํ โลกธาตุํ ปสฺเสยฺย, ติสหสฺสิมฺปิ โลกธาตุํ ปสฺเสยฺย, มหาสหสฺสิมฺปิ\csromansharedfootnote{719658848efec62e}{ติสหสฺสิํ มหาสหสฺสมฺปิ (ก)} โลกธาตุํ ปสฺเสยฺย, ยาวตกํ วา\csromansharedfootnote{eb6bfd2e1f3a2e9d}{ยาวตา (สี, ก)} ปน อากงฺเขยฺย, ตาวตกํ ปสฺเสยฺย. เอวํ ปริสุทฺธํ ภควโต ทิพฺพจกฺขุ. เอวํ ภควา ทิพฺเพน จกฺขุนาปิ จกฺขุมา.}

\prose{\csromanfolio{175}กถํ ภควา ปญฺญาจกฺขุนาปิ จกฺขุมา, ภควา มหาปญฺโญ ปุถุปญฺโญ ชวนปญฺโญ หาสปญฺโญ ติกฺขปญฺโญ นิพฺเพธิก\-ปญฺโญ ปญฺญาปเภท\-กุสโล ปภินฺนญาโณ อธิคต\-ปฏิสมฺภิท\-ปฺปตฺโต จตุเวสารชฺช\-ปฺปตฺโต ทสพลธารี ปุริสาสโภ ปุริสสีโห ปุริสนาโค ปุริสาชญฺโญ ปุริสโธรโยฺห อนนฺตญาโณ อนนฺตเตโช อนนฺตยโส อฑฺโฒ มหทฺธโน ธนวา เนตา วิเนตา อนุเนตา ปญฺญาเปตา นิชฺฌาเปตา เปกฺเขตา ปสาเทตา. โส หิ ภควา อนุปฺปนฺนสฺส มคฺคสฺส อุปฺปาเทตา, อสญฺชาตสฺส มคฺคสฺส สญฺชเนตา, อนกฺขาตสฺส มคฺคสฺส อกฺขาตา, มคฺคญฺญู มคฺควิทู มคฺคโกวิโท มคฺคานุคา จ ปน เอตรหิ สาวกา วิหรนฺติ ปจฺฉา สมนฺนาคตา.}

\prose{โส หิ ภควา ชานํ ชานาติ, ปสฺสํ ปสฺสติ, จกฺกุภูโต ญาณภูโต ธมฺมภูโต พฺรหฺมภูโต วตฺตา ปวตฺตา อตฺถสฺส นินฺเนตา อมตสฺส ทาตา ธมฺมสฺสามี\csromansharedfootnote{3fa2414d5b595851}{ธมฺมสามิ (สฺยา, ก)} ตถาคโต. นตฺถิ ตสฺส ภควโต อญฺญาตํ อทิฏฺฐํ อวิทิตํ อสจฺฉิกตํ อผสฺสิตํ\csromansharedfootnote{31d008d1502c9146}{อผุสิตํ (สฺยา, ก)} ปญฺญาย, อตีตํ อนาคตํ ปจฺจุปฺปนฺนํ\csromansharedfootnote{58c03379f58db169}{อตีตานาคตปจฺจุปฺปนฺนํ (สฺยา)} อุปาทาย สพฺเพ ธมฺมา สพฺพากาเรน พุทฺธสฺส ภควโต ญาณมุเข อาปาถํ อาคจฺฉนฺติ. ยํ กิญฺจิ เนยฺยํ นาม อตฺถิ ชานิตพฺพํ\csromansharedfootnote{371e3c006b69c05e}{อตฺถิ ธมฺมํ ชานิตพฺพํ (ก)} อตฺตตฺโถ วา ปรตฺโถ วา อุภยตฺโถ วา ทิฏฺฐิธมฺมฺมิโก วา อตฺโถ สมฺปรายิโก วา อตฺโถ อุตฺตาโน วา อตฺโถ คมฺภีโร วา อตฺโถ คุโฬฺห วา อตฺโถ ปฏิจฺฉนฺโน วา อตฺโถ เนยฺโย วา อตฺโถ นีโต วา อตฺโถ อนวชฺโช วา อตฺโถ นิกฺกิเลโส วา อตฺโถ โวทาโน วา อตฺโถ ปรมตฺโถ วา อตฺโถ\csromansharedfootnote{78f3b100b7a99dee}{ปรมตฺโถ วา อตฺโถ (ก)}, สพฺพํ ตํ อนฺโต พุทฺธญาเณ ปริวตฺตติ, สพฺพํ กายกมฺมํ พุทฺธสฺส ภควโต ญาณา\-นุปริวตฺติ, สพฺพํ วจีกมฺมํ ญาณา\-นุปริวตฺติ, สพฺพํ มโนกมฺมํ ญาณา\-นุปริวตฺติ. อตีเต พุทฺธสฺส ภควโต อปฺปฏิหตํ ญาณํ. อนาคเต อปฺปฏิหตํ ญาณํ, ปจฺจุปฺปนฺเน อปฺปฏิหตํ ญาณํ. ยาวตกํ เนยฺยํ ตาวตกํ ญาณํ. ยาวตกํ ญาณํ ตาวตกํ เนยฺยํ. เนยฺย\-ปริยนฺติกํ ญาณํ, ญาณ\-ปริยนฺติกํ เนยฺยํ, เนยฺยํ อติกฺกมิตฺวา ญาณํ นปฺปวตฺตติ, ญาณํ อติกฺกมิตฺวา เนยฺยปโถ นตฺถิ. อญฺญมญฺญ\-ปริยนฺต\-ฏฺฐายิโน เต ธมฺมา, ยถา ทฺวินฺนํ สมุคฺค\-ปฏลานํ สมฺมา\-ผุสิตานํ เหฏฺฐิมํ สมุคฺคปฏลํ}

\prose{\csromanfolio{176}อุปริมํ นาติวตฺตติ. อุปริมํ สมุคฺคปฏลํ เหฏฺฐิมํ นาติวตฺตติ. อญฺญมญฺญ\-ปริยนฺต\-ฏฺฐายิโน. เอวเมว พุทฺธสฺส ภควโต เนยฺยญฺจ ญาณญฺจ อญฺญมญฺญ\-ปริยนฺต\-ฏฺฐายิโน. ยาวตกํ เนยฺยํ ตาวตกํ ญาณํ. ยาวตกํ ญาณํ ตาวตกํ เนยฺยํ. เนยฺย\-ปริยนฺติกํ ญาณํ, ญาณ\-ปริยนฺติกํ เนยฺยํ. เนยฺยํ อติกฺกมิตฺวา ญาณํ นปฺปวตฺตติ, ญาณํ อติกฺกมิตฺวา เนยฺยปโถ นตฺถิ. อญฺญมญฺญ\-ปริยนฺต\-ฏฺฐายิโน เต ธมฺมา. สพฺพธมฺเมสุ พุทฺธสฺส ภควโต ญาณํ ปวตฺตติ. สพฺเพ ธมฺมา พุทฺธสฺส ภควโต อาวชฺชนปฏิพทฺธา อากงฺข\-ปฏิพทฺธา มนสิการ\-ปฏิพทฺธา จิตฺตุปฺปาท\-ปฏิพทฺธา. สพฺพสตฺเตสุ พุทฺธสฺส ภควโต ญาณํ ปวตฺตติ. สพฺเพสํ จ สตฺตานํ ภควา อาสยํ ชานาติ, อนุสยํ ชานาติ, จริตํ ชานาติ, อธิมุตฺติํ ชานาติ, อปฺปรชกฺเข มหารชกฺเข ติกฺขินฺทฺริเย มุทินฺทฺริเย สฺวากาเร ทฺวากาเร สุวิญฺญาปเย ทุวิญฺญาปเย ภพฺพาภพฺเพ สตฺเต ชานาติ. สเทวโก โลโก สมารโก สพฺรหฺมโก สสฺสมณ\-พฺราหฺมณี ปชา สเทวมนุสฺสา อนฺโตพุทฺธญาเณ ปริวตฺตติ.}

\prose{ยถา เย เกจิ มจฺฉกจฺฉปา อนฺตมโส ติมิติมิงฺคลํ\csromansharedfootnote{990b30dac5af811e}{ติมิติปิงฺคลํ (ก)} อุปาทาย อนฺโต\-มหาสมุทฺเท ปริวตฺตนฺติ. เอวเมว สเทวโก โลโก สมารโก โลโก สพฺรหฺมโก โลโก สสฺสมณ\-พฺราหฺมณี ปชา สเทวมนุสฺสา อนฺโตพุทฺธญาเณ ปริวตฺตติ.  ยถา เย เกจิ ปกฺขี อนฺตมโส ครุฬํ เวนเตยฺยํ อุปาทาย อากาสสฺส ปเทเส ปริวตฺตนฺติ. เอวเมว เยปิ เต สาริปุตฺตสมา ปญฺญาย สมนฺนาคตา, เตปิ พุทฺธญาณสฺส ปเทเส ปริวตฺตนฺติ. พุทฺธญาณํ เทวมนุสฺสานํ ปญฺญํ ผริตฺวา อภิภวิตฺวา ติฏฺฐติ.}

\prose{เยปิ เต ขตฺติย\-ปณฺฑิตา พฺราหฺมณ\-ปณฺฑิตา คหปติ\-ปณฺฑิตา สมณปณฺฑิตา นิปุณา กตปรปฺปวาทา วาลเวธิรูปา โวภินฺทนฺตา\csromansharedfootnote{8c0d862c67c1e117}{เต ภินฺทนฺตา (สฺยา, ก)} มญฺเญ จรนฺติ ปญฺญาคเตน ทิฏฺฐิคตานิ, เต ปเญฺห อภิส\-งฺขริ\-ตฺวา อภิส\-งฺขริ\-ตฺวา ตถาคตํ อุปสงฺกมิตฺวา ปุจฺฉนฺติ คูฬฺหานิ จ ปฏิจฺฉนฺนานิ. กถิตา วิสชฺชิตา จ เต ปญฺหา ภควตา\csromansharedfootnote{6ff0aa0ea04171ac}{ภควโต (ก)} โหนฺติ นิทฺทิฏฺฐ\-การณา. อุปกฺขิตฺตกา จ เต ภควโต สมฺปชฺชนฺติ. อถ โข ภควาว ตตฺถ อติโรจติ. ยทิทํ ปญฺญายาติ, เอวํ ภควา ปญฺญาจกฺขุนาปิ จกฺขุมา.}

\prosewithcloser{\csromanfolio{177}กถํ ภควา พุทฺธจกฺขุนาปิ จกฺขุมา, ภควา พุทฺธจกฺขุนา โลกํ โวโลเกนฺโต\csromansharedfootnote{5c03785ab2e84bc1}{โอโลเกนฺโต (ก)} อทฺทส สตฺเต อปฺปรชกฺเข มหารชกฺเข ติกฺขินฺทฺริเย มุทินฺทฺริเย สฺวากาเร ทฺวากาเร สุวิญฺญาปเย ทุวิญฺญาปเย อปฺเปกจฺเจ ปรโลก\-วชฺช\-ภยทสฺสาวิโน\csromansharedfootnote{4af99de183eccd47}{ปร...ทสฺสาวิเน (ก)} วิหรนฺเต. เสยฺยถาปิ นาม อุปฺปลินิยํ วา ปทุมินิยํ วา ปุณฺฑรีกินิยํ วา อปฺเปกจฺจานิ อุปฺปลานิ วา ปทุมานิ วา ปุณฺฑรีกานิ วา อุทเก ชาตานิ อุทเก สํวฑฺฒานิ อุทกานุคฺคตานิ อนฺโต\-นิมุคฺค\-โปสีนิ\csromansharedfootnote{6f9a46012b28ca50}{อนฺโตนิมฺมุคฺคโปสีนิ (ก)}, อปฺเปกจฺจานิ อุปฺปลานิ วา ปทุมานิ วา ปุณฺฑรีกานิ วา อุทเก ชาตานิ อุทเก สํวฑฺฒานิ สโมทกํ ฐิตานิ, อปฺเปกจฺจานิ อุปฺปลานิ วา ปทุมานิ วา ปุณฺฑรีกานิ วา อุทเก ชาตานิ อุทเก สํวฑฺฒานิ อุทกา อจฺจุคฺคมฺม ติฏฺฐนฺติ อนุปลิตฺตานิ อุทเกน. เอวเมวํ ภควา พุทฺธจกฺขุนา โลกํ โวโลเกนฺโต อทฺทส สตฺเต อปฺปรชกฺเข มหารชกฺเข ติกฺขินฺทฺริเย มุทินฺทฺริเย สฺวากาเร ทฺวากาเร สุวิญฺญาปเย ทุวิญฺญาปเย อปฺเปกจฺเจ ปรโลก\-วชฺช\-ภยทสฺสาวิโน วิหรนฺเต. ชานาติ ภควา “อยํ ปุคฺคโล ราคจริโต, อยํ โทสจริโต, อยํ โมหจริโต, อยํ วิตกฺกจริโต, อยํ สทฺธาจริโต, อยํ ญาณจริโต”ติ. ราคจริตสฺส ภควา ปุคฺคลสฺส อสุภกถํ กเถติ. โทสจริตสฺส ภควา ปุคฺคลสฺส เมตฺตาภาวนํ อาจิกฺขติ. โมหจริตสฺส ภควา ปุคฺคลสฺส อุทฺเทเส ปริปุจฺฉาย กาเลน ธมฺมสฺสวเน กาเลน ธมฺม\-สากจฺฉาย ครุสํวาเส นิเวเสติ. วิตกฺก\-จริตสฺส ภควา ปุคฺคลสฺส อานาปานสฺสติํ อาจิกฺขติ. สทฺธา\-จริตสฺส ภควา ปุคฺคลสฺส ปสาทนียํ นิมิตฺตํ อาจิกฺขติ พุทฺธสุโพธิํ\csromansharedfootnote{cd6ae893f340cf29}{พุทฺธสุพุทฺธตํ (ก)} ธมฺม\-สุธมฺมตํ สํฆสุปฺปฏิปตฺติํ สีลานิ จ อตฺตโน. ญาณจริตสฺส ภควา ปุคฺคลสฺส วิปสฺสนา\-นิมิตฺตํ อาจิกฺขติ อนิจฺจาการํ ทุกฺขาการํ อนตฺตาการํ.}{“เสเล ยถา ปพฺพตมุทฺธนิ\-ฏฺฐิโต,}
\prose{ยถาปิ ปสฺเส ชนตํ สมนฺตโต.}

\setlength{\csromangathapagewidth}{0pt}
\setlength{\csromangathawakwidth}{0pt}
\csromangathameasuregroup{}{ตถูปมํ ธมฺมมยํ สุเมธ, \\ ปาสาทมารุยฺห สมนฺตจกฺขุ. \\ โสกาวติณฺณํ\textsuperscript{1} ชนตมเปตโสโก, \\ อเวกฺขสฺสุ ชาติชราภิภูตนฺ”ติ.}
\csromangathameasurewak{ตถูปมํ ธมฺมมยํ สุเมธ, \\ ปาสาทมารุยฺห สมนฺตจกฺขุ. \\ โสกาวติณฺณํ\textsuperscript{1} ชนตมเปตโสโก, \\ อเวกฺขสฺสุ ชาติชราภิภูตนฺ”ติ.}
\setlength{\csromangathaleftcol}{0pt}
\csromangathagroup{\csromangathastanza{\csromangathawak{ตถูปมํ ธมฺมมยํ สุเมธ, \\ ปาสาทมารุยฺห สมนฺตจกฺขุ. \\ โสกาวติณฺณํ\csromansharedfootnote{1dc936e39dd63f21}{โสกาวกิณฺณํ (สฺยา)} ชนตม\-เปตโสโก, \\ อเวกฺขสฺสุ ชาติชราภิภูตนฺ”ติ.}}}

\prose{\csromanfolio{178}เอวํ ภควา พุทฺธจกฺขุนาปิ จกฺขุมา.}

\prose{กถํ ภควา สมนฺตจกฺขุนาปิ จกฺขุมา, \textbf{สมนฺตจกฺขุ} วุจฺจติ สพฺพญฺญุต\-ญาณํ, ภควา สพฺพญฺญุต\-ญาเณน อุเปโต สมุเปโต อุปาคโต สมุปาคโต อุปปนฺโน สมุปปนฺโน สมนฺนาคโต.}

\setlength{\csromangathapagewidth}{0pt}
\setlength{\csromangathawakwidth}{0pt}
\csromangathameasuregroup{}{“น ตสฺส อทฺทิฏฺฐมิธตฺถิ กิญฺจิ, \\ อโถ อวิญฺญาตมชานิตพฺพํ. \\ สพฺพํ อภิญฺญาสิ ยทตฺถิ เนยฺยํ, \\ ตถาคโต เตน สมนฺตจกฺขู”ติ. \\ เอวํ ภควา สมนฺตจกฺขุนาปิ จกฺขุมาติ น เม พฺยากาสิ จกฺขุมา.}
\csromangathameasurewak{“น ตสฺส อทฺทิฏฺฐมิธตฺถิ กิญฺจิ, \\ อโถ อวิญฺญาตมชานิตพฺพํ. \\ สพฺพํ อภิญฺญาสิ ยทตฺถิ เนยฺยํ, \\ ตถาคโต เตน สมนฺตจกฺขู”ติ. \\ เอวํ ภควา สมนฺตจกฺขุนาปิ จกฺขุมาติ น เม พฺยากาสิ จกฺขุมา.}
\setlength{\csromangathaleftcol}{0pt}
\csromangathagroup{\csromangathastanza{\csromangathawak{“น ตสฺส อทฺทิฏฺฐมิธตฺถิ กิญฺจิ, \\ อโถ อวิญฺญาตมชานิตพฺพํ. \\ สพฺพํ อภิญฺญาสิ ยทตฺถิ เนยฺยํ, \\ ตถาคโต เตน สมนฺตจกฺขู”ติ.}}\csromangathastanza{\csromangathawak{เอวํ ภควา สมนฺตจกฺขุนาปิ จกฺขุมาติ น เม พฺยากาสิ จกฺขุมา.}}}

\prose{\textbf{ยาวตติยญฺจ เทวีสิ, พฺยากโรตีติ เม สุตนฺ}ติ ยาวตติยํ พุทฺโธ สหธมฺมิกํ ปญฺหํ ปุฏฺโฐ พฺยากโรติ โน สํสาเรตีติ\csromansharedfootnote{b6a5824b5f469f82}{สมฺปายตีติ (สฺยา)} เอวํ มยา อุคฺคหิตํ, เอวํ มยา อุปธาริตํ, เอวํ มยา อุปลกฺขิตํ. \textbf{เทวีสี}ติ ภควา เจว อิสิ จาติ เทวีสิ, ยถา ราชา ปพฺพชิตา วุจฺจนฺติ ราชิสโย, พฺราหฺมณา ปพฺพชิตา วุจฺจนฺติ พฺราหฺมณิสโย, เอวเมว ภควา เจว อิสิ จาติ เทวีสิ.}

\prose{อถ วา ภควา ปพฺพชิโตติปิ \textbf{อิสิ,} มหนฺตํ สีลกฺขนฺธํ เอสี คเวสี ปริเยสีติปิ อิสิ, มหนฺตํ สมาธิ\-กฺขนฺธํ ฯเปฯ มหนฺตํ ปญฺญากฺขนฺธํ. มหนฺตํ วิมุตฺติ\-กฺขนฺธํ. มหนฺตํ วิมุตฺติ\-ญาณ\-ทสฺสนกฺขนฺธํ เอสี คเวสี ปริเยสีติปิ อิสิ, มหโต ตโมกายสฺส ปทาลนํ เอสี คเวสี ปริเยสีติปิ อิสิ, มหโต วิปลฺลาสสฺส ปเภทนํ เอสี คเวสี ปริเยสีติปิ อิสิ, มหโต ตณฺหาสลฺลสฺส อพฺพหนํ. มหโต ทิฏฺฐิ\-สงฺฆาฏสฺส วินิเวฐนํ. มหโต มานทฺธชสฺส ปปาตนํ. มหโต อภิสงฺขารสฺส วูปสมํ. มหโต โอฆสฺส นิตฺถรณํ. มหโต ภารสฺส นิกฺเขปนํ. มหโต สํสาร\-วฏฺฏสฺส อุปจฺเฉทํ. มหโต สนฺตาปสฺส นิพฺพาปนํ. มหโต ปริฬาหสฺส ปฏิปฺปสฺสทฺธิํ. มหโต ธมฺม\-ทฺธชสฺส อุสฺสาปนํ เอสี คเวสี ปริเยสีติปิ อิสิ, มหนฺเต สติปฏฺฐาเน. มหนฺเต สมฺมปฺปธาเน. มหนฺตานิ อินฺทฺริยานิ. มหนฺตานิ พลานิ. มหนฺเต โพชฺฌงฺเค. มหนฺตํ อริยํ อฏฺฐงฺคิกํ มคฺคํ. มหนฺตํ ปรมตฺถํ อมตํ นิพฺพานํ เอสี คเวสี ปริเยสีติปิ อิสิ, มเหสกฺเขหิ วา สตฺเตหิ เอสิโต คเวสิโต ปริเยสิโต}

\prose{\csromanfolio{179}“กหํ พุทฺโธ, กหํ ภควา, กหํ เทวเทโว, กหํ นราสโภ”ติปิ อิสีติ ยาวตติยญฺจ เทวีสิ, พฺยากโรตีติ เม สุตํ. เตนาห โส พฺราหฺมโณ\csromandash{}}

\prose{ทฺวาหํ สกฺกํ อปุจฺฉิสฺสํ, (อิจฺจายสฺมา โมฆราชา,)}

\prose{น เม พฺยากาสิ จกฺขุมา.}

\setlength{\csromangathapagewidth}{0pt}
\setlength{\csromangathawakwidth}{0pt}
\csromangathameasuregroup{}{ยาวตติยญฺจ เทวีสิ, \\ พฺยากโรตีติ เม สุตนฺติ. \\ \textbf{อยํ โลโก ปโร โลโก,} \\ \textbf{พฺรหฺมโลโก สเทวโก.} \\ \textbf{ทิฏฺฐิํ เต นาภิชานาติ,} \\ \textbf{โคตมสฺส ยสสฺสิโน.}}
\csromangathameasurewak{ยาวตติยญฺจ เทวีสิ, \\ พฺยากโรตีติ เม สุตนฺติ. \\ \textbf{อยํ โลโก ปโร โลโก,} \\ \textbf{พฺรหฺมโลโก สเทวโก.} \\ \textbf{ทิฏฺฐิํ เต นาภิชานาติ,} \\ \textbf{โคตมสฺส ยสสฺสิโน.}}
\setlength{\csromangathaleftcol}{0pt}
\csromangathagroup{\csromangathastanza{\csromangathawak{ยาวตติยญฺจ เทวีสิ, \\ พฺยากโรตีติ เม สุตนฺติ.}}\csromangathastanza{\csromangathawak{\csromangathaitemline{86}{\textbf{อยํ โลโก ปโร โลโก,}} \\ \csromangathacontline{\textbf{พฺรหฺมโลโก สเทวโก.}} \\ \csromangathacontline{\textbf{ทิฏฺฐิํ เต นาภิชานาติ,}} \\ \csromangathacontline{\textbf{โคตมสฺส ยสสฺสิโน.}}}}}

\prose{\textbf{อยํ โลโก ปโร โลโก}ติ \textbf{อยํ โลโก}ติ มนุสฺสโลโก. \textbf{ปโร โลโก}ติ มนุสฺสโลกํ ฐเปตฺวา สพฺโพ ปโร โลโกติ อยํ โลโก ปโร โลโก.}

\prose{\textbf{พฺรหฺมโลโก สเทวโก}ติ สเทวโก โลโก สมารโก สพฺรหฺมโก สสฺสมณ\-พฺราหฺมณี ปชา สเทว\-มนุสฺสาติ พฺรหฺมโลโก สเทวโก.}

\prose{\textbf{ทิฏฺฐิํ เต นาภิชานาตี}ติ ตุยฺหํ ทิฏฺฐิํ ขนฺติํ รุจิํ ลทฺธิํ อชฺฌาสยํ อธิปฺปายํ โลโก น ชานาติ, “อยํ เอวํทิฏฺฐิโก เอวํขนฺติโก เอวํรุจิโก เอวํลทฺธิโก เอวํอชฺฌาสโย เอวํอธิปฺปาโย”ติ น ชานาติ น ปสฺสติ น ทกฺขติ นาธิคจฺฉติ น วินฺทติ น ปฏิลภตีติ ทิฏฺฐิํ เต นาภิชานาติ.}

\prose{\textbf{โคตมสฺส ยสสฺสิโน}ติ ภควา ยสปฺปตฺโตติ ยสสฺสี. อถ วา ภควา สกฺกโต ครุกโต มานิโต ปูชิโต อปจิโต ลาภี จีวร\-ปิณฺฑปาต\-เสนาสนคิลาน- ปจฺจยเภสชฺชปริกฺขารานนฺติปิ ยสสฺสีติ โคตมสฺส ยสสฺสิโน. เตนาห โส พฺราหฺมโณ\csromandash{}}

\prose{\textbf{อยํ โลโก ปโร โลโก,}}

\prose{\textbf{พฺรหฺมโลโก สเทวโก.}}

\setlength{\csromangathapagewidth}{0pt}
\setlength{\csromangathawakwidth}{0pt}
\csromangathameasuregroup{\textbf{เอวํ อภิกฺกนฺต ท}สฺสาวิํ, \\ \textbf{กถํ โลกํ อเวกฺกฺ}หนฺตํ,}{\textbf{ทิฏฺฐิํ เต นาภิชานาติ,} \\ โคตมสฺส ยสสฺสิโนติ. \\ \csromangathabat{\textbf{เอวํ อภิกฺกนฺต ท}สฺสาวิํ,}{อตฺถิ ปเญฺหน อาคมํ.} \\ \csromangathabat{\textbf{กถํ โลกํ อเวกฺกฺ}หนฺตํ,}{มจฺจุราชา น ปสฺสติ.}}
\csromangathameasurewak{\textbf{ทิฏฺฐิํ เต นาภิชานาติ,} \\ โคตมสฺส ยสสฺสิโนติ.}
\csromangathasetleft{\textbf{เอวํ อภิกฺกนฺต ท}สฺสาวิํ, \\ \textbf{กถํ โลกํ อเวกฺกฺ}หนฺตํ,}
\csromangathagroup{\csromangathastanza{\csromangathawak{\textbf{ทิฏฺฐิํ เต นาภิชานาติ,} \\ โคตมสฺส ยสสฺสิโนติ.}}\csromangathastanza{\csromanfolio{180}\csromangathaitembat{87}{\textbf{เอวํ อภิกฺกนฺต ท}สฺสาวิํ,}{อตฺถิ ปเญฺหน อาคมํ.} \\ \csromangathacontbat{\textbf{กถํ โลกํ อเวกฺกฺ}หนฺตํ,}{มจฺจุราชา น ปสฺสติ.}}}

\prose{\textbf{เอวํ อภิกฺกนฺต}\-\textbf{ทสฺสาวินฺ}ติ เอวํ อภิกฺกนฺต\-ทสฺสาวิํ อคฺคทสฺสาวิํ เสฏฺฐทสฺสาวิํ วิเส\-ฏฺฐ\-ทสฺสาวิํ ปาโมกฺข\-ทสฺสาวิํ อุตฺตมทสฺสาวิํ ปรมทสฺสาวินฺติ เอวํ อภิกฺกนฺต\-ทสฺสาวิํ.}

\prose{\textbf{อตฺถิ ปเญฺหน อาคมนฺ}ติ ปเญฺหน อตฺถิโก อาคโตมฺหิ ฯเปฯ วหสฺเสตํ ภารนฺติ เอวมฺปิ อตฺถิ ปเญฺหน อาคมํ.}

\prose{\textbf{กถํ โลกํ อเวกฺขนฺตนฺ}ติ กถํ โลกํ อเวกฺขนฺตํ ปจฺจเวกฺขนฺตํ ตุลยนฺตํ ตีรยนฺตํ วิภาวยนฺตํ วิภูตํ กโรนฺตนฺติ กถํ โลกํ อเวกฺขนฺตํ.}

\prose{\textbf{มจฺจุราชา น ปสฺสตี}ติ มจฺจุราชา น ปสฺสติ น ทกฺขติ นาธิคจฺฉติ น วินฺทติ น ปฏิลภตีติ มจฺจุราชา น ปสฺสติ. เตนาห โส พฺราหฺมโณ\csromandash{}}

\setlength{\csromangathapagewidth}{0pt}
\setlength{\csromangathawakwidth}{0pt}
\csromangathameasuregroup{เอวํ อภิกฺกนฺตทสฺสาวิํ, \\ \textbf{กถํ โลกํ อเวกฺกฺ}หนฺตํ, \\ \textbf{สุญฺญโต โลกํ อเวกฺ}ขสฺสุ, \\ \textbf{อตฺตานุทิฏฺฐิํ อูหฺ}อจฺจ, \\ \textbf{เอวํ โลกํ อเวกฺข}นฺตํ,}{\csromangathabat{เอวํ อภิกฺกนฺตทสฺสาวิํ,}{อตฺถิ ปเญฺหน อาคมํ.} \\ \csromangathabat{\textbf{กถํ โลกํ อเวกฺกฺ}หนฺตํ,}{มจฺจุราชา น ปสฺสตีติ.} \\ \csromangathabat{\textbf{สุญฺญโต โลกํ อเวกฺ}ขสฺสุ,}{โมฆราช สทา สโต.} \\ \csromangathabat{\textbf{อตฺตานุทิฏฺฐิํ อูหฺ}อจฺจ,}{เอวํ มจฺจุตโร สิยา.} \\ \csromangathabat{\textbf{เอวํ โลกํ อเวกฺข}นฺตํ,}{มจฺจุราชา น ปสฺสติ.}}
\csromangathasetleft{เอวํ อภิกฺกนฺตทสฺสาวิํ, \\ \textbf{กถํ โลกํ อเวกฺกฺ}หนฺตํ, \\ \textbf{สุญฺญโต โลกํ อเวกฺ}ขสฺสุ, \\ \textbf{อตฺตานุทิฏฺฐิํ อูหฺ}อจฺจ, \\ \textbf{เอวํ โลกํ อเวกฺข}นฺตํ,}
\csromangathagroup{\csromangathastanza{\csromangathabat{เอวํ อภิกฺกนฺต\-ทสฺสาวิํ,}{อตฺถิ ปเญฺหน อาคมํ.} \\ \csromangathabat{\textbf{กถํ โลกํ อเวกฺกฺ}หนฺตํ,}{มจฺจุราชา น ปสฺสตีติ.}}\csromangathastanza{\csromangathaitembat{88}{\textbf{สุญฺญโต โลกํ อเวกฺ}ขสฺสุ,}{โมฆราช สทา สโต.} \\ \csromangathacontbat{\textbf{อตฺตานุทิฏฺฐิํ อูหฺ}อจฺจ,}{เอวํ มจฺจุตโร สิยา.} \\ \csromangathacontbat{\textbf{เอวํ โลกํ อเวกฺข}นฺตํ,}{มจฺจุราชา น ปสฺสติ.}}}

\prose{\textbf{สุญฺญโต โลกํ อเวกฺขสฺสู}ติ โลโกติ นิรย\textbf{โลโก} ติรจฺฉานโลโก เปตฺติวิสย\-โลโก มนุสฺสโลโก เทวโลโก ขนฺธโลโก ธาตุโลโก อายตนโลโก อยํ โลโก ปโร โลโก พฺรหฺมโลโก สเทวโก\csromansharedfootnote{5ac90af8bd744c7a}{สเทวโก โลโก (ก)}. อญฺญตโร ภิกฺขุ ภควนฺตํ เอตทโวจ “โลโก โลโก”ติ ภนฺเต วุจฺจติ, กิตฺตาวตา นุ โข ภนฺเต “โลโก”ติ วุจฺจตีติ. ลุชฺชตีติ โข ภิกฺขุ ตสฺมา “โลโก”ติ วุจฺจติ. กิญฺจ ลุชฺชติ. จกฺขุ โข ภิกฺขุ ลุชฺชติ, รูปา ลุชฺชนฺติ, จกฺขุวิญฺญาณํ ลุชฺชติ, จกฺขุ\-สมฺผสฺโส ลุชฺชติ, ยมฺปิทํ\csromansharedfootnote{85d85685b6b558b3}{ยมิทํ (ก) ปสฺส สํ 2. 278 ปิฏฺเฐ.} จกฺขุ\-สมฺผสฺส\-ปจฺจยา อุปฺปชฺชติ เวทยิตํ สุขํ วา ทุกฺขํ วา อทุกฺขมสุขํ วา, ตมฺปิ ลุชฺชติ. โสตํ ลุชฺชติ, คนฺธา ลุชฺชนฺติ ฯเปฯ กาโย \csromanfolio{181}ลุชฺชติ, โผฏฺฐพฺพา ลุชฺชนฺติ. มโน ลุชฺชติ, ธมฺมา ลุชฺชนฺติ, มโนวิญฺญาณํ ลุชฺชติ, มโนสมฺผสฺโส ลุชฺชติ, ยมฺปิทํ มโน\-สมฺผสฺส\-ปจฺจยา อุปฺปชฺชติ เวทยิตํ สุขํ วา ทุกฺขํ วา อทุกฺขมสุขํ วา, ตมฺปิ ลุชฺชติ. ลุชฺชตีติ โข ภิกฺขุ ตสฺมา “โลโก”ติ วุจฺจติ.}

\prose{\textbf{สุญฺญโต โลกํ อเวกฺขสฺสู}ติ ทฺวีหิ การเณหิ สุญฺญโต โลกํ อเวกฺขติ อวสิย\-ปวตฺต\-สลฺลกฺขณ\-วเสน\csromansharedfootnote{72df64a92a973f5b}{อวสฺสิยปวตฺต... (สฺยา)} วา ตุจฺฉ\-สงฺขาร\-สมนุปสฺสนา\-วเสน วา. กถํ อวสิย\-ปวตฺต\-สลฺลกฺขณ\-วเสน สุญฺญโต โลกํ อเวกฺขติ\csromandash{} รูเป วโส น ลพฺภติ. เวทนาย วโส น ลพฺภติ. สญฺญาย วโส น ลพฺภติ. สงฺขาเรสุ วโส น ลพฺภติ. วิญฺญาเณ วโส น ลพฺภติ. วุตฺตเญฺหตํ ภควตา\csromansharedfootnote{d016d19fdf0a2431}{ปสฺส สํ 2. 55 ปิฏฺเฐสุ.}\csromandash{}รูปํ ภิกฺขเว อนตฺตา, รูปญฺจ หิทํ ภิกฺขเว อตฺตา อภวิสฺส, นยิทํ รูปํ อาพาธาย สํวตฺเตยฺย, ลพฺเภถ จ รูเป “เอวํ เม รูปํ โหตุ, เอวํ เม รูปํ มา อโหสี”ติ. ยสฺมา จ โข ภิกฺขเว รูปํ อนตฺตา, ตสฺมา รูปํ อาพาธาย สํวตฺตติ. น จ ลพฺภติ รูเป “เอวํ เม รูปํ โหตุ, เอวํ เม รูปํ มา อโหสี”ติ.}

\prose{เวทนา อนตฺตา, เวทนา จ หิทํ ภิกฺขเว อตฺตา อภวิสฺส, นยิทํ เวทนา อาพาธาย สํวตฺเตยฺย, ลพฺเภถ จ เวทนาย “เอวํ เม เวทนา โหตุ, เอวํ เม เวทนา มา อโหสี”ติ. ยสฺมา จ โข ภิกฺขเว เวทนา อนตฺตา, ตสฺมา เวทนา อาพาธาย สํวตฺตติ, น จ ลพฺพติ เวทนาย “เอวํ เม เวทนา โหตุ, เอวํ เม เวทนา มา อโหสี”ติ.}

\prose{สญฺญา อนตฺตา, สญฺญา จ หิทํ ภิกฺขเว อตฺตา อภวิสฺส, นยิทํ สญฺญา อาพาธาย สํวตฺเตยฺย, ลพฺเภถ จ สญฺญาย “เอวํ เม สญฺญา โหตุ, เอวํ เม สญฺญา มา อโหสี”ติ. ยสฺมา จ โข ภิกฺขเว สญฺญา อนตฺตา, ตสฺมา สญฺญา อาพาธาย สํวตฺตติ, น จ ลพฺภติ สญฺญาย “เอวํ เม สญฺญา โหตุ, เอวํ เม สญฺญา มา อโหสี”ติ.}

\prose{สงฺขารา อนตฺตา, สงฺขารา จ หิทํ ภิกฺขเว อตฺตา อภวิสฺสํสุ, นยิทํ สงฺขารา อาพาธาย สํวตฺเตยฺยุํ, ลพฺเภถ จ สงฺขาเรสุ “เอวํ เม สงฺขารา โหนฺตุ, เอวํ เม สงฺขารา มา อเหสุนฺ”ติ. ยสฺมา จ โข \csromanfolio{182}ภิกฺขเว สงฺขารา อนตฺตา, ตสฺมา สงฺขารา อาพาธาย สํวตฺตนฺติ, น จ ลพฺภติ สงฺขาเรสุ “เอวํ เม สงฺขารา โหนฺตุ, เอวํ เม สงฺขารา มา อเหสุนฺ”ติ.}

\proseitem{15}{วิญฺญาณํ อนตฺตา, วิญฺญาณญฺจ หิทํ ภิกฺขเว อตฺตา อภวิสฺส, นยิทํ วิญฺญาณํ อาพาธาย สํวตฺเตยฺย, ลพฺเภถ จ วิญฺญาเณ “เอวํ เม วิญฺญาณํ โหตุ, เอวํ เม วิญฺญาณํ มา อโหสี”ติ. ยสฺมา จ โข ภิกฺขเว วิญฺญาณํ อนตฺตา, ตสฺมา วิญฺญาณํ อาพาธาย สํวตฺตติ, น จ ลพฺภติ วิญฺญาเณ “เอวํ เม วิญฺญาณํ โหตุ, เอวํ เม วิญฺญาณํ มา อโหสี”ติ.}

\prose{วุตฺตเญฺหตํ ภควตา\csromandash{}นายํ ภิกฺขเว กาโย ตุมฺหากํ นปิ อญฺเญสํ\csromansharedfootnote{b793bb54dac6ccb6}{ปเรสํ (สฺยา) ปสฺส สํ 1. 294 ปิฏฺเฐ.}, ปุราณมิทํ ภิกฺขเว กมฺมํ อภิสงฺขตํ อภิส\-ญฺเจตยิตํ เวทนิยํ ทฏฺฐพฺพํ. ตตฺร โข ภิกฺขเว สุตวา อริยสาวโก ปฏิจฺจ\-สมุปฺปาทํเยว สาธุกํ โยนิโส มนสิ กโรติ “อิติ อิมสฺมิํ สติ อิทํ โหติ, อิมสฺสุปฺปาทา อิทํ อุปฺปชฺชติ. อิมสฺมิํ อสติ อิทํ น โหติ, อิมสฺส นิโรธา อิทํ นิรุชฺฌติ. ยทิทํ อวิชฺชาปจฺจยา สงฺขารา, สงฺขาร\-ปจฺจยา วิญฺญาณํ, วิญฺญาณปจฺจยา นามรูปํ, นามรูป\-ปจฺจยา สฬายตนํ, สฬายตน\-ปจฺจยา ผสฺโส, ผสฺสปจฺจยา เวทนา, เวทนาปจฺจยา ตณฺหา, ตณฺหาปจฺจยา อุปาทานํ, อุปาทานปจฺจยา ภโว, ภวปจฺจยา ชาติ, ชาติปจฺจยา ชรามรณํ โสกปริเทว\-ทุกฺข- โทมนสฺสุ\-ปายาสา สมฺภวนฺติ. เอวเมตสฺส เกวลสฺส ทุกฺข\-กฺขนฺธสฺส สมุทโย โหติ.}

\prose{อวิชฺชาย เตฺวว อเสส\-วิราค\-นิโรธา สงฺขาร\-นิโรโธ, สงฺขาร\-นิโรธา วิญฺญาณนิโรโธ ฯเปฯ ชาตินิโรธา ชรามรณํ โสกปริเทว\-ทุกฺข- โทมนสฺสุ\-ปายาสา นิรุชฺฌนฺติ. เอวเมตสฺส เกวลสฺส ทุกฺข\-กฺขนฺธสฺส นิโรโธ โหตี”ติ. เอวํ อวสิย\-ปวตฺต\-สลฺลกฺขณ\-วเสน สุญฺญโต โลกํ อเวกฺขติ.}

\prose{กถํ ตุจฺฉ\-สงฺขาร\-สมนุปสฺสนา\-วเสน สุญฺญโต โลกํ อเวกฺขติ\csromandash{}รูเป สาโร น ลพฺภติ. เวทนาย สาโร น ลพฺภติ. สญฺญาย สาโร น ลพฺภติ. สงฺขาเรสุ สาโร น ลพฺภติ. วิญฺญาเณ สาโร น ลพฺภติ. รูปํ อสฺสารํ นิสฺสารํ สาราปคตํ นิจฺจสาร\-สาเรน วา สุขสาร\-สาเรน วา อตฺตสารสาเรน วา นิจฺเจน วา ธุเวน วา สสฺสเตน วา \csromanfolio{183}อวิปริณาม\-ธมฺเมน วา. เวทนา อสฺสารา นิสฺสารา สาราปคตา. สญฺญา อสฺสารา นิสฺสารา สาราปคตา. สงฺขารา อสฺสารา นิสฺสารา สาราปคตา. วิญฺญาณํ อสฺสารํ นิสฺสารํ สาราปคตํ นิจฺจสาร\-สาเรน วา สุขสาร\-สาเรน วา อตฺตสารสาเรน วา นิจฺเจน วา ธุเวน วา สสฺสเตน วา อวิปริณาม\-ธมฺเมน วา. ยถา นโฬ อสฺสาโร นิสฺสาโร สาราปคโต. ยถา จ เอรณฺโฑ ฯเปฯ. ยถา จ อุทุมฺพโร อสฺสาโร นิสฺสาโร สาราปคโต. ยถา จ เสตกจฺโฉ\csromansharedfootnote{cf9a7867ce2ba6ba}{เสตวจฺโฉ (ก)} อสฺสาโร นิสฺสาโร สาราปคโต. ยถา จ ปาลิภทฺทโก\csromansharedfootnote{c92351b4a9ca92bd}{ปาฬิภทฺทโก (ก)} อสฺสาโร นิสฺสาโร สาราปคโต. ยถา จ เผณปิณฺโฑ\csromansharedfootnote{0fd8238710cb606a}{เผณุปิณฺโฑ (สฺยา)} อสฺสาโร นิสฺสาโร สาราปคโต. ยถา จ อุทกปุพฺพุฬํ\csromansharedfootnote{cb9ddd5507089988}{ปุพฺพุลกํ (สฺยา)} อสฺสารํ นิสฺสารํ สาราปคตํ. ยถา จ มรีจิ อสฺสารา นิสฺสารา สาราปคตา. ยถา กทลิกฺขนฺโธ อสฺสาโร นิสฺสาโร สาราปคโต. ยถา มายา อสฺสารา นิสฺสารา สาราปคตา. เอวเมว รูปํ อสฺสารํ นิสฺสารํ สาราปคตํ นิจฺจสาร\-สาเรน วา สุขสาร\-สาเรน วา อตฺตสารสาเรน วา นิจฺเจน วา ธุเวน วา สสฺสเตน วา อวิปริณาม\-ธมฺเมน วา. เวทนา อสฺสารา นิสฺสารา สาราปคตา. สญฺญา อสฺสารา นิสฺสารา สาราปคตา. สงฺขารา อสฺสารา นิสฺสารา สาราปคตา. วิญฺญาณํ อสฺสารํ นิสฺสารํ สาราปคตํ นิจฺจสาร\-สาเรน วา สุขสาร\-สาเรน วา อตฺตสารสาเรน วา นิจฺเจน วา ธุเวน วา สสฺสเตน วา อวิปริณาม\-ธมฺเมน วา. เอวํ ตุจฺฉ\-สงฺขาร\-สมนุปสฺสนา\-วเสน สุญฺญโต โลกํ อเวกฺขติ. อิเมหิ ทฺวีหิ การเณหิ สุญฺญโต โลกํ อเวกฺขติ.}

\prose{อปิ จ ฉหากาเรหิ สุญฺญโต โลกํ อเวกฺขติ. จกฺขุ สุญฺญํ\csromansharedfootnote{76911ba71a0f3d3c}{สฺยา...โปตฺถเก อิมสฺมิํ ฐาเน อญฺญถา ทิสฺสติ.} อตฺเตน วา อตฺตนิเยน วา นิจฺเจน วา ธุเวน วา สสฺสเตน วา อวิปริณาม\-ธมฺเมน วา. โสตํ สุญฺญํ. ฆานํ สุญฺญํ. ชิวฺหา สุญฺญา. กาโย สุญฺโญ. มโน สุญฺโญ อตฺเตน วา อตฺตนิเยน วา นิจฺเจน วา ธุเวน วา สสฺสเตน วา อวิปริณาม\-ธมฺเมน วา. รูปา สุญฺญา. สทฺทา สุญฺญา. คนฺธา สุญฺญา. รสา สุญฺญา. โผฏฺฐพฺพา สุญฺญา. ธมฺมา สุญฺญา อตฺเตน วา อตฺตนิเยน วา นิจฺเจน วา ธุเวน วา สสฺสเตน วา อวิปริณาม\-ธมฺเมน วา. จกฺขุวิญฺญาณํ สุญฺญํ ฯเปฯ. มโนวิญฺญาณํ สุญฺญํ. จกฺขุ\-สมฺผสฺโส \csromanfolio{184}สุญฺโญ ฯเปฯ. มโนสมฺผสฺโส สุญฺโญ. จกฺขุ\-สมฺผสฺสชา เวทนา สุญฺญา ฯเปฯ. มโน\-สมฺผสฺสชา เวทนา สุญฺญา. รูปสญฺญา สุญฺญา ฯเปฯ. ธมฺมสญฺญา สุญฺญา. รูปสญฺเจตนา สุญฺญา ฯเปฯ. ธมฺม\-สญฺเจตนา สุญฺญา. รูปตณฺหา สุญฺญา. รูปวิตกฺโก สุญฺโญ. รูปวิจาโร สุญฺโญ ฯเปฯ. ธมฺมวิจาโร สุญฺโญ อตฺเตน วา อตฺตนิเยน วา นิจฺเจน วา ธุเวน วา สสฺสเตน วา อวิปริณาม\-ธมฺเมน วา. เอวํ ฉหากาเรหิ สุญฺญโต โลกํ อเวกฺขติ.}

\prose{อปิ จ ทสหากาเรหิ สุญฺญโต โลกํ อเวกฺขติ, รูปํ ริตฺตโต ตุจฺฉโต สุญฺญโต อนตฺตโต อสารกโต วธกโต วิภวโต อฆมูลโต สาสวโต สงฺขตโต. เวทนํ. สญฺญํ. สงฺขาเร. วิญฺญาณํ. จุติํ. อุปปตฺติํ. ปฏิสนฺธิํ. ภวํ. สํสารวฏฺฏํ ริตฺตโต ตุจฺฉโต สุญฺญโต อนตฺตโต อสารกโต วธกโต วิภวโต อฆมูลโต สาสวโต สงฺขตโต. เอวํ ทสหากาเรหิ สุญฺญโต โลกํ อเวกฺขติ.}

\prose{อปิ จ ทฺวาทสหา\-กาเรหิ สุญฺญโต โลกํ อเวกฺขติ. รูปํ น สตฺโต, น ชีโว, น นโร, น มานโว, น อิตฺถี, น ปุริโส, น อตฺตา, น อตฺตนิยํ, นาหํ, น มม, น โกจิ, น กสฺสจิ. เวทนา. สญฺญา. สงฺขารา. วิญฺญาณํ น สตฺโต, น ชีโว, น นโร, น มานโว, น อิตฺถี, น ปุริโส, น อตฺตา, น อตฺตนิยํ, นาหํ, น มม, น โกจิ, น กสฺสจิ. เอวํ ทฺวาทสหา\-กาเรหิ สุญฺญโต โลกํ อเวกฺขติ.}

\prose{วุตฺตเญฺหตํ ภควตา\csromandash{}“ยํ ภิกฺขเว น ตุมฺหากํ, ตํ ปชหถ, ตํ โว ปหีนํ ทีฆรตฺตํ หิตาย สุขาย ภวิสฺสติ. กิญฺจ ภิกฺขเว น ตุมฺหากํ. รูปํ ภิกฺขเว น ตุมฺหากํ, ตํ ปชหถ. ตํ โว ปหีนํ ทีฆรตฺตํ หิตาย สุขาย ภวิสฺสติ. เวทนา ภิกฺขเว น ตุมฺหากํ, ตํ ปชหถ, สา โว ปหีนา ทีฆรตฺตํ หิตาย สุขาย ภวิสฺสติ. สญฺญา ภิกฺขเว น ตุมฺหากํ, ตํ ปชหถ, สา โว ปหีนา ทีฆรตฺตํ หิตาย สุขาย ภวิสฺสติ. สงฺขารา ภิกฺขเว น ตุมฺหากํ, เต ปชหถ, เต โว ปหีนา ทีฆรตฺตํ หิตาย สุขาย ภวิสฺสนฺติ. วิญฺญาณํ ภิกฺขเว น ตุมฺหากํ, ตํ ปชหถ, ตํ โว ปหีนํ ทีฆรตฺตํ หิตาย สุขาย ภวิสฺสติ. \csromanfolio{185}เสยฺยถาปิ\csromansharedfootnote{1f640093a1a2eb51}{ตํ กิํ มญฺญถ (สฺยา, ก) ปสฺส สํ 2. 28 ปิฏฺเฐ.} ภิกฺขเว ยํ อิมสฺมิํ เชตวเน ติณกฏฺฐ\-สาขา\-ปลาสํ, ตํ ชโน หเรยฺย วา ฑเหยฺย วา ยถาปจฺจยํ วา กเรยฺย, อปิ นุ ตุมฺหากํ เอวมสฺส ‘อเมฺห ชโน ตรติ วา ฑหติ วา ยถาปจฺจยํ วา กโรตี’ติ. โน เหตํ ภนฺเต. ตํ กิสฺส เหตุ, น หิ โน เอตํ ภนฺเต อตฺตา วา อตฺตนิยํ วาติ. เอวเมว โข ภิกฺขเว ยํ น ตุมฺหากํ, ตํ ปชหถ, ตํ โว ปหีนํ ทีฆรตฺตํ หิตาย สุขาย ภวิสฺสติ. กิญฺจ ภิกฺขเว น ตุมฺหากํ. รูปํ ภิกฺขเว น ตุมฺหากํ, ตํ ปชหถ, ตํ โว ปหีนํ ทีฆรตฺตํ หิตาย สุขาย ภวิสฺสติ. เวทนา. สญฺญา. สงฺขารา. วิญฺญาณํ ภิกฺขเว น ตุมฺหากํ, ตํ ปชหถ, ตํ โว ปหีนํ ทีฆรตฺตํ หิตาย สุขาย ภวิสฺสตี”ติ. เอวมฺปิ สุญฺญโต โลกํ อเวกฺขติ.}

\prose{อายสฺมา อานนฺโท ภควนฺตํ เอตทโวจ “สุญฺโญ\csromansharedfootnote{7ff4cdbe500ce353}{สุญฺญโต (ก) ปสฺส สํ 2. 280 ปิฏฺเฐ.} โลโก สุญฺโญ โลโก”ติ ภนฺเต วุจฺจติ, กิตฺตาวตา นุ โข ภนฺเต “สุญฺโญ โลโก”ติ วุจฺจตีติ. ยสฺมา จ โข อานนฺท สุญฺญํ อตฺเตน วา อตฺตนิเยน วา, ตสฺมา “สุญฺโญ โลโก”ติ วุจฺจติ. กิญฺจานนฺท สุญฺญํ อตฺเตน วา อตฺตนิเยน วา. จกฺขุ โข อานนฺท สุญฺญํ อตฺเตน วา อตฺตนิเยน วา. รูปา สุญฺญา. จกฺขุวิญฺญาณํ สุญฺญํ. จกฺขุ\-สมฺผสฺโส สุญฺโญ, ยมฺปิทํ จกฺขุ\-สมฺผสฺส\-ปจฺจยา อุปฺปชฺชติ เวทยิตํ สุขํ วา ทุกฺขํ วา อทุกฺขมสุขํ วา. ตมฺปิ สุญฺญํ อตฺเตน วา อตฺตนิเยน วา. โสตํ สุญฺญํ. สทฺทา สุญฺญา. ฆานํ สุญฺญํ. คนฺธา สุญฺญา. ชิวฺหา สุญฺญา. รสา สุญฺญา. กาโย สุญฺโญ. โผฏฺฐพฺพา สุญฺญา. มโน สุญฺโญ. ธมฺมา สุญฺญา. มโนวิญฺญาณํ สุญฺญํ. มโนสมฺผสฺโส สุญฺโญ, ยมฺปิทํ สุญฺญํ มโน\-สมฺผสฺส\-ปจฺจยา อุปฺปชฺชติ เวทยิตํ สุขํ วา ทุขํ วา อทุกฺขมสุขํ วา, ตมฺปิ สุญฺญํ อตฺเตน วา อตฺตนิเยน วา. ยสฺมา จ โข อานนฺท สุญฺญํ อตฺเตน วา อตฺตนิเยน วา, ตสฺมา “สุญฺโญ โลโก”ติ วุจฺจตีติ. เอวมฺปิ สุญฺญโต โลกํ อเวกฺขติ.}

\setlength{\csromangathapagewidth}{0pt}
\setlength{\csromangathawakwidth}{0pt}
\csromangathameasuregroup{“สุทฺธํ ธมฺมสมุปฺปาทํ, \\ ปสฺสนฺตสฺส ยถาภูตํ, \\ ติณกฏฺฐสมํ โลกํ, \\ นาญฺญํ\textsuperscript{1} ปตฺถยเต กิญฺจิ,}{\csromangathabat{“สุทฺธํ ธมฺมสมุปฺปาทํ,}{สุทฺธสงฺขารสนฺตติํ.} \\ \csromangathabat{ปสฺสนฺตสฺส ยถาภูตํ,}{น ภยํ โหติ คามณิ.} \\ \csromangathabat{ติณกฏฺฐสมํ โลกํ,}{ยทา ปญฺญาย ปสฺสติ.} \\ \csromangathabat{นาญฺญํ\textsuperscript{1} ปตฺถยเต กิญฺจิ,}{อญฺญตฺรปฺปฏิสนฺธิยา”ติ.}}
\csromangathasetleft{“สุทฺธํ ธมฺมสมุปฺปาทํ, \\ ปสฺสนฺตสฺส ยถาภูตํ, \\ ติณกฏฺฐสมํ โลกํ, \\ นาญฺญํ\textsuperscript{1} ปตฺถยเต กิญฺจิ,}
\csromangathagroup{\csromangathastanza{\csromangathabat{“สุทฺธํ ธมฺม\-สมุปฺปาทํ,}{สุทฺธสงฺขารสนฺตติํ.} \\ \csromangathabat{ปสฺสนฺตสฺส ยถาภูตํ,}{น ภยํ โหติ คามณิ.}}\csromangathastanza{\csromanfolio{186}\csromangathabat{ติณกฏฺฐสมํ โลกํ,}{ยทา ปญฺญาย ปสฺสติ.} \\ \csromangathabat{นาญฺญํ\csromansharedfootnote{4eab81647c8be343}{น อญฺญํ (สี, สฺยา, ก)} ปตฺถยเต กิญฺจิ,}{อญฺญตฺรปฺปฏิสนฺธิยา”ติ.}}}

\prose{เอวมฺปิ สุญฺญโต โลกํ อเวกฺขติ.}

\prose{วุตฺตเญฺหตํ ภควตา\csromansharedfootnote{54e6e2288588d648}{ปสฺส สํ 2. 401 ปิฏฺเฐ.}\csromandash{}เอวเมว โข ภิกฺขเว ภิกฺขุ รูปํ สมนฺเนสติ ยาวตา รูปสฺส คติ, เวทนํ สมนฺเนสติ ยาวตา เวทนาย คติ, สญฺญํ สมนฺเนสติ ยาวตา สญฺญาย คติ, สงฺขาเร สมนฺเนสติ ยาวตา สงฺขารานํ คติ, วิญฺญาณํ สมนฺเนสติ ยาวตา วิญฺญาณสฺส คติ, ตสฺส รูปํ\csromansharedfootnote{14d01270b1dea3ea}{ตสฺส ภิกฺขุโน รูปํ (สฺยา)} สมนฺเนสโต ยาวตา รูปสฺส คติ. เวทนํ ฯเปฯ. สญฺญํ. สงฺขาเร. วิญฺญาณํ สมนฺเนสโต ยาวตา วิญฺญาณสฺส คติ, ยมฺปิสฺส ตํ โหติ “อหนฺ”ติ วา “มมนฺ”ติ วา “อสฺมี”ติ วา, ตมฺปิ ตสฺส น โหตีติ. เอวมฺปิ สุญฺญโต โลกํ อเวกฺขติ.}

\prose{\textbf{สุญฺญโต โลกํ อเวกฺขสฺสู}ติ สุญฺญโต โลกํ อเวกฺขสฺสุ ปจฺจเวกฺขสฺสุ ทกฺขสฺสุ ตุเลหิ ตีเรหิ วิภาเวหิ วิภูตํ กโรหีติ สุญฺญโต โลกํ อเวกฺขสฺสุ.}

\prose{\textbf{โมฆราช สทา สโต}ติ \textbf{โมฆราชา}ติ ภควา ตํ พฺราหฺมณํ นาเมน อาลปติ. \textbf{สทา}ติ สพฺพกาลํ ฯเปฯ สจฺฉิเม วโยขนฺเธ. \textbf{สโต}ติ จตูหิ การเณหิ สโต. กาเย กายานุปสฺสนา\-สติปฏฺฐานํ ภาเวนฺโต สโต ฯเปฯ โส วุจฺจติ สโตติ โมฆราช สทา สโต.}

\prose{\textbf{อตฺตานุทิฏฺฐิํ อูหจฺจา}ติ อตฺตานุทิฏฺฐิ วุจฺจติ วีสติวตฺถุกา สกฺกายทิฏฺฐิ. อิธ อสฺสุตวา ปุถุชฺชโน อริยานํ อทสฺสาวี อริยธมฺมสฺส อโกวิโท อริยธมฺเม อวินีโต, สปฺปุริสานํ อทสฺสาวี สปฺปุริส\-ธมฺมสฺส อโกวิโท สปฺปุริส\-ธมฺเม อวินีโต รูปํ อตฺตโต สมนุปสฺสติ, รูปวนฺตํ วา อตฺตานํ, อตฺตนิ วา รูปํ, รูปสฺมิํ วา อตฺตานํ. เวทนํ. สญฺญํ. สงฺขาเร. วิญฺญาณํ อตฺตโต สมนุปสฺสติ, วิญฺญาณวนฺตํ วา อตฺตานํ, อตฺตนิ วา วิญฺญาณํ, วิญฺญาณสฺมิํ วา อตฺตานํ. ยา เอวรูปา ทิฏฺฐิ ทิฏฺฐิคตํ ทิฏฺฐิคหนํ ทิฏฺฐิกนฺตาโร ทิฏฺฐิ\-วิสูกายิกํ ทิฏฺฐิ\-วิปฺผนฺทิตํ ทิฏฺฐิ\-สํโยชนํ คาโห ปฏิคฺคาโห อภินิเวโส ปรามาโส กุมฺมคฺโค มิจฺฉาปโถ มิจฺฉตฺตํ ติตฺถายตนํ วิปริยาส\-คฺคาโห วิปรีตคฺคาโห วิปลฺลาสคฺคาโห มิจฺฉาคาโห \csromanfolio{187}อยาถาวกสฺมิํ ยาถาวกนฺติ คาโห ยาวตา ทฺวาสฏฺฐิ ทิฏฺฐิคตานิ, อยํ อตฺตานุทิฏฺฐิ. \textbf{อตฺตานุทิฏฺฐิํ อูหจฺจา}ติ อตฺตานุทิฏฺฐิํ อูหจฺจ สมูหจฺจ\csromansharedfootnote{966c7ac49c15d1d1}{อุหจฺจ สมุหจฺจ (ก) สทฺทนีติยา ปน สเมติ.} อุทฺธริตฺวา สมุทฺธริตฺวา อุปฺปาฏยิตฺวา สมุปฺปาฏยิตฺวา ปชหิตฺวา วิโนเทตฺวา พฺยนฺตีกริตฺวา อนภาวํ คเมตฺวาติ อตฺตานุทิฏฺฐิํ อูหจฺจ.}

\prose{\textbf{เอวํ มจฺจุตโร สิยา}ติ เอวํ มจฺจุปิ ตเรยฺยาสิ, ชราปิ ตเรยฺยาสิ, มรณมฺปิ ตเรยฺยาสิ อุตฺตเรยฺยาสิ ปตเรยฺยาสิ สมติกฺกเมยฺยาสิ วีติวตฺเตยฺยาสีติ เอวํ มจฺจุตโร สิยา.}

\prose{\textbf{เอวํ โลกํ อเวกฺขนฺตนฺ}ติ เอวํ โลกํ อเวกฺขนฺตํ ปจฺจเวกฺขนฺตํ อุลยนฺตํ ตีรยนฺตํ วิภาวยนฺตํ วิภูตํ กโรนฺตนฺติ เอวํ โลกํ อเวกฺขนฺตํ.}

\prose{\textbf{มจฺจุ ราชา น ปสฺสตี}ติ มจฺจุปิ มจฺจุราชา, มาโรปิ มจฺจุราชา, มรณมฺปิ มจฺจุราชา. \textbf{น ปสฺสตี}ติ มจฺจุราชา น ปสฺสติ น ทกฺขติ นาธิคจฺฉติ น วินฺทติ น ปฏิลภติ. วุตฺตเญฺหตํ ภควตา\csromandash{} “เสยฺยถาปิ ภิกฺขเว อารญฺญิโก มิโค อรญฺเญ ปวเน จรมาโน วิสฺสตฺโถ คจฺฉติ, วิสฺสตฺโถ\csromansharedfootnote{a60694d0e195324d}{วิสฺสฏฺโฐ (ก)} ติฏฺฐติ, วิสฺสตฺโถ นิสีทติ, วิสฺสตฺโถ เสยฺยํ กปฺเปติ. ตํ กิสฺส เหตุ. อนาปาถคโต\csromansharedfootnote{b3d768e70175998b}{อนาปาตคโต (ก)} ภิกฺขเว ลุทฺทสฺส. เอวเมว โข ภิกฺขเว ภิกฺขุ วิวิจฺเจว กาเมหิ วิวิจฺจ อกุสเลหิ ธมฺเมหิ สวิตกฺกํ สวิจารํ วิเวกชํ ปีติสุขํ ปฐมํ ฌานํ อุปสมฺปชฺช วิหรติ. อยํ วุจฺจติ ภิกฺขเว ภิกฺขุ อนฺธมกาสิ\csromansharedfootnote{332f81f43e4c4669}{อนฺตมกาสิ (ก) ปสฺส ม 1. 214 ปิฏฺเฐ.}, มารํ อปทํ วธิตฺวา มารจกฺขุํ อทสฺสนํ คโต ปาปิมโต.}

\prose{ปุน จปรํ ภิกฺขเว ภิกฺขุ วิตกฺก\-วิจารานํ วูปสมา อชฺฌตฺตํ สมฺปสาทนํ เจตโส เอโกทิภาวํ อวิตกฺกํ อวิจารํ สมาธิชํ ปีติสุขํ ทุติยํ ฌานํ ฯเปฯ ตติยํ ฌานํ. จตุตฺถํ ฌานํ อุปสมฺปชฺช วิหรติ. อยํ วุจฺจติ ภิกฺขเว ภิกฺขุ อนฺธมกาสิ มารํ, อปทํ วธิตฺวา มารจกฺขุํ อทสฺสนํ คโต ปาปิมโต.}

\csromanmatikaanchor{mtk.37}
\csromantocmarkref{subsection}{16. ปิงฺคิยมาณวปุจฺฉานิทฺเทส}{mtk.37}
\bookmark[dest={mtk.37},level=2]{16. ปิงฺคิยมาณวปุจฺฉานิทฺเทส}
\prose{ปุน จปรํ ภิกฺขเว ภิกฺขุ สพฺพโส รูปสญฺญานํ สมติกฺกมา ปฏิฆ\-สญฺญานํ อตฺถงฺคมา นานตฺต\-สญฺญานํ อมนสิการา “อนนฺโต อากาโส”ติ อากาสา\-นญฺจา\-ยตนํ อุปสมฺปชฺช วิหรติ. อยํ วุจฺจติ ภิกฺขเว \csromanfolio{188}ภิกฺขุ อนฺธมกาสิ มารํ, อปทํ วธิตฺวา มารจกฺขุํ อทสฺสนํ \textbf{คโต ปาปิมโต.}}

\prose{ปุน จปรํ ภิกฺขเว ภิกฺขุ สพฺพโส อากาสา\-นญฺจา\-ยตนํ สมติกฺกมฺม “อนนฺตํ วิญฺญาณนฺ”ติ วิญฺญาณญฺจา\-ยตนํ อุปสมฺปชฺช วิหรติ ฯเปฯ.}

\prose{ปุน จปรํ ภิกฺขเว สพฺพโส วิญฺญาณญฺจา\-ยตนํ สมติกฺกมฺม “นตฺถิ กิญฺจี”ติ อากิญฺจญฺญา\-ยตนํ อุปสมฺปชฺช วิหรติ ฯเปฯ.}

\prose{ปุน จปรํ ภิกฺขเว สพฺพโส อากิญฺจญฺญา\-ยตนํ สมติกฺกมฺม เนว\-สญฺญา\-นา\-สญฺญา\-ยตนํ อุปสมฺปชฺช วิหรติ ฯเปฯ.}

\prose{ปุน จปรํ ภิกฺขเว สพฺพโส เนว\-สญฺญา\-นา\-สญฺญา\-ยตนํ สมติกฺกมฺม สญฺญาเวทยิต\-นิโรธํ อุปสมฺปชฺช วิหรติ. ปญฺญาย จสฺส ทิสฺวา อาสวา ปริกฺขีณา โหนฺติ. อยํ วุจฺจติ ภิกฺขเว ภิกฺขุ อนฺธมกาสิ มารํ, อปทํ วธิตฺวา มารจกฺขุํ อทสฺสนํ คโต ปาปิมโต, ติณฺโณ โลเก วิสตฺติกนฺ”ติ. โส วิสฺสตฺโถ คจฺฉติ, วิสฺสตฺโถ ติฏฺฐติ, วิสฺสตฺโถ นิสีทติ, วิสฺสตฺโถ เสยฺยํ กปฺเปติ. ตํ กิสฺส เหตุ, อนาปาถคโต ภิกฺขุ ปาปิมโตติ มจฺจุราชา น ปสฺสติ. เตนาห ภควา\csromandash{}}

\setlength{\csromangathapagewidth}{0pt}
\setlength{\csromangathawakwidth}{0pt}
\csromangathameasuregroup{สุญฺญโต โลกํ อเวกฺขสฺสุ, \\ อตฺตานุทิฏฺฐิํ อูหจฺจ, \\ เอวํ โลกํ อเวกฺขนฺตํ,}{\csromangathabat{สุญฺญโต โลกํ อเวกฺขสฺสุ,}{โมฆราช สทา สโต.} \\ \csromangathabat{อตฺตานุทิฏฺฐิํ อูหจฺจ,}{เอวํ มจฺจุตโร สิยา.} \\ \csromangathabat{เอวํ โลกํ อเวกฺขนฺตํ,}{มจฺจุราชา น ปสฺสตีติ.}}
\csromangathasetleft{สุญฺญโต โลกํ อเวกฺขสฺสุ, \\ อตฺตานุทิฏฺฐิํ อูหจฺจ, \\ เอวํ โลกํ อเวกฺขนฺตํ,}
\csromangathagroup{\csromangathastanza{\csromangathabat{สุญฺญโต โลกํ อเวกฺขสฺสุ,}{โมฆราช สทา สโต.} \\ \csromangathabat{อตฺตานุทิฏฺฐิํ อูหจฺจ,}{เอวํ มจฺจุตโร สิยา.} \\ \csromangathabat{เอวํ โลกํ อเวกฺขนฺตํ,}{มจฺจุราชา น ปสฺสตีติ.}}}

\prose{สห คาถา\-ปริโยสานา ฯเปฯ “สตฺถา เม ภนฺเต ภควา, สาวโกหมสฺมี”ติ.}

\vspace{\tipitakaparskip}
\csromancenter{โมฆราชมาณวปุจฺฉา\-นิทฺเทโส ปนฺนรสโม.}
\csromanhangingbegin
\hangingheaditem{89}{\textbf{ชิณฺโณหมสฺมิ อพโล วีตวณฺโณ}\csromansharedfootnote{c7a373ec68310e32}{วิวณฺโณ (สฺยา)}\textbf{, (อิจฺจายสฺมา ปิงฺคิโย,)}}
\hangingbody{\textbf{เนตฺตา น สุทฺธา สวนํ น ผาสุ.}}
\hangingbody{\textbf{มาหํ นสฺสํ โมมุโห} อนฺตราว2,}
\hangingbody{\textbf{อาจิกฺข ธมฺมํ ยมหํ วิชญฺญํ.}}
\hangingbody{\textbf{ชาติชราย อิธ วิปฺปหานํ.}}
\csromanhangingend

\prose{\csromanfolio{189}\textbf{ชิณฺโณหมสฺมิ อพโล วีตวณฺโณ}ติ \textbf{ชิณฺโณหมสฺมี}ติ ชิณฺโณ วุฑฺโฒ มหลฺลโก อทฺธคโต วโยอนุปฺปตฺโต วีสวสฺส\-สติโก ชาติโย. อพโลติ ทุพฺพโล อปฺปพโล อปฺปถาโม. วีตวณฺโณติ วีตวณฺโณ วิคตวณฺโณ วิคจฺฉิ\-ตว\-ณฺโณ, ยา สา ปุริมา สุภา วณฺณนิภา, สา อนฺตรหิตา อาทีนโว ปาตุภูโตติ ชิณฺโณหมสฺมิ อพโล วีตวณฺโณ.}

\prose{\textbf{อิจฺจายสฺมา ปิงฺคิโย}ติ \textbf{อิจฺจา}ติ ปทสทฺธิ ฯเปฯ. \textbf{อายสฺมา}ติ ปิยวจนํ ฯเปฯ. \textbf{ปิงฺคิโย}ติ ตสฺส พฺราหฺมณสฺส นามํ ฯเปฯ อภิลาโปติ อิจฺจายสฺมา ปิงฺคิโย.}

\prose{\textbf{เนตฺตา น สุทฺธา สวนํ น ผาสู}ติ เนตฺตา อสุทฺธา อวิสุทฺธา อปริสุทฺธา อโวทาตา, โน ตถา จกฺขุนา รูเป ปสฺสามีติ เนตฺตา น สุทฺธา. สฺ\textbf{อวนํ น ผาสู}ติ โสตํ อสุทฺธํ อวิสุทฺธํ อปริสุทฺธํ อโวทาตํ, โน ตถา โสเตน สทฺทํ สุโณมีติ เนตฺตา น สุทฺธา สวนํ น ผาสุ.}

\prose{\textbf{มาหํ นสฺสํ โมมุโห อนฺตราวา}ติ \textbf{มาหํ นสฺสนฺ}ติ มาหํ นสฺสํ มาหํ วินสฺสํ มาหํ ปนสฺสํ. \textbf{โมมุโห}ติ โมหมุโห อวิชฺชาคโต อญฺญาณี อวิภาวี ทุปฺปญฺโญ. \textbf{อนฺตราวา}ติ ตุยฺหํ ธมฺมํ ทิฏฺฐิํ ปฏิปทํ มคฺคํ อนญฺญาย อนธิคนฺตฺวา อวิทิตฺวา อปฺปฏิลภิตฺวา อผสฺสยิตฺวา อสจฺฉิกริตฺวา อนฺตราเยว กาลํ กเรยฺยนฺติ มาหํ นสฺสํ โมมุโห อนฺตราว.}

\prose{\textbf{อาจิกฺข ธมฺมํ ยมหํ วิชญฺญนฺ}ติ \textbf{ธมฺมนฺ}ติ อาทิกลฺยาณํ มชฺเฌกลฺยาณํ ปริโยสาน\-กลฺยาณํ สาตฺถํ สพฺยญฺชนํ เกวล\-ปริปุณฺณํ ปริสุทฺธํ พฺรหฺมจริยํ จตฺตาโร สติปฏฺฐาเน จตฺตาโร สมฺมปฺปทฺ.หาเน จตฺตาโร อิทฺธิปาเท ปญฺจินฺทฺริยานิ ปญฺจ พลานิ สตฺต โพชฺฌงฺเค อริยํ อฏฺฐงฺคิกํ มคฺคํ นิพฺพานญฺจ นิพฺพานคามินิญฺจ ปฏิปทํ อาจิกฺขาหิ เทเสหิ ปญฺญเปหิ ปฏฺฐเปหิ วิวราหิ วิภชาหิ อุตฺตานีกโรหิ ปกาเสหีติ อาจิกฺข ธมฺมํ. \textbf{ยมหํ วิชญฺญนฺ}ติ ยมหํ ชาเนยฺยํ อาชาเนยฺยํ วิชาเนยฺยํ ปฏิวิชาเนยฺยํ ปฏิวิชฺเฌยฺยํ อธิคจฺเฉยฺยํ ผสฺเสยฺยํ สจฺฉิ\-กเรยฺยนฺติ อาจิกฺข ธมฺมํ ยมหํ วิชญฺญํ.}

\prose{\textbf{ชาติชราย อิธ วิปฺปหานนฺ}ติ อิเธว ชาติ\-ชรา\-มรณสฺส ปหานํ วูปสมํ ปฏินิสฺสคฺคํ ปฏิปฺปสฺสทฺธิํ อมตํ นิพฺพานนฺติ ชาติชราย อิธ วิปฺปหานํ. เตนาห โส พฺราหฺมโณ\csromandash{}}

\prose{\csromanfolio{190}ชิณฺโณหมสฺมิ อพโล วีตวณฺโณ, (อิจฺจายสฺมา ปิงฺคิโย,)}

\prose{เนตฺตา น สุทฺธา สวนํ น ผาสุ.}

\setlength{\csromangathapagewidth}{0pt}
\setlength{\csromangathawakwidth}{0pt}
\csromangathameasuregroup{}{มาหํ นสฺสํ โมมุโห อนฺตราว, \\ อาจิกฺข ธมฺมํ ยมหํ วิชญฺญํ. \\ ชาติชราย อิธ วิปฺปหานนฺติ.}
\csromangathameasurewak{มาหํ นสฺสํ โมมุโห อนฺตราว, \\ อาจิกฺข ธมฺมํ ยมหํ วิชญฺญํ. \\ ชาติชราย อิธ วิปฺปหานนฺติ.}
\setlength{\csromangathaleftcol}{0pt}
\csromangathagroup{\csromangathastanza{\csromangathawak{มาหํ นสฺสํ โมมุโห อนฺตราว, \\ อาจิกฺข ธมฺมํ ยมหํ วิชญฺญํ. \\ ชาติชราย อิธ วิปฺปหานนฺติ.}}}

\csromanhangingbegin
\hangingheaditem{90}{\textbf{ทิสฺวาน รูเปสุ วิหญฺญมาเน, (ปิงฺคิยาติ ภควา,)}}
\hangingbody{\textbf{รุปฺปนฺติ รูเปสุ ชนา ปมตฺตา.}}
\hangingbody{\textbf{ตสฺมา ตุวํ ปิงฺคิย อปฺปมตฺโต,}}
\hangingbody{\textbf{ชหสฺสุ รูปํ อปุนพฺภวาย.}}
\csromanhangingend

\prose{\textbf{ทิสฺวาน รูเปสุ วิหญฺญมาเน}ติ \textbf{รูปนฺ}ติ จตฺตาโร จ มหาภูตา จตุนฺนญฺจ มหาภูตานํ อุปาทาย รูปํ, สตฺตา รูปเหตุ รูปปฺปจฺจยา รูปการณา หญฺญนฺติ วิหญฺญนฺติ อุปหญฺญนฺติ อุปฆาติยนฺติ\csromansharedfootnote{4383ad3fc241a7b4}{อุปฆาตยนฺติ (สฺยา, ก)}, รูเป สติ วิวิธ\-กมฺมการณา\csromansharedfootnote{9e12bbbf0919be82}{วิวิธกมฺมกรณานิ (ก)} กาเรนฺติ, กสาหิปิ ตาเฬนฺติ, เวตฺเตหิปิ ตาเฬนฺติ, อฑฺฒ\-ทณฺฑเกหิปิ ตาเฬนฺติ, หตฺถมฺปิ ฉินฺทนฺติ, ปาทมฺปิ ฉินฺทนฺติ, หตฺถปาทมฺปิ ฉินฺทนฺติ, กณฺณมฺปิ ฉินฺทนฺติ, นาสมฺปิ ฉินฺทนฺติ, กณฺณนาสมฺปิ ฉินฺทนฺติ, พิลงฺค\-ถาลิกมฺปิ กโรนฺติ, สงฺข\-มุณฺฑิกมฺปิ กโรนฺติ, ราหุมุขมฺปิ กโรนฺติ, โชติมาลิกมฺปิ กโรนฺติ, หตฺถ\-ปชฺโชติกมฺปิ กโรนฺติ, เอรก\-วตฺติกมฺปิ กโรนฺติ, จีรก\-วาสิกมฺปิ กโรนฺติ, เอเณยฺยกมฺปิ กโรนฺติ, พฬิส\-มํสิกมฺปิ กโรนฺติ, กหาปณิกมฺปิ กโรนฺติ, ขารา\-ปตจฺฉิกมฺปิ\csromansharedfootnote{afb224293fec740b}{ขาราปฏิจฺฉิกมฺปิ (ก)} กโรนฺติ, ปลิฆ\-ปริวตฺติกมฺปิ กโรนฺติ, ปลาล\-ปีฐกมฺปิ กโรนฺติ, ตตฺเตนปิ เตเลน โอสิญฺจนฺติ, สุนเขหิปิ ขาทาเปนฺติ, ชีวนฺตมฺปิ สูเล อุตฺตาเสนฺติ, อสินาปิ สีสํ ฉินฺทนฺติ, เอวํ สตฺตา รูปเหตุ รูปปฺปจฺจยา รูปการณา หญฺญนฺติ วิหญฺญนฺติ อุปหญฺญนฺติ อุปฆาติยนฺติ. เอวํ หญฺญมาเน วิหญฺญมาเน อุปหญฺญมาเน อุปฆาติยมาเน ทิสฺวา ปสฺสิตฺวา ตุลยิตฺวา ตีรยิตฺวา วิภาวยิตฺวา วิภูตํ กตฺวาติ ทิสฺวาน รูเปสุ วิหญฺญมาเน.}

\prose{\textbf{ปิงฺคิยาติ ภควา}ติ ปิงฺคิยาติ ภควา ตํ พฺราหฺมณํ นาเมน อาลปติ. \textbf{ภควา}ติ คารวาธิวจนเมตํ ฯเปฯ สจฺฉิกา ปญฺญตฺติ, ยทิทํ ภควาติ ปิงฺคิยาติ ภควา.}

\prose{\csromanfolio{191}\textbf{รุปฺปนฺติ รูเปสุ ชนา ปมตฺตา}ติ \textbf{รุปฺปนฺตี}ติ รุปฺปนฺติ กุปฺปนฺติ ปีฬยนฺติ\csromansharedfootnote{4d40df3d4391de1f}{ปีฬิยนฺติ (สฺยา, ก)} ฆฏฺฏยนฺติ, พฺยาธิตา โทมนสฺสิตา โหนฺติ. จกฺขุโรเคน รุปฺปนฺติ กุปฺปนฺติ ปีฬยนฺติ ฆฏฺฏยนฺติ, พฺยาธิตา โทมนสฺสิตา โหนฺติ. โสตโรเคน ฯเปฯ. กายโรเคน ฯเปฯ. qอํสมกสวาตาตปสรีสปสมฺผสฺเสหิ รุปฺปนฺติ กุปฺปนฺติ ปีฬยนฺติ ฆฏฺฏยนฺติ, พฺยาธิตา โทมนสฺสิตา โหนฺตีติ รุปฺปนฺติ รูเปสุ.}

\prose{อถ วา จกฺขุสฺมิํ หียมาเน หายมาเน ปริหายมาเน เวมาเน\csromansharedfootnote{de71020e83277c33}{วิหายมาเน (ก)} วิคจฺฉมาเน อนฺตารธายมาเน รุปฺปนฺติ ฯเปฯ โทมนสฺสิตา โหนฺติ. โสตสฺมิํ. ฆานสฺมิํ. ชิวฺหาย. กายสฺมิํ. รูปสฺมิํ. สทฺทสฺมิํ. คนฺธสฺมิํ. รสสฺมิํ. โผฏฺฐพฺพสฺมิํ. กุลสฺมิํ. คณสฺมิํ. อาวาสสฺมิํ. ลาภสฺมิํ. ยสสฺมิํ. ปสํสาย. สุขสฺมิํ. จีวรสฺมิํ. ปิณฺฑปาตสฺมิํ. เสนาสนสฺมิํ. คิลานปจฺจย\-เภสชฺชปริกฺขารสฺมิํ หียมาเน หายมาเน ปริหายมาเน เวมาเน วิคจฺฉมาเน อนฺตร\-ธายมาเน รุปฺปนฺติ กุปฺปนฺติ ปีฬยนฺติ ฆฏฺฏยนฺติ, พฺยาธิตา โทมนสฺสิตา โหนฺตีติ เอวมฺปิ รุปฺปนฺติ รูเปสุ.}

\prose{\textbf{ชนา}ติ ขตฺติยา จ พฺราหฺมณา จ เวสฺสา จ สุทฺทา จ คหฏฺฐา จ ปพฺพชิตา จ เทวา จ มนุสฺสา จ. \textbf{ปมตฺตา}ติ ปมาโท วตฺตพฺโพ กาย\-ทุจฺจริเตน วา วจีทุจฺจริเตน วา มโน\-ทุจฺจริเตน วา ปญฺจสุ กามคุเณสุ จิตฺตสฺส โวสคฺโค โวสคฺคา\-นุปฺปทานํ กุสลานํ วา ธมฺมานํ ภาวนาย อสกฺก\-จฺจ\-กิริ\-ยตา อสาตจฺจกิริยตา อนฏฺฐิต\-กิริยตา โอลีนวุตฺติตา นิกฺขิตฺต\-จฺฉนฺทตา นิกฺขิตฺตธุรตา อนาเสวนา อภาวนา อพหุลีกมฺมํ\csromansharedfootnote{920ccfce69942fda}{อพหุลิกมฺมํ (ก)} อนธิฏฺฐานํ อนนุโยโค ปมาโท, โย เอวรูโป ปมาโท ปมชฺชนา ปมชฺชิตตฺตํ, อยํ วุจฺจติ ปมาโท, อิมินา ปมาเทน สมนฺนาคตา ชนา ปมตฺตาติ รุปฺปนฺติ รูเปสุ ชนา ปมตฺตา.}

\prose{\textbf{ตสฺมา ตุวํ ปิงฺคิย อปฺปมตฺโต}ติ ตสฺมาติ \textbf{ตสฺมา} ตํการณา ตํเหตุ ตปฺปจฺจยา ตํนิทานํ เอวํ อาทีนวํ สมฺปสฺสมาโน รูเปสูติ ตสฺมา ตุวํ ปิงฺคิย. \textbf{อปฺปมตฺโต}ติ สกฺกจฺจการี สาตจฺจการี ฯเปฯ อปฺปมาโท กุสเลสุ ธมฺเมสูติ ตสฺมา ตุวํ ปิงฺคิย อปฺปมตฺโต.}

\prose{\textbf{ชหสฺสุ รูปํ อปุนพฺภวายา}ติ \textbf{รูปนฺ}ติ จตฺตาโร จ มหาภูตา จตุนฺนญฺจ มหาภูตานํ อุปาทาย รูปํ. \textbf{ชหสฺสุ รูปนฺ}ติ ชหสฺสุ รูปํ, ปชหสฺสุ รูปํ, วิโนเทหิ รูปํ, พฺยนฺตีกโรหิ รูปํ, อนภาวํ คเมหิ รูปํ. \textbf{อปุนพฺภ}}

\prose{\csromanfolio{192}\textbf{วายา}ติ ยถา เต รูปํ อิเธว นิรุชฺเฌยฺย, ปุน ปฏิสนฺธิโก ภโว น นิพฺพตฺเตยฺย กามธาตุยา วา รูปธาตุยา วา อรูปธาตุยา วา, กามภเว วา รูปภเว วา อรูปภเว วา, สญฺญาภเว วา อสญฺญาภเว วา เนว\-สญฺญา\-นา\-สญฺญาภเว วา, เอกโวการภเว วา จตุโวการภเว วา ปญฺจโวการภเว วา, ปุน คติยา วา อุปปตฺติยา วา ปฏิสนฺธิยา วา ภเว วา สํสาเร วา วฏฺเฏ วา น ชเนยฺย น สญฺชเนยฺย น นิพฺพตฺเตยฺย นาภินิพฺพตฺเตยฺย อิเธว นิรุชฺเฌย วูปสเมยฺย อตฺถํ คจฺเฉยฺย ปฏิปฺปสฺสมฺเภยฺยาติ ชหสฺสุ รูปํ อปุนพฺภวาย. เตนาห ภควา\csromandash{}}

\prose{ทิสฺวาน รูเปสุ วิหญฺญมาเน, (ปิงฺคิยาติ ภควา,)}

\prose{รุปฺปนฺติ รูเปสุ ชนา ปมตฺตา.}

\setlength{\csromangathapagewidth}{0pt}
\setlength{\csromangathawakwidth}{0pt}
\csromangathameasuregroup{}{ตสฺมา ตุวํ ปิงฺคิย อปฺปมตฺโต, \\ ชหสฺสุ รูปํ อปุนพฺภวายาติ. \\ \textbf{ทิสา จตสฺโส วิทิสา จตสฺโส,} \\ \textbf{อุทฺธํ อโธ ทส ทิสา อิมาโย.} \\ \textbf{น ตุยฺหํ อทิฏฺฐํ อสุตํ อมุตํ,} \\ \textbf{อโถ อวิญฺญาตํ กิญฺจนมตฺถิ โลเก.} \\ \textbf{อาจิกฺข ธมฺมํ ยมหํ วิชญฺญํ,} \\ \textbf{ชาติชราย อิธ วิปฺปหานํ.}}
\csromangathameasurewak{ตสฺมา ตุวํ ปิงฺคิย อปฺปมตฺโต, \\ ชหสฺสุ รูปํ อปุนพฺภวายาติ. \\ \textbf{ทิสา จตสฺโส วิทิสา จตสฺโส,} \\ \textbf{อุทฺธํ อโธ ทส ทิสา อิมาโย.} \\ \textbf{น ตุยฺหํ อทิฏฺฐํ อสุตํ อมุตํ,} \\ \textbf{อโถ อวิญฺญาตํ กิญฺจนมตฺถิ โลเก.} \\ \textbf{อาจิกฺข ธมฺมํ ยมหํ วิชญฺญํ,} \\ \textbf{ชาติชราย อิธ วิปฺปหานํ.}}
\setlength{\csromangathaleftcol}{0pt}
\csromangathagroup{\csromangathastanza{\csromangathawak{ตสฺมา ตุวํ ปิงฺคิย อปฺปมตฺโต, \\ ชหสฺสุ รูปํ อปุนพฺภวายาติ.}}\csromangathastanza{\csromangathawak{\csromangathaitemline{91}{\textbf{ทิสา จตสฺโส วิทิสา จตสฺโส,}} \\ \csromangathacontline{\textbf{อุทฺธํ อโธ ทส ทิสา อิมาโย.}} \\ \csromangathacontline{\textbf{น ตุยฺหํ อทิฏฺฐํ อสุตํ อมุตํ,}} \\ \csromangathacontline{\textbf{อโถ อวิญฺญาตํ กิญฺจนมตฺถิ โลเก.}} \\ \csromangathacontline{\textbf{อาจิกฺข ธมฺมํ ยมหํ วิชญฺญํ,}} \\ \csromangathacontline{\textbf{ชาติชราย อิธ วิปฺปหานํ.}}}}}

\prose{\textbf{ทิสาจตสฺโส วิทิสา จตสฺโส, อุทฺธํ อโธ ทส ทิสา อิมาโย}ติ ทส ทิสา.}

\prose{\textbf{น ตุยฺหํ อทิฏฺฐํ อสุตํ อมุตํ, อโถ อวิญฺญาตํ กิญฺจนมตฺถิ โลเก}ติ น ตุยฺหํ อทิฏฺฐํ อสุตํ อมุตํ อวิญฺญาตํ กิญฺจิ อตฺตตฺโถ วา ปรตฺโถ วา อุภยตฺโถ วา ทิฏฺฐธมฺมิโก วา อตฺโถ สมฺปรายิโก วา อตฺโถ อุตฺตาโน วา อตฺโถ คมฺภีโร วา อตฺโถ คูโฬฺห วา อตฺโถ ปฏิจฺฉนฺโน วา อตฺโถ เนยฺโย วา อตฺโถ นีโต วา อตฺโถ อนวชฺโช วา อตฺโถ นิกฺกิเลโส วา อตฺโถ โวทาโน วา อตฺโถ ปรมตฺโถ วา อตฺโถ นตฺถิ น อตฺถิ น สํวิชฺชติ นุปลพฺภตีติ น ตุยฺหํ อทิฏฺฐํ อสุตํ อมุตํ อโถ อวิญฺญาตํ กิญฺจนมตฺถิ โลเก.}

\prose{\textbf{อาจิกฺข ธมฺมํ ยมหํ วิชญฺญนฺ}ติ \textbf{ธมฺมนฺ}ติ อาทิกลฺยาณํ ฯเปฯ นิพฺพานญฺจ นิพฺพานคามินิญฺจ ปฏิปทํ อาจิกฺขาหิ เทเสหิ ปญฺญเปหิ ปฏฺฐเปหิ วิวราหิ วิภชาหิ \csromanfolio{193}อุตฺตานีกโรหิ ปกาเสหิ. \textbf{ยมหํ วิชญฺญนฺ}ติ ยมหํ ชาเนยฺยํ อาชาเนยฺยํ วิชาเนยฺยํ ปฏิวิชาเนยฺยํ ปฏิวิชฺเฌยฺยํ อธิคจฺเฉยฺยํ ผสฺเสยฺยํ สจฺฉิ\-กเรยฺยนฺติ อาจิกฺข ธมฺมํ ยมหํ วิชญฺญํ.}

\prose{\textbf{ชาติชราย อิธ วิปฺปหานนฺ}ติ อิเธว ชาติ\-ชรา\-มรณสฺส ปหานํ วูปสมํ ปฏินิสฺสคฺคํ ปฏิปฺปสฺสทฺธิํ อมตํ นิพฺพานนฺติ ชาติชราย อิธ วิปฺปหานํ. เตนาห โส พฺราหฺมโณ\csromandash{}}

\setlength{\csromangathapagewidth}{0pt}
\setlength{\csromangathawakwidth}{0pt}
\csromangathameasuregroup{}{ทิสา จตสฺโส วิทิสา จตสฺโส, \\ อุทฺธํ อโธ ทส ทิสา อิมาโย. \\ น ตุยฺหํ อทิฏฺฐํ อสุตํ อมุตํ, \\ อโถ อวิญฺญาตํ กิญฺจนมตฺถิ โลเก. \\ อาจิกฺข ธมฺมํ ยมหํ วิชญฺญํ,}
\csromangathameasurewak{ทิสา จตสฺโส วิทิสา จตสฺโส, \\ อุทฺธํ อโธ ทส ทิสา อิมาโย. \\ น ตุยฺหํ อทิฏฺฐํ อสุตํ อมุตํ, \\ อโถ อวิญฺญาตํ กิญฺจนมตฺถิ โลเก. \\ อาจิกฺข ธมฺมํ ยมหํ วิชญฺญํ,}
\setlength{\csromangathaleftcol}{0pt}
\csromangathagroup{\csromangathastanza{\csromangathawak{ทิสา จตสฺโส วิทิสา จตสฺโส, \\ อุทฺธํ อโธ ทส ทิสา อิมาโย. \\ น ตุยฺหํ อทิฏฺฐํ อสุตํ อมุตํ, \\ อโถ อวิญฺญาตํ กิญฺจนมตฺถิ โลเก.}}\csromangathastanza{\csromangathawak{อาจิกฺข ธมฺมํ ยมหํ วิชญฺญํ,}}}

\prose{ชาติชราย อิธ วิปฺปหานนฺติ.}

\csromanhangingbegin
\hangingheaditem{92}{\textbf{ตณฺหาธิปนฺเน มนุเช เปกฺขมาโน, (ปิงฺคิยาติ ภควา,)}}
\hangingbody{\textbf{สนฺตาปชาเต ชรสา ปเรเต.}}
\hangingbody{\textbf{ตสฺมา ตุวํ ปิงฺคิย อปฺปมตฺโต,}}
\hangingbody{\textbf{ชหสฺสุ ตณฺหํ อปุนพฺภวาย.}}
\csromanhangingend

\prose{\textbf{ตณาธิปนฺเน มนุเช เปกฺขมาโน}ติ ตณฺหาติ รูป\textbf{ตณฺหา} ฯเปฯ ธมฺมตณฺหา. \textbf{ตณฺหาธิปนฺเน}ติ ตณฺหาธิปนฺเน\csromansharedfootnote{278171e04c23dc2b}{ตณฺหาย อธิปนฺเน (ก)} ตณฺหานุเค ตณฺหานุคเต ตณฺหานุสเฏ ตณฺหาย ปนฺเน ปฏิปนฺเน อภิภูเต ปริยาทินฺน\-จิตฺเต. \textbf{มนุเช}ติ สตฺตา\-ธิวจนํ. \textbf{เปกฺขมาโน}ติ เปกฺขมาโน ทกฺขมาโน โอโลกยมาโน นิชฺฌายมาโน อุปปริกฺขมาโนติ ตณฺหาธิปนฺเน มนุเช เปกฺขมาโน. \textbf{ปิงฺคิยาติ ภควา}ติ ปิงฺคิยาติ ภควา ตํ พฺราหฺมณํ นาเมน อาลปติ. \textbf{ภควา}ติ คารวาธิวจนเมตํ ฯเปฯ สจฺฉิกา ปญฺญตฺติ, ยทิทํ ภควาติ ปิงฺคิยาติ ภควา.}

\prose{\textbf{สนฺตาปชาเต ชรสา ปเรเต}ติ \textbf{สนฺตาปชาเต}ติ ชาติยา สนฺตาปชาเต, ชราย สนฺตาปชาเต, พฺยาธินา สนฺตาปชาเต, มรเณน สนฺตาปชาเต, โสก\-ปริเทว\-ทุกฺข\-โทมนสฺสุปายาเสหิ สนฺตาปชาเต, เนรยิเกน ทุกฺเขน สนฺตาปชาเต ฯเปฯ ทิฏฺฐิ\-พฺยสเนน ทุกฺเขน สนฺตาปชาเต \csromanfolio{194}อีติชาเต อุปทฺทวชาเต อุปสคฺคชาเตติ สนฺตาปชาเต. \textbf{ชรสา ปเรเต}ติ ชราย ผุฏฺเฐ ปเรเต สโมหิเต สมนฺนาคเต. ชาติยา อนุคเต ชราย อนุสเฏ พฺยาธินา อภิภูเต มรเณน อพฺภาหเต อตาเณ อเลเณ อสรเณ อสรณีภูเตติ สนฺตาปชาเต ชรสา ปเรเต.}

\prose{\textbf{ตสฺมา ตุวํ ปิงฺคิย อปฺปมตฺโต}ติ ตสฺมาติ \textbf{ตสฺมา} ตํการณา ตํเหตุ ตปฺปจฺจยา ตํนิทานา เอวํ อาทีนวํ สมฺปสฺสมาโน ตณฺหายาติ ตสฺมา ตุวํ ปิงฺคิย. \textbf{อปฺปมตฺโต}ติ สกฺกจฺจการี ฯเปฯ อปฺปมาโท กุสเลสุ ธมฺเมสูติ ตสฺมา ตุวํ ปิงฺคิย อปฺปมตฺโต.}

\prose{\textbf{ชหสฺสุ ตณฺหํ อปุนพฺภวายา}ติ ตณฺหาติ รูป\textbf{ตณฺหา} ฯเปฯ ธมฺมตณฺหา. \textbf{ชหสฺสุ ตณฺหนฺ}ติ ชหสฺสุ ตณฺหํ, ปชหสฺสุ ตณฺหํ, วิโนเทหิ ตณฺหํ, พฺยนฺตีกโรหิ ตณฺหํ, อนภาวํ คเมหิ ตณฺหํ. \textbf{อปุนพฺภวายา}ติ ยถา เต ฯเปฯ ปุนปฏิสนฺธิโก ภโว น นิพฺพตฺเตยฺย กามธาตุยา วา รูปธาตุยา วา อรูปธาตุยา วา กามภเว วา รูปภเว วา อรูปภเว วา สญฺญาภเว วา อสญฺญาภเว วา เนว\-สญฺญา\-นา\-สญฺญาภเว วา, เอกโวการภเว วา จตุโวการภเว วา ปญฺจโวการภเว วา, ปุนคติยา วา อุปปตฺติยา วา ปฏิสนฺธิยา วา ภเว วา สํสาเร วา วฏฺเฏ วา น ชเนยฺย น สญฺชเนยฺย น นิพฺพตฺเตยฺย นาภินิพฺพตฺเตยฺย อิเธว นิรุชฺเฌยฺย วูปสเมยฺย อตฺถํ คจฺเฉยฺย ปฏิปฺปสฺสมฺเภยฺยาติ ชหสฺสุ ตณฺหํ อปุนพฺภวาย. เตนาห ภควา\csromandash{}}

\prose{ตณฺหาธิปนฺเน มนุเช เปกฺขมาโน, (ปิงฺคิยาติ ภควา,)}

\prose{สนฺตาปชาเต ชรสา ปเรเต.}

\prose{ตสฺมา ตุวํ ปิงฺคิย อปฺปมตฺโต,}

\prose{ชหสฺสุ ตณฺหํ อปุนพฺภวายาติ.}

\csromanmatikaanchor{mtk.38}
\csromantocmarkref{subsection}{17. ปารายนตฺถุติคาถานิทฺเทส}{mtk.38}
\bookmark[dest={mtk.38},level=2]{17. ปารายนตฺถุติคาถานิทฺเทส}
\prose{สห คาถา\-ปริโยสานา เย เต พฺราหฺมเณน สทฺธิํ เอกจฺฉนฺทา เอกปโยคา เอกาธิปฺปายา เอก\-วาสน\-วาสิตา, เตสํ อเนก\-ปาณสหสฺสานํ วิรชํ วีตมลํ ธมฺมจกฺขุํ อุทปาทิ “ยํ กิญฺจิ สมุทย\-ธมฺมํ, สพฺพํ ตํ นิโรธธมฺมนฺ”ติ. ตสฺส จ พฺราหฺมณสฺส วิรชํ วีตมลํ ธมฺมจกฺขุํ อุทปาทิ “ยํ กิญฺจิ สมุทย\-ธมฺมํ, สพฺพํ ตํ นิโรธธมฺมนฺ”ติ, สห ธมฺม\-จกฺขุสฺส ปฏิลาภา อชิน\-ชฏา\-วากจีร\-ติทณฺฑ\-กมณฺฑลุ\-เกสา จ มสฺสู จ อนฺตรหิตา. \csromanfolio{195}\textbf{ภณฺฑุ}\-\textbf{กาสายวตฺถวสโน สงฺฆาฏิปตฺตจีวร}\-\textbf{ธโร อนฺวตฺถ}\-\textbf{ปฏิปตฺติยา} \textbf{ปญฺชลิโก ภควนฺตํ นมสฺสมาโน นิสินฺโน โหติ “สตฺถา เม ภนฺเต} \textbf{ภควา, สาวโกหมสฺมี”ติ.}}

\vspace{\tipitakaparskip}
\csromancenter{ปิงฺคิยมาณวปุจฺฉา\-นิทฺเทโส\csromansharedfootnote{44432171f283b7e5}{สิงฺคิยปญฺหํ (ก)} โสฬสโม.}
\proseitem{93}{\textbf{อิทมโวจ ภควา มคเธสุ วิหรนฺโต ปาสาณเก เจติเย,} \textbf{ปริจารก}\-\textbf{โสฬสานํ}\csromansharedfootnote{4b9934135eca454c}{ปริจาริกโสฬสนฺนํ (สฺยา, ก)}\textbf{ พฺราหฺมณานํ อชฺฌิฏฺโฐ ปุฏฺโฐ ปุฏฺโฐ} \textbf{ปญฺหํ พฺยากาสิ.}}

\prose{\textbf{อิทมโวจ ภควา}ติ อิทํ ปารายนํ อโวจ. \textbf{ภควา}ติ คารวาธิวจนเมตํ ฯเปฯ สจฺฉิกา ปญฺญตฺติ, ยทิทํ ภควาติ อิทมโวจ ภควา. \textbf{มคเธสุ วิหรนฺโต}ติ มคธนามเก ชนปเท วิหรนฺโต อิริยโต วตฺเตนฺโต ปาเลนฺโต ยเปนฺโต ยาเปนฺโต. \textbf{ปาสาณเก เจติเย}ติ ปาสาณก\-เจติยํ วุจฺจติ พุทฺธาสนนฺติ มคเธสุ วิหรนฺโต ปาสาณเก เจติเย. \textbf{ปริจารก}\-\textbf{โสฬสานํ พฺราหฺมณานนฺ}ติ ปิงฺคิโย\csromansharedfootnote{e1f2c2961de02dc0}{สิงฺคิโย (ก)} พฺราหฺมโณ พาวริสฺส พฺราหฺมณสฺส ปทฺโธปทฺธจโร ปริจารโก\csromansharedfootnote{00d93d2b2be417b1}{ปริจาริโก (สฺยา, ก)} สิสฺโส, ปิงฺคิเยน\csromansharedfootnote{55011d6bf07c9fb1}{เตน (ก)} เต โสฬสาติ เอวมฺปิ ปริจารก\-โสฬสานํ พฺราหฺมณานํ. อถ วา เต โสฬส พฺราหฺมณา พุทฺธสฺส ภควโต ปทฺธา ปทฺธจรา ปริจารกา สิสฺสาติ เอวมฺปิ ปริจารก\-โสฬสานํ พฺราหฺมณานํ.}

\prose{\textbf{อชฺฌิฏฺโฐ ปุฏฺโฐ ปุฏฺโฐ ปญฺหํ พฺยากาสี}ติ อชฺฌิฏฺโฐติ \textbf{อชฺฌิฏฺโฐ} อชฺเฌสิโต. \textbf{ปุฏฺโฐ ปุฏฺโฐ}ติ ปุฏฺโฐ ปุฏฺโฐ ปุจฺฉิโต ปุจฺฉิโต ยาจิโต ยาจิโต อชฺเฌสิโต อชฺเฌสิโต ปสาทิโต ปสาทิโต. \textbf{ปญฺหํ พฺยากาสี}ติ ปญฺหํ พฺยากาสิ อาจิกฺขิ เทเสสิ ปญฺญเปสิ ปฏฺฐเปสิ วิวริ วิภชิ อุตฺตานีอากาสิ ปกาเสสีติ อชฺฌิฏฺโฐ ปุฏฺโฐ ปุฏฺโฐ ปญฺหํ พฺยากาสิ. เตเนตํ วุจฺจติ\csromandash{}}

\prose{\textbf{อิทมโวจ ภควา มคเธสุ วิหรนฺโต ปาสาณเก เจติเย,}}

\prose{ปริจารก\-โสฬสานํ พฺราหฺมณานํ อชฺฌิฏฺโฐ ปุฏฺโฐ ปุฏฺโฐ ปญฺหํ}

\prose{พฺยากาสีติ.}

\proseitem{94}{\csromanfolio{196}\textbf{เอกเมกสฺส เจปิ ปญฺหสฺส อตฺถมญฺญาย ธมฺมมญฺญาย ธมฺมา}\-\textbf{นุธมฺมํ ปฏิปชฺเชยฺย, คจฺเฉยฺเยว ชรามรณสฺส ปารํ, ปารงฺคมนียา อิเม ธมฺมาติ, ตสฺมา อิมสฺส ธมฺม}\-\textbf{ปริยายสฺส ปารายนนฺเตว อธิวจนํ.}}

\prose{\textbf{เอกเมกสฺส เจปิ ปญฺหสฺสา}ติ เอกเมกสฺส เจปิ อชิตปญฺหสฺส เอกเมกสฺส เจปิ ติสฺส\-เมตฺเตยฺยปญฺหสฺส เอกเมกสฺส เจปิ ปุณฺณก\-ปญฺหสฺส เอกเมกสฺส เจปิ เมตฺตคู\-ปญฺหสฺส เอกเมกสฺส เจปิ โธตก\-ปญฺหสฺส เอกเมกสฺส เจปิ อุปสีว\-ปญฺหสฺส เอกเมกสฺส เจปิ นนฺทก\-ปญฺหสฺส เอกเมกสฺส เจปิ เหมกปญฺหสฺส เอกเมกสฺส เจปิ โตเทยฺย\-ปญฺหสฺส เอกเมกสฺส เจปิ กปฺปปญฺหสฺส เอกเมกสฺส เจปิ ชตุกณฺณิกปญฺหสฺส เอกเมกสฺส เจปิ ภทฺราวุธ\-ปญฺหสฺส เอกเมกสฺส เจปิ อุทยปญฺหสฺส เอกเมกสฺส เจปิ โปสาลปญฺหสฺส เอกเมกสฺส เจปิ โมฆราช\-ปญฺหสฺส เอกเมกสฺส เจปิ ปิงฺคิย\-ปญฺหสฺสาติ เอกเมกสฺส เจปิ ปญฺหสฺส.}

\prose{\textbf{อตฺถมญฺญาย ธมฺมมญฺญายา}ติ เสฺวว ปโญฺห ธมฺโม, วิสชฺชนํ\csromansharedfootnote{0000864cffdec293}{วิสฺสชฺชนํ (ก)} อตฺโถติ อตฺถํ อญฺญาย ชานิตฺวา ตุลยิตฺวา ตีรยิตฺวา วิภาวยิตฺวา วิภูตํ กตฺวาติ อตฺถมญฺญาย. \textbf{ธมฺมมญฺญายา}ติ ธมฺมํ อญฺญาย ชานิตฺวา ตุลยิตฺวา ตีรยิตฺวา วิภาวยิตฺวา วิภูตํ กตฺวาติ ธมฺมมญฺญายาติ อตฺถมญฺญาย ธมฺมมญฺญาย. \textbf{ธมฺมา}\-\textbf{นุธมฺมํ ปฏิปชฺเชยฺยา}ติ สมฺมาปฏิปทํ อนุโลม\-ปฏิปทํ อปจฺจนีกปฏิปทํ อนฺวตฺถ\-ปฏิปทํ ธมฺมานุธมฺม\-ปฏิปทํ ปฏิปชฺเชยฺยาติ ธมฺมา\-นุธมฺมํ ปฏิปชฺเชยฺย. \textbf{คจฺเฉยฺเยว ชรามรณสฺส ปารนฺ}ติ ชรามรณสฺส ปารํ วุจฺจติ อมตํ นิพฺพานํ. โย โส สพฺพ\-สงฺขาร\-สมโถ สพฺพู\-ปธิ\-ปฺปฏินิสฺสคฺโค ตณฺหกฺขโย วิราโค นิโรโธ นิพฺพานํ. \textbf{คจฺเฉยฺเยว ชรามรณสฺส ปารนฺ}ติ ชรามรณสฺส ปารํ คจฺเฉยฺย, ปารํ อธิคจฺเฉยฺย, ปารํ อธิผสฺเสยฺย, ปารํ สจฺฉิกเรยฺยาติ คจฺเฉยฺเยว ชรามรณสฺส ปารํ. \textbf{ปารงฺคมนียา อิเม ธมฺมา}ติ อิเม ธมฺมา ปารงฺคมนียา, ปารํ ปาเปนฺติ, ปารํ สมฺปาเปนฺติ, ปารํ สมนุปาเปนฺติ, ชรามรณสฺส ตรณาย\csromansharedfootnote{36320f34793877a8}{ตารณาย (สฺยา)} สํวตฺตนฺตีติ ปารงฺคมนียา อิเม ธมฺมาติ.}

\prose{\textbf{ตสฺมา อิมสฺส ธมฺมปริยายสฺสา}ติ ตสฺมาติ \textbf{ตสฺมา} ตํการณา ตํเหตุ ตปฺปจฺจยา ตํนิทานาติ ตสฺมา. \textbf{อิมสฺส ธมฺมปริยายสฺสา}ติ \csromanfolio{197}อิมสฺส ปารายนสฺสาติ ตสฺมา อิมสฺส ธมฺม\-ปริยายสฺส. \textbf{ปารายนนฺเตว อธิวจนนฺ}ติ ปารํ วุจฺจติ อมตํ นิพฺพานํ ฯเปฯ นิโรโธ นิพฺพานํ. อยนํ วุจฺจติ มคฺโค. เสยฺยถิทํ, สมฺมาทิฏฺฐิ ฯเปฯ สมฺมาสมาธิ. \textbf{อธิวจนนฺ}ติ สงฺขา สมญฺญา ปญฺญตฺติ โวหาโร นามํ นามกมฺมํ นามเธยฺยํ นิรุตฺติ พฺยญฺชนํ อภิลาโปติ ปารายนนฺเตว อธิวจนํ. เตเนตํ วุจฺจติ\csromandash{}}

\prose{เอกเมกสฺส เจปิ ปญฺหสฺส, อตฺถมญฺญาย ธมฺมมญฺญาย}

\prose{ธมฺมา\-นุธมฺมํ ปฏิปชฺเชยฺย, คจฺเฉยฺเยว ชรามรณสฺส}

\prose{ปารํ, ปารงฺคมนียา อิเม ธมฺมาติ, ตสฺมา อิมสฺส}

\prose{ธมฺม\-ปริยายสฺส ปารายนนฺเตว อธิวจนนฺติ.}

\setlength{\csromangathapagewidth}{0pt}
\setlength{\csromangathawakwidth}{0pt}
\csromangathameasuregroup{\textbf{อชิโต ติสฺสเมตฺเต}ยฺโย, \\ \textbf{โธตโก อุปสีโว} จ, \\ โตเทยฺยกปฺปา ทุภโย, \\ \textbf{ภทฺราวุโธ อุทโย} จ, \\ \textbf{โมฆราชา จ เมทฺ}หาวี, \\ \textbf{เอเต พุทฺธํ อุปาคจฺ}ฉุํ, \\ \textbf{ปุจฺฉนฺตา นิปุเณ ปฺ}อเญฺห,}{\csromangathabat{\textbf{อชิโต ติสฺสเมตฺเต}ยฺโย,}{ปุณฺณโก อถ เมตฺตคู.} \\ \csromangathabat{\textbf{โธตโก อุปสีโว} จ,}{นนฺโท จ อถ เหมโก.} \\ \csromangathabat{โตเทยฺยกปฺปา ทุภโย,}{ชตุกณฺณี จ ปณฺฑิโต.} \\ \csromangathabat{\textbf{ภทฺราวุโธ อุทโย} จ,}{โปสาโล จาปิ พฺราหฺมโณ.} \\ \csromangathabat{\textbf{โมฆราชา จ เมทฺ}หาวี,}{ปิงฺคิโย จ มหาอิสิ.} \\ \csromangathabat{\textbf{เอเต พุทฺธํ อุปาคจฺ}ฉุํ,}{สมฺปนฺนจรณํ อิสิํ.} \\ \csromangathabat{\textbf{ปุจฺฉนฺตา นิปุเณ ปฺ}อเญฺห,}{พุทฺธเสฏฺฐํ อุปาคมุํ.}}
\csromangathasetleft{\textbf{อชิโต ติสฺสเมตฺเต}ยฺโย, \\ \textbf{โธตโก อุปสีโว} จ, \\ โตเทยฺยกปฺปา ทุภโย, \\ \textbf{ภทฺราวุโธ อุทโย} จ, \\ \textbf{โมฆราชา จ เมทฺ}หาวี, \\ \textbf{เอเต พุทฺธํ อุปาคจฺ}ฉุํ, \\ \textbf{ปุจฺฉนฺตา นิปุเณ ปฺ}อเญฺห,}
\csromangathagroup{\csromangathastanza{\csromangathaitembat{95}{\textbf{อชิโต ติสฺสเมตฺเต}ยฺโย,}{ปุณฺณโก อถ เมตฺตคู.} \\ \csromangathacontbat{\textbf{โธตโก อุปสีโว} จ,}{นนฺโท จ อถ เหมโก.}}\csromangathastanza{\csromangathaitembat{96}{โตเทยฺยกปฺปา ทุภโย,}{ชตุกณฺณี จ ปณฺฑิโต.} \\ \csromangathacontbat{\textbf{ภทฺราวุโธ อุทโย} จ,}{โปสาโล จาปิ พฺราหฺมโณ.} \\ \csromangathacontbat{\textbf{โมฆราชา จ เมทฺ}หาวี,}{ปิงฺคิโย จ มหาอิสิ.}}\csromangathastanza{\csromangathaitembat{97}{\textbf{เอเต พุทฺธํ อุปาคจฺ}ฉุํ,}{สมฺปนฺน\-จรณํ อิสิํ.} \\ \csromangathacontbat{\textbf{ปุจฺฉนฺตา นิปุเณ ปฺ}อเญฺห,}{พุทฺธเสฏฺฐํ อุปาคมุํ.}}}

\prose{\textbf{เอเต พุทฺธํ อุปาคจฺฉุนฺ}ติ \textbf{เอเต}ติ โสฬส ปารายนิยา พฺราหฺมณา. พุทฺโธติ โย โส ภควา สยมฺภู อนาจริยโก ปุพฺเพ อนนุสฺสุเตสุ ธมฺเมสุ สามํ สจฺจานิ อภิสมฺพุชฺฌิ, ตตฺถ จ สพฺพญฺญุตํ ปตฺโต พเลสุ จ วสีภาวํ. พุทฺโธติ เกนฏฺเฐน พุทฺโธ, พุชฺฌิตา สจฺจานีติ พุทฺโธ, โพเธตา ปชายาติ พุทฺโธ, สพฺพญฺญุตาย พุทฺโธ, สพฺพทสฺสาวิ\-ตาย พุทฺโธ, อภิญฺเญยฺยตาย พุทฺโธ, วิสวิตาย พุทฺโธ, ขีณาสว\-สงฺขาเตน พุทฺโธ, นิรุปกฺกิเลส\-สงฺขาเตน พุทฺโธ, เอกนฺตวีตราโคติ พุทฺโธ, เอกนฺตวีตโทโสติ พุทฺโธ, เอกนฺตวีตโมโหติ พุทฺโธ, เอกนฺตนิกฺกิเลโสติ พุทฺโธ, เอกายนมคฺคํ คโตติ พุทฺโธ, เอโก อนุตฺตรํ สมฺมาสมฺโพธิํ อภิสมฺพุทฺโธติ พุทฺโธ, อพุทฺธิ\-วิห\-ตตฺตา พุทฺธิ\-ปฏิลาภาติ พุทฺโธ. พุทฺโธติ เนตํ นามํ มาตรา กตํ, น ปิตรา กตํ, น ภาตรา กตํ, น ภคินิยา กตํ, น มิตฺตามจฺเจหิ กตํ, น ญาติสาโลหิเตหิ กตํ, น สมณ\-พฺราหฺมเณหิ กตํ, น เทวตาหิ กตํ, วิโมกฺข\-นฺติกเมตํ พุทฺธานํ ภควนฺตานํ โพธิยา มูเล สห สพฺพญฺญุต\-ญาณสฺส ปฏิลาภา สจฺฉิกา \csromanfolio{198}ปญฺญตฺติ, ยทิทํ พุทฺโธติ. \textbf{เอเต พุทฺธํ อุปาคจฺฉุนฺ}ติ เอเต พุทฺธํ อุปาคมิํสุ อุปสงฺกมิํสุ ปยิรุปาสิํสุ ปริปุจฺฉิํสุ ปริปญฺหิํสูติ เอเต พุทฺธํ อุปาคจฺฉุํ.}

\prose{\textbf{สมฺปนฺน}\-\textbf{จรณํ อิสินฺ}ติ \textbf{จรณํ} วุจฺจติ สีลา\-จาร\-นิพฺพตฺติ. สีลสํวโรปิ จรณํ, อินฺทฺริย\-สํวโรปิ จรณํ, โภชเน มตฺตญฺญุตาปิ จรณํ, ชาคริยานุโยโคปิ จรณํ, สตฺตปิ สทฺธมฺมา จรณํ, จตฺตาริปิ ฌานานิ จรณํ. \textbf{สมฺปนฺน}\-\textbf{จรณนฺ}ติ สมฺปนฺน\-จรณํ เสฏฺฐจรณํ วิเส\-ฏฺฐ\-จรณํ\csromansharedfootnote{2a3d5c2416e57895}{วิสิฏฺฐจรณํ (ก)} ปาโมกฺข\-จรณํ อุตฺตมจรณํ ปวรจรณํ. \textbf{อิสี}ติ อิสิ, ภควา มหนฺตํ สีลกฺขนฺธํ เอสี คเวสี ปริเยสีติ อิสิ ฯเปฯ มเหสกฺเขหิ วา สตฺเตหิ เอสิโต คเวสิโต ปริเยสิโต “กหํ พุทฺโธ กหํ ภควา กหํ เทวเทโว กหํ นราสโภ”ติ อิสีติ สมฺปนฺน\-จรณํ อิสิํ.}

\prose{\textbf{ปุจฺฉนฺตา นิปุเณ ปเญฺห}ติ \textbf{ปุจฺฉนฺตา}ติ ปุจฺฉนฺตา ยาจนฺตา อชฺเฌสนฺตา ปสาเทนฺตา. \textbf{นิปุเณ ปเญฺห}ติ คมฺภีเร ทุทฺทเส ทุรนุโพเธ สนฺเต ปณีเต อตกฺกาวจเร นิปุเณ ปณฺฑิต\-เวทนีเย ปเญฺหติ ปุจฺฉนฺตา นิปุเณ ปเญฺห.}

\prose{\textbf{พุทฺธเสฏฺฐํ อุปาคมุนฺ}ติ พุทฺโธติ โย โส ภควา ฯเปฯ สจฺฉิกา ปญฺญตฺติ, ยทิทํ พุทฺโธติ. \textbf{เสฏฺฐนฺ}ติ อคฺคํ เสฏฺฐํ วิเสฏฺฐํ ปาโมกฺขํ อุตฺตมํ ปวรํ พุทฺธํ อุปาคมุํ อุปาคมิํสุ อุปสงฺกมิํสุ ปยิรุปาสิํสุ ปริปุจฺฉิํสุ ปริปญฺหิํสูติ พุทฺธเสฏฺฐํ อุปาคมุํ. เตเนตํ วุจฺจติ\csromandash{}}

\setlength{\csromangathapagewidth}{0pt}
\setlength{\csromangathawakwidth}{0pt}
\csromangathameasuregroup{เอเต พุทฺธํ อุปาคจฺฉุํ, \\ ปุจฺฉนฺตา นิปุเณ ปเญฺห, \\ \textbf{เตสํ พุทฺโธ ปพฺยากา}สิ, \\ \textbf{ปญฺหานํ เวยฺยากร}เณน,}{\csromangathabat{เอเต พุทฺธํ อุปาคจฺฉุํ,}{สมฺปนฺน\textbf{จรณํ} อิสิํ.} \\ \csromangathabat{ปุจฺฉนฺตา นิปุเณ ปเญฺห,}{\textbf{พุทฺธเสฏฺฐํ อุปาคมุนฺ}ติ.} \\ \csromangathabat{\textbf{เตสํ พุทฺโธ ปพฺยากา}สิ,}{\textbf{ปญฺหํ ปุฏฺโฐ} ยถาตถํ.} \\ \csromangathabat{\textbf{ปญฺหานํ เวยฺยากร}เณน,}{โตเสสิ พฺราหฺมเณ มุนิ.}}
\csromangathasetleft{เอเต พุทฺธํ อุปาคจฺฉุํ, \\ ปุจฺฉนฺตา นิปุเณ ปเญฺห, \\ \textbf{เตสํ พุทฺโธ ปพฺยากา}สิ, \\ \textbf{ปญฺหานํ เวยฺยากร}เณน,}
\csromangathagroup{\csromangathastanza{\csromangathabat{เอเต พุทฺธํ อุปาคจฺฉุํ,}{สมฺปนฺน\-\textbf{จรณํ} อิสิํ.} \\ \csromangathabat{ปุจฺฉนฺตา นิปุเณ ปเญฺห,}{\textbf{พุทฺธเสฏฺฐํ อุปาคมุนฺ}ติ.}}\csromangathastanza{\csromangathaitembat{98}{\textbf{เตสํ พุทฺโธ ปพฺยากา}สิ,}{\textbf{ปญฺหํ ปุฏฺโฐ} ยถาตถํ.} \\ \csromangathacontbat{\textbf{ปญฺหานํ เวยฺยากร}เณน,}{โตเสสิ พฺราหฺมเณ มุนิ.}}}

\prose{\textbf{เตสํ พุทฺโธ ปพฺยากาสี}ติ \textbf{เตสนฺ}ติ โสฬสานํ ปารายนิยานํ พฺราหฺมณานํ. พุทฺโธติ โย โส ภควา ฯเปฯ สจฺฉิกา ปญฺญตฺติ, ยทิทํ พุทฺโธติ. \textbf{ปพฺยากาสี}ติ เตสํ พุทฺโธ ปพฺยากาสิ อาจิกฺขิ เทเสสิ ปญฺญเปสิ ปฏฺฐเปสิ วิวริ วิภชิ อุตฺตานีอกาสิ ปกาเสสีติ เตสํ พุทฺโธ ปพฺยากาสิ.}

\prose{\textbf{ปญฺหํ ปุฏฺโฐ ยถาตถนฺ}ติ ปญฺหํ ปุฏฺโฐติ \textbf{ปญฺหํ ปุฏฺโฐ} ปุจฺฉิโต ยาจิโต อชฺเฌสิโต ปสาทิโต. \textbf{ยถาตถนฺ}ติ ยถา อาจิกฺขิตพฺพํ, ตถา \csromanfolio{199}อาจิกฺขิ. ยถา เทสิตพฺพํ, ตถา เทเสสิ. ยถา ปญฺญเปตพฺพํ, ตถา ปญฺญเปสิ. ยถา ปฏฺฐเปตพฺพํ, ตถา ปฏฺฐเปสิ. ยถา วิวริตพฺพํ, ตถา วิวริ. ยถา วิภชิตพฺพํ, ตถา วิภชิ. ยถา อุตฺตานีกาตพฺพํ, ตถา อุตฺตานีอกาสิ. ยถา ปกาสิตพฺพํ, ตถา ปกาเสสีติ ปญฺหํ ปุฏฺโฐ ยถาตถํ.}

\prose{\textbf{ปญฺหานํ เวยฺยากรเณนา}ติ ปญฺหานํ เวยฺยากรเณน อาจิกฺขเนน เทสเนน ปญฺญปเนน ปฏฺฐปเนน วิวรเณน วิภชเนน อุตฺตานีกมฺเมน ปกาสเนนาติ ปญฺหานํ เวยฺยากรเณน.}

\prose{\textbf{โตเสสิ พฺราหฺมเณ มุนี}ติ \textbf{โตเสสี}ติ โตเสสิ วิโตเสสิ ปสาเทสิ อาราเธสิ อตฺตมเน อกาสิ. \textbf{พฺราหฺมเณ}ติ โสฬส ปารายนิเย พฺราหฺมเณ. \textbf{มุนี}ติ โมนํ วุจฺจติ ญาณํ ฯเปฯ สงฺค\-ชาลม\-ติจฺจ โส มุนีติ โตเสสิ พฺราหฺมเณ มุนิ. เตเนตํ วุจฺจติ\csromandash{}}

\setlength{\csromangathapagewidth}{0pt}
\setlength{\csromangathawakwidth}{0pt}
\csromangathameasuregroup{เตสํ พุทฺโธ ปพฺยากาสิ, \\ ปญฺหานํ เวยฺยากรเณน, \\ \textbf{เต โตสิตา จกฺขุมฺ}อตา, \\ \textbf{พฺรหฺมจริยมจรฺ}อิํสุ,}{\csromangathabat{เตสํ พุทฺโธ ปพฺยากาสิ,}{ปญฺหํ ปุฏฺโฐ ยถาตถํ.} \\ \csromangathabat{ปญฺหานํ เวยฺยากรเณน,}{โตเสสิ พฺราหฺมเณ \textbf{มุนี}ติ.} \\ \csromangathabat{\textbf{เต โตสิตา จกฺขุมฺ}อตา,}{\textbf{พุทฺเธนา}\textbf{ทิจฺจพนฺธุนา}.} \\ \csromangathabat{\textbf{พฺรหฺมจริยมจรฺ}อิํสุ,}{วรปญฺญสฺส สนฺติเก.}}
\csromangathasetleft{เตสํ พุทฺโธ ปพฺยากาสิ, \\ ปญฺหานํ เวยฺยากรเณน, \\ \textbf{เต โตสิตา จกฺขุมฺ}อตา, \\ \textbf{พฺรหฺมจริยมจรฺ}อิํสุ,}
\csromangathagroup{\csromangathastanza{\csromangathabat{เตสํ พุทฺโธ ปพฺยากาสิ,}{ปญฺหํ ปุฏฺโฐ ยถาตถํ.} \\ \csromangathabat{ปญฺหานํ เวยฺยากรเณน,}{โตเสสิ พฺราหฺมเณ \textbf{มุนี}ติ.}}\csromangathastanza{\csromangathaitembat{99}{\textbf{เต โตสิตา จกฺขุมฺ}อตา,}{\textbf{พุทฺเธนา}\-\textbf{ทิจฺจพนฺธุนา}.} \\ \csromangathacontbat{\textbf{พฺรหฺมจริยมจรฺ}อิํสุ,}{วรปญฺญสฺส สนฺติเก.}}}

\prose{\textbf{เต โตสิตา จกฺขุมตา}ติ \textbf{เต}ติ โสฬส ปารายนิยา พฺราหฺมณา. \textbf{โตสิตา}ติ โตสิตา วิโตสิตา ปสาทิตา อาราธิตา อตฺตมนา กตาติ เต โตสิตา. \textbf{จกฺขุมตา}ติ ภควา ปญฺจหิ จกฺขูหิ จกฺขุมา, มํสจกฺขุนาปิ จกฺขุมา, ทิพฺพจกฺขุนาปิ จกฺขุมา, ปญฺญาจกฺขุนาปิ จกฺขุมา, พุทฺธจกฺขุนาปิ จกฺขุมา, สมนฺตจกฺขุนาปิ จกฺขุมา. กถํ ภควา มํสจกฺขุนาปิ จกฺขุมา ฯเปฯ. เอวํ ภควา สมนฺตจกฺขุนาปิ จกฺขุมาติ เต โตสิตา จกฺขุมตา.}

\prose{\textbf{พุทฺเธนา}\-\textbf{ทิจฺจพนฺธุนา}ติ พุทฺโธติ โย โส ภควา ฯเปฯ สจฺฉิกา ปญฺญตฺติ, ยทิทํ พุทฺโธติ. \textbf{อาทิจฺจพนฺธุนา}ติ \textbf{อาทิจฺโจ} วุจฺจติ สูริโย, โส โคตโม โคตฺเตน, ภควาปิ โคตโม โคตฺเตน, ภควา สูริยสฺส โคตฺตญาตโก โคตฺตพนฺธุ, ตสฺมา พุทฺโธ อาทิจฺจพนฺธูติ พุทฺเธนา\-ทิจฺจพนฺธุนา.}

\prose{\csromanfolio{200}\textbf{พฺรหฺมจริยมจริํสู}ติ พฺรหฺมจริยํ วุจฺจติ อสทฺธมฺม\-สมาปตฺติยา อารติ วิรติ ปฏิวิรติ เวรมณี วิรมณํ อกิริยา อกรณํ อนชฺฌาปตฺติ เวลาอนติกฺกโม เสตุฆาโต. อปิ จ นิปฺปริยาย\-วเสน พฺรหฺมจริยํ วุจฺจติ อริโย อฏฺฐงฺคิโก มคฺโค. เสยฺยถิทํ, สมฺมาทิฏฺฐิ สมฺมาสงฺกปฺโป สมฺมาวาจา สมฺมากมฺมนฺโต สมฺมาอาชีโว สมฺมาวายาโม สมฺมาสติ สมฺมาสมาธิ. พฺรหฺมจริยมจริํสูติ พฺรหฺมจริยํ จริํสุ อจริํสุ สมาทาย วตฺติํสูติ พฺรหฺมจริยม\-จริํสุ.}

\prose{\textbf{วรปญฺญสฺส สนฺติเก}ติ วรปญฺญสฺส อคฺคปญฺญสฺส เสฏฺฐปญฺญสฺส วิเส\-ฏฺฐ\-ปญฺญสฺส ปาโมกฺข\-ปญฺญสฺส อุตฺตมปญฺญสฺส ปวรปญฺญสฺส. \textbf{สนฺติเก}ติ สนฺติเก สามนฺตา อาสนฺเน อวิทูเร อุปกฏฺเฐติ วรปญฺญสฺส สนฺติเก. เตเนตํ วุจฺจติ\csromandash{}}

\setlength{\csromangathapagewidth}{0pt}
\setlength{\csromangathawakwidth}{0pt}
\csromangathameasuregroup{เต โตสิตา จกฺขุมตา, \\ พฺรหฺมจริยมจริํสุ, \\ \textbf{เอกเมกสฺส ปญฺหสฺ}ส, \\ \textbf{ตถา โย ปฏิปชฺชฺ}เอยฺย,}{\csromangathabat{เต โตสิตา จกฺขุมตา,}{พุทฺเธนาทิจฺจพนฺธุนา.} \\ \csromangathabat{พฺรหฺมจริยมจริํสุ,}{วรปญฺญสฺส สนฺติเกติ.} \\ \csromangathabat{\textbf{เอกเมกสฺส ปญฺหสฺ}ส,}{ยถา พุทฺเธน เทสิตํ.} \\ \csromangathabat{\textbf{ตถา โย ปฏิปชฺชฺ}เอยฺย,}{คจฺเฉ ปารํ อปารโต.}}
\csromangathasetleft{เต โตสิตา จกฺขุมตา, \\ พฺรหฺมจริยมจริํสุ, \\ \textbf{เอกเมกสฺส ปญฺหสฺ}ส, \\ \textbf{ตถา โย ปฏิปชฺชฺ}เอยฺย,}
\csromangathagroup{\csromangathastanza{\csromangathabat{เต โตสิตา จกฺขุมตา,}{พุทฺเธนา\-ทิจฺจพนฺธุนา.} \\ \csromangathabat{พฺรหฺมจริยม\-จริํสุ,}{วรปญฺญสฺส สนฺติเกติ.}}\csromangathastanza{\csromangathaitembat{100}{\textbf{เอกเมกสฺส ปญฺหสฺ}ส,}{ยถา พุทฺเธน เทสิตํ.} \\ \csromangathacontbat{\textbf{ตถา โย ปฏิปชฺชฺ}เอยฺย,}{คจฺเฉ ปารํ อปารโต.}}}

\prose{\textbf{เอกเมกสฺส ปญฺหสฺสา}ติ เอกเมกสฺส อชิตปญฺหสฺส เอกเมกสฺส ติสฺส\-เมตฺเตยฺยปญฺหสฺส ฯเปฯ เอกเมกสฺส ปิงฺคิย\-ปญฺหสฺสาติ เอกเมกสฺส ปญฺหสฺส.}

\prose{\textbf{ยถา พุทฺเธน เทสิตนฺ}ติ พุทฺโธติ โย โส ภควา สยมฺภู ฯเปฯ สจฺฉิกา ปญฺญตฺติ, ยทิทํ พุทฺโธติ. \textbf{ยถา พุทฺเธน เทสิตนฺ}ติ ยถา พุทฺเธน อาจิกฺขิตํ เทสิตํ ปญฺญปิตํ ปฏฺฐปิตํ วิวริตํ วิภชิตํ\csromansharedfootnote{478b891f8ebeec3d}{วิภตฺตํ (สฺยา)} อุตฺตานีกตํ ปกาสิตนฺติ ยถา พุทฺเธน เทสิตํ.}

\prose{\textbf{ตถา โย ปฏิปชฺเชยฺยา}ติ สมฺมาปฏิปทํ อนุโลม\-ปฏิปทํ อปจฺจนีกปฏิปทํ อนฺวตฺถ\-ปฏิปทํ ธมฺมานุธมฺม\-ปฏิปทํ ปฏิปชฺเชยฺยาติ ตถา โย ปฏิปชฺเชยฺย.}

\prose{\textbf{คจฺเฉ ปารํ อปารโต}ติ ปารํ วุจฺจติ อมตํ นิพฺพานํ ฯเปฯ นิโรโธ นิพฺพานํ. อปารํ วุจฺจนฺติ กิเลสา จ ขนฺธา จ อภิสงฺขารา จ. \textbf{คจฺเฉ ปารํ}}

\setlength{\csromangathapagewidth}{0pt}
\setlength{\csromangathawakwidth}{0pt}
\csromangathameasuregroup{}{เอกเมกสฺส ปญฺหสฺส, \\ ยถา พุทฺเธน เทสิตํ. \\ ตถา โย ปฏิปชฺเชยฺย,}
\csromangathameasurewak{เอกเมกสฺส ปญฺหสฺส, \\ ยถา พุทฺเธน เทสิตํ. \\ ตถา โย ปฏิปชฺเชยฺย,}
\setlength{\csromangathaleftcol}{0pt}
\csromangathagroup{\csromangathastanza{\csromangathawak{\csromanfolio{201}เอกเมกสฺส ปญฺหสฺส, \\ ยถา พุทฺเธน เทสิตํ. \\ ตถา โย ปฏิปชฺเชยฺย,}}}

\prose{\textbf{อปารโต}ติ อปารโต ปารํ คจฺเฉยฺย, ปารํ อธิคจฺเฉยฺย, ปารํ ผสฺเสยฺย, ปารํ สจฺฉิกเรยฺยาติ คจฺเฉ ปารํ อปารโต.}

\prose{เตเนตํ วุจฺจติ\csromandash{}}

\setlength{\csromangathapagewidth}{0pt}
\setlength{\csromangathawakwidth}{0pt}
\csromangathameasuregroup{\textbf{อปารา ปารํ คจฺเฉ}ยฺย, \\ \textbf{มคฺโค โส ปารํ คมฺ}อนาย,}{คจฺเฉ ปารํ อปารโตติ. \\ \csromangathabat{\textbf{อปารา ปารํ คจฺเฉ}ยฺย,}{ภาเวนฺโต มคฺคมุตฺตมํ.} \\ \csromangathabat{\textbf{มคฺโค โส ปารํ คมฺ}อนาย,}{ตสฺมา ปารายนํ อิติ.}}
\csromangathasetleft{\textbf{อปารา ปารํ คจฺเฉ}ยฺย, \\ \textbf{มคฺโค โส ปารํ คมฺ}อนาย,}
\csromangathagroup{\csromangathastanza{คจฺเฉ ปารํ อปารโตติ.}\csromangathastanza{\csromangathaitembat{101}{\textbf{อปารา ปารํ คจฺเฉ}ยฺย,}{ภาเวนฺโต มคฺคมุตฺตมํ.} \\ \csromangathacontbat{\textbf{มคฺโค โส ปารํ คมฺ}อนาย,}{ตสฺมา ปารายนํ อิติ.}}}

\prose{\textbf{อปารา ปารํ คจฺเฉยฺยา}ติ อปารํ วุจฺจนฺติ กิเลสา จ ขนฺธา จ อภิสงฺขารา จ. ปารํ วุจฺจติ อมตํ นิพฺพานํ ฯเปฯ ตณฺหกฺขโย วิราโค นิโรโธ นิพฺพานํ. \textbf{อปารา ปารํ คจฺเฉยฺยา}ติ อปารา ปารํ คจฺเฉยฺย ปารํ อธิคจฺเฉยฺย, ปารํ ผสฺเสยฺย, ปารํ สจฺฉิกเรยฺยาติ อปารา ปารํ คจฺเฉยฺย.}

\prose{\textbf{ภาเวนฺโต มคฺค}\-\textbf{มุตฺต}\-\textbf{มนฺ}ติ มคฺคมุตฺตมํ วุจฺจติ อริโย อฏฺฐงฺคิโก มคฺโค, เสยฺยถิทํ, สมฺมาทิฏฺฐิ ฯเปฯ สมฺมาสมาธิ. \textbf{มคฺค}\-\textbf{มุตฺต}\-\textbf{มนฺ}ติ มคฺคํ อคฺคํ เสฏฺฐํ วิเสฏฺฐํ ปาโมกฺขํ อุตฺตมํ ปวรํ. \textbf{ภาเวนฺโต}ติ ภาเวนฺโต อาเสวนฺโต พหุลีกโรนฺโตติ ภาเวนฺโต มคฺคมุตฺตมํ.}

\prose{\textbf{มคฺโค โส ปารํ คมนายา}ติ\csromandash{}}

\setlength{\csromangathapagewidth}{0pt}
\setlength{\csromangathawakwidth}{0pt}
\csromangathameasuregroup{มคฺโค ปนฺโถ ปโถ ปชฺโช\textsuperscript{1}, \\ นาวา อุตฺตรเสตุ จ,}{\csromangathabat{มคฺโค ปนฺโถ ปโถ ปชฺโช\textsuperscript{1},}{อญฺชสํ วฏุมา’ยนํ.} \\ \csromangathabat{นาวา อุตฺตรเสตุ จ,}{กุลฺโล จ ภิสิ สงฺกโม\textsuperscript{1}.}}
\csromangathasetleft{มคฺโค ปนฺโถ ปโถ ปชฺโช\textsuperscript{1}, \\ นาวา อุตฺตรเสตุ จ,}
\csromangathagroup{\csromangathastanza{\csromangathabat{มคฺโค ปนฺโถ ปโถ ปชฺโช\csromansharedfootnote{8c3f93ae1a367e32}{อทฺโธ (ก)},}{อญฺชสํ วฏุมา’ยนํ.} \\ \csromangathabat{นาวา อุตฺตรเสตุ จ,}{กุลฺโล จ ภิสิ สงฺกโม\csromansharedfootnote{8c3f93ae1a367e32}{อทฺโธ (ก)}.}}}

\prose{\textbf{ปารํ คมนายา}ติ ปารํ คมนาย ปารํ สมฺปาปนาย ปารํ สมนุปาปนาย ชรามรณสฺส ตรณายาติ มคฺโค โส ปารํ คมนาย.}

\prose{\textbf{ตสฺมา ปารายนํ อิตี}ติ \textbf{ตสฺมา}ติ ตสฺมา ตํการณา ตํเหตุ ตปฺปจฺจยา ตํนิทานา. ปารํ วุจฺจติ อมตํ นิพฺพานํ ฯเปฯ นิโรโธ นิพฺพานํ. อยนํ วุจฺจติ มคฺโค. \textbf{อิตี}ติ ปทสนฺธิ ฯเปฯ ปทา\-นุปุพฺพตา\-เปตํ อิตีติ ตสฺมา ปารายนํ อิติ. เตเนตํ วุจฺจติ\csromandash{}}

\prose{\textbf{อปารา ปารํ คจฺเฉยฺย, ภาเวนฺโต มคฺคมุตฺตมํ.}}

\setlength{\csromangathapagewidth}{0pt}
\setlength{\csromangathawakwidth}{0pt}
\csromangathameasuregroup{\textbf{มคฺโค โส ปารํ คมฺ}อนาย,}{\csromangathabat{\textbf{มคฺโค โส ปารํ คมฺ}อนาย,}{\textbf{ตสฺมา} ปารายนํ อิตีติ.}}
\csromangathasetleft{\textbf{มคฺโค โส ปารํ คมฺ}อนาย,}
\csromangathagroup{\csromangathastanza{\csromangathabat{\textbf{มคฺโค โส ปารํ คมฺ}อนาย,}{\textbf{ตสฺมา} ปารายนํ อิตีติ.}}}

\vspace{\tipitakaparskip}
\csromancenter{ปารายนตฺถุติคาถา\-นิทฺเทโส สตฺตรสโม.}
\csromanmatikaanchor{mtk.39}
\csromantocmarkref{subsection}{18. ปารายนานุคีติคาถานิทฺเทส}{mtk.39}
\bookmark[dest={mtk.39},level=2]{18. ปารายนานุคีติคาถานิทฺเทส}
\csromanhangingbegin
\hangingheaditem{102}{\csromanfolio{202}\textbf{ปารายนม}\-\textbf{นุคายิสฺสํ, (อิจฺจายสฺมา ปิงฺคิโย,)}}
\hangingbody{\textbf{ยถา’ทฺทกฺขิ ตถา’กฺขาสิ, วิมโล ภูริเมธโส.}}
\hangingbody{\textbf{นิกฺกาโม นิพฺพโน นาโค, กิสฺส เหตุ มุสา ภเณ.}}
\csromanhangingend

\prose{\textbf{ปารายนม}\-\textbf{นุคายิสฺสนฺ}ติ คีตม\-นุคายิสฺสํ, กถิตม\-นุ\-กถยิสฺสํ\csromansharedfootnote{98b4050d27795e90}{กถิตมนุคายิสฺสํ (สฺยา) เอวํ สพฺพปเทสุ อนุคายิสฺสนฺติ อาคตํ.}, ภณิตม\-นุ\-ภณิสฺสํ, ลปิตม\-นุ\-ลปิสฺสํ, ภาสิตมนุภาสิสฺสนฺติ ปารายนม\-นุคายิสฺสํ. \textbf{อิจฺจายสฺมา ปิงฺคิโย}ติ \textbf{อิจฺจา}ติ ปทสนฺธิ ฯเปฯ ปทา\-นุปุพฺพตา\-เปตํ อิจฺจาติ. \textbf{อายสฺมา}ติ ปิยวจนํ ครุวจนํ สคารวสปฺปติสฺสาธิวจนเมตํ อายสฺมาติ. \textbf{ปิงฺคิโย}ติ ตสฺส เถรสฺส นามํ สงฺขา สมญฺญา ปญฺญตฺติ โวหาโร นามํ นามกมฺมํ นามเธยฺยํ นิรุตฺติ พฺยญฺชนํ อภิลาโปติ อิจฺจายสฺมา ปิงฺคิโย.}

\prose{\textbf{ยถา’ทฺทกฺขิ ตถา’กฺขาสี}ติ ยถา อทฺทกฺขิ, ตถา อกฺขาสิ อาจิกฺขิ เทเสสิ ปญฺญเปสิ ปฏฺฐเปสิ วิวริ วิภชิ อุตฺตานีอกาสิ ปกาเสสิ. “สพฺเพ สงฺขารา อนิจฺจา”ติ ยถา อทฺทกฺขิ, ตถา อกฺขาสิ อาจิกฺขิ เทเสสิ ปญฺญเปสิ ปฏฺฐเปสิ วิวริ วิภชิ อุตฺตานีอกาสิ ปกาเสสิ. “สพฺเพ สงฺขารา ทุกฺขา”ติ ฯเปฯ. “สพฺเพ ธมฺมา อนตฺตา”ติ. “ยํ กิญฺจิ สมุทย\-ธมฺมํ, สพฺพํ ตํ นิโรธธมฺมนฺ”ติ ยถา อทฺทกฺขิ, ตถา อกฺขาสิ อาจิกฺขิ เทเสสิ ปญฺญเปสิ ปฏฺฐเปสิ วิวริ วิภชิ อุตฺตานีอกาสิ ปกาเสสีติ ยถา’ทฺทกฺขิ ตถา’กฺขาสิ.}

\prose{\textbf{วิมโล ภูริเมธโส}ติ \textbf{วิมโล}ติ ราโค มลํ, โทโส มลํ, โมโห มลํ, โกโธ, อุปนาโห ฯเปฯ สพฺพา\-กุสลา\-ภิสงฺขารา มลา. เต มลา พุทฺธสฺส ภควโต ปหีนา อุจฺฉินฺนมูลา ตาลาวตฺถุกตา อนภาวํกตา อายติํ อนุปฺปาทธมฺมา. อมโล พุทฺโธ วิมโล นิมฺมโล มลาปคโต มลวิปฺปหีโน มลวิมุตฺโต สพฺพ\-มล\-วีติวตฺโต, ภูริ วุจฺจติ ปถวี. ภควา ตาย\csromansharedfootnote{8e74af9172b625d1}{ภควา อิมาย (สฺยา)} ปถวิสมาย ปญฺญาย วิปุลาย วิตฺถตาย สมนฺนาคโต, เมธา วุจฺจติ ปญฺญา, ยา ปญฺญา ปชานนา ฯเปฯ อโมโห ธมฺมวิจโย สมฺมาทิฏฺฐิ. ภควา อิมาย เมธาย ปญฺญาย อุเปโต \csromanfolio{203}สมุเปโต อุปาคโต สมุปาคโต อุปปนฺโน สมุปปนฺโน สมนฺนาคโต, ตสฺมา พุทฺโธ สุเมธโสติ วิมโล ภูริเมธโส.}

\prose{\textbf{นิกฺกาโม นิพฺพโน นาโค}ติ กามาติ อุทฺทานโต เทฺว \textbf{กามา} วตฺถุกามา จ กิเลสกามา จ ฯเปฯ. อิเม วุจฺจนฺติ วตฺถุกามา ฯเปฯ. อิเม วุจฺจนฺติ กิเลสกามา. พุทฺธสฺส ภควโต วตฺถุกามา ปริญฺญาตา, กิเลสกามา ปหีนา. วตฺถุกามานํ ปริญฺญาตตฺตา กิเลสกามานํ ปหีนตฺตา ภควา น กาเม กาเมติ, น กาเม อิจฺฉติ, น กาเม ปตฺเถติ, น กาเม ปิเหติ, น กาเม อภิชปฺปติ. เย กาเม กาเมนฺติ, กาเม อิจฺฉนฺติ, กาเม ปตฺเถนฺติ, กาเม ปิเหนฺติ, กาเม อภิชปฺปนฺติ, เต กามกามิโน ราคราคิโน สญฺญสญฺญิโน. ภควา น กาเม กาเมติ, น กาเม อิจฺฉติ, น กาเม ปตฺเถติ, น กาเม ปิเหติ, น กาเม อภิชปฺปติ. ตสฺมา พุทฺโธ อกาโม นิกฺกาโม จตฺตกาโม วนฺตกาโม มุตฺตกาโม ปหีนกาโม ปฏินิสฺสฏฺฐ\-กาโม วีตราโค วิคตราโค จตฺตราโค วนฺตราโค มุตฺตราโค ปหีนราโค ปฏินิสฺสฏฺฐ\-ราโค นิจฺฉาโต นิพฺพุโต สีติภูโต สุขปฺปฏิสํเวที พฺรหฺมภูเตน อตฺตนา วิหรตีติ นิกฺกาโม.}

\prose{\textbf{นิพฺพโน}ติ ราโค วนํ, โทโส วนํ, โมโห วนํ, โกโธ วนํ, อุปนาโห วนํ ฯเปฯ สพฺพา\-กุสลา\-ภิสงฺขารา วนา. เต วนา พุทฺธสฺส ภควโต ปหีนา อุจฺฉินฺนมูลา ตาลาวตฺถุกตา อนภาวํกตา อายติํ อนุปฺปาทธมฺมา. ตสฺมา พุทฺโธ อวโน วิวโน นิพฺพโน วนาปคโต วนวิปฺปหีโน วนวิมุตฺโต สพฺพวนวีติวตฺโตติ นิพฺพโน. \textbf{นาโค}ติ นาโค, ภควา อาคุํ น กโรตีติ นาโค. น คจฺฉตีติ นาโค. น อาคจฺฉตีติ นาโค ฯเปฯ. เอวํ ภควา น อาคจฺฉตีติ นาโคติ นิกฺกาโม นิพฺพโน นาโค.}

\prose{\textbf{กิสฺส เหตุ มุสา ภเณ}ติ \textbf{กิสฺส เหตู}ติ กิสฺส เหตุ กิํเหตุ กิํการณา กิํนิทานา กิํปจฺจยาติ กิสฺส เหตุ. \textbf{มุสา ภเณ}ติ มุสา ภเณยฺย กเถยฺย ทีเปยฺย โวหเรยฺย. \textbf{มุสา ภเณ}ติ โมสวชฺชํ ภเณยฺย. มุสาวาทํ ภเณยฺย. อนริยวาทํ ภเณยฺย. อิเธกจฺโจ สภาคโต\csromansharedfootnote{cd53de9580c27b7f}{สภคฺคโต (สฺยา)} วา ปริสาคโต\csromansharedfootnote{dfcec68090cc8967}{ปริสคฺคโต (สฺยา)} วา ญาติมชฺฌคโต วา ปูคมชฺฌคโต วา ราชกุล\-มชฺฌคโต วา อภินีโต สกฺขิปุฏฺโฐ “เอหมฺโภ\csromansharedfootnote{c583ebb7f0aaa586}{เอหิ โภ (สฺยา) ปสฺส ม 3. 95 ปิฏฺเฐ.} ปุริส ยํ ชานาสิ, ตํ วเทหี”ติ โส อชานํ วา อาห \csromanfolio{204}“ชานามี”ติ, ชานํ วา อาห “น ชานามี”ติ. อปสฺสํ วา อาห “ปสฺสามี”ติ, ปสฺสํ วา อาห “น ปสฺสามี”ติ อิติ อตฺตเหตุ วา ปรเหตุ วา อามิส\-กิญฺจิกฺข\-เหตุ วา สมฺปชานมุสา ภาสติ. อิทํ วุจฺจติ โมสวชฺชํ.}

\proseitem{18}{อปิ จ ตีหากาเรหิ มุสาวาโท โหติ, ปุพฺเพวสฺส โหติ “มุสา ภณิสฺสนฺ”ติ. ภณนฺตสฺส โหติ “มุสา ภณามี”ติ. ภณิตสฺส โหติ “มุสา มยา ภณิตนฺ”ติ. อิเมหิ ตีหากาเรหิ มุสาวาโท โหติ.  อปิ จ จตูหากาเรหิ มุสาวาโท โหติ ปุพฺเพวสฺส โหติ “มุสา ภณิสฺสนฺ”ติ. ภณนฺตสฺส โหติ “มุสา ภณามี”ติ. ภณิตสฺส โหติ “มุสา มยา ภณิตนฺ”ติ วินิธาย ทิฏฺฐิํ. อิเมหิ จตูหากาเรหิ มุสาวาโท โหติ.  อปิ จ ปญฺจหากาเรหิ. ฉหากาเรหิ. สตฺตหากาเรหิ. อฏฺฐหากาเรหิ มุสาวาโท โหติ, ปุพฺเพวสฺส โหติ “มุสา ภณิสฺสนฺ”ติ. ภณนฺตสฺส โหติ “มุสา ภณามี”ติ. ภณิตสฺส โหติ “มุสา มยา ภณิตนฺ”ติ วินิธาย ทิฏฺฐิํ, วินิธาย ขนฺติํ, วินิธาย รุจิํ, วินิธาย สญฺญํ, วินิธาย ภาวํ. อิเมหิ อฏฺฐหากาเรหิ มุสาวาโท โหติ. โมสวชฺชํ กิสฺส เหตุ มุสา ภเณยฺย กถายฺย ทีเปยฺย โวหเรยฺยาติ กิสฺส เหตุ มุสา ภเณ. เตนาห เถโร ปิงฺคิโย\csromandash{}}

\prose{ปารายนม\-นุคายิสฺสํ, (อิจฺจายสฺมา ปิงฺคิโย,)}

\setlength{\csromangathapagewidth}{0pt}
\setlength{\csromangathawakwidth}{0pt}
\csromangathameasuregroup{ยถา’ทฺทกฺขิ ตถา’กฺขาสิ, \\ นิกฺกาโม นิพฺพโน นาโค, \\ \textbf{ปหีน}\textbf{มล}\textbf{โมหสฺ}ส, \\ \textbf{หนฺทาหํ กิตฺตยิสฺสฺ}อามิ,}{\csromangathabat{ยถา’ทฺทกฺขิ ตถา’กฺขาสิ,}{วิมโล ภูริเมธโส.} \\ \csromangathabat{นิกฺกาโม นิพฺพโน นาโค,}{กิสฺส เหตุ มุสา ภเณติ.} \\ \csromangathabat{\textbf{ปหีน}\textbf{มล}\textbf{โมหสฺ}ส,}{มานมกฺขปฺปหายิโน.} \\ \csromangathabat{\textbf{หนฺทาหํ กิตฺตยิสฺสฺ}อามิ,}{คิรํ วณฺณูปสํหิตํ.}}
\csromangathasetleft{ยถา’ทฺทกฺขิ ตถา’กฺขาสิ, \\ นิกฺกาโม นิพฺพโน นาโค, \\ \textbf{ปหีน}\textbf{มล}\textbf{โมหสฺ}ส, \\ \textbf{หนฺทาหํ กิตฺตยิสฺสฺ}อามิ,}
\csromangathagroup{\csromangathastanza{\csromangathabat{ยถา’ทฺทกฺขิ ตถา’กฺขาสิ,}{วิมโล ภูริเมธโส.} \\ \csromangathabat{นิกฺกาโม นิพฺพโน นาโค,}{กิสฺส เหตุ มุสา ภเณติ.}}\csromangathastanza{\csromangathaitembat{103}{\textbf{ปหีน}\-\textbf{มล}\-\textbf{โมหสฺ}ส,}{มาน\-มกฺข\-ปฺปหายิโน.} \\ \csromangathacontbat{\textbf{หนฺทาหํ กิตฺตยิสฺสฺ}อามิ,}{คิรํ วณฺณู\-ปสํหิตํ.}}}

\prose{\textbf{ปหีน}\-\textbf{มล}\-\textbf{โมหสฺสา}ติ \textbf{มลนฺ}ติ ราโค มลํ, โทโส มลํ, โมโห มลํ, มาโน มลํ, ทิฏฺฐิ มลํ, กิเลโส มลํ, สพฺพ\-ทุจฺจริตํ มลํ, สพฺพ\-ภวคามิ\-กมฺมํ มลํ.}

\prose{\textbf{โมโห}ติ ยํ ทุกฺเข อญฺญาณํ ฯเปฯ อวิชฺชาลงฺคี โมโห อกุสลมูลํ. อยํ วุจฺจติ โมโห. มลญฺจ โมโห จ พุทฺธสฺส ภควโต ปหีนา อุจฺฉินฺนมูลา ตาลาวตฺถุกตา อนภาวํกตา \csromanfolio{205}อายติํ อนุปฺปาทธมฺมา. ตสฺมา พุทฺโธ ปหีนมลโมโหติ ปหีน\-มล\-โมหสฺส.}

\prose{\textbf{มานมกฺขปฺปหายิโน}ติ \textbf{มาโน}ติ เอกวิเธน มาโน, ยา จิตฺตสฺส อุนฺนติ\csromansharedfootnote{ff3c29e799fcfacf}{อุณฺณติ (สฺยา, ก)}. ทุวิเธน มาโน, อตฺตุกฺกํสน\-มาโน, ปรวมฺภน\-มาโน. ติวิเธน มาโน, “เสยฺโยหมสฺมี”ติ มาโน, “สทิโสหมสฺมี”ติ มาโน, “หีโนหมสฺมี”ติ มาโน. จตุพฺพิเธน มาโน, ลาเภน มานํ ชเนติ, ยเสน มานํ ชเนติ, ปสํสาย มานํ ชเนติ, สุเขน มานํ ชเนติ. ปญฺจวิเธน มาโน, “ลาภิมฺหิ มนาปิกานํ รูปานนฺ”ติ มานํ ชเนติ, “ลาภิมฺหิ มนาปิกานํ สทฺทานํ. คนฺธานํ. รสานํ. โผฏฺฐพฺพานนฺ”ติ มานํ ชเนติ. ฉพฺพิเธน มาโน, จกฺขุ\-สมฺปทาย มานํ ชเนติ, โสตสมฺปทาย. ฆานสมฺปทาย. ชิวฺหาสมฺปทาย. กายสมฺปทาย. มโนสมฺปทาย มานํ ชเนติ. สตฺตวิเธน มาโน, มาโน อติมาโน มานาติมาโน โอมาโน อวมาโน อสฺมิมาโน มิจฺฉามาโน. อฏฺฐวิเธน มาโน, ลาเภน มานํ ชเนติ, อลาเภน โอมานํ ชเนติ, ยเสน มานํ ชเนติ, อยเสน โอมานํ ชเนติ, ปสํสาย มานํ ชเนติ, นินฺทาย โอมานํ ชเนติ, สุเขน มานํ ชเนติ, ทุกฺเขน โอมานํ ชเนติ. นววิเธน มาโน, เสยฺยสฺส “เสยฺโยหมสฺมี”ติ มาโน, เสยฺยสฺส “สทิโสหมสฺมี”ติ มาโน, เสยฺยสฺส “หีโนหมสฺมี”ติ มาโน. สทิสสฺส “เสยฺโยหมสฺมี”ติ มาโน, สทิสสฺส “สทิโสหมสฺมี”ติ มาโน, สทิสสฺส “หีโนหมสฺมี”ติ มาโน. หีนสฺส “เสยฺโยหมสฺมี”ติ มาโน. หีนสฺส “สทิโสหมสฺมี”ติ มาโน, หีนสฺส “หีโนหมสฺมี”ติ มาโน. ทสวิเธน มาโน, อิเธกจฺโจ มานํ ชเนติ ชาติยา วา โคตฺเตน วา โกลปุตฺติเยน วา วณฺณ\-โปกฺขร\-ตาย วา ธเนน วา อชฺเฌเนน วา กมฺมายตเนน วา สิปฺปายตเนน วา วิชฺชาฏฺฐาเนน\csromansharedfootnote{156d35656bcbc5a0}{วิชฺชาฐาเนน (ก)} วา สุเตน วา ปฏิภาเนน วา อญฺญตร\-ญฺญตเรน วา วตฺถุนา. โย เอวรูโป มาโน มญฺญนา มญฺญิตตฺตํ อุนฺนติ อุนฺนโม ธโช สมฺปคฺคาโห เกตุกมฺยตา จิตฺตสฺส. อยํ วุจฺจติ มาโน.}

\prose{\textbf{มกฺโข}ติ โย มกฺโข มกฺขายนา มกฺขายิตตฺตํ นิฏฺฐุริยํ นิฏฺฐุริย\-กมฺมํ\csromansharedfootnote{99d8356c4055768d}{นิตฺถุริยกมฺมํ (ก) ปสฺส อภิ 2. 371 ปิฏฺเฐ.}. อยํ วุจฺจติ มกฺโข. พุทฺธสฺส ภควโต มาโน จ มกฺโข จ ปหีนา \csromanfolio{206}อุจฺฉินฺนมูลา ตาลาวตฺถุกตา อนภาวํกตา อายติํ อนุปฺปาทธมฺมา. ตสฺมา พุทฺโธ มาน\-มกฺข\-ปฺป\-หายีติ มาน\-มกฺข\-ปฺปหายิโน.}

\proseitem{18}{\textbf{หนฺทาหํ กิตฺตยิสฺสามิ, คิรํ วณฺณูปสญฺหิตนฺ}ติ \textbf{หนฺทาหนฺ}ติ ปทสนฺธิ ปทสํสคฺโค ปทปาริปูรี อกฺขร\-สมวาโย พฺยญฺชน\-สิลิฏฺฐตา ปทา\-นุปุพฺพตา\-เปตํ หนฺทาหนฺติ. \textbf{กิตฺตยิสฺสามิ คิรํ วณฺณูปสํหิตนฺ}ติ วณฺเณน อุเปตํ สมุเปตํ อุปาคตํ สมุปาคตํ อุปปนฺนํ สมุปปนฺนํ สมนฺนาคตํ วาจํ คิรํ พฺยปฺปถํ อุทีรณํ\csromansharedfootnote{550ef23873fe19e3}{โอทีรณํ (สฺยา)} กิตฺตยิสฺสามิ เทเสสฺสามิ ปญฺญเปสฺสามิ ปฏฺฐเปสฺสามิ วิวริสฺสามิ วิภชิสฺสามิ อุตฺตานีกริสฺสามิ ปกาเสสฺสามีติ หนฺทาหํ กิตฺตยิสฺสามิ, คิรํ วณฺณู\-ปสํหิตํ. เตนาห เถโร ปิงฺคิโย\csromandash{}}

\setlength{\csromangathapagewidth}{0pt}
\setlength{\csromangathawakwidth}{0pt}
\csromangathameasuregroup{ปหีนมลโมหสฺส,}{\csromangathabat{ปหีนมลโมหสฺส,}{มานมกฺขปฺปหายิโน.}}
\csromangathasetleft{ปหีนมลโมหสฺส,}
\csromangathagroup{\csromangathastanza{\csromangathabat{ปหีน\-มล\-โมหสฺส,}{มาน\-มกฺข\-ปฺปหายิโน.}}}

\prose{หนฺทาหํ กิตฺตยิสฺสามิ, คิรํ วณฺณูปสญฺหิตนฺติ.}

\setlength{\csromangathapagewidth}{0pt}
\setlength{\csromangathawakwidth}{0pt}
\csromangathameasuregroup{}{\textbf{ตโมนุโท พุทฺโธ สมนฺตจกฺขุ,} \\ \textbf{โลกนฺตคู สพฺพภวาติวตฺโต.} \\ \textbf{อนาสโว สพฺพทุกฺข}\textbf{ปฺปหีโน,} \\ \textbf{สจฺจวฺหโย พฺรหฺเม อุปาสิโต เม.}}
\csromangathameasurewak{\textbf{ตโมนุโท พุทฺโธ สมนฺตจกฺขุ,} \\ \textbf{โลกนฺตคู สพฺพภวาติวตฺโต.} \\ \textbf{อนาสโว สพฺพทุกฺข}\textbf{ปฺปหีโน,} \\ \textbf{สจฺจวฺหโย พฺรหฺเม อุปาสิโต เม.}}
\setlength{\csromangathaleftcol}{0pt}
\csromangathagroup{\csromangathastanza{\csromangathawak{\csromangathaitemline{104}{\textbf{ตโมนุโท พุทฺโธ สมนฺตจกฺขุ,}} \\ \csromangathacontline{\textbf{โลกนฺตคู สพฺพภวาติวตฺโต.}} \\ \csromangathacontline{\textbf{อนาสโว สพฺพทุกฺข}\-\textbf{ปฺปหีโน,}} \\ \csromangathacontline{\textbf{สจฺจวฺหโย พฺรหฺเม อุปาสิโต เม.}}}}}

\prose{\textbf{ตโมนุโท พุทฺโธ สมนฺต}\-\textbf{จกฺขู}ติ \textbf{ตโมนุโท}ติ ราคตมํ โทสตมํ โมหตมํ มานตมํ ทิฏฺฐิตมํ กิเลสตมํ ทุจฺจริตตมํ อนฺธกรณํ อญฺญาณกรณํ ปญฺญา\-นิโรธิกํ วิฆาต\-ปกฺขิกํ อนิพฺพานสํวตฺตนิกํ นุทิ ปนุทิ ปชหิ วิโนเทสิ พฺยนฺตีอกาสิ อนภาวํ คเมสิ. พุทฺโธติ โย โส ภควา ฯเปฯ สจฺฉิกา ปญฺญตฺติ, ยทิทํ พุทฺโธติ. สมนฺตจกฺขุ วุจฺจติ สพฺพญฺญุต\-ญาณํ ฯเปฯ ตถาคโต เตน สมนฺต\-จกฺขูติ ตโมนุโท พุทฺโธ สมนฺตจกฺขุ.}

\prose{\textbf{โลกนฺตคู สพฺพภวาติวตฺโต}ติ โลโกติ เอโก \textbf{โลโก}, ภวโลโก. เทฺว โลกา, ภวโลโก จ สมฺภวโลโก จ. สมฺปตฺติ\-ภว\-โลโก จ สมฺปตฺติ\-สมฺภว\-โลโก จ. วิปตฺติ\-ภว\-โลโก จ วิปตฺติ\-สมฺภว\-โลโก จ\csromansharedfootnote{d3d437af9c36886d}{เทฺว โลกา สมฺปตฺติ จ ภวโลโก วิปตฺติ จ ภวโลโก (สฺยา)}. ตโย โลกา. ติสฺโส เวทนา. จตฺตาโร โลกา, จตฺตาโร อาหารา. ปญฺจ โลกา, ปญฺจุ\-ปาทาน\-กฺขนฺธา. ฉ โลกา, ฉ อชฺฌตฺติกานิ อายตนานิ. สตฺต โลกา, \csromanfolio{207}สตฺต\-วิญฺญาณ\-ฏฺฐิติโย. อฏฺฐ โลกา, อฏฺฐ โลกธมฺมา. นว โลกา, นว สตฺตาวาสา. ทส โลกา, ทส อายตนานิ. ทฺวาทส โลกา, ทฺวาทสา\-ยตนานิ. อฏฺฐารส โลกา, อฏฺฐารส ธาตุโย. \textbf{โลกนฺตคู}ติ ภควา โลกสฺส อนฺตคโต อนฺตปฺปตฺโต โกฏิคโต โกฏิปฺปตฺโต ฯเปฯ นิพฺพานคโต นิพฺพานปฺปตฺโต. โส วุตฺถวาโส จิณฺณจรโณ ฯเปฯ ชาติ\-มรณ\-สํสาโร, นตฺถิ ตสฺส ปุนพฺภโวติ \textbf{โลกนฺตคู.}}

\prose{\textbf{สพฺพภวาติวตฺโต}ติ \textbf{ภวา}ติ เทฺว ภวา กมฺมภโว จ ปฏิสนฺธิโก จ ปุนพฺภโว. กตโม กมฺมภโว. ปุญฺญา\-ภิสงฺขาโร อปุญฺญา\-ภิสงฺขาโร อาเนญฺชา\-ภิสงฺขาโร, อยํ กมฺมภโว. กตโม ปฏิสนฺธิโก ปุนพฺภโว. ปฏิสนฺธิกา รูปา เวทนา สญฺญา สงฺขารา วิญฺญาณํ, อยํ ปฏิสนฺธิโก ปุนพฺภโว. ภควา กมฺมภวญฺจ ปฏิสนฺธิ\-กญฺจ ปุนพฺภวํ อติวตฺโต\csromansharedfootnote{cdd5d9dbd0850257}{อุปาติวตฺโต (ก)} อติกฺกนฺโต วีติวตฺโตติ โลกนฺตคู สพฺพภวาติวตฺโต.}

\prose{\textbf{อนาสโว สพฺพทุกฺข}\-\textbf{ปฺปหีโน}ติ \textbf{อนาสโว}ติ จตฺตาโร อาสวา กามาสโว ภวาสโว ทิฏฺฐาสโว อวิชฺชาสโว. เต อาสวา พุทฺธสฺส ภควโต ปหีนา อุจฺฉินฺนมูลา ตาลาวตฺถุกตา อนภาวํกตา อายติํ อนุปฺปาทธมฺมา. ตสฺมา พุทฺโธ อนาสโว. \textbf{สพฺพทุกฺข}\-\textbf{ปฺปหีโน}ติ สพฺพํ ตสฺส ปฏิสนฺธิกํ ชาติทุกฺขํ ชราทุกฺขํ พฺยาธิทุกฺขํ มรณทุกฺขํ โสกปริเทวทุกฺขโทมนสฺสุปายาส\-ทุกฺขํ ฯเปฯ ทิฏฺฐิพฺยสน\-ทุกฺขํ ปหีนํ สมุจฺฉินฺนํ วูปสนฺตํ ปฏิปฺปสฺสทฺธํ อภพฺพุ\-ปฺปตฺติกํ ญาณคฺคินา ทฑฺฒํ. ตสฺมา พุทฺโธ สพฺพทุกฺข\-ปฺปหีโนติ อนาสโว สพฺพทุกฺข\-ปฺปหีโน.}

\prose{\textbf{สจฺจวฺหโย พฺรหฺเม อุปาสิโต เม}ติ สจฺจวฺหโยติ \textbf{สจฺจวฺหโย} สทิสนาโม สทิสวฺหโย สจฺจ\-สทิส\-วฺหโย. วิปสฺสี ภควา สิขี ภควา เวสฺสภู ภควา กกุสนฺโธ ภควา โกณาคมโน ภควา กสฺสโป ภควา. เต พุทฺธา ภควนฺโต สทิสนามา สทิสวฺหยา. ภควาปิ สกฺยมุนิ เตสํ พุทฺธานํ ภควนฺตานํ สทิสนาโม สทิสวฺหโยติ ตสฺมา พุทฺโธ สจฺจวฺหโย.}

\prose{\textbf{พฺรหฺเม อุปาสิโต เม}ติ โส มยา ภควา อาสิโต อุปาสิโต ปยิรุปาสิโต ปริปุจฺฉิโต ปริปญฺหิโตติ สจฺจวฺหโย พฺรหฺเม อุปาสิโต เม. เตนาห เถโร ปิงฺคิโย\csromandash{}}

\setlength{\csromangathapagewidth}{0pt}
\setlength{\csromangathawakwidth}{0pt}
\csromangathameasuregroup{}{ตโมนุโท พุทฺโธ สมนฺตจกฺขุ, \\ โลกนฺตคู สพฺพภวาติวตฺโต. \\ อนาสโว สพฺพทุกฺขปฺปหีโน, \\ สจฺจวฺหโย พฺรหฺเม อุปาสิโต เมติ. \\ \textbf{ทิโช ยถา กุพฺพนกํ ปหาย,} \\ \textbf{พหุปฺผลํ กานนมา}\textbf{วเสยฺย.} \\ \textbf{เอวมหํ อปฺปทสฺเส ปหาย,} \\ \textbf{มโหทธิํ หํโสริว อชฺฌปตฺโต}\textsuperscript{1}\textbf{.}}
\csromangathameasurewak{ตโมนุโท พุทฺโธ สมนฺตจกฺขุ, \\ โลกนฺตคู สพฺพภวาติวตฺโต. \\ อนาสโว สพฺพทุกฺขปฺปหีโน, \\ สจฺจวฺหโย พฺรหฺเม อุปาสิโต เมติ. \\ \textbf{ทิโช ยถา กุพฺพนกํ ปหาย,} \\ \textbf{พหุปฺผลํ กานนมา}\textbf{วเสยฺย.} \\ \textbf{เอวมหํ อปฺปทสฺเส ปหาย,} \\ \textbf{มโหทธิํ หํโสริว อชฺฌปตฺโต}\textsuperscript{1}\textbf{.}}
\setlength{\csromangathaleftcol}{0pt}
\csromangathagroup{\csromangathastanza{\csromangathawak{\csromanfolio{208}ตโมนุโท พุทฺโธ สมนฺตจกฺขุ, \\ โลกนฺตคู สพฺพภวาติวตฺโต. \\ อนาสโว สพฺพทุกฺข\-ปฺปหีโน, \\ สจฺจวฺหโย พฺรหฺเม อุปาสิโต เมติ.}}\csromangathastanza{\csromangathawak{\csromangathaitemline{105}{\textbf{ทิโช ยถา กุพฺพนกํ ปหาย,}} \\ \csromangathacontline{\textbf{พหุปฺผลํ กานนมา}\-\textbf{วเสยฺย.}} \\ \csromangathacontline{\textbf{เอวมหํ อปฺปทสฺเส ปหาย,}} \\ \csromangathacontline{\textbf{มโหทธิํ หํโสริว อชฺฌปตฺโต}\csromansharedfootnote{b946337fc6ab7312}{อชฺช ปตฺโต (ก)}\textbf{.}}}}}

\prose{\textbf{ทิโช ยถา กุพฺพนกํ ปหาย, พหุปฺผลํ กานนมา}\-\textbf{วเสยฺยา}ติ ทิโช วุจฺจติ ปกฺขี. กิํการณา ทิโช วุจฺจติ ปกฺขี, ทฺวิกฺขตฺตุํ ชายตีติ ทิโช มาตุกุจฺฉิมฺหา จ อณฺฑโกสมฺหา จ. ตํการณา ทิโช วุจฺจติ ปกฺขีติ ทิโช. \textbf{ยถา กุพฺพนกํ ปหายา}ติ ยถา ทิโช กุพฺพนกํ ปริตฺตวนกํ อปฺปผลํ อปฺปภกฺขํ อปฺโปทกํ ปหาย ชหิตฺวา อติกฺกมิตฺวา สมติกฺกมิตฺวา วีติวตฺเตตฺวา อญฺญํ พหุปฺผลํ พหุภกฺขํ พหูทกํ\csromansharedfootnote{79a6aa94a145bfa7}{พหุรุกฺขํ (สฺยา)} มหนฺตํ กานนํ วนสณฺฑํ อธิคจฺเฉยฺย วินฺเทยฺย ปฏิลเภยฺย ตสฺมิํ จ วนสณฺเฑ วาสํ กปฺเปยฺยาติ ทิโช ยถา กุพฺพนกํ ปหาย, พหุปฺผลํ กานนํ อาวเสยฺย.}

\prose{\textbf{เอวมหํ อปฺปทสฺเส ปหาย, มโหทธิํ หํโสริว อชฺฌปตฺโต}ติ \textbf{เอวนฺ}ติ โอปมฺม\-สมฺปฏิปาทนํ. \textbf{อปฺปทสฺเส ปหายา}ติ โย จ พาวรี พฺราหฺมโณ เย จญฺเญ ตสฺส อาจริยา พุทฺธํ ภควนฺตํ อุปาทาย อปฺปทสฺสา ปริตฺตทสฺสา โถกทสฺสา โอมกทสฺสา ลามกทสฺสา ฉตุกฺกทสฺสา\csromansharedfootnote{f5e378c16b9df6ea}{ชตุกฺกทสฺสา (สฺยา), ชตุกทสฺสา (สีฏฺฐ)} วา. เต อปฺปทสฺเส ปริตฺตทสฺเส โถกทสฺเส โอมกทสฺเส ลามกทสฺเส ฉตุกฺกทสฺเส ปหาย ปชหิตฺวา อติกฺกมิตฺวา สมติกฺกมิตฺวา วีติวตฺเตตฺวา พุทฺธํ ภควนฺตํ อปฺปมาณทสฺสํ อคฺคทสฺสํ เสฏฺฐทสฺสํ วิเสฏฺฐทสฺสํ ปาโมกฺขทสฺสํ อุตฺตมทสฺสํ ปวรทสฺสํ อสมํ อสมสมํ อปฺปฏิสมํ อปฺปฏิภาคํ อปฺปฏิปุคฺคลํ เทวาติเทวํ นราสภํ ปุริสสีหํ ปุริสนาคํ ปุริสาชญฺญํ ปุริสนิสภํ ปุริส\-โธรยฺหํ ทสพลธาริํ\csromansharedfootnote{a9918d0976626c17}{ทสพลํ ตาทิํ (สฺยา)} อธิคจฺฉิํ วินฺทิํ ปฏิลภิํ. ยถา จ หํโส \csromanfolio{209}มหนฺตํ มานสกํ\csromansharedfootnote{346c895ad903749f}{มานุสกตํ (สฺยา)} วา สรํ อโนตตฺตํ วา ทหํ มหาสมุทฺทํ วา อกฺโขภํ อมิโตทกํ ชลราสิํ อธิคจฺเฉยฺย วินฺเทยฺย ปฏิลเภยฺย, เอวเมว พุทฺธํ ภควนฺตํ อกฺโขภํ อมิตเตชํ ปภินฺนญาณํ วิวฏจกฺขุํ ปญฺญาปเภท\-กุสลํ อธิคต\-ปฏิสมฺภิทํ จตุเวสารชฺชปฺปตฺตํ สุทฺธา\-ธิมุตฺตํ เสตปจฺจตฺตํ อทฺวยภาณิํ ตาทิํ ตถาปฏิญฺญํ อปริตฺตํ มหนฺตํ คมฺภีรํ อปฺปเมยฺยํ ทุปฺปริโยคาหํ ปหูตรตนํ สาครสมํ ฉฬงฺคุ\-เปกฺขาย สมนฺนาคตํ อตุลํ วิปุลํ อปฺปเมยฺยํ, ตํ ตาทิสํ ปวทตํ มคฺควาทินํ\csromansharedfootnote{9bb47eba2c2b7879}{ปวรมคฺควาทินํ (ก)} เมรุมิว นคานํ ครุฬมิว ทิชานํ สีหมิว มิคานํ อุทธิมิว อณฺณวานํ อธิคจฺฉิํ, ตํ สตฺถารํ ชินปวรํ มเหสินฺติ เอวมหํ อปฺปทสฺเส ปหาย, มโหทธิํ หํโสริว อชฺฌปตฺโต. เตนาห เถโร ปิงฺคิโย\csromandash{}}

\setlength{\csromangathapagewidth}{0pt}
\setlength{\csromangathawakwidth}{0pt}
\csromangathameasuregroup{}{ทิโช ยถา กุพฺพนกํ ปหาย, \\ พหุปฺผลํ กานนมาวเสยฺย. \\ เอวมหํ อปฺปทสฺเส ปหาย, \\ มโหทธิํ หํโสริว อชฺฌปตฺโตติ.}
\csromangathameasurewak{ทิโช ยถา กุพฺพนกํ ปหาย, \\ พหุปฺผลํ กานนมาวเสยฺย. \\ เอวมหํ อปฺปทสฺเส ปหาย, \\ มโหทธิํ หํโสริว อชฺฌปตฺโตติ.}
\setlength{\csromangathaleftcol}{0pt}
\csromangathagroup{\csromangathastanza{\csromangathawak{ทิโช ยถา กุพฺพนกํ ปหาย, \\ พหุปฺผลํ กานนมา\-วเสยฺย. \\ เอวมหํ อปฺปทสฺเส ปหาย, \\ มโหทธิํ หํโสริว อชฺฌปตฺโตติ.}}}

\csromanhangingbegin
\hangingheaditem{106}{\textbf{เย เม ปุพฺเพ วิยากํสุ,}}
\hangingbody{\textbf{หุรํ โคตมสาสนา, “อิจฺจาสิ อิติ ภวิสฺสติ.}”}
\hangingbody{\textbf{สพฺพํ ตํ อิติหีติหํ, สพฺพํ ตํ ตกฺกวฑฺฒนํ. (5)}}
\csromanhangingend

\prose{\textbf{เย เม ปุพฺเพ วิยากํสู}ติ \textbf{เย}ติ โย จ พาวรี พฺราหฺมโณ เย จญฺเญ ตสฺส อาจริยา, เต สกํ ทิฏฺฐิํ สกํ ขนฺติํ สกํ รุจิํ สกํ ลทฺธิํ สกํ อชฺฌาสยํ สกํ อธิปฺปายํ พฺยากํสุ อาจิกฺขิํสุ เทสยิํสุ ปญฺญปิํสุ ปฏฺฐปิํสุ วิวริํสุ วิภชิํสุ อุตฺตานีอกํสุ ปกาเสสุนฺติ เย เม ปุพฺเพ วิยากํสุ.}

\prose{\textbf{หุรํ โคตมสาสนา}ติ หุรํ โคตมสาสนา ปรํ โคตมสาสนา ปุเร โคตมสาสนา ปฐมตรํ โคตมสาสนา พุทฺธสาสนา ชินสาสนา ตถาคต\-สาสนา\csromansharedfootnote{46b7ae5aeff15680}{ตถาคตสาสนา เทวสาสนา (สฺยา, ก)} อรหนฺต\-สาสนาติ หุรํ โคตมสาสนา.}

\prose{\textbf{อิจฺจาสิ อิติ ภวิสฺสตี}ติ เอวํ กิร อาสิ, เอวํ กิร ภวิสฺสตีติ อิจฺจาสิ อิติ ภวิสฺสติ.}

\prose{\csromanfolio{210}\textbf{สพฺพํ ตํ อิติหีติหนฺ}ติ สพฺพํ ตํ อิติหีติหํ อิติกิราย ปรมฺปราย ปิฏก\-สมฺปทาย ตกฺกเหตุ นยเหตุ อาการ\-ปริวิตกฺเกน ทิฏฺฐิ\-นิชฺฌาน\-กฺขนฺติยา น สามํ สยม\-ภิญฺญาตํ น อตฺตปจฺจกฺขํ ธมฺมํ ยํ กถยิํสูติ สพฺพํ ตํ อิติหีติหํ.}

\prose{\textbf{สพฺพํ ตํ ตกฺก}\-\textbf{วฑฺฒนนฺ}ติ สพฺพํ ตํ ตกฺกวฑฺฒนํ วิตกฺก\-วฑฺฒนํ สงฺกปฺป\-วฑฺฒนํ กาม\-วิตกฺก\-วฑฺฒนํ พฺยาปาท\-วิตกฺก\-วฑฺฒนํ วิหิํสา\-วิตกฺก\-วฑฺฒนํ ญาติ\-วิตกฺก\-วฑฺฒนํ ชนปท\-วิตกฺก\-วฑฺฒนํ อมรา\-วิตกฺก\-วฑฺฒนํ ปรานุทยตาปฏิสํยุตฺต\-วิตกฺก\-วฑฺฒนํ ลาภ\-สกฺการ\-สิโลก- ปฏิสํยุตฺต\-วิตกฺก\-วฑฺฒนํ อนวญฺญตฺติปฏิสํยุตฺตวิตกฺก\-วฑฺฒนนฺติ สพฺพํ ตํ ตกฺกวฑฺฒนํ. เตนาห เถโร ปิงฺคิโย\csromandash{}}

\setlength{\csromangathapagewidth}{0pt}
\setlength{\csromangathawakwidth}{0pt}
\csromangathameasuregroup{เย เม ปุพฺเพ วิยากํสุ, \\ \textbf{เอโก ตโมนุทาสีนฺ}โอ, \\ \textbf{โคตโม ภูริปญฺญฺ}อาโณ,}{\csromangathabat{เย เม ปุพฺเพ วิยากํสุ,}{หุรํ โคตมสาสนา,} \\ “อิจฺจาสิ อิติ ภวิสฺสติ.” \\ สพฺพํ ตํ อิติหีติหํ, \\ สพฺพํ ตํ ตกฺกวฑฺฒนนฺติ. \\ \csromangathabat{\textbf{เอโก ตโมนุทาสีนฺ}โอ,}{ชุติมา โส ปภงฺกโร.} \\ \csromangathabat{\textbf{โคตโม ภูริปญฺญฺ}อาโณ,}{โคตโม ภูริเมธโส.}}
\csromangathasetleft{เย เม ปุพฺเพ วิยากํสุ, \\ \textbf{เอโก ตโมนุทาสีนฺ}โอ, \\ \textbf{โคตโม ภูริปญฺญฺ}อาโณ,}
\csromangathagroup{\csromangathastanza{\csromangathabat{เย เม ปุพฺเพ วิยากํสุ,}{หุรํ โคตมสาสนา,} \\ “อิจฺจาสิ อิติ ภวิสฺสติ.” \\ สพฺพํ ตํ อิติหีติหํ,}\csromangathastanza{สพฺพํ ตํ ตกฺก\-วฑฺฒนนฺติ.}\csromangathastanza{\csromangathaitembat{107}{\textbf{เอโก ตโมนุทาสีนฺ}โอ,}{ชุติมา โส ปภงฺกโร.} \\ \csromangathacontbat{\textbf{โคตโม ภูริปญฺญฺ}อาโณ,}{โคตโม ภูริเมธโส.}}}

\prose{\textbf{เอโก ตโมนุทาสีโน}ติ \textbf{เอโก}ติ ภควา ปพฺพชฺช\-สงฺขาเตน เอโก, อทุติยฏฺเฐน เอโก, ตณฺหาย ปหานฏฺเฐน เอโก, เอกนฺตวีตราโคติ เอโก, เอกนฺตวีตโทโสติ เอโก, เอกนฺตวีตโมโหติ เอโก, เอกนฺตนิกฺกิเลโสติ เอโก, เอกายนมคฺคํ คโตติ เอโก, เอโก อนุตฺตรํ สมฺมาสมฺโพธิํ อภิสมฺพุทฺโธติ เอโก.}

\prose{กถํ ภควา ปพฺพชฺช\-สงฺขาเตน เอโก, ภควา ทหโรว สมาโน สุสุ กาฬเกโส ภเทฺรน โยพฺพเนน สมนฺนาคโต ปฐเมน วยสา อกามกานํ มาตาปิตูนํ อสฺสุมุขานํ โรทนฺตานํ วิลปนฺตานํ ญาติสํฆํสพฺพํ ฆราวาส\-ปลิโพธํ ฉินฺทิตฺวา ปุตฺตทาร\-ปลิโพธํ ฉินฺทิตฺวา ญาติปลิโพธํ ฉินฺทิตฺวา มิตฺตามจฺจ\-ปลิโพธํ ฉินฺทิตฺวา สนฺนิธิ\-ปลิโพธํ ฉินฺทิตฺวา เกสมสฺสุํ โอหาเรตฺวา กาสายานิ วตฺถานิ อจฺฉาเทตฺวา อคารสฺมา อนคาริยํ ปพฺพชิตฺวา อกิญฺจนภาวํ อุปคนฺตฺวา เอโก จรติ วิหรติ อิริยติ วตฺเตติ ปาเลติ ยเปติ ยาเปติ. เอวํ ภควา ปพฺพชฺช\-สงฺขาเตน เอโก.}

\prose{กถํ ภควา อทุติยฏฺเฐน เอโก, เอวํ ปพฺพชิโต สมาโน เอโก อรญฺญ\-วนปตฺถานิ ปนฺตานิ เสนาสนานิ ปฏิเสวติ อปฺปสทฺทานิ \csromanfolio{211}อปฺปนิคฺโฆสานิ วิชนวาตานิ มนุสฺส\-ราหสฺเสยฺยกานิ\csromansharedfootnote{6f415e1a4f030cf0}{มนุสฺสราหเสยฺยกานิ (สฺยา, ก)} ปฏิสลฺลาน\-สารุปฺปานิ\csromansharedfootnote{1dde5abc50fa62a8}{ปฏิสลฺลาณสารุปฺปานิ (ก)}, โส เอโก คจฺฉติ, เอโก ติฏฺฐติ, เอโก นิสีทติ, เอโก เสยฺยํ กปฺเปติ, เอโก คามํ ปิณฺฑาย ปวิสติ, เอโก อภิกฺกมติ, เอโก ปฏิกฺกมติ, เอโก รโห นิสีทติ, เอโก จงฺกมํ อธิฏฺฐาติ, เอโก จรติ วิหรติ อิริยติ วตฺเตติ ปาเลติ ยเปติ ยาเปติ เอวํ ภควา อทุติยฏฺเฐน เอโก.}

\prose{กถํ ภควา ตณฺหาย ปหานฏฺเฐน เอโก, โส เอวํ เอโก อทุติโย อปฺปมตฺโต อาตาปี ปหิตตฺโต วิหรนฺโต นชฺชา เนรญฺชราย ตีเร โพธิรุกฺข\-มูเล มหาปธานํ ปทหนฺโต มารํ สเสนํ กณฺหํ นมุจิํ ปมตฺตพนฺธุํ วิธมิตฺวา ตณฺหาชาลินิํ\csromansharedfootnote{ba1347b486c90232}{ตณฺหํ ชาลินิํ (สฺยา)} วิสฏํ\csromansharedfootnote{91d348722152ecdc}{สริตํ (สฺยา) ขุ 7. 362 ปิฏฺเฐปิ.} วิสตฺติกํ ปชหิ วิโนเทสิ พฺยนฺตีอกาสิ อนภาวํ คเมสิ.}

\setlength{\csromangathapagewidth}{0pt}
\setlength{\csromangathawakwidth}{0pt}
\csromangathameasuregroup{“ตณฺหาทุติโย ปุริโส, \\ อิตฺถภาวญฺญถาภาวํ, \\ เอตมาทีนวํ\textsuperscript{1} ญ ตฺวา, \\ วีตตโณฺห อนาทาโน,}{\csromangathabat{“ตณฺหาทุติโย ปุริโส,}{ทีฆมทฺธาน สํสรํ.} \\ \csromangathabat{อิตฺถภาวญฺญถาภาวํ,}{สํสารํ นาติวตฺตติ.} \\ \csromangathabat{เอตมาทีนวํ\textsuperscript{1} ญ ตฺวา,}{ตณฺหํ\textsuperscript{1} ทุกฺขสฺส สมฺภวํ.} \\ \csromangathabat{วีตตโณฺห อนาทาโน,}{สโต ภิกฺขุ ปริพฺพเช”ติ.}}
\csromangathasetleft{“ตณฺหาทุติโย ปุริโส, \\ อิตฺถภาวญฺญถาภาวํ, \\ เอตมาทีนวํ\textsuperscript{1} ญ ตฺวา, \\ วีตตโณฺห อนาทาโน,}
\csromangathagroup{\csromangathastanza{\csromangathabat{“ตณฺหาทุติโย ปุริโส,}{ทีฆมทฺธาน สํสรํ.} \\ \csromangathabat{อิตฺถ\-ภาว\-ญฺญถา\-ภาวํ,}{สํสารํ นาติวตฺตติ.}}\csromangathastanza{\csromangathabat{เอตมาทีนวํ\csromansharedfootnote{82482eab8b1a700d}{เอวมาทีนวํ (ก) ปสฺส ขุ 1. 201 ปิฏฺเฐ.} ญ ตฺวา,}{ตณฺหํ\csromansharedfootnote{fa9c36cb8bd1277f}{ตณฺหา (สฺยา, ก) ขุ 7. 362 ปิฏฺเฐ.} ทุกฺขสฺส สมฺภวํ.} \\ \csromangathabat{วีตตโณฺห อนาทาโน,}{สโต ภิกฺขุ ปริพฺพเช”ติ.}}}

\prose{เอวํ ภควา ตณฺหาย ปหานฏฺเฐน เอโก.}

\prose{กถํ ภควา เอกนฺตวีตราโคติ เอโก, ราคสฺส ปหีนตฺตา เอกนฺตวีตราโคติ เอโก, โทสสฺส ปหีนตฺตา เอกนฺตวีตโทโสติ เอโก, โมหสฺส ปหีนตฺตา เอกนฺตาวีตโมโหติ เอโก, กิเลสานํ ปหีนตฺตา เอกนฺตนิกฺกิเลโสติ เอโก.}

\prose{กถํ ภควา เอกายนมคฺคํ คโตติ เอโก, เอกายนมคฺโค วุจฺจติ จตฺตาโร สติปฏฺฐานา ฯเปฯ อริโย อฏฺฐงฺคิโก มคฺโค.}

\setlength{\csromangathapagewidth}{0pt}
\setlength{\csromangathawakwidth}{0pt}
\csromangathameasuregroup{}{“เอกายนํ ชาติขยนฺตทสฺสี, \\ มคฺคํ ปชานาติ หิตานุกมฺปี. \\ เอเตน มคฺเคน ตริํสุ\textsuperscript{1} ปุพฺเพ, \\ ตริสฺสนฺติ เย จ ตรนฺติ โอฆนฺ”ติ.}
\csromangathameasurewak{“เอกายนํ ชาติขยนฺตทสฺสี, \\ มคฺคํ ปชานาติ หิตานุกมฺปี. \\ เอเตน มคฺเคน ตริํสุ\textsuperscript{1} ปุพฺเพ, \\ ตริสฺสนฺติ เย จ ตรนฺติ โอฆนฺ”ติ.}
\setlength{\csromangathaleftcol}{0pt}
\csromangathagroup{\csromangathastanza{\csromangathawak{\csromanfolio{212}“เอกายนํ ชาติ\-ขยนฺต\-ทสฺสี, \\ มคฺคํ ปชานาติ หิตานุกมฺปี. \\ เอเตน มคฺเคน ตริํสุ\csromansharedfootnote{382fc189ad86788c}{อตริํสุ (ก) ปสฺส สํ 3. 162 ปิฏฺเฐ, ขุ 7. 362 ปิฏฺเฐ จ.} ปุพฺเพ, \\ ตริสฺสนฺติ เย จ ตรนฺติ โอฆนฺ”ติ.}}}

\proseitem{18}{เอวํ ภควา เอกายนมคฺคํ คโตติ เอโก.}

\prose{กถํ ภควา เอโก อนุตฺตรํ สมฺมาสมฺโพธิํ อภิสมฺพุทฺโธติ เอโก. โพธิ วุจฺจติ จตูสุ มคฺเคสุ ญาณํ, ปญฺญา ปญฺญินฺทฺริยํ ปญฺญาพลํ ธมฺมวิจย\-สมฺโพชฺฌงฺโค วีมํสา วิปสฺสนา สมฺมาทิฏฺฐิ, ภควา เตน โพธิญาเณน “สพฺเพ สงฺขารา อนิจฺจา”ติ พุชฺฌิ, “สพฺเพ สงฺขารา ทุกฺขา”ติ พุชฺฌิ ฯเปฯ “สพฺเพ ธมฺมา อนตฺตา”ติ พุชฺฌิ,. “ยํ กิญฺจิ สมุทย\-ธมฺมํ, สพฺพํ ตํ นิโรธธมฺมนฺ”ติ พุชฺฌิ. อถ วา ยํ พุชฺฌิตพฺพํ อนุพุชฺฌิตพฺพํ ปฏิพุชฺฌิตพฺพํ สมฺพุชฺฌิตพฺพํ อธิคนฺตพฺพํ ผสฺสิตพฺพํ สจฺฉิกาตพฺพํ, สพฺพํ ตํ เตน โพธิญาเณน พุชฺฌิ อนุพุชฺฌิ ปฏิพุชฺฌิ สมฺพุชฺฌิ อธิคจฺฉิ ผสฺเสสิ สจฺฉากาสิ. เอวํ ภควา เอโก อนุตฺตรํ สมฺมาสมฺโพธิํ อภิสมฺพุทฺโธติ เอโก.}

\prose{\textbf{ตโมนุโท}ติ ภควา ราคตมํ โทสตมํ โมหตมํ ทิฏฺฐิตมํ กิเลสตมํ ทุจฺจริตตมํ อนฺธกรณํ อจกฺขุกรณํ อญฺญาณกรณํ ปญฺญา\-นิโรธิกํ วิฆาต\-ปกฺขิกํ อนิพฺพานสํวตฺตนิกํ นุทิ ปนุทิ ปชหิ วิโนเทสิ พฺยนฺตีอกาสิ อนภาวํ คเมสิ. \textbf{อาสีโน}ติ นิสินฺโน ภควา ปาสาณเก เจติเยติ อาสีโน\csromansharedfootnote{07e457fdbaed8967}{อาสิโน (ก)}.}

\setlength{\csromangathapagewidth}{0pt}
\setlength{\csromangathawakwidth}{0pt}
\csromangathameasuregroup{นคสฺส ปสฺเส อาสีนํ, \\ สาวกา ปยิรุปาสนฺติ,}{\csromangathabat{นคสฺส ปสฺเส อาสีนํ,}{มุนิํ ทุกฺขสฺส ปารคุํ.} \\ \csromangathabat{สาวกา ปยิรุปาสนฺติ,}{เตวิชฺชา มจฺจุหายิโนติ.}}
\csromangathasetleft{นคสฺส ปสฺเส อาสีนํ, \\ สาวกา ปยิรุปาสนฺติ,}
\csromangathagroup{\csromangathastanza{\csromangathabat{นคสฺส ปสฺเส อาสีนํ,}{มุนิํ ทุกฺขสฺส ปารคุํ.} \\ \csromangathabat{สาวกา ปยิรุปาสนฺติ,}{เตวิชฺชา มจฺจุหายิโนติ.}}}

\prose{เอวมฺปิ ภควา อาสีโน.  อถ วา ภควา สพฺโพ\-สฺสุกฺก\-ปฏิปฺปสฺสทฺธ\-ตฺตา อาสีโน, โส วุตฺถวาโส จิณฺณจรโณ ฯเปฯ ชาติ\-มรณ\-สํสาโร, นตฺถิ ตสฺส ปุนพฺภโวติ. เอวมฺปิ ภควา อาสีโนติ เอโก ตโมนุทาสีโน.}

\prose{\textbf{ชุติมา โส ปภงฺกโร}ติ \textbf{ชุติมา}ติ ชุติมา มติมา ปณฺฑิโต ปญฺญวา พุทฺธิมา ญาณี วิภาวี เมธาวิ. \textbf{ปภงฺกโร}ติ ปภงฺกโร อาโลกกโร โอภาสกโร ทีปงฺกโร ปทีปกโร อุชฺโชตกโร ปชฺโชตกโรติ ชุติมา โส ปภงฺกโร.}

\prose{\csromanfolio{213}\textbf{โคตโม ภูริปญฺญาโณ}ติ โคตโม ภูริปญฺญาโณ ญาณปญฺญาโณ ปญฺญาธโช ปญฺญาเกตุ ปญฺญา\-ธิปเตยฺโย วิจยพหุโล ปวิจยพหุโล โอกฺขายน\-พหุโล สโมกฺขายน\-ธมฺโม วิภูตวิหารี ตจฺจริโต ตพฺพหุโล ตคฺครุโก ตนฺนินฺโน ตปฺโปโณ ตปฺปพฺภาโร ตทธิมุตฺโต ตทธิปเตยฺโย.}

\setlength{\csromangathapagewidth}{0pt}
\setlength{\csromangathawakwidth}{0pt}
\csromangathameasuregroup{ธโช รถสฺส ปญฺญาณํ, \\ ราชา รฏฺฐสฺส ปญฺญาณํ,}{\csromangathabat{ธโช รถสฺส ปญฺญาณํ,}{ธูโม\textsuperscript{1} ปญฺญาณมคฺคิโน.} \\ \csromangathabat{ราชา รฏฺฐสฺส ปญฺญาณํ,}{ภตฺตา ปญฺญาณมิตฺถิยาติ.}}
\csromangathasetleft{ธโช รถสฺส ปญฺญาณํ, \\ ราชา รฏฺฐสฺส ปญฺญาณํ,}
\csromangathagroup{\csromangathastanza{\csromangathabat{ธโช รถสฺส ปญฺญาณํ,}{ธูโม\csromansharedfootnote{d3455bc877588eb8}{ธุโม (สฺยา)} ปญฺญาณมคฺคิโน.} \\ \csromangathabat{ราชา รฏฺฐสฺส ปญฺญาณํ,}{ภตฺตา ปญฺญาณมิ\-ตฺถิยาติ.}}}

\prose{เอวเมว โคตโม ภูริปญฺญาโณ ญาณปญฺญาโณ ปญฺญาธโช ปญฺญาเกตุ ปญฺญา\-ธิปเตยฺโย วิจยพหุโล ปวิจยพหุโล โอกฺขายน\-พหุโล สโมกฺขายน\-ธมฺโม วิภูตวิหารี ตจฺจริโต ตพฺพหุโล ตคฺครูโก ตนฺนินฺโน ตปฺโปโณ ตปฺปพฺภาโร ตทธิมุตฺโต ตทธิปเตยฺโยติ โคตโม ภูริปญฺญาโณ.}

\prose{\textbf{โคตโม ภูริเมธโส}ติ ภูริ วุจฺจติ ปถวี, ภควา ตาย ปถวิสมาย ปญฺญาย วิปุลาย วิตฺถตาย สมนฺนาคโต. เมธา วุจฺจติ ปญฺญา, ยา ปญฺญา ปชานนา ฯเปฯ อโมโห ธมฺมวิจโย สมฺมาทิฏฺฐิ. ภควา อิมาย เมธาย อุเปโต สมุเปโต อุปาคโต สมุปาคโต อุปปนฺโน สมุปปนฺโน สมนฺนาคโต, ตสฺมา พุทฺโธ สุเมธโสติ\csromansharedfootnote{f432bbdc3aa43dd8}{ภูริเมธโสติ (สฺยา) เอวมุปริปิ.} โคตโม ภูริเมธโส. เตนาห เถโร ปิงฺคิโย\csromandash{}}

\setlength{\csromangathapagewidth}{0pt}
\setlength{\csromangathawakwidth}{0pt}
\csromangathameasuregroup{เอโก ตโมนุทาสีโน, \\ โคตโม ภูริปญฺญาโณ, \\ \textbf{โย เม ธมฺมมเทสฺ}เอสิ, \\ \textbf{ตณฺหกฺขยมนีติกฺ}อํ,}{\csromangathabat{เอโก ตโมนุทาสีโน,}{ชุติมา โส ปภงฺกโร.} \\ \csromangathabat{โคตโม ภูริปญฺญาโณ,}{\textbf{โคตโม ภูริเมธโส}ติ.} \\ \csromangathabat{\textbf{โย เม ธมฺมมเทสฺ}เอสิ,}{สนฺทิฏฺฐิกมกาลิกํ.} \\ \csromangathabat{\textbf{ตณฺหกฺขยมนีติกฺ}อํ,}{ยสฺส นตฺถิ อุปมา กฺวจิ.}}
\csromangathasetleft{เอโก ตโมนุทาสีโน, \\ โคตโม ภูริปญฺญาโณ, \\ \textbf{โย เม ธมฺมมเทสฺ}เอสิ, \\ \textbf{ตณฺหกฺขยมนีติกฺ}อํ,}
\csromangathagroup{\csromangathastanza{\csromangathabat{เอโก ตโมนุทาสีโน,}{ชุติมา โส ปภงฺกโร.} \\ \csromangathabat{โคตโม ภูริปญฺญาโณ,}{\textbf{โคตโม ภูริเมธโส}ติ.}}\csromangathastanza{\csromangathaitembat{108}{\textbf{โย เม ธมฺมมเทสฺ}เอสิ,}{สนฺทิฏฺฐิกม\-กาลิกํ.} \\ \csromangathacontbat{\textbf{ตณฺหกฺขยมนีติกฺ}อํ,}{ยสฺส นตฺถิ อุปมา กฺวจิ.}}}

\prose{\textbf{โย เม ธมฺมมเทเสสี}ติ \textbf{โย}ติ โย โส ภควา สยมฺภู อนาจริยโก ปุพฺเพ อนนุสฺสุเตสุ ธมฺเมสุ สามํ สจฺจานิ อภิสมฺพุชฺฌิ, ตตฺถ จ สพฺพญฺญุตํ ปตฺโต พเลสุ จ วสีภาวํ. \textbf{ธมฺมมเทเสสี}ติ \textbf{ธมฺมนฺ}ติ อาทิกลฺยาณํ มชฺเฌกลฺยาณํ ปริโยสาน\-กลฺยาณํ สาตฺถํ สพฺยญฺชนํ เกวล\-ปริปุณฺณํ ปริสุทฺธํ พฺรหฺมจริยํ จตฺตาโร สติปฏฺฐเน ฯเปฯ อริยํ อฏฺฐงฺคิกํ มคฺคํ นิพฺพานญฺจ นิพฺพานคามินิญฺจ ปฏิปทํ อาจิกฺขิ เทเสสิ ปญฺญเปสิ \csromanfolio{214}ปฏฺฐเปสิ วิวริ วิภชิ อุตฺตานีอกาสิ ปกาเสสีติ โย เม ธมฺมมเทเสสิ.}

\prose{\textbf{สนฺทิฏฺฐิกมกาลิกนฺ}ติ สนฺทิฏฺฐิกํ อกาลิกํ เอหิปสฺสิกํ โอปเนยฺยิกํ ปจฺจตฺตํ เวทิตพฺพํ วิญฺญูหีติ เอวํ สนฺทิฏฺฐิกํ. อถ วา โย ทิฏฺเฐว ธมฺเม อริยํ อฏฺฐงฺคิกํ มคฺคํ ภาเวติ, ตสฺส มคฺคสฺส อนนฺตรา สมนนฺตรา อธิคจฺฉเตว ผลํ วินฺทติ ปฏิลภตีติ เอวมฺปิ สนฺทิฏฺฐิกม\-กาลิกํ. ยถา มนุสฺสา กาลิกํ ธนํ ทตฺวา อนนฺตรา น ลภนฺติ กาลํ อาคเมนฺติ, เนวายํ ธมฺโม. โย ทิฏฺเฐว ธมฺเม อริยํ อฏฺฐงฺคิกํ มคฺคํ ภาเวติ, ตสฺส มคฺคสฺส อนนฺตรา สมนนฺตรา อธิคจฺฉเตว ผลํ วินฺทติ ปฏิลภติ, น ปรตฺถ น ปรโลเก, เอวํ อกาลิกนฺติ สนฺทิฏฺฐิกม\-กาลิกํ.}

\prose{\textbf{ตณฺหกฺขยมนีติกนฺ}ติ ตณฺหาติ รูป\textbf{ตณฺหา} ฯเปฯ ธมฺมตณฺหา. \textbf{ตณฺหกฺขยนฺ}ติ ตณฺหกฺขยํ ราคกฺขยํ โทสกฺขยํ โมหกฺขยํ คติกฺขยํ อุปปตฺติ\-กฺขยํ ปฏิสนฺธิ\-กฺขยํ ภวกฺขยํ สํสารกฺขยํ วฏฺฏกฺขยํ. \textbf{อนีติกนฺ}ติ \textbf{อีติ}วุจฺจนฺติ กิเลสา จ ขนฺธา จ อภิสงฺขารา จ. ¢ติปฺปหานํ อีติวูปสมํ อีติปฏินิสฺสคฺคํ อีติปฏิปฺปสฺสทฺธิํ อมตํ นิพฺพานนฺติ ตณฺหกฺขยม\-นีติกํ.}

\prose{\textbf{ยสฺส นตฺถิ อุปมา กฺวจี}ติ \textbf{ยสฺสา}ติ นิพฺพานสฺส. \textbf{นตฺถิ อุปมา}ติ อุปมา นตฺถิ, อุปนิธา นตฺถิ, สทิสํ นตฺถิ, ปฏิภาโค นตฺถิ น อตฺถิ น สํวิชฺชติ นุปลพฺภติ. \textbf{กฺวจี}ติ กฺวจิ กิมฺหิจิ กตฺถจิ อชฺฌตฺตํ วา พหิทฺธา วา อชฺฌตฺต\-พหิทฺธา วาติ ยสฺส นตฺถิ อุปมา กฺวจิ. เตนาห เถโร ปิงฺคิโย\csromandash{}}

\setlength{\csromangathapagewidth}{0pt}
\setlength{\csromangathawakwidth}{0pt}
\csromangathameasuregroup{โย เม ธมฺมมเทเสสิ, \\ ตณฺหกฺขยมนีติกํ, \\ \textbf{กิํ นุ ตมฺหา วิปฺป}วสสิ, \\ \textbf{โคตมา ภูริปญฺ}ญาณา,}{\csromangathabat{โย เม ธมฺมมเทเสสิ,}{สนฺทิฏฺฐิกมกาลิกํ.} \\ \csromangathabat{ตณฺหกฺขยมนีติกํ,}{ยสฺส นตฺถิ อุปมา กฺวจฺ\textbf{อีติ}.} \\ \csromangathabat{\textbf{กิํ นุ ตมฺหา วิปฺป}วสสิ,}{มุหุตฺตมปิ ปิงฺคิย.} \\ \csromangathabat{\textbf{โคตมา ภูริปญฺ}ญาณา,}{โคตมา ภูริเมโส.}}
\csromangathasetleft{โย เม ธมฺมมเทเสสิ, \\ ตณฺหกฺขยมนีติกํ, \\ \textbf{กิํ นุ ตมฺหา วิปฺป}วสสิ, \\ \textbf{โคตมา ภูริปญฺ}ญาณา,}
\csromangathagroup{\csromangathastanza{\csromangathabat{โย เม ธมฺมมเทเสสิ,}{สนฺทิฏฺฐิกม\-กาลิกํ.} \\ \csromangathabat{ตณฺหกฺขยม\-นีติกํ,}{ยสฺส นตฺถิ อุปมา กฺวจฺ\textbf{อีติ}.}}\csromangathastanza{\csromangathaitembat{109}{\textbf{กิํ นุ ตมฺหา วิปฺป}วสสิ,}{มุหุตฺตมปิ ปิงฺคิย.} \\ \csromangathacontbat{\textbf{โคตมา ภูริปญฺ}ญาณา,}{โคตมา ภูริเมโส.}}}

\prose{\textbf{กิํนุ ตมฺหา วิปฺปวสสี}ติ กิํนุ พุทฺธมฺหา วิปฺปวสสิ อเปสิ อปคจฺฉสิ\csromansharedfootnote{35dc2d331b743188}{อปคจฺฉิ (สี)} วินา โหสีติ กิํนุ ตมฺหา วิปฺปวสสิ.}

\prose{\textbf{มุหุตฺตมปิ ปิงฺคิยา}ติ มุหุตฺตมฺปิ ขณมฺปิ ลยมฺปิ วยมฺปิ อทฺธมฺปีติ มุหุตฺตมปิ. \textbf{ปิงฺคิยา}ติ พาวรี ตํ นตฺตารํ นาเมน อาลปติ.}

\prose{\csromanfolio{215}\textbf{โคตมา ภูริปญฺญาณา}ติ โคตมา ภูริปญฺญาณา ญาณปญฺญาณา ปญฺญาธชา ปญฺญาเกตุมฺหา ปญฺญา\-ธิปเตยฺยมฺหา วิจยพหุลา ปวิจยพหุลา โอกฺขายน\-พหุลา สโมกฺขายน\-ธมฺมา วิภูต\-วิหาริมฺหา ตจฺจริตา ตพฺพหุลา ตคฺครุกา ตนฺนินฺนา ตปฺโปณา ตปฺปพฺภารา ตทธิมุตฺตา ตทธิปเตยฺยมฺหาติ โคตมา ภูริปญฺญาณา.}

\prose{\textbf{โคตมา ภูริเมธส}ติ ภูริ วุจฺจติ ปถวี, ภควา ตาย ปถวิสมาย ปญฺญาย วิปุลาย วิตฺถตาย สมนฺนาคโต. เมธา วุจฺจติ ปญฺญา, ยา ปญฺญา ปชานนา ฯเปฯ อโมโห ธมฺมวิจโย สมฺมาทิฏฺฐิ. ภควา อิมาย เมธาย ปญฺญาย อุเปโต สมุเปโต อุปาคโต สมุปาคโต อุปปนฺโน สมุปปนฺโน สมนฺนาคโต, ตสฺมา พุทฺโธ สุเมธโสติ โคตมา ภูริเมธส. เตนาห โส พฺราหฺมโณ\csromandash{}}

\setlength{\csromangathapagewidth}{0pt}
\setlength{\csromangathawakwidth}{0pt}
\csromangathameasuregroup{กิํนุ ตมฺหา วิปฺปวสิ, \\ โคตมา ภูริปญฺญาณา,}{\csromangathabat{กิํนุ ตมฺหา วิปฺปวสิ,}{มุหุตฺตมปิ ปิงฺคิย.} \\ \csromangathabat{โคตมา ภูริปญฺญาณา,}{โคตมา ภูริเมธสาติ.}}
\csromangathasetleft{กิํนุ ตมฺหา วิปฺปวสิ, \\ โคตมา ภูริปญฺญาณา,}
\csromangathagroup{\csromangathastanza{\csromangathabat{กิํนุ ตมฺหา วิปฺปวสิ,}{มุหุตฺตมปิ ปิงฺคิย.} \\ \csromangathabat{โคตมา ภูริปญฺญาณา,}{โคตมา ภูริเมธสาติ.}}}

\proseitem{110}{\textbf{โย เต ธมฺมมเทเสสิ, สนฺทิฏฺฐิกม}\-\textbf{กาลิกํ.}}

\prose{\textbf{ตณฺหกฺขยม}\-\textbf{นีติกํ, ยสฺส นตฺถิ อุปมา กฺวจิ.} (9)}

\prose{\textbf{โย เต ธมฺมมเทเสสี}ติ โย โส ภควา ฯเปฯ ตตฺถ จ สพฺพญฺญุตํ ปตฺโต พเลสุ จ วสีภาวํ. \textbf{ธมฺมมเทเสสี}ติ \textbf{ธมฺมนฺ}ติ อาทิกลฺยาณํ มชฺเฌกลฺยาณํ ฯเปฯ นิพฺพานญฺจ นิพฺพานคามินิญฺจ ปฏิปทํ อาจิกฺขิ เทเสสิ ปญฺญเปสิ ปฏฺฐเปสิ วิวริ วิภชิ อุตฺตานีอกาสิ ปกาเสสีติ โย เต ธมฺมมเทเสสิ.}

\prose{\textbf{สนฺทิฏฺฐิกมกาลิกนฺ}ติ สนฺทิฏฺฐิกํ อกาลิกํ เอหิปสฺสิกํ โอปเนยฺยิกํ ปจฺจตฺตํ เวทิตพฺพํ วิญฺญูหีติ เอวํ สนฺทิฏฺฐิกํ. อถ วา โย ทิฏฺเฐว ธมฺเม อริยํ อฏฺฐงฺคิกํ มคฺคํ ภาเวติ, ตสฺส มคฺคสฺส อนนฺตรา สมนนฺตรา อธิคจฺฉเตว ผลํ วินฺทติ ปฏิลภตีติ เอวมฺปิ สนฺทิฏฺฐิกํ. \textbf{อกาลิกนฺ}ติ ยถา มนุสฺสา กาลิกํ ธนํ ทตฺวา อนนฺตรา น ลภนฺติ, กาลํ อาคเมนฺติ, เนวายํ ธมฺโม. โย ทิฏฺเฐว ธมฺเม อริยํ อฏฺฐงฺคิกํ มคฺคํ ภาเวติ. ตสฺส มคฺคสฺส อนนฺตรา สมนนฺตรา อธิคจฺฉเตว ผลํ วินฺทติ ปฏิลภติ, น ปรตฺถ น ปรโลเก, เอวํ อกาลิกนฺติ สนฺทิฏฺฐิกม\-กาลิกํ.}

\prose{\csromanfolio{216}\textbf{ตณฺหกฺขยมนีติกนฺ}ติ ตณฺหาติ รูป\textbf{ตณฺหา} ฯเปฯ ธมฺมตณฺหา. \textbf{ตณฺหกฺขยนฺ}ติ ตณฺหกฺขยํ ราคกฺขยํ โทสกฺขยํ โมหกฺขยํ คติกฺขยํ อุปปตฺติ\-กฺขยํ ปฏิสนฺธิ\-กฺขยํ ภวกฺขยํ สํสารกฺขยํ วฏฺฏกฺขยํ. \textbf{อนีติกนฺ}ติ \textbf{อีติ} วุจฺจนฺติ กิเลสา จ ขนฺธา จ อภิสงฺขารา จ. ¢ติปฺปหานํ อีติวูปสมํ อีติปฏินิสฺสคฺคํ อีติปฏิปฺปสฺสทฺธิํ อมตํ นิพฺพานนฺติ ตณฺหกฺขยม\-นีติกํ.}

\prose{\textbf{ยสฺส นตฺถิ อุปมา กฺวจี}ติ \textbf{ยสฺสา}ติ นิพฺพานสฺส. \textbf{นตฺถิ อุปมา}ติ อุปมา นตฺถิ, อุปนิธา นตฺถิ, สทิสํ นตฺถิ, ปฏิภาโค นตฺถิ น สติ น สํวิชฺชติ นุปลพฺภติ. \textbf{กฺวจี}ติ กฺวจิ กิมฺหิจิ กตฺถจิ อชฺฌตฺตํ วา พหิทฺธา วา อชฺฌตฺต\-พหิทฺธา วาติ ยสฺส นตฺถิ อุปมา กฺวจิ. เตนาห โส พฺราหฺมโณ\csromandash{}}

\setlength{\csromangathapagewidth}{0pt}
\setlength{\csromangathawakwidth}{0pt}
\csromangathameasuregroup{โย เต ธมฺมมเทเสสิ, \\ ตณฺหกฺขยมนีติกํ,}{\csromangathabat{โย เต ธมฺมมเทเสสิ,}{สนฺทิฏฺฐิกมกาลิกํ.} \\ \csromangathabat{ตณฺหกฺขยมนีติกํ,}{ยสฺส นตฺถิ อุปมา กฺวจฺ\textbf{อีติ}.}}
\csromangathasetleft{โย เต ธมฺมมเทเสสิ, \\ ตณฺหกฺขยมนีติกํ,}
\csromangathagroup{\csromangathastanza{\csromangathabat{โย เต ธมฺมมเทเสสิ,}{สนฺทิฏฺฐิกม\-กาลิกํ.} \\ \csromangathabat{ตณฺหกฺขยม\-นีติกํ,}{ยสฺส นตฺถิ อุปมา กฺวจฺ\textbf{อีติ}.}}}

\proseitem{111}{\textbf{นาหํ ตมฺหา วิปฺปวสามิ, มุหุตฺตมปิ พฺราหฺมณ.}}

\prose{\textbf{โคตมา ภูริปญฺญาณา, โคตมา ภูริเมธสา.} (10)}

\prose{\textbf{นาหํ ตมฺหา วิปฺปวสามี}ติ นาหํ พุทฺธมฺหา วิปฺปวสามิ อเปมิ อปคจฺฉามิ วินา โหมีติ นาหํ ตมฺหา วิปฺปวสามิ.}

\prose{\textbf{มุหุตฺตมปิ พฺราหฺมณา}ติ มุหุตฺตมฺปิ ขณมฺปิ ลยมฺปิ วยมฺปิ อทฺธมฺปีติ มุหุตฺตมปิ. \textbf{พฺราหฺมณา}ติ คารเวน มาตุลํ อาลปติ.}

\prose{\textbf{โคตมา ภูริปญฺญาณา}ติ โคตมา ภูริปญฺญาณา ญาณปญฺญาณา ปญฺญาธชา ปญฺญาเกตุมฺหา ปญฺญา\-ธิปเตยฺยมฺหา วิจยพหุลา ปวิจยพหุลา โอกฺขายน\-พหุลา สโมกฺขายน\-ธมฺมา วิภูต\-วิหาริมฺหา ตจฺจริตา ตพฺพหุลา ตคฺครุกา ตนฺนินฺนา ตปฺโปณา ตปฺปพฺภารา ตทธิมุตฺตา ตทธิปเตยฺยมฺหาติ โคตมา ภูริปญฺญาณา.}

\prose{\textbf{โคตมา ภูริเมธสา}ติ ภูริ วุจฺจติ ปถวี, ภควา ตาย ปถวิสมาย ปญฺญาย วิปุลาย วิตฺถตาย สมนฺนาคโต. เมธา วุจฺจติ ปญฺญา, ยา ปญฺญา ปชานนา ฯเปฯ อโมโห ธมฺมวิจโย สมฺมาทิฏฺฐิ. ภควา อิมาย เมธาย ปญฺญาย อุเปโต สมุเปโต อุปาคโต สมุปาคโต อุปปนฺโน สมุปปนฺโน สมนฺนาคโต, ตสฺมา พุทฺโธ สุเมธโสติ โคตมา ภูริเมธสา. เตนาห เถโร ปิงฺคิโย\csromandash{}}

\setlength{\csromangathapagewidth}{0pt}
\setlength{\csromangathawakwidth}{0pt}
\csromangathameasuregroup{นาหํ ตมฺหา วิปฺปวสามิ, \\ โคตมา ภูริปญฺญาณา,}{\csromangathabat{นาหํ ตมฺหา วิปฺปวสามิ,}{มุหุตฺตมปิ พฺราหฺมณ.} \\ \csromangathabat{โคตมา ภูริปญฺญาณา,}{โคตมา ภูริเมธสาติ.}}
\csromangathasetleft{นาหํ ตมฺหา วิปฺปวสามิ, \\ โคตมา ภูริปญฺญาณา,}
\csromangathagroup{\csromangathastanza{\csromanfolio{217}\csromangathabat{นาหํ ตมฺหา วิปฺปวสามิ,}{มุหุตฺตมปิ พฺราหฺมณ.} \\ \csromangathabat{โคตมา ภูริปญฺญาณา,}{โคตมา ภูริเมธสาติ.}}}

\proseitem{112}{\textbf{โย เม ธมฺมมเทเสสิ, สนฺทิฏฺฐิกม}\-\textbf{กาลิกํ.}}

\prose{\textbf{ตณฺหกฺขยม}\-\textbf{นีติกํ, ยสฺส นตฺถิ อุปมา กฺวจิ.} (11)}

\prose{\textbf{โย เม ธมฺมมเทเสสี}ติ โย โส ภควา สยมฺภู อนาจริยโก ปุพฺเพ อนนุสฺสุเตสุ ธมฺเมสุ สามํ สจฺจานิ อภิสมฺพุชฺฌิ, ตตฺถ จ สพฺพญฺญุตํ ปตฺโต พเลสุ จ วสีภาวํ. \textbf{ธมฺมมเทเสสี}ติ \textbf{ธมฺมนฺ}ติ อาทิกลฺยาณํ มชฺเฌกลฺยาณํ ปริโยสาน\-กลฺยาณํ สาตฺถํ สพฺยญฺชนํ เกวล\-ปริปุณฺณํ ปริสุทฺธํ พฺรหฺมจริยํ จตฺตาโร สติปฏฺฐาเน จตฺตาโร สมฺมปฺปธาเน จตฺตาโร อิทฺธิปาเท ปญฺจินฺทฺริยานิ ปญฺจ พลานิ สตฺต โพชฺฌงฺเค อริยํ อฏฺฐงฺคิกํ มคฺคํ นิพฺพานญฺจ นิพฺพานคามินิญฺจ ปฏิปทํ อาจิกฺขิ เทเสสิ ปญฺญเปสิ ปฏฺฐเปสิ วิวริ วิภชิ อุตฺตานีอกาสิ ปกาเสสีติ โย เม ธมฺมมเทเสสิ.}

\prose{\textbf{สนฺทิฏฺฐิกมกาลิกนฺ}ติ สนฺทิฏฺฐิกํ อกาลิกํ เอหิปสฺสิกํ โอปเนยฺยิกํ ปจฺจตฺตํ เวทิตพฺพํ วิญฺญูหีติ เอวํ สนฺทิฏฺฐิกํ. อถ วา โย ทิฏฺเฐว ธมฺเม อริยํ อฏฺฐงฺคิกํ มคฺคํ ภาเวติ, ตสฺส มคฺคสฺส อนนฺตรา สมนนฺตรา อธิคจฺฉเตว ผลํ วินฺทติ ปฏิลภตีติ เอวมฺปิ สนฺทิฏฺฐิกํ. \textbf{อกาลิกนฺ}ติ ยถา มนุสฺสา กาลิกํ ธนํ ทตฺวา อนนฺตรา น ลภนฺติ, กาลํ อาคเมนฺติ, เนวายํ ธมฺโม. โย ทิฏฺเฐว ธมฺเม อริยํ อฏฺฐงฺคิกํ มคฺคํ ภาเวติ, ตสฺส มคฺคสฺส อนนฺตรา สมนนฺตรา อธิคจฺฉเตว ผลํ วินฺทติ ปฏิลภติ, น ปรตฺถ น ปรโลเก, เอวํ อกาลิกนฺติ สนฺทิฏฺฐิกม\-กาลิกํ.}

\prose{\textbf{ตณฺหกฺขยมนีติกนฺ}ติ ตณฺหาติ รูป\textbf{ตณฺหา} ฯเปฯ ธมฺมตณฺหา. \textbf{ตณฺหกฺขยนฺ}ติ ตณฺหกฺขยํ ราคกฺขยํ โทสกฺขยํ โมหกฺขยํ คติกฺขยํ อุปปตฺติ\-กฺขยํ ปฏิสนฺธิ\-กฺขยํ ภวกฺขยํ สํสารกฺขยํ วฏฺฏกฺขยํ. \textbf{อนีติกนฺ}ติ \textbf{อีติ} วุจฺจนฺติ กิเลสา จ ขนฺธา จ อภิสงฺขารา จ. ¢ติปฺปหานํ อีติวูปสมํ อีติปฏิปฺปสฺสทฺธิํ อมตํ นิพฺพานนฺติ ตณฺหกฺขยม\-นีติกํ.}

\prose{\textbf{ยสฺส นตฺถิ อุปมา กฺวจี}ติ \textbf{ยสฺสา}ติ นิพฺพานสฺส. \textbf{นตฺถิ อุปมา}ติ อุปมา นตฺถิ, อุปนิธา นตฺถิ, สทิสํ นตฺถิ, ปฏิภาโค นตฺถิ น สติ น สํวิชฺชติ นุปลพฺภติ. \textbf{กฺวจี}ติ กฺวจิ กิมฺหิจิ กตฺถจิ อชฺฌตฺตํ วา พหิทฺธา วา อชฺฌตฺต\-พหิทฺธา วาติ ยสฺส นตฺถิ อุปมา กฺวจิ. เตนาห เถโร ปิงฺคิโย\csromandash{}}

\setlength{\csromangathapagewidth}{0pt}
\setlength{\csromangathawakwidth}{0pt}
\csromangathameasuregroup{โย เม ธมฺมมเทเสสิ, \\ ตณฺหกฺขยมนีติกํ,}{\csromangathabat{โย เม ธมฺมมเทเสสิ,}{สนฺทิฏฺฐิกมกาลิกํ.} \\ \csromangathabat{ตณฺหกฺขยมนีติกํ,}{ยสฺส นตฺถิ อุปมา กฺวจีติ.}}
\csromangathasetleft{โย เม ธมฺมมเทเสสิ, \\ ตณฺหกฺขยมนีติกํ,}
\csromangathagroup{\csromangathastanza{\csromanfolio{218}\csromangathabat{โย เม ธมฺมมเทเสสิ,}{สนฺทิฏฺฐิกม\-กาลิกํ.} \\ \csromangathabat{ตณฺหกฺขยม\-นีติกํ,}{ยสฺส นตฺถิ อุปมา กฺวจีติ.}}}

\proseitem{113}{\textbf{ปสฺสามิ นํ มนสา จกฺขุนาว,}}

\prose{\textbf{รตฺตินฺทิวํ พฺราหฺมณ อปฺปมตฺโต.}}

\prose{\textbf{นมสฺสมาโน วิวเสมิ}\csromansharedfootnote{3b334bef3300dd3e}{นมสฺสมาโนว วเสมิ (สีฏฺฐ) ...วิวสามิ (สฺยา)}\textbf{ รตฺติํ,}}

\prose{\textbf{เตเนว มญฺญามิ อวิปฺปวาสํ.} (12)}

\prose{\textbf{ปสฺสามิ นํ มนสา จกฺขุนาวา}ติ ยถา จกฺขุมา ปุริโส อาโลเก รูปคตานิ ปสฺเสยฺย ทกฺเขยฺย โอโลเกยฺย นิชฺฌาเยยฺย อุปปริกฺเขยฺย, เอวเมวาหํ พุทฺธํ ภควนฺตํ มนสา ปสฺสามิ ทกฺขามิ โอโลเกมิ นิชฺฌายามิ อุปปริกฺขามีติ ปสฺสามิ นํ มนสา จกฺขุนาว.}

\prose{\textbf{รตฺตินฺทิวํ พฺราหฺมณ อปฺปมตฺโต}ติ รตฺติญฺจ ทิวา จ พุทฺธานุสฺสติํ มนสา ภาเวนฺโต อปฺปมตฺโตติ รตฺตินฺทิวํ พฺราหฺมณ อปฺปมตฺโต.}

\prose{\textbf{นมสฺสมาโน วิวเสมิ รตฺตินฺ}ติ \textbf{นมสฺสมาโน}ติ กาเยน วา นมสฺสมาโน, วาจาย วา นมสฺสมาโน, จิตฺเตน วา นมสฺสมาโน, อนฺวตฺถ\-ปฏิปตฺติยา วา นมสฺสมาโน, ธมฺมานุธมฺม\-ปฏิปตฺติยา วา นมสฺสมาโน สกฺการมาโน ครุการมาโน มานยมาโน ปูชยมาโน รตฺตินฺทิวํ วิวเสมิ อตินาเมมิ อติกฺกเมมีติ นมสฺสมาโน วิวเสมิ รตฺติํ.}

\prose{\textbf{เตเนว มญฺญามิ อวิปฺปวาสนฺ}ติ ตาย พุทฺธา\-นุสฺสติยา ภาเวนฺโต อวิปฺปวาโสติ ตํ มญฺญามิ, อวิปฺปวุฏฺโฐติ ตํ มญฺญามิ ชานามิ, เอวํ ชานามิ เอวํ อาชานามิ เอวํ วิชานามิ เอวํ ปฏิวิชานามิ เอวํ ปฏิวิชฺฌามีติ เตเนว มญฺญามิ อวิปฺปวาสํ. เตนาห เถโร ปิงฺคิโย\csromandash{}}

\prose{\textbf{ปสฺสามิ นํ มนสา จกฺขุนาว,}}

\prose{\textbf{รตฺตินฺทิวํ พฺราหฺมณ อปฺปมตฺโต.}}

\prose{\textbf{นมสฺสมาโน วิวเสมิ รตฺติํ,}}

\prose{เตเนว มญฺญามิ อวิปฺปวาสนฺติ.}

\proseitem{114}{\csromanfolio{219}\textbf{สทฺธา จ ปีติ จ มโน สติ จ,}}

\prose{\textbf{นา’เปนฺติ’เม โคตม}\-\textbf{สาสนมฺหา.}}

\prose{\textbf{ยํ ยํ ทิสํ วชติ ภูริปญฺโญ,}}

\prose{\textbf{ส เตน เตเนว นโตหมสฺมิ. (13)}}

\prose{\textbf{สทฺธา จ ปีติ จ มโน สติ จา}ติ \textbf{สทฺธา}ติ ยา จ ภควนฺตํ อารพฺภ สทฺธา สทฺทหนา\csromansharedfootnote{44ccea940da1d798}{สทฺธหนา (ก)} โอกปฺปนา อภิปฺปสาโท สทฺธา สทฺธินฺทฺริยํ สทฺธาพลํ. \textbf{ปีตี}ติ ยา ภควนฺตํ อารพฺภ ปีติ ปาโมชฺชํ\csromansharedfootnote{898b8ed305ba8a74}{ปามุชฺชํ (สฺยา)} โมทนา อาโมทนา ปโมทนา หาโส ปหาโส วิตฺติ ตุฏฺฐิ โอทคฺยํ อตฺตมนตา จิตฺตสฺส. \textbf{มโน}ติ ยญฺจ ภควนฺตํ อารพฺภ จิตฺตํ มโน มานสํ หทยํ ปณฺฑรํ มโน มนายตนํ มนินฺทฺริยํ วิญฺญาณํ วิญฺญาณ\-กฺขนฺโธ ตชฺชา มโนวิญฺญาณ\-ธาตุ. \textbf{สตี}ติ ยา ภควนฺตํ อารพฺภ สติ อนุสฺสติ สมฺมาสตีติ สทฺธา จ ปีติ จ มโน สติ จ.}

\prose{\textbf{นา’เปนฺติ’เม โคตม}\-\textbf{สาสนมฺหา}ติ อิเม จตฺตาโร ธมฺมา โคตมสาสนา พุทฺธสาสนา ชินสาสนา ตถาคต\-สาสนา อรหนฺตสาสนา นาเปนฺติ น คจฺฉนฺติ น วิชหนฺติ น วินาเสนฺตีติ นาเปนฺติเม โคตม\-สาสนมฺหา.}

\prose{\textbf{ยํ ยํ ทิสํ วชติ ภูริปญฺโญ}ติ \textbf{ยํ ยํ ทิสนฺ}ติ ปุรตฺถิมํ วา ทิสํ ปจฺฉิมํ วา ทิสํ ทกฺขิณํ วา ทิสํ อุตฺตรํ วา ทิสํ วชติ คจฺฉติ กมติ อภิกฺกมติ. ภูริปญฺโญติ \textbf{ภูริปญฺโญ} มหาปญฺโญ ติกฺขปญฺโญ ปุถุปญฺโญ หาสปญฺโญ ชวนปญฺโญ นิพฺเพธิก\-ปญฺโญ. ภูริ วุจฺจติ ปถวี, ภควา ตาย ปถวิสมาย ปญฺญาย วิปุลาย วิตฺถตาย สมนฺนาคโตติ ยํ ยํ ทิสํ วชติ ภูริปญฺโญ.}

\prose{\textbf{ส เตน เตเนว นโตหมสฺมี}ติ โส เยน พุทฺโธ เตน เตเนว นโต ตนฺนินฺโน ตปฺโปโณ ตปฺปพฺภาโร ตทธิมุตฺโต ตทธิปเตยฺโยติ ส เตน เตเนว นโตหมสฺมิ. เตนาห เถโร ปิงฺคิโย\csromandash{}}

\setlength{\csromangathapagewidth}{0pt}
\setlength{\csromangathawakwidth}{0pt}
\csromangathameasuregroup{}{สทฺธา จ ปีติ จ มโน สติ จ, \\ นา’เปนฺติ’เม โคตมสาสนมฺหา. \\ ยํ ยํ ทิสํ วชติ ภูริปญฺโญ, \\ ส เตน เตเนว นโตหมสฺมีติ.}
\csromangathameasurewak{สทฺธา จ ปีติ จ มโน สติ จ, \\ นา’เปนฺติ’เม โคตมสาสนมฺหา. \\ ยํ ยํ ทิสํ วชติ ภูริปญฺโญ, \\ ส เตน เตเนว นโตหมสฺมีติ.}
\setlength{\csromangathaleftcol}{0pt}
\csromangathagroup{\csromangathastanza{\csromangathawak{\csromanfolio{220}สทฺธา จ ปีติ จ มโน สติ จ, \\ นา’เปนฺติ’เม โคตม\-สาสนมฺหา. \\ ยํ ยํ ทิสํ วชติ ภูริปญฺโญ, \\ ส เตน เตเนว นโตหมสฺมีติ.}}}

\proseitem{115}{\textbf{ชิณฺณสฺส เม ทุพฺพล}\-\textbf{ถามกสฺส,}}

\prose{\textbf{เตเนว กาโย น ปเลติ ตตฺถ.}}

\prose{\textbf{สงฺกปฺป}\-\textbf{ยนฺตาย วชามิ นิจฺจํ,}}

\prose{\textbf{มโน หิ เม พฺราหฺมณ เตน ยุตฺโต. (14)}}

\prose{\textbf{ชิณฺณสฺส เม ทุพฺพล}\-\textbf{ถามกสฺสา}ติ ชิณฺณสฺส วุฑฺฒสฺส มหลฺลกสฺส อทฺธคตสฺส วโย\-อนุปฺปตฺตสฺส. ทุพฺพล\-ถามกสฺสาติ ทุพฺพล\-ถามกสฺส อปฺปถามกสฺส ปริตฺต\-ถามกสฺสาติ ชิณฺณสฺส เม ทุพฺพล\-ถามกสฺส.}

\prose{\textbf{เตเนว กาโย น ปเลติ ตตฺถา}ติ กาโย เยน พุทฺโธ, เตน น ปเลติ น วชติ น คจฺฉติ นาติกฺกมตีติ เตเนว กาโย น ปเลติ ตตฺถ.}

\prose{\textbf{สงฺกปฺป}\-\textbf{ยนฺตาย วชามิ นิจฺจนฺ}ติ สงฺกปฺป\-คมเนน วิตกฺก\-คมเนน ญาณคมเนน ปญฺญาคมเนน พุทฺธิคมเนน วชามิ คจฺฉามิ อติกฺกมามีติ สงฺกปฺป\-ยนฺตาย วชามิ นิจฺจํ.}

\prose{\textbf{มโน หิ เม พฺราหฺมณ เตน ยุตฺโต}ติ มโนติ ยํ จิตฺตํ \textbf{มโน} มานสํ ฯเปฯ ตชฺชา มโนวิญฺญาณ\-ธาตุ. \textbf{มโน หิ เม พฺราหฺมณ เตน ยุตฺโต}ติ มโน เยน พุทฺโธ เตน ยุตฺโต ปยุตฺโต สํยุตฺโตติ มโน หิ เม พฺราหฺมณ เตน ยุตฺโต. เตนาห เถโร ปิงฺคิโย\csromandash{}}

\prose{\textbf{ชิณฺณสฺส เม ทุพฺพล}\-\textbf{ถามกสฺส,}}

\prose{\textbf{เตเนว กาโย น ปเลติ ตตฺถ.}}

\setlength{\csromangathapagewidth}{0pt}
\setlength{\csromangathawakwidth}{0pt}
\csromangathameasuregroup{\textbf{ปงฺเก สยาโน ปริผ}นฺทมาโน, \\ \textbf{อถทฺทสาสิํ สมฺพฺ}อุทฺธํ,}{\textbf{สงฺกปฺป}\textbf{ยนฺตาย วชามิ นิจฺจํ,} \\ มโน หิ เม พฺราหฺมณ เตน ยุตฺโตติ. \\ \csromangathabat{\textbf{ปงฺเก สยาโน ปริผ}นฺทมาโน,}{ทีปา ทีปํ อุปลฺลวิํ.} \\ \csromangathabat{\textbf{อถทฺทสาสิํ สมฺพฺ}อุทฺธํ,}{โอฆติณฺณมนาสวํ.}}
\csromangathameasurewak{\textbf{สงฺกปฺป}\textbf{ยนฺตาย วชามิ นิจฺจํ,} \\ มโน หิ เม พฺราหฺมณ เตน ยุตฺโตติ.}
\csromangathasetleft{\textbf{ปงฺเก สยาโน ปริผ}นฺทมาโน, \\ \textbf{อถทฺทสาสิํ สมฺพฺ}อุทฺธํ,}
\csromangathagroup{\csromangathastanza{\csromangathawak{\textbf{สงฺกปฺป}\-\textbf{ยนฺตาย วชามิ นิจฺจํ,} \\ มโน หิ เม พฺราหฺมณ เตน ยุตฺโตติ.}}\csromangathastanza{\csromangathaitembat{116}{\textbf{ปงฺเก สยาโน ปริผ}นฺทมาโน,}{ทีปา ทีปํ อุปลฺลวิํ.} \\ \csromangathacontbat{\textbf{อถทฺทสาสิํ สมฺพฺ}อุทฺธํ,}{โอฆติณฺณม\-นาสวํ.}}}

\prose{(15)}

\prose{\csromanfolio{221}\textbf{ปงฺเก สยาโน ปริผนฺท}\-\textbf{มาโน}ติ \textbf{ปงฺเก สยาโน}ติ กามปงฺเก กามกทฺทเม กามกิเลเส กามพฬิเส กามปริฬาเห กามปลิโพเธ เสมาโน สยมาโน วสมาโน อาวสมาโน ปริวสมาโน\csromansharedfootnote{7a836dafc4f53730}{อวเสมาโน ปริเสมาโน (สฺยา)}ติ ปงฺเก สยาโน. \textbf{ปริผนฺท}\-\textbf{มาโน}ติ ตณฺหา\-ผนฺทนาย ผนฺทมาโน, ทิฏฺฐิ\-ผนฺทนาย ผนฺทมาโน, กิเลส\-ผนฺทนาย ผนฺทมาโน, ปโยค\-ผนฺทนาย ผนฺทมาโน, วิปาก\-ผนฺทนาย ผนฺทมาโน, มโนทุจฺจริต\-ผนฺทนาย ผนฺทมาโน, รตฺโต ราเคน ผนฺทมาโน, ทุฏฺโฐ โทเสน ผนฺทมาโน, มูโฬฺห โมเหน ผนฺทมาโน, วินิพนฺโธ มาเนน ผนฺทมาโน, ปรามฏฺโฐ ทิฏฺฐิยา ผนฺทมาโน, วิกฺเขปคโต อุทฺธจฺเจน ผนฺทมาโน, อนิฏฺฐงฺคโต วิจิกิจฺฉาย ผนฺทมาโน, ถามคโต อนุสเยหิ ผนฺทมาโน, ลาเภน ผนฺทมาโน, อลาเภน ผนฺทมาโน, ยเสน ผนฺทมาโน, อยเสน ผนฺทมาโน, ปสํสาย ผนฺทมาโน, นินฺทาย ผนฺทมาโน, สุเขน ผนฺทมาโน, ทุกฺเขน ผนฺทมาโน, ชาติยา ผนฺทมาโน, ชราย ผนฺทมาโน, พฺยาธินา ผนฺทมาโน, มรเณน ผนฺทมาโน, โสก\-ปริเทว\-ทุกฺข\-โทมนสฺสุปายาเสหิ ผนฺทมาโน, เนรยิเกน ทุกฺเขน ผนฺทมาโน, ติรจฺฉาน\-โยนิ\-เกน ทุกฺเขน ผนฺทมาโน, เปตฺติวิสยิ\-เกน ทุกฺเขน ผนฺทมาโน, มานุสิเกน ทุกฺเขน. คพฺโภ\-กฺกนฺติ\-มูลเกน ทุกฺเขน. คพฺภ\-ฏฺฐิติ\-มูลเกน ทุกฺเขน. คพฺภวุฏฺฐาน\-มูลเกน ทุกฺเขน. ชาตสฺสู\-ปนิพนฺธ\-เกน ทุกฺเขน. ชาตสฺส ปราเธยฺยเกน ทุกฺเขน. อตฺตูปกฺกเมน ทุกฺเขน. ปรูปกฺกเมน ทุกฺเขน. สงฺขาร\-ทุกฺเขน. วิปริณาม\-ทุกฺเขน. จกฺขุโรเคน ทุกฺเขน. โสตโรเคน ทุกฺเขน. ฆานโรเคน ทุกฺเขน. ชิวฺหาโรเคน ทุกฺเขน. กายโรเคน ทุกฺเขน. สีสโรเคน ทุกฺเขน. กณฺณโรเคน ทุกฺเขน. มุขโรเคน ทุกฺเขน. ทนฺตโรเคน ทุกฺเขน. โอฏฺฐโรเคน ทุกฺเขน. กาเสน. สาเสน. ปินาเสน. qอาเหน\csromansharedfootnote{d652d2f94ee59b7a}{ทาเหน (ก) ขุ 7. 35 ปิฏฺเฐปิ.}. ชเรน. กุจฺฉิโรเคน. มุจฺฉาย. ปกฺขนฺทิกาย. สูลาย. วิสูจิกาย. กุฏฺเฐน. คณฺเฑน. กิลาเสน. โสเสน. อปมาเรน. ททฺทุยา. กณฺฑุยา. กจฺฉุยา. รขสาย. วิตจฺฉิกาย. โลหิตปิตฺเตน\csromansharedfootnote{400aa94e55fe63d7}{โลหิเตน. ปิตฺเตน (สฺยา, ก)}. มธุเมเหน. อํสาย. ปิฬกาย. ภคนฺทลาย\csromansharedfootnote{9fa0c04364fa3384}{ภคนฺทลาย (สฺยา)}. ปิตฺต\-สมุฏฺฐาเนน อาพาเธน. เสมฺห\-สมุฏฺฐาเนน อาพาเธน. วาต\-สมุฏฺฐาเนน อาพาเธน. สนฺนิปาติเกน อาพาเธน. อุตุปริณาม\-เชน อาพาเธน. วิสม\-ปริหาร\-เชน อาพาเธน. โอปกฺกมิเกน อาพาเธน. กมฺม\-วิปาก\-เชน อาพาเธน. สีเตน. อุเณฺหน.}

\prose{\csromanfolio{222}ชิฆจฺฉาย. ปิปาสาย. อุจฺจาเรน. ปสฺสาเวน. qอํสมกสวาตาตปสรีสปสมฺผสฺเสน ทุกฺเขน. มาตุมรเณน ทุกฺเขน. ปิตุมรเณน ทุกฺเขน. ปุตฺตมรเณน ทุกฺเขน. ธีตุมรเณน ทุกฺเขน. ญาติพฺยสเนน ทุกฺเขน. โภคพฺยสเนน ทุกฺเขน. โรคพฺยสเนน ทุกฺเขน. สีลพฺยสเนน ทุกฺเขน. ทิฏฺฐิพฺยเนน ทุกฺเขน ผนฺทมาโน ปริผนฺทมาโน ปเวธมาโน สมฺปเวธมาโนติ ปงฺเก สยาโน ปริผนฺทมาโน.}

\prose{\textbf{ทีปา ทีปํ อุปลฺลวินฺ}ติ สตฺถารโต สตฺถารํ ธมฺมกฺขานโต ธมฺมกฺขานํ คณโต คณํ ทิฏฺฐิยา ทิฏฺฐิํ ปฏิปทาย ปฏิปทํ มคฺคโต มคฺคํ ปลฺลวิํ อุปลฺลวิํ สมฺปลฺลวินฺติ ทีปา ทีปํ อุปลฺลวิํ.}

\prose{\textbf{อถทฺทสาสิํ สมฺพุทฺธนฺ}ติ \textbf{อถา}ติ ปทสนฺธิ ปทสํสคฺโค ปทปาริปูรี อกฺขร\-สมวาโย พฺยญฺชน\-สิลิฏฺฐตา ปทา\-นุปุพฺพตา\-เปตํ อถาติ. \textbf{อทฺทสาสินฺ}ติ อทฺทสํ อทฺทกฺขิํ อปสฺสิํ ปฏิวิชฺฌิํ. \textbf{พุทฺโธ}ติ โย โส ภควา สยมฺภู อนาจริยโก ฯเปฯ สจฺฉิกา ปญฺญตฺติ, ยทิทํ พุทฺโธติ อถทฺทสาสิํ สมฺพุทฺธํ.}

\prose{\textbf{โอฆติณฺณ}\-\textbf{มนา}\-\textbf{สวนฺ}ติ \textbf{โอฆติณฺณนฺ}ติ ภควา กาโมฆํ ติณฺโณ, ภโวฆํ ติณฺโณ, ทิฏฺโฐฆํ ติณฺโณ, อวิชฺโชฆํ ติณฺโณ, สพฺพ\-สํสารปถํ ติณฺโณ อุตฺติณฺโณ นิตฺถิณฺโณ อติกฺกนฺโต สมติกฺกนฺโต วีติวตฺโต, โส วุตฺถวาโส จิณฺณจรโณ ฯเปฯ ชาติ\-มรณ\-สํสาโร, นตฺถิ ตสฺส ปุนพฺภโวติ โอฆติณฺณํ. \textbf{อนาสวนฺ}ติ จตฺตาโร อาสวา กามาสโว ภวาสโว ทิฏฺฐาสโว อวิชฺชาสโว, เต อาสวา พุทฺธสฺส ภควโต ปหีนา อุจฺฉินฺนมูลา ตาลาวตฺถุกตา อนภาวํกตา อายติํ อนุปฺปาทธมฺมา, ตสฺมา พุทฺโธ อนาสโวติ โอฆติณฺณม\-นาสวํ. เตนาห เถโร ปิงฺคิโย\csromandash{}}

\setlength{\csromangathapagewidth}{0pt}
\setlength{\csromangathawakwidth}{0pt}
\csromangathameasuregroup{ปงฺเก สยาโน ปริผนฺทมาโน, \\ อถทฺทสาสิํ สมฺพุทฺธํ,}{\csromangathabat{ปงฺเก สยาโน ปริผนฺทมาโน,}{ทีปา ทีปํ อุปลฺลวิํ.} \\ \csromangathabat{อถทฺทสาสิํ สมฺพุทฺธํ,}{โอฆติณฺณมนาสวนฺติ.}}
\csromangathasetleft{ปงฺเก สยาโน ปริผนฺทมาโน, \\ อถทฺทสาสิํ สมฺพุทฺธํ,}
\csromangathagroup{\csromangathastanza{\csromangathabat{ปงฺเก สยาโน ปริผนฺทมาโน,}{ทีปา ทีปํ อุปลฺลวิํ.} \\ \csromangathabat{อถทฺทสาสิํ สมฺพุทฺธํ,}{โอฆติณฺณ\-มนา\-สวนฺติ.}}}

\proseitem{117}{\textbf{ยถา อหู วกฺกลิ มุตฺตสทฺโธ,}}

\prose{\textbf{ภทฺราวุโธ อาฬวิโคตโม จ.}}

\prose{\textbf{เอวเมว ตฺวมฺปิ ปมุญฺจสฺสุ สทฺธํ,}}

\prose{\textbf{คมิสฺสสิ ตฺวํ ปิงฺคิย มจฺจุเธยฺยสฺส ปารํ.} (16)}

\prose{\csromanfolio{223}\textbf{ยถา อหู วกฺกลิ มุตฺตสทฺโธ, ภทฺราวุโธ อาฬวิโคตโม จา}ติ ยถา วกฺกลิตฺเถโร\csromansharedfootnote{34197337ac3bb451}{วกฺกลิ (สฺยา)} สทฺโธ สทฺธาครุโก สทฺธา\-ปุพฺพงฺคโม สทฺธาธิมุตฺโต สทฺธา\-ธิปเตยฺโย อรหตฺตปฺปตฺโต, ยถา ภทฺราวุโธ เถโร สทฺโธ สทฺธาครุโก สทฺธา\-ปุพฺพงฺคโม สทฺธาธิมุตฺโต สทฺธา\-ธิปเตยฺโย อรหตฺตปฺปตฺโต, ยถา อาฬวิโคตโม เถโร สทฺโธ สทฺธาครุโก สทฺธา\-ปุพฺพงฺคโม สทฺธาธิมุตฺโต สทฺธา\-ธิปเตยฺโย อรหตฺต\-ปฺปตฺโตติ ยถา อหู วกฺกลิ มุตฺตสทฺโธ, ภทฺราวุโธ อาฬวิโคตโม จ.}

\prose{\textbf{เอวเมว ตฺวมฺปิ ปมุญฺจสฺสุ สทฺธนฺ}ติ เอวเมว ตฺวํ สทฺธํ มุญฺจสฺสุ ปมุญฺจสฺสุ สมฺปมุญฺจสฺสุ อธิมุญฺจสฺสุ โอกปฺเปหิ. “สพฺเพ สงฺขารา อนิจฺจา”ติ สทฺธํ มุญฺจสฺสุ ปมุญฺจสฺสุ สมฺปมุญฺจสฺสุ อธิมุญฺจสฺสุ โอกปฺเปหิ. “สพฺเพ สงฺขารา ทุกฺขา”ติ. “สพฺเพ ธมฺมา อนตฺตา”ติ สทฺธํ มุญฺจสฺสุ ปมุญฺจสฺสุ สมฺปมุญฺจสฺสุ อธิมุญฺจสฺสุ โอกปฺเปหิ ฯเปฯ. “ยํ กิญฺจิ สมุทย\-ธมฺมํ, สพฺพํ ตํ นิโรธธมฺมนฺ”ติ สทฺธํ มุญฺจสฺสุ ปมุญฺจสฺสุ สมฺปมุญฺจสฺสุ อธิมุญฺจสฺสุ โอกปฺเปหีติ เอวเมว ตฺวมฺปิ ปมุญฺจสฺสุ สทฺธํ.}

\prose{\textbf{คมิสฺสสิ ตฺวํ ปิงฺคิย มจฺจุเธยฺยสฺส ปารนฺ}ติ มจฺจุเธยฺยํ วุจฺจนฺติ กิเลสา จ ขนฺธา จ อภิสงฺขารา จ. มจฺจุเธยฺยสฺส ปารํ วุจฺจติ อมตํ นิพฺพานํ, โย โส สพฺพ\-สงฺขาร\-สมโถ สพฺพู\-ปธิ\-ปฏินิสฺสคฺโค ตณฺหกฺขโย วิราโค นิโรโธ นิพฺพานํ. \textbf{คมิสฺสสิ ตฺวํ ปิงฺคิย มจฺจุเธยฺยสฺส ปารนฺ}ติ ตฺวํ ปารํ คมิสฺสสิ, ปารํ อธิคมิสฺสสิ, ปารํ ผสฺสิสฺสสิ, ปารํ สจฺฉิกริสฺสสีติ คมิสฺสสิ ตฺวํ ปิงฺคิย มจฺจุเธยฺยสฺส ปารํ. เตนาห ภควา\csromandash{}}

\prose{\textbf{ยถา อหู วกฺกลิ มุตฺตสทฺ}โธ,}

\prose{ภทฺราวุโธ อาฬวิโคตโม จ.}

\prose{เอวเมว ตฺวมฺปิ ปมุญฺจสฺสุ สทฺธํ,}

\prose{คมิสฺสสิ ตฺวํ ปิงฺคิย มจฺจุเธยฺยสฺส ปารนฺติ.}

\proseitem{118}{\textbf{เอส ภิยฺโย ปสีทามิ, สุตฺวาน มุนิโน วโจ.}}

\prose{\textbf{วิวฏฺฏจฺฉโท}\csromansharedfootnote{bbfae0fab4a41a60}{วิวฏจฺฉทโน (ก) สทฺทนีติปทมาลา โอโลเกตพฺพา.}\textbf{ สมฺพุทฺโธ, อขิโล ปฏิภานวา}\csromansharedfootnote{32f2191d413669c3}{ปฏิภาณวา (สฺยา)}\textbf{.} (17)}

\prose{\textbf{เอส ภิยฺโย ปสีทามี}ติ เอส ภิยฺโย ปสีทามิ, ภิยฺโย ภิยฺโย สทฺทหามิ, ภิยฺโย ภิยฺโย โอกปฺเปมิ, ภิยฺโย ภิยฺโย อธิมุจฺจามิ, “สพฺเพ สงฺขารา อนิจฺจา”ติ ภิยฺโย ภิยฺโย ปสีทามิ, ภิยฺโย ภิยฺโย \csromanfolio{224}สทฺทหามิ, ภิยฺโย ภิยฺโย โอกปฺเปมิ, ภิยฺโย ภิยฺโย อธิมุจฺจามิ, “สพฺเพ สงฺขารา ทุกฺขา”ติ ภิยฺโย ภิยฺโย ปสีทามิ ฯเปฯ “สพฺเพ ธมฺมา อนตฺตา”ติ ภิยฺโย ภิยฺโย ปสีทามิ ฯเปฯ. “ยํ กิญฺจิ สมุทย\-ธมฺมํ, สพฺพํ ตํ นิโรธธมฺมนฺ”ติ ภิยฺโย ภิยฺโย ปสีทามิ, ภิยฺโย ภิยฺโย สทฺทหามิ, ภิยฺโย ภิยฺโย โอกปฺเปมิ, ภิยฺโย ภิยฺโย อธิมุจฺจามีติ เอส ภิยฺโย ปสีทามิ.}

\prose{\textbf{สุตฺวาน มุนิโน วโจ}ติ \textbf{มุนี}ติ โมนํ วุจฺจติ ญาณํ ฯเปฯ สงฺค\-ชาลม\-ติจฺจ โส มุนิ. \textbf{สุตฺวาน มุนิโน วโจ}ติ ตุยฺหํ วจนํ พฺยปฺปถํ เทสนํ อนุสาสนํ อนุสิฏฺฐํ สุตฺวาน อุคฺคเหตฺวาน อุปธารยิตฺวาน อุปลกฺขยิตฺวานาติ สุตฺวาน มุนิโน วโจ.}

\prose{\textbf{วิวฏฺฏจฺฉโท สมฺพุทฺโธ}ติ \textbf{ฉทนนฺ}ติ ปญฺจ ฉทนานิ ตณฺหาฉทนํ ทิฏฺฐิฉทนํ กิเลสฉทนํ ทุจฺจริต\-ฉทนํ อวิชฺชาฉทนํ, ตานิ ฉทนานิ พุทฺธสฺส ภควโต วิวฏานิ วิทฺธํสิตานิ สมุคฺฆาฏิตานิ ปหีนานิ สมุจฺฉินฺนานิ วูปสนฺตานิ ปฏิปฺปสฺสทฺธานิ อภพฺพุ\-ปฺปตฺติกานิ ญาณคฺคินา ทฑฺฒานิ, ตสฺมา พุทฺโธ วิวฏฺฏจฺฉโท. \textbf{พุทฺโธ}ติ โย โส ภควา ฯเปฯ สจฺฉิกา ปญฺญตฺติ, ยทิทํ พุทฺโธติ วิวฏฺฏจฺฉโท สมฺพุทฺโธ.}

\prose{\textbf{อขิโล ปฏิภานวา}ติ \textbf{อขิโล}ติ ราโค ขิโล, โทโส ขิโล, โมโห ขิโล, โกโธ ขิโล, อุปนาโห ฯเปฯ สพฺพา\-กุสลา\-ภิสงฺขารา ขิลา, เต ขิลา พุทฺธสฺส ภควโต ปหีนา อุจฺฉินฺนมูลา ตาลาวตฺถุกตา อนภาวํกตา อายติํ อนุปฺปาทธมฺมา, ตสฺมา พุทฺโธ อขิโล.}

\prose{\textbf{ปฏิภานวา}ติ ตโย ปฏิภานวนฺโต ปริยตฺติ ปฏิภานวา ปริปุจฺฉา\-ปฏิภานวา อธิคม\-ปฏิภานวา. กตโม ปริยตฺติ\-ปฏิภานวา, อิเธกจฺจสฺส พุทฺธวจนํ\csromansharedfootnote{3793d616c16ab90b}{ปกติยา (ก)} ปริยาปุตํ\csromansharedfootnote{3abab0119a198666}{ปริยาปุฏํ (สฺยา, ก)} โหติ สุตฺตํ เคยฺยํ เวยฺยากรณํ คาถา อุทานํ อิติวุตฺตกํ ชาตกํ อพฺภุตธมฺมํ เวทลฺลํ, ตสฺส ปริยตฺติํ นิสฺสาย ปฏิภาติ\csromansharedfootnote{6200954c17a7c9b3}{ปฏิภายติ (ก)}, อยํ ปริยตฺติ\-ปฏิภานวา.}

\prose{กตโม ปริปุจฺฉา\-ปฏิภานวา, อิเธกจฺโจ ปริปุจฺฉิตา โหติ อตฺเถ จ ญาเย จ ลกฺขเณ จ การเณ จ ฐานาฐาเน จ, ตสฺส ปริปุจฺฉํ นิสฺสาย ปฏิภาติ, อยํ ปริปุจฺฉา\-ปฏิภานวา.}

\prose{\csromanfolio{225}กตโม อธิคม\-ปฏิภานวา, อิเธกจฺจสฺส อธิคตา โหนฺติ จตฺตาโร สติปฏฺฐานา จตฺตาโร สมฺมปฺปธานา จตฺตาโร อิทฺธิปาทา ปญฺจินฺทฺริยานิ ปญฺจ พลานิ สตฺต โพชฺฌงฺคา อริโย อฏฺฐงฺคิโก มคฺโค จตฺตาโร อริยมคฺคา จตฺตาริ สามญฺญผลานิ จตสฺโส ปฏิสมฺภิทาโย ฉ อภิญฺญาโย, ตสฺส อตฺโถ ญาโต ธมฺโม ญาโต นิรุตฺติ ญาตา, อตฺเถ ญาเต อตฺโถ ปฏิภาติ, ธมฺเม ญาเต ธมฺโม ปฏิภาติ, นิรุตฺติยา ญาตาย นิรุตฺติ ปฏิภาติ, อิเมสุ ตีสุ ญาเณสุ ญาณํ ปฏิภาน\-ปฏิสมฺภิทา. ภควา อิมาย ปฏิภาน\-ปฏิสมฺภิทาย อุเปโต สมุเปโต อุปาคโต สมุปาคโต อุปปนฺโน สมุปปนฺโน สมนฺนาคโต, ตสฺมา พุทฺโธ ปฏิภานวา. ยสฺส ปริยตฺติ นตฺถิ, ปริปุจฺฉา นตฺถิ, อธิคโม นตฺถิ, กิํ ตสฺส ปฏิภายิสฺสตีติ อขิโล ปฏิภานวา. เตนาห เถโร ปิงฺคิโย\csromandash{}}

\setlength{\csromangathapagewidth}{0pt}
\setlength{\csromangathawakwidth}{0pt}
\csromangathameasuregroup{เอส ภิยฺโย ปสีทามิ, \\ วิวฏฺฏจฺฉโท สมฺพุทฺโธ,}{\csromangathabat{เอส ภิยฺโย ปสีทามิ,}{สุตฺวาน มุนิโน วโจ\textbf{.}} \\ \csromangathabat{วิวฏฺฏจฺฉโท สมฺพุทฺโธ,}{อขิโล ปฏิภานวาติ\textbf{.}}}
\csromangathasetleft{เอส ภิยฺโย ปสีทามิ, \\ วิวฏฺฏจฺฉโท สมฺพุทฺโธ,}
\csromangathagroup{\csromangathastanza{\csromangathabat{เอส ภิยฺโย ปสีทามิ,}{สุตฺวาน มุนิโน วโจ\textbf{.}} \\ \csromangathabat{วิวฏฺฏจฺฉโท สมฺพุทฺโธ,}{อขิโล ปฏิภานวาติ\textbf{.}}}}

\proseitem{119}{อธิเทเว\csromansharedfootnote{09a71fa0cfbf037f}{อติเทเว (ก)} \textbf{อภิญฺญาย, สพฺพํ เวทิ ปโรปรํ}\csromansharedfootnote{6f971e4ce60bd24c}{ปโรวรํ (สีฏฺฐ)}\textbf{.}}

\prose{\textbf{ปญฺหานนฺตกโร สตฺถา, กงฺขีนํ ปฏิชานตํ. (18)}}

\prose{\textbf{อธิเทเว อภิญฺญายา}ติ เทวาติ ตโย \textbf{เทวา} สมฺมุติเทวา\csromansharedfootnote{0494d7ee35a8515b}{สมฺมติเทวา (สฺยา) 85 ปิฏฺเฐปิ.} อุปปตฺติเทวา วิสุทฺธิเทวา. กตเม สมฺมุติเทวา, สมฺมุติเทวา วุจฺจนฺติ ราชาโน\csromansharedfootnote{11451ffbb9044467}{กตเม สมฺมติเทวา ราชาโน (สฺยา) เอวมุปริปิ.} จ ราชกุมารา จ เทวิโย จ, อิเม วุจฺจนฺติ สมฺมุติเทวา. กตเม อุปปตฺติเทวา, อุปปตฺติเทวา วุจฺจนฺติ จาตุมหาราชิกา เทวา ตาวติํสา เทวา ฯเปฯ พฺรหฺมกายิกา เทวา, เย จ เทวา ตทุตฺตริ, อิเม วุจฺจนฺติ อุปปตฺติเทวา. กตเม วิสุทฺธิเทวา, วิสุทฺธิเทวา วุจฺจนฺติ ตถาคตา ตถาคต\-สาวกา อรหนฺโต ขีณาสวา, เย จ ปจฺเจก\-สมฺพุทฺธา, อิเม วุจฺจนฺติ วิสุทฺธิเทวา. ภควา สมฺมุติเทเว อธิเทวาติ อภิญฺญาย, อุปปตฺติเทเว อธิเทวาติ อภิญฺญาย, วิสุทฺธิเทเว อธิเทวาติ อภิญฺญาย ชานิตฺวา ตุลยิตฺวา ตีรยิตฺวา วิภาวยิตฺวา วิภูตํ กตฺวาติ อธิเทเว อภิญฺญาย.}

\prose{\csromanfolio{226}\textbf{สพฺพํ เวทิ ปโรปรนฺ}ติ ภควา อตฺตโน จ ปเรสญฺจ อธิเทวกเร ธมฺเม เวทิ อญฺญาสิ อผสฺสิ ปฏิวิชฺฌิ. กตเม อตฺตโน อธิเทวกรา ธมฺมา, สมฺมาปฏิปทา อนุโลม\-ปฏิปทา อปจฺจนีกปฏิปทา อนฺวตฺถ\-ปฏิปทา ธมฺมานุธมฺม\-ปฏิปทา สีเลสุ ปริปูรการิตา อินฺทฺริเยสุ คุตฺตทฺวารตา โภชเน มตฺตญฺญุตา ชาคริยานุโยโค สติสมฺปชญฺญํ จตฺตาโร สติปฏฺฐานา ฯเปฯ อริโย อฏฺฐงฺคิโก มคฺโค, อิเม วุจฺจนฺติ อตฺตโน อธิเทวกรา ธมฺมา.}

\prose{กตเม ปเรสํ อธิเทวกรา ธมฺมา, สมฺมาปฏิปทา ฯเปฯ อริโย อฏฺฐงฺคิโก มคฺโค, อิเม วุจฺจนฺติ ปเรสํ อธิเทวกรา ธมฺมา. เอวํ ภควา อตฺตโน จ ปเรสญฺจ อธิเทวกเร ธมฺเม เวทิ อญฺญาสิ อผสฺสิ ปฏิวิชฺฌีติ สพฺพํ เวทิ ปโรปรํ.}

\prose{\textbf{ปญฺหานนฺตกโร สตฺถา}ติ ภควา ปาราย\-นิก\-ปญฺหานํ อนฺตกโร ปริยนฺตกโร ปริจฺเฉทกโร ปริวฏุมกโร. สภิย\-ปญฺหานํ\csromansharedfootnote{8cb387d00e0bd785}{ปริสปญฺหานํ (สฺยา), ปิงฺคิยปญฺหานํ (ก)} อนฺตกโร ปริยนฺตกโร ปริจฺเฉทกโร ปริวฏุมกโร. สกฺกปญฺหานํ. สุยามปญฺหานํ. ภิกฺขุ\-ปญฺหานํ. ภิกฺขุนี\-ปญฺหานํ. อุปาสก\-ปญฺหานํ. อุปาสิกา\-ปญฺหานํ. ราชปญฺหานํ. ขตฺติย\-ปญฺหานํ. พฺราหฺมณ\-ปญฺหานํ. เวสฺสปญฺหานํ. สุทฺทปญฺหานํ. เทวปญฺหานํ. พฺรหฺมปญฺหานํ อนฺตกโร ปริยนฺตกโร ปริจฺเฉทกโร ปริวฏุม\-กโรติ ปญฺหานนฺตกโร. \textbf{สตฺถา}ติ ภควา สตฺถวาโห. ยถา สตฺถวาโห สตฺเถ กนฺตารํ ตาเรติ, โจรกนฺตารํ ตาเรติ, วาฬกนฺตารํ ตาเรติ, ทุพฺภิกฺข\-กนฺตารํ ตาเรติ, นิรุทก\-กนฺตารํ ตาเรติ อุตฺตาเรติ นิตฺตาเรติ\csromansharedfootnote{926d9e24bcc9fc43}{นิตฺถาเรติ (สฺยา, ก)} ปตาเรติ, เขมนฺตภูมิํ สมฺปาเปติ. เอวเมว ภควา สตฺถวาโห สตฺเต กนฺตารํ ตาเรติ. ชาติกนฺตารํ ตาเรติ. ชรากนฺตารํ. พฺยาธิกนฺตารํ. มรณกนฺตารํ. โสกปริเทวทุกฺขโทมนสฺสุปายาส\-กนฺตารํ ตาเรติ, ราคกนฺตารํ ตาเรติ, โทสกนฺตารํ. โมหกนฺตารํ. มานกนฺตารํ. ทิฏฺฐิกนฺตารํ. กิเลสกนฺตารํ. ทุจฺจริต\-กนฺตารํ ตาเรติ, ราคคหนํ ตาเรติ, โทสคหนํ. โมหคหนํ. ทิฏฺฐิคหนํ. กิเลสคหนํ. ทุจฺจริต\-คหนํ ตาเรติ อุตฺตาเรติ นิตฺตาเรติ ปตาเรติ, เขมนฺตํ อมตํ นิพฺพานํ สมฺปาเปตีติ เอวมฺปิ ภควา สตฺถวาโห.}

\prose{อถ วา ภควา เนตา วิเนตา อนุเนตา ปญฺญาเปตา นิชฺฌาเปตา เปกฺขตา ปสาเทตาติ เอวํ ภควา สตฺถวาโห.  อถ วา ภควา \csromanfolio{227}อนุปฺปนฺนสฺส มคฺคสฺส อุปฺปาเทตา, อสญฺชาตสฺส มคฺคสฺส สญฺชเนตา, อนกฺขาตสฺส มคฺคสฺส อกฺขาตา, มคฺคญฺญู มคฺควิทู มคฺคโกวิโท มคฺคานุคา จ ปน เอตรหิ สาวกา วิหรนฺติ ปจฺฉา สมนฺนาคตาติ เอวมฺปิ ภควา สตฺถวาโหติ ปญฺหานนฺตกโร สตฺถา.}

\prose{\textbf{กงฺขีนํ ปฏิชานตนฺ}ติ สกงฺขา อาคนฺตฺวา นิกฺกงฺขา สมฺปชฺชนฺติ, สลฺเลขา อาคนฺตฺวา นิลฺเลขา สมฺปชฺชนฺติ, สเทฺวฬฺหกา อาคนฺตฺวา นิเทฺวฬฺหกา สมฺปชฺชนฺติ, สวิจิกิจฺฉา อาคนฺตฺวา นิพฺพิจิกิจฺฉา สมฺปชฺชนฺติ, สราคา อาคนฺตฺวา วีตราคา สมฺปชฺชนฺติ, สโทสา อาคนฺตฺวา วีตโทสา สมฺปชฺชนฺติ, สโมหา อาคนฺตฺวา วีตโมหา สมฺปชฺชนฺติ, สกิเลสา อาคนฺตฺวา นิกฺกิเลสา สมฺปชฺชนฺตีติ กงฺขีนํ ปฏิชานตํ. เตนาห เถโร ปิงฺคิโย\csromandash{}}

\setlength{\csromangathapagewidth}{0pt}
\setlength{\csromangathawakwidth}{0pt}
\csromangathameasuregroup{อธิเทเว อภิญฺญาย, \\ ปญฺหานนฺตกโร สตฺถา,}{\csromangathabat{อธิเทเว อภิญฺญาย,}{สพฺพํ เวทิ ปโรปรํ.} \\ \csromangathabat{ปญฺหานนฺตกโร สตฺถา,}{กงฺขีนํ ปฏิชานตนฺติ.}}
\csromangathasetleft{อธิเทเว อภิญฺญาย, \\ ปญฺหานนฺตกโร สตฺถา,}
\csromangathagroup{\csromangathastanza{\csromangathabat{อธิเทเว อภิญฺญาย,}{สพฺพํ เวทิ ปโรปรํ.} \\ \csromangathabat{ปญฺหานนฺตกโร สตฺถา,}{กงฺขีนํ ปฏิชานตนฺติ.}}}

\proseitem{120}{\textbf{อสํหีรํ อสํกุปฺปํ,}}

\prose{\textbf{ยสฺส นตฺถิ อุปมา กฺวจิ.}}

\prose{\textbf{อทฺธา คมิสฺสามิ น เมตฺถ กงฺขา,}}

\prose{\textbf{เอวํ มํ ธาเรหิ อธิมุตฺต}\-\textbf{จิตฺตํ.} (19)}

\prose{\textbf{อสํหีรํ อสํกุปฺปนฺ}ติ อสํหีรํ วุจฺจติ อมตํ นิพฺพานํ. โย โส สพฺพ\-สงฺขาร\-สมโถ สพฺพู\-ปธิ\-ปฏินิสฺสคฺโค ตณฺหกฺขโย วิราโค นิโรโธ นิพฺพานํ. \textbf{อสํหีรนฺ}ติ ราเคน โทเสน โมเหน โกเธน อุปนาเหน มกฺเขน ปฬาเสน อิสฺสาย มจฺฉริเยน มายาย สาเฐยฺเยน ถมฺเภน สารมฺเภน มาเนน อติมาเนน มเทน ปมาเทน สพฺพกิเลเสหิ สพฺพ\-ทุจฺจริเตหิ สพฺพ\-ปริฬาเหหิ สพฺพาสเวหิ สพฺพทรเถหิ สพฺพสนฺตาเปหิ สพฺพากุสลา\-ภิสงฺขาเรหิ อสํหาริยํ นิพฺพานํ นิจฺจํ ธุวํ สสฺสตํ อวิปริณาม\-ธมฺมนฺติ อสํหีรํ.}

\prose{\textbf{อสํกุปฺปนฺ}ติ อสํกุปฺปํ วุจฺจติ อมตํ นิพฺพานํ. โย โส สพฺพ\-สงฺขาร\-สมโถ ฯเปฯ นิโรโธ นิพฺพานํ. นิพฺพานสฺส\csromansharedfootnote{1e4b297b6c15dcb5}{ยสฺส (สฺยา)} น อุปฺปาโท ปญฺญายติ, วโย นตฺถิ, น ตสฺส อญฺญถตฺตํ\csromansharedfootnote{7943ecab8b581184}{ตสฺส อญฺญทตฺถุ (สฺยา)} ปญฺญายติ, นิพฺพานํ นิจฺจํ ธุวํ สสฺสตํ อวิปริณาม\-ธมฺมนฺติ อสํหีรํ อสํกุปฺปํ.}

\prose{\csromanfolio{228}\textbf{ยสฺส นตฺถิ อุปมา กฺวจี}ติ \textbf{ยสฺสา}ติ นิพฺพานสฺส. \textbf{นตฺถิ อุปมา}ติ อุปมา นตฺถิ, อุปนิธา นตฺถิ, สทิสํ นตฺถิ, ปฏิภาโค นตฺถิ น สติ น สํวิชฺชติ นุปลพฺภติ. \textbf{กฺวจี}ติ กฺวจิ กิมฺหิจิ กตฺถจิ อชฺฌตฺตํ วา พหิทฺธา วา อชฺฌตฺต\-พหิทฺธา วาติ ยสฺส นตฺถิ อุปมา กฺวจิ.}

\prose{\textbf{อทฺธา คมิสฺสามิ น เมตฺถ กงฺขา}ติ \textbf{อทฺธา}ติ เอกํสวจนํ นิสฺสํสยวจนํ นิกฺกงฺข\-วจนํ อเทฺวชฺฌวจนํ อเทฺวฬฺหก\-วจนํ นิโยควจนํ อปณฺณก\-วจนํ อวิรทฺธ\-วจนํ อวตฺถาปน\-วจนเมตํ อทฺธาติ. \textbf{คมิสฺสามี}ติ คมิสฺสามิ อธิคมิสฺสามิ ผสฺสิสฺสามิ สจฺฉิกริสฺสามีติ อทฺธา คมิสฺสามิ. \textbf{น เมตฺถ กงฺขา}ติ \textbf{เอตฺถา}ติ นิพฺพาเน กงฺขา นตฺถิ, วิจิกิจฺฉา นตฺถิ, เทฺวฬฺหกํ นตฺถิ, สํสโย นตฺถิ น อตฺถิ น สํวิชฺชติ นุปลพฺภติ, ปหีโน สมุจฺฉินฺโน วูปสนฺโต ปฏิปฺปสฺสทฺโธ อภพฺพุ\-ปฺปตฺติโก ญาณคฺคินา ทฑฺโฒติ อทฺธา คมิสฺสามิ น เมตฺถ กงฺขา.}

\prose{\textbf{เอวํ มํ ธาเรหิ อธิมุตฺต}\-\textbf{จิตฺตนฺ}ติ \textbf{เอวํ มํ ธาเรหี}ติ เอวํ มํ อุปลกฺเขหิ. \textbf{อธิมุตฺต}\-\textbf{จิตฺตนฺ}ติ นิพฺพานนินฺนํ นิพฺพานโปณํ นิพฺพาน\-ปพฺภารํ นิพฺพานา\-ธิมุตฺตนฺติ เอวํ มํ ธาเรหิ อธิมุตฺต\-จิตฺตนฺติ. เตนาห เถโร ปิงฺคิโย\csromandash{}}

\setlength{\csromangathapagewidth}{0pt}
\setlength{\csromangathawakwidth}{0pt}
\csromangathameasuregroup{}{อสํหีรํ อสํกุปฺปํ, \\ ยสฺส นตฺถิ อุปมา กฺวจิ.}
\csromangathameasurewak{อสํหีรํ อสํกุปฺปํ, \\ ยสฺส นตฺถิ อุปมา กฺวจิ.}
\setlength{\csromangathaleftcol}{0pt}
\csromangathagroup{\csromangathastanza{\csromangathawak{อสํหีรํ อสํกุปฺปํ, \\ ยสฺส นตฺถิ อุปมา กฺวจิ.}}}

\vspace{\tipitakaparskip}
\csromancenter{อทฺธา คมิสฺสามิ น เมตฺถ กงฺขา,}
\csromancenter{เอวํ มํ ธาเรหิ อธิมุตฺต\-จิตฺตนฺติ.}
\csromancenter{ปารายนานุคีติคาถา\-นิทฺเทโส อฏฺฐารสโม.}
\nitthitammid{\textbf{ปารายนวคฺโค สมตฺโต.}}
\csromanmatikaanchor{mtk.40}
\csromantocmarkref{subsection}{19. ขคฺควิสาณสุตฺตนิทฺเทส}{mtk.40}
\bookmark[dest={mtk.40},level=2]{19. ขคฺควิสาณสุตฺตนิทฺเทส}
\csromanheader{2}{19. \textbf{ขคฺควิสาณสุตฺต}\-\textbf{นิทฺเทส}}
\chapterhead{\textbf{ปฐมวคฺค}}
\setlength{\csromangathapagewidth}{0pt}
\setlength{\csromangathawakwidth}{0pt}
\csromangathameasuregroup{}{\textbf{สพฺเพสุ ภูเตสุ นิธาย ทณฺฑํ,} \\ \textbf{อวิเหฐยํ อญฺญตรมฺปิ เตสํ.} \\ \textbf{น ปุตฺตมิจฺเฉยฺย กุโต สหายํ,} \\ \textbf{เอโก จเร ขคฺควิสาณ}\textbf{กปฺโป.}}
\csromangathameasurewak{\textbf{สพฺเพสุ ภูเตสุ นิธาย ทณฺฑํ,} \\ \textbf{อวิเหฐยํ อญฺญตรมฺปิ เตสํ.} \\ \textbf{น ปุตฺตมิจฺเฉยฺย กุโต สหายํ,} \\ \textbf{เอโก จเร ขคฺควิสาณ}\textbf{กปฺโป.}}
\setlength{\csromangathaleftcol}{0pt}
\csromangathagroup{\csromangathastanza{\csromangathawak{\csromanfolio{229}\csromangathaitemline{211}{\textbf{สพฺเพสุ ภูเตสุ นิธาย ทณฺฑํ,}} \\ \csromangathacontline{\textbf{อวิเหฐยํ อญฺญตรมฺปิ เตสํ.}} \\ \csromangathacontline{\textbf{น ปุตฺตมิจฺเฉยฺย กุโต สหายํ,}} \\ \csromangathacontline{\textbf{เอโก จเร ขคฺควิสาณ}\-\textbf{กปฺโป.}}}}}

\prose{\textbf{สพฺเพสุ ภูเตสุ นิธาย ทณฺฑนฺติ สพฺเพสู}ติ สพฺเพน สพฺพํ สพฺพถา สพฺพํ อเสสํ นิสฺเสสํ ปริยาทิย\-น\-วจนเมตํ สพฺเพสูติ. \textbf{ภูเตสู}ติ ภูตา วุจฺจนฺติ ตสา จ ถาวรา จ. ตสาติ เยสํ ตสิตตณฺหา อปฺปหีนา, เยสญฺจ ภยเภรวา อปฺปหีนา. กิํการณา วุจฺจนฺติ ตสา, เต ตสนฺติ อุตฺตสนฺติ ปริตสนฺติ ภายนฺติ สนฺตาสํ อาปชฺชนฺติ, ตํการณา วุจฺจนฺติ ตสา. ถาวราติ เยสํ ตสิตตณฺหา ปหีนา, เยสญฺจ ภยเภรวา ปหีนา. กิํการณา วุจฺจนฺติ ถาวรา, เต น ตสนฺติ น อุตฺตสนฺติ น ปริตสนฺติ น ภายนฺติ น สนฺตาสํ อาปชฺชนฺติ, ตํการณา วุจฺจนฺติ ถาวรา. \textbf{ทณฺฑนฺ}ติ ตโย ทณฺฑา กายทณฺโฑ วจีทณฺโฑ มโนทณฺโฑ. ติวิธํ กายทุจฺจริตํ กายทณฺโฑ, จตุพฺพิธํ วจีทุจฺจริตํ วจีทณฺโฑ, ติวิธํ มโนทุจฺจริตํ มโนทณฺโฑ. \textbf{สพฺเพสุ ภูเตสุ นิธาย ทณฺฑนฺ}ติ สพฺเพสุ ภูเตสุ ทณฺฑํ นิธาย นิทหิตฺวา โอโรปยิตฺวา สโมโรปยิตฺวา นิกฺขิปิตฺวา ปฏิปฺปสฺสมฺภิตฺวาติ สพฺเพสุ ภูเตสุ นิธาย ทณฺฑํ.}

\prose{\textbf{อวิเหฐยํ อญฺญตรมฺปิ เตสนฺ}ติ เอกเมกมฺปิ สตฺตํ ปาณินา วา เลฑฺฑุนา วา ทณฺเฑน วา สตฺเถน วา อนฺทุยา\csromansharedfootnote{627ed5935345d133}{อรุยา (สฺยา), อทฺทุยา (ก)} วา รชฺชุยา วา อวิเหฐยนฺโต, สพฺเพปิ สตฺเต ปาณินา วา เลฑฺฑุนา วา ทณฺเฑน วา สตฺเถน วา อนฺทุยา วา รชฺชุยา วา อวิเหฐยนฺโตติ อวิเหฐยํ อญฺญตรมฺปิ เตสํ.}

\prose{\textbf{น ปุตฺตมิจฺเฉยฺย กุโต สหายนฺ}ติ \textbf{นา}ติ ปฏิกฺเขโป. \textbf{ปุตฺตา}ติ จตฺตาโร ปุตฺตา อตฺรโช ปุตฺโต เขตฺตโช ปุตฺโต ทินฺนโก ปุตฺโต อนฺเตวาสิโก ปุตฺโต. \textbf{สหายนฺ}ติ สหายา วุจฺจนฺติ เยหิ สห อาคมนํ ผาสุ, คมนํ ผาสุ, คมนาคมนํ ผาสุ, ฐานํ ผาสุ, นิสชฺชนํ ผาสุ, สยนํ\csromansharedfootnote{5d9d618f1f0b147d}{นิปชฺชนํ (สฺยา)} ผาสุ, อาลปนํ ผาสุ, สลฺลปนํ ผาสุ, อุลฺลปนํ ผาสุ, สมุลฺลปนํ \csromanfolio{230}ผาสุ. \textbf{น ปุตฺตมิจฺเฉยฺย กุโต สหายนฺ}ติ ปุตฺตมฺปิ น อิจฺเฉยฺย น สาทิเยยฺย น ปตฺถเยยฺย น ปิหเยยฺย นาภิชปฺเปยฺย, กุโต มิตฺตํ วา สนฺทิฏฺฐํ วา สมฺภตฺตํ วา สหายํ วา อิจฺเฉยฺย\csromansharedfootnote{637f5334ef52017e}{อิจฺฉิสฺสติ (สฺยา) เอวมีทิเสสุ ปเทสุ อนาคตวิภตฺติยา.} สาทิเยยฺย ปตฺถเยยฺย ปิหเยยฺย อภิชปฺเปยฺยาติ น ปุตฺตมิจฺเฉยฺย กุโต สหายํ.}

\prose{\textbf{เอโก จเร ขคฺควิสาณ}\-\textbf{กปฺโป}ติ \textbf{เอโก}ติ โส ปจฺเจก\-สมฺพุทฺโธ ปพฺพชฺชา\-สงฺขาเตน เอโก, อทุติยฏฺเฐน เอโก, ตณฺหาย ปหานฏฺเฐน\csromansharedfootnote{d86e6acc2a331330}{ตณฺหาปหานฏฺเฐน (สฺยา) ขุ 7. 361 ปิฏฺเฐ.} เอโก, เอกนฺตวีตราโคติ เอโก, เอกนฺตวีตโทโสติ เอโก, เอกนฺตวีตโมโหติ เอโก, เอกนฺตนิกฺกิเลโสติ เอโก, เอกายนมคฺคํ คโตติ เอโก, เอโก อนุตฺตรํ ปจฺเจก\-สมฺโพธิํ อภิสมฺพุทฺโธติ เอโก.}

\prose{กถํ โส ปจฺเจก\-สมฺพุทฺโธ ปพฺพชฺชา\-สงฺขาเตน เอโก, โส ปจฺเจก\-สมฺพุทฺโธ สพฺพํ ฆราวาส\-ปลิโพธํ ฉินฺทิตฺวา ปุตฺตทาร\-ปลิโพธํ ฉินฺทิตฺวา ญาติพลิโพธํ ฉินฺทิตฺวา สนฺนิธิ\-ปลิโพธํ ฉินฺทิตฺวา เกสมสฺสุํ โอหาเรตฺวา กาสายานิ วตฺถานิ อจฺฉาเทตฺวา อคารสฺมา อนคาริยํ ปพฺพชิตฺวา อกิญฺจนภาวํ อุปคนฺตฺวา เอโก จรติ วิหรติ อิริยติ วตฺเตติ ปาเลติ ยเปติ ยาเปตีติ เอวํ โส ปจฺเจก\-สมฺพุทฺโธ ปพฺพชฺชา\-สงฺขาเตน เอโก.}

\prose{กถํ โส ปจฺเจก\-สมฺพุทฺโธ อทุติยฏฺเฐน เอโก, โส เอวํ ปพฺพชิโต สมาโน เอโก อรญฺญ\-วนปตฺถานิ ปนฺตานิ เสนาสนานิ ปฏิเสวติ อปฺปสทฺทานิ อปฺปนิคฺโฆสานิ วิชนวาตานิ มนุสฺส\-ราหสฺเสยฺยกานิ ปฏิสลฺลาน\-สารุปฺปานิ, โส เอโก คจฺฉติ, เอโก ติฏฺฐติ, เอโก นิสีทติ, เอโก เสยฺยํ กปฺเปติ, เอโก คามํ ปิณฺฑาย ปวิสติ, เอโก อภิกฺกมติ, เอโก ปฏิกฺกมติ, เอโก รโห นิสีทติ, เอโก จงฺกมํ อธิฏฺฐาติ, เอโก จรติ วิหรติ อิริยติ วตฺเตติ ปาเลติ ยเปติ ยาเปตีติ เอวํ โส ปจฺเจก\-สมฺพุทฺโธ อทุติยฏฺเฐน เอโก.}

\prose{กถํ โส ปจฺเจก\-สมฺพุทฺโธ ตณฺหาย ปหานฏฺเฐน เอโก, โส เอวํ เอโก อทุติโย อปฺปมตฺโต อาตาปี ปหิตตฺโต วิหรนฺโต มหาปธานํ ปทหนฺโต มารํ สเสนกํ นมุจิํ กณฺหํ ปมตฺตพนฺธุํ วิธเมตฺวา จ ตณฺหาชาลินิํ วิสริตํ วิสตฺติกํ ปชหิ วิโนเทสิ พฺยนฺตีอกาสิ อนภาวํ\-คเมสิ.}

\setlength{\csromangathapagewidth}{0pt}
\setlength{\csromangathawakwidth}{0pt}
\csromangathameasuregroup{“ตณฺหาทุติโย ปุริโส, \\ อิตฺถภาวญฺญถาภาวํ, \\ เอตมาทีนวํ ญ ตฺวา, \\ วีตตโณฺห อนาทาโน,}{\csromangathabat{“ตณฺหาทุติโย ปุริโส,}{ทีฆมทฺธาน สํสรํ.} \\ \csromangathabat{อิตฺถภาวญฺญถาภาวํ,}{สํสารํ นาติวตฺตติ.} \\ \csromangathabat{เอตมาทีนวํ ญ ตฺวา,}{ตณฺหํ ทุกฺขสฺส สมฺภวํ.} \\ \csromangathabat{วีตตโณฺห อนาทาโน,}{สโต ภิกฺขุ ปริพฺพเช”ติ.}}
\csromangathasetleft{“ตณฺหาทุติโย ปุริโส, \\ อิตฺถภาวญฺญถาภาวํ, \\ เอตมาทีนวํ ญ ตฺวา, \\ วีตตโณฺห อนาทาโน,}
\csromangathagroup{\csromangathastanza{\csromanfolio{231}\csromangathabat{“ตณฺหาทุติโย ปุริโส,}{ทีฆมทฺธาน สํสรํ.} \\ \csromangathabat{อิตฺถ\-ภาว\-ญฺญถา\-ภาวํ,}{สํสารํ นาติวตฺตติ.}}\csromangathastanza{\csromangathabat{เอตมาทีนวํ ญ ตฺวา,}{ตณฺหํ ทุกฺขสฺส สมฺภวํ.} \\ \csromangathabat{วีตตโณฺห อนาทาโน,}{สโต ภิกฺขุ ปริพฺพเช”ติ.}}}

\prose{เอวํ โส ปจฺเจก\-สมฺพุทฺโธ ตณฺหาย ปหานฏฺเฐน เอโก.}

\prose{กถํ โส ปจฺเจก\-สมฺพุทฺโธ เอกนฺตวีตราโคติ เอโก, ราคสฺส ปหีนตฺตา เอกนฺตวีตราโคติ เอโก, โทสสฺส ปหีนตฺตา เอกนฺตวีตโทโสติ เอโก, โมหสฺส ปหีนตฺตา เอกนฺตวีตโมโหติ เอโก, กิเลสานํ ปหีนตฺตา เอกนฺตนิกฺกิเลโสติ เอโก, เอวํ โส ปจฺเจก\-สมฺพุทฺโธ เอกนฺตวีตราโคติ เอโก.}

\prose{กถํ โส ปจฺเจก\-สมฺพุทฺโธ เอกายนมคฺคํ คโตติ เอโก, เอกายนมคฺโค วุจฺจติ จตฺตาโร สติปฏฺฐานา จตฺตาโร สมฺมปฺปธานา จตฺตาโร อิทฺธิปาทา ปญฺจินฺทฺริยานิ ปญฺจ พลานิ สตฺต โพชฺฌงฺคา อริโย อฏฺฐงฺคิโก มคฺโค.}

\setlength{\csromangathapagewidth}{0pt}
\setlength{\csromangathawakwidth}{0pt}
\csromangathameasuregroup{}{“เอกายนํ ชาติขยนฺตทสฺสี, \\ มคฺคํ ปชานาติ หิตานุกมฺปี. \\ เอเตน มคฺเคน ตริํสุ ปุพฺเพ, \\ ตริสฺสนฺติ เย จ ตรนฺติ โอฆนฺ”ติ.}
\csromangathameasurewak{“เอกายนํ ชาติขยนฺตทสฺสี, \\ มคฺคํ ปชานาติ หิตานุกมฺปี. \\ เอเตน มคฺเคน ตริํสุ ปุพฺเพ, \\ ตริสฺสนฺติ เย จ ตรนฺติ โอฆนฺ”ติ.}
\setlength{\csromangathaleftcol}{0pt}
\csromangathagroup{\csromangathastanza{\csromangathawak{“เอกายนํ ชาติ\-ขยนฺต\-ทสฺสี, \\ มคฺคํ ปชานาติ หิตานุกมฺปี. \\ เอเตน มคฺเคน ตริํสุ ปุพฺเพ, \\ ตริสฺสนฺติ เย จ ตรนฺติ โอฆนฺ”ติ.}}}

\prose{เอวํ โส ปจฺเจก\-สมฺพุทฺโธ เอกายนมคฺคํ คโตติ เอโก.}

\prose{กถํ โส ปจฺเจก\-สมฺพุทฺโธ เอโก อนุตฺตรํ ปจฺเจก\-สมฺโพธิํ อภิสมฺพุทฺโธติ เอโก, โพธิ วุจฺจติ จตูสุ มคฺเคสุ ญาณํ. ปญฺญา ปญฺญินฺทฺริยํ ปญฺญาพลํ ธมฺมวิจย\-สมฺโพชฺฌงฺโค วีมํสา วิปสฺสนา สมฺมาทิฏฺฐิ. โส ปจฺเจก\-สมฺพุทฺโธ มคฺค\-ปจฺเจก\-สมฺพุทฺโธ ญาณ\-ปจฺเจก\-สมฺพุทฺโธ “สพฺเพ สงฺขารา อนิจฺจา”ติ พุชฺฌิ, “สพฺเพ สงฺขารา ทุกฺขา”ติ พุชฺฌิ, “สพฺเพ ธมฺมา อนตฺตา”ติ พุชฺฌิ. “อวิชฺชาปจฺจยา สงฺขารา”ติ พุชฺฌิ, “สงฺขาร\-ปจฺจยา วิญฺญาณนฺ”ติ พุชฺฌิ, “วิญฺญาณปจฺจยา นามรูปนฺ”ติ พุชฺฌิ, “นามรูป\-ปจฺจยา สฬายตนนฺ”ติ พุชฺฌิ, “สฬายตน\-ปจฺจยา ผสฺโส”ติ พุชฺฌิ, “ผสฺสปจฺจยา เวทนา”ติ พุชฺฌิ, “เวทนาปจฺจยา ตณฺหา”ติ พุชฺฌิ, “ตณฺหาปจฺจยา อุปาทานนฺ”ติ พุชฺฌิ, “อุปาทานปจฺจยา ภโว”ติ พุชฺฌิ, “ภวปจฺจยา ชาตี”ติ พุชฺฌิ,}

\prose{\csromanfolio{232}“ชาติปจฺจยา ชรามรณนฺ”ติ พุชฺฌิ. “อวิชฺชานิโรธา สงฺขาร\-นิโรโธ”ติ พุชฺฌิ, “สงฺขาร\-นิโรธา วิญฺญาณนิโรโธ”ติ พุชฺฌิ, “วิญฺญาณนิโรธา นามรูป\-นิโรโธ”ติ พุชฺฌิ, “นามรูป\-นิโรธา สฬายตน\-นิโรโธ”ติ พุชฺฌิ, “สฬายตน\-นิโรธา ผสฺสนิโรโธ”ติ พุชฺฌิ, “ผสฺสนิโรธา เวทนานิโรโธ”ติ พุชฺฌิ, “เวทนานิโรธา ตณฺหานิโรโธ”ติ พุชฺฌิ, “ตณฺหานิโรธา อุปาทานนิโรโธ”ติ พุชฺฌิ, “อุปาทานนิโรธา ภวนิโรโธ”ติ พุชฺฌิ, “ภวนิโรธา ชาตินิโรโธ”ติ พุชฺฌิ, “ชาตินิโรธา ชรามรณ\-นิโรโธ”ติ พุชฺฌิ. “อิทํ ทุกฺขนฺ”ติ พุชฺฌิ, “อยํ ทุกฺขสมุทโย”ติ พุชฺฌิ, “อยํ ทุกฺขนิโรโธ”ติ พุชฺฌิ, “อยํ ทุกฺขนิโรธ\-คามินี ปฏิปทา”ติ พุชฺฌิ. “อิเม อาสวา”ติ พุชฺฌิ, “อยํ อาสวสมุทโย”ติ พุชฺฌิ ฯเปฯ “อยํ อาสว\-นิโรธ\-คามินี ปฏิปทา”ติ พุชฺฌิ. “อิเม ธมฺมา อภิญฺเญยฺยา”ติ พุชฺฌิ, “อิเม ธมฺมา ปหาตพฺพา”ติ พุชฺฌิ, “อิเม ธมฺมา สจฺฉิกาตพฺพา”ติ พุชฺฌิ, “อิเม ธมฺมา ภาเวตพฺพา”ติ พุชฺฌิ. ฉนฺนํ ผสฺสา\-ยตนานํ สมุทยญฺจ อตฺถงฺคมญฺจ อสฺสาทญฺจ อาทีนวญฺจ นิสฺสรณญฺจ พุชฺฌิ, ปญฺจนฺนํ อุปาทาน\-กฺขนฺธานํ สมุทยญฺจ ฯเปฯ นิสฺสรณญฺจ พุชฺฌิ, จตุนฺนํ มหาภูตานํ สมุทยญฺจ อตฺถงฺคมญฺจ อสฺสาทญฺจ อาทีนวญฺจ นิสฺสรณญฺจ พุชฺฌิ, “ยํ กิญฺจิ สมุทย\-ธมฺมํ, สพฺพํ ตํ นิโรธธมฺมนฺ”ติ พุชฺฌิ.}

\prose{อถ วา ยํ พุชฺฌิตพฺพํ อนุพุชฺฌิตพฺพํ ปฏิพุชฺฌิตพฺพํ สมฺพุชฺฌิตพฺพํ อธิคนฺตพฺพํ ผสฺสิตพฺพํ สจฺฉิกาตพฺพํ, สพฺพํ ตํ เตน ปจฺเจกโพธิ\-ญาเณน พุชฺฌิ อนุพุชฺฌิ ปฏิพุชฺฌิ สมฺพุชฺฌิ อธิคจฺฉิ ผสฺเสสิ สจฺฉากาสีติ เอวํ โส ปจฺเจก\-สมฺพุทฺโธ เอโก อนุตฺตรํ ปจฺเจก\-สมฺโพธิํ อภิสมฺพุทฺโธติ เอโก.}

\prose{\textbf{จเร}ติ อฏฺฐ จริยาโย อิริยาปถ\-จริยา อายตนจริยา สติจริยา สมาธิจริยา ญาณจริยา มคฺคจริยา ปตฺติจริยา โลกตฺถจริยา. \textbf{อิริยาปถ}\-\textbf{จริยา}ติ จตูสุ อิริยาปเถสุ. \textbf{อายตน}\-\textbf{จริยา}ติ ฉสุ อชฺฌตฺติก\-พาหิเรสุ อายตเนสุ. \textbf{สติจริยา}ติ จตูสุ สติปฏฺฐาเนสุ. \textbf{สมาธิจริยา}ติ จตูสุ ฌาเนสุ. \textbf{ญาณจริยา}ติ จตูสุ อริยสจฺเจสุ. \textbf{มคฺคจริยา}ติ จตูสุ อริยมคฺเคสุ. \textbf{ปตฺติจริยา}ติ จตูสุ สามญฺญผเลสุ. \textbf{โลกตฺถจริยา}ติ ตถาคเตสุ อรหนฺเตสุ สมฺมา\-สมฺพุทฺเธสุ ปเทสโต ปจฺเจก\-สมฺพุทฺเธสุ ปเทสโต สาวเกสุ. อิริยาปถ\-จริยา จ ปณิธิ\-สมฺปนฺนานํ, อายตนจริยา จ อินฺทฺริเยสุ}

\prose{\csromanfolio{233}คุตฺตทฺวารานํ, สติจริยา จ อปฺปมาท\-วิหารีนํ, สมาธิจริยา จ อธิจิตฺตม\-นุยุตฺตานํ, ญาณจริยา จ พุทฺธิ\-สมฺปนฺนานํ, มคฺคจริยา จ สมฺมา\-ปฏิปนฺนานํ, ปตฺติจริยา จ อธิคต\-ผลานํ, โลกตฺถจริยา จ ตถาคตานํ อรหนฺตานํ สมฺมา\-สมฺพุทฺธานํ ปเทสโต ปจฺเจก\-พุทฺธานํ ปเทสโต สาวกานํ, อิมา อฏฺฐ จริยาโย. อปราปิ อฏฺฐ จริยาโย อธิมุจฺจนฺโต สทฺธาย จรติ, ปคฺคณฺหนฺโต วีริเยน จรติ, อุปฏฺฐาเปนฺโต สติยา จรติ, อวิกฺเขปํ กโรนฺโต สมาธินา จรติ, ปชานนฺโต ปญฺญาย จรติ, วิชานนฺโต วิญฺญาณ\-จริยาย จรติ. “เอวํ ปฏิปนฺนสฺส กุสลา ธมฺมา อายาเปนฺตี”ติ อายตน\-จริยาย จรติ. “เอวํ ปฏิปนฺโน วิเสสมธิคจฺฉตี”ติ วิเสสจริยาย จรติ. อิมา อฏฺฐ จริยาโย.}

\prose{อปราปิ อฏฺฐ จริยาโย ทสฺสนจริยา จ สมฺมาทิฏฺฐิยา, อภินิโรปน\-จริยา จ สมฺมา\-สงฺกปฺปสฺส, ปริคฺคห\-จริยา จ สมฺมาวาจาย, สมุฏฺฐาน\-จริยา จ สมฺมา\-กมฺมนฺตสฺส, โวทานจริยา จ สมฺมาอาชีวสฺส, ปคฺคหจริยา จ สมฺมาวายามสฺส, อุปฏฺฐาน\-จริยา จ สมฺมาสติยา, อวิกฺเขป\-จริยา จ สมฺมา\-สมาธิสฺส, อิมา อฏฺฐ จริยาโย.}

\prose{\textbf{ขคฺควิสาณ}\-\textbf{กปฺโป}ติ ยถา ขคฺคสฺส นาม วิสาณํ เอกํ โหติ อทุติยํ, เอวเมว โส ปจฺเจก\-สมฺพุทฺโธ ตกฺกปฺโป ตสฺสทิโส ตปฺปฏิภาโค. ยถา อติโลณํ วุจฺจติ โลณกปฺโป, อติติตฺตกํ วุจฺจติ ติตฺตกปฺโป, อติมธุรํ วุจฺจติ มธุรกปฺโป, อติอุณฺหํ วุจฺจติ อคฺคิกปฺโป, อติสีตลํ วุจฺจติ หิมกปฺโป, มหา\-อุทกกฺขนฺโธ วุจฺจติ สมุทฺทกปฺโป, มหาภิญฺญา\-พลปฺปตฺโต สาวโก วุจฺจติ สตฺถุกปฺโปติ. เอวเมว โส ปจฺเจก\-สมฺพุทฺโธ ตตฺถ ตกฺกปฺโป ตสฺสทิโส ตปฺปฏิภาโค เอโก อทุติโย มุตฺตพนฺธโน สมฺมา โลเก จรติ วิหรติ อิริยติ วตฺเตติ ปาเลติ ยเปติ ยาเปตีติ เอโก จเร ขคฺควิสาณ\-กปฺโป. เตนาห โส ปจฺเจก\-สมฺพุทฺโธ\csromandash{}}

\setlength{\csromangathapagewidth}{0pt}
\setlength{\csromangathawakwidth}{0pt}
\csromangathameasuregroup{}{สพฺเพสุ ภูเตสุ นิธาย ทณฺฑํ, \\ อวิเหฐยํ อญฺญตรมฺปิ เตสํ. \\ น ปุตฺตมิจฺเฉยฺย กุโต สหายํ, \\ เอโก จเร ขคฺควิสาณกปฺโปติ. \\ \textbf{สํสคฺค}\textbf{ชาตสฺส ภวนฺติ สฺเนหา,} \\ \textbf{สฺเนหนฺวยํ ทุกฺขมิทํ ปโหติ.} \\ \textbf{อาทีนวํ สฺเนหชํ เปกฺขมาโน,} \\ \textbf{เอโก จเร ขคฺควิสาณ}\textbf{กปฺโป.}}
\csromangathameasurewak{สพฺเพสุ ภูเตสุ นิธาย ทณฺฑํ, \\ อวิเหฐยํ อญฺญตรมฺปิ เตสํ. \\ น ปุตฺตมิจฺเฉยฺย กุโต สหายํ, \\ เอโก จเร ขคฺควิสาณกปฺโปติ. \\ \textbf{สํสคฺค}\textbf{ชาตสฺส ภวนฺติ สฺเนหา,} \\ \textbf{สฺเนหนฺวยํ ทุกฺขมิทํ ปโหติ.} \\ \textbf{อาทีนวํ สฺเนหชํ เปกฺขมาโน,} \\ \textbf{เอโก จเร ขคฺควิสาณ}\textbf{กปฺโป.}}
\setlength{\csromangathaleftcol}{0pt}
\csromangathagroup{\csromangathastanza{\csromangathawak{สพฺเพสุ ภูเตสุ นิธาย ทณฺฑํ, \\ อวิเหฐยํ อญฺญตรมฺปิ เตสํ. \\ น ปุตฺตมิจฺเฉยฺย กุโต สหายํ, \\ เอโก จเร ขคฺควิสาณ\-กปฺโปติ.}}\csromangathastanza{\csromangathawak{\csromanfolio{234}\csromangathaitemline{122}{\textbf{สํสคฺค}\-\textbf{ชาตสฺส ภวนฺติ สฺเนหา,}} \\ \csromangathacontline{\textbf{สฺเนหนฺวยํ ทุกฺขมิทํ ปโหติ.}} \\ \csromangathacontline{\textbf{อาทีนวํ สฺเนหชํ เปกฺขมาโน,}} \\ \csromangathacontline{\textbf{เอโก จเร ขคฺควิสาณ}\-\textbf{กปฺโป.}}}}}

\prose{\textbf{สํสคฺค}\-\textbf{ชาตสฺส ภวนฺติ สฺเนหา}ติ สํสคฺคาติ เทฺว \textbf{สํสคฺคา} ทสฺสน\-สํสคฺโค จ สวนสํสคฺโค จ. กตโม ทสฺสน\-สํสคฺโค, อิเธกจฺโจ ปสฺสติ อิตฺถิํ วา กุมาริํ วา อภิรูปํ ทสฺสนียํ ปาสาทิกํ ปรมาย วณฺณ\-โปกฺขร\-ตาย สมนฺนาคตํ, ทิสฺวา ปสฺสิตฺวา อนุพฺยญฺชนโส นิมิตฺตํ คณฺหาติ เกสา วา โสภนา\csromansharedfootnote{a689849e0b8836ef}{โสภณา (สฺยา)}, มุขํ วา โสภนํ, อกฺขี วา โสภนา, กณฺณา วา โสภนา, นาสา วา โสภนา, โอฏฺฐา วา โสภนา, ทนฺตา วา โสภนา, มุขํ วา โสภนํ, คีวา วา โสภนา, ถนา วา โสภนา, อุรํ วา โสภนํ, อุทรํ วา โสภนํ, กฏิ วา โสภนา, อูรู วา โสภนา, ชงฺฆา วา โสภนา, หตฺถา วา โสภนา, ปาทา วา โสภนา, องฺคุลิโย วา โสภนา, นขา วา โสภนาติ ทิสฺวา ปสฺสิตฺวา อภินนฺทติ อภิวทติ อภิปตฺเถติ อนุปฺปาเทติ\csromansharedfootnote{3b707ec5ce4f04eb}{อนุสฺสรติ (ก)} อนุพนฺธติ ราคพนฺธนํ, อยํ ทสฺสน\-สํสคฺโค.}

\prose{กตโม สวนสํสคฺโค, อิเธกจฺโจ สุณาติ “อสุกสฺมิํ นาม คาเม วา นิคเม วา อิตฺถี วา กุมารี วา อภิรูปา ทสฺสนียา ปาสาทิกา ปรมาย วณฺณ\-โปกฺขร\-ตาย สมนฺนาคตา”ติ สุตฺวา สุณิตฺวา อภินนฺทติ อภิวทติ อภิปตฺเถติ อนุปฺปาเทติ อนุพนฺธติ ราคพนฺธนํ, อยํ สวนสํสคฺโค.}

\prose{\textbf{สฺเนหา}ติ เทฺว สฺเนหา ตณฺหาสฺเนโห จ ทิฏฺฐิสฺเนโห จ. กตโม ตณฺหาสฺเนโห, ยาวตา ตณฺหา\-สงฺขาเตน สีมกตํ มริยาทิกตํ โอธิกตํ ปริยนฺตกตํ ปริคฺคหิตํ มมายิตํ, อิทํ มม, เอตํ มม, เอตฺตกํ มม, เอตฺตาวตา มม, มม รูปา สทฺทา คนฺธา รสา โผฏฺฐพฺพา อตฺถรณา ปาวุรณา ทาสิทาสา อเชฬกา กุกฺกุฏสูกรา หตฺถิ\-ควา\-สฺส\-วฬวา เขตฺตํ วตฺถุ หิรญฺญํ สุวณฺณํ คามนิคม\-ราชธานิโย รฏฺฐญฺจ ชนปโท จ โกโส จ โกฏฺฐาคารญฺจ เกวลมฺปิ มหาปถวิํ ตณฺหาวเสน มมายติ, ยาวตา อฏฺฐสตํ ตณฺหาวิจริตํ, อยํ ตณฺหาสฺเนโห.}

\prose{\csromanfolio{235}กตโม ทิฏฺฐิสฺเนโห, วีสติวตฺถุกา สกฺกายทิฏฺฐิ, ทสวตฺถุกา มิจฺฉาทิฏฺฐิ, ทสวตฺถุกา อนฺตคฺคาหิกา ทิฏฺฐิ. ยา เอวรูปา ทิฏฺฐิ ทิฏฺฐิคตํ ทิฏฺฐิคหนํ ทิฏฺฐิกนฺตาโร ทิฏฺฐิ\-วิสูกายิกํ ทิฏฺฐิ\-วิปฺผนฺทิตํ ทิฏฺฐิ\-สํโยชนํ คาโห ปฏิคฺคาโห\csromansharedfootnote{399566196e5c6340}{ปติฏฺฐาโห (ก)} อภินิเวโส ปรามาโส กุมฺมคฺโค มิจฺฉาปโถ มิจฺฉตฺตํ ติตฺถายตนํ วิปริยาส\-คฺคาโห วิปรีตคฺคาโห วิปลฺลาสคฺคาโห มิจฺฉาคาโห อยาถาวกสฺมิํ “ยาถาวกนฺ”ติ คาโห ยาวตา ทฺวาสฏฺฐิ ทิฏฺฐิคตานิ. อยํ ทิฏฺฐิสฺเนโห.}

\prose{\textbf{สํสคฺค}\-\textbf{ชาตสฺส ภวนฺติ สฺเนหา}ติ ทสฺสน\-สํสคฺค\-ปจฺจยา จ สวน\-สํสคฺค\-ปจฺจยา จ ตณฺหาสฺเนโห จ ทิฏฺฐิสฺเนโห จ ภวนฺติ สมฺภวนฺติ ชายนฺติ สญฺชายนฺติ นิพฺพตฺตนฺติ อภินิพฺพตฺตนฺติ ปาตุภวนฺตีติ สํสคฺค\-ชาตสฺส ภวนฺติ สฺเนหา.}

\prose{\textbf{สฺเนหนฺวยํ ทุกฺขมิทํ ปโหตี}ติ \textbf{สฺเนโห}ติ เทฺว สฺเนหา ตณฺหาสฺเนโย จ ทิฏฺฐิสฺเนโห จ ฯเปฯ อยํ ตณฺหาสฺเนโห ฯเปฯ อยํ ทิฏฺฐิสฺเนโห. \textbf{ทุกฺขมิทํ ปโหตี}ติ อิเธกจฺโจ กาเยน ทุจฺจริตํ จรติ, วาจาย ทุจฺจริตํ จรติ, มนสา ทุจฺจริตํ จรติ, ปาณมฺปิ หนติ, อทินฺนมฺปิ อาทิยติ, สนฺธิมฺปิ ฉินฺทติ, นิลฺโลปมฺปิ\csromansharedfootnote{909cc6ca74ad2031}{วิโลปมฺปิ (สฺยา) ปสฺส ขุ 7. 315 ปิฏฺเฐ.} หรติ, เอกาคาริกมฺปิ กโรติ, ปริปนฺเถปิ ติฏฺฐติ, ปรทารมฺปิ คจฺฉติ, มุสาปิ ภณติ, ตเมนํ คเหตฺวา รญฺโญ ทสฺเสนฺติ “อยํ เทว โจโร อาคุจารี, อิมสฺส ยํ อิจฺฉสิ, ตํ ทณฺฑํ ปเณหี”ติ. ตเมนํ ราชา ตํ ปริภาสติ, โส ปริภาส\-ปจฺจยาปิ ทุกฺขํ โทมนสฺสํ\csromansharedfootnote{c248ed5bb979034f}{ทุกฺขโทมนสฺสํ (สฺยา)} ปฏิสํเวเทติ. เอตํ ภยํ ทุกฺขํ โทมนสฺสํ กุโต ชาตํ, ตสฺส สฺเนหปจฺจยา จ นนฺทิปจฺจยา จ ราคปจฺจยา จ นนฺทิ\-ราคปจฺจยา จ ชาตํ.}

\prose{เอตฺตเกนปิ ราชา น ตุสฺสติ, ตเมนํ ราชา พนฺธาเปติ อนฺทุพนฺธเนน วา รชฺชุ\-พนฺธเนน วา สงฺขาลิกพนฺธเนน วา เวตฺต\-พนฺธเนน วา ลตาพนฺธเนน วา ปกฺเขป\-พนฺธเนน วา ปริกฺเขป\-พนฺธเนน วา คามพนฺธเนน วา นิคม\-พนฺธเนน วา รฏฺฐ\-พนฺธเนน วา ชนปท\-พนฺธเนน วา อนฺตมโส สวจนียมฺปิ กโรติ “น เต ลพฺภา อิโต ปกฺกมิตุนฺ”ติ, โส พนฺธน\-ปจฺจยาปิ ทุกฺขํ โทมนสฺสํ ปฏิสํเวเทติ. เอตํ ภยํ ทุกฺขํ โทมนสฺสํ กุโต ชาตํ, ตสฺส สฺเนหปจฺจยา จ นนฺทิปจฺจยา จ ราคปจฺจยา จ นนฺทิ\-ราคปจฺจยา จ ชาตํ.}

\prose{\csromanfolio{236}เอตฺตเกนปิ ราชา น ตุสฺสติ, ตเมนํ ราชา ตสฺเสว\csromansharedfootnote{0e9f8360f7a3e995}{ตสฺส (สฺยา)} ธนํ อาหราเปติ สตํ วา สหสฺสํ วา สตสหสฺสํ วา, โส ธนชานิ\-ปจฺจยาปิ ทุกฺขํ โทมนสฺสํ ปฏิสํเวเทติ. เอตํ ภยํ ทุกฺขํ โทมนสฺสํ กุโตชาตํ, ตสฺส สฺเนหปจฺจยา จ นนฺทิปจฺจยา จ ราคปจฺจยา จ นนฺทิ\-ราคปจฺจยา จ ชาตํ.}

\prose{เอตฺตเกนปิ ราชา น ตุสฺสติ, ตเมนํ ราชา วิวิธา กมฺมการณา\csromansharedfootnote{7989c8fc7b9a6a4f}{วิวิธานิ กมฺมกรณานิ (ก)} การาเปติ, กสาหิปิ ตาเฬติ, เวตฺเตหิปิ ตาเฬติ, อฑฺฒทณฺเฑหิปิ ตาเฬติ, หตฺถมฺปิ ฉินฺทติ, ปาทมฺปิ ฉินฺทติ, หตฺถปาทมฺปิ ฉินฺทติ, กณฺณมฺปิ ฉินฺทติ, นาสมฺปิ ฉินฺทติ, กณฺณนาสมฺปิ ฉินฺทติ, พิลงฺค\-ถาลิกมฺปิ กโรติ, สงฺข\-มุณฺฑิกมฺปิ กโรติ, ราหุมุขมฺปิ กโรติ, โชติมาลิกมฺปิ กโรติ, หตฺถ\-ปชฺโชติกมฺปิ กโรติ, เอรก\-วตฺติกมฺปิ\csromansharedfootnote{5ea562081dd84d1b}{เอรกวฏฺฏิกมฺปิ (สฺยา, ก) ปสฺส ม 1. 122 ปิฏฺเฐ.} กโรติ, จีรก\-วาสิกมฺปิ กโรติ, เอเณยฺยกมฺปิ กโรติ, พฬิส\-มํสิกมฺปิ กโรติ, กหาปณิกมฺปิ กโรติ, ขารา\-ปตจฺฉิกมฺปิ\csromansharedfootnote{ba0fc14a82a674ca}{ขาราปฏิจฺฉิกมฺปิ (ก) เหฏฺฐา 190 ปิฏฺเฐ.} กโรติ, ปลิฆ\-ปริวตฺติกมฺปิ กโรติ, ปลาล\-ปีฐกมฺปิ กโรติ, ตตฺเตนปิ เตเลน โอสิญฺจติ, สุนเขหิปิ ขาทาเปติ, ชีวนฺตมฺปิ สูเล อุตฺตาเสติ, อสินาปิ สีสํ ฉินฺทติ, โส กมฺมการณ\-ปจฺจยาปิ ทุกฺขํ โทมนสฺสํ ปฏิสํเวเทติ. เอตํ ภยํ ทุกฺขํ โทมนสฺสํ กุโต ชาตํ, ตสฺส สฺเนหปจฺจยา จ นนฺทิปจฺจยา จ ราคปจฺจยา จ นนฺทิ\-ราคปจฺจยา จ ชาตํ. ราชา อิเมสํ จตุนฺนํ ทณฺฑานํ อิสฺสโร.}

\prose{โส สเกน กมฺเมน กายสฺส เภทา ปรํ มรณา อปายํ ทุคฺคติํ วินิปาตํ นิรยํ อุปปชฺชติ. ตเมนํ นิรยปาลา ปญฺจ\-วิธ\-พนฺธนํ นาม กมฺมการณํ กโรนฺติ\csromansharedfootnote{70ced0ca97109850}{กาเรนฺติ (สฺยา, ก) ปสฺส ม 3. 221 ปิฏฺเฐ.}, ตตฺตํ อโยขิลํ หตฺเถ คเมนฺติ, ตตฺตํ อโยขิลํ ทุติเย หตฺเถ คเมนฺติ, ตตฺตํ อโยขิลํ ปาเท คเมนฺติ, ตตฺตํ อโยขิลํ ทุติเย ปาเท คเมนฺติ, ตตฺตํ อโยขิลํ มชฺเฌ อุรสฺมิํ คเมนฺติ. โส ตตฺถ ทุกฺขา ติพฺพา\csromansharedfootnote{4070fd8c92cc8b2b}{ติปฺปา (สฺยา)} ขรา กฏุกา เวทนา เวเทติ, น จ ตาว กาลํ กโรติ, ยาว น ตํ ปาปกมฺมํ พฺยนฺตีโหติ. เอตํ ภยํ ทุกฺขํ โทมนสฺสํ กุโต ชาตํ, ตสฺส สฺเนหปจฺจยา จ นนฺทิปจฺจยา จ ราคปจฺจยา จ นนฺทิ\-ราคปจฺจยา จ ชาตํ.}

\prose{\csromanfolio{237}ตเมนํ นิรยปาลา สํเวเสตฺวา\csromansharedfootnote{a55aa3e6c1e6f5a1}{สํเวสิตฺวา (สฺยา, ก) ปสฺส ม 3. 221 ปิฏฺเฐ.} กุฐารีหิ\csromansharedfootnote{1af5086b3d661632}{กุธารีหิ (สฺยา, ก)} ตจฺฉนฺติ ฯเปฯ. ตเมนํ นิรยปาลา อุทฺธํปาทํ อโธสิรํ คเหตฺวา วาสีหิ ตจฺฉนฺติ. ตเมนํ นิรยปาลา รเถ โยเชตฺวา อาทิตฺตาย ปถวิยา สมฺปชฺชลิตาย สโชติภูตาย\csromansharedfootnote{fa34c419341e9ba4}{สญฺโชติภูตาย (สฺยา)} สาเรนฺติปิ ปจฺจาสาเรนฺติปิ. ตเมนํ นิรยปาลา มหนฺตํ องฺคารปพฺพตํ อาทิตฺตํ สมฺปชฺชลิตํ สโชติภูตํ อาโรเปนฺติปิ โอโรเปนฺติปิ. ตเมนํ นิรยปาลา อุทฺธํปาทํ อโธสิรํ คเหตฺวา ตตฺตาย โลหกุมฺภิยา ปกฺขิปนฺติ อาทิตฺตาย สมฺปชฺชลิตาย สโชติภูตาย, โส ตตฺถ เผณุทฺเทหกํ ปจฺจติ, โส ตตฺถ เผณุทฺเทหกํ ปจฺจมาโน สกิมฺปิ อุทฺธํ คจฺฉติ, สกิมฺปิ อโธ คจฺฉติ, สกิมฺปิ ติริยํ คจฺฉติ. โส ตตฺถ ทุกฺขา ติพฺพา ขรา กฏุกา เวทนา เวเทติ, น จ ตาว กาลํ กโรติ, ยาว น ตํ ปาปกมฺมํ พฺยนฺตีโหติ. เอตํ ภยํ ทุกฺขํ โทมนสฺสํ กุโต ชาตํ, ตสฺส สฺเนหปจฺจยา จ นนฺทิปจฺจยา จ ราคปจฺจยา จ นนฺทิ\-ราคปจฺจยา จ ชาตํ.}

\prose{ตเมนํ นิรยปาลา มหานิรเย ปกฺขิปนฺติ, โส โข ปน มหานิรโย\csromandash{}}

\setlength{\csromangathapagewidth}{0pt}
\setlength{\csromangathawakwidth}{0pt}
\csromangathameasuregroup{จตุกฺกณฺโณ จตุทฺวาโร, \\ อโยปาการปริยนฺโต, \\ ตสฺส อโยมยา ภูมิ, \\ สมนฺตา โยชนสตํ, \\ กทริยาตปนา\textsuperscript{1} โฆรา, \\ โลมหํสนรูปา จ, \\ ปุรตฺถิมาย จ ภิตฺติยา, \\ qอหนฺโต ปาปกมฺมนฺเต, \\ ปจฺฉิมาย จ ภิตฺติยา, \\ qอหนฺโต ปาปกมฺมนฺเต, \\ ทกฺขิณาย จ ภิตฺติยา, \\ qอหนฺโต ปาปกมฺมนฺเต, \\ อุตฺตราย จ ภิตฺติยา, \\ qอหนฺโต ปาปกมฺมนฺเต, \\ เหฏฺฐโต จ สมุฏฺฐาย, \\ qอหนฺโต ปาปกมฺมนฺเต, \\ ฉทนมฺหา สมุฏฺฐาย, \\ qอหนฺโต ปาปกมฺมนฺเต, \\ อโยกปาลมาทิตฺตํ, \\ เอวํ อวีจินิรโย, \\ ตตฺถ สตฺตา มหาลุทฺทา, \\ อจฺจนฺตปาปกมฺมนฺตา, \\ ชาตเวทสโม กาโย, \\ ปสฺส กมฺมานํ ทฬฺหตฺตํ, \\ ปุรตฺถิเมนปิ ธาวนฺติ, \\ อุตฺตเรนปิ ธาวนฺติ, \\ ยํ ยํ ทิสํ ปธาวนฺติ, \\ อภินิกฺขมิตาสา เต, \\ น เต ตโต นิกฺขมิตุํ, \\ เตสญฺจ ปาปกมฺมนฺตํ,}{\csromangathabat{จตุกฺกณฺโณ จตุทฺวาโร,}{วิภตฺโต ภาคโส มิโต.} \\ \csromangathabat{อโยปาการปริยนฺโต,}{อยสา ปฏิกุชฺชิโต.} \\ \csromangathabat{ตสฺส อโยมยา ภูมิ,}{ชลิตา เตชสา ยุตา.} \\ \csromangathabat{สมนฺตา โยชนสตํ,}{ผริตฺวา ติฏฺฐติ สพฺพทา.} \\ \csromangathabat{กทริยาตปนา\textsuperscript{1} โฆรา,}{อจฺจิมนฺโต ทุราสทา.} \\ \csromangathabat{โลมหํสนรูปา จ,}{เภสฺมา ปฏิภยา ทุขา.} \\ \csromangathabat{ปุรตฺถิมาย จ ภิตฺติยา,}{อจฺจิกฺขนฺโธ สมุฏฺฐิโต.} \\ \csromangathabat{qอหนฺโต ปาปกมฺมนฺเต,}{ปจฺฉิมาย ปฏิหญฺญติ.} \\ \csromangathabat{ปจฺฉิมาย จ ภิตฺติยา,}{อจฺจิกฺขนฺโธ สมุฏฺฐิโต.} \\ \csromangathabat{qอหนฺโต ปาปกมฺมนฺเต,}{ปุริมาย ปฏิหญฺญติ.} \\ \csromangathabat{ทกฺขิณาย จ ภิตฺติยา,}{อจฺจิกฺขนฺโธ สมุฏฺฐิโต.} \\ \csromangathabat{qอหนฺโต ปาปกมฺมนฺเต,}{อุตฺตราย ปฏิหญฺญติ.} \\ \csromangathabat{อุตฺตราย จ ภิตฺติยา,}{อจฺจิกฺขนฺโธ สมุฏฺฐิโต.} \\ \csromangathabat{qอหนฺโต ปาปกมฺมนฺเต,}{ทกฺขิณาย ปฏิหญฺญติ.} \\ \csromangathabat{เหฏฺฐโต จ สมุฏฺฐาย,}{อจฺจิกฺขนฺโธ ภยานโก.} \\ \csromangathabat{qอหนฺโต ปาปกมฺมนฺเต,}{ฉทนสฺมิํ ปฏิหญฺญติ.} \\ \csromangathabat{ฉทนมฺหา สมุฏฺฐาย,}{อจฺจิกฺขนฺโธ ภยานโก.} \\ \csromangathabat{qอหนฺโต ปาปกมฺมนฺเต,}{ภูมิยํ ปฏิหญฺญติ.} \\ \csromangathabat{อโยกปาลมาทิตฺตํ,}{สนฺตตฺตํ ชลิตํ ยถา.} \\ \csromangathabat{เอวํ อวีจินิรโย,}{เหฏฺฐา อุปริ ปสฺสโต.} \\ \csromangathabat{ตตฺถ สตฺตา มหาลุทฺทา,}{มหากิพฺพิสการิโน.} \\ \csromangathabat{อจฺจนฺตปาปกมฺมนฺตา,}{ปจฺจนฺติ น จ มิยฺยเร\textsuperscript{1}.} \\ \csromangathabat{ชาตเวทสโม กาโย,}{เตสํ นิรยวาสินํ.} \\ \csromangathabat{ปสฺส กมฺมานํ ทฬฺหตฺตํ,}{น ภสฺมา โหติ นปิ มสิ.} \\ \csromangathabat{ปุรตฺถิเมนปิ ธาวนฺติ,}{ตโต ธาวนฺติ ปจฺฉิมํ.} \\ \csromangathabat{อุตฺตเรนปิ ธาวนฺติ,}{ตโต ธาวนฺติ ทกฺขิณํ.} \\ \csromangathabat{ยํ ยํ ทิสํ ปธาวนฺติ,}{ตํ ตํ ทฺวารํ ปิธียติ.} \\ \csromangathabat{อภินิกฺขมิตาสา เต,}{สตฺตา โมกฺขคเวสิโน.} \\ \csromangathabat{น เต ตโต นิกฺขมิตุํ,}{ลภนฺติ กมฺมปจฺจยา.} \\ \csromangathabat{เตสญฺจ ปาปกมฺมนฺตํ,}{อวิปกฺกํ กตํ พหุนฺติ.}}
\csromangathasetleft{จตุกฺกณฺโณ จตุทฺวาโร, \\ อโยปาการปริยนฺโต, \\ ตสฺส อโยมยา ภูมิ, \\ สมนฺตา โยชนสตํ, \\ กทริยาตปนา\textsuperscript{1} โฆรา, \\ โลมหํสนรูปา จ, \\ ปุรตฺถิมาย จ ภิตฺติยา, \\ qอหนฺโต ปาปกมฺมนฺเต, \\ ปจฺฉิมาย จ ภิตฺติยา, \\ qอหนฺโต ปาปกมฺมนฺเต, \\ ทกฺขิณาย จ ภิตฺติยา, \\ qอหนฺโต ปาปกมฺมนฺเต, \\ อุตฺตราย จ ภิตฺติยา, \\ qอหนฺโต ปาปกมฺมนฺเต, \\ เหฏฺฐโต จ สมุฏฺฐาย, \\ qอหนฺโต ปาปกมฺมนฺเต, \\ ฉทนมฺหา สมุฏฺฐาย, \\ qอหนฺโต ปาปกมฺมนฺเต, \\ อโยกปาลมาทิตฺตํ, \\ เอวํ อวีจินิรโย, \\ ตตฺถ สตฺตา มหาลุทฺทา, \\ อจฺจนฺตปาปกมฺมนฺตา, \\ ชาตเวทสโม กาโย, \\ ปสฺส กมฺมานํ ทฬฺหตฺตํ, \\ ปุรตฺถิเมนปิ ธาวนฺติ, \\ อุตฺตเรนปิ ธาวนฺติ, \\ ยํ ยํ ทิสํ ปธาวนฺติ, \\ อภินิกฺขมิตาสา เต, \\ น เต ตโต นิกฺขมิตุํ, \\ เตสญฺจ ปาปกมฺมนฺตํ,}
\csromangathagroup{\csromangathastanza{\csromangathabat{จตุกฺกณฺโณ จตุทฺวาโร,}{วิภตฺโต ภาคโส มิโต.} \\ \csromangathabat{อโยปาการปริยนฺโต,}{อยสา ปฏิกุชฺชิโต.}}\csromangathastanza{\csromangathabat{ตสฺส อโยมยา ภูมิ,}{ชลิตา เตชสา ยุตา.} \\ \csromangathabat{สมนฺตา โยชนสตํ,}{ผริตฺวา ติฏฺฐติ สพฺพทา.}}\csromangathastanza{\csromangathabat{กทริยาตปนา\csromansharedfootnote{bfc1d370003df029}{กทริยา ตาปนา (ก) ขุ 7. 317 ปิฏฺเฐปิ.} โฆรา,}{อจฺจิมนฺโต ทุราสทา.} \\ \csromangathabat{โลมหํสนรูปา จ,}{เภสฺมา ปฏิภยา ทุขา.}}\csromangathastanza{\csromangathabat{ปุรตฺถิมาย จ ภิตฺติยา,}{อจฺจิกฺขนฺโธ สมุฏฺฐิโต.} \\ \csromangathabat{qอหนฺโต ปาปกมฺมนฺเต,}{ปจฺฉิมาย ปฏิหญฺญติ.}}\csromangathastanza{\csromangathabat{ปจฺฉิมาย จ ภิตฺติยา,}{อจฺจิกฺขนฺโธ สมุฏฺฐิโต.} \\ \csromangathabat{qอหนฺโต ปาปกมฺมนฺเต,}{ปุริมาย ปฏิหญฺญติ.}}\csromangathastanza{\csromangathabat{ทกฺขิณาย จ ภิตฺติยา,}{อจฺจิกฺขนฺโธ สมุฏฺฐิโต.} \\ \csromangathabat{qอหนฺโต ปาปกมฺมนฺเต,}{อุตฺตราย ปฏิหญฺญติ.}}\csromangathastanza{\csromanfolio{238}\csromangathabat{อุตฺตราย จ ภิตฺติยา,}{อจฺจิกฺขนฺโธ สมุฏฺฐิโต.} \\ \csromangathabat{qอหนฺโต ปาปกมฺมนฺเต,}{ทกฺขิณาย ปฏิหญฺญติ.}}\csromangathastanza{\csromangathabat{เหฏฺฐโต จ สมุฏฺฐาย,}{อจฺจิกฺขนฺโธ ภยานโก.} \\ \csromangathabat{qอหนฺโต ปาปกมฺมนฺเต,}{ฉทนสฺมิํ ปฏิหญฺญติ.}}\csromangathastanza{\csromangathabat{ฉทนมฺหา สมุฏฺฐาย,}{อจฺจิกฺขนฺโธ ภยานโก.} \\ \csromangathabat{qอหนฺโต ปาปกมฺมนฺเต,}{ภูมิยํ ปฏิหญฺญติ.}}\csromangathastanza{\csromangathabat{อโยกปาลมาทิตฺตํ,}{สนฺตตฺตํ ชลิตํ ยถา.} \\ \csromangathabat{เอวํ อวีจินิรโย,}{เหฏฺฐา อุปริ ปสฺสโต.}}\csromangathastanza{\csromangathabat{ตตฺถ สตฺตา มหาลุทฺทา,}{มหากิพฺพิสการิโน.} \\ \csromangathabat{อจฺจนฺตปาปกมฺมนฺตา,}{ปจฺจนฺติ น จ มิยฺยเร\csromansharedfootnote{3334732f3063e03f}{มียเร (ก)}.}}\csromangathastanza{\csromangathabat{ชาตเวทสโม กาโย,}{เตสํ นิรยวาสินํ.} \\ \csromangathabat{ปสฺส กมฺมานํ ทฬฺหตฺตํ,}{น ภสฺมา โหติ นปิ มสิ.}}\csromangathastanza{\csromangathabat{ปุรตฺถิเมนปิ ธาวนฺติ,}{ตโต ธาวนฺติ ปจฺฉิมํ.} \\ \csromangathabat{อุตฺตเรนปิ ธาวนฺติ,}{ตโต ธาวนฺติ ทกฺขิณํ.}}\csromangathastanza{\csromangathabat{ยํ ยํ ทิสํ ปธาวนฺติ,}{ตํ ตํ ทฺวารํ ปิธียติ.} \\ \csromangathabat{อภินิกฺขมิตาสา เต,}{สตฺตา โมกฺขคเวสิโน.}}\csromangathastanza{\csromangathabat{น เต ตโต นิกฺขมิตุํ,}{ลภนฺติ กมฺมปจฺจยา.} \\ \csromangathabat{เตสญฺจ ปาปกมฺมนฺตํ,}{อวิปกฺกํ กตํ พหุนฺติ.}}}

\prose{เอตํ ภยํ ทุกฺขํ โทมนสฺสํ กุโต ชาตํ, ตสฺส สฺเนหปจฺจยา จ นนฺทิปจฺจยา จ ราคปจฺจยา จ นนฺทิ\-ราคปจฺจยา จ ชาตํ.}

\prose{ยานิ จ เนรยิกานิ ทุกฺขานิ ยานิ จ ติรจฺฉาน\-โยนิ\-กานิ ทุกฺขานิ ยานิ จ เปตฺติวิสยิ\-กานิ ทุกฺขานิ ยานิ จ มานุสิกานิ ทุกฺขานิ, ตานิ กุโต ชาตานิ กุโต สญฺชาตานิ กุโต นิพฺพตฺตานิ กุโต อภินิพฺพตฺตานิ กุโต ปาตุภูตานิ. ตสฺส สฺเนหปจฺจยา จ นนฺทิปจฺจยา จ ราคปจฺจยา จ นนฺทิ\-ราคปจฺจยา จ ภวนฺติ สมฺภวนฺติ ชายนฺติ สญฺชายนฺติ นิพฺพตฺตนฺติ อภินิพฺพตฺตนฺติ ปาตุภวนฺตีติ สฺเนหนฺวยํ ทุกฺขมิทํ ปโหติ.}

\prose{\csromanfolio{239}\textbf{อาทีนวํ สฺเนหชํ เปกฺขมาโน}ติ สฺเนโหติ เทฺว สฺเนหา ตณฺหา\textbf{สฺเนโห} จ ทิฏฺฐิสฺเนโห จ ฯเปฯ อยํ ตณฺหาสฺเนโห ฯเปฯ อยํ ทิฏฺฐิสฺเนโห. \textbf{อาทีนวํ สฺเนหชํ เปกฺขมาโน}ติ ตณฺหาสฺเนโห จ ทิฏฺฐิสฺเนโห จ อาทีนวํ สฺเนหชํ เปกฺขมาโน ทกฺขมาโน โอโลกยมาโน นิชฺฌายมาโน อุปปริกฺขมาโนติ อาทีนวํ สฺเนหชํ เปกฺขมาโน, เอโก จเร ขคฺควิสาณ\-กปฺโป. เตนาห โส ปจฺเจก\-สมฺพุทฺโธ\csromandash{}}

\setlength{\csromangathapagewidth}{0pt}
\setlength{\csromangathawakwidth}{0pt}
\csromangathameasuregroup{}{สํสคฺคชาตสฺส ภวนฺติ สฺเนหา, \\ สฺเนหนฺวยํ ทุกฺขมิทํ ปโหติ.}
\csromangathameasurewak{สํสคฺคชาตสฺส ภวนฺติ สฺเนหา, \\ สฺเนหนฺวยํ ทุกฺขมิทํ ปโหติ.}
\setlength{\csromangathaleftcol}{0pt}
\csromangathagroup{\csromangathastanza{\csromangathawak{สํสคฺค\-ชาตสฺส ภวนฺติ สฺเนหา, \\ สฺเนหนฺวยํ ทุกฺขมิทํ ปโหติ.}}}

\prose{อาทีนวํ สฺเนหชํ เปกฺขมาโน,}

\prose{เอโก จเร ขคฺควิสาณ\-กปฺโปติ.}

\setlength{\csromangathapagewidth}{0pt}
\setlength{\csromangathawakwidth}{0pt}
\csromangathameasuregroup{}{\textbf{มิตฺเต สุหชฺเช อนุกมฺปมาโน,} \\ \textbf{หาเปติ อตฺถํ ปฏิพทฺธ}\textbf{จิตฺโต}\textsuperscript{1}\textbf{.} \\ \textbf{เอตํ ภยํ สนฺถเว}\textsuperscript{1}\textbf{ เปกฺขมาโน,} \\ \textbf{เอโก จเร ขคฺควิสาณ}\textbf{กปฺโป.}}
\csromangathameasurewak{\textbf{มิตฺเต สุหชฺเช อนุกมฺปมาโน,} \\ \textbf{หาเปติ อตฺถํ ปฏิพทฺธ}\textbf{จิตฺโต}\textsuperscript{1}\textbf{.} \\ \textbf{เอตํ ภยํ สนฺถเว}\textsuperscript{1}\textbf{ เปกฺขมาโน,} \\ \textbf{เอโก จเร ขคฺควิสาณ}\textbf{กปฺโป.}}
\setlength{\csromangathaleftcol}{0pt}
\csromangathagroup{\csromangathastanza{\csromangathawak{\csromangathaitemline{123}{\textbf{มิตฺเต สุหชฺเช อนุกมฺปมาโน,}} \\ \csromangathacontline{\textbf{หาเปติ อตฺถํ ปฏิพทฺธ}\-\textbf{จิตฺโต}\csromansharedfootnote{33ff988ecfa0187c}{ปฏิพนฺธจิตฺโต (ก)}\textbf{.}} \\ \csromangathacontline{\textbf{เอตํ ภยํ สนฺถเว}\csromansharedfootnote{28e298e9c065f376}{สนฺธเว (ก)}\textbf{ เปกฺขมาโน,}} \\ \csromangathacontline{\textbf{เอโก จเร ขคฺควิสาณ}\-\textbf{กปฺโป.}}}}}

\prose{\textbf{มิตฺเต สุหชฺเช อนุกมฺปมาโน, หาเปติ อตฺถํ ปฏิพทฺธ}\-\textbf{จิตฺโต}ติ \textbf{มิตฺตา}ติ เทฺว มิตฺตา อคาริกมิตฺโต จ อนาคาริกมิตฺโต\csromansharedfootnote{d5f37b1142d9dc49}{ปพฺพชิตมิตฺโต (ก) เอวมุปริปิ.} จ. กตโม อคาริกมิตฺโต, อิเธกจฺโจ ทุทฺททํ ททาติ, ทุจฺจชํ จชติ, ทุกฺกรํ กโรติ, ทุกฺขมํ ขมติ, คุยฺหมสฺส อาจิกฺขติ, คุยฺหมสฺส ปริคูหติ\csromansharedfootnote{27a551ce7e839add}{ปริคุยฺหติ (สฺยา, ก)}, อาปทาสุ น วิชหติ, ชีวิตมฺปิสฺส อตฺถาย ปริจฺจตฺตํ โหติ, ขีเณ นาติมญฺญติ. อยํ อคาริกมิตฺโต.}

\prose{กตโม อนาคาริกมิตฺโต, อิธ ภิกฺขุ ปิโย จ โหติ มนาโป จ ครุ จ ภาวนีโย จ วตฺตา จ วจนกฺขโม จ คมฺภีรญฺจ กถํ กตฺตา โน จ อฏฺฐาเน นิโยเชติ, อธิสีเล สมาทเปติ, จตุนฺนํ สติปฏฺฐานานํ ภาวนานุโยเค สมาทเปติ, จตุนฺนํ สมฺม\-ปฺปธานานํ ฯเปฯ จตุนฺนํ อิทฺธิปาทานํ. ปญฺจนฺนํ อินฺทฺริยานํ. ปญฺจนฺนํ พลานํ. สตฺตนฺนํ โพชฺฌงฺคานํ. อริยสฺส อฏฺฐงฺคิกสฺส มคฺคสฺส ภาวนานุโยเค สมาทเปติ. อยํ อนาคาริกมิตฺโต.}

\prose{\csromanfolio{240}สุหชฺชา วุจฺจนฺติ เยหิ สห อาคมนํ ผาสุ, คมนํ ผาสุ, ฐานํ ผาสุ, นิสชฺชนํ\csromansharedfootnote{f4ca47ecaaab4341}{นิสชฺชา (ก)} ผาสุ, สยนํ ผาสุ, อาลปนํ ผาสุ, สลฺลปนํ ผาสุ, อุลฺลปนํ ผาสุ, สมุลฺลปนํ ผาสุ. \textbf{มิตฺเต สุหชฺเช อนุกมฺปมาโน หาเปติ อตฺถนฺ}ติ มิตฺเต จ สุหชฺเช จ สนฺทิฏฺเฐ จ สมฺภตฺเต จ สหาเย จ อนุกมฺปมาโน อนุเปกฺขมาโน อนุคยฺหมาโน อตฺตตฺถมฺปิ ปรตฺถมฺปิ อุภยตฺถมฺปิ หาเปติ, ทิฏฺฐ\-ธมฺมิกมฺปิ อตฺถํ หาเปติ, สมฺปรายิกมฺปิ อตฺถํ หาเปติ, ปรมตฺถมฺปิ หาเปติ ปหาเปติ ปริหาเปติ ปริธํเสติ ปริวชฺเชติ\csromansharedfootnote{cf3e802be774915c}{ปริสชฺเชติ (สฺยา)} อนฺตรธาเปตีติ มิตฺเต สุหชฺเช อนุกมฺปมาโน หาเปติ อตฺถํ.}

\prose{\textbf{ปฏิพทฺธ}\-\textbf{จิตฺโต}ติ ทฺวีหิ การเณหิ ปฏิพทฺธ\-จิตฺโต โหติ อตฺตานํ วา นีจํ ฐเปนฺโต ปรํ อุจฺจํ ฐเปนฺโต ปฏิพทฺธ\-จิตฺโต โหติ, อตฺตานํ วา อุจฺจํ ฐเปนฺโต ปรํ นีจํ ฐเปนฺโต ปฏิพทฺธ\-จิตฺโต โหติ. กถํ อตฺตานํ นีจํ ฐเปนฺโต ปรํ อุจฺจํ ฐเปนฺโต ปฏิพทฺธ\-จิตฺโต โหติ, ตุเมฺห เม พหูปการา, อหํ ตุเมฺห นิสฺสาย ลภามิ จีวร ปิณฺฑปาต เสนาสนคิลาน ปจฺจย เภสชฺช ปริกฺขารํ, ยมฺปิ\csromansharedfootnote{f05c9d60d27377f4}{เยปิ (ก) เอวํ 293 ปิฏฺเฐปิ.} เม อญฺเญ ทาตุํ วา กาตุํ วา มญฺญนฺติ ตุเมฺห นิสฺสาย ตุเมฺห สมฺปสฺสนฺตา, ยมฺปิ เม โปราณํ มาตาเปตฺติกํ นามโคตฺตํ, ตมฺปิ เม อนฺตรหิตํ, ตุเมฺหหิ อหํ ญายามิ “อสุกสฺส กุลุปโก อสุกาย กุลุปโก”ติ. เอวํ อตฺตานํ นีจํ ฐเปนฺโต ปรํ อุจฺจํ ฐเปนฺโต ปฏิพทฺธ\-จิตฺโต โหติ.}

\prose{กถํ อตฺตานํ อุจฺจํ ฐเปนฺโต ปรํ นีจํ ฐเปนฺโต ปฏิพทฺธ\-จิตฺโต โหติ, อหํ ตุมฺหากํ พหูปกาโร, ตุเมฺห มํ อาคมฺม พุทฺธํ สรณํ คตา, ธมฺมํ สรณํ คตา, สํฆํ สรณํ คตา, ปาณาติปาตา ปฏิวิรตา, อทินฺนาทานา ปฏิวิรตา, กาเมสุ\-มิจฺฉาจารา ปฏิวิรตา, มุสาวาทา ปฏิวิรตา, สุรา\-เมรย\-มชฺช\-ปมาทฏฺฐานา ปฏิวิรตา, อหํ ตุมฺหากํ อุทฺเทสํ เทมิ, ปริปุจฺฉํ เทมิ, อุโปสถํ อาจิกฺขามิ, นวกมฺมํ อธิฏฺฐามิ, อถ ปน ตุเมฺห มํ อุชฺฌิตฺวา\csromansharedfootnote{1c2d6dd278085d7a}{ปริจฺจชิตฺวา (สฺยา)} อญฺเญ สกฺกโรถ ครุํ กโรถ มาเนถ ปูเชถาติ เอวํ อตฺตานํ อุจฺจํ ฐเปนฺโต ปรํ นีจํ ฐเปนฺโต ปฏิพทฺธ\-จิตฺโต โหตีติ มิตฺเต สุหชฺเช อนุกมฺปมาโน, หาเปติ อตฺถํ ปฏิพทฺธ\-จิตฺโต.}

\prose{\textbf{เอตํ ภยํ สนฺถเว เปกฺขมาโน}ติ \textbf{ภยนฺ}ติ ชาติภยํ ชราภยํ พฺยาธิภยํ มรณภยํ ราชภยํ โจรภยํ อคฺคิภยํ อุทกภยํ}

\prose{\csromanfolio{241}อตฺตานุวาทภยํ ปรานุวาทภยํ ทณฺฑภยํ ทุคฺคติภยํ อูมิภยํ กุมฺภิลภยํ อาวฏฺฏภยํ สุสุมารภยํ\csromansharedfootnote{00fc95dcb5552371}{สุํสุมารภยํ (สฺยา)} อาชีวิกภยํ อสิโลกภยํ ปริส\-สารชฺชภยํ มทนภยํ ภยานกํ ฉมฺภิตตฺตํ โลมหํโส เจตโส อุพฺเพโค อุตฺราโส. \textbf{สนฺถเว}ติ เทฺว สนฺถวา ตณฺหาสนฺถโว จ ทิฏฺฐิสนฺถโว จ ฯเปฯ อยํ ตณฺหาสนฺถโว ฯเปฯ อยํ ทิฏฺฐิสนฺถโว. \textbf{เอตํ ภยํ สนฺถเว เปกฺขมาโน}ติ เอตํ ภยํ สนฺถเว เปกฺขมาโน ทกฺขมาโน โอโลกยมาโน นิชฺฌายมาโน อุปปริกฺขมาโนติ เอตํ ภยํ สนฺถเว เปกฺขมาโน, เอโก จเร ขคฺควิสาณ\-กปฺโป. เตนาห โส ปจฺเจก\-สมฺพุทฺโธ\csromandash{}}

\setlength{\csromangathapagewidth}{0pt}
\setlength{\csromangathawakwidth}{0pt}
\csromangathameasuregroup{}{มิตฺเต สุหชฺเช อนุกมฺปมาโน, \\ หาเปติ อตฺถํ ปฏิพทฺธจิตฺโต. \\ เอตํ ภยํ สนฺถเว เปกฺขมาโน, \\ เอโก จเร ขคฺควิสาณกปฺโปติ. \\ \textbf{วํโส วิสาโลว ยถา วิสตฺโต,} \\ \textbf{ปุตฺเตสุ ทาเรสุ จ ยา อเปกฺขา.} \\ \textbf{วํสกฺกฬีโรว}\textsuperscript{1}\textbf{ อสชฺชมาโน,} \\ \textbf{เอโก จเร ขคฺควิสาณ}\textbf{กปฺโป.}}
\csromangathameasurewak{มิตฺเต สุหชฺเช อนุกมฺปมาโน, \\ หาเปติ อตฺถํ ปฏิพทฺธจิตฺโต. \\ เอตํ ภยํ สนฺถเว เปกฺขมาโน, \\ เอโก จเร ขคฺควิสาณกปฺโปติ. \\ \textbf{วํโส วิสาโลว ยถา วิสตฺโต,} \\ \textbf{ปุตฺเตสุ ทาเรสุ จ ยา อเปกฺขา.} \\ \textbf{วํสกฺกฬีโรว}\textsuperscript{1}\textbf{ อสชฺชมาโน,} \\ \textbf{เอโก จเร ขคฺควิสาณ}\textbf{กปฺโป.}}
\setlength{\csromangathaleftcol}{0pt}
\csromangathagroup{\csromangathastanza{\csromangathawak{มิตฺเต สุหชฺเช อนุกมฺปมาโน, \\ หาเปติ อตฺถํ ปฏิพทฺธ\-จิตฺโต. \\ เอตํ ภยํ สนฺถเว เปกฺขมาโน, \\ เอโก จเร ขคฺควิสาณ\-กปฺโปติ.}}\csromangathastanza{\csromangathawak{\csromangathaitemline{124}{\textbf{วํโส วิสาโลว ยถา วิสตฺโต,}} \\ \csromangathacontline{\textbf{ปุตฺเตสุ ทาเรสุ จ ยา อเปกฺขา.}} \\ \csromangathacontline{\textbf{วํสกฺกฬีโรว}\csromansharedfootnote{0a5870e586a91330}{วํเส กฬีโรว (ก)}\textbf{ อสชฺชมาโน,}} \\ \csromangathacontline{\textbf{เอโก จเร ขคฺควิสาณ}\-\textbf{กปฺโป.}}}}}

\prose{\textbf{วํโส วิสาโลว ยถา วิสตฺโต}ติ วํโส วุจฺจติ เวฬุคุมฺโพ. ยถา เวฬุคุมฺพสฺมิํ โปราณกา วํสา สตฺตา\csromansharedfootnote{435914665d7fae80}{เวฬุคุมฺพสฺมิํ กณฺฏกา ชฏิตา สํสิพฺพิตา (สฺยา)} วิสตฺตา อาสตฺตา ลคฺคา ลคฺคิตา ปลิพุทฺธา, เอวเมว วิสตฺติกา วุจฺจติ ตณฺหา. โย ราโค สาราโค อนุนโย อนุโรโธ นนฺที นนฺทิราโค จิตฺตสฺส สาราโค อิจฺฉา มุจฺฉา อชฺโฌสานํ เคโธ ปลิเคโธ สงฺโค ปงฺโก เอชา มายา ชนิกา สญฺชนนี สิพฺพินี ชาลินี สริตา วิสตฺติกา สุตฺตํ วิสฏา อายูหนี ทุติยา ปณิธิ ภวเนตฺติ วนํ วนโถ สนฺถโว สิเนโห อเปกฺขา ปฏิพนฺธุ อาสา อาสีสนา อาสีสิตตฺตํ รูปาสา สทฺทาสา คนฺธาสา รสาสา โผฏฺฐพฺพาสา ลาภาสา ธนาสา ปุตฺตาสา ชีวิตาสา ชปฺปา ปชปฺปา อภิชปฺปา ชปฺปนา ชปฺปิตตฺตํ โลลุปฺปํ โลลุปฺปายนา โลลุปฺปา\-ยิตตฺตํ ปุจฺฉญฺชิกตา สาธุกมฺยตา อธมฺมราโค วิสมโลโภ นิกนฺติ นิกามนา ปตฺถนา ปิหนา สมฺปตฺถนา กามตณฺหา \csromanfolio{242}ภวตณฺหา วิภวตณฺหา รูปตณฺหา อรูปตณฺหา นิโรธตณฺหา รูปตณฺหา สทฺทตณฺหา คนฺธตณฺหา รสตณฺหา โผฏฺฐพฺพ\-ตณฺหา ธมฺมตณฺหา โอโฆ โยโค คนฺโถ อุปาทานํ อาวรณํ นีวรณํ ฉทนํ พนฺธนํ อุปกฺกิเลโส อนุสโย ปริยุฏฺฐานํ ลตา เววิจฺฉํ ทุกฺขมูลํ ทุกฺขนิทานํ ทุกฺขปฺปภโว มารปาโส มารพฬิสํ มารวิสโย มารนิวาโส มารพนฺธนํ ตณฺหานที ตณฺหาชาลํ ตณฺหาคทฺทุลํ ตณฺหาสมุทฺโท อภิชฺฌา โลโภ อกุสลมูลํ.}

\prose{\textbf{วิสตฺติกา}ติ เกนฏฺเฐน วิสตฺติกา, วิสาลาติ วิสตฺติกา, วิสตาติ วิสตฺติกา, วิสฏาติ วิสตฺติกา, วิสมาติ วิสตฺติกา, วิสกฺกตีติ วิสตฺติกา, วิสํหรตีติ วิสตฺติกา, วิสํวาทิกาติ วิสตฺติกา, วิสมูลาติ วิสตฺติกา, วิสผลาติ วิสตฺติกา, วิสปริโภคาติ วิสตฺติกา, วิสาลา วา ปน ตณฺหา รูเป สทฺเท คนฺเธ รเส โผฏฺฐพฺเพ กุเล คเณ อาวาเส ลาเภ ยเส ปสํสาย สุเข จีวเร ปิณฺฑปาเต เสนาสเน คิลานปจฺจย\-เภสชฺชปริกฺขาเร กามธาตุยา รูปธาตุยา อรูปธาตุยา กามภเว รูปภเว อรูปภเว สญฺญาภเว อสญฺญาภเว เนว\-สญฺญา\-นา\-สญฺญาภเว เอกโวการภเว จตุโวการภเว ปญฺจโวการภเว อตีเต อนาคเต ปจฺจุปฺปนฺเน ทิฏฺฐ\-สุต\-มุต\-วิญฺญาตพฺเพสุ ธมฺเมสุ วิสตา วิตฺถตาติ วิสตฺติกาติ วํโส วิสาโลว ยถา วิสตฺโต.}

\prose{\textbf{ปุตฺเตสุ ทาเรสุ จ ยา อเปกฺขา}ติ \textbf{ปุตฺตา}ติ จตฺตาโร ปุตฺตา อตฺรโช ปุตฺโต เขตฺตโช ปุตฺโต ทินฺนโก ปุตฺโต อนฺเตวาสิโก ปุตฺโต. \textbf{ทารา} วุจฺจนฺติ ภริยาโย. \textbf{อเปกฺขา} วุจฺจนฺติ ตณฺหา. โย ราโค สาราโค ฯเปฯ อภิชฺฌา โลโภ อกุสลมูลนฺติ ปุตฺเตสุ ทาเรสุ จ ยา อเปกฺขา.}

\prose{\textbf{วํสกฺกฬีโรว อสชฺชมาโน}ติ วํโส วุจฺจติ เวฬุคุมฺโพ. ยถา เวฬุคุมฺพสฺมิํ ตรุณกา กฬีรกา\csromansharedfootnote{c9a8c683e8b6977e}{ตรุณกฬีรา (สฺยา)} อสตฺตา อลคฺคา อคธิตา\csromansharedfootnote{ae62d467cb6d692e}{อปลิเวฏฺฐิตา (สฺยา)} อปลิพุทฺธา นิกฺขนฺตา นิสฺสฏา วิปฺปมุตฺตา เอวเมว. \textbf{สชฺชา}ติ เทฺว สชฺชนา ตณฺหาสชฺชนา จ ทิฏฺฐิสชฺชนา จ ฯเปฯ อยํ ตณฺหาสชฺชนา ฯเปฯ อยํ ทิฏฺฐิสชฺชนา. ตสฺส}

\prose{\csromanfolio{243}ปจฺเจก\-สมฺพุทฺธสฺส ตณฺหาสชฺชนา ปหีนา, ทิฏฺฐิสชฺชนา ปฏินิสฺสฏฺฐา. ตณฺหาสชฺชนาย ปหีนตฺตา ทิฏฺฐิ\-สชฺชนาย ปฏินิสฺสฏฺฐ\-ตฺตา โส ปจฺเจก\-สมฺพุทฺโธ รูเป น สชฺชติ. สทฺเท น สชฺชติ. คนฺเธ น สชฺชติ. รเส น สชฺชติ. โผฏฺฐพฺเพ น สชฺชติ. กุเล. คเณ. อาวาเส. ลาเภ. ยเส. ปสํสาย. สุเข. จีวเร. ปิณฺฑปาเต. เสนาสเน. คิลานปจฺจย\-เภสชฺชปริกฺขาเร. กามธาตุยา. รูปธาตุยา. อรูปธาตุยา. กามภเว. รูปภเว. อรูปภเว. สญฺญาภเว. อสญฺญาภเว. เนว\-สญฺญา\-นา\-สญฺญาภเว. เอกโวการภเว. จตุโวการภเว. ปญฺจโวการภเว. อตีเต. อนาคเต. ปจฺจุปฺปนฺเน. ทิฏฺฐ\-สุต\-มุต\-วิญฺญาตพฺเพสุ ธมฺเมสุ น สชฺชติ น คณฺหาติ น พชฺฌติ น ปลิพชฺฌติ น มุจฺฉติ นิกฺขนฺโต นิสฺสโฏ วิปฺปมุตฺโต วิสญฺญุตฺโต วิมริยาทิกเตน เจตสา วิหรตีติ วํสกฺกฬีโรว อสชฺชมาโน, เอโก จเร ขคฺควิสาณ\-กปฺโป. เตนาห โส ปจฺเจก\-สมฺพุทฺโธ\csromandash{}}

\setlength{\csromangathapagewidth}{0pt}
\setlength{\csromangathawakwidth}{0pt}
\csromangathameasuregroup{\textbf{มิโค อรญฺญมฺหิ ยถา อพทฺโธ}\textsuperscript{1}\textbf{,}}{วํโส วิสาโลว ยถา วิสตฺโต, \\ ปุตฺเตสุ ทาเรสุ จ ยา อเปกฺขา. \\ วํสกฺกฬีโรว อสชฺชมาโน, \\ เอโก จเร ขคฺควิสาณกปฺโปติ. \\ \csromangathabat{\textbf{มิโค อรญฺญมฺหิ ยถา อพทฺโธ}\textsuperscript{1}\textbf{,}}{\textbf{เยนิจฺฉกํ คจฺฉติ โคจราย.}} \\ \textbf{วิญฺญู นโร เสริตํ เปกฺขมาโน.} \\ \textbf{เอโก จเร ขคฺควิสาณ}\textbf{กปฺโป.}}
\csromangathameasurewak{วํโส วิสาโลว ยถา วิสตฺโต, \\ ปุตฺเตสุ ทาเรสุ จ ยา อเปกฺขา. \\ วํสกฺกฬีโรว อสชฺชมาโน, \\ เอโก จเร ขคฺควิสาณกปฺโปติ.}
\csromangathasetleft{\textbf{มิโค อรญฺญมฺหิ ยถา อพทฺโธ}\textsuperscript{1}\textbf{,}}
\csromangathagroup{\csromangathastanza{\csromangathawak{วํโส วิสาโลว ยถา วิสตฺโต, \\ ปุตฺเตสุ ทาเรสุ จ ยา อเปกฺขา. \\ วํสกฺกฬีโรว อสชฺชมาโน, \\ เอโก จเร ขคฺควิสาณ\-กปฺโปติ.}}\csromangathastanza{\csromangathaitembat{125}{\textbf{มิโค อรญฺญมฺหิ ยถา อพทฺโธ}\csromansharedfootnote{5119c6d1ac61d290}{อพนฺโธ (สฺยา, ก)}\textbf{,}}{\textbf{เยนิจฺฉกํ คจฺฉติ โคจราย.}} \\ \csromangathacontline{\textbf{วิญฺญู นโร เสริตํ เปกฺขมาโน.}} \\ \csromangathacontline{\textbf{เอโก จเร ขคฺควิสาณ}\-\textbf{กปฺโป.}}}}

\prose{\textbf{มิโค อรญฺญมฺหิ ยถา อพทฺโธ, เยนิจฺฉกํ คจฺฉติ โคจรายา}ติ \textbf{มิโค}ติ เทฺว มิคา เอณิมิโค จ ปสทมิโค จ. ยถา อารญฺญิโก\csromansharedfootnote{9ac6775075b9ed08}{อารญฺญโก (สฺยา, ก)} มิโค อรญฺเญ ปวเน จรมาโน\csromansharedfootnote{2905eda195e2c813}{อรญฺเญ วสมาโน (สฺยา)} วิสฺสตฺโถ คจฺฉติ, วิสฺสตฺโถ ติฏฺฐติ, วิสฺสตฺโถ นิสีทติ, วิสฺสตฺโถ เสยฺยํ กปฺเปติ.}

\prose{วุตฺตํเหตํ ภควตา\csromandash{}เสยฺยถาปิ ภิกฺขเว อารญฺญิโก มิโค อรญฺเญ ปวเน จรมาโน วิสฺสตฺโถ คจฺฉติ, วิสฺสตฺโถ ติฏฺฐติ, วิสฺสตฺโถ นิสีทติ, วิสฺสตฺโถ เสยฺยํ กปฺเปติ. ตํ กิสฺส เหตุ, อนาปาถคโต ภิกฺขเว ลุทฺทสฺส. เอวเมว โข ภิกฺขเว ภิกฺขุ \csromanfolio{244}วิวิจฺเจว กาเมหิ วิวิจฺจ อกุสเลหิ ธมฺเมหิ สวิตกฺกํ สวิจารํ วิเวกชํ ปีติสุขํ ปฐมํ ฌานํ อุปสมฺปชฺช วิหรติ. อยํ วุจฺจติ ภิกฺขเว ภิกฺขุ อนฺธมกาสิ มารํ, อปทํ\csromansharedfootnote{53bcf530441efe3f}{ปรํ (ก) ม 1. 214 ปิฏฺเฐ.} วธิตฺวา มารจกฺขุํ อทสฺสนํ คโต ปาปิมโต.}

\prose{ปุน จปรํ ภิกฺขเว ภิกฺขุ วิตกฺก\-วิจารานํ วูปสมา อชฺฌตฺตํ สมฺปสาทนํ เจตโส เอโกทิภาวํ อวิตกฺกํ อวิจารํ สมาธิชํ ปีติสุขํ ทุติยํ ฌานํ อุปสมฺปชฺช วิหรติ. อยํ วุจฺจติ ภิกฺขเว ภิกฺขุ อนฺธมกาสิ มารํ, อปทํ วธิตฺวา มารจกฺขุํ อทสฺสนํ คโต ปาปิมโต.}

\prose{ปุน จปรํ ภิกฺขเว ภิกฺขุ ปีติยา จ วิราคา อุเปกฺขโก จ วิหรติ, สโต จ สมฺปชาโน สุขญฺจ กาเยน ปฏิสํเวเทติ. ยํ ตํ อริยา อาจิกฺขนฺติ “อุเปกฺขโก สติมา สุขวิหารี”ติ ตติยํ ฌานํ อุปสมฺปชฺช วิหรติ, อยํ วุจฺจติ ภิกฺขเว ภิกฺขุ อนฺธมกาสิ มารํ, อปทํ วธิตฺวา มารจกฺขุํ อทสฺสนํ คโต ปาปิมโต.}

\prose{ปุน จปรํ ภิกฺขเว ภิกฺขุ สุขสฺส จ ปหานา ทุกฺขสฺส จ ปหานา ปุพฺเพว โสมนสฺส\-โทมนสฺสานํ อตฺถงฺคมา อทุกฺขมสุขํ อุเปกฺขา\-สติ\-ปาริสุทฺธิํ จตุตฺถํ ฌานํ อุปสมฺปชฺช วิหรติ, อยํ วุจฺจติ ภิกฺขเว ภิกฺขุ อนฺธมกาสิ มารํ, อปทํ วธิตฺวา มารจกฺขุํ อทสฺสนํ คโต ปาปิมโต.}

\prose{ปุน จปรํ ภิกฺขเว ภิกฺขุ สพฺพโส รูปสญฺญานํ สมติกฺกมา ปฏิฆ\-สญฺญานํ อตฺถงฺคมา นานตฺต\-สญฺญานํ อมนสิการา “อนนฺโต อากาโส”ติ อากาสา\-นญฺจา\-ยตนํ อุปสมฺปชฺช วิหรติ, อยํ วุจฺจติ ภิกฺขเว ภิกฺขุ อนฺธมกาสิ มารํ, อปทํ วธิตฺวา มารจกฺขุํ อทสฺสนํ คโต ปาปิมโต.}

\prose{ปุน จปรํ ภิกฺขเว ภิกฺขุ สพฺพโส อากาสา\-นญฺจา\-ยตนํ สมติกฺกมฺม “อนนฺตํ วิญฺญาณนฺ”ติ วิญฺญาณญฺจา\-ยตนํ อุปสมฺปชฺช วิหรติ ฯเปฯ.}

\prose{ปุน จปรํ ภิกฺขเว สพฺพโส วิญฺญาณญฺจา\-ยตนํ สมติกฺกมฺม “นตฺถิ กิญฺจี”ติ อากิญฺจญฺญา\-ยตนํ อุปสมฺปชฺช วิหรติ ฯเปฯ.}

\prose{ปุน จปรํ ภิกฺขเว สพฺพโส อากิญฺจญฺญา\-ยตนํ สมติกฺกมฺม เนว\-สญฺญา\-นา\-สญฺญา\-ยตนํ อุปสมฺปชฺช วิหรติ ฯเปฯ.}

\prose{\csromanfolio{245}ปุน จปรํ ภิกฺขเว สพฺพโส เนว\-สญฺญา\-นา\-สญฺญา\-ยตนํ สมติกฺกมฺม สญฺญาเวทยิต\-นิโรธํ อุปสมฺปชฺช วิหรติ. ปญฺญาย จสฺส ทิสฺวา อาสวา ปริกฺขีณา โหนฺติ. อยํ วุจฺจติ ภิกฺขเว ภิกฺขุ อนฺธมกาสิ มารํ, อปทํ วธิตฺวา มารจกฺขุํ อทสฺสนํ คโต ปาปิมโต, ติณฺโณ โลเก วิสตฺติกํ, โส วิสฺสตฺโถ คจฺฉติ, วิสฺสตฺโถ ติฏฺฐติ, วิสฺสตฺโถ นิสีทติ, วิสฺสตฺโถ เสยฺยํ กปฺเปติ. ตํ กิสฺส เหตุ, อนาปาถคโต ภิกฺขเว ปาปิมโตติ มิโค อรญฺญมฺหิ ยถา อพทฺโธ, เยนิจฺฉกํ คจฺฉติ โคจราย.}

\prose{\textbf{วิญฺญู นโร เสริตํ เปกฺขมาโน}ติ \textbf{วิญฺญู}ติ ปณฺฑิโต ปญฺญวา พุทฺธิมา ญาณี วิภาวี เมธาวี. \textbf{นโร}ติ สตฺโต มานโว โปโส ปุคฺคโล ชีโว ชาคุ ชนฺตุ อินฺทคุ มนุโช. \textbf{เสรี}ติ เทฺว เสรี ธมฺโมปิ เสรี ปุคฺคโลปิ เสรี. กตโม ธมฺโม เสรี, จตฺตาโร สติปฏฺฐานา จตฺตาโร สมฺมปฺปธานา จตฺตาโร อิทฺธิปาทา ปญฺจินฺทฺริยานิ ปญฺจ พลานิ สตฺต โพชฺฌงฺคา อริโย อฏฺฐงฺคิโก มคฺโค. อยํ ธมฺโม เสรี. กตโม ปุคฺคโล เสรี. โย อิมินา เสรินา ธมฺเมน สมนฺนาคโต. โส วุจฺจติ ปุคฺคโล เสรี. \textbf{วีญฺญู นโร เสริตํ เปกฺขมาโน}ติ วิญฺญู นโร เสริตํ ธมฺมํ เปกฺขมาโน ทกฺขมาโน โอโลกยมาโน นิชฺฌายมาโน อุปปริกฺขมาโนติ วิญฺญู นโร เสริตํ เปกฺขมาโน, เอโก จเร ขคฺควิสาณ\-กปฺโป. เตนาห โส ปจฺเจก\-สมฺพุทฺโธ\csromandash{}}

\setlength{\csromangathapagewidth}{0pt}
\setlength{\csromangathawakwidth}{0pt}
\csromangathameasuregroup{}{มิโค อรญฺญมฺหิ ยถา อพทฺโธ, \\ เยนิจฺฉกํ คจฺฉติ โคจราย.}
\csromangathameasurewak{มิโค อรญฺญมฺหิ ยถา อพทฺโธ, \\ เยนิจฺฉกํ คจฺฉติ โคจราย.}
\setlength{\csromangathaleftcol}{0pt}
\csromangathagroup{\csromangathastanza{\csromangathawak{มิโค อรญฺญมฺหิ ยถา อพทฺโธ, \\ เยนิจฺฉกํ คจฺฉติ โคจราย.}}}

\prose{วิญฺญู นโร เสริตํ เปกฺขมาโน,}

\prose{เอโก จเร ขคฺควิสาณ\-กปฺโปติ.}

\setlength{\csromangathapagewidth}{0pt}
\setlength{\csromangathawakwidth}{0pt}
\csromangathameasuregroup{}{\textbf{อามนฺตนา โหติ สหายมชฺเฌ,} \\ \textbf{วาเส ฐาเน คมเน จาริกาย.} \\ \textbf{อนภิชฺฌิตํ เสริตํ เปกฺขมาโน,} \\ \textbf{เอโก จเร ขคฺควิสาณ}\textbf{กปฺโป.}}
\csromangathameasurewak{\textbf{อามนฺตนา โหติ สหายมชฺเฌ,} \\ \textbf{วาเส ฐาเน คมเน จาริกาย.} \\ \textbf{อนภิชฺฌิตํ เสริตํ เปกฺขมาโน,} \\ \textbf{เอโก จเร ขคฺควิสาณ}\textbf{กปฺโป.}}
\setlength{\csromangathaleftcol}{0pt}
\csromangathagroup{\csromangathastanza{\csromangathawak{\csromangathaitemline{126}{\textbf{อามนฺตนา โหติ สหายมชฺเฌ,}} \\ \csromangathacontline{\textbf{วาเส ฐาเน คมเน จาริกาย.}} \\ \csromangathacontline{\textbf{อนภิชฺฌิตํ เสริตํ เปกฺขมาโน,}} \\ \csromangathacontline{\textbf{เอโก จเร ขคฺควิสาณ}\-\textbf{กปฺโป.}}}}}

\prose{\textbf{อามนฺตนา โหติ สหายมชฺเฌ, วาเส ฐาเน คมเน จาริกายา}ติ สหายา วุจฺจนฺติ เยหิ สห อาคมนํ ผาสุ, คมนํ ผาสุ, คมนาคมนํ ผาสุ, ฐานํ ผาสุ, นิสชฺชนํ ผาสุ, สยนํ ผาสุ, อาลปนํ \csromanfolio{246}ผาสุ, สลฺลปนํ ผาสุ, อุลฺลปนํ ผาสุ, สมุลฺลปนํ ผาสุ. \textbf{อามนฺตนา โหติ สหายมชฺเฌ, วาเส ฐาเน คมเน จาริกายา}ติ สหายมชฺเฌ วาเสปิ ฐาเนปิ คมเนปิ จาริกายปิ อตฺตตฺถ\-มนฺตนา ปรตฺถ\-มนฺตนา อุภยตฺถ\-มนฺตนา ทิฏฺฐธมฺมิก\-ตฺถ\-มนฺตนา สมฺปรายิกตฺถ\-มนฺตนา ปรมตฺถ\-มนฺตนา\csromansharedfootnote{5bcef0cf577dc7ea}{อุภยตฺถมนฺตนา (ก)}ติ อามนฺตนา โหติ สหายมชฺเฌ, วาเส ฐาเน คมเน จาริกาย.}

\prose{\textbf{อนภิชฺฌิตํ เสริตํ เปกฺขมาโน}ติ อนภิชฺฌิตํ เอตํ วตฺถุ พาลานํ อสปฺปุริสานํ ติตฺถิยานํ ติตฺถิย\-สาวกานํ, ยทิทํ ภณฺฑุ\-กาสายวตฺถวสนตา. อภิชฺฌิตํ เอตํ วตฺถุ ปณฺฑิตานํ สปฺปุริสานํ พุทฺธ\-สาวกานํ ปจฺเจก\-พุทฺธานํ, ยทิทํ ภณฺฑุ\-กาสายวตฺถวสนตา. \textbf{เสรี}ติ เทฺว เสรี ธมฺโมปิ เสรี ปุคฺคโลปิ เสรี. กตโม ธมฺโม เสรี. จตฺตาโร สติปฏฺฐานา ฯเปฯ อริโย อฏฺฐงฺคิโก มคฺโค, อยํ ธมฺโม เสรี. กตโม ปุคฺคโล เสรี. โย อิมินา เสรินา ธมฺเมน สมนฺนาคโต, โส วุจฺจติ ปุคฺคโล เสรี. \textbf{อนภิชฺฌิตํ เสริตํ เปกฺขมาโน}ติ เสริตํ ธมฺมํ เปกฺขมาโน ทกฺขมาโน โอโลกยมาโน นิชฺฌายมาโน อุปปริกฺขมาโนติ อนภิชฺฌิตํ เสริตํ เปกฺขมาโน, เอโก จเร ขคฺควิสาณ\-กปฺโป. เตนาห โส ปจฺเจก\-สมฺพุทฺโธ\csromandash{}}

\setlength{\csromangathapagewidth}{0pt}
\setlength{\csromangathawakwidth}{0pt}
\csromangathameasuregroup{}{อามนฺตนา โหติ สหายมชฺเฌ, \\ วาเส ฐาเน คมเน จาริกาย.}
\csromangathameasurewak{อามนฺตนา โหติ สหายมชฺเฌ, \\ วาเส ฐาเน คมเน จาริกาย.}
\setlength{\csromangathaleftcol}{0pt}
\csromangathagroup{\csromangathastanza{\csromangathawak{อามนฺตนา โหติ สหายมชฺเฌ, \\ วาเส ฐาเน คมเน จาริกาย.}}}

\prose{อนภิชฺฌิตํ เสริตํ เปกฺขมาโน,}

\prose{เอโก จเร ขคฺควิสาณ\-กปฺโปติ.}

\setlength{\csromangathapagewidth}{0pt}
\setlength{\csromangathawakwidth}{0pt}
\csromangathameasuregroup{}{ขิฑฺฑา\textsuperscript{1} \textbf{รตี โหติ สหายมชฺเฌ,} \\ \textbf{ปุตฺเตสุ จ วิปุลํ โหติ เปมํ.} \\ \textbf{ปิยวิปฺปโยคํ}\textsuperscript{1}\textbf{ วิชิคุจฺฉมาโน,} \\ \textbf{เอโก จเร ขคฺควิสาณ}\textbf{กปฺโป.}}
\csromangathameasurewak{ขิฑฺฑา\textsuperscript{1} \textbf{รตี โหติ สหายมชฺเฌ,} \\ \textbf{ปุตฺเตสุ จ วิปุลํ โหติ เปมํ.} \\ \textbf{ปิยวิปฺปโยคํ}\textsuperscript{1}\textbf{ วิชิคุจฺฉมาโน,} \\ \textbf{เอโก จเร ขคฺควิสาณ}\textbf{กปฺโป.}}
\setlength{\csromangathaleftcol}{0pt}
\csromangathagroup{\csromangathastanza{\csromangathawak{\csromangathaitemline{127}{ขิฑฺฑา\csromansharedfootnote{83f744067dd76b9e}{ขิฏฺฏา (ก)} \textbf{รตี โหติ สหายมชฺเฌ,}} \\ \csromangathacontline{\textbf{ปุตฺเตสุ จ วิปุลํ โหติ เปมํ.}} \\ \csromangathacontline{\textbf{ปิยวิปฺปโยคํ}\csromansharedfootnote{2a597c0ecbbf901c}{วิปฺปโยคมฺปิ (ก)}\textbf{ วิชิคุจฺฉมาโน,}} \\ \csromangathacontline{\textbf{เอโก จเร ขคฺควิสาณ}\-\textbf{กปฺโป.}}}}}

\prose{\textbf{ขิฑฺฑา รตี โหติ สหายมชฺเฌ}ติ ขิฑฺฑาติ เทฺว \textbf{ขิฑฺฑา} กายิกา จ ขิฏฺฏา วาจสิกา จ ขิฑฺฑา. กตมา กายิกา ขิฑฺฑา. หตฺถีหิปิ กีฬนฺติ, อสฺเสหิปิ กีฬนฺติ, รเถหิปิ กีฬนฺติ, ธนุหิปิ กีฬนฺติ, ถรูหิปิ กีฬนฺติ, อฏฺฐปเทหิปิ กีฬนฺติ, ทสปเทหิปิ กีฬนฺติ, อากาเสปิ กีฬนฺติ, ปริหาร\-ปเถปิ กีฬนฺติ, สนฺติกายปิ กีฬนฺติ, ขลิกายปิ กีฬนฺติ,}

\prose{\csromanfolio{247}ฆฏิกายปิ กีฬนฺติ, สลากหตฺเถนปิ กีฬนฺติ, อกฺเขนปิ กีฬนฺติ, ปงฺคจีเรนปิ กีฬนฺติ, วงฺกเกนปิ กีฬนฺติ, โมกฺขจิกายปิ กีฬนฺติ, จิงฺคุลเกนปิ กีฬนฺติ, ปตฺตาฬฺหเกนปิ กีฬนฺติ, รถเกนปิ กีฬนฺติ, ธนุเกนปิ กีฬนฺติ, อกฺขริกายปิ กีฬนฺติ, มเนสิกายปิ กีฬนฺติ, ยถาวชฺเชนปิ กีฬนฺติ, อยํ กายิกา ขิฑฺฑา.}

\prose{กตมา วาจสิกา ขิฑฺฑา. มุขเภริกํ มุขาลมฺพรํ มุขฑิณฺฑิมกํ\csromansharedfootnote{c4b5c19970101c83}{มุขเทณฺฑิมกํ (สฺยา), มุขทินฺทิมกํ (ก)} มุขจลิมกํ มุขเกรกํ\csromansharedfootnote{356821c6d01170ca}{มุขเภรุกํ (สฺยา)} มุขททฺทริกํ นาฏกํ ลาสํ คีตํ ทวกมฺมํ. อยํ วาจสิกา ขิฑฺฑา.}

\prose{\textbf{รตี}ติ อนุกฺกณฺฐิตาธิวจนเมตํ รตีติ. \textbf{สหายา} วุจฺจนฺติ เยหิ สห อาคมนํ ผาสุ, คมนํ ผาสุ, คมนาคมนํ ผาสุ, ฐานํ ผาสุ, นิสชฺชนํ ผาสุ, สยนํ ผาสุ, อาลปนํ ผาสุ, สลฺลปนํ ผาสุ, อุลฺลปนํ ผาสุ, สมุลฺลปนํ ผาสุ. \textbf{ขิฑฺฑา รตี โหติ สหายมชฺเฌ}ติ ขิฑฺฑา จ รติ จ สหายมชฺเฌ โหตีติ ขิฑฺฑา รตี โหติ สหายมชฺเฌ.}

\prose{\textbf{ปุตฺเตสุ จ วิปุลํ โหติ เปมนฺ}ติ ปุตฺตาติ จตฺตาโร \textbf{ปุตฺตา} อตฺรโช ปุตฺโต เขตฺตโช ปุตฺโต ทินฺนโก ปุตฺโต อนฺเตวาสิโก ปุตฺโต. \textbf{ปุตฺเตสุ จ วิปุลํ โหติ เปมนฺ}ติ ปุตฺเตสุ จ อธิมตฺตํ โหติ เปมนฺติ ปุตฺเตสุ จ วิปุลํ โหติ เปมํ.}

\prose{\textbf{ปิยวิปฺปโยคํ วิชิคุจฺฉ}\-\textbf{มาโน}ติ เทฺว ปิยา สตฺตา วา สงฺขารา วา. กตเม สตฺตา ปิยา. อิธ ย’สฺส เต โหนฺติ อตฺถกามา หิตกามา ผาสุกามา โยคกฺเขมกามา มาตา วา ปิตา วา ภาตา วา ภคินี วา ปุตฺโต วา ธีตา วา มิตฺตา วา อมจฺจา วา ญาตี วา สาโลหิตา วา. อิเม สตฺตา ปิยา.}

\prose{กตเม สงฺขารา ปิยา. มนาปิกา รูปา มนาปิกา สทฺทา มนาปิกา คนฺธา มนาปิกา รสา มนาปิกา โผฏฺฐพฺพา. อิเม สงฺขารา ปิยา. \textbf{ปิยวิปฺปโยคํ วิชิคุจฺฉ}\-\textbf{มาโน}ติ ปิยานํ วิปฺปโยคํ วิชิคุจฺฉมาโน อฏฺฏียมาโน หรายมาโนติ ปิยวิปฺปโยคํ วิชิคุจฺฉมาโน, เอโก จเร ขคฺควิสาณ\-กปฺโป. เตนาห โส ปจฺเจก\-สมฺพุทฺโธ\csromandash{}}

\setlength{\csromangathapagewidth}{0pt}
\setlength{\csromangathawakwidth}{0pt}
\csromangathameasuregroup{}{ขิฑฺฑา รตี โหติ สหายมชฺเฌ, \\ ปุตฺเตสุ จ วิปุลํ โหติ เปมํ. \\ ปิยวิปฺปโยคํ วิชิคุจฺฉมาโน, \\ เอโก จเร ขคฺควิสาณกปฺโปติ. \\ \textbf{จาตุทฺทิโส อปฺปฏิโฆ จ โหติ,} \\ \textbf{สนฺตุสฺสมาโน อิตรีตเรน.} \\ \textbf{ปริสฺสยานํ สหิตา อฉมฺภี,} \\ \textbf{เอโก จเร ขคฺควิสาณ}\textbf{กปฺโป.}}
\csromangathameasurewak{ขิฑฺฑา รตี โหติ สหายมชฺเฌ, \\ ปุตฺเตสุ จ วิปุลํ โหติ เปมํ. \\ ปิยวิปฺปโยคํ วิชิคุจฺฉมาโน, \\ เอโก จเร ขคฺควิสาณกปฺโปติ. \\ \textbf{จาตุทฺทิโส อปฺปฏิโฆ จ โหติ,} \\ \textbf{สนฺตุสฺสมาโน อิตรีตเรน.} \\ \textbf{ปริสฺสยานํ สหิตา อฉมฺภี,} \\ \textbf{เอโก จเร ขคฺควิสาณ}\textbf{กปฺโป.}}
\setlength{\csromangathaleftcol}{0pt}
\csromangathagroup{\csromangathastanza{\csromangathawak{\csromanfolio{248}ขิฑฺฑา รตี โหติ สหายมชฺเฌ, \\ ปุตฺเตสุ จ วิปุลํ โหติ เปมํ. \\ ปิยวิปฺปโยคํ วิชิคุจฺฉมาโน, \\ เอโก จเร ขคฺควิสาณ\-กปฺโปติ.}}\csromangathastanza{\csromangathawak{\csromangathaitemline{128}{\textbf{จาตุทฺทิโส อปฺปฏิโฆ จ โหติ,}} \\ \csromangathacontline{\textbf{สนฺตุสฺสมาโน อิตรีตเรน.}} \\ \csromangathacontline{\textbf{ปริสฺสยานํ สหิตา อฉมฺภี,}} \\ \csromangathacontline{\textbf{เอโก จเร ขคฺควิสาณ}\-\textbf{กปฺโป.}}}}}

\prose{\textbf{จาตุทฺทิโส อปฺปฏิโฆ จ โหตี}ติ \textbf{จาตุทฺทิโส}ติ โส ปจฺเจก\-สมฺพุทฺโธ เมตฺตา\-สหคเตน เจตสา เอกํ ทิสํ ผริตฺวา วิหรติ. ตถา ทุติยํ. ตถา ตติยํ. ตถา จตุตฺถํ. อิติ อุทฺธมโธ ติริยํ สพฺพธิ สพฺพตฺตตาย สพฺพาวนฺตํ โลกํ เมตฺตา\-สหคเตน เจตสา วิปุเลน มหคฺคเตน อปฺปมาเณน อเวเรน อพฺยาปชฺเชน\csromansharedfootnote{13294d9b12ee0e91}{อพฺยาปชฺเฌน (สฺยา) ปสฺส ที 3. 187 ปิฏฺเฐ.} ผริตฺวา วิหรติ. กรุณา\-สหคเตน ฯเปฯ. มุทิตา\-สหคเตน ฯเปฯ. อุเปกฺขา\-สหคเตน เจตสา เอกํ ทิสํ ผริตฺวา วิหรติ. ตถา ทุติยํ. ตถา ตติยํ ฯเปฯ อพฺยาปชฺเชน ผริตฺวา วิหรติ. \textbf{จาตุทฺทิโส อปฺปฏิโฆ จ โหตี}ติ เมตฺตาย ภาวิตตฺตา เย ปุรตฺถิมาย ทิสาย สตฺตา, เต อปฺปฏิกูลา\csromansharedfootnote{dacfe310ee0f1ebb}{อปฺปฏิกุลา (พหูสุ)} โหนฺติ, เย ทกฺขิณาย ทิสาย สตฺตา, เต อปฺปฏิกูลา โหนฺติ. เย ปจฺฉิมาย ทิสาย สตฺตา, เต อปฺปฏิกูลา โหนฺติ. เย อุตฺตราย ทิสาย สตฺตา, เต อปฺปฏิกูลา โหนฺติ. เย ปุรตฺถิมาย อนุทิสาย สตฺตา, เต อปฺปฏิกูลา โหนฺติ. เย ทกฺขิณาย อนุทิสาย สตฺตา, เต อปฺปฏิกูลา โหนฺติ. เย ปจฺฉิมาย อนุทิสาย สตฺตา, เต อปฺปฏิกูลา โหนฺติ. เย อุตฺตราย อนุทิสาย สตฺตา, เต อปฺปฏิกูลา โหนฺติ. เย เหฏฺฐิมาย ทิสาย สตฺตา, เต อปฺปฏิกูลา โหนฺติ. เย อุปริมาย ทิสาย สตฺตา, เต อปฺปฏิกูลา โหนฺติ. เย ทิสาสุ วิทิสาสุ สตฺตา, เต อปฺปฏิกูลา โหนฺติ. กรุณาย ภาวิตตฺตา. มุทิตาย ภาวิตตฺตา. อุเปกฺขาย ภาวิตตฺตา เย ปุรตฺถิมาย ทิสาย สตฺตา, เต อปฺปฏิกูลา โหนฺติ ฯเปฯ. เย ทิสาสุ วิทิสาสุ สตฺตา, เต อปฺปฏิกูลา โหนฺตีติ จาตุทฺทิโส อปฺปฏิโฆ จ โหติ.}

\prose{\csromanfolio{249}\textbf{สนฺตุสฺสมาโน อิตรีตเรนา}ติ โส ปจฺเจก\-สมฺพุทฺโธ สนฺตุฏฺโฐ โหติ อิตรีตเรน จีวเรน, อิตรีตร\-จีวร\-สนฺตุฏฺฐิยา จ วณฺณวาที, น จ จีวรเหตุ อเนสนํ อปฺปติรูปํ\csromansharedfootnote{2daa5533648ae6ab}{อปฺปฏิรูปํ (สฺยา)} อาปชฺชติ. อลทฺธา จ จีวรํ น ปริตสฺสติ, ลทฺธา จ จีวรํ อคธิโต อมุจฺฉิโต อนชฺโฌสนฺโน\csromansharedfootnote{e611116e8ad840aa}{อนชฺโฌปนฺโน (สฺยา)} อาทีนวทสฺสาวี นิสฺสรณปญฺโญ ปริภุญฺชติ. ตาย จ ปน อิตรีตร\-จีวร\-สนฺตุฏฺฐิยา เนว\-ตฺตานุกฺกํเสติ, น ปรํ วมฺเภติ, โย หิ ตตฺถ ทกฺโข อนลโส สมฺปชาโน ปติสฺสโต. อยํ วุจฺจติ ปจฺเจก\-สมฺพุทฺโธ โปราเณ อคฺคญฺเญ อริยวํเส ฐิโต. สนฺตุฏฺโฐ โหติ อิตรีตเรน ปิณฺฑปาเตน ฯเปฯ. สนฺตุฏฺโฐ โหติ อิตรีตเรน เสนาสเนน ฯเปฯ. สนฺตุฏฺโฐ โหติ อิตรีตเรน คิลานปจฺจย\-เภสชฺช\-ปริกฺขาเรน. อิตรีตร\-คิลานปจฺจย\-เภสชฺช- ปริกฺขาร\-สนฺตุฏฺฐิยา จ วณฺณวาที, น จ คิลานปจฺจย\-เภสชฺช- ปริกฺขารเหตุ อเนสนํ อปฺปติรูปํ อาปชฺชติ. อลทฺธา จ คิลานปจฺจย\-เภสชฺช\-ปริกฺขารานํ น ปริตสฺสติ, ลทฺธา จ คิลานปจฺจย\-เภสชฺช\-ปริกฺขารํ อคธิโต อมุจฺฉิโต อนชฺโฌสนฺโน อาทีนวทสฺสาวี นิสฺสรณปญฺโญ ปริภุญฺชติ. ตาย จ อิตรีตร\-คิลานปจฺจย\-เภสชฺชปริกฺขาร\-สนฺตุฏฺฐิยา เนว\-ตฺตานุกฺกํเสติ, น ปรํ วมฺเภติ, โย หิ ตตฺถ ทกฺโข อนลโส สมฺปชาโน ปติสฺสโต. อยํ วุจฺจติ ปจฺเจก\-สมฺพุทฺโธ โปราเณ อคฺคญฺเญ อริยวํเส ฐิโตติ สนฺตุสฺสมาโน อิตรีตเรน.}

\prose{\textbf{ปริสฺสยานํ สหิตา อฉมฺภี}ติ ปริสฺสยาติ เทฺว \textbf{ปริสฺสยา} ปากฏ\-ปริสฺสยา จ ปฏิจฺฉนฺน\-ปริสฺสยา จ. กตเม \textbf{ปากฏ}\-\textbf{ปริสฺสยา.} สีหา พฺยคฺฆาทีปี อจฺฉา ตรจฺฉา โกกา มหิํสา\csromansharedfootnote{0bbe9667ea9a7851}{โคมหิสา (สฺยา) ขุ 7. 10 ปิฏฺเฐสุปิ.} หตฺถี อหี วิจฺฉิกา สตปที โจรา วา อสฺสุ มานวา วา กตกมฺมา วา อกตกมฺมา วา, จกฺขุโรโค โสตโรโค ฆานโรโค ชิวฺหาโรโค กายโรโค สีสโรโค กณฺณโรโค มุขโรโค ทนฺตโรโค กาโส สาโส ปินาโส ฑาโห ชโร กุจฺฉิโรโค มุจฺฉา ปกฺขนฺทิกา สูลา วิสูจิกา กุฏฺฐํ คณฺโฑ กิลาโส โสโส อปมาโร ททฺทุ กณฺฑุ กจฺฉุ รขสา วิตจฺฉิกา โลหิตปิตฺตํ มธุเมโห อํสา ปิฬกา ภคนฺทลา ปิตฺต\-สมุฏฺฐานา อาพาธา, เสมฺห\-สมุฏฺฐานา อาพาธา, วาตสมุฏฺฐานา อาพาธา, สนฺนิปาติกา อาพาธา, อุตุปริณามชา อาพาธา,}

\prose{\csromanfolio{250}วิสม\-ปริหารชา อาพาธา, โอปกฺกมิกา อาพาธา กมฺมวิปากชา อาพาธา, สีตํ อุณฺหํ ชิฆจฺฉา ปิปาสา อุจฺจาโร ปสฺสาโว ฑํสมกส\-วาตา\-ตป\-สรีสป\-สมฺผสฺสา อิติ วา, อิเม วุจฺจนฺติ ปากฏปริสฺสาย.}

\prose{กตเม \textbf{ปฏิจฺฉนฺน}\-\textbf{ปริสฺสยา.} กายทุจฺจริตํ วจีทุจฺจริตํ มโนทุจฺจริตํ กามจฺฉนฺท\-นีวรณํ พฺยาปาท\-นีวรณํ ถินมิทฺธ\-นีวรณํ อุทฺธจฺจ\-กุกฺกุจฺจ\-นีวรณํ วิจิกิจฺฉา\-นีวรณํ ราโค โทโส โมโห โกโธ อุปนาโห มกฺโข ปฬาโส อิสฺสา มจฺฉริยํ มายา สาเฐยฺยํ ถมฺโภ สารมฺโภ มาโน อติมาโน มโท ปมาโท สพฺเพ กิเลสา สพฺเพ ทุจฺจริตา สพฺเพ ทรถา สพฺเพ ปริฬาหา สพฺเพ สนฺตาปา สพฺพา\-กุสลา\-ภิสงฺขารา, อิเม วุจฺจนฺติ ปฏิจฺฉนฺน\-ปริสฺสยา.}

\prose{\textbf{ปริสฺสยา}ติ เกนฏฺเฐน ปริสฺสยา, ปริสหนฺตีติ ปริสฺสยา, ปริหานาย สํวตฺตนฺตีติ ปริสฺสยา, ตตฺราสยาติ ปริสฺสยา. กถํ ปริสหนฺตีติ ปริสฺสยา, เต ปริสฺสยา ตํ ปุคฺคลํ สหนฺติ ปริสหนฺติ อภิภวนฺติ อชฺโฌตฺถรนฺติ ปริยาทิยนฺติ มทฺทนฺติ, เอวํ ปริสหนฺตีติ ปริสฺสยา.}

\prose{กถํ ปริหานาย สํวตฺตนฺตีติ ปริสฺสยา, เต ปริสฺสยา กุสลานํ ธมฺมานํ อนฺตรายาย ปริหานาย สํวตฺตนฺติ, กตเมสํ กุสลานํ ธมฺมานํ, สมฺมา\-ปฏิปทาย อนุโลม\-ปฏิปทาย อปจฺจนีกปฏีปทาย อนฺวตฺถ\-ปฏิปทาย ธมฺมานุธมฺม\-ปฏิปทาย สีเลสุ ปริปูร\-การิ\-ตาย อินฺทฺริเยสุ คุตฺตทฺวารตาย โภชเน มตฺตญฺญุตาย ชาคริยานุโยคสฺส สติสมฺปชญฺญสฺส จตุนฺนํ สติปฏฺฐานานํ ภาวนา\-นุโยคสฺส จตุนฺนํ สมฺม\-ปฺปธานานํ. จตุนฺนํ อิทฺธิปาทานํ. ปญฺจนฺนํ อินฺทฺริยานํ. ปญฺจนฺนํ พลานํ. สตฺตนฺนํ โพชฺฌงฺคานํ. อริยสฺส อฏฺฐงฺคิกสฺส มคฺคสฺส ภาวนา\-นุโยคสฺส อิเมสํ กุสลานํ ธมฺมานํ อนฺตรายาย ปริหานาย สํวตฺตนฺติ, เอวํ ปริหานาย สํวตฺตนฺตีติ ปริสฺสยา.}

\prose{กถํ ตตฺราสยาติ ปริสฺสยา, ตตฺเถเต\csromansharedfootnote{83f6805b026944ff}{ตตฺร เต (ก)} ปาปกา อกุสลา ธมฺมา อุปฺปชฺชนฺติ อตฺตภาว\-สนฺนิสฺสยา, ยถา พิเล พิลาสยา ปาณา สยนฺติ, ทเก ทกาสยา\csromansharedfootnote{d39fdd4344a13c57}{อุทเก อุทกาสยา (สฺยา)} ปาณา สยนฺติ, วเน วนาสยา ปาณา สยนฺติ, รุกฺเข รุกฺขาสยา ปาณา สยนฺติ. เอวเมว ตตฺเถเต ปาปกา อกุสลา ธมฺมา อุปฺปชฺชนฺติ อตฺตภาว\-สนฺนิสฺสยาติ, เอวมฺปิ ตตฺรสยาติ ปริสฺสยา.}

\prose{\csromanfolio{251}วุตฺตํเหตํ ภควตา\csromandash{}สานฺเตวาสิโก ภิกฺขเว ภิกฺขุ สาจริยโก ทุกฺขํ น ผาสุ วิหรติ. กถญฺจ ภิกฺขเว ภิกฺขุ สานฺเตวาสิโก สาจริยโก ทุกฺขํ น ผาสุ วิหรติ. อิธ ภิกฺขเว จกฺขุนา รูปํ ทิสฺวา อุปฺปชฺชนฺติ เย ปาปกา อกุสลา ธมฺมา สรสงฺกปฺปา สญฺโญชนียา\csromansharedfootnote{a9a8330869c83bed}{สญฺโญชนิกา (ก)}, ตฺยสฺส อนฺโต วสนฺติ อนฺวาสฺสวนฺติ ปาปกา อกุสลา ธมฺมาติ ตสฺมา สานฺเตวาสิโก วุจฺจติ, เตน สมุทาจเรน สมุทาจรนฺติ นํ “ปาปกา อกุสลา ธมฺมา”ติ ตสฺมา “สาจริยโก”ติ วุจฺจติ.}

\prose{ปุน จปรํ ภิกฺขเว ภิกฺขุโน โสเตน สทฺทํ สุตฺวา ฯเปฯ ฆาเนน คนฺธํ ฆายิตฺวา ฯเปฯ ชิวฺหาย รสํ สายิตฺวา ฯเปฯ กาเยน โผฏฺฐพฺพํ ผุสิตฺวา ฯเปฯ มนสา ธมฺมํ วิญฺญาย อุปฺปชฺชนฺติ เย ปาปกา อกุสลา ธมฺมา สรสงฺกปฺปา สญฺโญชนียา. ตฺยสฺส อนฺโต วสนฺติ อนฺวาสฺสวนฺติ ปาปกา อกุสลา ธมฺมาติ ตสฺมา “สานฺเตวาสิโก”ติ วุจฺจติ. เตน สมุทาจเรน สมุทาจรนฺติ นํ “ปาปกา อกุสลา ธมฺมา”ติ ตสฺมา “สาจริยโก”ติ วุจฺจติ. เอวํ โข ภิกฺขเว ภิกฺขุ สานฺเตวาสิโก สาจริยโก ทุกฺขํ น ผาสุ วิหรตีติ. เอวมฺปิ ตตฺราสยาติ ปริสฺสยา.}

\prose{วุตฺตํ เหตํ ภควตา\csromandash{}ตโยเม ภิกฺขเว อนฺตรามลา อนฺตราอมิตฺตา อนฺตราสปตฺตา อนฺตราวธกา อนฺตรา\-ปจฺจตฺถิกา. กตเม ตโย. โลโภ ภิกฺขเว อนฺตรามโล\csromansharedfootnote{8f900410c5621126}{อนฺตรามลํ (สฺยา, ก) ปสฺส ขุ 1. 252 ปิฏฺเฐ.} อนฺตราอมิตฺโต อนฺตราสปตฺโต อนฺตราวธโก อนฺตรา\-ปจฺจตฺถิโก, โทโส ภิกฺขเว ฯเปฯ โมโห ภิกฺขเว อนฺตรามโล อนฺตรา- อมิตฺโต อนฺตราสปตฺโต อนฺตราวธโก อนฺตรา\-ปจฺจตฺถิโก. อิเม โข ภิกฺขเว ตโย อนฺตรามลา อนฺตราอมิตฺตา อนฺตาราสปตฺตา อนฺตราวธกา อนฺตรา\-ปจฺจตฺถิกาติ.}

\setlength{\csromangathapagewidth}{0pt}
\setlength{\csromangathawakwidth}{0pt}
\csromangathameasuregroup{อนตฺถชนโน โลโภ, \\ ภยมนฺตรโต ชาตํ, \\ ลุทฺโธ อตฺถํ น ชานาติ, \\ อนฺธตมํ ตทา โหติ, \\ อนตฺถชนโน โทโส, \\ ภยมนฺตรโต ชาตํ, \\ ทุฏฺโฐ อตฺถํ น ชานาติ, \\ อนฺธตมํ ตทา โหติ, \\ อนตฺถชนโน โมโห, \\ ภยมนฺตรโต ชาตํ, \\ มูโฬฺห อตฺถํ น ชานาติ, \\ อนฺธตมํ ตทา โหติ,}{\csromangathabat{อนตฺถชนโน โลโภ,}{โลโภ จิตฺตปฺปโกปโน.} \\ \csromangathabat{ภยมนฺตรโต ชาตํ,}{ตํ ชโน นาวพุชฺฌติ.} \\ \csromangathabat{ลุทฺโธ อตฺถํ น ชานาติ,}{ลุทฺโธ ธมฺมํ น ปสฺสติ.} \\ \csromangathabat{อนฺธตมํ ตทา โหติ,}{ยํ โลโภ สหเต นรํ.} \\ \csromangathabat{อนตฺถชนโน โทโส,}{โทโส จิตฺตปฺปโกปโน.} \\ \csromangathabat{ภยมนฺตรโต ชาตํ,}{ตํ ชโน นาวพุชฺฌติ.} \\ \csromangathabat{ทุฏฺโฐ อตฺถํ น ชานาติ,}{ทุฏฺโฐ ธมฺมํ น ปสฺสติ.} \\ \csromangathabat{อนฺธตมํ ตทา โหติ,}{ยํ โทโส สหเต นรํ.} \\ \csromangathabat{อนตฺถชนโน โมโห,}{โมโห จิตฺตปฺปโกปโน.} \\ \csromangathabat{ภยมนฺตรโต ชาตํ,}{ตํ ชโน นาวพุชฺฌติ.} \\ \csromangathabat{มูโฬฺห อตฺถํ น ชานาติ,}{มูโฬฺห ธมฺมํ น ปสฺสติ.} \\ \csromangathabat{อนฺธตมํ ตทา โหติ,}{ยํ โมโห สหเต นรนฺติ.}}
\csromangathasetleft{อนตฺถชนโน โลโภ, \\ ภยมนฺตรโต ชาตํ, \\ ลุทฺโธ อตฺถํ น ชานาติ, \\ อนฺธตมํ ตทา โหติ, \\ อนตฺถชนโน โทโส, \\ ภยมนฺตรโต ชาตํ, \\ ทุฏฺโฐ อตฺถํ น ชานาติ, \\ อนฺธตมํ ตทา โหติ, \\ อนตฺถชนโน โมโห, \\ ภยมนฺตรโต ชาตํ, \\ มูโฬฺห อตฺถํ น ชานาติ, \\ อนฺธตมํ ตทา โหติ,}
\csromangathagroup{\csromangathastanza{\csromangathabat{อนตฺถชนโน โลโภ,}{โลโภ จิตฺตปฺปโกปโน.} \\ \csromangathabat{ภยมนฺตรโต ชาตํ,}{ตํ ชโน นาวพุชฺฌติ.}}\csromangathastanza{\csromangathabat{ลุทฺโธ อตฺถํ น ชานาติ,}{ลุทฺโธ ธมฺมํ น ปสฺสติ.} \\ \csromangathabat{อนฺธตมํ ตทา โหติ,}{ยํ โลโภ สหเต นรํ.}}\csromangathastanza{\csromangathabat{อนตฺถชนโน โทโส,}{โทโส จิตฺตปฺปโกปโน.} \\ \csromangathabat{ภยมนฺตรโต ชาตํ,}{ตํ ชโน นาวพุชฺฌติ.}}\csromangathastanza{\csromanfolio{252}\csromangathabat{ทุฏฺโฐ อตฺถํ น ชานาติ,}{ทุฏฺโฐ ธมฺมํ น ปสฺสติ.} \\ \csromangathabat{อนฺธตมํ ตทา โหติ,}{ยํ โทโส สหเต นรํ.}}\csromangathastanza{\csromangathabat{อนตฺถชนโน โมโห,}{โมโห จิตฺตปฺปโกปโน.} \\ \csromangathabat{ภยมนฺตรโต ชาตํ,}{ตํ ชโน นาวพุชฺฌติ.}}\csromangathastanza{\csromangathabat{มูโฬฺห อตฺถํ น ชานาติ,}{มูโฬฺห ธมฺมํ น ปสฺสติ.} \\ \csromangathabat{อนฺธตมํ ตทา โหติ,}{ยํ โมโห สหเต นรนฺติ.}}}

\prose{เอวมฺปิ ตตฺราสยาติ ปริสฺสยา.}

\prose{วุตฺตเญฺหตํ ภควตา\csromandash{}ตโย โข มหาราช ปุริสสฺส ธมฺมา อชฺฌตฺตํ อุปฺปชฺชมานา อุปฺปชฺชนฺติ อหิตาย ทุกฺขาย อผาสุวิหาราย. กตเม ตโย, โลโภ โข มหาราช ปุริสสฺส ธมฺโม อชฺฌตฺตํ อุปฺปชฺชมาโน อุปฺปชฺชติ อหิตาย ทุกฺขาย อผาสุวิหาราย, โทโส โข มหาราช ฯเปฯ โมโห โข มหาราช ปุริสสฺส ธมฺโม อชฺฌตฺตํ อุปฺปชฺชมาโน อุปฺปชฺชติ อหิตาย ทุกฺขาย อผาสุวิหาราย, อิเม โข มหาราช ตโย ปุริสสฺส ธมฺมา อชฺฌตฺตํ อุปฺปชฺชมานา อุปฺปชฺชนฺติ อหิตาย ทุกฺขาย อผาสุวิหาราย.}

\setlength{\csromangathapagewidth}{0pt}
\setlength{\csromangathawakwidth}{0pt}
\csromangathameasuregroup{โลโภ โทโส จ โมโห จ, \\ หิํสนฺติ อตฺตสมฺภูตา,}{\csromangathabat{โลโภ โทโส จ โมโห จ,}{ปุริสํ ปาปเจตสํ.} \\ \csromangathabat{หิํสนฺติ อตฺตสมฺภูตา,}{ตจสารํว สมฺผลนฺติ\textsuperscript{1}.}}
\csromangathasetleft{โลโภ โทโส จ โมโห จ, \\ หิํสนฺติ อตฺตสมฺภูตา,}
\csromangathagroup{\csromangathastanza{\csromangathabat{โลโภ โทโส จ โมโห จ,}{ปุริสํ ปาปเจตสํ.} \\ \csromangathabat{หิํสนฺติ อตฺตสมฺภูตา,}{ตจสารํว สมฺผลนฺติ\csromansharedfootnote{9a0648cdcee826d1}{สผลนฺติ (ก) ขุ 1. 226 ปิฏฺเฐ.}.}}}

\prose{เอวมฺปิ ตตฺราสยาติ ปริสฺสยา.}

\prose{วุตฺตเญฺหตํ ภควตา\csromandash{}}

\setlength{\csromangathapagewidth}{0pt}
\setlength{\csromangathawakwidth}{0pt}
\csromangathameasuregroup{}{ราโค จ โทโส จ อิโตนิทานา, \\ อรตี รตี โลมหํโส อิโตชา. \\ อิโต สมุฏฺฐาย มโนวิตกฺกา, \\ กุมารกา ธงฺกมิโวสฺสชนฺตีติ\textsuperscript{1}.}
\csromangathameasurewak{ราโค จ โทโส จ อิโตนิทานา, \\ อรตี รตี โลมหํโส อิโตชา. \\ อิโต สมุฏฺฐาย มโนวิตกฺกา, \\ กุมารกา ธงฺกมิโวสฺสชนฺตีติ\textsuperscript{1}.}
\setlength{\csromangathaleftcol}{0pt}
\csromangathagroup{\csromangathastanza{\csromangathawak{ราโค จ โทโส จ อิโตนิทานา, \\ อรตี รตี โลมหํโส อิโตชา. \\ อิโต สมุฏฺฐาย มโนวิตกฺกา, \\ กุมารกา ธงฺกมิ\-โว\-สฺสชนฺตีติ\csromansharedfootnote{40cb1247d295cbfb}{ธงฺกมิโวสฺสชฺชนฺตีติ (สฺยา, ก) ปสฺส ขุ 1. 320 ปิฏฺเฐ.}.}}}

\prose{เอวมฺปิ ตตฺราสยาติ ปริสฺสยา.}

\prose{\textbf{ปริสฺสยานํ สหิตา}ติ ปริสฺสเย สหิตา อาราธิตา อชฺโฌตฺถริตา ปริยาทิตา ปฏินิสฺสตาติ ปริสฺสยานํ สหิตา. \textbf{อฉมฺภี}ติ โส}

\prose{\csromanfolio{253}ปจฺเจก\-สมฺพุทฺโธ อภีรู อจฺฉมฺภี อนุตฺราสี อปลายี ปหีน\-ภยเภรโว วิคต\-โลมหํโส วิหรตีติ ปริสฺสยานํ สหิตา อฉมฺภี, เอโก จเร ขคฺควิสาณ\-กปฺโป. เตนาห โส ปจฺเจก\-สมฺพุทฺโธ\csromandash{}}

\setlength{\csromangathapagewidth}{0pt}
\setlength{\csromangathawakwidth}{0pt}
\csromangathameasuregroup{}{จาตุทฺทิโส อปฺปฏิโฆ จ โหติ, \\ สนฺตุสฺสมาโน อิตรีตเรน. \\ ปริสฺสยานํ สหิตา อฉมฺภี, \\ เอโก จเร ขคฺควิสาณกปฺโปติ.}
\csromangathameasurewak{จาตุทฺทิโส อปฺปฏิโฆ จ โหติ, \\ สนฺตุสฺสมาโน อิตรีตเรน. \\ ปริสฺสยานํ สหิตา อฉมฺภี, \\ เอโก จเร ขคฺควิสาณกปฺโปติ.}
\setlength{\csromangathaleftcol}{0pt}
\csromangathagroup{\csromangathastanza{\csromangathawak{จาตุทฺทิโส อปฺปฏิโฆ จ โหติ, \\ สนฺตุสฺสมาโน อิตรีตเรน. \\ ปริสฺสยานํ สหิตา อฉมฺภี, \\ เอโก จเร ขคฺควิสาณ\-กปฺโปติ.}}}

\proseitem{129}{\textbf{ทุสฺสงฺคหา ปพฺพชิตาปิ เอเก,}}

\prose{\textbf{อโถ คหฏฺฐา ฆรมาวสนฺตา.}}

\prose{\textbf{อปฺโปสฺสุกฺโก ปรปุตฺเตสุ หุตฺวา,}}

\prose{\textbf{เอโก จเร ขคฺควิสาณ}\-\textbf{กปฺโป.} (9)}

\prose{\textbf{ทุสฺสงฺคหา ปพฺพชิตาปิ เอเก}ติ ปพฺพชิตาปิ อิเธกจฺเจ นิสฺสเยปิ ทิยฺยมาเน อุทฺเทเสปิ ทิยฺยมาเน ปริปุจฺฉายปิ\csromansharedfootnote{fdf04176d7b2e12c}{ปริปุจฺเฉปิ (ก)} ทิยฺยมาเน จีวเรปิ ทิยฺยมาเน ปตฺเตปิ ทิยฺยมาเน โลหถาลเกปิ ทิยฺยมาเน ธมฺมกรเณปิ\csromansharedfootnote{98c29d4a56f5f691}{ธมฺมกรเกปิ (สฺยา)} ทิยฺยมาเน ปริสฺสาวเนปิ ทิยฺยมาเน ถวิเกปิ ทิยฺยมาเน อุปาหเนปิ ทิยฺยมาเน กายพนฺธเนปิ ทิยฺยมาเน น สุณนฺติ, น โสตํ โอทหนฺติ, น อญฺญา จิตฺตํ อุปฏฺฐเปนฺติ, อนสฺสวา อวจนกรา ปฏิโลม\-วุตฺติโน อญฺเญเนว มุขํ กโรนฺตีติ ทุสฺสงฺคหา ปพฺพชิตาปิ เอเก.}

\prose{\textbf{อโถ คหฏฺฐา ฆรมาวสนฺตา}ติ คหฏฺฐาปิ อิเธกจฺเจ หตฺถิมฺหิปิ ทิยฺยมาเน ฯเปฯ. รเถปิ. เขตฺเตปิ. วตฺถุมฺหิปิ. หิรญฺเญปิ. สุวณฺเณปิ ทิยฺยมาเน คาเมปิ. นิคเมปิ. นคเรปิ. รฏฺเฐปิ. ชนปเทปิ ทิยฺยมาเน น สุณนฺติ น โสตํ โอทหนฺติ, น อญฺญาจิตฺตํ อุปฏฺฐเปนฺติ, อนสฺสวา อวจนกรา ปฏิโลม\-วุตฺติโน อญฺเญเนว มุขํ กโรนฺตีติ อโถ คหฏฺฐา ฆรมาวสนฺตา.}

\prose{\textbf{อปฺโปสฺสุกฺโก ปรปุตฺเตสุ หุตฺวา}ติ อตฺตานํ ฐเปตฺวา สพฺเพ อิมสฺมิํ อตฺเถ ปรปุตฺตา เตสุ ปรปุตฺเตสุ อปฺโปสฺสุกฺโก หุตฺวา อพฺยาวโฏ หุตฺวา อนเปกฺโข หุตฺวาติ อปฺโปสฺสุกฺโก ปรปุตฺเตสุ หุตฺวา, เอโก จเร ขคฺควิสาณ\-กปฺโป. เตนาห โส ปจฺเจก\-สมฺพุทฺโธ\csromandash{}}

\setlength{\csromangathapagewidth}{0pt}
\setlength{\csromangathawakwidth}{0pt}
\csromangathameasuregroup{}{ทุสฺสงฺคหา ปพฺพชิตาปิ เอเก, \\ อโถ คหฏฺฐา ฆรมาวสนฺตา. \\ อปฺโปสฺสุกฺโก ปรปุตฺเตสุ หุตฺวา, \\ เอโก จเร ขคฺควิสาณกปฺโปติ.}
\csromangathameasurewak{ทุสฺสงฺคหา ปพฺพชิตาปิ เอเก, \\ อโถ คหฏฺฐา ฆรมาวสนฺตา. \\ อปฺโปสฺสุกฺโก ปรปุตฺเตสุ หุตฺวา, \\ เอโก จเร ขคฺควิสาณกปฺโปติ.}
\setlength{\csromangathaleftcol}{0pt}
\csromangathagroup{\csromangathastanza{\csromangathawak{\csromanfolio{254}ทุสฺสงฺคหา ปพฺพชิตาปิ เอเก, \\ อโถ คหฏฺฐา ฆรมาวสนฺตา. \\ อปฺโปสฺสุกฺโก ปรปุตฺเตสุ หุตฺวา, \\ เอโก จเร ขคฺควิสาณ\-กปฺโปติ.}}}

\proseitem{130}{โอโรปยิตฺวา\csromansharedfootnote{0564f64f067ff7df}{โวโรปยิตฺวา (สฺยา)} \textbf{คิหิพฺยญฺชนานิ,}}

\prose{\textbf{สญฺฉินฺนปตฺโต}\csromansharedfootnote{a4628aae05eb63ed}{สํสีนปตฺโต (สีฏฺฐ)}\textbf{ ยถา โกวิฬาโร.}}

\prose{\textbf{เฉตฺวาน วีโร คิหิพนฺธนานิ,}}

\prose{\textbf{เอโก จเร ขคฺควิสาณ}\-\textbf{กปฺโป.} (10)}

\prose{\textbf{โอโรปยิตฺวา คิหิพฺยญฺชนานี}ติ คิหิพฺยญฺชนานิ วุจฺจนฺติ เกสา จ มสฺสู จ มาลา จ คนฺธญฺจ วิเลปนญฺจ อาภรณญฺจ ปิลนฺธนญฺจ วตฺถญฺจ ปารุปนญฺจ เวฐนญฺจ อุจฺฉาทนํ ปริมทฺทนํ นฺหาปนํ\csromansharedfootnote{6a1ada551fb69100}{นหาปนํ (สฺยา)} สมฺพาหนํ อาทาสํ อญฺชนํ มาลา\-คนฺธ\-วิเลปนํ มุขจุณฺณํ มุขเลปนํ หตฺถพนฺธํ สิขาพนฺธํ ทณฺฑํ นาฬิกํ\csromansharedfootnote{fa1932d0527cd1dd}{นาลิกํ (ก) ปสฺส ที 1. 7 ปิฏฺเฐ.} ขคฺคํ ฉตฺตํ จิตฺรูปาหนํ อุณฺหีสํ มณิํ วาลพีชนิํ โอทาตานิ วตฺถานิ ทีฆทสานิ\csromansharedfootnote{ef8edf8aaa127b01}{ทีฆรสฺสานิ (สฺยา)} อิติ วา. \textbf{โอโรปยิตฺวา คิหิพฺยญฺชนานี}ติ คิหิพฺยญฺชนานิ โอโรปยิตฺวา สโมโรปยิตฺวา นิกฺขิปิตฺวา ปฏิปสฺสมฺภยิตฺวาติ โอโรปยิตฺวา คิหิพฺยญฺชนานิ.}

\prose{\textbf{สญฺฉินฺนปตฺโต ยถา โกวิฬาโร}ติ ยถา โกวิฬารสฺส ปตฺตานิ ฉินฺนานิ สญฺฉินฺนานิ ปติตานิ ปริปติตานิ, เอวเมว ตสฺส ปจฺเจก\-สมฺพุทฺธสฺส คิหิพฺยญฺชนานิ ฉินฺนานิ สญฺฉินฺนานิ ปติตานีติ สญฺฉินฺนปตฺโต ยถา โกวิฬาโร.}

\prose{\textbf{เฉตฺวาน วีโร คิหิพนฺธนานี}ติ \textbf{วีโร}ติ วีริยวาติ วีโร, ปหูติ วีโร, วิสวีติ วีโร, อลมตฺโตติ วีโร, สูโรติ วีโร, วิกฺกนฺโต อภิรู อจฺฉมฺภี อนุตฺราสี อปลายี ปหีน\-ภยเภรโวติ วีโร, วิคต\-โลมหํโสติ วีโร.}

\setlength{\csromangathapagewidth}{0pt}
\setlength{\csromangathawakwidth}{0pt}
\csromangathameasuregroup{}{วิรโต อิธ\textsuperscript{1} สพฺพปาเกหิ, \\ นิรยทุกฺขํ อติจฺจ วีริยวาโส. \\ โส วีริยวา ปธานวา, \\ ธีโร\textsuperscript{1} ตาทิ ปวุจฺจเต ตถตฺตา.}
\csromangathameasurewak{วิรโต อิธ\textsuperscript{1} สพฺพปาเกหิ, \\ นิรยทุกฺขํ อติจฺจ วีริยวาโส. \\ โส วีริยวา ปธานวา, \\ ธีโร\textsuperscript{1} ตาทิ ปวุจฺจเต ตถตฺตา.}
\setlength{\csromangathaleftcol}{0pt}
\csromangathagroup{\csromangathastanza{\csromangathawak{วิรโต อิธ\csromansharedfootnote{8921d0fc02e422b1}{อารโต อิเธว (สฺยา) ปสฺส ขุ 1. 360 ปิฏฺเฐ.} สพฺพปาเกหิ, \\ นิรยทุกฺขํ อติจฺจ วีริยวาโส. \\ โส วีริยวา ปธานวา, \\ ธีโร\csromansharedfootnote{0984a0af328bfcec}{วีโร (สฺยา, ก) ปสฺส ขุ 1. 360 ปิฏฺเฐ.} ตาทิ ปวุจฺจเต ตถตฺตา.}}}

\prose{\csromanfolio{255}คิหิพนฺธนานิ วุจฺจนฺติ ปุตฺตา จ ภริยา จ ทาสา จ ทาสี จ อเชฬกา จ กุกฺกุฏสูกรา จ หตฺถิ\-ควา\-สฺส\-วฬวา จ เขตฺตญฺจ วตฺถุ จ หิรญฺญญฺจ สุวณฺณญฺจ คามนิคม\-ราชธานิโย จ รฏฺฐญฺจ ชนปโท จ โกโส จ โกฏฺฐาคารญฺจ ยํ กิญฺจิ รชนียวตฺถุ.}

\prose{\textbf{เฉตฺวาน วีโร คิหิพนฺธนานี}ติ โส ปจฺเจก\-สมฺพุทฺโธ วีโร คิหิพนฺธนานิ ฉินฺทิตฺวา สมุจฺฉินฺทิตฺวา ชหิตฺวา วิโนเทตฺวา พฺยนฺตีกริตฺวา อนภาวํ คเมตฺวาติ เฉตฺวาน วีโร คิหิพนฺธนานิ, เอโก จเร ขคฺควิสาณ\-กปฺโป. เตนาห โส ปจฺเจก\-สมฺพุทฺโธ\csromandash{}}

\setlength{\csromangathapagewidth}{0pt}
\setlength{\csromangathawakwidth}{0pt}
\csromangathameasuregroup{}{โอโรปยิตฺวา คิหิพฺยญฺชนานิ, \\ สญฺฉินฺนปตฺโต ยถา โกวิฬาโร. \\ เฉตฺวาน วีโร คิหิพนฺธนานิ, \\ เอโก จเร ขคฺควิสาณกปฺโปติ.}
\csromangathameasurewak{โอโรปยิตฺวา คิหิพฺยญฺชนานิ, \\ สญฺฉินฺนปตฺโต ยถา โกวิฬาโร. \\ เฉตฺวาน วีโร คิหิพนฺธนานิ, \\ เอโก จเร ขคฺควิสาณกปฺโปติ.}
\setlength{\csromangathaleftcol}{0pt}
\csromangathagroup{\csromangathastanza{\csromangathawak{โอโรปยิตฺวา คิหิพฺยญฺชนานิ, \\ สญฺฉินฺนปตฺโต ยถา โกวิฬาโร. \\ เฉตฺวาน วีโร คิหิพนฺธนานิ, \\ เอโก จเร ขคฺควิสาณ\-กปฺโปติ.}}}

\vspace{\tipitakaparskip}
\csromancenter{ปฐโม วคฺโค.}
\chapterhead{\textbf{ทุติยวคฺค}}
\setlength{\csromangathapagewidth}{0pt}
\setlength{\csromangathawakwidth}{0pt}
\csromangathameasuregroup{}{\textbf{สเจ ลเภถ นิปกํ สหายํ,} \\ \textbf{สทฺธิํ จรํ สาธุวิหาริ ธีรํ.} \\ \textbf{อภิภุยฺย สพฺพานิ ปริสฺสยานิ,} \\ \textbf{จเรยฺย เตน’ตฺตมโน สตีมา.}}
\csromangathameasurewak{\textbf{สเจ ลเภถ นิปกํ สหายํ,} \\ \textbf{สทฺธิํ จรํ สาธุวิหาริ ธีรํ.} \\ \textbf{อภิภุยฺย สพฺพานิ ปริสฺสยานิ,} \\ \textbf{จเรยฺย เตน’ตฺตมโน สตีมา.}}
\setlength{\csromangathaleftcol}{0pt}
\csromangathagroup{\csromangathastanza{\csromangathawak{\csromangathaitemline{131}{\textbf{สเจ ลเภถ นิปกํ สหายํ,}} \\ \csromangathacontline{\textbf{สทฺธิํ จรํ สาธุวิหาริ ธีรํ.}} \\ \csromangathacontline{\textbf{อภิภุยฺย สพฺพานิ ปริสฺสยานิ,}} \\ \csromangathacontline{\textbf{จเรยฺย เตน’ตฺตมโน สตีมา.}}}}}

\prose{\textbf{สเจ ลเภถ นิปกํ สหายนฺ}ติ สเจ นิปกํ ปณฺฑิตํ ปญฺญวนฺตํ พุทฺธิมนฺตํ ญาณิํ วิภาวิํ เมธาวิํ สหายํ ลเภยฺย ปฏิลเภยฺย อธิคจฺเฉยฺย วินฺเทยฺยาติ สเจ ลเภถ นิปกํ สหายํ.}

\prose{\textbf{สทฺธิํจรํ สาธุวิหาริ ธีรนฺ}ติ \textbf{สทฺธิํจรนฺ}ติ เอกโต จรํ. \textbf{สาธุวิหารินฺ}ติ ปฐเมนปิ ฌาเนน สาธุวิหาริํ. ทุติเยนปิ ฌาเนน. ตติเยนปิ ฌาเนน. จตุตฺเถนปิ ฌาเนน สาธุวิหาริํ. เมตฺตายปิ เจโตวิมุตฺติยา สาธุวิหาริํ. กรุณายปิ. มุทิตายปิ. อุเปกฺขายปิ เจโตวิมุตฺติยา สาธุวิหาริํ. อากาสา\-นญฺจา\-ยตนสมาปตฺติยาปิ สาธุวิหาริํ. วิญฺญาณญฺจายตน\-สมาปตฺติยาปิ ฯเปฯ. อากิญฺจญฺญายตน\-สมาปตฺติยาปิ ฯเปฯ. \csromanfolio{256}\textbf{เนว}\-\textbf{สญฺญา}\-\textbf{นา}\-\textbf{สญฺญา}\-\textbf{ยตนสมาปตฺติยาปิ สาธุวิหาริํ นิโรธ}\-\textbf{สมาปตฺติยาปิ} \textbf{สาธุวิหาริํ. ผลสมาปตฺติยาปิ สาธุวิหาริํ. ธีรนฺติ ธีรํ ปณฺฑิตํ} ปญฺญวนฺตํ พุทฺธิมนฺตํ ญาณิํ วิภาวิํ เมธาวินฺติ สทฺธิํจรํ สาธุวิหาริ ธีรํ.}

\prose{\textbf{อภิภุยฺย สพฺพานิ ปริสฺสยานี}ติ \textbf{ปริสฺสยา}ติ เทฺว ปริสฺสยา ปากฏ\-ปริสฺสยา จ ปฏิจฺฉนฺน\-ปริสฺสยา จ ฯเปฯ อิเม วุจฺจนฺติ ปากฏ\-ปริสฺสยา ฯเปฯ อิเม วุจฺจนฺติ ปฏิจฺฉนฺน\-ปริสฺสยา ฯเปฯ. เอวมฺปิ ตตฺราสยาติ ปริสฺสยา. \textbf{อภิภุยฺย สพฺพานิ ปริสฺสยานี}ติ สพฺเพ ปริสฺสเย อภิภุยฺย อภิภวิตฺวา อชฺโฌตฺถริตฺวา ปริยาทิยิตฺวา มทฺทิตฺวาติ อภิภุยฺย สพฺพานิ ปริสฺสยานิ.}

\prose{\textbf{จเรยฺย เตน’ตฺตมโน สตีมา}ติ โส ปจฺเจก\-สมฺพุทฺโธ เตน นิปเกน ปณฺฑิเตน ปญฺญวนฺเตน พุทฺธิมนฺเตน ญาณินา วิภาวินา เมธาวินา สหาเยน สทฺธิํ อตฺตมโน ตุฏฺฐมโน หฏฺฐมโน ปหฏฺฐมโน อุทคฺคมโน มุทิตมโน จเรยฺย วิหเรยฺย อิริเยยฺย วตฺเตยฺย ปาเลยฺย ยเปยฺย ยาเปยฺยาติ จเรยฺย เตนตฺตมโน. \textbf{สตีมา}ติ โส ปจฺเจก\-สมฺพุทฺโธ สติมา โหติ ปรเมน สติเนปกฺเกน สมนฺนาคโต จิรกตมฺปิ จิรภาสิตมฺปิ สริตา อนุสฺสริตาติ จเรยฺย เตน’ตฺตมโน สตีมา. เตนาห โส ปจฺเจก\-สมฺพุทฺโธ\csromandash{}}

\setlength{\csromangathapagewidth}{0pt}
\setlength{\csromangathawakwidth}{0pt}
\csromangathameasuregroup{}{สเจ ลเภถ นิปกํ สหายํ, \\ สทฺธิํจรํ สาธุวิหาริ ธีรํ.}
\csromangathameasurewak{สเจ ลเภถ นิปกํ สหายํ, \\ สทฺธิํจรํ สาธุวิหาริ ธีรํ.}
\setlength{\csromangathaleftcol}{0pt}
\csromangathagroup{\csromangathastanza{\csromangathawak{สเจ ลเภถ นิปกํ สหายํ, \\ สทฺธิํจรํ สาธุวิหาริ ธีรํ.}}}

\prose{อภิภุยฺย สพฺพานิ ปริสฺสยานิ,}

\prose{จเรยฺย เตนตฺตมโน สตีมาติ.}

\setlength{\csromangathapagewidth}{0pt}
\setlength{\csromangathawakwidth}{0pt}
\csromangathameasuregroup{}{\textbf{โน เจ ลเภถ นิปกํ สหายํ,} \\ สทฺธิํจรํ สาธุวิหาริ ธีรํ. \\ \textbf{ราชาว รฏฺฐํ วิชิตํ ปหาย,} \\ \textbf{เอโก จเร ขคฺควิสาณ}\textbf{กปฺโป.}}
\csromangathameasurewak{\textbf{โน เจ ลเภถ นิปกํ สหายํ,} \\ สทฺธิํจรํ สาธุวิหาริ ธีรํ. \\ \textbf{ราชาว รฏฺฐํ วิชิตํ ปหาย,} \\ \textbf{เอโก จเร ขคฺควิสาณ}\textbf{กปฺโป.}}
\setlength{\csromangathaleftcol}{0pt}
\csromangathagroup{\csromangathastanza{\csromangathawak{\csromangathaitemline{132}{\textbf{โน เจ ลเภถ นิปกํ สหายํ,}} \\ \csromangathacontline{สทฺธิํจรํ สาธุวิหาริ ธีรํ.} \\ \csromangathacontline{\textbf{ราชาว รฏฺฐํ วิชิตํ ปหาย,}} \\ \csromangathacontline{\textbf{เอโก จเร ขคฺควิสาณ}\-\textbf{กปฺโป.}}}}}

\prose{\textbf{โน เจ ลเภถ นิปกํ สหายนฺ}ติ โน เจ นิปกํ ปณฺฑิตํ ปญฺญวนฺตํ พุทฺธิมนฺตํ ญาณิํ วิภาวิํ เมธาวิํ สหายํ ลเภยฺย ปฏิลเภยฺย อธิคจฺเฉยฺย วินฺเทยฺยาติ โน เจ ลเภถ นิปกํ สหายํ.}

\prose{\csromanfolio{257}\textbf{สทฺธิํจรํ สาธุวิหาริ ธีรนฺ}ติ \textbf{สทฺธิํจรนฺ}ติ เอกโต จรํ. \textbf{สาธุวิหารินฺ}ติ ปฐเมนปิ ฌาเนน สาธุวิหาริํ ฯเปฯ นิโรธ\-สมาปตฺติยาปิ สาธุวิหาริํ ผลสมาปตฺติยาปิ สาธุวิหาริํ. \textbf{ธีรนฺ}ติ ธีรํ ปณฺฑิตํ ปญฺญวนฺตํ พุทฺธิมนฺตํ ญาณิํ วิภาวิํ เมธาวินฺติ สทฺธิํจรํ สาธุวิหาริ ธีรํ.}

\prose{\textbf{ราชาว รฏฺฐํ วิชิตํ ปหายา}ติ ราชา ขตฺติโย มุทฺธาภิสิตฺโต วิชิตสงฺคาโม นิหต\-ปจฺจามิตฺโต ลทฺธาธิปฺปาโย ปริปุณฺณ\-โกส\-โกฏฺฐาคาโร รฏฺฐญฺจ ชนปทญฺจ โกสญฺจ โกฏฺฐาคารญฺจ ปหูต\-หิรญฺญ\-สุวณฺณํ นครญฺจ ปริจฺจชิตฺวา เกสมสฺสุํ โอหาเรตฺวา กาสายานิ วตฺถานิ อจฺฉาเทตฺวา อคารสฺมา อนคาริยํ ปพฺพชิตฺวา อกิญฺจนภาวํ อุปคนฺตฺวา เอโก จรติ วิหรติ อิริยติ วตฺเตติ ปาเลติ ยเปติ ยาเปติ. เอวํ ปจฺเจก\-สมฺพุทฺโธปิ สพฺพํ ฆราวาส\-ปลิโพธํ ฉินฺทิตฺวา ปุตฺตทาร\-ปลิโพธํ ฉินฺทิตฺวา ญาติปลิโพธํ ฉินฺทิตฺวา มิตฺตามจฺจ\-ปลิโพธํ ฉินฺทิตฺวา เกสมสฺสุํ โอหาเรตฺวา กาสายานิ วตฺถานิ อจฺฉาเทตฺวา อคารสฺมา อนคาริยํ ปพฺพชิตฺวา อกิญฺจนภาวํ อุปคนฺตฺวา เอโก จรติ วิหรติ อิริยติ วตฺเตติ ปาเลติ ยเปติ ยาเปตีติ ราชาว รฏฺฐํ วิชิตํ ปหาย, เอโก จเร ขคฺควิสาณ\-กปฺโป. เตนาห โส ปจฺเจก\-สมฺพุทฺโธ\csromandash{}}

\prose{โน เจ ลเภถ นิปกํ สหายํ,}

\prose{สทฺธิํจรํ สาธุวิหาริ ธีรํ.}

\prose{ราชาว รฏฺฐํ วิชิตํ ปหาย,}

\prose{เอโก จเร ขคฺควิสาณ\-กปฺโปติ.}

\setlength{\csromangathapagewidth}{0pt}
\setlength{\csromangathawakwidth}{0pt}
\csromangathameasuregroup{}{\textbf{อทฺธา ปสํสาม สหายสมฺปทํ,} \\ \textbf{เสฏฺฐา สมา เสวิตพฺพา สหายา.} \\ \textbf{เอเต อลทฺธา อนวชฺชโภชี,} \\ \textbf{เอโก จเร ขคฺควิสาณ}\textbf{กปฺโป.}}
\csromangathameasurewak{\textbf{อทฺธา ปสํสาม สหายสมฺปทํ,} \\ \textbf{เสฏฺฐา สมา เสวิตพฺพา สหายา.} \\ \textbf{เอเต อลทฺธา อนวชฺชโภชี,} \\ \textbf{เอโก จเร ขคฺควิสาณ}\textbf{กปฺโป.}}
\setlength{\csromangathaleftcol}{0pt}
\csromangathagroup{\csromangathastanza{\csromangathawak{\csromangathaitemline{133}{\textbf{อทฺธา ปสํสาม สหายสมฺปทํ,}} \\ \csromangathacontline{\textbf{เสฏฺฐา สมา เสวิตพฺพา สหายา.}} \\ \csromangathacontline{\textbf{เอเต อลทฺธา อนวชฺชโภชี,}} \\ \csromangathacontline{\textbf{เอโก จเร ขคฺควิสาณ}\-\textbf{กปฺโป.}}}}}

\prose{\textbf{อทฺธา ปสํสาม สหาย}\-\textbf{สมฺปทนฺ}ติ \textbf{อทฺธา}ติ เอกํสวจนํ นิสฺสํสยวจนํ นิกฺกงฺข\-วจนํ อเทฺวชฺฌวจนํ อเทฺวฬฺหก\-วจนํ นิโยควจนํ อปณฺณก\-วจนํ อวิรทฺธ\-วจนํ อวตฺถาปน\-วจนเมตํ อทฺธาติ. \textbf{สหาย}\-\textbf{สมฺปทนฺ}ติ สหายสมฺปทา วุจฺจติ โย โส สหาโย อเสกฺเขน สีลกฺขนฺเธน สมนฺนาคโต โหติ. อเสกฺเขน สมาธิ\-กฺขนฺเธน. อเสกฺเขน \csromanfolio{258}ปญฺญา\-กฺขนฺเธน. อเสกฺเขน วิมุตฺติ\-กฺขนฺเธน. อเสกฺเขน วิมุตฺติ\-ญาณ\-ทสฺสนกฺขนฺเธน สมนฺนาคโต โหติ. \textbf{อทฺธา ปสํสาม สหาย}\-\textbf{สมฺปทนฺ}ติ สหายสมฺปทํ ปสํสาม โถเมม กิตฺเตม วณฺเณมาติ อทฺธา ปสํสาม สหายสมฺปทํ.}

\prose{\textbf{เสฏฺฐา สมา เสวิตพฺพา สหายา}ติ เสฏฺฐา โหนฺติ สหายา สีเลน สมาธินา ปญฺญาย วิมุตฺติยา วิมุตฺติ\-ญาณ\-ทสฺสเนน, สมา สทิสา โหนฺติ สหาย สีเลน สมาธินา ปญฺญาย วิมุตฺติยา วิมุตฺติ\-ญาณ\-ทสฺสเนน, เสฏฺฐา วา สหายา สทิสา วา สหายา เสวิตพฺพา ภชิตพฺพา ปยิรุปาสิตพฺพา ปริปุจฺฉิตพฺพา ปริปญฺหิตพฺพาติ เสฏฺฐา สมา เสวิตพฺพา สหายา.}

\prose{\textbf{เอเต อลทฺธา อนวชฺชโภชี}ติ อตฺถิ ปุคฺคโล สาวชฺชโภชี อตฺถิ ปุคฺคโล อนวชฺชโภชีติ. กตโม จ ปุคฺคโล สาวชฺชโภชี. อิเธกจฺโจ ปุคฺคโล กุหนาย ลปนาย เนมิตฺติกตาย นิปฺเปสิกตาย ลาเภน ลาภํ นิชิคีสนตาย\csromansharedfootnote{72b4a8b08af3cf74}{นิชิคิํสนตาย (สฺยา)} ทารุทาเนน เวฬุทาเนน ปตฺตทาเนน ปุปฺผทาเนน ผลทาเนน สินานทาเนน จุณฺณทาเนน มตฺติกาทาเนน ทนฺตกฏฺฐ\-ทาเนน มุโขทก\-ทาเนน จาฏุกมฺยตาย\csromansharedfootnote{271139561e1b3693}{ปาตุกมฺยตาย (สฺยา) ขุ 7. 290 ปิฏฺเฐปิ.} มุคฺคสูปฺยตาย\csromansharedfootnote{b88f2659563e5f1b}{มุคฺคสูปตาย (สฺยา)} ปาริภฏฺยตาย ปีฐมทฺทิก\-ตาย วตฺถุวิชฺชาย ติรจฺฉาน\-วิชฺชาย องฺควิชฺชาย นกฺขตฺต\-วิชฺชาย ทูตคมเนน ปหิณคมเนน ชงฺฆ\-เปสนิเยน เวชฺชกมฺเมน นวกมฺเมน\csromansharedfootnote{deec4e79e0f7c887}{ทูตกมฺเมน (สฺยา, ก) ขุ 7. 290 ปิฏฺเฐ.} ปิณฺฑ\-ปฏิปิณฺฑเกน ทานา\-นุปฺปทาเนน อธมฺเมน วิสเมน ลทฺธา ลภิตฺวา อธิคนฺตฺวา วินฺทิตฺวา ปฏิลภิตฺวา ชีวิกํ\csromansharedfootnote{f8e93fdd53c8c38b}{ชีวิตํ (ก)} กปฺเปติ. อยํ วุจฺจติ ปุคฺคโล สาวชฺชโภชี.}

\prose{กตโม จ ปุคฺคโล อนวชฺชโภชี. อิเธกจฺโจ ปุคฺคโล น กุหนาย น ลปนาย น เนมิตฺติกตาย น นิปฺเปสิกตาย น ลาเภน ลาภํ นิชิคีสนตาย น ทารุทาเนน น เวฬุทาเนน น ปตฺตทาเนน น ปุปฺผทาเนน น ผลทาเนน น สินานทาเนน น จุณฺณทาเนน น มตฺติกาทาเนน น ทนฺตกฏฺฐ\-ทาเนน น มุโขทก\-ทาเนน น จาฏุกมฺยตาย น มุคฺคสูปฺยตาย น ปาริภฏฺยตาย น ปีฐมทฺทิก\-ตาย น วตฺถุวิชฺชาย น ติรจฺฉาน\-วิชฺชาย}

\prose{\csromanfolio{259}น องฺควิชฺชาย น นกฺขตฺต\-วิชฺชาย น ทูตคมเนน น ปหิณคมเนน น ชงฺฆ\-เปสนิเยน น เวชฺชกมฺเมน น นวกมฺเมน น ปิณฺฑ\-ปฏิปิณฺฑเกน น ทานา\-นุปฺปทาเนน ธมฺเมน สเมน ลทฺธา ลภิตฺวา อธิคนฺตฺวา วินฺทิตฺวา ปฏิลภิตฺวา ชีวิกํ กปฺเปติ. อยํ วุจฺจติ ปุคฺคโล อนวชฺชโภชี.}

\prose{\textbf{เอเต อลทฺธา อนวชฺชโภชี}ติ เอเต อนวชฺชโภชี อลทฺธา อลภิตฺวา อนธิคนฺตฺวา อวินฺทิตฺวา อปฺปฏิลภิตฺวาติ เอเต อลทฺธา อนวชฺชโภชี, เอโก จเร ขคฺควิสาณ\-กปฺโป. เตนาห โส ปจฺเจก\-สมฺพุทฺโธ\csromandash{}}

\setlength{\csromangathapagewidth}{0pt}
\setlength{\csromangathawakwidth}{0pt}
\csromangathameasuregroup{}{อทฺธา ปสํสาม สหายสมฺปทํ, \\ เสฏฺฐา สมา เสวิตพฺพา สหายา. \\ เอเต อลทฺธา อนวชฺชโภชี, \\ เอโก จเร ขคฺควิสาณกปฺโปติ.}
\csromangathameasurewak{อทฺธา ปสํสาม สหายสมฺปทํ, \\ เสฏฺฐา สมา เสวิตพฺพา สหายา. \\ เอเต อลทฺธา อนวชฺชโภชี, \\ เอโก จเร ขคฺควิสาณกปฺโปติ.}
\setlength{\csromangathaleftcol}{0pt}
\csromangathagroup{\csromangathastanza{\csromangathawak{อทฺธา ปสํสาม สหายสมฺปทํ, \\ เสฏฺฐา สมา เสวิตพฺพา สหายา. \\ เอเต อลทฺธา อนวชฺชโภชี, \\ เอโก จเร ขคฺควิสาณ\-กปฺโปติ.}}}

\proseitemwithcloser{134}{\textbf{ทิสฺวา สุวณฺณสฺส ปภสฺสรานิ,}}{\textbf{กมฺมารปุตฺเตน สุนิฏฺฐิตานิ.}}
\setlength{\csromangathapagewidth}{0pt}
\setlength{\csromangathawakwidth}{0pt}
\csromangathameasuregroup{}{\textbf{สํฆฏฺฏ}\textbf{ยนฺตานิ}\textsuperscript{1}\textbf{ ทุเว ภุชสฺมิํ,} \\ \textbf{เอโก จเร ขคฺควิสาณ}\textbf{กปฺโป.}}
\csromangathameasurewak{\textbf{สํฆฏฺฏ}\textbf{ยนฺตานิ}\textsuperscript{1}\textbf{ ทุเว ภุชสฺมิํ,} \\ \textbf{เอโก จเร ขคฺควิสาณ}\textbf{กปฺโป.}}
\setlength{\csromangathaleftcol}{0pt}
\csromangathagroup{\csromangathastanza{\csromangathawak{\textbf{สํฆฏฺฏ}\-\textbf{ยนฺตานิ}\csromansharedfootnote{b129124ff7f77eb7}{สํฆฏฺฏมานานิ (ขุ 1. 286 ปิฏฺเฐ.)}\textbf{ ทุเว ภุชสฺมิํ,} \\ \textbf{เอโก จเร ขคฺควิสาณ}\-\textbf{กปฺโป.}}}}

\prose{\textbf{ทิสฺวา สุวณฺณสฺส ปภสฺสรานี}ติ ทิสฺวา ปสฺสิตฺวา ตุลยิตฺวา ตีรยิตฺวา วิภาวยิตฺวา วิภูตํ กตฺวา. \textbf{สุวณฺณสฺสา}ติ ชาตรูปสฺส. \textbf{ปภสฺสรานี}ติ ปริสุทฺธานิ ปริโยทาตานีติ ทิสฺวา สุวณฺณสฺส ปภสฺสรานิ.}

\prose{\textbf{กมฺมารปุตฺเตน สุนิฏฺฐิตานี}ติ กมฺมารปุตฺโต วุจฺจติ สุวณฺณกาโร. \textbf{กมฺมารปุตฺเตน สุนิฏฺฐิตานี}ติ กมฺมารปุตฺเตน สุนิฏฺฐิตานิ สุกตานิ สุปริกมฺม\-กตานีติ กมฺมารปุตฺเตน สุนิฏฺฐิตานิ.}

\prose{\textbf{สํฆฏฺฏ}\-\textbf{ยนฺตานิ ทุเว ภุชสฺมินฺ}ติ ภุโช วุจฺจติ หตฺโถ. ยถา เอกสฺมิํ หตฺเถ เทฺว นูปุรานิ\csromansharedfootnote{a485e9f498581d89}{ธุวรานิ (สฺยา)} ฆฏฺเฏนฺติ\csromansharedfootnote{cfa912c9296e6950}{ฆเฏนฺติ (สฺยา)}, เอวเมว สตฺตา ตณฺหาวเสน ทิฏฺฐิวเสน นิรเย ฆฏฺเฏนฺติ, ติรจฺฉาน\-โยนิยํ ฆฏฺเฏนฺติ, เปตฺติวิสเย ฆฏฺเฏนฺติ, มนุสฺสาโลเก ฆฏฺเฏนฺติ, เทวโลเก ฆฏฺเฏนฺติ, คติยา คติํ}

\prose{\csromanfolio{260}อุปปตฺติยา อุปปตฺติํ ปฏิสนฺธิยา ปฏิสนฺธิํ ภเวน ภวํ สํสาเรน สํสารํ วฏฺเฏน วฏฺฏํ ฆฏฺเฏนฺติ สํฆฏฺเฏนฺติ สํฆฏฺเฏนฺตา จรนฺติ วิหรนฺติ อิริยนฺติ วตฺเตนฺติ ปาเลนฺติ ยเปนฺติ ยาเปนฺตีติ สํฆฏฺฏ\-ยนฺตานิ ทุเว ภุชสฺมิํ, เอโก จเร ขคฺควิสาณ\-กปฺโป. เตนาห โส ปจฺเจก\-สมฺพุทฺโธ\csromandash{}}

\prosewithcloser{ทิสฺวา สุวณฺณสฺส ปภสฺสรานิ,}{กมฺมารปุตฺเตน สุนิฏฺฐิตานิ.}
\setlength{\csromangathapagewidth}{0pt}
\setlength{\csromangathawakwidth}{0pt}
\csromangathameasuregroup{}{สํฆฏฺฏยนฺตานิ ทุเว ภุชสฺมิํ, \\ เอโก จเร ขคฺควิสาณกปฺโป. \\ \textbf{เอวํ ทุตีเยน สห มมสฺส,} \\ \textbf{วาจาภิลาโป อภิสชฺชนา วา.} \\ \textbf{เอตํ ภยํ อายติํ เปกฺขมาโน,} \\ เอโก จเร ขคฺควิสาณกปฺโป.}
\csromangathameasurewak{สํฆฏฺฏยนฺตานิ ทุเว ภุชสฺมิํ, \\ เอโก จเร ขคฺควิสาณกปฺโป. \\ \textbf{เอวํ ทุตีเยน สห มมสฺส,} \\ \textbf{วาจาภิลาโป อภิสชฺชนา วา.} \\ \textbf{เอตํ ภยํ อายติํ เปกฺขมาโน,} \\ เอโก จเร ขคฺควิสาณกปฺโป.}
\setlength{\csromangathaleftcol}{0pt}
\csromangathagroup{\csromangathastanza{\csromangathawak{สํฆฏฺฏ\-ยนฺตานิ ทุเว ภุชสฺมิํ, \\ เอโก จเร ขคฺควิสาณ\-กปฺโป.}}\csromangathastanza{\csromangathawak{\csromangathaitemline{135}{\textbf{เอวํ ทุตีเยน สห มมสฺส,}} \\ \csromangathacontline{\textbf{วาจาภิลาโป อภิสชฺชนา วา.}} \\ \csromangathacontline{\textbf{เอตํ ภยํ อายติํ เปกฺขมาโน,}} \\ \csromangathacontline{เอโก จเร ขคฺควิสาณ\-กปฺโป.}}}}

\prose{\textbf{เอวํ ทุตีเยน สห มมสฺสา}ติ ตณฺหาทุติโย วา โหติ ปุคฺคลทุติโย วา. กถํ ตณฺหาทุติโย โหติ, \textbf{ตณฺหา}ติ รูปตณฺหา ฯเปฯ ธมฺมตณฺหา. ยสฺเสสา ตณฺหา อปฺปหีนา, โส วุจฺจติ ตณฺหาทุติโย.}

\setlength{\csromangathapagewidth}{0pt}
\setlength{\csromangathawakwidth}{0pt}
\csromangathameasuregroup{ตณฺหาทุติโย ปุริโส, \\ อิตฺถภาวญฺญถาภาวํ,}{\csromangathabat{ตณฺหาทุติโย ปุริโส,}{ทีฆมทฺธาน สํสรํ.} \\ \csromangathabat{อิตฺถภาวญฺญถาภาวํ,}{สํสารํ นาติวตฺตตีติ.}}
\csromangathasetleft{ตณฺหาทุติโย ปุริโส, \\ อิตฺถภาวญฺญถาภาวํ,}
\csromangathagroup{\csromangathastanza{\csromangathabat{ตณฺหาทุติโย ปุริโส,}{ทีฆมทฺธาน สํสรํ.} \\ \csromangathabat{อิตฺถ\-ภาว\-ญฺญถา\-ภาวํ,}{สํสารํ นาติวตฺตตีติ.}}}

\prose{เอวํ ตณฺหาทุติโย วา โหติ.}

\prose{กถํ ปุคฺคลทุติโย โหติ, อิเธกจฺโจ น อตฺถเหตุ\csromansharedfootnote{0c65729098ba3205}{อตฺตเหตุ (สฺยา)} น การณเหตุ อุทฺธโต อวูปสนฺต\-จิตฺโต เอกสฺส วา ทุติโย โหติ ทฺวินฺนํ วา ตติโย โหติ, ติณฺณํ วา จตุตฺโถ โหติ. ตตฺถ พหุํ สมฺผปฺปลาปํ ปลปติ. เสยฺยถิทํ, ราชกถํ โจรกถํ มหามตฺตกถํ เสนากถํ ภยกถํ ยุทฺธกถํ อนฺนกถํ ปานกถํ วตฺถกถํ สยนกถํ มาลากถํ คนฺธกถํ ญาติกถํ ยานกถํ คามกถํ นิคมกถํ นครกถํ ชนปทกถํ อิตฺถิกถํ\csromansharedfootnote{a50102d1d99f3adf}{อิตฺถิกถํ ปุริสกถํ (สฺยา, ก) ที 1. 7, 62 ปิฏฺเฐ; สํ 3. 368 ปิฏฺเฐ จ ปสฺสิตพฺพํ.} สูรกถํ วิสิขากถํ กุมฺภ\-ฏฺฐานกถํ ปุพฺพเปตกถํ}

\prose{\csromanfolio{261}นานตฺตกถํ โลกกฺขายิกํ สมุทฺท\-กฺขายิกํ อิติภวา\-ภวกถํ กเถติ. เอวํ ปุคฺคลทุติโย โหตีติ เอวํ ทุตีเยน สห มมสฺส.}

\prose{\textbf{วาจาภิลาโป อภิสชฺชนา วา}ติ วาจาภิลาโป วุจฺจติ พาตฺติํส\-ติรจฺฉาน\-กถา. เสยฺยถิทํ, ราชกถํ ฯเปฯ อิติภวา\-ภวกถํ. \textbf{อภิสชฺชนา วา}ติ เทฺว สชฺชนา, ตณฺหาสชฺชนา จ ทิฏฺฐิสชฺชนา จ ฯเปฯ อยํ ตณฺหาสชฺชนา ฯเปฯ อยํ ทิฏฺฐิ\-สชฺชนาติ วาจาภิลาโป อภิสชฺชนา วา.}

\prose{\textbf{เอตํ ภยํ อายติํ เปกฺขมาโน}ติ \textbf{ภยนฺ}ติ ชาติภยํ ชราภยํ พฺยาธิภยํ มรณภยํ ราชภยํ โจรภยํ อคฺคิภยํ อุทกภยํ อตฺตานุวาทภยํ ปรานุวาทภยํ ทณฺฑภยํ ทุคฺคติภยํ อูมิภยํ กุมฺภิลภยํ อาวฏฺฏภยํ สุสุมารภยํ อาชีวิกภยํ อสิโลกภยํ ปริส\-สารชฺชภยํ มทนภยํ ภยานกํ ฉมฺภิตตฺตํ โลมหํโส เจตโส อุพฺเพโค อุตฺราโส. เอตํ ภยํ อายติํ เปกฺขมาโนติ เอตํ ภยํ อายติํ เปกฺขมาโน ทกฺขมาโน โอโลกยมาโน นิชฺฌายมาโน อุปปริกฺขมาโนติ เอตํ ภยํ อายติํ เปกฺขมาโน, เอโก จเร ขคฺควิสาณ\-กปฺโป. เตนาห โส ปจฺเจก\-สมฺพุทฺโธ\csromandash{}}

\setlength{\csromangathapagewidth}{0pt}
\setlength{\csromangathawakwidth}{0pt}
\csromangathameasuregroup{}{เอวํ ทุตีเยน สห มมสฺส, \\ วาจาภิลาโป อภิสชฺชนา วา.}
\csromangathameasurewak{เอวํ ทุตีเยน สห มมสฺส, \\ วาจาภิลาโป อภิสชฺชนา วา.}
\setlength{\csromangathaleftcol}{0pt}
\csromangathagroup{\csromangathastanza{\csromangathawak{เอวํ ทุตีเยน สห มมสฺส, \\ วาจาภิลาโป อภิสชฺชนา วา.}}}

\prose{เอตํ ภยํ อายติํ เปกฺขมาโน,}

\prose{เอโก จเร ขคฺควิสาณ\-กปฺโป.}

\setlength{\csromangathapagewidth}{0pt}
\setlength{\csromangathawakwidth}{0pt}
\csromangathameasuregroup{}{\textbf{กามา หิ จิตฺรา มธุรา มโนรมา,} \\ \textbf{วิรูปรูเปน มเถนฺติ จิตฺตํ.} \\ \textbf{อาทีนวํ กามคุเณสุ ทิสฺวา,} \\ เอโก จเร ขคฺควิสาณกปฺโป.}
\csromangathameasurewak{\textbf{กามา หิ จิตฺรา มธุรา มโนรมา,} \\ \textbf{วิรูปรูเปน มเถนฺติ จิตฺตํ.} \\ \textbf{อาทีนวํ กามคุเณสุ ทิสฺวา,} \\ เอโก จเร ขคฺควิสาณกปฺโป.}
\setlength{\csromangathaleftcol}{0pt}
\csromangathagroup{\csromangathastanza{\csromangathawak{\csromangathaitemline{136}{\textbf{กามา หิ จิตฺรา มธุรา มโนรมา,}} \\ \csromangathacontline{\textbf{วิรูปรูเปน มเถนฺติ จิตฺตํ.}} \\ \csromangathacontline{\textbf{อาทีนวํ กามคุเณสุ ทิสฺวา,}} \\ \csromangathacontline{เอโก จเร ขคฺควิสาณ\-กปฺโป.}}}}

\prose{\textbf{กามา หิ จิตฺรา มธุรา มโนรมา}ติ กามาติ อุทฺทานโต เทฺว \textbf{กามา} วตฺถุกามา จ กิเลสกามา จ ฯเปฯ. อิเม วุจฺจนฺติ วตฺถุกามา ฯเปฯ. อิเม วุจฺจนฺติ กิเลสกามา. \textbf{จิตฺรา}ติ นานาวณฺณา รูปา นานาวณฺณา สทฺทา นานาวณฺณา คนฺธา นานาวณฺณา รสา นานาวณฺณา โผฏฺฐพฺพา อิฏฺฐา กนฺตา มนาปา ปิยรูปา กามูปสํหิตา รชนียา. \textbf{มธุรา}ติ วุตฺตํ เหตํ ภควตา\csromansharedfootnote{7b3306c9781f9024}{ม 1. 120 ปิฏฺเฐ; ม 3. 276 ปิฏฺเฐ จ ปสฺสิตพฺพํ.} \csromanfolio{262}ปญฺจิเม ภิกฺขเว กามคุณา. กตเม ปญฺจ, จกฺขุวิญฺเญยฺยา รูปา อิฏฺฐา กนฺตา มนาปา ปิยรูปา กามูปสํหิตา รชนียา. โสตวิญฺเญยฺยา สทฺทา ฯเปฯ. ฆานวิญฺเญยฺยา คนฺธา. ชิวฺหาวิญฺเญยฺยา รสา. กายวิญฺเญยฺยา โผฏฺฐพฺพา อิฏฺฐา กนฺตา มนาปา ปิยรูปา กามูปสํหิตา รชนียา. อิเม โข ภิกฺขเว ปญฺจ กามคุณา. ยํ โข ภิกฺขเว อิเม ปญฺจ กามคุเณ ปฏิจฺจ อุปฺปชฺชติ สุขํ โสมนสฺสํ, อิทํ วุจฺจติ กามสุขํ มิฬฺหสุขํ\csromansharedfootnote{8343b26228d7f778}{มีฬฺหสุขํ (ปสฺสม 3. 276 ปิฏฺเฐ.)} ปุถุชฺชนสุขํ อนริยสุขํ “น เสวิตพฺพํ น ภาเวตพฺพํ น พหุลีกาตพฺพํ ภายิตพฺพํ เอตสฺส สุขสฺสา”ติ วทามีติ กามา หิ จิตฺรา มธุรา. \textbf{มโนรมา}ติ \textbf{มโน}ติ ยํ จิตฺตํ ฯเปฯ ตชฺชา มโนวิญฺญาณ\-ธาตุ, มโน รเมนฺติ โถเมนฺติ โตเสนฺติ ปหาเสนฺตีติ กามา หิ จิตฺรา มธุรา มโนรมา.}

\prose{\textbf{วิรูปรูเปน มเถนฺติ จิตฺตนฺ}ติ นานาวณฺเณหิ รูเปหิ ฯเปฯ นานาวณฺเณหิ โผฏฺฐพฺเพหิ จิตฺตํ มเถนฺติ โตเสนฺติ ปหาเสนฺตีติ วิรูปรูเปน มเถนฺติ จิตฺตํ.}

\prose{\textbf{อาทีนวํ กามคุเณสุ ทิสฺวา}ติ วุตฺตํ เหตํ ภควตา\csromandash{}โก จ ภิกฺขเว กามานํ อาทีนโว. อิธ ภิกฺขเว กุลปุตฺโต เยน สิปฺปฏฺฐาเนน ชีวิกํ กปฺเปติ, ยทิ มุทฺทาย, ยทิ คณนาย, ยทิ สงฺขาเนน, ยทิ กสิยา, ยทิ วณิชฺชาย, ยทิ โครกฺเขน, ยทิ อิสฺสตฺเถน,\csromansharedfootnote{a9e8937e88471c50}{อิสฺสฏฺเฐน (สฺยา), อิสฺสตฺเตน (ก) ปสฺส ม 1. 120 ปิฏฺเฐ.} ยทิ ราชโปริเสน, ยทิ สิปฺปญฺญตเรน, สีตสฺส ปุรกฺขโต อุณฺหสฺส ปุรกฺขโต ฑํสมกส\-วาตา\-ตป\-สรีสป\-สมฺผสฺเสหิ สมฺผสฺสมาโน\csromansharedfootnote{41501397fb2ba420}{ริสฺสมาโน (ม 1. 120 ปิฏฺเฐ.)} ขุปฺปิปาสาย มียมาโน, อยํ ภิกฺขเว กามานํ อาทีนโว สนฺทิฏฺฐิโก ทุกฺขกฺขนฺโธ กามเหตุ กามนิทานํ กามาธิกรณํ กามานเมว เหตุ.}

\prose{ตสฺส เจ ภิกฺขเว กุลปุตฺตสฺส เอวํ อุฏฺฐหโต ฆฏโต วายมโต เต โภคา นาภิ\-นิ\-ปฺผชฺชนฺติ, โส โสจติ กิลมติ ปริเทวติ, อุรตฺตาฬิํ กนฺทติ, สมฺโมหํ อาปชฺชติ “โมฆํ วต เม อุฏฺฐานํ, อผโล วต เม วายาโม”ติ, อยมฺปิ ภิกฺขเว กามานํ อาทีนโว สนฺทิฏฺฐิโก ทุกฺขกฺขนฺโธ กามเหตุ กามนิทานํ กามาธิกรณํ กามานเมว เหตุ.}

\prose{\csromanfolio{263}ตสฺส เจ ภิกฺขเว กุลปุตฺตสฺส เอวํ อุฏฺฐหโต ฆฏโต วายมโต เต โภคา อภินิ\-ปฺผชฺชนฺติ, โส เตสํ โภคานํ อารกฺขาธิกรณํ ทุกฺขํ โทมนสฺสํ\csromansharedfootnote{2dfc8562dc7e98a9}{ทุกฺขโทมนสฺสํ (สฺยา, ก)} ปฏิสํเวเทติ “กินฺติ เม โภเค เนว ราชาโน หเรยฺยุํ, น โจรา หเรยฺยุํ, น อคฺคิ ฑเหยฺย, น อุทกํ วเหยฺย, น อปฺปิยา ทายาทา หเรยฺยุนฺ”ติ. ตสฺส เอวํ อารกฺขโต โคปยโต เต โภเค ราชาโน วา หรนฺติ, โจรา วา หรนฺติ, อคฺคิ วา ฑหติ, อุทกํ วา วหติ, อปฺปิยา วา ทายาทา หรนฺติ, โส โสจติ ฯเปฯ สมฺโมหํ อาปชฺชติ “ยมฺปิ เม อโหสิ, ตมฺปิ โน นตฺถี”ติ. อยมฺปิ ภิกฺขเว กามานํ อาทีนโว สนฺทิฏฺฐิโก ทุกฺขกฺขนฺโธ กามเหตุ กามนิทานํ กามาธิกรณํ กามนเมว เหตุ.}

\prose{ปุน จปรํ ภิกฺขเว กามเหตุ กามนิทานํ กามาธิกรณํ กามานเมว เหตุ ราชาโนปิ ราชูหิ วิวทนฺติ, ขตฺติยาปิ ขตฺติเยหิ วิวทนฺติ, พฺราหฺมณาปิ พฺราหฺมเณหิ วิวทนฺติ, คหปตีปิ คหปตีหิ วิวทนฺติ, มาตาปิ ปุตฺเตน วิวทติ, ปุตฺโตปิ มาตรา วิวทติ, ปิตาปิ ปุตฺเตน วิวทติ, ปุตฺโตปิ ปิตรา วิวทติ, ภาตาปิ ภคินิยา วิวทติ, ภคินีปิ ภาตรา วิวทติ, สหาโยปิ สหาเยน วิวทติ. เต ตตฺถ กลห\-วิคฺคห\-วิวาทา\-ปนฺนา อญฺญมญฺญํ ปาณีหิปิ อุปกฺกมนฺติ, เลฑฺฑูหิปิ อุปกฺกมนฺติ, ทณฺเฑหิปิ อุปกฺกมนฺติ, สตฺเถหิปิ อุปกฺกมนฺติ, เต ตตฺถ มรณมฺปิ นิคจฺฉนฺติ มรณมตฺตมฺปิ ทุกฺขํ. อยมฺปิ ภิกฺขเว กามานํ อาทีนโว สนฺทิฏฺฐิโก ทุกฺขกฺขนฺโธ กามเหตุ กามนิทานํ กามาธิกรณํ กามานเมว เหตุ.}

\prose{ปุน จปรํ ภิกฺขเว กามเหตุ กามนิทานํ กามาธิกรณํ กามานเมว เหตุ อสิจมฺมํ คเหตฺวา ธนุกลาปํ สนฺนยฺหิตฺวา อุภโตพฺยูฬฺหํ\csromansharedfootnote{dc98fa0db2e39d3b}{อุภโตวิยูฬฺหํ (สฺยา) ปสฺส ม 1. 121 ปิฏฺเฐ.} สงฺคามํ ปกฺขนฺทนฺติ อุสูสุปิ ขิปฺปมาเนสุ สตฺตีสุปิ ขิปฺปมานาสุ อสีสุปิ วิชฺโชตลนฺเตสุ, เต ตตฺถ อุสูหิปิ วิชฺฌนฺติ, สตฺตีหิปิ\csromansharedfootnote{924175bfc0b4dbcf}{สตฺติยาปิ (ม 1. 121, 128 ปิฏฺเฐ.)} วิชฺฌนฺติ, อสินาปิ สีสํ ฉินฺทนฺติ. เต ตตฺถ มรณมฺปิ นิคจฺฉนฺติ มรณมตฺตมฺปิ ทุกฺขํ. อยมฺปิ ภิกฺขเว กามานํ อาทีนโว สนฺทิฏฺฐิโก ทุกฺขกฺขนฺโธ กามเหตุ กามนิทานํ กามาธิกรณํ กามานเมว เหตุ.}

\prose{ปุน จปรํ ภิกฺขเว กามเหตุ กามนิทานํ กามาธิกรณํ กามานเมว เหตุ อสิจมฺมํ คเหตฺวา ธนุกลาปํ สนฺนยฺหิตฺวา อทฺทาวเลปนา\csromansharedfootnote{c811db1421fe94d3}{อทฺธาวเลปนา (สฺยา)}}

\prose{\csromanfolio{264}อุปการิโย ปกฺขนฺทนฺติ อุสูสุปิ ขิปฺปมาเนสุ สตฺตีสุปิ ขิปฺปมานาสุ อสีสุปิ วิชฺโชตลนฺเตสุ. เต ตตฺถ อุสูหิปิ วิชฺฌนฺติ, สตฺตีหิปิ วิชฺฌนฺติ, ฉกณกายปิ\csromansharedfootnote{ea85bc0bfe38960a}{ฉกณฏิยาปิ (สฺยา), ปกฺกฏฺฐิยา (สีฏฺฐ) ปสฺส ม 1. 121 ปิฏฺเฐ.} โอสิญฺจนฺติ. อภิวคฺเคนปิ โอมทฺทนฺติ, อสินาปิ สีสํ ฉินฺทนฺติ. เต ตตฺถ มรณมฺปิ นิคจฺฉนฺติ มรณมตฺตมฺปิ ทุกฺขํ. อยมฺปิ ภิกฺขเว กามานํ อาทีนโว สนฺทิฏฺฐิโก ทุกฺขกฺขนฺโธ กามเหตุ กามนิทานํ กามาธิกรณํ กามานเมว เหตุ.}

\prose{ปุน จปรํ ภิกฺขเว กามเหตุ กามนิทานํ กามาธิกรณํ กามานเมว เหตุ สนฺธิมฺปิ ฉินฺทนฺติ, นิลฺโลปมฺปิ หรนฺติ, เอกาคาริกมฺปิ กโรนฺติ, ปริปนฺเถปิ ติฏฺฐนฺติ, ปรทารมฺปิ คจฺฉนฺติ. ตเมนํ ราชาโน คเหตฺวา วิวิธา กมฺมการณา กาเรนฺติ, กสาหิปิ ตาเฬนฺติ, วตฺเตหิปิ ตาเฬนฺติ, อฑฺฒ\-ทณฺฑเกหิปิ ตาเฬนฺติ, หตฺถมฺปิ ฉินฺทนฺติ ฯเปฯ อสินาปิ สีสํ ฉินฺทนฺติ. เต ตตฺถ มรณมฺปิ นิคจฺฉนฺติ มรณมตฺตมฺปิ ทุกฺขํ. อยมฺปิ ภิกฺขเว กามานํ อาทีนโว สนฺทิฏฺฐิโก ทุกฺขกฺขนฺโธ กามเหตุ กามนิทานํ กามาธิกรณํ กามานเมว เหตุ.}

\prose{ปุน จปรํ ภิกฺขเว กามเหตุ กามนิทานํ กามาธิกรณํ กามานเมว เหตุ กาเยน ทุจฺจริตํ จรนฺติ, วาจาย ทุจฺจริตํ จรนฺติ, มนสา ทุจฺจริตํ จรนฺติ. เต กาเยน ทุจฺจริตํ จริตฺวา วาจาย ทุจฺจริตํ จริตฺวา มนสา ทุจฺจริตํ จริตฺวา กายสฺส เภทา ปรํ มรณา อปายํ ทุคฺคติํ วินิปาตํ นิรยํ อุปปชฺชนฺติ. อยมฺปิ ภิกฺขเว กามานํ อาทีนโว สมฺปรายิโก ทุกฺขกฺขนฺโธ กามเหตุ กามนิทานํ กามาธิกรณํ กามานเมว เหตุ.}

\prose{\textbf{อาทีนวํ กามคุเณสุ ทิสฺวา}ติ กามคุเณสุ อาทีนวํ ทิสฺวา ปสฺสิตฺวา ตุลยิตฺวา ตีรยิตฺวา วิภาวยิตฺวา วิภูตํ กตฺวาติ อาทีนวํ กามคุเณสุ ทิสฺวา, เอโก จเร ขคฺควิสาณ\-กปฺโป. เตนาห โส ปจฺเจก\-สมฺพุทฺโธ\csromandash{}}

\setlength{\csromangathapagewidth}{0pt}
\setlength{\csromangathawakwidth}{0pt}
\csromangathameasuregroup{}{กามา หิ จิตฺรา มธุรา มโนรมา, \\ วิรูปรูเปน มเถนฺติ จิตฺตํ. \\ อาทีนวํ กามคุเณสุ ทิสฺวา, \\ เอโก จเร ขคฺควิสาณกปฺโปติ. \\ \textbf{¢ตี จ คณฺโฑ จ อุปทฺทโว จ,} \\ \textbf{โรโค จ สลฺลญฺจ ภยญฺจ เมตํ.} \\ \textbf{เอตํ ภยํ กามคุเณสุ ทิสฺวา,} \\ \textbf{เอโก จเร ขคฺควิสาณ}\textbf{กปฺโป.}}
\csromangathameasurewak{กามา หิ จิตฺรา มธุรา มโนรมา, \\ วิรูปรูเปน มเถนฺติ จิตฺตํ. \\ อาทีนวํ กามคุเณสุ ทิสฺวา, \\ เอโก จเร ขคฺควิสาณกปฺโปติ. \\ \textbf{¢ตี จ คณฺโฑ จ อุปทฺทโว จ,} \\ \textbf{โรโค จ สลฺลญฺจ ภยญฺจ เมตํ.} \\ \textbf{เอตํ ภยํ กามคุเณสุ ทิสฺวา,} \\ \textbf{เอโก จเร ขคฺควิสาณ}\textbf{กปฺโป.}}
\setlength{\csromangathaleftcol}{0pt}
\csromangathagroup{\csromangathastanza{\csromangathawak{กามา หิ จิตฺรา มธุรา มโนรมา, \\ วิรูปรูเปน มเถนฺติ จิตฺตํ. \\ อาทีนวํ กามคุเณสุ ทิสฺวา, \\ เอโก จเร ขคฺควิสาณ\-กปฺโปติ.}}\csromangathastanza{\csromangathawak{\csromanfolio{265}\csromangathaitemline{137}{\textbf{¢ตี จ คณฺโฑ จ อุปทฺทโว จ,}} \\ \csromangathacontline{\textbf{โรโค จ สลฺลญฺจ ภยญฺจ เมตํ.}} \\ \csromangathacontline{\textbf{เอตํ ภยํ กามคุเณสุ ทิสฺวา,}} \\ \csromangathacontline{\textbf{เอโก จเร ขคฺควิสาณ}\-\textbf{กปฺโป.}}}}}

\prose{\textbf{¢ตี จ คณฺโฑ จ อุปทฺทโว จ, โรโค จ สลฺลญฺจ ภยญฺจ เมตนฺ}ติ วุตฺตํ เหตํ ภควตา\csromandash{} “ภยนฺ”ติ ภิกฺขเว กามานเมตํ อธิวจนํ, “ทุกฺขนฺ”ติ ภิกฺขเว กามานเมตํ อธิวจนํ, “โรโค”ติ ภิกฺขเว กามานเมตํ อธิวจนํ, “คณฺโฑ”ติ ภิกฺขเว กามานเมตํ อธิวจนํ. “สลฺลนฺ”ติ ภิกฺขเว กามานเมตํ อธิวจนํ, “สงฺโค”ติ ภิกฺขเว กามานเมตํ อธิวจนํ, “ปงฺโก”ติ ภิกฺขเว กามานเมตํ อธิวจนํ, “คพฺโภ”ติ ภิกฺขเว กามานเมตํ อธิวจนํ. กสฺมา จ ภิกฺขเว “ภยนฺ”ติ กามานเมตํ อธิวจนํ. ยสฺมา จ กามราครตฺตายํ ภิกฺขเว ฉนฺทราค\-วินิพทฺโธ\csromansharedfootnote{0e9fbc64144b77c6}{ฉนฺทราควินิพนฺโธ (สฺยา, ก) ปสฺส อํ 3. 114 ปิฏฺเฐ.} ทิฏฺฐธมฺมิกาปิ ภยา น ปริมุจฺจติ, สมฺปรายิกาปิ ภยา น ปริมุจฺจติ, ตสฺมา “ภยนฺ”ติ กามานเมตํ อธิวจนํ. กสฺมา จ ภิกฺขเว “ทุกฺขนฺ”ติ. “โรโค”ติ. “คณฺโฑ”ติ. “สลฺลนฺ”ติ. “สงฺโค”ติ. “ปงฺโก”ติ. “คพฺโภ”ติ กามานเมตํ อธิวจนํ. ยสฺมา จ กามราครตฺตายํ ภิกฺขเว ฉนฺทราค\-วินิพทฺโธ ทิฏฺฐธมฺมิกาปิ คพฺภา น ปริมุจฺจติ, สมฺปรายิกาปิ คพฺภา น ปริมุจฺจติ, ตสฺมา “คพฺโภ”ติ กามานเมตํ อธิวจนนฺติ.}

\setlength{\csromangathapagewidth}{0pt}
\setlength{\csromangathawakwidth}{0pt}
\csromangathameasuregroup{ภยํ ทุกฺขญฺจ โรโค จ, \\ ปงฺโก คพฺโภ จ อุภยํ, \\ โอติณฺโณ สาตรูเปน, \\ ยโต จ ภิกฺขุ อาตาปี, \\ โส อิมํ ปลิปถํ ทุคฺคํ,}{\csromangathabat{ภยํ ทุกฺขญฺจ โรโค จ,}{คณฺโฑ สลฺลญฺจ สงฺโค จ.} \\ \csromangathabat{ปงฺโก คพฺโภ จ อุภยํ,}{เอเต กามา ปวุจฺจนฺติ.} \\ ยตฺถ สตฺโต ปุถุชฺชโน. \\ \csromangathabat{โอติณฺโณ สาตรูเปน,}{ปุน คพฺภาย คจฺฉติ.} \\ \csromangathabat{ยโต จ ภิกฺขุ อาตาปี,}{สมฺปชญฺญํ น ริจฺจติ\textsuperscript{1}.} \\ \csromangathabat{โส อิมํ ปลิปถํ ทุคฺคํ,}{อติกฺกมฺม ตถาวิโธ.}}
\csromangathameasurewak{ยตฺถ สตฺโต ปุถุชฺชโน.}
\csromangathasetleft{ภยํ ทุกฺขญฺจ โรโค จ, \\ ปงฺโก คพฺโภ จ อุภยํ, \\ โอติณฺโณ สาตรูเปน, \\ ยโต จ ภิกฺขุ อาตาปี, \\ โส อิมํ ปลิปถํ ทุคฺคํ,}
\csromangathagroup{\csromangathastanza{\csromangathabat{ภยํ ทุกฺขญฺจ โรโค จ,}{คณฺโฑ สลฺลญฺจ สงฺโค จ.} \\ \csromangathabat{ปงฺโก คพฺโภ จ อุภยํ,}{เอเต กามา ปวุจฺจนฺติ.}}\csromangathastanza{\csromangathawak{ยตฺถ สตฺโต ปุถุชฺชโน.}}\csromangathastanza{\csromangathabat{โอติณฺโณ สาตรูเปน,}{ปุน คพฺภาย คจฺฉติ.} \\ \csromangathabat{ยโต จ ภิกฺขุ อาตาปี,}{สมฺปชญฺญํ น ริจฺจติ\csromansharedfootnote{61a4a57c16d3ad5e}{น ริญฺจติ (สฺยา, ก) สํ 2. 408 ปิฏฺเฐปิ.}.} \\ \csromangathabat{โส อิมํ ปลิปถํ ทุคฺคํ,}{อติกฺกมฺม ตถาวิโธ.}}}

\prose{ปชํ ชาติชรูเปตํ, ผนฺทมานํ อเวกฺขตีติ\csromandash{}}

\prose{\textbf{¢ตี จ คณฺโฑ จ อุปทฺทโว จ,}}

\prose{\textbf{โรโค จ สลฺลญฺจ ภยญฺจ เมตํ.}}

\prose{\csromanfolio{266}\textbf{เอตํ ภยํ กามคุเณสุ ทิสฺวา}ติ เอตํ ภยํ กามคุเณสุ ทิสฺวา ปสฺสิตฺวา ตุลยิตฺวา ตีรยิตฺวา วิภาวยิตฺวา วิภูตํ กตฺวาติ เอตํ ภยํ กามคุเณสุ ทิสฺวา, เอโก จเร ขคฺควิสาณ\-กปฺโป. เตนาห โส ปจฺเจก\-สมฺพุทฺโธ\csromandash{}}

\setlength{\csromangathapagewidth}{0pt}
\setlength{\csromangathawakwidth}{0pt}
\csromangathameasuregroup{}{¢ตี จ คณฺโฑ จ อุปทฺทโว จ, \\ โรโค จ สลฺลญฺจ ภยญฺจ เมตํ.}
\csromangathameasurewak{¢ตี จ คณฺโฑ จ อุปทฺทโว จ, \\ โรโค จ สลฺลญฺจ ภยญฺจ เมตํ.}
\setlength{\csromangathaleftcol}{0pt}
\csromangathagroup{\csromangathastanza{\csromangathawak{¢ตี จ คณฺโฑ จ อุปทฺทโว จ, \\ โรโค จ สลฺลญฺจ ภยญฺจ เมตํ.}}}

\prose{เอตํ ภยํ กามคุเณสุ ทิสฺวา,}

\prose{เอโก จเร ขคฺควิสาณ\-กปฺโปติ.}

\setlength{\csromangathapagewidth}{0pt}
\setlength{\csromangathawakwidth}{0pt}
\csromangathameasuregroup{}{\textbf{สีตญฺจ อุณฺหญฺจ ขุทํ ปิปาสํ,} \\ \textbf{วาตาตเป ฑํสสรีสเป}\textsuperscript{1}\textbf{ จ.} \\ \textbf{สพฺพานิเป’ตานิ อภิสมฺ}\textbf{ภวิ}\textbf{ตฺวา,} \\ \textbf{เอโก จเร ขคฺควิสาณ}\textbf{กปฺโป.}}
\csromangathameasurewak{\textbf{สีตญฺจ อุณฺหญฺจ ขุทํ ปิปาสํ,} \\ \textbf{วาตาตเป ฑํสสรีสเป}\textsuperscript{1}\textbf{ จ.} \\ \textbf{สพฺพานิเป’ตานิ อภิสมฺ}\textbf{ภวิ}\textbf{ตฺวา,} \\ \textbf{เอโก จเร ขคฺควิสาณ}\textbf{กปฺโป.}}
\setlength{\csromangathaleftcol}{0pt}
\csromangathagroup{\csromangathastanza{\csromangathawak{\csromangathaitemline{138}{\textbf{สีตญฺจ อุณฺหญฺจ ขุทํ ปิปาสํ,}} \\ \csromangathacontline{\textbf{วาตาตเป ฑํสสรีสเป}\csromansharedfootnote{0afe0943badb6cf9}{qอํสสิริํสเป (สฺยา), ฑํสมกสสรีสเป (ก)}\textbf{ จ.}} \\ \csromangathacontline{\textbf{สพฺพานิเป’ตานิ อภิสมฺ}\-\textbf{ภวิ}\-\textbf{ตฺวา,}} \\ \csromangathacontline{\textbf{เอโก จเร ขคฺควิสาณ}\-\textbf{กปฺโป.}}}}}

\prose{\textbf{สีตญฺจ อุณฺหญฺจ ขุทํ ปิปาสนฺ}ติ \textbf{สีตนฺ}ติ ทฺวีหิ การเณหิ สีตํ โหติ, อพฺภนฺตร\-ธาตุ\-ปฺปโกป\-วเสน วา สีตํ โหติ, พหิทฺธา อุตุวเสน วา สีตํ โหติ. \textbf{อุณฺหนฺ}ติ ทฺวีหิ การเณหิ อุณฺหํ โหติ, อพฺภนฺตร\-ธาตุ\-ปฺปโกป\-วเสน วา อุณฺหํ โหติ, พหิทฺธา อุตุวเสน วา อุณฺหํ โหติ. \textbf{ขุทา}\csromansharedfootnote{a0f00a3852ccbd68}{ขุทฺทา (สฺยา, ก)} วุจฺจติ ฉาตโก. \textbf{ปิปาสา} วุจฺจติ อุทกปิปาสาติ สีตญฺจ อุณฺหญฺจ ขุทํ ปิปาสํ.}

\prose{\textbf{วาตาตเป ฑํสสรีสเป จา}ติ วาตาติ ปุรตฺถิมา \textbf{วาตา} ปจฺฉิมา วาตา อุตฺตรา วาตา ทกฺขิณา วาตา สรชา วาตา อรชา วาตา สีตา วาตา อุณฺหา วาตา ปริตฺตา วาตา อธิมตฺตา วาตา เวรมฺภวาตา ปกฺขวาตา สุปณฺณวาตา ตาลปณฺณวาตา\csromansharedfootnote{1da33735307b1318}{ตาลวณฺฑวาตา (ก) 109 ปิฏฺเฐ นตฺถิ ปาฐนานตฺตํ.} วิธูปนวาตา. \textbf{อาตโป} วุจฺจติ สูริยสนฺตาโป. q\textbf{อํสา} วุจฺจนฺติ ปิงฺคลมกฺขิกา. \textbf{สรีสปา} วุจฺจนฺติ อหีติ วาตาตเป ฑํสสรีสเป จ.}

\prose{\textbf{สพฺพานิเป’ตานิ อภิสมฺ}\-\textbf{ภวิ}\-\textbf{ตฺวา}ติ อภิภวิตฺวา อชฺโฌตฺถริตฺวา ปริยาทิยิตฺวา มทฺทิตฺวาติ สพฺพานิเป’ตานิ อภิสมฺ\-ภวิ\-ตฺวา, เอโก จเร ขคฺควิสาณ\-กปฺโป. เตนาห โส ปจฺเจก\-สมฺพุทฺโธ\csromandash{}}

\setlength{\csromangathapagewidth}{0pt}
\setlength{\csromangathawakwidth}{0pt}
\csromangathameasuregroup{}{สีตญฺจ อุณฺหญฺจ ขุทํ ปิปาสํ, \\ วาตาตเป ฑํสสรีสเป จ. \\ สพฺพานิเป’ตานิ อภิสมฺภวิตฺวา, \\ เอโก จเร ขคฺควิสาณกปฺโปติ.}
\csromangathameasurewak{สีตญฺจ อุณฺหญฺจ ขุทํ ปิปาสํ, \\ วาตาตเป ฑํสสรีสเป จ. \\ สพฺพานิเป’ตานิ อภิสมฺภวิตฺวา, \\ เอโก จเร ขคฺควิสาณกปฺโปติ.}
\setlength{\csromangathaleftcol}{0pt}
\csromangathagroup{\csromangathastanza{\csromangathawak{\csromanfolio{267}สีตญฺจ อุณฺหญฺจ ขุทํ ปิปาสํ, \\ วาตาตเป ฑํสสรีสเป จ. \\ สพฺพานิเป’ตานิ อภิสมฺ\-ภวิ\-ตฺวา, \\ เอโก จเร ขคฺควิสาณ\-กปฺโปติ.}}}

\proseitem{139}{\textbf{นาโคว ยูถานิ วิวชฺชยิตฺวา,}}

\prose{\textbf{สญฺชาตขนฺโธ ปทุมี อุฬาโร.}}

\prose{\textbf{ยถาภิรนฺตํ วิหเร}\csromansharedfootnote{314c72c774eb37d3}{วีหรํ (ขุ 1. 287 ปิฏฺเฐ.)}\textbf{ อรญฺเญ,}}

\prose{\textbf{เอโก จเร ขคฺควิสาณ}\-\textbf{กปฺโป.} (9)}

\prose{\textbf{นาโคว ยูถานิ วิวชฺชยิตฺวา}ติ นาโค วุจฺจติ หตฺถินาโค, ปจฺเจก\-สมฺพุทฺโธปิ นาโค. กิํการณา ปจฺเจก\-สมฺพุทฺโธ นาโค. อาคุํ น กโรตีติ นาโค. น คจฺฉตีติ นาโค. น อาคจฺฉตีติ นาโค. กถํ โส ปจฺเจก\-สมฺพุทฺโธ อาคุํ น กโรตีติ นาโค. อาคุ วุจฺจติ ปาปกา อกุสลา ธมฺมา สํกิเลสิกา โปโนภวิกา สทรา ทุกฺขวิปากา อายติํ ชาติ\-ชรา\-มรณิยา.}

\setlength{\csromangathapagewidth}{0pt}
\setlength{\csromangathawakwidth}{0pt}
\csromangathameasuregroup{}{อาคุํ น กโรติ กิญฺจิ โลเก, \\ สพฺพสํโยเค วิสชฺช พนฺธนานิ. \\ สพฺพตฺถ น สชฺชตี วิมุตฺโต, \\ นาโค ตาทิ ปวุจฺจเต ตถตฺตา.}
\csromangathameasurewak{อาคุํ น กโรติ กิญฺจิ โลเก, \\ สพฺพสํโยเค วิสชฺช พนฺธนานิ. \\ สพฺพตฺถ น สชฺชตี วิมุตฺโต, \\ นาโค ตาทิ ปวุจฺจเต ตถตฺตา.}
\setlength{\csromangathaleftcol}{0pt}
\csromangathagroup{\csromangathastanza{\csromangathawak{อาคุํ น กโรติ กิญฺจิ โลเก, \\ สพฺพสํโยเค วิสชฺช พนฺธนานิ. \\ สพฺพตฺถ น สชฺชตี วิมุตฺโต, \\ นาโค ตาทิ ปวุจฺจเต ตถตฺตา.}}}

\prose{เอวํ โส ปจฺเจก\-สมฺพุทฺโธ อาคุํ น กโรตีติ นาโค.}

\prose{กถํ โส ปจฺเจก\-สมฺพุทฺโธ น คจฺฉตีติ นาโค, โส ปจฺเจก\-สมฺพุทฺโธ น ฉนฺทาคติํ คจฺฉติ, น โทสาคติํ คจฺฉติ, น โมหาคติํ คจฺฉติ, น ภยาคติํ คจฺฉติ, น ราควเสน คจฺฉติ, น โทสวเสน คจฺฉติ, น โมหวเสน คจฺฉติ, น มานวเสน คจฺฉติ, น ทิฏฺฐิวเสน คจฺฉติ, น อุทฺธจฺจวเสน คจฺฉติ, น วิจิกิจฺฉา\-วเสน คจฺฉติ, น อนุสยวเสน คจฺฉติ, น วคฺเคหิ กปฺเปหิ ยายติ นียติ\csromansharedfootnote{ee843794401f93eb}{นิยฺยติ (สฺยา, ก)} วุยฺหติ สํหรียติ. เอวํ โส ปจฺเจก\-สมฺพุทฺโธ น คจฺฉตีติ นาโค.}

\prose{\csromanfolio{268}กถํ โส ปจฺเจก\-สมฺพุทฺโธ น อาคจฺฉตีติ นาโค, โสตาปตฺติ มคฺเคน เย กิเลสา ปหีนา, เต กิเลเส น ปุเนติ น ปจฺเจติ น ปจฺจาคจฺฉติ. สกทาคามิ\-มคฺเคน ฯเปฯ อนาคามิมคฺเคน ฯเปฯ อรหตฺต\-มคฺเคน เย กิเลสา ปหีนา, เต กิเลเส น ปุเนติ น ปจฺเจติ น ปจฺจาคจฺฉติ. เอวํ โส ปจฺเจก\-สมฺพุทฺโธ น อาคจฺฉตีติ นาโค.}

\prose{\textbf{นาโคว ยูถานิ วิวชฺชยิตฺวา}ติ ยถา โส หตฺถินาโค ยูถานิ วิวชฺเชตฺวา ปริวชฺเชตฺวา อภินิวชฺเชตฺวา เอโกว อรญฺญวนม\-ชฺโฌคาเหตฺวา\csromansharedfootnote{5189d98623fdd687}{อรญฺเญ วนมชฺฌสฺส อชฺโฌคาเหตฺวา (สฺยา)} จรติ วิหรติ อิริยติ วตฺเตติ ปาเลติ ยเปติ ยาเปติ. ปจฺเจก\-สมฺพุทฺโธปิ คณํ วิวชฺเชตฺวา ปริวชฺเชตฺวา อภิวชฺเชตฺวา เอโก\csromansharedfootnote{153f28cf7b4a99c8}{เอโก จเร ขคฺควิสาณกปฺโป (สฺยา)} อรญฺญ\-วนปตฺถานิ ปนฺตานิ เสนาสนานิ ปฏิเสวติ อปฺปสทฺทานิ อปฺปนิคฺโฆสานิ วิชนวาตานิ มนุสฺส\-ราหสฺเสยฺยกานิ ปฏิสลฺลาน\-สารุปฺปานิ. โส เอโก คจฺฉติ, เอโก ติฏฺฐติ, เอโก นิสีทติ, เอโก เสยฺยํ กปฺเปติ, เอโก คามํ ปิณฺฑาย ปวิสติ, เอโก ปฏิกฺกมติ, เอโก รโห นิสีทติ, เอโก จงฺกมํ อธิฏฺฐาติ, เอโก จรติ วิหรติ อิริยติ วตฺเตติ ปาเลติ ยเปติ ยาเปตีติ นาโคว ยูถานิ วิวชฺชยิตฺวา.}

\prose{\textbf{สญฺชาตขนฺโธ ปทุมี อุฬาโร}ติ ยถา โส หตฺถินาโค สญฺชาต\-กฺขนฺโธ สตฺตรตโน วา โหติ อฏฺฐรตโน วา. ปจฺเจก\-สมฺพุทฺโธปิ สญฺชาต\-กฺขนฺโธ อเสกฺเขน สีลกฺขนฺเธน, อเสกฺเขน สมาธิ\-กฺขนฺเธน, อเสกฺเขน ปญฺญา\-กฺขนฺเธน, อเสกฺเขน วิมุตฺติ\-กฺขนฺเธน, อเสกฺเขน วิมุตฺติ\-ญาณ\-ทสฺสนกฺขนฺเธน. ยถา โส หตฺถินาโค ปทุมี, ปจฺเจก\-สมฺพุทฺโธปิ สตฺตหิ โพชฺฌงฺค\-ปุปฺเผหิ ปทุมี สติสมฺโพชฺฌงฺค\-ปุปฺเผน ธมฺมวิจยสมฺโพชฺฌงฺค\-ปุปฺเผน วีริยสมฺโพชฺฌงฺค\-ปุปฺเผน ปีติสมฺโพชฺฌงฺค\-ปุปฺเผน ปสฺสทฺธิสมฺโพชฺฌงฺค\-ปุปฺเผน สมาธิสมฺโพชฺฌงฺค\-ปุปฺเผน อุเปกฺขาสมฺโพชฺฌงฺค\-ปุปฺเผน. ยถา โส หตฺถินาโค อุฬาโร ถาเมน พเลน ชเวน สูเรน, ปจฺเจก\-สมฺพุทฺโธปิ อุฬาโร สีเลน สมาธินา ปญฺญาย วิมุตฺติยา วิมุตฺติ\-ญาณ\-ทสฺสเนนาติ สญฺชาตขนฺโธ ปทุมี อุฬาโร.}

\prose{\textbf{ยถาภิรนฺตํ วิหเร อรญฺเญ}ติ ยถา โส หตฺถินาโค ยถาภิรนฺตํ อรญฺเญ วิหรติ. ปจฺเจก\-สมฺพุทฺโธปิ ยถาภิรนฺตํ อรญฺเญ วิหรติ.}

\prose{\csromanfolio{269}ปฐเมนปิ ฌาเนน ยถาภิรนฺตํ อรญฺเญ วิหรติ. ทุติเยนปิ ฌาเนน. ตติเยนปิ ฌาเนน. จตุตฺเถนปิ ฌาเนน ยถาภิรนฺตํ อรญฺเญ วิหรติ. เมตฺตายปิ เจโตวิมุตฺติยา ยถาภิรนฺตํ อรญฺเญ วิหรติ. กรุณายปิ เจโตวิมุตฺติยา. มุทิตายปิ เจโตวิมุตฺติยา. อุเปกฺขายปิ เจโตวิมุตฺติยา ยถาภิรนฺตํ อรญฺเญ วิหรติ. อากาสานญฺจายตนมาปตฺติยาปิ ยถาภิรนฺตํ อรญฺเญ วิหรติ. วิญฺญาณญฺจายตน\-สมาปตฺติยาปิ. อากิญฺจญฺญายตน\-สมาปตฺติยาปิ. เนว\-สญฺญา\-นา\-สญฺญา\-ยตนสมาปตฺติยาปิ. นิโรธ\-สมาปตฺติยาปิ. ผลสมาปตฺติยาปิ ยถาภิรนฺตํ อรญฺเญ วิหรตีติ ยถาภิรนฺตํ วิหเร อรญฺเญ, เอโก จเร ขคฺควิสาณ\-กปฺโป. เตนาห โส ปจฺเจก\-สมฺพุทฺโธ\csromandash{}}

\setlength{\csromangathapagewidth}{0pt}
\setlength{\csromangathawakwidth}{0pt}
\csromangathameasuregroup{}{นาโคว ยูถานิ วิวชฺชยิตฺวา, \\ สญฺชาตขนฺโธ ปทุมี อุฬาโร. \\ ยถาภิรนฺตํ วิหเร อรญฺเญ, \\ เอโก จเร ขคฺควิสาณกปฺโปติ.}
\csromangathameasurewak{นาโคว ยูถานิ วิวชฺชยิตฺวา, \\ สญฺชาตขนฺโธ ปทุมี อุฬาโร. \\ ยถาภิรนฺตํ วิหเร อรญฺเญ, \\ เอโก จเร ขคฺควิสาณกปฺโปติ.}
\setlength{\csromangathaleftcol}{0pt}
\csromangathagroup{\csromangathastanza{\csromangathawak{นาโคว ยูถานิ วิวชฺชยิตฺวา, \\ สญฺชาตขนฺโธ ปทุมี อุฬาโร. \\ ยถาภิรนฺตํ วิหเร อรญฺเญ, \\ เอโก จเร ขคฺควิสาณ\-กปฺโปติ.}}}

\proseitem{140}{\textbf{อฏฺฐาน’ตํ สงฺคณิการตสฺส,}}

\prose{\textbf{ยํ ผสฺสเย}\csromansharedfootnote{42b4d0a431f05529}{ผุสฺสเย (สฺยา, ก)}\textbf{ สามยิกํ}\csromansharedfootnote{329b68bc1e65c9bf}{อาสามายิกํ (ก)}\textbf{ วิมุตฺติํ.}}

\prose{\textbf{อาทิจฺจ}\-\textbf{พนฺธุสฺส วโจ นิสมฺม,}}

\prose{\textbf{เอโก จเร ขคฺควิสาณ}\-\textbf{กปฺโป.} (10)}

\prose{\textbf{อฏฺฐาน’ตํ สงฺคณิการตสฺส, ยํ ผสฺสเย สามยิกํ วิมุตฺตินฺ}ติ วุตฺตํ เหตํ ภควตา\csromandash{}ยาวตานนฺท\csromansharedfootnote{61a65db86faade4a}{โส วตานนฺท (ม 3. 152 ปิฏฺเฐ.)} ภิกฺขุ สงฺคณิการาโม สงฺคณิกรโต สงฺคณิกา\-รามตํ อนุยุตฺโต คณาราโม คณรโต คณสมฺมุทิโต (คณารามตํ อนุยุตฺโต)\csromansharedfootnote{81e0cf3b802dff59}{( ) นตฺถิ ม 3. 152 ปิฏฺเฐ.}, ยํ ตํ เนกฺขมฺมสุขํ ปวิเวกสุขํ อุปสมสุขํ สมฺโพธิสุขํ, ตสฺส สุขสฺส นิกามลาภี ภวิสฺสติ อกิจฺฉลาภี อกสิรลาภีติ เนตํ ฐานํ วิชฺชติ. โย จ โข โส อานนฺท ภิกฺขุ เอโก คณสฺมา วูปกฏฺโฐ วิหรติ, ตสฺเสตํ ภิกฺขุโน ปาฏิกงฺขํ. ยํ ตํ เนกฺขมฺมสุขํ ปวิเวกสุขํ อุปสมสุขํ สมฺโพธิสุขํ, ตสฺส สุขสฺส นิกามลาภี ภวิสฺสติ อกิจฺฉลาภี อกสิรลาภีติ ฐานเมตํ วิชฺชติ. ยาวตานนฺท ภิกฺขุ สงฺคณิการาโม สงฺคณิกรโต}

\prose{\csromanfolio{270}สงฺคณิกา\-รามตํ อนุยุตฺโต คณาราโม คณรโต คณสมฺมุทิโต (คณารามตํ อนุยุตฺโต,) สามายิกํ\csromansharedfootnote{0071b83476255243}{สามยิกํ (สฺยา, ก)} วา กนฺตํ เจโตวิมุตฺติํ อุปสมฺปชฺช วิหริสฺสติ, อสามายิกํ วา อกุปฺปนฺติ เนตํ ฐานํ วิชฺชติ. โย จ โข โส อานนฺท ภิกฺขุ เอโก คณสฺมา วูปกฏฺโฐ วิหรติ, ตสฺเสตํ ภิกฺขุโน ปาฏิกงฺขํ สามายิกํ วา กนฺตํ เจโตวิมุตฺติํ อุปสมฺปชฺช วิหริสฺสติ, อสามายิกํ วา อกุปฺปนฺติ ฐานเมตํ วิชฺชตีติ อฏฺฐาน’ตํ สงฺคณิการตสฺส, ยํ ผสฺสเย สามยิกํ วิมุตฺติํ.}

\prose{\textbf{อาทิจฺจ}\-\textbf{พนฺธุสฺส วโจ นิสมฺมา}ติ อาทิจฺโจ วุจฺจติ สูริโย, โส โคตโม โคตฺเตน. ปจฺเจก\-สมฺพุทฺโธปิ โคตโม โคตฺเตน. โส ปจฺเจก\-สมฺพุทฺโธ สูริยสฺส โคตฺตญาตโก โคตฺตพนฺธุ, ตสฺมา ปจฺเจก\-สมฺพุทฺโธ อาทิจฺจพนฺธุ. \textbf{อาทิจฺจ}\-\textbf{พนฺธุสฺส วโจ นิสมฺมา}ติ อาทิจฺจ\-พนฺธุสฺส วจนํ พฺยปฺปถํ เทสนํ อนุสาสนํ อนุสิฏฺฐํ สุตฺวา สุณิตฺวา อุคฺคเหตฺวา อุปธารยิตฺวา อุปลกฺขยิตฺวาติ อาทิจฺจ\-พนฺธุสฺส วโจ นิสมฺม, เอโก จเร ขคฺควิสาณ\-กปฺโป. เตนาห โส ปจฺเจก\-สมฺพุทฺโธ\csromandash{}}

\setlength{\csromangathapagewidth}{0pt}
\setlength{\csromangathawakwidth}{0pt}
\csromangathameasuregroup{}{อฏฺฐาน’ตํ สงฺคณิการตสฺส, \\ ยํ ผสฺสเย สามยิกํ วิมุตฺติํ. \\ อาทิจฺจพนฺธุสฺส วโจ นิสมฺม, \\ เอโก จเร ขคฺควิสาณกปฺโปติ.}
\csromangathameasurewak{อฏฺฐาน’ตํ สงฺคณิการตสฺส, \\ ยํ ผสฺสเย สามยิกํ วิมุตฺติํ. \\ อาทิจฺจพนฺธุสฺส วโจ นิสมฺม, \\ เอโก จเร ขคฺควิสาณกปฺโปติ.}
\setlength{\csromangathaleftcol}{0pt}
\csromangathagroup{\csromangathastanza{\csromangathawak{อฏฺฐาน’ตํ สงฺคณิการตสฺส, \\ ยํ ผสฺสเย สามยิกํ วิมุตฺติํ. \\ อาทิจฺจ\-พนฺธุสฺส วโจ นิสมฺม, \\ เอโก จเร ขคฺควิสาณ\-กปฺโปติ.}}}

\vspace{\tipitakaparskip}
\csromancenter{ทุติโย วคฺโค.}
\chapterhead{\textbf{ตติยวคฺค}}
\setlength{\csromangathapagewidth}{0pt}
\setlength{\csromangathawakwidth}{0pt}
\csromangathameasuregroup{}{\textbf{ทิฏฺฐีวิสูกานิ อุปาติวตฺโต,} \\ \textbf{ปตฺโต นิยามํ ปฏิลทฺธ}\textbf{มคฺโค.} \\ \textbf{อุปฺปนฺนญาโณมฺหิ อนญฺญเนยฺโย,} \\ \textbf{เอโก จเร ขคฺควิสาณ}\textbf{กปฺโป.}}
\csromangathameasurewak{\textbf{ทิฏฺฐีวิสูกานิ อุปาติวตฺโต,} \\ \textbf{ปตฺโต นิยามํ ปฏิลทฺธ}\textbf{มคฺโค.} \\ \textbf{อุปฺปนฺนญาโณมฺหิ อนญฺญเนยฺโย,} \\ \textbf{เอโก จเร ขคฺควิสาณ}\textbf{กปฺโป.}}
\setlength{\csromangathaleftcol}{0pt}
\csromangathagroup{\csromangathastanza{\csromangathawak{\csromangathaitemline{141}{\textbf{ทิฏฺฐีวิสูกานิ อุปาติวตฺโต,}} \\ \csromangathacontline{\textbf{ปตฺโต นิยามํ ปฏิลทฺธ}\-\textbf{มคฺโค.}} \\ \csromangathacontline{\textbf{อุปฺปนฺนญาโณมฺหิ อนญฺญเนยฺโย,}} \\ \csromangathacontline{\textbf{เอโก จเร ขคฺควิสาณ}\-\textbf{กปฺโป.}}}}}

\prose{\textbf{ทิฏฺฐีวิสูกานิ อุปาติวตฺโต}ติ ทิฏฺฐิวิสูกานิ วุจฺจนฺติ วีสติ วตฺถุกา สกฺกายทิฏฺฐี. อิธ อสฺสุตวา ปุถุชฺชโน อริยานํ อทสฺสาวี อริยธมฺมสฺส}

\prose{\csromanfolio{271}อโกวิโท อริยธมฺเม อวินีโต สปฺปุริสานํ อทสฺสาวี สปฺปุริส\-ธมฺมสฺส อโกวิโท สปฺปุริส\-ธมฺเม อวินีโต รูปํ อตฺตโต สมนุปสฺสติ รูปวนฺตํ วา อตฺตานํ, อตฺตนิ วา รูปํ, รูปสฺมิํ วา อตฺตานํ. เวทนํ. สญฺญํ. สงฺขาเร. วิญฺญาณํ อตฺตโต สมนุปสฺสติ, วิญฺญาณวนฺตํ วา อตฺตานํ, อตฺตนิ วา วิญฺญาณํ, วิญฺญาณสฺมิํ วา อตฺตนํ. ยา เอวรูปา ทิฏฺฐิ ทิฏฺฐิคตํ ทิฏฺฐิคหนํ ทิฏฺฐิกนฺตาโร ทิฏฺฐิ\-วิสูกายิกํ ทิฏฺฐิ\-วิปฺผนฺทิตํ ทิฏฺฐิ\-สํโยชนํ คาโห ปฏิคฺคาโห อภินิเวโส ปรามาโส กุมฺมคฺโค มิจฺฉาปโถ มิจฺฉตฺตํ ติตฺถายตนํ วิปริเยส\-คฺคาโห วิปรีตคฺคาโห วิปลฺลาสคฺคาโห มิจฺฉาคาโห อยาถาวกสฺมิํ “ยาถาวกนฺ”ติ คาโห, ยาวตา ทฺวาสฏฺฐิ ทิฏฺฐิคตานิ, อิมานิ ทิฏฺฐิวิสูกานิ. \textbf{ทิฏฺฐีวิสูกานิ อุปาติวตฺโต}ติ ทิฏฺฐิวิสูกานิ อุปาติวตฺโต อติกฺกนฺโต สมติกฺกนฺโต วีติวตฺโตติ ทิฏฺฐีวิสูกานิ อุปาติวตฺโต.}

\prose{\textbf{ปตฺโต นิยามํ ปฏิลทฺธ}\-\textbf{มคฺโค}ติ นิยามา วุจฺจนฺติ จตฺตาโร มคฺคา. อริโย อฏฺฐงฺคิโก มคฺโค. เสยฺยถิทํ, สมฺมาทิฏฺฐิสมฺมาสงฺกปฺโป สมฺมาวาจา สมฺมากมฺมนฺโต สมฺมาอาชีโว สมฺมาวายาโม สมฺมาสติ สมฺมาสมาธิ. จตูหิ อริยมคฺเคหิ สมนฺนาคโต นิยามํ ปตฺโต สมฺปตฺโต อธิคโต ผสฺสิโต สจฺฉิกโตติ ปตฺโต นิยามํ. \textbf{ปฏิลทฺธ}\-\textbf{มคฺโค}ติ ลทฺธมคฺโค ปฏิลทฺธ\-มคฺโค อธิคตมคฺโค ผสฺสิตมคฺโค สจฺฉิกตมคฺโคติ ปตฺโต นิยามํ ปฏิลทฺธ\-มคฺโค.}

\prose{\textbf{อุปฺปนฺนญาโณมฺหิ อนญฺญเนยฺโย}ติ ตสฺส ปจฺเจก\-สมฺพุทฺธสฺส ญาณํ อุปฺปนฺนํ สมุปฺปนฺนํ นิพฺพตฺตํ อภินิพฺพตฺตํ ปาตุภูตํ “สพฺเพ สงฺขารา อนิจฺจา”ติ ญาณํ อุปฺปนฺนํ สมุปฺปนฺนํ นิพฺพตฺตํ อภินิพฺพตฺตํ ปาตุภูตํ. “สพฺเพ สงฺขารา ทุกฺขา”ติ ฯเปฯ. “สพฺเพ ธมฺมา อนตฺตา”ติ ฯเปฯ. “ยํ กิญฺจิ สมุทย\-ธมฺมํ, สพฺพํ ตํ นิโรธธมฺมนฺ”ติ ญาณํ อุปฺปนฺนํ สมุปฺปนฺนํ นิพฺพตฺตํ อภินิพฺพตฺตํ ปาตุภูตนฺติ อุปฺปนฺนญาโณมฺหิ. \textbf{อนญฺญเนยฺโย}ติ โส ปจฺเจก\-สมฺพุทฺโธ น ปรเนยฺโย น ปรปฺปตฺติโย น ปรปฺปจฺจโย น ปรปฏิพทฺธคู, ยถาภูตํ\csromansharedfootnote{672b3cfb0804b7fc}{ตํ (ก)} ชานาติ ปสฺสติ อสมฺมูโฬฺห สมฺปชาโน ปฏิสฺสโต. “สพฺเพ สงฺขารา อนิจฺจา”ติ น ปรเนยฺโย น ปรปฺปตฺติโย น ปรปฺปจฺจโย น ปรปฏิพทฺธคู, ยถาภูตํ\csromansharedfootnote{672b3cfb0804b7fc}{ตํ (ก)} ชานาติ ปสฺสติ อสมฺมูโฬฺห สมฺปชาโน ปฏิสฺสโต. “สพฺเพ สงฺขารา ทุกฺขา”ติ ฯเปฯ. “สพฺเพ ธมฺมา อนตฺตา”ติ ฯเปฯ. “ยํ กิญฺจิ สมุทย\-ธมฺมํ, สพฺพํ ตํ นิโรธธมฺมนฺ”ติ น ปรเนยฺโย}

\prose{\csromanfolio{272}น ปรปฺปตฺติโย น ปรปฺปจฺจโย น ปรปฏิพทฺธคู, ยถาภูตํ\csromansharedfootnote{20031aa27d049afb}{นิราสโย (สีฏฺฐ) ขุ 1. 288 ปิฏฺเฐปิ.} ชานาติ ปสฺสติ อสมฺมูโฬฺห สมฺปชาโน ปฏิสฺสโตติ อุปฺปนฺนญาโณมฺหิ อนญฺญเนยฺโย, เอโก จเร ขคฺควิสาณ\-กปฺโป. เตนาห โส ปจฺเจก\-สมฺพุทฺโธ\csromandash{}}

\setlength{\csromangathapagewidth}{0pt}
\setlength{\csromangathawakwidth}{0pt}
\csromangathameasuregroup{}{ทีฏฺฐีวิสูกานิ อุปาติวตฺโต, \\ ปตฺโต นิยามํ ปฏิลทฺธมคฺโค. \\ อุปฺปนฺนญาโณมฺหิ อนญฺญเนยฺโย, \\ เอโก จเร ขคฺควิสาณกปฺโปติ. \\ \textbf{นิลฺโลลุโป นิกฺกุโห นิปฺปิปาโส,} \\ \textbf{นิมฺมกฺโข นิทฺธนฺตกสาวโมโห.} \\ \textbf{นิราสโส}\textsuperscript{1}\textbf{ สพฺพโลเก ภวิตฺวา,} \\ \textbf{เอโก จเร ขคฺควิสาณ}\textbf{กปฺโป.}}
\csromangathameasurewak{ทีฏฺฐีวิสูกานิ อุปาติวตฺโต, \\ ปตฺโต นิยามํ ปฏิลทฺธมคฺโค. \\ อุปฺปนฺนญาโณมฺหิ อนญฺญเนยฺโย, \\ เอโก จเร ขคฺควิสาณกปฺโปติ. \\ \textbf{นิลฺโลลุโป นิกฺกุโห นิปฺปิปาโส,} \\ \textbf{นิมฺมกฺโข นิทฺธนฺตกสาวโมโห.} \\ \textbf{นิราสโส}\textsuperscript{1}\textbf{ สพฺพโลเก ภวิตฺวา,} \\ \textbf{เอโก จเร ขคฺควิสาณ}\textbf{กปฺโป.}}
\setlength{\csromangathaleftcol}{0pt}
\csromangathagroup{\csromangathastanza{\csromangathawak{ทีฏฺฐีวิสูกานิ อุปาติวตฺโต, \\ ปตฺโต นิยามํ ปฏิลทฺธ\-มคฺโค. \\ อุปฺปนฺนญาโณมฺหิ อนญฺญเนยฺโย, \\ เอโก จเร ขคฺควิสาณ\-กปฺโปติ.}}\csromangathastanza{\csromangathawak{\csromangathaitemline{142}{\textbf{นิลฺโลลุโป นิกฺกุโห นิปฺปิปาโส,}} \\ \csromangathacontline{\textbf{นิมฺมกฺโข นิทฺธนฺตกสาวโมโห.}} \\ \csromangathacontline{\textbf{นิราสโส}\csromansharedfootnote{20031aa27d049afb}{นิราสโย (สีฏฺฐ) ขุ 1. 288 ปิฏฺเฐปิ.}\textbf{ สพฺพโลเก ภวิตฺวา,}} \\ \csromangathacontline{\textbf{เอโก จเร ขคฺควิสาณ}\-\textbf{กปฺโป.}}}}}

\prose{\textbf{นิลฺโลลุโป นิกฺกุโห นิปฺปิปาโส}ติ \textbf{โลลุปฺปํ} วุจฺจติ ตณฺหา. โย ราโค สาราโค ฯเปฯ อภิชฺฌา โลโภ อกุสลมูลํ. สา โลลุปฺปา ตณฺหา ตสฺส ปจฺเจก\-สมฺพุทฺธสฺส ปหีนา อุจฺฉินฺนมูลา ตาลาวตฺถุกตา อนภาวํกตา อายติํ อนุปฺปาทธมฺมา. ตสฺมา ปจฺเจก\-สมฺพุทฺโธ นิลฺโลลุโป.}

\prose{\textbf{นิกฺกุโห}ติ ตีณิ กุหนวตฺถูนิ ปจฺจย\-ปฏิเสวน\-สงฺขาตํ กุหนวตฺถุ อิริยาปถ\-สงฺขาตํ กุหนวตฺถุ สามนฺต\-ชปฺปน\-สงฺขาตํ กุหนวตฺถุ.}

\prose{กตมํ ปจฺจย\-ปฏิเสวน\-สงฺขาตํ กุหนวตฺถุ, อิธ คหปติกา ภิกฺขุํ\csromansharedfootnote{58037c0bc16d8ace}{ภิกฺขู (ก) ขุ 7. 172 ปิฏฺเฐ.} นิมนฺเตนฺติ จีวรปิณฺฑปาตเสนาสน\-คิลานปจฺจย\-เภสชฺช- ปริกฺขาเรหิ, โส ปาปิจฺโฉ อิจฺฉาปกโต อตฺถิโก จีวร\-ปิณฺฑปาต\-เสนาสน\-คิลานปจฺจย\-เภสชฺช\-ปริกฺขารานํ ภิยฺโยกมฺยตํ อุปาทาย จีวรํ ปจฺจกฺขาติ, ปิณฺฑปาตํ ปจฺจกฺขาติ, เสนาสนํ ปจฺจกฺขาติ, คิลานปจฺจย\-เภสชฺช\-ปริกฺขารํ ปจฺจกฺขาติ, โส เอวมาห “กิํ สมณสฺส มหคฺเฆน จีวเรน, เอตํ สารุปฺปํ, ยํ สมโณ สุสานา วา สงฺการกูฏา วา ปาปณิกา วา นนฺตกานิ อุจฺจินิตฺวา สงฺฆาฏิกํ กริตฺวา ธาเรยฺย. กิํ สมณสฺส มหคฺเฆน ปิณฺฑปาเตน, เอตํ สารุปฺปํ ยํ สมโณ อุญฺฉาจริยาย ปิณฺฑิยาโลเปน ชีวิกํ กปฺเปยฺย. กิํ สมณสฺส มหคฺเฆน เสนาสเนน, เอตํ สารุปฺปํ, ยํ สมโณ}

\prose{\csromanfolio{273}รุกฺขมูลิโก วา อสฺส โสสานิโก วา อพฺโภกาสิโก วา, กิํ สมณสฺส มหคฺเฆน คิลานปจฺจย\-เภสชฺช\-ปริกฺขาเรน, เอตํ สารุปฺปํ, ยํ สมโณ ปูติมุตฺเตน วา หริตกี\-ขณฺเฑน วา โอสธํ กเรยฺยา”ติ. ตทุปาทาย ลูขํ จีวรํ ธาเรติ, ลูขํ ปิณฺฑปาตํ ปริภุญฺชติ, ลูขํ เสนาสนํ ปฏิเสวติ, ลูขํ คิลานปจฺจย\-เภสชฺช\-ปริกฺขารํ ปฏิเสวติ. ตเมนํ คหปติกา เอวํ ชานนฺติ “อยํ สมโณ อปฺปิจฺโฉ สนฺตุฏฺโฐ ปวิวิตฺโต อสํสฏฺโฐ อารทฺธวีริโย ธุตวาโท”ติ ภิยฺโย ภิยฺโย นิมนฺเตนฺติ จีวร\-ปิณฺฑปาต\-เสนาสนคิลาน ปจฺจย เภสชฺช ปริกฺขาเรหิ, โส เอวมาห “ติณฺณํ สมฺมุขีภาวา สทฺโธ กุลปุตฺโต พหุํ ปุญฺญํ ปสวติ, สทฺธาย สมฺมุขีภาวา สทฺโธ กุลปุตฺโต พหุํ ปุญฺญํ ปสวติ, เทยฺยธมฺมสฺส สมฺมุขีภาวา สทฺโธ กุลปุตฺโต พหุํ ปุญฺญํ ปสวติ, ทกฺขิเณยฺยานํ สมฺมุขีภาวา สทฺโธ กุลปุตฺโต พหุํ ปุญฺญํ ปสวติ. ตุมฺหา\-กญฺเจ\-วายํ สทฺธา อตฺถิ, เทยฺยธมฺโม จ สํวิชฺชติ, อหญฺจ ปฏิคฺคาหโก. สเจหํ น ปฏิคฺคเหสฺสามิ, เอวํ ตุเมฺห ปุญฺเญน ปริพาหิรา ภวิสฺสถ\csromansharedfootnote{b431f803fdb06ea5}{ภวิสฺสนฺติ (ขุ 7. 173 ปิฏฺเฐ)}, น มยฺหํ อิมินา อตฺโถ, อปิ จ ตุมฺหากํเยว อนุกมฺปาย ปฏิคฺคณฺหามี”ติ ตทุปาทาย พหุมฺปิ จีวรํ ปฏิคฺคณฺหาติ, พหุมฺปิ ปิณฺฑปาตํ ปฏิคฺคณฺหาติ, พหุมฺปิ เสนาสนํ ปฏิคฺคณฺหาติ, พหุมฺปิ คิลานปจฺจย\-เภสชฺช\-ปริกฺขารํ ปฏิคฺคณฺหาติ. ยา เอวรูปา ภากุฏิกา ภากุฏิยํ กุหนา กุหายนา กุหิตตฺตํ, อิทํ ปจฺจย\-ปฏิเสวน\-สงฺขาตํ กุหนวตฺถุ.}

\prose{กตมํ อิริยาปถ\-สงฺขาตํ กุหนวตฺถุ, อิเธกจฺโจ ปาปิจฺโฉ อิจฺฉาปกโต สมฺภาวนา\-ธิปฺปาโย “เอวํ มํ ชโน สมฺภาเวสฺสตี”ติ คมนํ สณฺฐเปติ, ฐานํ สณฺฐเปติ, นิสชฺชํ\csromansharedfootnote{68f80a24d10d9c32}{นิสชฺชนํ (ก)} สณฺฐเปติ, สยนํ สณฺฐเปติ. ปณิธาย คจฺฉติ, ปณิธาย ติฏฺฐติ, ปณิธาย นิสีทติ, ปณิธาย เสยฺยํ กปฺเปติ. สมาหิโต วิย คจฺฉติ, สมาหิโต วิย ติฏฺฐติ, สมาหิโต วิย นิสีทติ, สมาหิโต วิย เสยฺยํ กปฺเปติ. อาปาถกชฺฌายีว\csromansharedfootnote{04eb128eab667b8d}{อาปาตกฌายี จ (ก) 4. กุหายิตตฺตํ (สฺยา, ก), วิสุทฺธิมคฺคฏฺฐกถา โอโลเกตพฺพา.} โหติ. ยา เอวรูปา อิริยาปถสฺส อาฐปนา ฐปนา สณฺฐปนา ภากุฏิกา ภากุฏิยํ กุหนา กุหายนา กุหิตตฺตํ\csromansharedfootnote{9e6132b7c3808bb9}{กุหายิตตฺตํ (สฺยา, ก), วิสุทฺธิมคฺคฏฺฐกถา โอโลเกตพฺพา.}. อิทํ อิริยาปถ\-สงฺขาตํ กุหนวตฺถุ.}

\prose{\csromanfolio{274}กตมํ สามนฺต\-ชปฺปน\-สงฺขาตํ กุหนวตฺถุ, อิเธกจฺโจ ปาปิจฺโฉ อิจฺฉาปกโต สมฺภาวนา\-ธิปฺปาโย “เอวํ มํ ชโน สมฺภาเวสฺสตี”ติ อริยธมฺเม สนฺนิสฺสิต\-วาจํ ภาสติ, “โย เอวรูปํ จีวรํ ธาเรติ, โส สมโณ มเหสกฺโข”ติ ภณติ, “โย เอวรูปํ ปตฺตํ ธาเรติ, โลหถาลกํ ธาเรติ, ธมฺมกรณํ ธาเรติ. ปริสฺสาวนํ ธาเรติ, กุญฺจิกํ ธาเรติ, อุปาหนํ ธาเรติ, กายพนฺธนํ ธาเรติ, อาโยคํ\csromansharedfootnote{64f8de3612a97593}{อาโยคพนฺธนํ (สฺยา, ก) ปสฺส ขุ 7. 173 ปิฏฺเฐ.} ธาเรติ, โส สมโณ มเหสกฺโข”ติ ภณติ. “ยสฺส เอวรูโป อุปชฺฌาโย, โส สมโณ มเหสกฺโข”ติ ภณติ. “ยสฺส เอวรูโป อาจริโย. เอวรูปา สมานุ\-ปชฺฌายกา. สมานาจริยกา. มิตฺตา. สนฺทิฏฺฐา. สมฺภตฺตา. สหายา, โส สมโณ มเหสกฺโข”ติ ภณติ. “โย เอวรูเป วิหาเร วสติ. อฑฺฒโยเค วสติ. ปาสาเท วสติ. หมฺมิเย วสติ. คุหายํ วสติ. เลเณ วสติ. กุฏิยํ วสติ. กูฏาคาเร วสติ. อฏฺเฏ วสติ. มาเฬ วสติ. อุทฺทณฺเฑ\csromansharedfootnote{204e26d2dcd2095c}{อุฏฺฏณฺเฑ (ก)} วสติ. อุปฏฺฐาน\-สาลายํ วสติ. มณฺฑเป วสติ. รุกฺขมูเล วสติ. โส สมโณ มเหสกฺโข”ติ ภณติ.}

\prose{อถ วา โกรชิก\-โกรชิโก ภากุฏิก\-ภากุฏิโก กุหกกุหโก ลปกลปโก มุขสมฺภาวิโก “อยํ สมโณ อิมาสํ เอวรูปานํ สนฺตานํ วิหาร\-สมาปตฺตีนํ ลาภี”ติ ตาทิสํ คมฺภีรํ คูฬฺหํ นิปุณํ ปฏิจฺฉนฺนํ โลกุตฺตรํ สุญฺญตา\-ปฏิสญฺญุตฺตํ กถํ กเถติ. ยา เอวรูปา ภากุฏิกา ภากุฏิยํ กุหนา กุหายนา กุหิตตฺตํ, อิทํ สามนฺต\-ชปฺปน\-สงฺขาตํ กุหนวตฺถุ. ตสฺส ปจฺเจก\-สมฺพุทฺธสฺส อิมานิ ตีณิ กุหนวตฺถุนิ ปหีนานิ สมุจฺฉินฺนานิ วูปสนฺตานิ ปฏิปฺปสฺสทฺธานิ อภพฺพุ\-ปฺปตฺติกานิ ญาณคฺคินา ทฑฺฒานิ. ตสฺมา โส ปจฺเจก\-สมฺพุทฺโธ นิกฺกุโห.}

\prose{\textbf{นิปฺปิปาโส}ติ \textbf{ปิปาสา} วุจฺจติ ตณฺหา. โย ราโค สาราโค ฯเปฯ อภิชฺฌา โลโภ อกุสลมูลํ. สา ปิปาสา ตณฺหา ตสฺส ปจฺเจก\-สมฺพุทฺธสฺส ปหีนา อุจฺฉินฺนมูลา ตาลาวตฺถุกตา อนภาวํกตา อายติํ อนุปฺปาทธมฺมา. ตสฺมา ปจฺเจก\-สมฺพุทฺโธ นิปฺปิปาโสติ นิลฺโลลุโป นิกฺกุโห นิปฺปิปาโส.}

\prose{\csromanfolio{275}\textbf{นิมฺมกฺโข นิทฺธนฺตกสาวโมโห}ติ \textbf{มกฺโข}ติ โย มกฺโข มกฺขายนา\csromansharedfootnote{d37b3c907f36219c}{มกฺขิยนา (ก) ปสฺส อภิ 2. 371 ปิฏฺเฐ.} มกฺขายิตตฺตํ นิฏฺฐุริยํ นิฏฺฐุริย\-กมฺมํ. \textbf{กสาโว}ติ ราโค กสาโว, โทโส กสาโว, โมโห กสาโว, โกโธ. อุปนาโห. มกฺโข. ปฬาโส ฯเปฯ สพฺพา\-กุสลา\-ภิสงฺขารา กสาวา. \textbf{โมโห}ติ ทุกฺเข อญฺญาณํ, ทุกฺขสมุทเย อญฺญาณํ, ทุกฺขนิโรเธ อญฺญาณํ, ทุกฺขนิโรธ\-คามินิยา ปฏิปทาย อญฺญาณํ, ปุพฺพนฺเต อญฺญาณํ, อปรนฺเต อญฺญาณํ, ปุพฺพนฺตา\-ปรนฺเต อญฺญาณํ. อิทปฺปจฺจยตาปฏิจฺจสมุปฺปนฺเนยุ ธมฺเมสุ อญฺญาณํ. ยํ เอวรูปํ อญฺญาณํ อทสฺสนํ อนภิสมโย อนนุโพโธ อปฺปฏิเวโธ อสํคาหนา อปริโยคาหนา อสมเปกฺขนา อปจฺจเวกฺขณา อปจฺจกฺขกมฺมํ ทุมฺเมชฺฌํ พาลฺยํ อสมฺปชญฺญํ โมโห ปโมโห สมฺโมโห อวิชฺชา อวิชฺโชโฆ อวิชฺชาโยโค อวิชฺชานุสโย อวิชฺชา\-ปริยุฏฺฐานํ อวิชฺชาลงฺคี โมโห อกุสลมูลํ. ตสฺส ปจฺเจก\-สมฺพุทฺธสฺส มกฺโข จ กสาโว จ โมโห จ วนฺตา สํวนฺตา นิทฺธนฺตา ปหีนา สมุจฺฉินฺนา วูปสนฺตา ปฏิปฺปสฺสทฺธา อภพฺพุ\-ปฺปตฺติกา ญาณคฺคินา ทฑฺฒาติ โส ปจฺเจก\-สมฺพุทฺโธ นิมฺมกฺโข นิทฺธนฺตกสาวโมโห.}

\prose{\textbf{นิราสโส สพฺพโลเก ภวิตฺวา}ติ \textbf{อาสา} วุจฺจติ ตณฺหา. โย ราโค สาราโค ฯเปฯ อภิชฺฌา โลโภ อกุสลมูลํ. \textbf{สพฺพโลเก}ติ สพฺพอปายโลเก สพฺพ\-มนุสฺสโลเก สพฺพเทวโลเก สพฺพ\-ขนฺธโลเก สพฺพธาตุโลเก สพฺพ- อายตนโลเก. \textbf{นิราสโส สพฺพโลเก ภวิตฺวา}ติ สพฺพโลเก นิราสโส ภวิตฺวา นิตฺตโณฺห ภวิตฺวา นิปฺปิปาโส ภวิตฺวาติ นิราสโส สพฺพโลเก ภวิตฺวา, เอโก จเร ขคฺควิสาณ\-กปฺโป. เตนาห โส ปจฺเจก\-สมฺพุทฺโธ\csromandash{}}

\prose{นิลฺโลลุโป นิกฺกุโห นิปฺปิปาโส,}

\prose{นิมฺมกฺโข นิทฺธนฺตกสาวโมโห.}

\setlength{\csromangathapagewidth}{0pt}
\setlength{\csromangathawakwidth}{0pt}
\csromangathameasuregroup{}{นิราสโส สพฺพโลเก ภวิตฺวา, \\ เอโก จเร ขคฺควิสาณกปฺโปติ. \\ \textbf{ปาปํ สหายํ ปริวชฺชเยถ,} \\ \textbf{อนตฺถทสฺสิํ วิสเม นิวิฏฺฐํ.} \\ \textbf{สยํ น เสเว ปสุตํ ปมตฺตํ,} \\ \textbf{เอโก จเร ขคฺควิสาณ}\textbf{กปฺโป.}}
\csromangathameasurewak{นิราสโส สพฺพโลเก ภวิตฺวา, \\ เอโก จเร ขคฺควิสาณกปฺโปติ. \\ \textbf{ปาปํ สหายํ ปริวชฺชเยถ,} \\ \textbf{อนตฺถทสฺสิํ วิสเม นิวิฏฺฐํ.} \\ \textbf{สยํ น เสเว ปสุตํ ปมตฺตํ,} \\ \textbf{เอโก จเร ขคฺควิสาณ}\textbf{กปฺโป.}}
\setlength{\csromangathaleftcol}{0pt}
\csromangathagroup{\csromangathastanza{\csromangathawak{นิราสโส สพฺพโลเก ภวิตฺวา, \\ เอโก จเร ขคฺควิสาณ\-กปฺโปติ.}}\csromangathastanza{\csromangathawak{\csromangathaitemline{143}{\textbf{ปาปํ สหายํ ปริวชฺชเยถ,}} \\ \csromangathacontline{\textbf{อนตฺถทสฺสิํ วิสเม นิวิฏฺฐํ.}} \\ \csromangathacontline{\textbf{สยํ น เสเว ปสุตํ ปมตฺตํ,}} \\ \csromangathacontline{\textbf{เอโก จเร ขคฺควิสาณ}\-\textbf{กปฺโป.}}}}}

\prose{\csromanfolio{276}\textbf{ปาปํ สหายํ ปริวชฺชเยถา}ติ ปาปสหาโย วุจฺจติ โย โส สหาโย ทสวตฺถุกาย มิจฺฉาทิฏฺฐิยา สมนฺนาคโต “นตฺถิ ทินฺนํ, นตฺถิ ยิฏฺฐํ, นตฺถิ หุตํ, นตฺถิ สุกต\-ทุกฺกฏานํ กมฺมานํ ผลํ วิปาโก, นตฺถิ อยํ โลโก, นตฺถิ ปโร โลโก, นตฺถิ มาตา, นตฺถิ ปิตา, นตฺถิ สตฺตา โอปปาติกา, นตฺถิ โลเก สมณพฺราหฺมณา สมฺมคฺคตา สมฺมาปฏิปนฺนา. เย อิมญฺจ โลกํ ปรญฺจ โลกํ สยํ อภิญฺญา สจฺฉิกตฺวา ปเวเทนฺตี”ติ. อยํ ปาปสหาโย, ปาปํ สหยํ ปริวชฺชเยถาติ ปาปํ สหายํ วชฺเชยฺย ปริวชฺเชยฺยาติ ปาปํ สหายํ ปริวชฺชเยถ.}

\prose{\textbf{อนตฺถทสฺสิํ วิสเม นิวิฏฺฐนฺ}ติ อนตฺถทสฺสี วุจฺจติ โย โส สหาโย ทสวตฺถุกาย มิจฺฉาทิฏฺฐิยา สมนฺนาคโต “นตฺถิ ทินฺนํ, นตฺถิ ยิฏฺฐํ ฯเปฯ. เย อิมญฺจ โลกํ ปรญฺจ โลกํ สยํ อภิญฺญา สจฺฉิกตฺวา ปเวเทนฺตี”ติ. \textbf{วิสเม นิวิฏฺฐนฺ}ติ วิสเม กายกมฺเม นิวิฏฺฐํ. วิสเม วจีกมฺเม นิวิฏฺฐํ. วิสเม มโนกมฺเม นิวิฏฺฐํ. วิสเม ปาณาติปาเต นิวิฏฺฐํ. วิสเม อทินฺนาทาเน นิวิฏฺฐํ. วิสเม กาเมสุ\-มิจฺฉาจาเร นิวิฏฺฐํ. วิสเม มุสาวาเท นิวิฏฺฐํ. วิสมาย ปิสุณาย วาจาย\csromansharedfootnote{3138c50f07bb79d4}{ปิสุณวาจาย (ก)} นิวิฏฺฐํ. วิสมาย ผรุสาย วาจาย\csromansharedfootnote{78583c96faf971a7}{ผรุสวาจาย (ก)} นิวิฏฺฐํ. วิสเม สมฺผปฺปลาเป นิวิฏฺฐํ. วิสมาย อภิชฺฌาย นิวิฏฺฐํ. วิสเม พฺยาปาเท นิวิฏฺฐํ. วิสมาย มิจฺฉาทิฏฺฐิยา นิวิฏฺฐํ. วิสเมสุ สงฺขาเรสุ นิวิฏฺฐํ. วิสเมสุ ปญฺจสุ กามคุเณสุ นิวิฏฺฐํ. วิสเมสุ ปญฺจสุ นีวรเณสุ นิวิฏฺฐํ วินิวิฏฺฐํ สตฺตํ อลฺลีนํ อุปคตํ อชฺโฌสิตํ อธิมุตฺตนฺติ อนตฺถทสฺสิํ วิสเม นิวิฏฺฐํ.}

\prose{\textbf{สยํ น เสเว ปสุตํ ปมตฺตนฺ}ติ \textbf{ปสุตนฺ}ติ โยปิ กาเม เอสติ คเวสติ ปริเยสติ ตจฺจริโต ตพฺพหุโล ตคฺครุโก ตนฺนินฺโน ตปฺโปโณ ตปฺปพฺภาโร ตทธิมุตฺโต ตทธิปเตยฺโย, โสปิ กามปฺปสุโต. โยปิ ตณฺหาวเสน รูเป ปริเยสติ, สทฺเท. คนฺเธ. รเส. โผฏฺฐพฺเพ ปริเยสติ ตจฺจริโต ตพฺพหุโล ตคฺครุโก ตนฺนินฺโน ตปฺโปโณ ตปฺปพฺภาโร ตทธิมุตฺโต ตทธิปเตยฺโย, โสปิ กามปฺปสุโต. โยปิ ตณฺหาวเสน รูเป ปฏิลภติ, สทฺเท. คนฺเธ. รเส. โผฏฺฐพฺเพ ปฏิลภติ ตจฺจริโต ตพฺพหุโล ตคฺครุโก ตนฺนินฺโน ตปฺโปโณ ตปฺปพฺภาโร ตทธิมุตฺโต ตทธิปเตยฺโย, โสปิ กามปฺปสุโต. โยปิ ตณฺหาวเสน รูเป ปริภุญฺชติ, สทฺเท. คนฺเธ. รเส.}

\prose{\csromanfolio{277}โผฏฺฐพฺเพ ปริภุญฺชติ ตจฺจริโต ตพฺพหุโล ตคฺครุโก ตนฺนินฺโน ตปฺโปโณ ตปฺปพฺภาโร ตทธิมุตฺโต ตทธิปเตยฺโย, โสปิ กามปฺปสุโต. ยถา กลหการโก กลหปฺปสุโต กมฺมการโก กมฺมปฺปสุโต โคจเร จรนฺโต โคจรปฺปสุโต ฌายี ฌานปฺปสุโต, เอวเมว โย กาเม เอสติ คเวสติ ปริเยสติ ตจฺจริโต ตพฺพหุโล ตคฺครุโก ตนฺนินฺโน ตปฺโปโณ ตปฺปพฺภาโร ตทธิมุตฺโต ตทธิปเตยฺโย, โสปิ กามปฺปสุโต. โยปิ ตณฺหาวเสน รูเป ปริเยสติ ฯเปฯ. โยปิ ตณฺหาวเสน รูเป ปฏิลภติ ฯเปฯ. โยปิ ตณฺหาวเสน รูเป ปริภุญฺชติ, สทฺเท. คนฺเธ. รเส. โผฏฺฐพฺเพ ปริภุญฺชติ ตจฺจริโต ตพฺพหุโล ตคฺครุโก ตนฺนินฺโน ตปฺโปโณ ตปฺปพฺภาโร ตทธิมุตฺโต ตทธิปเตยฺโย, โสปิ กามปฺปสุโต. \textbf{ปมตฺตนฺ}ติ ปมาโท วตฺตพฺโพ กายทุจฺจริเต วา วจีทุจฺจริเต วา มโนทุจฺจริเต วา ปญฺจสุ กามคุเณสุ วา จิตฺตสฺส โวสคฺโค โวสคฺคา\-นุปฺปทานํ กุสลานํ ธมฺมานํ ภาวนาย อสกฺก\-จฺจ\-กิริ\-ยตา อสาตจฺจกิริยตา อนฏฺฐิต\-กิริยตา โอลีนวุตฺติตา นิกฺขิตฺต\-จฺฉนฺทตา นิกฺขิตฺตธุรตา อนาเสวนา อภาวนา อพหุลีกมฺมํ อนธิฏฺฐานํ อนนุโยโค, โย เอวรูโป มาโท ปมชฺชนา ปมชฺชิตตฺตํ, อยํ วุจฺจติ ปมาโท.}

\prose{\textbf{สยํ น เสเว ปสุตํ ปมตฺตนฺ}ติ ปสุตํ น เสเวยฺย ปมตฺตญฺจ สยํ น เสเวยฺย สามํ น เสเวยฺย น นิเสเวยฺย น สํเสเวยฺย น ปริสํเสเวยฺย น อาจเรยฺย น สมาจเรยฺย น สมาทาย วตฺเตยฺยาติ สยํ น เสเว ปสุตํ ปมตฺตํ, เอโก จเร ขคฺควิสาณ\-กปฺโป. เตนาห โส ปจฺเจก\-สมฺพุทฺโธ\csromandash{}}

\setlength{\csromangathapagewidth}{0pt}
\setlength{\csromangathawakwidth}{0pt}
\csromangathameasuregroup{}{ปาปํ สหายํ ปริวชฺชเยถ, \\ อนตฺถทสฺสิํ วิสเม นิวิฏฺฐํ.}
\csromangathameasurewak{ปาปํ สหายํ ปริวชฺชเยถ, \\ อนตฺถทสฺสิํ วิสเม นิวิฏฺฐํ.}
\setlength{\csromangathaleftcol}{0pt}
\csromangathagroup{\csromangathastanza{\csromangathawak{ปาปํ สหายํ ปริวชฺชเยถ, \\ อนตฺถทสฺสิํ วิสเม นิวิฏฺฐํ.}}}

\prose{สยํ น เสเว ปสุตํ ปมตฺตํ,}

\prose{เอโก จเร ขคฺควิสาณ\-กปฺโปติ.}

\setlength{\csromangathapagewidth}{0pt}
\setlength{\csromangathawakwidth}{0pt}
\csromangathameasuregroup{}{\textbf{พหุสฺสุตํ ธมฺมธรํ ภเชถ,} \\ \textbf{มิตฺตํ อุฬารํ ปฏิภานวนฺตํ.} \\ \textbf{อญฺญาย อตฺถานิ วิเนยฺย กงฺขํ,} \\ \textbf{เอโก จเร ขคฺควิสาณ}\textbf{กปฺโป.}}
\csromangathameasurewak{\textbf{พหุสฺสุตํ ธมฺมธรํ ภเชถ,} \\ \textbf{มิตฺตํ อุฬารํ ปฏิภานวนฺตํ.} \\ \textbf{อญฺญาย อตฺถานิ วิเนยฺย กงฺขํ,} \\ \textbf{เอโก จเร ขคฺควิสาณ}\textbf{กปฺโป.}}
\setlength{\csromangathaleftcol}{0pt}
\csromangathagroup{\csromangathastanza{\csromangathawak{\csromangathaitemline{144}{\textbf{พหุสฺสุตํ ธมฺมธรํ ภเชถ,}} \\ \csromangathacontline{\textbf{มิตฺตํ อุฬารํ ปฏิภานวนฺตํ.}} \\ \csromangathacontline{\textbf{อญฺญาย อตฺถานิ วิเนยฺย กงฺขํ,}} \\ \csromangathacontline{\textbf{เอโก จเร ขคฺควิสาณ}\-\textbf{กปฺโป.}}}}}

\prose{\csromanfolio{278}\textbf{พหุสฺสุตํ ธมฺมธรํ ภเชถา}ติ พหุสฺสุโต โหติ มิตฺโต สุตธโร สุตสนฺนิจโย, เย เต ธมฺมา อาทิกลฺยาณา มชฺเฌกลฺยาณา ปริโยสาน\-กลฺยาณา สาตฺถํ สพฺยญฺชนํ\csromansharedfootnote{b554163f6b24f6b0}{สตฺถา สพฺยญฺชนา (สฺยา) ปสฺส ม 3. 62 ปิฏฺเฐ.} เกวล\-ปริปุณฺณํ ปริสุทฺธํ พฺรหฺมจริยํ อภิวทนฺติ, ตถารูปาสฺส ธมฺมา พหุสฺสุตา โหนฺติ ธาตา\csromansharedfootnote{a4c017d0895ce78b}{ธตา (สฺยา)} วจสา ปริจิตา มนสา\-นุเปกฺขิตา ทิฏฺฐิยา สุปฺปฏิวิทฺธา. \textbf{ธมฺมธรนฺ}ติ ธมฺมํ ธาเรนฺตํ\csromansharedfootnote{dee24bf1887bd0ff}{ธาเรติ (ก) 4. ปริยาปุฏํ (สฺยา, ก) ปสฺส ที 3. 164ปิฏฺเฐ. 5. ฉฬภิญฺญาโย (สฺยา)} สุตฺตํ เคยฺยํ เวยฺยากรณํ คาถํ อุทานํ อิติวุตฺตกํ ชาตกํ อพฺภุตธมฺมํ เวทลฺลํ. \textbf{พหุสฺสุตํ ธมฺมธรํภเชถา}ติ พหุสฺสุตญฺจ ธมฺมธรญฺจ มิตฺตํ ภเชยฺย สํภเชยฺย เสเวยฺย นิเสเวยฺย สํเสเวยฺย ปฏิเสเวยฺยาติ พหุสฺสุตํ ธมฺมธรํ ภเชถ.}

\prose{\textbf{มิตฺตํ อุฬารํ ปฏิภานวนฺตนฺ}ติ อุฬาโร โหติ มิตฺโต สีเลน สมาธินา ปญฺญาย วิมุตฺติยา วิมุตฺติ\-ญาณ\-ทสฺสเนน. \textbf{ปฏิภานวนฺตนฺ}ติ ตโย ปฏิภานวนฺโต ปริยตฺติ\-ปฏิภานวา ปริปุจฺฉา\-ปฏิภานวา อธิคม\-ปฏิภานวา. กตโม ปริยตฺติ\-ปฏิภานวา, อิเธกจฺจสฺส พุทฺธวจนํ ปริยาปุตํ\csromansharedfootnote{a71c20bbdb93e147}{ปริยาปุฏํ (สฺยา, ก) ปสฺส ที 3. 164ปิฏฺเฐ.} โหติ สุตฺตํ เคยฺยํ เวยฺยากรณํ คาถา อุทานํ อิติวุตฺตกํ ชาตกํ อพฺภุตธมฺมํ เวทลฺลํ, ตสฺส ปริยตฺติํ นิสฺสาย ปฏิภาติ, อยํ ปริยตฺติ\-ปฏิภานวา.}

\prose{กตโม ปริปุจฺฉา\-ปฏิภานวา, อิเธกจฺโจ ปริปุจฺฉิโตปิ โหติ อตฺเถ จ ญาเย จ ลกฺขเณ จ การเณ จ ฐานาฏฺฐาเน จ, ตสฺส ปริปุจฺฉํ นิสฺสาย ปฏิภาติ, อยํ ปริปุจฺฉา\-ปฏิภานวา.}

\prose{กตโม อธิคม\-ปฏิภานวา, อิเธกจฺจสฺส อธิคตา โหนฺติ จตฺตาโร สติปฏฺฐานา จตฺตาโร สมฺมปฺปธานา จตฺตาโร อิทฺธิปาทา ปญฺจินฺทฺริยานิ ปญฺจ พลานิ สตฺต โพชฺฌงฺคา อริโย อฏฺฐงฺคิโก มคฺโค จตฺตาโร อริยมคฺคา จตฺตาริ สามญฺญผลานิ จตสฺโส ปฏิสมฺภิ ทาโย ฉ อภิญฺญาโย\csromansharedfootnote{78abf8addb2be13c}{ฉฬภิญฺญาโย (สฺยา)}. ตสฺส อตฺโถ ญาโต, ธมฺโม ญาโต, นิรุตฺติ ญาตา. อตฺเถ ญาเต อตฺโถ ปฏิภาติ, ธมฺเม ญาเต ธมฺโม ปฏิภาติ, นิรุตฺติยา ญาตาย นิรุตฺติ ปฏิภาติ, อิเมสุ ตีสุ ญาณํ ปฏิภาน\-ปฏิสมฺภิทา. โส ปจฺเจก\-สมฺพุทฺโธ อิมาย ปฏิภาน\-ปฏิสมฺภิทาย อุเปโต สมุเปโต อุปาคโต สมุปาคโต อุปปนฺโน สมุปปนฺโน สมนฺนาคโต, ตสฺมา ปจฺเจก\-สมฺพุทฺโธ}

\prose{\csromanfolio{279}ปฏิภานวา. ยสฺส ปริยตฺติ นตฺถิ ปริปุจฺฉา นตฺถิ อธิคโม นตฺถิ, กิํ ตสฺส ปฏิภายิสฺสตีติ มิตฺตํ อุฬารํ ปฏิภานวนฺตํ.}

\prose{\textbf{อญฺญาย อตฺถานิ วิเนยฺย กงฺขนฺ}ติ อตฺตตฺถํ อญฺญาย ปรตฺถํ อญฺญาย อุภยตฺถํ อญฺญาย ทิฏฺฐธมฺมิกตฺถํ อญฺญาย สมฺปรายิกตฺถํ อญฺญาย ปรมตฺถํ อญฺญาย อภิญฺญาย ชานิตฺวา ตุลยิตฺวา ตีรยิตฺวา วิภาวยิตฺวา วิภูตํ กตฺวา กงฺขํ วิเนยฺย ปฏิวิเนยฺย ปชเหยฺย วิโนเทยฺย พฺยนฺตีกเรยฺย อนภาวํ คเมยฺยาติ อญฺญาย อตฺถานิ วิเนยฺย กงฺขํ, เอโก จเร ขคฺควิสาณ\-กปฺโป. เตนาห โส ปจฺเจก\-สมฺพุทฺโธ\csromandash{}}

\setlength{\csromangathapagewidth}{0pt}
\setlength{\csromangathawakwidth}{0pt}
\csromangathameasuregroup{}{พหุสฺสุตํ ธมฺมธรํ ภเชถ, \\ มิตฺตํ อุฬารํ ปฏิภานวนฺตํ.}
\csromangathameasurewak{พหุสฺสุตํ ธมฺมธรํ ภเชถ, \\ มิตฺตํ อุฬารํ ปฏิภานวนฺตํ.}
\setlength{\csromangathaleftcol}{0pt}
\csromangathagroup{\csromangathastanza{\csromangathawak{พหุสฺสุตํ ธมฺมธรํ ภเชถ, \\ มิตฺตํ อุฬารํ ปฏิภานวนฺตํ.}}}

\prose{อญฺญาย อตฺถานิ วิเนยฺย กงฺขํ,}

\prose{เอโก จเร ขคฺควิสาณ\-กปฺโปติ.}

\setlength{\csromangathapagewidth}{0pt}
\setlength{\csromangathawakwidth}{0pt}
\csromangathameasuregroup{}{ขิฑฺฑํ รติํ\textsuperscript{1} \textbf{กามสุขญฺจ โลเก,} \\ \textbf{อนลงฺกริตฺวา อนเปกฺขมาโน.} \\ \textbf{วิภูสฏฺฐานา}\textsuperscript{1}\textbf{ วิรโต สจฺจวาที,} \\ \textbf{เอโก จเร ขคฺควิสาณ}\textbf{กปฺโป.}}
\csromangathameasurewak{ขิฑฺฑํ รติํ\textsuperscript{1} \textbf{กามสุขญฺจ โลเก,} \\ \textbf{อนลงฺกริตฺวา อนเปกฺขมาโน.} \\ \textbf{วิภูสฏฺฐานา}\textsuperscript{1}\textbf{ วิรโต สจฺจวาที,} \\ \textbf{เอโก จเร ขคฺควิสาณ}\textbf{กปฺโป.}}
\setlength{\csromangathaleftcol}{0pt}
\csromangathagroup{\csromangathastanza{\csromangathawak{\csromangathaitemline{145}{ขิฑฺฑํ รติํ\csromansharedfootnote{49c9bb5a23d3a799}{รตี (สฺยา)} \textbf{กามสุขญฺจ โลเก,}} \\ \csromangathacontline{\textbf{อนลงฺกริตฺวา อนเปกฺขมาโน.}} \\ \csromangathacontline{\textbf{วิภูสฏฺฐานา}\csromansharedfootnote{b4284648ba5219d0}{วิภูสนฏฺฐานา (สฺยา, ก), วิภูสณฏฺฐานา (สีฏฺฐ)}\textbf{ วิรโต สจฺจวาที,}} \\ \csromangathacontline{\textbf{เอโก จเร ขคฺควิสาณ}\-\textbf{กปฺโป.}}}}}

\prose{\textbf{ขิฑฺฑํรติํ กามสุขญฺจ โลเก}ติ ขิฑฺฑาติ เทฺว \textbf{ขิฑฺฑา}, กายิกา ขิฑฺฑา จ วาจสิกา ขิฑฺฑา จ ฯเปฯ อยํ กายิกา ขิฑฺฑา ฯเปฯ อยํ วาจสิกา ขิฑฺฑา. \textbf{รตี}ติ อนุกฺกณฺฐิตาธิวจนเมตํ รตีติ. \textbf{กามสุขนฺ}ติ วุตฺตํ เหตํ ภควตา\csromansharedfootnote{489bb279798c0d6c}{ปสฺส ม 2. 233 ปิฏฺเฐ.}\csromandash{}ปญฺจิเม ภิกฺขเว กามคุณา. กตเม ปญฺจ, จกฺขุวิญฺเญยฺยา รูปา อิฏฺฐา กนฺตา มนาปา ปิยรูปา กามูปสํหิตา รชนียา, โสตวิญฺเญยฺยา สทฺทา ฯเปฯ ฆานวิญฺเญยฺยา คนฺธา. ชิวฺหาวิญฺเญยฺยา รสา. กายวิญฺเญยฺยา โผฏฺฐพฺพา อิฏฺฐา กนฺตา มนาปา ปิยรูปา กามูปสํหิตา รชนียา, อิเม โข ภิกฺขเว ปญฺจ กามคุณา. ยํ โข ภิกฺขเว อิเม ปญฺจ กามคุเณ ปฏิจฺจ อุปฺปชฺชติ สุขํ โสมนสฺสํ, อิทํ วุจฺจติ กามสุขํ. โลเกติ มนุสฺสโลเกติ ขิฑฺฑํ รติํ กามสุขญฺจ โลเก.}

\prose{\csromanfolio{280}\textbf{อนลงฺกริตฺวา อนเปกฺขมาโน}ติ ขิฑฺฑญฺจ รติญฺจ กามสุขญฺจ โลเก อนลงฺกริตฺวา อนเปกฺโข หุตฺวา ปชหิตฺวา วิโนเทตฺวา พฺยนฺตีกริตฺวา อนภาวํ คเมตฺวาติ อนลงฺกริตฺวา อนเปกฺขมาโน.}

\prose{\textbf{วิภูสฏฺฐานา วิรโต สจฺจวาที}ติ วิภูสาติ เทฺว \textbf{วิภูสา}, อตฺถิ อคาริกวิภูสา\csromansharedfootnote{7bc3e1bcf46a5df4}{อคาริกสฺส วิภูสา (ก) เอวมุปริปิ.}, อตฺถิ อนา\-คาริก\-วิภูสา. กตมา อคาริกวิภูสา, เกสา จ มสฺสู จ มาลาคนฺธญฺจ วิเลปนญฺจ อาภรณญฺจ ปิลนฺธนญฺจ วตฺถญฺจ ปารุปนญฺจ\csromansharedfootnote{4f09044f3b9f9b85}{ปสาธนญฺจ (สฺยา), สยนาสนญฺจ (ก) 254 ปิฏฺเฐ นตฺถิ ปาฐนานตฺตํ.} เวฐนญฺจ อุจฺฉาทนํ ปริมทฺทนํ นฺหาปนํ สมฺพาหนํ อาทาสํ อญฺชนํ มาลา\-คนฺธ\-วิเลปนํ มุขจุณฺณํ มุขเลปนํ หตฺถพนฺธํ สิขาพนฺธํ ทณฺฑํ นาฬิกํ ขคฺคํ ฉตฺตํ จิตฺรุปาหนํ อุณฺหีสํ มณิํ วาฬพีชนิํ โอทาตานิ วตฺถานิ ทีฆทสานิ อิติ วา, อยํ อคาริกวิภูสา.}

\prose{กตมา อนา\-คาริก\-วิภูสา, จีวรมณฺฑนา ปตฺตมณฺฑนา เสนาสน\-มณฺฑนา อิมสฺส วา ปูติกายสฺส พาหิรานํ วา ปริกฺขารานํ มณฺฑนา วิภูสนา เกฬนา ปริเกฬนา คทฺธิกตา คทฺธิกตฺตํ\csromansharedfootnote{b9a7fa0587ebf45d}{เคธิกตา เคธิกตฺตํ (สฺยา) ปสฺส อภิ 2. 365 ปิฏฺเฐ.} จปลตา\csromansharedfootnote{79cd816cdf0949ca}{จปลนา (สฺยา)} จาปลฺยํ, อยํ อนา\-คาริก\-วิภูสา.}

\prose{\textbf{สจฺจวาที}ติ โส ปจฺเจก\-สมฺพุทฺโธ สจฺจวาที สจฺจสนฺโธ เถโต ปจฺจยิโก อวิสํวาทโก โลกสฺส วิภูสฏฺฐานา อารโต วิรโต ปฏิวิรโต นิกฺขนฺโต นิสฺสโฏ วิปฺปมุตฺโต วิสญฺญุตฺโต วิมริยาทิกเตน เจตสา วิหรตีติ วิภูสฏฺฐานา วิรโต สจฺจวาที, เอโก จเร ขคฺควิสาณ\-กปฺโป. เตนาห โส ปจฺเจก\-สมฺพุทฺโธ\csromandash{}}

\setlength{\csromangathapagewidth}{0pt}
\setlength{\csromangathawakwidth}{0pt}
\csromangathameasuregroup{}{ขิฑฺฑํ รติํ กามสุขญฺจ โลเก, \\ อนลงฺกริตฺวา อนเปกฺขมาโน.}
\csromangathameasurewak{ขิฑฺฑํ รติํ กามสุขญฺจ โลเก, \\ อนลงฺกริตฺวา อนเปกฺขมาโน.}
\setlength{\csromangathaleftcol}{0pt}
\csromangathagroup{\csromangathastanza{\csromangathawak{ขิฑฺฑํ รติํ กามสุขญฺจ โลเก, \\ อนลงฺกริตฺวา อนเปกฺขมาโน.}}}

\prose{วิภูสฏฺฐานา วิรโต สจฺจวาที,}

\prose{เอโก จเร ขคฺควิสาณ\-กปฺโปติ.}

\setlength{\csromangathapagewidth}{0pt}
\setlength{\csromangathawakwidth}{0pt}
\csromangathameasuregroup{}{\textbf{ปุตฺตญฺจ ทารํ ปิตรญฺจ มาตรํ,} \\ \textbf{ธนานิ ธญฺญานิ จ พนฺธวานิ.} \\ \textbf{หิตฺวาน กามานิ ยโถธิกานิ,} \\ \textbf{เอโก จเร ขคฺควิสาณ}\textbf{กปฺโป.}}
\csromangathameasurewak{\textbf{ปุตฺตญฺจ ทารํ ปิตรญฺจ มาตรํ,} \\ \textbf{ธนานิ ธญฺญานิ จ พนฺธวานิ.} \\ \textbf{หิตฺวาน กามานิ ยโถธิกานิ,} \\ \textbf{เอโก จเร ขคฺควิสาณ}\textbf{กปฺโป.}}
\setlength{\csromangathaleftcol}{0pt}
\csromangathagroup{\csromangathastanza{\csromangathawak{\csromangathaitemline{146}{\textbf{ปุตฺตญฺจ ทารํ ปิตรญฺจ มาตรํ,}} \\ \csromangathacontline{\textbf{ธนานิ ธญฺญานิ จ พนฺธวานิ.}} \\ \csromangathacontline{\textbf{หิตฺวาน กามานิ ยโถธิกานิ,}} \\ \csromangathacontline{\textbf{เอโก จเร ขคฺควิสาณ}\-\textbf{กปฺโป.}}}}}

\prose{\csromanfolio{281}\textbf{ปุตฺตญฺจ ทารํ ปิตรญฺจ มาตรนฺ}ติ ปุตฺตาติ จตฺตาโร \textbf{ปุตฺตา} อตฺรโช ปุตฺโต เขตฺตโช\csromansharedfootnote{31457e795f2181f6}{เขตฺรโช (สฺยา, ก)} ปุตฺโต ทินฺนโก ปุตฺโต อนฺเตวาสิโก ปุตฺโต. \textbf{ทารา} วุจฺจนฺติ ภริยาโย. \textbf{ปิตา}ติ โย โส ชนโก. \textbf{มาตา}ติ ยา สา ชนิกาติ ปุตฺตญฺจ ทารํ ปิตรญฺจ มาตรํ.}

\prose{\textbf{ธนานิ ธญฺญานิ จ พนฺธวานี}ติ ธนานิ วุจฺจนฺติ หิรญฺญํ สุวณฺณํ มุตฺตา มณิ เวฬุริโย สงฺโข สิลา ปวาฬํ รชตํ ชาตรูปํ โลหิตงฺโค\csromansharedfootnote{72cde618a10fc3b8}{โลหิตโก (?)} มสารคลฺลํ. ธญฺญานิ วุจฺจนฺติ ปุพฺพณฺณํ อปรณฺณํ. \textbf{ปุพฺพณฺณํ} นาม สาลิ วีหิ ยโว โคธุโม กงฺคุ วรโก กุทฺรูสโก\csromansharedfootnote{e06a9e1830aede7e}{กุทฺรุสโก (สฺยา)}. \textbf{อปรณฺณํ} นาม สูเปยฺยํ. \textbf{พนฺธวานี}ติ จตฺตาโร พนฺธวา ญาติพนฺธวาปิ พนฺธุ โคตฺต\-พนฺธวาปิ พนฺธุ มิตฺต\-พนฺธวาปิ พนฺธุ สิปฺป\-พนฺธวาปิ พนฺธูติ ธนานิ ธญฺญานิ จ พนฺธวานิ.}

\prose{\textbf{หิตฺวาน กามานิ ยโถธิกานี}ติ \textbf{กามา}ติ อุทฺทานโต เทฺว กามา วตฺถุกามา จ กิเลสกามา จ ฯเปฯ. อิเม วุจฺจนฺติ วตฺถุกามา ฯเปฯ. อิเม วุจฺจนฺติ กิเลสกามา. \textbf{หิตฺวาน กามานี}ติ วตฺถุกาเม ปริชานิตฺวา กิเลสกาเม ปหายา ปชหิตฺวา วิโนเทตฺวา พฺยนฺตีกริตฺวา อนภาวํ คเมตฺวา. \textbf{หิตฺวาน กามานิ ยโถธิกานี}ติ โสตาปตฺติ\-มคฺเคน เย กิเลสา ปหีนา, เต กิเลเส น ปุเนติ น ปจฺเจติ น ปจฺจาคจฺฉติ. สกทาคามิ\-มคฺเคน เย กิเลสา ปหีนา. อนาคามิมคฺเคน เย กิเลสา ปหีนา. อรหตฺต\-มคฺเคน เย กิเลสา ปหีนา, เต กิเลเส น ปุเนติ น ปจฺเจติ น ปจฺจาคจฺฉตีติ หิตฺวาน กามานิ ยโถธิกานิ, เอโก จเร ขคฺควิสาณ\-กปฺโป. เตนาห โส ปจฺเจก\-สมฺพุทฺโธ\csromandash{}}

\prose{ปุตฺตญฺจ ทารํ ปิตรญฺจ มาตรํ,}

\prose{ธนานิ ธญฺญานิ จ พนฺธวานิ.}

\setlength{\csromangathapagewidth}{0pt}
\setlength{\csromangathawakwidth}{0pt}
\csromangathameasuregroup{}{หิตฺวาน กามานิ ยโถธิกานิ, \\ เอโก จเร ขคฺควิสาณกปฺโปติ. \\ \textbf{สงฺโค เอโส ปริตฺตเมตฺถ โสขฺยํ,} \\ \textbf{อปฺปสฺสาโท ทุกฺขเมตฺถ ภิยฺโย.} \\ \textbf{คโฬ เอโส อิติ ญ ตฺวา มติมา}\textsuperscript{1}\textbf{,} \\ \textbf{เอโก จเร ขคฺควิสาณ}\textbf{กปฺโป.}}
\csromangathameasurewak{หิตฺวาน กามานิ ยโถธิกานิ, \\ เอโก จเร ขคฺควิสาณกปฺโปติ. \\ \textbf{สงฺโค เอโส ปริตฺตเมตฺถ โสขฺยํ,} \\ \textbf{อปฺปสฺสาโท ทุกฺขเมตฺถ ภิยฺโย.} \\ \textbf{คโฬ เอโส อิติ ญ ตฺวา มติมา}\textsuperscript{1}\textbf{,} \\ \textbf{เอโก จเร ขคฺควิสาณ}\textbf{กปฺโป.}}
\setlength{\csromangathaleftcol}{0pt}
\csromangathagroup{\csromangathastanza{\csromangathawak{หิตฺวาน กามานิ ยโถธิกานิ, \\ เอโก จเร ขคฺควิสาณ\-กปฺโปติ.}}\csromangathastanza{\csromangathawak{\csromanfolio{282}\csromangathaitemline{147}{\textbf{สงฺโค เอโส ปริตฺตเมตฺถ โสขฺยํ,}} \\ \csromangathacontline{\textbf{อปฺปสฺสาโท ทุกฺขเมตฺถ ภิยฺโย.}} \\ \csromangathacontline{\textbf{คโฬ เอโส อิติ ญ ตฺวา มติมา}\csromansharedfootnote{6a0467155b431f0c}{มุตีมา (ขุ 1. 288 ปิฏฺเฐ).}\textbf{,}} \\ \csromangathacontline{\textbf{เอโก จเร ขคฺควิสาณ}\-\textbf{กปฺโป.}}}}}

\prose{\textbf{สงฺโค เอโส ปริตฺตเมตฺถ โสขฺยนฺ}ติ สงฺโคติ วา พฬิสนฺติ วา อามิสนฺติ วา ลคฺคนนฺติ วา ปลิโพโธติ วา ปญฺจนฺเนตํ กามคุณานํ อธิวจนํ. \textbf{ปริตฺตเมตฺถ โสขฺยนฺ}ติ วุตฺตํ เหตํ ภควตา\csromandash{}“ปญฺจิเม ภิกฺขเว กามคุณา. กตเม ปญฺจ, จกฺขุวิญฺเญยฺยา รูปา อิฏฺฐา กนฺตา มนาปา ปิยรูปา กามูปสํหิตา รชนียา ฯเปฯ กายวิญฺเญยฺยา โผฏฺฐพฺพา อิฏฺฐา กนฺตา มนาปา ปิยรูปา กามูปสํหิตา รชนียา, อิเม โข ภิกฺขเว ปญฺจ กามคุณา. ยํ โข ภิกฺขเว อิเม ปญฺจ กามคุเณ ปฏิจฺจ อุปฺปชฺชติ สุขํ โสมนสฺสํ, อิทํ วุจฺจติ กามสุขํ. อปฺปกํ เอตํ สุขํ, ปริตฺตกํ เอตํ สุขํ, โถกกํ เอตํ สุขํ, โอมกํ เอตํ สุขํ, ลามกํ เอตํ สุขํ, ฉตุกฺกํ เอตํ สุขนฺ”ติ สงฺโค เอโส ปริตฺตเมตฺถ โสขฺยํ.}

\prose{\textbf{อปฺปสฺสาโท ทุกฺขเมตฺถ ภิยฺโย}ติ อปฺปสฺสาทา กามา วุตฺตา ภควตา\csromansharedfootnote{c168abe05435a11a}{ปสฺส ม 1. 185 ปิฏฺเฐ.} พหุทุกฺขา พหุปายาสา\csromansharedfootnote{ddcf7ed00fa308d3}{พหูปายาสา (สฺยา)}, อาทีนโว เอตฺถ ภิยฺโย. อฏฺฐิ\-กงฺกลู\-ปมา กามา วุตฺตา ภควตา. มํสเปสูปมา กามา วุตฺตา ภควตา. ติณุกฺกูปมา กามา วุตฺตา ภควตา. องฺคารกาสูปมา กามา วุตฺตา ภควตา. สุปินกูปมา กามา วุตฺตา ภควตา. ยาจิตกูปมา กามา วุตฺตา ภควตา. รุกฺข\-ผลู\-ปมา กามา วุตฺตา ภควตา. อสิสูนูปมา กามา วุตฺตา ภควตา, สตฺติสูลูปมา กามา วุตฺตา ภควตา. สปฺปสิรูปมา กามา วุตฺตา ภควตา พหุทุกฺขา พหุปายาสา. อาทีนโว เอตฺถ ภิยฺโยติ อปฺปสฺสาโท ทุกฺขเมตฺถ ภิยฺโย.}

\prose{\textbf{คโฬ เอโส อิติ ญ ตฺวา มติมา}ติ คโฬติ วา พฬิสนฺติ วา อามิสนฺติ วา ลคฺคนนฺติ วา พนฺธนนฺติ วา ปลิโพโธติ วา ปญฺจนฺเนตํ กามคุณานํ อธิวจนํ. \textbf{อิตี}ติ ปทสนฺธิ ปทสํสคฺโค ปทปาริปูรี อกฺขร\-สมวาโย พฺยญฺชน\-สิลิฏฺฐตา ปทา\-นุปุพฺพตา\-เปตํ อิตีติ. \textbf{มติมา}ติ ปณฺฑิโต ปญฺญวา พุทฺธิมา ญาณี วิภาวี เมธาวี. \textbf{คโฬ เอโส อิติ ญ ตฺวา มติมา}ติ มติมา คโฬติ ญ ตฺวา พฬิสนฺติ ญ ตฺวา อามิสนฺติ}

\prose{\csromanfolio{283}ญตฺวา ลคฺคนนฺติ ญ ตฺวา พนฺธนนฺติ ญ ตฺวา ปลิโพโธติ ญ ตฺวา ชานิตฺวา ตุลยิตฺวา ตีรยิตฺวา วิภาวยิตฺวา วิภูตํ กตฺวาติ คโฬ เอโส อิติ ญ ตฺวา มติมา, เอโก จเร ขคฺควิสาณ\-กปฺโป. เตนาห โส ปจฺเจก\-สมฺพุทฺโธ\csromandash{}}

\setlength{\csromangathapagewidth}{0pt}
\setlength{\csromangathawakwidth}{0pt}
\csromangathameasuregroup{}{สงฺโค เอโส ปริตฺตเมตฺถ โสขฺยํ, \\ อปฺปสฺสาโท ทุกฺขเมตฺถ ภิยฺโย. \\ คโฬ เอโส อิติ ญตฺวา มติมา, \\ เอโก จเร ขคฺควิสาณกปฺโปติ. \\ \textbf{สนฺทาลยิตฺวาน สญฺโญชนานิ,} \\ \textbf{ชาลํว เภตฺวา สลิลมฺพุจารี.} \\ \textbf{อคฺคีว ทฑฺฒํ อนิวตฺตมาโน,} \\ \textbf{เอโก จเร ขคฺควิสาณ}\textbf{กปฺโป.}}
\csromangathameasurewak{สงฺโค เอโส ปริตฺตเมตฺถ โสขฺยํ, \\ อปฺปสฺสาโท ทุกฺขเมตฺถ ภิยฺโย. \\ คโฬ เอโส อิติ ญตฺวา มติมา, \\ เอโก จเร ขคฺควิสาณกปฺโปติ. \\ \textbf{สนฺทาลยิตฺวาน สญฺโญชนานิ,} \\ \textbf{ชาลํว เภตฺวา สลิลมฺพุจารี.} \\ \textbf{อคฺคีว ทฑฺฒํ อนิวตฺตมาโน,} \\ \textbf{เอโก จเร ขคฺควิสาณ}\textbf{กปฺโป.}}
\setlength{\csromangathaleftcol}{0pt}
\csromangathagroup{\csromangathastanza{\csromangathawak{สงฺโค เอโส ปริตฺตเมตฺถ โสขฺยํ, \\ อปฺปสฺสาโท ทุกฺขเมตฺถ ภิยฺโย. \\ คโฬ เอโส อิติ ญตฺวา มติมา, \\ เอโก จเร ขคฺควิสาณ\-กปฺโปติ.}}\csromangathastanza{\csromangathawak{\csromangathaitemline{148}{\textbf{สนฺทาลยิตฺวาน สญฺโญชนานิ,}} \\ \csromangathacontline{\textbf{ชาลํว เภตฺวา สลิลมฺพุจารี.}} \\ \csromangathacontline{\textbf{อคฺคีว ทฑฺฒํ อนิวตฺตมาโน,}} \\ \csromangathacontline{\textbf{เอโก จเร ขคฺควิสาณ}\-\textbf{กปฺโป.}}}}}

\prose{\textbf{สนฺทาลยิตฺวาน สญฺโญชนานี}ติ ทส สญฺโญชนานิ, กามราค\-สญฺโญชนํ ปฏิฆ\-สญฺโญชนํ มานสญฺโญชนํ ทิฏฺฐิ\-สญฺโญชนํ วิจิกิจฺฉา\-สญฺโญชนํ สีลพฺพตปรามาส\-สญฺโญชนํ ภวราค\-สญฺโญชนํ อิสฺสาสญฺโญชนํ มจฺฉริย\-สญฺโญชนํ อวิชฺชา\-สญฺโญชนํ. \textbf{สนฺทาลยิตฺวาน สญฺโญชนานี}ติ ทส สญฺโญชนานิ ทาลยิตฺวา สนฺทาลยิตฺวา ปชหิตฺวา วิโนเทตฺวา พฺยนฺตีกริตฺวา อนภาวํ คเมตฺวาติ สนฺทาลยิตฺวาน สญฺโญชนานิ.}

\prose{\textbf{ชาลํว เภตฺวา สลิลมฺพุจารี}ติ ชาลํ วุจฺจติ สุตฺตชาลํ. สลิลํ วุจฺจติ อุทกํ, อมฺพุจารี วุจฺจติ มจฺโฉ, ยถา มจฺโฉ ชาลํ ภินฺทิตฺวา ปภินฺทิตฺวา ทาลยิตฺวา ปทาลยิตฺวา สมฺปทาลยิตฺวา จรติ วิหรติ อิริยติ วตฺเตติ ปาเลติ ยเปติ ยาเปติ, เอวเมว เทฺว ชาลา ตณฺหาชาลญฺจ ทิฏฺฐิชาลญฺจ ฯเปฯ อิทํ ตณฺหาชาลํ ฯเปฯ อิทํ ทิฏฺฐิชาลํ. ตสฺส ปจฺเจก\-สมฺพุทฺธสฺส ตณฺหาชาลํ ปหีนํ, ทิฏฺฐิชาลํ ปฏินิสฺสฏฺฐํ. ตณฺหาชาลสฺส ปหีนตฺตา ทิฏฺฐิชาลสฺส ปฏินิสฺสฏฺฐ\-ตฺตา โส ปจฺเจก\-สมฺพุทฺโธ รูเป น สชฺชติ, สทฺเท น สชฺชติ, คนฺเธ น สชฺชติ ฯเปฯ ทิฏฺฐ\-สุต\-มุต\-วิญฺญาตพฺเพสุ ธมฺเมสุ น สชฺชติ น คณฺหาติ น พชฺฌติ น ปลิพชฺฌติ นิกฺขนฺโต นิสฺสโฏ วิปฺปมุตฺโต วิสญฺญุตฺโต วิมริยาทิกเตน เจตสา วิหรตีติ ชาลํว เภตฺวา สลิลมฺพุจารี.}

\prose{\csromanfolio{284}\textbf{อคฺคีว ทฑฺฒํ อนิวตฺตมาโน}ติ ยถา อคฺคิ ติณกฏฺฐุ\-ปาทานํ ทหนฺโต คจฺฉติ อนิวตฺตนฺโต, เอวเมว ตสฺส ปจฺเจก\-สมฺพุทฺธสฺส โสตาปตฺติ\-มคฺเคน เย กิเลสา ปหีนา, เต กิเลเส น ปุเนติ น ปจฺเจติ น ปจฺจาคจฺฉติ. สกทาคามิ\-มคฺเคน. อนาคามิมคฺเคน. อรหตฺต\-มคฺเคน เย กิเลสา ปหีนา, เต กิเลเส น ปุเนติ น ปจฺเจติ น ปจฺจาคจฺฉตีติ อคฺคีว ทฑฺฒํ อนิวตฺตมาโน, เอโก จเร ขคฺควิสาณ\-กปฺโป. เตนาห โส ปจฺเจก\-สมฺพุทฺโธ\csromandash{}}

\setlength{\csromangathapagewidth}{0pt}
\setlength{\csromangathawakwidth}{0pt}
\csromangathameasuregroup{}{สนฺทาลยิตฺวาน สญฺโญชนานิ, \\ ชาลํว เภตฺวา สลิลมฺพุจารี. \\ อคฺคีว ทฑฺฒํ อนิวตฺตมาโน, \\ เอโก จเร ขคฺควิสาณกปฺโปติ.}
\csromangathameasurewak{สนฺทาลยิตฺวาน สญฺโญชนานิ, \\ ชาลํว เภตฺวา สลิลมฺพุจารี. \\ อคฺคีว ทฑฺฒํ อนิวตฺตมาโน, \\ เอโก จเร ขคฺควิสาณกปฺโปติ.}
\setlength{\csromangathaleftcol}{0pt}
\csromangathagroup{\csromangathastanza{\csromangathawak{สนฺทาลยิตฺวาน สญฺโญชนานิ, \\ ชาลํว เภตฺวา สลิลมฺพุจารี. \\ อคฺคีว ทฑฺฒํ อนิวตฺตมาโน, \\ เอโก จเร ขคฺควิสาณ\-กปฺโปติ.}}}

\proseitem{149}{\textbf{โอกฺขิตฺตจกฺขุ น จ ปาทโลโล,}}

\prose{\textbf{คุตฺตินฺทฺริโย รกฺขิต}\-\textbf{มานสาโน.}}

\prose{\textbf{อนวสฺสุโต อปริฑยฺหมาโน}\csromansharedfootnote{6b91851573898a12}{อปริทยฺหมาโน (ก)}\textbf{,}}

\prose{\textbf{เอโก จเร ขคฺควิสาณ}\-\textbf{กปฺโป.} (9)}

\prose{\textbf{โอกฺขิตฺตจกฺขุ น จ ปาทโลโล}ติ กถํ ขิตฺตจกฺขุ โหติ. อิเธกจฺโจ ภิกฺขุ\csromansharedfootnote{0d9c036a54b89ba9}{นตฺถิ สฺยาโปตฺถเก, ขุ 7. 285 ปิฏฺเฐ จ.} จกฺขุโลโล จกฺขุโลลิเยน สมนฺนาคโต โหติ, “อทิฏฺฐํ ทกฺขิตพฺพํ, ทิฏฺฐํ สมติกฺกมิตพฺพนฺ”ติ อาราเมน อารามํ อุยฺยาเนน อุยฺยานํ คาเมน คามํ นิคเมน นิคมํ นคเรน นครํ รฏฺเฐน รฏฺฐํ ชนปเทน ชนปทํ ทีฆจาริกํ อนวฏฺฐิต\-จาริกํ\csromansharedfootnote{65b0826962c4344b}{อนวตฺถจาริกํ (สฺยา)} อนุยุตฺโต โหติ รูปทสฺสนาย, เอวํ ขิตฺตจกฺขุ โหติ.}

\prose{อถ วา ภิกฺขุ อนฺตรฆรํ ปวิฏฺโฐ วีถิํ ปฏิปนฺโน อสํวุโต คจฺฉติ, หตฺถิํ โอโลเกนฺโต อสฺสํ โอโลเกนฺโต รถํ โอโลเกนฺโต ปตฺติํ โอโลเกนฺโต กุมารเก โอโลเกนฺโต กุมาริกาโย โอโลเกนฺโต อิตฺถิโย โอโลเกนฺโต ปุริเส โอโลเกนฺโต อนฺตราปณํ โอโลเกนฺโต ฆรมุขานิ โอโลเกนฺโต อุทฺธํ โอโลเกนฺโต อโธ โอโลเกนฺโต ทิสาวิทิสํ วิเปกฺขมาโน\csromansharedfootnote{52c3052422e12ed3}{เปกฺขมาโน (สฺยา, ก)} คจฺฉติ, เอวมฺปิ ขิตฺตจกฺขุ โหติ.}

\prose{\csromanfolio{285}อถ วา ภิกฺขุ จกฺขุนา รูปํ ทิสฺวา นิมิตฺตคฺคาหี โหติ อนุพฺยญฺชนคฺคาหี, ยตฺวาธิกรณเมนํ จกฺขุนฺทฺริยํ อสํวุตํ วิหรนฺตํ อภิชฺฌา\-โทมนสฺสา ปาปกา อกุสลา ธมฺมา อนฺวาสฺสเวยฺยุํ, ตสฺส สํวราย น ปฏิปชฺชติ, น รกฺขติ จกฺขุนฺทฺริยํ, จกฺขุนฺทฺริเย น สํวรํ อาปชฺชติ, เอวมฺปิ ขิตฺตจกฺขุ โหติ.}

\prose{ยถา วา ปเนเก โภนฺโต สมณพฺราหฺมณา สทฺธาเทยฺยานิ โภชนานิ ภุญฺชิตฺวา เต เอวรูปํ วิสูกทสฺสนํ อนุยุตฺตา วิหรนฺติ. เสยฺยถิทํ, นจฺจํ คีตํ วาทิตํ เปกฺขํ อกฺขานํ ปาณิสฺสรํ เวตาฬํ กุมฺภถูณํ โสภนกํ\csromansharedfootnote{1fbcfc3d18e407f2}{โสภนครกํ (สฺยา), โสภนกรณํ (ก)} จณฺฑาลํ วํสํ โธวนํ หตฺถิยุทฺธํ อสฺสยุทฺธํ มหิํสยุทฺธํ\csromansharedfootnote{abbdbe958b34474e}{มหิสยุทฺธํ (สฺยา)} อุสภยุทฺธํ อชยุทฺธํ เมณฺฑยุทฺธํ กุกฺกุฏยุทฺธํ วฏฺฏกยุทฺธํ ทณฺฑยุทฺธํ มุฏฺฐิยุทฺธํ นิพฺพุทฺธํ อุยฺโยธิกํ พลคฺคํ เสนาพฺยูหํ อนีกทสฺสนํ อิติ วา, อิติ เอวรูปํ วิสูกทสฺสนํ อนุยุตฺโต โหติ, เอวมฺปิ ขิตฺตจกฺขุ โหติ.}

\prose{กถํ โอกฺขิตฺตจกฺขุ โหติ. อิเธกจฺโจ ภิกฺขุ น จกฺขุโลโล น จกฺขุโลลิเยน สมนฺนาคโต โหติ, “อทิฏฺฐํ ทกฺขิตพฺพํ, ทิฏฺฐํ สมติกฺกมิตพฺพนฺ”ติ น อาราเมน อารามํ น อุยฺยาเนน อุยฺยานํ น คาเมน คามํ น นิคเมน นิคมํ น นคเรน นครํ น รฏฺเฐน รฏฺฐํ น ชนปเทน ชนปทํ ทีฆจาริกํ อนวฏฺฐิตจารีกํ อนุยุตฺโต โหติ รูปทสฺสนาย, เอวํ โอกฺขิตฺตจกฺขุ โหติ.}

\prose{อถ วา ภิกฺขุ อนฺตรฆรํ ปวิฏฺโฐ วีถิํ ปฏิปนฺโน สํวุโต คจฺฉติ, น หตฺถิํ โอโลเกนฺโต น อสฺสํ โอโลเกนฺโต น รถํ โอโลเกนฺโต น ปตฺติํ โอโลเกนฺโต น กุมารเก โอโลเกนฺโต น กุมาริกาโย โอโลเกนฺโต น อิตฺถิโย โอโลเกนฺโต น ปุริเส โอโลเกนฺโต น อนฺตราปณํ โอโลเกนฺโต น ฆรมุขานิ โอโลเกนฺโต น อุทฺธํ โอโลเกนฺโต น อโธ โอโลเกนฺโต น ทิสาวิทิสํ วิเปกฺขมาโน คจฺฉติ, เอวมฺปิ โอกฺขิตฺตจกฺขุ โหติ.}

\prose{อถ วา ภิกฺขุ จกฺขุนา รูปํ ทิสฺวา น นิมิตฺตคฺคาหี โหติ นานุพฺยญฺชน\-คฺคาหี, ยตฺวาธิกรณเมนํ จกฺขุนฺทฺริยํ อสํวุตํ วิหรนฺตํ อภิชฺฌา\-โทมนสฺสา ปาปกา อกุสลา ธมฺมา อนฺวาสฺสเวยฺยุํ, ตสฺส สํวราย ปฏิปชฺชติ, รกฺขติ จกฺขุนฺทฺริยํ, จกฺขุนฺทฺริเย สํวรํ อาปชฺชติ, เอวมฺปิ โอกฺขิตฺตจกฺขุ โหติ.}

\prose{\csromanfolio{286}ยถา วา ปเนเก โภนฺโต สมณพฺราหฺมณา สทฺธาเทยฺยานิ โภชนานิ ภุญฺชิตฺวา เต เอวรูปํ วิสูกทสฺสนํ อนุยุตฺตา วิหรนฺติ. เสยฺยถิทํ, นจฺจํ คีตํ วาทิตํ ฯเปฯ อนีกทสฺสนํ อิติ วา, อิติ เอวรูปา วิสูกทสฺสนา ปฏิวิรโต, เอวมฺปิ โอกฺขิตฺตจกฺขุ โหติ.}

\prose{\textbf{น จ ปาทโลโล}ติ กถํ ปาทโลโล โหติ. อิเธกจฺโจ ภิกฺขุ ปาทโลโล ปาทโลลิเยน สมนฺนาคโต โหติ, อาราเมน อารามํ อุยฺยาเนน อุยฺยานํ คาเมน คามํ นิคเมน นิคมํ นคเรน นครํ รฏฺเฐน รฏฺฐํ ชนปเทน ชนปทํ ทีฆจาริกํ อนวฏฺฐิต\-จาริกํ อนุยุตฺโต โหติ, เอวมฺปิ ปาทโลโล โหติ.}

\prose{อถ วา ภิกฺขุ อนฺโตสํฆาราเม\csromansharedfootnote{7e642708f0c727ca}{อนฺโตปิ สํฆาราเม (ก)} ปาทโลโล ปาทโลลิเยน สมนฺนาคโต โหติ, น อตฺถเหตุ น การณเหตุ อุทฺธโต อวูปสนฺต\-จิตฺโต ปริเวณโต ปริเวณํ คจฺฉติ, วิหารโต วิหารํ คจฺฉติ, อฑฺฒโยคโต อฑฺฒโยคํ คจฺฉติ, ปาสาทโต ปาสาทํ คจฺฉติ, หมฺมิยโต หมฺมิยํ คจฺฉติ, คุหโต คุหํ คจฺฉติ, เลณโต เลณํ คจฺฉติ, กุฏิยา กุฏิํ คจฺฉติ, กูฏาคารโต กูฏาคารํ คจฺฉติ, อฏฺฏโต อฏฺฏํ คจฺฉติ, มาฬโต มาฬํ คจฺฉติ, อุทฺทณฺฑโต อุทฺทณฺฑํ คจฺฉติ\csromansharedfootnote{ff42177fea2c93f6}{อุทฺทณฺฑํ คจฺฉติ, อุทฺโธสิตโต อุทฺโธสิตํ คจฺฉติ (สฺยา) ปสฺส ขุ 7. 51 ปิฏฺเฐ.}, อุปฏฺฐาน\-สาลโต อุปฏฺฐานสาลํ คจฺฉติ, มณฺฑปโต มณฺฑปํ คจฺฉติ, รุกฺขมูลโต รุกฺขมูลํ คจฺฉติ. ยตฺถ วา ปน ภิกฺขู นิสีทนฺติ วา คจฺฉนฺติ วา, ตตฺถ เอกสฺส วา ทุติโย โหติ, ทฺวินฺนํ วา ตติโย โหติ, ติณฺณํ วา จตุตฺโถ โหติ, ตตฺถ พหุํ สมฺผปฺปลาปํ ปลปติ. เสยฺยถิทํ, ราชกถํ โจรกถํ ฯเปฯ อิติภวา\-ภวกถํ กเถติ, เอวมฺปิ ปาทโลโล โหติ.}

\prose{\textbf{น จ ปาทโลโล}ติ โส ปจฺเจก\-สมฺพุทฺโธ ปาทโลลิยา อารโต วิรโต ปฏิวิรโต นิกฺขนฺโต นิสฺสโฏ วิปฺปมุตฺโต วิสญฺญุตฺโต วิมริยาทิกเตน เจตสา, ปฏิสลฺลานา\-ราโม โหติ ปฏิสลฺลาน\-รโต อชฺฌตฺตํ เจโตสมถม\-นุยุตฺโต อนิรากต\-ชฺฌาโน วิปสฺสนาย สมนฺนาคโต พฺรูเหตา สุญฺญาคารํ ฌายี ฌานรโต เอกตฺตม\-นุยุตฺโต สทตฺถครุโกติ โอกฺขิตฺตจกฺขุ น จ ปาทโลโล.}

\prose{\csromanfolio{287}\textbf{คุตฺตินฺทฺริโย รกฺขิต}\-\textbf{มานสาโน}ติ \textbf{คุตฺตินฺทฺริโย}ติ โส ปจฺเจก\-สมฺพุทฺโธ จกฺขุนา รูปํ ทิสฺวา น นิมิตฺตคฺคาหี โหติ นานุพฺยญฺชน\-คฺคาหี, ยตฺวาธิกรณเมนํ จกฺขุนฺทฺริยํ อสํวุตํ วิหรนฺตํ อภิชฺฌา\-โทมนสฺสา ปาปกา อกุสลา ธมฺมา อนฺวาสฺสเวยฺยุํ, ตสฺส สํวราย ปฏิปชฺชติ, รกฺขติ จกฺขุนฺทฺริยํ, จกฺขุนฺทฺริเย สํวรํ อาปชฺชติ, โสเตน สทฺทํ สุตฺวา ฯเปฯ ฆาเนน คนฺธํ ฆายิตฺวา. ชิวฺหาย รสํ สายิตฺวา. กาเยน โผฏฺฐพฺพํ ผุสิตฺวา. มนสา ธมฺมํ วิญฺญาย น นิมิตฺตคฺคาหี โหติ นานุพฺยญฺชน\-คฺคาหี, ยตฺวาธิกรณเมนํ มนินฺทฺริยํ อสํวุตํ วิหรนฺตํ อภิชฺฌา\-โทมนสฺสา ปาปกา อกุสลา ธมฺมา อนฺวาสฺสเวยฺยุํ, ตสฺส สํวราย ปฏิปชฺชติ, รกฺขติ มนินฺทฺริยํ, มนินฺทฺริเย สํวรํ อาปชฺชตีติ คุตฺตินฺทฺริโย. \textbf{รกฺขิต}\-\textbf{มานสาโน}ติ โคปิตมานสาโนติ คุตฺตินฺทฺริโย รกฺขิต\-มานสาโน.}

\prose{\textbf{อนวสฺสุโต ปริฑยฺห}\-\textbf{มาโน}ติ วุตฺตํ เหตํ อายสฺมตา มหา\-โมคฺคลฺลาเนน\csromandash{}อวสฺสุตปริยายญฺจ\csromansharedfootnote{4c70575f7ee4adf2}{ปสฺส สํ 2. 390 ปิฏฺเฐ.} โว อาวุโส เทเสสฺสามิ อนวสฺสุตปริยายญฺจ, ตํ สุณาถ สาธุกํ มนสิกโรถ, ภาสิสฺสามีติ. “เอวมาวุโส”ติ โข เต ภิกฺขู อายสฺมโต มหา\-โมคฺคลฺลานสฺส ปจฺจสฺโสสุํ. อายสฺมา มหาโมคฺคลฺลาโน\csromansharedfootnote{c938078d1438d43e}{มหาโมคฺคลาโน (ก)} เอตทโวจ\csromandash{}}

\prose{กถญฺจาวุโส อวสฺสุโต โหติ, อิธาวุโส ภิกฺขุ จกฺขุนา รูปํ ทิสฺวา ปิยรูเป รูเป อธิมุจฺจติ, อปฺปิยรูเป รูเป พฺยาปชฺชติ, อนุปฏฺฐิต\-กายสฺสติ\csromansharedfootnote{af2773ae6ebd4241}{อนุปฏฺฐิตกายสติ (สฺยา, ก)} จ วิหรติ ปริตฺตเจตโส, ตญฺจ เจโตวิมุตฺติํ ปญฺญาวิมุตฺติํ ยถาภูตํ นปฺปชานาติ, ยตฺถสฺส\csromansharedfootnote{e3403b9e29e8c651}{ตตฺถ เย (ก) สํ 2. 390 ปิฏฺเฐ.} เต อุปฺปนฺนา ปาปกา อกุสลา ธมฺมา อปริเสสา นิรุชฺฌนฺติ. โสเตน สทฺทํ สุตฺวา ฯเปฯ. มนสา ธมฺมํ วิญฺญาย ปิยรูเป ธมฺเม อธิมุจฺจติ, อปฺปิยรูเป ธมฺเม พฺยาปชฺชติ, อนุปฏฺฐิต\-กายสฺสติ จ วิหรติ ปริตฺตเจตโส, ตญฺจ เจโตวิมุตฺติํ ปญฺญาวิมุตฺติํ ยถาภูตํ นปฺปชานาติ, ยตฺถสฺส เต อุปฺปนฺนา ปาปกา อกุสลา ธมฺมา อปริเสสา นิรุชฺฌนฺติ. อยํ วุจฺจตาวุโส ภิกฺขุ อวสฺสุโต จกฺขุ\-วิญฺเญยฺเยสุ รูเปสุ ฯเปฯ อวสฺสุโต มโนวิญฺเญยฺเยสุ ธมฺเมสุ เอวํวิหาริํ จาวุโส ภิกฺขุํ จกฺขุโต เจปิ นํ มาโร อุปสงฺกมติ, ลเภเถว\csromansharedfootnote{81acec1baa6099d3}{ลภเตว (สฺยา, ก) เอวมุปริปิ.} มาโร}

\prose{\csromanfolio{288}โอตารํ, ลเภถ\csromansharedfootnote{885a64ca00d218de}{ลภติ (สฺยา, ก) เอวมุปริปิ.} มาโร อารมฺมณํ. โสตโต เจปิ นํ ฯเปฯ. มนโต เจปิ นํ มาโร อุปสงฺกมติ, ลเภเถว มาโร โอตารํ, ลเภถ มาโร อารมฺมณํ.}

\prose{เสยฺยถาปิ อาวุโส นฬาคารํ วา ติณาคารํ วา สุกฺขํ โกฬาปํ\csromansharedfootnote{4f18a912827f13b7}{โกลาปํ (สฺยา) สํ 2. 391 ปิฏฺเฐปิ.} เตโรวสฺสิกํ. ปุรตฺถิมาย เจปิ นํ ทิสาย ปุริโส อาทิตฺตาย ติณุกฺกาย อุปสงฺกเมยฺย, ลเภเถว อคฺคิ โอตารํ, เลเภถ อคฺคิ อารมฺมณํ. ปจฺฉิมาย เจปิ นํ ทิสาย ฯเปฯ. อุตฺตราย เจปิ นํ ทิสาย. ทกฺขิณาย เจปิ นํ ทิสาย. เหฏฺฐิมโต\csromansharedfootnote{db6f0d5644116dab}{ปจฺฉโต (สฺยา), เหฏฺฐิมาย (ก)} เจปิ นํ ทิสาย. อุปริมโต\csromansharedfootnote{a2aa690251a2a6b4}{อุปริโต (สฺยา), อุปริมาย (ก)} เจปิ นํ ทิสาย. ยโต กุโตจิ เจปิ นํ ปุริโส อาทิตฺตาย ติณุกฺกาย อุปสงฺกเมยฺย, ลเภเถว อคฺคิ โอตารํ. ลเภถ อคฺคิ อารมฺมณํ. เอวเมว โข อาวุโส เอวํวิหาริํ ภิกฺขุํ จกฺขุโต เจปิ นํ มาโร อุปสงฺกมติ. ลเภเถว มาโร โอตารํ, ลเภถ มาโร อารมฺมณํ. โสตโต เจปิ นํ ฯเปฯ. มนโต เจปิ นํ มาโร อุปสงฺกมติ, ลเภเถว มาโร โอตารํ, ลเภถ มาโร อารมฺมณํ.}

\prose{เอวํวิหาริํ จาวุโส ภิกฺขุํ รูปา อธิภํสุ\csromansharedfootnote{2d0cb0ac59df74d9}{อภิภวิํสุ (สฺยา), อภิภํสุ (ก) เอวมุปริปิ.}, น ภิกฺขุ รูเป อธิโภสิ. สทฺทา ภิกฺขุํ อธิภํสุ, น ภิกฺขุ สทฺเท อธิโภสิ. คนฺธา ภิกฺขุํ อธิภํสุ, น ภิกฺขุ คนฺเธ อธิโภสิ. รสา ภิกฺขุํ อธิภํสุ, น ภิกฺขุ รเส อธิโภสิ. โผฏฺฐพฺพา ภิกฺขุํ อธิภํสุ, น ภิกฺขุ โผฏฺฐพฺเพ อธิโภสิ. ธมฺมา ภิกฺขุํ อธิภํสุ, น ภิกฺขุ ธมฺเม อธิโภสิ. อยํ วุจฺจตาวุโส ภิกฺขุ รูปาธิภูโต สทฺทาธิภูโต คนฺธาธิภูโต รสาธิภูโต โผฏฺฐพฺพา\-ธิภูโต ธมฺมาธิภูโต อธิภูโต อนธิภู, อธิภํสุ นํ ปาปกา อกุสลา ธมฺมา สํกิเลสิกา โปโนภวิกา สทรา ทุกฺขวิปากา อายติํ ชาติ\-ชรา\-มรณิยา, เอวํ โข อาวุโส อวสฺสุโต โหติ.}

\prose{กถญฺจาวุโส อนวสฺสุโต โหติ. อิธาวุโส ภิกฺขุ จกฺขุนา รูปํ ทิสฺวา ปิยรูเป รูเป นาธิมุจฺจติ, อปฺปิยรูเป รูเป น พฺยาปชฺชติ, อุปฏฺฐิต\-กายสฺสติ จ วิหรติ อปฺปมาณเจตโส. ตญฺจ เจโตวิมุตฺติํ}

\prose{\csromanfolio{289}ปญฺญาวิมุตฺติํ ยถาภูตํ ปชานาติ, ยตฺถสฺส เต อุปฺปนฺนา ปาปกา อกุสลา ธมฺมา อปริเสสา นิรุชฺฌนฺติ. โสเตน สทฺทํ สุตฺวา ฯเปฯ. มนสา ธมฺมํ วิญฺญาย ปิยรูเป ธมฺเม นาธิมุจฺจติ, อปฺปิยรูเป ธมฺเม น พฺยาปชฺชติ, อุปฏฺฐิต\-กายสฺสติ จ วิหรติ อปฺปมาณเจตโส. ตญฺจ เจโตวิมุตฺติํ ปญฺญาวิมุตฺติํ ยถาภูตํ ปชานาติ. ยตฺถสฺส เต อุปฺปนฺนา ปาปกา อกุสลา ธมฺมา อปริเสสา นิรุชฺฌนฺติ. อยํ วุจฺจตาวุโส ภิกฺขุ อนวสฺสุโต จกฺขุ\-วิญฺเญยฺเยสุ รูเปสุ ฯเปฯ อนวสฺสุโต มโนวิญฺเญยฺเยสุ ธมฺเมสุ. เอวํวิหาริํ จาวุโส ภิกฺขุํ จกฺขุโต เจปิ นํ มาโร อุปสงฺกมติ, เนว ลเภถ มาโร โอตารํ, น ลเภถ มาโร อารมฺมณํ. โสตโต เจปิ นํ ฯเปฯ. มนโต เจปิ นํ มาโร อุปสงฺกมติ, เนว ลเภถ มาโร โอตารํ, น ลเภถ มาโร อารมฺมณํ.}

\prose{เสยฺยถาปิ อาวุโส กูฏาคารา วา กูฏาคารสาลา\csromansharedfootnote{174b9728c6263164}{สนฺถาคารสาลา (สฺยา) ปสฺส สํ 2. 392 ปิฏฺเฐ.} วา พหลมตฺติกา อทฺทาวเลปนา\csromansharedfootnote{597334dffa5cafad}{อลฺลาวเลปนา (สฺยา)}. ปุรตฺถิมาย เจปิ นํ ทิสาย ปุริโส อาทิตฺตาย ติณุกฺกาย อุปสงฺกเมยฺย, เนว ลเภถ อคฺคิ โอตารํ, น ลเภถ อคฺคิ อารมฺมณํ. ปจฺฉิมาย เจปิ นํ ทิสาย. อุตฺตราย เจปิ นํ ทิสาย. ทกฺขิณาย เจปิ นํ ทิสาย. เหฏฺฐิมโต เจปิ นํ ทิสาย. อุปริมโต เจปิ นํ ทิสาย. ยโต กุโตจิ เจปิ นํ ปุริโส อาทิตฺตาย ติณุกฺกาย อุปสงฺกเมยฺย, เนว ลเภถ อคฺคิ โอตารํ, น ลเภถ อคฺคิ อารมฺมณํ. เอวเมว โข อาวุโส เอวํวิหาริํ ภิกฺขุํ จกฺขุโต เจปิ นํ มาโร อุปสงฺกมติ, เนว ลเภถ มาโร โอตารํ, น ลเภถ มาโร อารมฺมณํ. โสตโต เจปิ นํ ฯเปฯ. มนโต เจปิ นํ มาโร อุปสงฺกมติ, เนว ลเภถ มาโร โอตารํ, น ลเภถ มาโร อารมฺมณํ.}

\prose{เอวํวิหารี จาวุโส ภิกฺขุ รูเป อธิโภสิ\csromansharedfootnote{dffaa0a16aaa319e}{อภิภวิํสุ (สฺยา), อภิโภสิ (ก) ปสฺส สํ 2. 392 ปิฏฺเฐ.}, น รูปา ภิกฺขุํ อธิภํสุ. สทฺเท ภิกฺขุ อธิโภสิ, น สทฺทา ภิกฺขุํ อธิภํสุ. คนฺเธ ภิกฺขุ อธิโภสิ, น คนฺธา ภิกฺขุํ อธิภํสุ. รเส ภิกฺขุ อธิโภสิ, น รสา ภิกฺขุํ อธิภํสุ. โผฏฺฐพฺเพ ภิกฺขุ อธิโภสิ, น โผฏฺฐพฺพา ภิกฺขุํ อธิภํสุ. ธมฺเม ภิกฺขุ อธิโภสิ, น ธมฺมา ภิกฺขุํ อธิภํสุ. อยํ}

\prose{\csromanfolio{290}วุจฺจตาวุโส ภิกฺขุ รูปาธิภู สทฺทาธิภู คนฺธาธิภู รสาธิภู โผฏฺฐพฺพาธิภู ธมฺมาธิภู อธิภู อนธิภูโต\csromansharedfootnote{78a751b3e7b08ef9}{อนภิภูโต เกหิ จิ กิเลเสหิ (ก) ปสฺส สํ 2. 392 ปิฏฺเฐ.} อธิโภสิ เต ปาปเก อกุสเล ธมฺเม สํกิเลสิเก โปโนภวิเก สทเร ทุกฺขวิปาเก อายติํ ชาติ\-ชรา\-มรณิเย. เอวํ โข อาวุโส อนวสฺสุโต โหตีติ อนวสฺสุโต.}

\prose{\textbf{อปริฑยฺหมาโน}ติ ราคเชน ปริฬาเหน อปริฑยฺหมาโน, โทสเชน ปริฬาเหน อปริฑยฺหมาโน, โมหเชน ปริฬาเหน อปริฑยฺหมาโนติ อนวสฺสุโต อปริฑยฺหมาโน, เอโก จเร ขคฺควิสาณ\-กปฺโป. เตนาห โส ปจฺเจก\-สมฺพุทฺโธ\csromandash{}}

\setlength{\csromangathapagewidth}{0pt}
\setlength{\csromangathawakwidth}{0pt}
\csromangathameasuregroup{}{โอกฺขิตฺตจกฺขุ น จ ปาทโลโล, \\ คุตฺตินฺทฺริโย รกฺขิตมานสาโน. \\ อนวสฺสุโต อปริฑยฺหมาโน, \\ เอโก จเร ขคฺควิสาณกปฺโปติ.}
\csromangathameasurewak{โอกฺขิตฺตจกฺขุ น จ ปาทโลโล, \\ คุตฺตินฺทฺริโย รกฺขิตมานสาโน. \\ อนวสฺสุโต อปริฑยฺหมาโน, \\ เอโก จเร ขคฺควิสาณกปฺโปติ.}
\setlength{\csromangathaleftcol}{0pt}
\csromangathagroup{\csromangathastanza{\csromangathawak{โอกฺขิตฺตจกฺขุ น จ ปาทโลโล, \\ คุตฺตินฺทฺริโย รกฺขิต\-มานสาโน. \\ อนวสฺสุโต อปริฑยฺหมาโน, \\ เอโก จเร ขคฺควิสาณ\-กปฺโปติ.}}}

\proseitem{150}{\textbf{โอหารยิตฺวา คิหิพฺยญฺชนานิ,}}

\prose{\textbf{สญฺฉนฺนปตฺโต ยถา ปาริฉตฺตโก.}}

\prose{\textbf{กาสายวตฺโถ อภินิ}\-\textbf{กฺขมิ}\-\textbf{ตฺวา,}}

\prose{\textbf{เอโก จเร ขคฺควิสาณ}\-\textbf{กปฺโป. (10)}}

\prose{\textbf{โอหารยิตฺวา คิหิพฺยญฺชนานี}ติ คิหิพฺยญฺชนานิ วุจฺจนฺติ เกสา จ มสฺสู จ ฯเปฯ ทีฆทสานิ อิติ วา. \textbf{โอหารยิตฺวา คิหิพฺยญฺชนานี}ติ คิหิพฺยญฺชนานิ โอโรปยิตฺวา สโมโรปยิตฺวา นิกฺขิปิตฺวา ปฏิปฺปสฺสมฺภิตฺวาติ โอหารยิตฺวา คิหิพฺยญฺชนานิ.}

\prose{\textbf{สญฺฉนฺนปตฺโต ยถา ปาริฉตฺตโก}ติ ยถา โส ปาริจฺฉตฺตโก โกวิฬาโร พหล\-ปตฺต\-ปลาโส\csromansharedfootnote{1ee9c2b3169cf864}{สาขปตฺตปลาโส (ก)} สนฺทจฺฉาโย\csromansharedfootnote{338c1577c43c5b7b}{สณฺฑจฺฉาโย (สฺยา), สนฺตจฺฉาโย (ก) ปสฺส ม 1. 108 ปิฏฺเฐ.}. เอวเมว โส ปจฺเจก\-สมฺพุทฺโธ ปริปุณฺณปตฺตจีวรธโรติ สญฺฉนฺนปตฺโต ยถา ปาริฉตฺตโก.}

\prose{\csromanfolio{291}\textbf{กาสายวตฺโถ อภินิ}\-\textbf{กฺขมิ}\-\textbf{ตฺวา}ติ โส ปจฺเจก\-สมฺพุทฺโธ สพฺพํ ฆราวาส\-ปลิโพธํ ฉินฺทิตฺวา ปุตฺตทาร\-ปลิโพธํ ฉินฺทิตฺวา ญาติปลิโพธํ ฉินฺทิตฺวา มิตฺตามจฺจ\-ปลิโพธํ ฉินฺทิตฺวา สนฺนิธิ\-ปลิโพธํ ฉินฺทิตฺวา เกสมสฺสุํ โอหาเรตฺวา\csromansharedfootnote{4c0d3cff101453a6}{โอหารยิตฺวา (สฺยา, ก)} กาสายานิ วตฺถานิ อจฺฉาเทตฺวา อคารสฺมา อนคาริยํ ปพฺพชิตฺวา อกิญฺจนภาวํ อุปคนฺตฺวา เอโก จรติ วิหรติ อิริยติ วตฺเตติ ปาเลติ ยเปติ ยาเปตีติ กาสายวตฺโถ อภินิ\-กฺขมิ\-ตฺวา, เอโก จเร ขคฺควิสาณ\-กปฺโป. เตนาห โส ปจฺเจก\-สมฺพุทฺโธ\csromandash{}}

\setlength{\csromangathapagewidth}{0pt}
\setlength{\csromangathawakwidth}{0pt}
\csromangathameasuregroup{}{โอหารยิตฺวา คิหิพฺยญฺชนานิ, \\ สญฺฉนฺนปตฺโต ยถา ปาริฉตฺตโก.}
\csromangathameasurewak{โอหารยิตฺวา คิหิพฺยญฺชนานิ, \\ สญฺฉนฺนปตฺโต ยถา ปาริฉตฺตโก.}
\setlength{\csromangathaleftcol}{0pt}
\csromangathagroup{\csromangathastanza{\csromangathawak{โอหารยิตฺวา คิหิพฺยญฺชนานิ, \\ สญฺฉนฺนปตฺโต ยถา ปาริฉตฺตโก.}}}

\prose{กาสายวตฺโถ อภินิ\-กฺขมิ\-ตฺวา,}

\prose{เอโก จเร ขคฺควิสาณ\-กปฺโปติ.}

\proseitem{151}{\textbf{รเสสุ เคธํ อกรํ อโลโล,}}

\prose{\textbf{อนญฺญโปสี สปทานจารี.}}

\prose{\textbf{กุเล กุเล อปฺปฏิพทฺธจิตฺโต,}}

\prose{\textbf{เอโก จเร ขคฺควิสาณ}\-\textbf{กปฺโป.} (11)}

\prose{\textbf{รเสสุ เคธํ อกรํ อโลโล}ติ รโสติ มูล\textbf{รโส} ขนฺธรโส ตจรโส ปตฺตรโส ปุปฺผรโส ผลรโส อมฺพิลํ มธุรํ ติตฺตกํ กฏุกํ โลณิกํ ขาริกํ ลมฺพิกํ\csromansharedfootnote{f6d5b871bbb44440}{ลมฺพิลํ (สฺยา, ก) ปสฺส อายตนวิภงฺเค 73 ปิฏฺเฐ.} กสาโว สาทุ อสาทุ สีตํ อุณฺหํ. สนฺเตเก สมณพฺราหฺมณา รสคิทฺธา, เต ชิวฺหคฺเคน รสคฺคานิ ปริเยสนฺตา อาหิณฺฑนฺติ, เต อมฺพิลํ ลภิตฺวา อนมฺพิลํ ปริเยสนฺติ, อนมฺพิลํ ลภิตฺวา อมฺพิลํ ปริเยสนฺติ, มธุรํ ลภิตฺวา อมธุรํ ปริเยสนฺติ, อมธุรํ ลภิตฺวา มธุรํ ปริเยสนฺติ, ติตฺตกํ ลภิตฺวา อติตฺตกํ ปริเยสนฺติ, อติกฺกกํ ลภิตฺวา ติตฺตกํ ปริเยสนฺติ, กฏุกํ ลภิตฺวา อกฏุกํ ปริเยสนฺติ, อกฏุกํ ลภิตฺวา กฏุกํ ปริเยสนฺติ, โลณิกํ ลภิตฺวา อโลณิกํ ปริเยสนฺติ, อโลณิกํ ลภิตฺวา โลณิกํ ปริเยสนฺติ, ขาริกํ ลภิตฺวา อขาริกํ ปริเยสนฺติ, อขาริกํ ลภิตฺวา ขาริกํ ปริเยสนฺติ, กสาวํ ลภิตฺวา อกสาวํ ปริเยสนฺติ, อกสาวํ ลภิตฺวา กสาวํ ปริเยสนฺติ, ลมฺพิกํ ลภิตฺวา อลมฺพิกํ ปริเยสนฺติ, อลมฺพิกํ ลภิตฺวา ลมฺพิกํ ปริเยสนฺติ, สาธุํ ลภิตฺวา อสาทุํ ปริเยสนฺติ, อสาทุํ}

\prose{\csromanfolio{292}ลภิตฺวา สาทุํ ปริเยสนฺติ, สีตํ ลภิตฺวา อุณฺหํ ปริเยสนฺติ, อุณฺหํ ลภิตฺวา สีตํ ปริเยสนฺติ. เต ยํ ยํ ลภนฺติ เตน เตน น ตุสฺสนฺติ, อปราปรํ ปริเยสนฺติ มนาปิเกสุ รเสสุ รตฺตา คิทฺธา คถิตา มุจฺฉิตา อชฺโฌสนฺนา ลคฺคา ลคฺคิตา ปลิพุทฺธา. สา รสตณฺหา ตสฺส ปจฺเจก\-สมฺพุทฺธสฺส ปหีนา อุจฺฉินฺนมูลา ตาลาวตฺถุกตา อนภาวํกตา อายติํ อนุปฺปาทมฺโม, ตสฺมา โส ปจฺเจก\-สมฺพุทฺโธ ปฏิสงฺขา โยนิโส อาหารํ อาหาเรติ “เนว ทวาย, น มทาย, น มณฺฑนาย, น วิภูสนาย, ยาวเทว อิมสฺส กายสฺส ฐิติยา, ยาปนาย, วิหิํสู\-ปรติยา, พฺรหฺมจริยา\-นุคฺคหาย, อิติ ปุราณญฺจ เวทนํ ปฏิหงฺขามิ, นวญฺจ เวทนํ น อุปฺปาเทสฺสามิ, ยาตฺรา จ เม ภวิสฺสติ อนวชฺชตา จ ผาสุวิหาโร จา”ติ.}

\prose{ยถา วณํ อาลิมฺเปยฺย ยาวเทว อารุหณตฺถาย\csromansharedfootnote{ad9aabdaaadee1b5}{โรปนตฺถาย (สฺยา)}, ยถา วา อกฺขํ อพฺภญฺเชยฺย ยาวเทว ภารสฺส นิตฺถรณ\-ตฺถาย, ยถา ปุตฺตมํสํ อาหารํ อาหาเรยฺย ยาวเทว กนฺตารสฺส นิตฺถรณ\-ตฺถาย. เอวเมว โส ปจฺเจก\-สมฺพุทฺโธ ปฏิสงฺขา โยนิโส อาหารํ อาหาเรติ “เนว ทวาย, น มทาย, น มณฺฑนาย, น วิภูสนาย, ยาวเทว อิมสฺส กายสฺส ฐิติยา, ยาปนาย, วิหิํสู\-ปรติยา, พฺรหฺมจริยา\-นุคฺคหาย, อิติ ปุราณญฺจ เวทนํ ปฏิหงฺขามิ, นวญฺจ เวทนํ น อุปฺปาเทสฺสามิ, ยาตฺรา จ เม ภวิสฺสติ อนวชฺชตา จ ผาสุวิหาโร จา”ติ. รสตณฺหาย อารโต วิรโต ปฏิวิรโต นิกฺขนฺโต นิสฺสโฏ วิปฺปมุตฺโต วิสญฺญุตฺโต วิมริยาทิกเตน เจตสา วิหรตีติ รเสสุ เคธํ อกรํ.}

\prose{\textbf{อโลโล}ติ \textbf{โลลุปฺปํ} วุจฺจติ ตณฺหา, โย ราโค สาราโค ฯเปฯ อภิชฺฌา โลโภ อกุสลมูลํ. สา โลลุปฺปา ตณฺหา ตสฺส ปจฺเจก\-สมฺพุทฺธสฺส ปหีนา อุจฺฉินฺนมูลา ตาลาวตฺถุกตา อนภาวํกตา อายติํ อนุปฺปาทธมฺมา, ตสฺมา ปจฺเจก\-สมฺพุทฺโธ อโลโลติ รเสสุ เคธํ อกรํ อโลโล.}

\prose{\textbf{อนญฺญโปสี สปทานจารี}ติ \textbf{อนญฺญโปสี}ติ โส ปจฺเจก\-สมฺพุทฺโธ อตฺตานญฺเญว โปเสติ น ปรนฺติ.}

\setlength{\csromangathapagewidth}{0pt}
\setlength{\csromangathawakwidth}{0pt}
\csromangathameasuregroup{อนญฺญโปสิ’มญฺญาตํ, \\ ขีณาสวํ วนฺตโทสํ,}{\csromangathabat{อนญฺญโปสิ’มญฺญาตํ,}{ทนฺตํ สาเร ปติฏฺฐิตํ\textsuperscript{1}.} \\ \csromangathabat{ขีณาสวํ วนฺตโทสํ,}{ตมหํ พฺรูมิ พฺราหฺมณนฺติ.}}
\csromangathasetleft{อนญฺญโปสิ’มญฺญาตํ, \\ ขีณาสวํ วนฺตโทสํ,}
\csromangathagroup{\csromangathastanza{\csromanfolio{293}\csromangathabat{อนญฺญโปสิ’มญฺญาตํ,}{ทนฺตํ สาเร ปติฏฺฐิตํ\csromansharedfootnote{0c5de68d13044e96}{สาเรสุ สุปติฏฺฐิตํ (สฺยา, ก) ปสฺส ขุ 1. 81 ปิฏฺเฐ.}.} \\ \csromangathabat{ขีณาสวํ วนฺตโทสํ,}{ตมหํ พฺรูมิ พฺราหฺมณนฺติ.}}}

\prose{\textbf{อนญฺญโปสี สปทานจารี}ติ โส ปจฺเจก\-สมฺพุทฺโธ ปุพฺพณฺหสมยํ\csromansharedfootnote{3da6bc654f28e8d1}{ปุพฺพนฺหสมยํ (ก)} นิวาเสตฺวา ปตฺต\-จีวรมา\-ทาย คามํ วา นิคมํ วา ปิณฺฑาย ปวิสติ, รกฺขิเตเนว กาเยน รกฺขิตาย วาจาย รกฺขิเตน จิตฺเตน อุปฏฺฐิตาย สติยา สํวุเตหิ อินฺทฺริเยหิ โอกฺขิตฺตจกฺขุ อิริยาปถ\-สมฺปนฺโน กุลา กุลํ อนติกฺกมนฺโต ปิณฺฑาย จรตีติ อนญฺญโปสี สปทานจารี.}

\prose{\textbf{กุเล กุเล อปฺปฏิพทฺธ}\-\textbf{จิตฺโต}ติ ทฺวีหิ การเณหิ ปฏิพทฺธ\-จิตฺโต โหติ, อตฺตานํ วา นีจํ ฐเปนฺโต ปรํ อุจฺจํ ฐเปนฺโต ปฏิพทฺธ\-จิตฺโต โหติ, อตฺตานํ วา อุจฺจํ ฐเปนฺโต ปรํ นีจํ ฐเปนฺโต ปฏิพทฺธ\-จิตฺโต โหติ. กถํ อตฺตานํ นีจํ ฐเปนฺโต ปรํ อุจฺจํ ฐเปนฺโต ปฏิพทฺธ\-จิตฺโต โหติ, ตุเมฺห เม พหูปการา, อหํ ตุเมฺห นิสฺสาย ลภามิ จีวร\-ปิณฺฑปาต\-เสนาสน\-คิลานปจฺจย\-เภสชฺช\-ปริกฺขารํ. ยมฺปิ เม อญฺเญ ทาตุํ วา กาตุํ วา มญฺญนฺติ ตุเมฺห นิสฺสาย ตุเมฺห ปสฺสนฺตา. ยมฺปิ เม โปราณํ มาตาเปตฺติกํ นามโคตฺตํ, ตมฺปิ เม อนฺตรหิตํ ตุเมฺหหิ อหํ ญายามิ “อสุกสฺส กุลุปโก, อสุกาย กุลุปโก”ติ, เอวํ อตฺตานํ นีจํ ฐเปนฺโต ปรํ อุจฺจํ ฐเปนฺโต ปฏิพทฺธ\-จิตฺโต โหติ.}

\prose{กถํ อตฺตานํ อุจฺจํ ฐเปนฺโต ปรํ นีจํ ฐเปนฺโต ปฏิพทฺธ\-จิตฺโต โหติ, อหํ ตุมฺหากํ พหูปกาโร, ตุเมฺห มํ อาคมฺม พุทฺธํ สรณํ คตา, ธมฺมํ สรณํ คตา, สํฆํ สรณํ คตา, ปาณาติปาตา ปฏิวิรตา, อทินฺนาทานา ปฏิวิรตา, กาเมสุ\-มิจฺฉาจารา ปฏิวิรตา, มุสาวาทา ปฏิวิรตา, สุรา\-เมรย\-มชฺช\-ปมาทฏฺฐานา ปฏิวิรตา. ตุมฺหากํ อหํ อุทฺเทสํ เทมิ, ปริปุจฺฉํ เทมิ, อุโปสถํ อาจิกฺขามิ, นวกมฺมํ อธิฏฺฐามิ. อถ จ ปน ตุเมฺห มํ อุชฺฌิตฺวา อญฺเญ สกฺกโรถ ครุํ กโรถ มาเนถ ปูเชถาติ เอวํ อตฺตานํ อุจฺจํ ฐเปนฺโต ปรํ นีจํ ฐเปนฺโต ปฏิพทฺธ\-จิตฺโต โหติ.}

\prose{\textbf{กุเล กุเล อปฺปฏิพทฺธจิตฺโต}ติ โส ปจฺเจก\-สมฺพุทฺโธ กุลปลิโพเธน อปฺปฏิพทฺธจิตฺโต โหติ, คณปลิโพเธน อปฺปฏิพทฺธจิตฺโต โหติ, อาวาส\-ปลิโพเธน อปฺปฏิพทฺธจิตฺโต โหติ, จีวร\-ปลิโพเธน อปฺปฏิพทฺธจิตฺโต โหติ, ปิณฺฑปาต\-ปลิโพเธน อปฺปฏิพทฺธจิตฺโต โหติ,}

\prose{\csromanfolio{294}เสนาสน\-ปลิโพเธน อปฺปฏิพทฺธจิตฺโต โหติ, คิลานปจฺจยเภสชฺชปริกฺขาร\-ปลิโพเธน อปฺปฏิพทฺธจิตฺโต โหตีติ กุเล กุเล อปฺปฏิพทฺธจิตฺโต โหติ, เอโก จเร ขคฺควิสาณ\-กปฺโป. เตนาห โส ปจฺเจก\-สมฺพุทฺโธ\csromandash{}}

\setlength{\csromangathapagewidth}{0pt}
\setlength{\csromangathawakwidth}{0pt}
\csromangathameasuregroup{}{รเสสุ เคธํ อกรํ อโลโล, \\ อนญฺญโปสี สปทานจารี. \\ กุเล กุเล อปฺปฏิพทฺธจิตฺโต, \\ เอโก จเร ขคฺควิสาณกปฺโปติ.}
\csromangathameasurewak{รเสสุ เคธํ อกรํ อโลโล, \\ อนญฺญโปสี สปทานจารี. \\ กุเล กุเล อปฺปฏิพทฺธจิตฺโต, \\ เอโก จเร ขคฺควิสาณกปฺโปติ.}
\setlength{\csromangathaleftcol}{0pt}
\csromangathagroup{\csromangathastanza{\csromangathawak{รเสสุ เคธํ อกรํ อโลโล, \\ อนญฺญโปสี สปทานจารี. \\ กุเล กุเล อปฺปฏิพทฺธจิตฺโต, \\ เอโก จเร ขคฺควิสาณ\-กปฺโปติ.}}}

\vspace{\tipitakaparskip}
\csromancenter{ตติโย วคฺโค.}
\chapterhead{\textbf{จตุตฺถวคฺค}}
\setlength{\csromangathapagewidth}{0pt}
\setlength{\csromangathawakwidth}{0pt}
\csromangathameasuregroup{}{\textbf{ปหาย ปญฺจาวรณานิ เจตโส,} \\ \textbf{อุปกฺกิเลเส พฺยปนุชฺช สพฺเพ.} \\ \textbf{อนิสฺสิโต เฉตฺว}\textsuperscript{1}\textbf{ สิเนหโทสํ}\textsuperscript{1}\textbf{,} \\ \textbf{เอโก จเร ขคฺควิสาณ}\textbf{กปฺโป.}}
\csromangathameasurewak{\textbf{ปหาย ปญฺจาวรณานิ เจตโส,} \\ \textbf{อุปกฺกิเลเส พฺยปนุชฺช สพฺเพ.} \\ \textbf{อนิสฺสิโต เฉตฺว}\textsuperscript{1}\textbf{ สิเนหโทสํ}\textsuperscript{1}\textbf{,} \\ \textbf{เอโก จเร ขคฺควิสาณ}\textbf{กปฺโป.}}
\setlength{\csromangathaleftcol}{0pt}
\csromangathagroup{\csromangathastanza{\csromangathawak{\csromangathaitemline{152}{\textbf{ปหาย ปญฺจาวรณานิ เจตโส,}} \\ \csromangathacontline{\textbf{อุปกฺกิเลเส พฺยปนุชฺช สพฺเพ.}} \\ \csromangathacontline{\textbf{อนิสฺสิโต เฉตฺว}\csromansharedfootnote{1aeeba1c9fc426fe}{เฉตฺวา (สฺยา, ก)}\textbf{ สิเนหโทสํ}\csromansharedfootnote{0fc053cd0c5828eb}{สฺเนหโทสํ (สฺยา, ก)}\textbf{,}} \\ \csromangathacontline{\textbf{เอโก จเร ขคฺควิสาณ}\-\textbf{กปฺโป.}}}}}

\prose{\textbf{ปหาย ปญฺจาวรณานิ เจตโส}ติ โส ปจฺเจก\-สมฺพุทฺโธ กามจฺฉนฺท\-นีวรณํ ปหาย ปชหิตฺวา วิโนเทตฺวา พฺยนฺตีกริตฺวา อนภาวํ คเมตฺวา, พฺยาปาท\-นีวรณํ. ถินมิทฺธ\-นีวรณํ. อุทฺธจฺจ\-กุกฺกุจฺจ\-นีวรณํ. วิจิกิจฺฉา\-นีวรณํ ปหาย ปชหิตฺวา วิโนเทตฺวา พฺยนฺตีกริตฺวา อนภาวํ คเมตฺวา, วิวิจฺเจว กาเมหิ, วิวิจฺจ อกุสเลหิ ธมฺเมหิ สวิตกฺกํ สวิจารํ วิเวกชํ ปีติสุขํ ปฐมํ ฌานํ อุปสมฺปชฺช วิหรตีติ ปหาย ปญฺจาวรณานิ เจตโส.}

\prose{\textbf{อุปกฺกิเลเส พฺยปนุชฺช สพฺเพ}ติ ราโค จิตฺตสฺส อุปกฺกิเลโส, โทโส จิตฺตสฺส อุปกฺกิเลโส, โมโห จิตฺตสฺส อุปกฺกิเลโส, โกโธ. อุปนาโห ฯเปฯ สพฺพา\-กุสลา\-ภิสงฺขารา จิตฺตสฺส อุปกฺกิเลสา. \textbf{อุปกฺกิเลเส พฺยปนุชฺช สพฺเพ}ติ สพฺเพ จิตฺตสฺส อุปกฺกิเลเส พฺยปนุชฺช ปนุทิตฺวา ปชหิตฺวา วิโนเทตฺวา พฺยนฺตีกริตฺวา อนภาวํ คเมตฺวาติ อุปกฺกิเลเส พฺยปนุชฺช สพฺเพ.}

\prose{\csromanfolio{295}\textbf{อนิสฺสิโต เฉตฺว สิเนหโทสนฺ}ติ \textbf{อนิสฺสิโต}ติ เทฺว นิสฺสยา ตณฺหานิสฺสโย จ ทิฏฺฐินิสฺสโย จ ฯเปฯ อยํ ตณฺหานิสฺสโย ฯเปฯ อยํ ทิฏฺฐินิสฺสโย. \textbf{สิเนโห}ติ เทฺว สฺเนหา ตณฺหาสฺเนโห จ ทิฏฺฐิสฺเนโห จ ฯเปฯ อยํ ตณฺหาสฺเนโห ฯเปฯ อยํ ทิฏฺฐิสฺเนโห. \textbf{โทโส}ติ โย จิตฺตสฺส อาฆาโต ปฏิฆาโต ปฏิฆํ ปฏิวิโรโธ โกโป ปโกโป สมฺปโกโป โทโส ปโทโส สมฺปโทโส จิตฺตสฺส พฺยาปตฺติ มโนปโทโส โกโธ กุชฺฌนา กุชฺฌิตตฺตํ โทโส ทุสฺสนา ทุสฺสิตตฺตํ พฺยาปตฺติ พฺยาปชฺชนา พฺยาปชฺชิตตฺตํ จณฺฑิกฺกํ อสุโรโป อนตฺตมนตา จิตฺตสฺส. \textbf{อนิสฺสิโต เฉตฺว สิเนหโทสนฺ}ติ โส ปจฺเจก\-สมฺพุทฺโธ ตณฺหาสฺเนหญฺจ ทิฏฺฐิสฺเนหญฺจ โทสญฺจ เฉตฺวา อุจฺฉินฺทิตฺวา สมุจฺฉินฺทิตฺวา ปชหิตฺวา วิโนเทตฺวา พฺยนฺตีกริตฺวา อนภาวํ คเมตฺวา จกฺขุํ อนิสฺสิโต, โสตํ อนิสฺสิโต ฯเปฯ. ทิฏฺฐ\-สุต\-มุต\-วิญฺญาตพฺเพ ธมฺเม อนิสฺสิโต อนลฺลีโน อนุปคโต อนชฺโฌสิโต อนธิมุตฺโต นิกฺขนฺโต นิสฺสโฏ วิปฺปมุตฺโต วิสญฺญุตฺโต วิมริยาทิกเตน เจตสา วิหรตีติ อนิสฺสิโต เฉตฺว สิเนหโทสํ, เอโก จเร ขคฺควิสาณ\-กปฺโป. เตนาห โส ปจฺเจก\-สมฺพุทฺโธ\csromandash{}}

\setlength{\csromangathapagewidth}{0pt}
\setlength{\csromangathawakwidth}{0pt}
\csromangathameasuregroup{}{ปหาย ปญฺจาวรณานิ เจตโส, \\ อุปกฺกิเลเส พฺยปนุชฺช สพฺเพ. \\ อนิสฺสิโต เฉตฺว สิเนหโทสํ, \\ เอโก จเร ขคฺควิสาณกปฺโปติ. \\ \textbf{วิปิฏฺฐิ}\textbf{กตฺวาน สุขํ ทุขญฺจ,} \\ \textbf{ปุพฺเพว จ โสมนสฺส}\textbf{โทมนสฺสํ.} \\ \textbf{ลทฺธานุ}\textbf{เปกฺขํ สมถํ วิสุทฺธํ,} \\ \textbf{เอโก จเร ขคฺควิสาณ}\textbf{กปฺโป.}}
\csromangathameasurewak{ปหาย ปญฺจาวรณานิ เจตโส, \\ อุปกฺกิเลเส พฺยปนุชฺช สพฺเพ. \\ อนิสฺสิโต เฉตฺว สิเนหโทสํ, \\ เอโก จเร ขคฺควิสาณกปฺโปติ. \\ \textbf{วิปิฏฺฐิ}\textbf{กตฺวาน สุขํ ทุขญฺจ,} \\ \textbf{ปุพฺเพว จ โสมนสฺส}\textbf{โทมนสฺสํ.} \\ \textbf{ลทฺธานุ}\textbf{เปกฺขํ สมถํ วิสุทฺธํ,} \\ \textbf{เอโก จเร ขคฺควิสาณ}\textbf{กปฺโป.}}
\setlength{\csromangathaleftcol}{0pt}
\csromangathagroup{\csromangathastanza{\csromangathawak{ปหาย ปญฺจาวรณานิ เจตโส, \\ อุปกฺกิเลเส พฺยปนุชฺช สพฺเพ. \\ อนิสฺสิโต เฉตฺว สิเนหโทสํ, \\ เอโก จเร ขคฺควิสาณ\-กปฺโปติ.}}\csromangathastanza{\csromangathawak{\csromangathaitemline{153}{\textbf{วิปิฏฺฐิ}\-\textbf{กตฺวาน สุขํ ทุขญฺจ,}} \\ \csromangathacontline{\textbf{ปุพฺเพว จ โสมนสฺส}\-\textbf{โทมนสฺสํ.}} \\ \csromangathacontline{\textbf{ลทฺธานุ}\-\textbf{เปกฺขํ สมถํ วิสุทฺธํ,}} \\ \csromangathacontline{\textbf{เอโก จเร ขคฺควิสาณ}\-\textbf{กปฺโป.}}}}}

\prose{\textbf{วิปิฏฺฐิ}\-\textbf{กตฺวาน สุขํ ทุขญฺจ, ปุพฺเพว จ โสมนสฺส}\-\textbf{โทมนสฺสนฺ}ติ โส ปจฺเจก\-สมฺพุทฺโธ สุขสฺส จ ปหานา ทุกฺขสฺส จ ปหานา ปุพฺเพว โสมนสฺส\-โทมนสฺสานํ อตฺถงฺคมา อทุกฺขมสุขํ อุเปกฺขา\-สติ\-ปาริสุทฺธิํ จตุตฺถํ ฌานํ อุปสมฺปชฺช วิหรตีติ วิปิฏฺฐิ\-กตฺวาน สุขํ ทุขญฺจ, ปุพฺเพว จ โสมนสฺส\-โทมนสฺสํ.}

\prose{\csromanfolio{296}\textbf{ลทฺธานุ}\-\textbf{เปกฺขํ สมถํ วิสุทฺธ}นฺติ \textbf{อุเปกฺขา}ติ ยา จตุตฺถชฺฌาเน อุเปกฺขา อุเปกฺขนา อชฺฌุเปกฺขนา จิตฺตสมตา จิตฺต\-ปฺปสฺสทฺธตา\csromansharedfootnote{084cbcf17760241f}{จิตฺตวิสฏตา (ก) ปสฺส ขุ 7. 402 ปิฏฺเฐ.} มชฺฌตฺตตา จิตฺตสฺส. \textbf{สมโถ}ติ ยา จิตฺตสฺส ฐิติ สณฺฐิติ อวฏฺฐิติ อวิสาหาโร\csromansharedfootnote{c85a598f7a298e95}{อวิสํหาโร (ก) ปสฺส อภิ 1. 19 ปิฏฺเฐ.} อวิกฺเขโป อวิสาหฏ\-มานสตา\csromansharedfootnote{902ca285204c6ebc}{อวิสํหฏมานสตา (ก)} สมโถ สมาธินฺทฺริยํ สมาธิพลํ สมฺมาสมาธิ, จตุตฺถชฺฌาเน อุเปกฺขา จ สมโถ จ สุทฺธา โหนฺติ วิสุทฺธา ปริโยทาตา อนงฺคณา วิคตู\-ปกฺกิเลสา มุทุภูตา กมฺมนิยา ฐิตา อาเนญฺชปฺปตฺตา. \textbf{ลทฺธานุ}\-\textbf{เปกฺขํ สมถํ วิสุทฺธนฺ}ติ จตุตฺถ\-ชฺฌานํ อุเปกฺขญฺจ สมถญฺจ ลทฺธา ลภิตฺวา วินฺทิตฺวา ปฏิลภิตฺวาติ ลทฺธานุ\-เปกฺขํ สมถํ วิสุทฺธํ, เอโก จเร ขคฺควิสาณ\-กปฺโป. เตนาห โส ปจฺเจก\-สมฺพุทฺโธ\csromandash{}}

\setlength{\csromangathapagewidth}{0pt}
\setlength{\csromangathawakwidth}{0pt}
\csromangathameasuregroup{}{วิปิฏฺฐิกตฺวาน สุขํ ทุขญฺจ, \\ ปุพฺเพว จ โสมนสฺสโทมนสฺสํ.}
\csromangathameasurewak{วิปิฏฺฐิกตฺวาน สุขํ ทุขญฺจ, \\ ปุพฺเพว จ โสมนสฺสโทมนสฺสํ.}
\setlength{\csromangathaleftcol}{0pt}
\csromangathagroup{\csromangathastanza{\csromangathawak{วิปิฏฺฐิ\-กตฺวาน สุขํ ทุขญฺจ, \\ ปุพฺเพว จ โสมนสฺส\-โทมนสฺสํ.}}}

\prose{ลทฺธานุ\-เปกฺขํ สมถํ วิสุทฺธํ,}

\prose{เอโก จเร ขคฺควิสาณ\-กปฺโปติ.}

\proseitem{154}{\textbf{อารทฺธวิริโย ปรมตฺถ}\-\textbf{ปตฺติยา,}}

\prose{\textbf{อลีนจิตฺโต อกุสีตวุตฺติ.}}

\prose{\textbf{ทฬฺหนิกฺกโม ถาม}\-\textbf{พลู}\-\textbf{ปปนฺโน,}}

\prose{\textbf{เอโก จเร ขคฺควิสาณ}\-\textbf{กปฺโป.(3)}}

\prose{\textbf{อารทฺธวิริโย ปรมตฺตปตฺติยา}ติ \textbf{ปรมตฺถํ} วุจฺจติ อมตํ นิพฺพานํ. โย โส สพฺพ\-สงฺขาร\-สมโถ สพฺพู\-ปธิ\-ปฏินิสฺสคฺโค ตณฺหกฺขโย วิราโค นิโรโธ นิพฺพานํ. ปรมตฺถสฺส ปตฺติยา ลาภาย ปฏิลาภาย อธิคมาย ผสฺสนาย สจฺฉิกิริยาย อารทฺธวีริโย วิหรติ, อกุสลานํ ธมฺมานํ ปหานาย กุสลานํ ธมฺมานํ สมฺปทาย ถามวา\csromansharedfootnote{6a0d768f35360f84}{ถามสา (ก)} ทฬฺหปรกฺกโม อนิกฺขิตฺตธุโร กุสเลสุ ธมฺเมสูติ อารทฺธวิริโย ปรมตฺถ\-ปตฺติยา.}

\prose{\textbf{อลีนจิตฺโต อกุสีตวุตฺตี}ติ โส ปจฺเจก\-สมฺพุทฺโธ อนุปฺปนฺนานํ ปาปกานํ อกุสลานํ ธมฺมานํ อนุปฺปาทาย ฉนฺทํ ชเนติ วายมติ, วีริยํ อารภติ, จิตฺตํ ปคฺคณฺหาติ ปทหติ. อุปฺปนฺนานํ ปาปกานํ อกุสลานํ ธมฺมานํ ปหานาย ฯเปฯ. อนุปฺปนฺนานํ กุสลานํ ธมฺมานํ อุปฺปาทาย ฯเปฯ. อุปฺปนฺนานํ}

\prose{\csromanfolio{297}กุสลานํ ธมฺมานํ ฐิติยา อสมฺโมสาย ภิยฺโยภาวาย เวปุลฺลาย ภาวนาย ปาริปูริยา\csromansharedfootnote{6ce35d8e98d5b640}{ภาวนาปาริปูริยา (สฺยา)} ฉนฺทํ ชเนติ วายมติ, วีริยํ อารภติ, จิตฺตํ ปคฺคณฺหาติ ปทหตีติ, เอวํ อลีนจิตฺโต อกุสีตวุตฺติ.}

\prose{อถ วา “กามํ ตโจ จ นฺหารุ จ อฏฺฐิ จ อวสิสฺสตุ\csromansharedfootnote{423c17b91ee0fb05}{อวสุสฺสตุ (สฺยา) ม 2. 146 ปิฏฺเฐ; สํ 1. 267 ปิฏฺเฐ จ ปสฺสิตพฺพํ.}, สรีเร อุปสุสฺสตุ มํสโลหิตํ. ยํ ตํ ปุริสถาเมน ปุริสพเลน ปุริสวีริเยน ปุริส\-ปรกฺกเมน ปตฺตพฺพํ, น ตํ อปาปุณิตฺวา วีริยสฺส สณฺฐานํ\csromansharedfootnote{22677d555ae706e2}{วิริยสฺส ฐานํ (สฺยา)} ภวิสฺสตี”ติ จิตฺตํ ปคฺคณฺหาติ ปทหติ, เอวมฺปิ อลีนจิตฺโต อกุสีตวุตฺติ.}

\setlength{\csromangathapagewidth}{0pt}
\setlength{\csromangathawakwidth}{0pt}
\csromangathameasuregroup{“นาสิสฺสํ น ปิวิสฺสามิ, \\ นปิ ปสฺสํ นิปาเตสฺสํ,}{\csromangathabat{“นาสิสฺสํ น ปิวิสฺสามิ,}{วิหารโต น นิกฺขเม.} \\ \csromangathabat{นปิ ปสฺสํ นิปาเตสฺสํ,}{ตณฺหาสลฺเล อนูหเต”ติ.}}
\csromangathasetleft{“นาสิสฺสํ น ปิวิสฺสามิ, \\ นปิ ปสฺสํ นิปาเตสฺสํ,}
\csromangathagroup{\csromangathastanza{\csromangathabat{“นาสิสฺสํ น ปิวิสฺสามิ,}{วิหารโต น นิกฺขเม.} \\ \csromangathabat{นปิ ปสฺสํ นิปาเตสฺสํ,}{ตณฺหาสลฺเล อนูหเต”ติ.}}}

\prose{จิตฺตํ ปคฺคณฺหาติ ปทหติ, เอวมฺปิ อลีนจิตฺโต อกุสีตวุตฺติ.}

\prose{“น ตาวาหํ อิมํ ปลฺลงฺกํ ภินฺทิสฺสามิ, ยาว เม น อนุปาทาย อาสเวหิ จิตฺตํ วิมุจฺจิสฺสตี”ติ จิตฺตํ ปคฺคณฺหาติ ปทหติ, เอวมฺปิ อลีนจิตฺโต อกุสีตวุตฺติ.อ ตาวาหํ อิมมฺหา อาสนา วุฏฺฐหิสฺสามิ, ยาว เม น อนุปาทาย อาสเวหิ จิตฺตํ วิมุจฺจิสฺสาตี”ติ จิตฺตํ ปคฺคณฺหาติ ปทหติ, เอวมฺปิ อลีนจิตฺโต อกุสีตวุตฺติ.}

\prose{“น ตาวาหํ อิมมฺหา จงฺกมา โอโรหิสฺสามิ. วิหารา นิกฺขมิสฺสามิ. อฑฺฒโยคา นิกฺขมิสฺสามิ. ปาสาทา นิกฺขมิสฺสามิ. หมฺมิยา นิกฺขมิสฺสามิ. คุหาย นิกฺขมิสฺสามิ. เลณา นิกฺขมิสฺสามิ. กุฏิยา นิกฺขมิสฺสามิ. กูฏาคารา นิกฺขมิสฺสามิ. อฏฺฏา นิกฺขมิสฺสามิ. มาฬา นิกฺขมิสฺสามิ. อุทฺทณฺฑา นิกฺขมิสฺสามิ. อุปฏฺฐาน\-สาลาย นิกฺขมิสฺสามิ. มณฺฑปา นิกฺขมิสฺสามิ. รุกฺขมูลา นิกฺขมิสฺสามิ, ยาว เม น อนุปาทาย อาสเวหิ จิตฺตํ วิมุจฺจิสฺสตี”ติ จิตฺตํ ปคฺคณฺหาติ ปทหติ, เอวมฺปิ อลีนจิตฺโต อกุสีตวุตฺติ.}

\prose{“อิมสฺมิํ เยว\csromansharedfootnote{83a75885567999a3}{อิมสฺมิํเยว อาคมสนฺธิ \csromandash{} ม.พ.ป.} ปุพฺพณฺหสมเย\csromansharedfootnote{a3b16a4159439bbb}{ปุพฺพณฺหสมยํ (ก) ขุ 7. 51 ปิฏฺเฐปิ.} อริยธมฺมํ อาหริสฺสามิ สมาหริสฺสามิ อธิคจฺฉิสฺสามิ ผสฺสยิสฺสามิ สจฺฉิกริสฺสามีติ จิตฺตํ ปคฺคณฺหาติ ปทหติ, เอวมฺปิ อลีนจิตฺโต อกุสีตวุตฺติ. อิมสฺมิํ เยว\csromansharedfootnote{83a75885567999a3}{อิมสฺมิํเยว อาคมสนฺธิ \csromandash{} ม.พ.ป.} มชฺฌนฺหิก\-สมเย ฯเปฯ.}

\prose{\csromanfolio{298}สายนฺหสมเย. ปุเรภตฺตํ. ปจฺฉาภตฺตํ. ปุริมยามํ. มชฺฌิมยามํ. ปจฺฉิมยามํ. กาเฬ. ชุเณฺห. วสฺเส. เหมนฺเต. คิเมฺห. ปุริเม วโยขนฺเธ. มชฺฌิเม วโยขนฺเธ. ปจฺฉิเม วโยขนฺเธ อริยธมฺมํ อาหริสฺสามิ สมาหริสฺสามิ อธิคจฺฉิสฺสามิ ผสฺสยิสฺสามิ สจฺฉิกริสฺสามี”ติ จิตฺตํ ปคฺคณฺหาติ ปทหติ, เอวมฺปิ อลีนจิตฺโต อกุสีตวุตฺติ.}

\prose{\textbf{ทฬฺหนิกฺกโม ถาม}\-\textbf{พลู}\-\textbf{ปปนฺโน}ติ โส ปจฺเจก\-สมฺพุทฺโธ ทฬฺหสมาทาโน อโหสิ กุสเลสุ ธมฺเมสุ อวฏฺฐิต\-สมาทาโน กายสุจริเต วจีสุจริเต มโนสุจริเต ทานสํวิภาเค สีลสมาทาเน อุโปสถุปวาเส มตฺเตยฺยตาย\csromansharedfootnote{4396bb32ef25e84a}{เมตฺเตยฺยตาย (สฺยา, ก) โมคฺคลฺลานพฺยากรเณ 4. 39 สุตฺตํ โอโลเกตพฺพํ.} เปตฺเตยฺยตาย สามญฺญตาย พฺรหฺมญฺญตาย กุเล\-เชฏฺฐา\-ปจายิ\-ตาย อญฺญตร\-ญฺญตเรสุ อธิกุสเลสุ ธมฺเมสูติ ทฬฺหนิกฺกโม. \textbf{ถาม}\-\textbf{พลู}\-\textbf{ปปนฺโน}ติ โส ปจฺเจก\-สมฺพุทฺโธ ถาเมน จ พเลน จ วีริเยน จ ปรกฺกเมน จ ปญฺญาย จ อุเปโต โหติ สมุเปโต อุปาคโต สมุปาคโต อุปปนฺโน สมุปปนฺโน สมนฺนาคโตติ ทฬฺหนิกฺกโม ถาม\-พลู\-ปปนฺโน, เอโก จเร ขคฺควิสาณ\-กปฺโป. เตนาห โส ปจฺเจก\-สมฺพุทฺโธ\csromandash{}}

\setlength{\csromangathapagewidth}{0pt}
\setlength{\csromangathawakwidth}{0pt}
\csromangathameasuregroup{}{อารทฺธวิริโย ปรมตฺถปตฺติยา, \\ อลีนจิตฺโต อกุสีตวุตฺติ.}
\csromangathameasurewak{อารทฺธวิริโย ปรมตฺถปตฺติยา, \\ อลีนจิตฺโต อกุสีตวุตฺติ.}
\setlength{\csromangathaleftcol}{0pt}
\csromangathagroup{\csromangathastanza{\csromangathawak{อารทฺธวิริโย ปรมตฺถ\-ปตฺติยา, \\ อลีนจิตฺโต อกุสีตวุตฺติ.}}}

\prose{ทฬฺหนิกฺกโม ถาม\-พลู\-ปปนฺโน,}

\prose{เอโก จเร ขคฺควิสาณ\-กปฺโปติ.}

\setlength{\csromangathapagewidth}{0pt}
\setlength{\csromangathawakwidth}{0pt}
\csromangathameasuregroup{}{\textbf{ปฏิสลฺลานํ ฌานม}\textbf{ริญฺจมาโน,} \\ \textbf{ธมฺเมสุ นิจฺจํ อนุธมฺมจารี.} \\ \textbf{อาทีนวํ สมฺมสิตา ภเวสุ,} \\ \textbf{เอโก จเร ขคฺควิสาณ}\textbf{กปฺโป.}}
\csromangathameasurewak{\textbf{ปฏิสลฺลานํ ฌานม}\textbf{ริญฺจมาโน,} \\ \textbf{ธมฺเมสุ นิจฺจํ อนุธมฺมจารี.} \\ \textbf{อาทีนวํ สมฺมสิตา ภเวสุ,} \\ \textbf{เอโก จเร ขคฺควิสาณ}\textbf{กปฺโป.}}
\setlength{\csromangathaleftcol}{0pt}
\csromangathagroup{\csromangathastanza{\csromangathawak{\csromangathaitemline{155}{\textbf{ปฏิสลฺลานํ ฌานม}\-\textbf{ริญฺจมาโน,}} \\ \csromangathacontline{\textbf{ธมฺเมสุ นิจฺจํ อนุธมฺมจารี.}} \\ \csromangathacontline{\textbf{อาทีนวํ สมฺมสิตา ภเวสุ,}} \\ \csromangathacontline{\textbf{เอโก จเร ขคฺควิสาณ}\-\textbf{กปฺโป.}}}}}

\prose{\textbf{ปฏิสลฺลานํ ฌานม}\-\textbf{ริญฺจมาโน}ติ โส ปจฺเจก\-สมฺพุทฺโธ ปฏิสลฺลานา\-ราโม โหติ ปฏิสลฺลาน\-รโต อชฺฌตฺตํ เจโตสมถม\-นุยุตฺโต อนิรากต\-ชฺฌาโน วิปสฺสนาย สมนฺนาคโต พฺรูเหตา สุญฺญาคารํ ฌายี ฌานรโต เอกตฺตม\-นุยุตฺโต สทตฺถครุโกติ ปฏิสลฺลานํ. \textbf{ฌานม}\-\textbf{ริญฺจมาโน}ติ โส ปจฺเจก\-สมฺพุทฺโธ ทฺวีหิ การเณหิ ฌานํ น ริญฺจติ อนุปฺปนฺนสฺส วา ปฐมสฺส ฌานสฺส อุปฺปาทาย ยุตฺโต ปยุตฺโต}

\prose{\csromanfolio{299}สํยุตฺโต อายุตฺโต สมายุตฺโต, อนุปฺปนฺนสฺส วา ทุติยสฺส ฌานสฺส. อนุปฺปนฺนสฺส วา ตติยสฺส ฌานสฺส. อนุปฺปนฺนสฺส วา จตุตฺถสฺส ฌานสฺส อุปฺปาทาย ยุตฺโต ปยุตฺโต สํยุตฺโต อายุตฺโต สมายุตฺโตติ เอวมฺปิ ฌานํ น ริญฺจติ.}

\prose{อถ วา อุปฺปนฺนํ วา ปฐมํ ฌานํ อาเสวติ ภาเวติ พหุลีกโรติ, อุปฺปนฺนํ วา ทุติยํ ฌานํ. อุปฺปนฺนํ วา ตติยํ ฌานํ. อุปฺปนฺนํ วา จตุตฺถํ ฌานํ อาเสวติ ภาเวติ พหุลีกโรติ, เอวมฺปิ ฌานํ น ริญฺจตีติ ปฏิสลฺลานํ ฌานม\-ริญฺจมาโน.}

\prose{\textbf{ธมฺเมสุ นิจฺจํ อนุธมฺมจารี}ติ ธมฺมา วุจฺจนฺติ จตฺตาโร สติปฏฺฐานา ฯเปฯ อริโย อฏฺฐงฺคิโก มคฺโค. กตเม อนุธมฺมา, สมฺมาปฏิปทา อปจฺจนีกปฏิปทา อนฺวตฺถ\-ปฏิปทา ธมฺมานุธมฺม\-ปฏิปทา สีเลสุ ปริปูรการิตา อินฺทฺริเยสุ คุตฺตทฺวารตา โภชเน มตฺตญฺญุตา ชาคริยานุโยโค สติสมฺปชญฺญํ, อิเม วุจฺจนฺติ อนุธมฺมา. \textbf{ธมฺเมสุ นิจฺจํ อนุธมฺมจารี}ติ ธมฺเมสุ นิจฺจกาลํ ธุวกาลํ สตตํ สมิตํ อโวกิณฺณํ โปงฺขานุโปงฺขํ อุทกูมิกชาตํ\csromansharedfootnote{446cbb3de7255fd3}{อุทกุมฺมิชาตํ (สฺยา), อุทกุมฺมิกชาตํ (ก) 43 ปิฏฺเฐปิ.} อวีจิ สนฺตติ สหิตํ ผสฺสิตํ ปุเรภตฺตํ ปจฺฉาภตฺตํ ปุริมยามํ มชฺฌิมยามํ ปจฺฉิมยามํ กาเฬ ชุเณฺห วสฺเส เหมนฺเต คิเมฺห ปุริเม วโยขนฺเธ มชฺฌิเม วโยขนฺเธ ปจฺฉิเม วโยขนฺเธ จรติ วิหรติ อิริยติ วตฺเตติ ปาเลติ ยเปติ ยาเปตีติ ธมฺเมสุ นิจฺจํ อนุธมฺมจารี.}

\prose{\textbf{อาทีนวํ สมฺมสิตา ภเวสู}ติ “สพฺเพ สงฺขารา อนิจฺจา”ติ อาทีนวํ สมฺมสิตา ภเวสุ, “สพฺเพ สงฺขารา ทุกฺขา”ติ. “สพฺเพ ธมฺมา อนตฺตา”ติ ฯเปฯ. “ยํ กิญฺจิ สมุทย\-ธมฺมํ, สพฺพํ ตํ นิโรธธมฺมนฺ”ติ อาทีนวํ สมฺมสิตา ภเวสูติ อาทีนวํ สมฺมสิตา ภเวสุ, เอโก จเร ขคฺควิสาณ\-กปฺโป. เตนาห โส ปจฺเจก\-สมฺพุทฺโธ\csromandash{}}

\setlength{\csromangathapagewidth}{0pt}
\setlength{\csromangathawakwidth}{0pt}
\csromangathameasuregroup{}{ปฏิสลฺลานํ ฌานมริญฺจมาโน,}
\csromangathameasurewak{ปฏิสลฺลานํ ฌานมริญฺจมาโน,}
\setlength{\csromangathaleftcol}{0pt}
\csromangathagroup{\csromangathastanza{\csromangathawak{ปฏิสลฺลานํ ฌานม\-ริญฺจมาโน,}}}

\vspace{\tipitakaparskip}
\csromancenter{ธมฺเมสุ นิจฺจํ อนุธมฺมจารี.}
\setlength{\csromangathapagewidth}{0pt}
\setlength{\csromangathawakwidth}{0pt}
\csromangathameasuregroup{}{อาทีนวํ สมฺมสิตา ภเวสุ, \\ เอโก จเร ขคฺควิสาณกปฺโปติ. \\ \textbf{ตณฺหกฺขยํ ปตฺถยม}\textbf{ปฺปมตฺโต,} \\ \textbf{อเนฬมูโค}\textsuperscript{1}\textbf{ สุตวา สตีมา.} \\ \textbf{สงฺขาตธมฺโม นิยโต ปธานวา,} \\ \textbf{เอโก จเร ขคฺควิสาณ}\textbf{กปฺโป.}}
\csromangathameasurewak{อาทีนวํ สมฺมสิตา ภเวสุ, \\ เอโก จเร ขคฺควิสาณกปฺโปติ. \\ \textbf{ตณฺหกฺขยํ ปตฺถยม}\textbf{ปฺปมตฺโต,} \\ \textbf{อเนฬมูโค}\textsuperscript{1}\textbf{ สุตวา สตีมา.} \\ \textbf{สงฺขาตธมฺโม นิยโต ปธานวา,} \\ \textbf{เอโก จเร ขคฺควิสาณ}\textbf{กปฺโป.}}
\setlength{\csromangathaleftcol}{0pt}
\csromangathagroup{\csromangathastanza{\csromangathawak{อาทีนวํ สมฺมสิตา ภเวสุ, \\ เอโก จเร ขคฺควิสาณ\-กปฺโปติ.}}\csromangathastanza{\csromangathawak{\csromanfolio{300}\csromangathaitemline{156}{\textbf{ตณฺหกฺขยํ ปตฺถยม}\-\textbf{ปฺปมตฺโต,}} \\ \csromangathacontline{\textbf{อเนฬมูโค}\csromansharedfootnote{ad3734ca1deda7f1}{อเนลมูโค (สฺยา, ก)}\textbf{ สุตวา สตีมา.}} \\ \csromangathacontline{\textbf{สงฺขาตธมฺโม นิยโต ปธานวา,}} \\ \csromangathacontline{\textbf{เอโก จเร ขคฺควิสาณ}\-\textbf{กปฺโป.}}}}}

\prose{\textbf{ตณฺหกฺขยํ ปตฺถยม}\-\textbf{ปฺปมตฺโต}ติ ตณฺหาติ รูป\textbf{ตณฺหา} ฯเปฯ ธมฺมตณฺหา. \textbf{ตณฺหกฺขยนฺ}ติ ราคกฺขยํ โทสกฺขยํ โมหกฺขยํ คติกฺขยํ อุปปตฺติ\-กฺขยํ ปฏิสนฺธิ\-กฺขยํ ภวกฺขยํ สํสารกฺขยํ วฏฺฏกฺขยํ ปตฺถยนฺโต อิจฺฉนฺโต สาทิยนฺโต ปิหยนฺโต อภิชปฺปนฺโตติ ตณฺหกฺขยํ ปตฺถยํ. \textbf{อปฺปมตฺโต}ติ โส ปจฺเจก\-สมฺพุทฺโธ สกฺกจฺจการี สาตจฺจการี ฯเปฯ อปฺปมาโท กุสเลสุ ธมฺเมสูติ ตณฺหกฺขยํ ปตฺถยม\-ปฺปมตฺโต.}

\prose{\textbf{อเนฬมูโค สุตวา สตีมา}ติ \textbf{อเนฬมูโค}ติ โส ปจฺเจก\-สมฺพุทฺโธ ปณฺฑิโต ปญฺญวา พุทฺธิมา ญาณี วิภาวี เมธาวี. \textbf{สุตวา}ติ โส ปจฺเจก\-สมฺพุทฺโธ พหุสฺสุโต โหติ สุตธโร สุตสนฺนิจฺจโย, เย เต ธมฺมา อาทิกลฺยาณา มชฺเฌกลฺยาณา ปริโยสาน\-กลฺยาณา สาตฺถํ สพฺยญฺชนํ เกวล\-ปริปุณฺณํ ปริสุทฺธํ พฺรหฺมจริยํ อภิวทนฺติ, ตถารูปาสฺส ธมฺมา พหุสฺสุตา โหนฺติ ธาตา วจสา ปริจิตา มนสา\-นุเปกฺขิตา ทิฏฺฐิยา สุปฺปฏิวิทฺธา. \textbf{สตีมา}ติ โส ปจฺเจก\-สมฺพุทฺโธ สติมา โหติ ปรเมน สติเนปกฺเกน สมนฺนาคตตฺตา จิรกตมฺปิ จิรภาสิตมฺปิ สริตา อนุสฺสริตาติ อเนฬมูโค สุตวา สตีมา.}

\prose{\textbf{สงฺขาตธมฺโม นิยโต ปธานวา}ติ สงฺขาตธมฺโม วุจฺจติ ญาณํ. ยา ปญฺญา ปชานนา ฯเปฯ อโมโห ธมฺมวิจโย สมฺมาทิฏฺฐิ. \textbf{สงฺขาต}\-\textbf{ธมฺโม}ติ โส ปจฺเจก\-สมฺพุทฺโธ สงฺขาตธมฺโม ญาตธมฺโม ตุลิตธมฺโม ตีริตธมฺโม วิภูตธมฺโม วิภาวิต\-ธมฺโม “สพฺเพ สงฺขารา อนิจฺจา”ติ สงฺขาตธมฺโม ฯเปฯ. “ยํ กิญฺจิ สมุทย\-ธมฺมํ, สพฺพํ ตํ นิโรธธมฺมนฺ”ติ สงฺขาตธมฺโม ญาตธมฺโม ตุลิตธมฺโม ตีริตธมฺโม วิภูตธมฺโม วิภาวิต\-ธมฺโม. อถ วา ตสฺส ปจฺเจก\-สมฺพุทฺธสฺส จ ขนฺธา สํขิตฺตา, ธาตุโย สํขิตฺตา, อายตนานิ สํขิตฺตานิ, คติโย สํขิตฺตา, อุปปตฺติโย สํขิตฺตา, ปฏิสนฺธิโย สํขิตฺตา, ภวา สํขิตฺตา, สํสารา สํขิตฺตา, วฏฺฏา สํขิตฺตา. อถ วา โส ปจฺเจก\-สมฺพุทฺโธ ขนฺธ\-ปริยนฺเต ฐิโต,}

\prose{\csromanfolio{301}ธาตุปริยนฺเต ฐิโต, อายตน\-ปริยนฺเต ฐิโต, คติปริยนฺเต ฐิโต, อุปปตฺติ\-ปริยนฺเต ฐิโต, ปฏิสนฺธิ\-ปริยนฺเต ฐิโต, ภวปริยนฺเต ฐิโต, สํสาร\-ปริยนฺเต ฐิโต, วฏฺฏปริยนฺเต ฐิโต, สงฺขาร\-ปริยนฺเต ฐิโต, อนฺติมภเว\csromansharedfootnote{59be0c17f76a575d}{อนฺติเม ภเว (ก) 86 ปิฏฺเฐปิ.} ฐิโต, อนฺติม\-สมุสฺสเย ฐิโต อนฺติม\-เทห\-ธโร ปจฺเจก\-สมฺพุทฺโธ.}

\setlength{\csromangathapagewidth}{0pt}
\setlength{\csromangathawakwidth}{0pt}
\csromangathameasuregroup{ตสฺสายํ ปจฺฉิมโก ภโว, \\ ชาติมรณสํสาโร\textsuperscript{1},}{\csromangathabat{ตสฺสายํ ปจฺฉิมโก ภโว,}{จริโมยํ สมุสฺสโย.} \\ \csromangathabat{ชาติมรณสํสาโร\textsuperscript{1},}{นตฺถิ ตสฺส ปุนพฺภโวติ.}}
\csromangathasetleft{ตสฺสายํ ปจฺฉิมโก ภโว, \\ ชาติมรณสํสาโร\textsuperscript{1},}
\csromangathagroup{\csromangathastanza{\csromangathabat{ตสฺสายํ ปจฺฉิมโก ภโว,}{จริโมยํ สมุสฺสโย.} \\ \csromangathabat{ชาติ\-มรณ\-สํสาโร\csromansharedfootnote{672b5ef60e0ba722}{ชาติชรามรณสํสาโร (สฺยา) เอวมีทิเสสุ ฐาเนสุ.},}{นตฺถิ ตสฺส ปุนพฺภโวติ.}}}

\prose{ตํการณา ปจฺเจก\-สมฺพุทฺโธ สงฺขาตธมฺโม. \textbf{นิยโต}ติ นิยามา วุจฺจนฺติ จตฺตาโร อริยมคฺคา. จตูหิ อริยมคฺเคหิ สมนฺนาคโตติ นิยโต, นิยามํ ปตฺโต สมฺปตฺโต อธิคโต ผสฺสิโต สจฺฉิกโต ปตฺโต นิยามํ. \textbf{ปธานวา}ติ ปธานํ วุจฺจติ วีริยํ, โส เจตโส วีริยารมฺโภ นิกฺกโม ปรกฺกโม อุยฺยาโม วายาโม อุสฺสาโห อุสฺโสฬฺหี ถาโม ธิติ อสิถิล\-ปรกฺกโม อนิกฺขิตฺตจฺฉนฺทตา อนิกฺขิตฺตธุรตา ธุรสมฺปคฺคาโห วีริยํ วีริยินฺทฺริยํ วีริยพลํ สมฺมาวายาโม. โส ปจฺเจก\-สมฺพุทฺโธ อิมินา ปธาเนน อุเปโต สมุเปโต อุปาคโต สมุปาคโต อุปปนฺโน สมุปปนฺโน สมนฺนาคโต, ตสฺมา โส ปจฺเจก\-สมฺพุทฺโธ ปธานวาติ สงฺขาตธมฺโม นิยโต ปธานวา, เอโก จเร ขคฺควิสาณ\-กปฺโป. เตนาห โส ปจฺเจก\-สมฺพุทฺโธ\csromandash{}}

\setlength{\csromangathapagewidth}{0pt}
\setlength{\csromangathawakwidth}{0pt}
\csromangathameasuregroup{}{ตณฺหกฺขยํ ปตฺถยมปฺปมตฺโต, \\ อเนฬมูโค สุตวา สตีมา. \\ สงฺขาตธมฺโม นิยโต ปธานวา, \\ เอโก จเร ขคฺควิสาณกปฺโปติ. \\ \textbf{สีโหว สทฺเทสุ อสนฺตสนฺโต,} \\ \textbf{วาโตว ชาลมฺหิ อสชฺชมาโน.} \\ \textbf{ปทุมํว โตเยน อลิมฺปมาโน}\textsuperscript{1}\textbf{,} \\ \textbf{เอโก จเร ขคฺควิสาณ}\textbf{กปฺโป.}}
\csromangathameasurewak{ตณฺหกฺขยํ ปตฺถยมปฺปมตฺโต, \\ อเนฬมูโค สุตวา สตีมา. \\ สงฺขาตธมฺโม นิยโต ปธานวา, \\ เอโก จเร ขคฺควิสาณกปฺโปติ. \\ \textbf{สีโหว สทฺเทสุ อสนฺตสนฺโต,} \\ \textbf{วาโตว ชาลมฺหิ อสชฺชมาโน.} \\ \textbf{ปทุมํว โตเยน อลิมฺปมาโน}\textsuperscript{1}\textbf{,} \\ \textbf{เอโก จเร ขคฺควิสาณ}\textbf{กปฺโป.}}
\setlength{\csromangathaleftcol}{0pt}
\csromangathagroup{\csromangathastanza{\csromangathawak{ตณฺหกฺขยํ ปตฺถยม\-ปฺปมตฺโต, \\ อเนฬมูโค สุตวา สตีมา. \\ สงฺขาตธมฺโม นิยโต ปธานวา, \\ เอโก จเร ขคฺควิสาณ\-กปฺโปติ.}}\csromangathastanza{\csromangathawak{\csromangathaitemline{157}{\textbf{สีโหว สทฺเทสุ อสนฺตสนฺโต,}} \\ \csromangathacontline{\textbf{วาโตว ชาลมฺหิ อสชฺชมาโน.}} \\ \csromangathacontline{\textbf{ปทุมํว โตเยน อลิมฺปมาโน}\csromansharedfootnote{372e87aa7a4a52e1}{อลิปฺปมาโน ขุ 1. 290 ปิฏฺเฐ.}\textbf{,}} \\ \csromangathacontline{\textbf{เอโก จเร ขคฺควิสาณ}\-\textbf{กปฺโป.}}}}}

\prose{\textbf{สีโหว สทฺเทสุ อสนฺตสนฺโต}ติ ยถา สีโห มิคราชา สทฺเทสุ อสนฺตาสี อปริสนฺตาสี อนุตฺราสี อนุพฺพิคฺโค อนุสฺสงฺกี\csromansharedfootnote{19bec89302a7f050}{อนุสฺสุกี (สฺยา)} \csromanfolio{302}อนุตฺราโส อภีรู อจฺฉมฺภี อนุตฺราสี อปลายี. ปจฺเจก\-สมฺพุทฺโธปิ สทฺเทสุ อสนฺตาสี อปริสนฺตาสี อนุตฺราสี อนุพฺพิคฺโค อนุสฺสงฺกี อนุตฺราโส อภีรู อจฺฉมฺภี อนุตฺราสี อปลายี ปหีน\-ภยเภรโว วิคต\-โลมหํโส วิหรตีติ สีโหว สทฺเทสุ อสนฺตสนฺโต.}

\prose{\textbf{วาโตว ชาลมฺหิ อสชฺชมาโน}ติ \textbf{วาโต}ติ ปุรตฺถิมาวาตา ปจฺฉิมา วาตา อุตฺตารา วาตา ทกฺขิณา วาตา สรชา วาตา อรชา วาตา สีตา วาตา อุณฺหา วาตา ปริตฺตา วาตา อธิมตฺตา วาตา เวรมฺภวาตา ปกฺขวาตา\csromansharedfootnote{36cc28841703bd54}{ปกฺขิวาตา (สฺยา) 109, 266 ปิฏฺเฐ, นตฺถิ ปาฐนานตฺตํ.} สุปณฺณวาตา ตาลปณฺณวาตา วิธูปนวาตา. ชาลํ วุจฺจติ สุตฺตชาลํ. ยถา วาโต ชาลมฺหิ น สชฺชติ น คณฺหาติ น พชฺฌติ น ปลิพชฺฌติ, เอวเมว เทฺว ชาลา ตณฺหาชาลญฺจ ทิฏฺฐิชาลญฺจ ฯเปฯ อิทํ ตณฺหาชาลํ ฯเปฯ อิทํ ทิฏฺฐิชาลํ. ตสฺส ปจฺเจก\-สมฺพุทฺธสฺส ตณฺหาชาลํ ปหีนํ, ทิฏฺฐิชาลํ ปฏินิสฺสฏฺฐํ, ตณฺหาชาลสฺส ปหีนตฺตา ทิฏฺฐิชาลสฺส ปฏินิสฺสฏฺฐ\-ตฺตา โส ปจฺเจก\-สมฺพุทฺโธ รูเป น สชฺชติ, สทฺเท น สชฺชติ ฯเปฯ ทิฏฺฐ\-สุต\-มุต\-วิญฺญาตพฺเพสุ ธมฺเมสุ น สชฺชติ น คณฺหาติ น พชฺฌติ น ปลิพชฺฌติ นิกฺขนฺโต นิสฺสโฏ วิปฺปมุตฺโต วิสญฺญุตฺโต วิมริยาทิกเตน เจตสา วิหรตีติ วาโตว ชาลมฺหิ อสชฺชมาโน.}

\prose{\textbf{ปทุมํว โตเยน อลิมฺปมาโน}ติ ปทุมํ วุจฺจติ ปทุมปุปฺผํ. โตยํ วุจฺจติ อุทกํ. ยถา ปทุมปุปฺผํ โตเยน น ลิปฺปติ น ปลิปฺปติ น อุปลิปฺปติ, อลิตฺตํ อปลิตฺตํ อนุปลิตฺตํ, เอวเมว เทฺว เลปา ตณฺหาเลโป จ ทิฏฺฐิเลโป จ ฯเปฯ อยํ ตณฺหาเลโป ฯเปฯ อยํ ทิฏฺฐิเลโป. ตสฺส ปจฺเจก\-สมฺพุทฺธสฺส ตณฺหาเลโป ปหีโน, ทิฏฺฐิเลโป ปฏินิสฺสฏฺโฐ, ตณฺหาเลปสฺส ปหีนตฺตา ทิฏฺฐิเลปสฺส ปฏินิสฺสฏฺฐ\-ตฺตา โส ปจฺเจก\-สมฺพุทฺโธ รูเป น ลิมฺปติ, สทฺเท น ลิมฺปติ ฯเปฯ ทิฏฺฐ\-สุต\-มุต\-วิญฺญาตพฺเพสุ ธมฺเมสุ น ลิมฺปติ น ปลิมฺปติ น อุปลิมฺปติ, อลิตฺโต อปลิตฺโต อนุปลิตฺโต นิกฺขนฺโต นิสฺสโฏ วิปฺปมุตฺโต วิสญฺญุตฺโต วิมริยาทิกเตน เจตสา วิหรตีติ ปทุมํว โตเยน อลิมฺปมาโน, เอโก จเร ขคฺควิสาณ\-กปฺโป. เตนาห โส ปจฺเจก\-สมฺพุทฺโธ\csromandash{}}

\setlength{\csromangathapagewidth}{0pt}
\setlength{\csromangathawakwidth}{0pt}
\csromangathameasuregroup{}{สีโหว สทฺเทสุ อสนฺตสนฺโต, \\ วาโตว ชาลมฺหิ อสชฺชมาโน. \\ ปทุมํว โตเยน อลิมฺปมาโน, \\ เอโก จเร ขคฺควิสาณกปฺโปติ. \\ \textbf{สีโห ยถา ทาฐพลี ปสยฺห,} \\ \textbf{ราชา มิคานํ อภิภุยฺย จารี.} \\ \textbf{เสเวถ ปนฺตานิ เสนาสนานิ,} \\ \textbf{เอโก จเร ขคฺควิสาณ}\textbf{กปฺโป.}}
\csromangathameasurewak{สีโหว สทฺเทสุ อสนฺตสนฺโต, \\ วาโตว ชาลมฺหิ อสชฺชมาโน. \\ ปทุมํว โตเยน อลิมฺปมาโน, \\ เอโก จเร ขคฺควิสาณกปฺโปติ. \\ \textbf{สีโห ยถา ทาฐพลี ปสยฺห,} \\ \textbf{ราชา มิคานํ อภิภุยฺย จารี.} \\ \textbf{เสเวถ ปนฺตานิ เสนาสนานิ,} \\ \textbf{เอโก จเร ขคฺควิสาณ}\textbf{กปฺโป.}}
\setlength{\csromangathaleftcol}{0pt}
\csromangathagroup{\csromangathastanza{\csromangathawak{\csromanfolio{303}สีโหว สทฺเทสุ อสนฺตสนฺโต, \\ วาโตว ชาลมฺหิ อสชฺชมาโน. \\ ปทุมํว โตเยน อลิมฺปมาโน, \\ เอโก จเร ขคฺควิสาณ\-กปฺโปติ.}}\csromangathastanza{\csromangathawak{\csromangathaitemline{158}{\textbf{สีโห ยถา ทาฐพลี ปสยฺห,}} \\ \csromangathacontline{\textbf{ราชา มิคานํ อภิภุยฺย จารี.}} \\ \csromangathacontline{\textbf{เสเวถ ปนฺตานิ เสนาสนานิ,}} \\ \csromangathacontline{\textbf{เอโก จเร ขคฺควิสาณ}\-\textbf{กปฺโป.}}}}}

\prose{\textbf{สีโห ยถา ทาฐพลี ปสยฺห, ราชา มิคานํ อภิภุยฺย จารี}ติ ยถา สีโห มิคราชา ทาฐพลี ทาฐาวุโธ สพฺเพ ติรจฺฉานคเต ปาเณ อภิภุยฺย อภิภวิตฺวา อชฺโฌตฺถริตฺวา ปริยาทิยิตฺวา มทฺทิตฺวา จรติ วิหรติ อิริยติ วตฺเตติ ปาเลติ ยเปติ ยาเปติ, ปจฺเจก\-สมฺพุทฺโธปิ ปญฺญาพลี ปญฺญาวุโธ สพฺพปาณภูเต ปุคฺคเล ปญฺญาย อภิภุยฺย อภิภวิตฺวา อชฺโฌตฺถริตฺวา ปริยาทิยิตฺวา มทฺทิตฺวา จรติ วิหรติ อิริยติ วตฺเตติ ปาเลติ ยเปติ ยาเปตีติ สีโห ยถา ทาฐพลี ปสยฺห, ราชา มิคานํ อภิภุยฺย จารี.}

\prose{\textbf{เสเวถ ปนฺตานิ เสนาสนานี}ติ ยถา สีโห มิคราชา อรญฺญวนม\-ชฺโฌคาเหตฺวา จรติ วิหรติ อิริยติ วตฺเตติ ปาเลติ ยเปติ ยาเปติ, ปจฺเจก\-สมฺพุทฺโธปิ อรญฺญ\-วนปตฺถานิ ปนฺตานิ เสนาสนานิ ปฏิเสวติ อปฺปสทฺทานิ อปฺปนิคฺโฆสานิ วิชนวาตานิ มนุสฺส\-ราหสฺเสยฺยกานิ ปฏิสลฺลาน\-สารุปฺปานิ, โส เอโก คจฺฉติ, เอโก ติฏฺฐติ, เอโก นิสีทติ, เอโก เสยฺยํ กปฺเปติ, เอโก คามํ ปิณฺฑาย ปวิสติ, เอโก ปฏิกฺกมติ, เอโก รโห นิสีทติ, เอโก จงฺกมํ อธิฏฺฐาติ, เอโก จรติ วิหรติ อิริยติ วตฺเตติ ปาเลติ ยเปติ ยาเปตีติ เสเวถ ปนฺตานิ เสนาสนานิ, เอโก จเร ขคฺควิสาณ\-กปฺโป. เตนาห โส ปจฺเจก\-สมฺพุทฺโธ\csromandash{}}

\prose{\textbf{สีโห ยถา ทาฐพลี ปสยฺห,}}

\prose{ราชา มิคานํ อภิภูยฺย จารี.}

\setlength{\csromangathapagewidth}{0pt}
\setlength{\csromangathawakwidth}{0pt}
\csromangathameasuregroup{}{\textbf{เสเวถ ปนฺตานิ เสนาสนานิ,} \\ เอโก จเร ขคฺควิสาณกปฺโปติ. \\ \textbf{เมตฺตํ อุเปกฺขํ กรุณํ วิมุตฺติํ,} \\ \textbf{อาเสวมาโน มุทิตญฺจ กาเล.} \\ \textbf{สพฺเพน โลเกน อวิรุชฺฌมาโน,} \\ \textbf{เอโก จเร ขคฺควิสาณ}\textbf{กปฺโป.}}
\csromangathameasurewak{\textbf{เสเวถ ปนฺตานิ เสนาสนานิ,} \\ เอโก จเร ขคฺควิสาณกปฺโปติ. \\ \textbf{เมตฺตํ อุเปกฺขํ กรุณํ วิมุตฺติํ,} \\ \textbf{อาเสวมาโน มุทิตญฺจ กาเล.} \\ \textbf{สพฺเพน โลเกน อวิรุชฺฌมาโน,} \\ \textbf{เอโก จเร ขคฺควิสาณ}\textbf{กปฺโป.}}
\setlength{\csromangathaleftcol}{0pt}
\csromangathagroup{\csromangathastanza{\csromangathawak{\textbf{เสเวถ ปนฺตานิ เสนาสนานิ,} \\ เอโก จเร ขคฺควิสาณ\-กปฺโปติ.}}\csromangathastanza{\csromangathawak{\csromanfolio{304}\csromangathaitemline{159}{\textbf{เมตฺตํ อุเปกฺขํ กรุณํ วิมุตฺติํ,}} \\ \csromangathacontline{\textbf{อาเสวมาโน มุทิตญฺจ กาเล.}} \\ \csromangathacontline{\textbf{สพฺเพน โลเกน อวิรุชฺฌมาโน,}} \\ \csromangathacontline{\textbf{เอโก จเร ขคฺควิสาณ}\-\textbf{กปฺโป.}}}}}

\prose{\textbf{เมตฺตํ อุเปกฺขํ กรุณํ วิมุตฺติํ, อาเสวมาโน มุทิตญฺจ กาเล}ติ โส ปจฺเจก\-สมฺพุทฺโธ เมตฺตา\-สหคเตน เจตสา เอกํ ทิสํ ผริตฺวา วิหรติ. ตถา ทุติยํ. ตถา ตติยํ. ตถา จตุตฺถํ. อิติ อุทฺธมโธ ติริยํ สพฺพธิ สพฺพตฺตตาย สพฺพาวนฺตํ โลกํ เมตฺตา\-สหคเตน เจตสา วิปุเลน มหคฺคเตน อปฺปมาเณน อเวเรน อพฺยาปชฺเชน ผริตฺวา วิหรติ. กรุณา\-สหคเตน เจตสา ฯเปฯ. มุทิตา\-สหคเตน เจตสา ฯเปฯ. อุเปกฺขา\-สหคเตน เจตสา วิปุเลน มหคฺคเตน อปฺปมาเณน อเวเรน อพฺยาปชฺเชน ผริตฺวา วิหรตีติ เมตฺตํ อุเปกฺขํ กรุณํ วิมุตฺติํ, อาเสวมาโน มุทิตญฺจ กาเล.}

\prose{\textbf{สพฺเพน โลเกน อวิรุชฺฌมาโน}ติ เมตฺตาย ภาวิตตฺตา เย ปุรตฺถิมาย ทิสาย สตฺตา, เต อปฺปฏิกูลา โหนฺติ. เย ปจฺฉิมาย ทิสาย สตฺตา. เย อุตฺตราย ทิสาย สตฺตา. เย ทกฺขิณาย ทิสาย สตฺตา. เย ปุรตฺถิมาย อนุทิสาย สตฺตา. เย ปจฺฉิมาย อนุทิสาย สตฺตา. เย อุตฺตราย อนุทิสาย สตฺตา. เย ทกฺขิณาย อนุทิสาย สตฺตา. เย เหฏฺฐิมาย\csromansharedfootnote{128d47bc7fbc52ab}{อโธคมาย (สฺยา)} ทิสาย สตฺตา. เย อุปริมาย ทิสาย สตฺตา. เย ทสสุ ทิสาสุ สตฺตา, เต อปฺปฏิกูลา โหนฺติ. กรุณาย ภาวิตตฺตา. มุทิตาย ภาวิตตฺตา. อุเปกฺขาย ภาวิตตฺตา เย ปุรตฺถิมาย ทิสาย สตฺตา ฯเปฯ เย ทสสุ ทิสาสุ สตฺตา, เต อปฺปฏิกูลา โหนฺติ. \textbf{สพฺเพน โลเกน อวิรุชฺฌมาโน}ติ สพฺเพน โลเกน อวิรุชฺฌมาโน อปฺปฏิวิรุชฺฌมาโน อนาฆาติยมาโน\csromansharedfootnote{971c5c0c7843df7c}{อฆฏฺฏิยมาโน (สฺยา)} อปฺปฏิหญฺญมาโนติ สพฺเพน โลเกน อวิรุชฺฌมาโน, เอโก จเร ขคฺควิสาณ\-กปฺโป. เตนาห โส ปจฺเจก\-สมฺพุทฺโธ\csromandash{}}

\prose{\textbf{เมตฺตํ อุเปกฺขํ กรุณํ วิมุตฺติํ,}}

\prose{\textbf{อาเสวมาโน มุทิตญฺจ กาเล.}}

\setlength{\csromangathapagewidth}{0pt}
\setlength{\csromangathawakwidth}{0pt}
\csromangathameasuregroup{}{\textbf{สพฺเพน โลเกน อวิรุชฺฌมาโน,} \\ เอโก จเร ขคฺควิสาณกปฺโปติ.}
\csromangathameasurewak{\textbf{สพฺเพน โลเกน อวิรุชฺฌมาโน,} \\ เอโก จเร ขคฺควิสาณกปฺโปติ.}
\setlength{\csromangathaleftcol}{0pt}
\csromangathagroup{\csromangathastanza{\csromangathawak{\textbf{สพฺเพน โลเกน อวิรุชฺฌมาโน,} \\ เอโก จเร ขคฺควิสาณ\-กปฺโปติ.}}}

\proseitem{160}{\csromanfolio{305}\textbf{ราคญฺจ โทสญฺจ ปหาย โมหํ,}}

\prose{\textbf{สนฺทาลยิตฺวาน สญฺโญชนานิ}\csromansharedfootnote{67b8bec3dac4ec71}{สํโยชนานิ (ก)}\textbf{.}}

\prose{\textbf{อสนฺตสํ ชีวิต}\-\textbf{สงฺขยมฺหิ,}}

\prose{\textbf{เอโก จเร ขคฺควิสาณ}\-\textbf{กปฺโป.} (9)}

\prose{\textbf{ราคญฺจ โทสญฺจ ปหาย โมหนฺ}ติ ราโคติ โย \textbf{ราโค} สาราโค ฯเปฯ อภิชฺฌา โลโภ อกุสลมูลํ. \textbf{โทโส}ติ โย จิตฺตสฺส อาฆาโต ฯเปฯ จณฺฑิกฺกํ อสุโรโป อนตฺตมนตา จิตฺตสฺส. \textbf{โมโห}ติ ทุกฺเข อญฺญาณํ ฯเปฯ อวิชฺชาลงฺคี โมโห อกุสลมูลํ. \textbf{ราคญฺจ โทสญฺจ ปหาย โมหนฺ}ติ โส ปจฺเจก\-สมฺพุทฺโธ ราคญฺจ โทสญฺจ โมหญฺจ ปหาย ปชหิตฺวา วิโนเทตฺวา พฺยนฺตีกริตฺวา อนภาวํ คเมตฺวาติ ราคญฺจ โทสญฺจ ปหาย โมหํ.}

\prose{\textbf{สนฺทาลยิตฺวาน สญฺโญชนานี}ติ ทส สญฺโญชนานิ กามราค\-สญฺโญชนํ ปฏิฆ\-สญฺโญชนํ ฯเปฯ อวิชฺชา\-สญฺโญชนํ. \textbf{สนฺทาลยิตฺวาน สญฺโญชนานี}ติ ทส สญฺโญชนานิ สนฺทาลยิตฺวา ปทาลยิตฺวา สมฺปทาลยิตฺวา ปชหิตฺวา วิโนเทตฺวา พฺยนฺตีกริตฺวา อนภาวํ คเมตฺวาติ สนฺทาลยิตฺวาน สญฺโญชนานิ.}

\prose{\textbf{อสนฺตสํ ชีวิตสงฺขยมฺหี}ติ โส ปจฺเจก\-สมฺพุทฺโธ ชีวิต\-ปริโยสาเน อสนฺตาสี อนุตฺราสี อนุพฺพิคฺโค อนุสฺสงฺกี อนุตฺราโส อภีรู อจฺฉมฺภี อนุตฺราสี อปลายี ปหีน\-ภยเภรโว วิคต\-โลมหํโส วิหรตีติ อสนฺตสํ ชีวิต\-สงฺขยมฺหิ, เอโก จเร ขคฺควิสาณ\-กปฺโป. เตนาห โส ปจฺเจก\-สมฺพุทฺโธ\csromandash{}}

\prose{\textbf{ราคญฺจ โทสญฺจ ปหาย โมหํ,}}

\prose{\textbf{สนฺทาลยิตฺวาน สญฺโญชนานิ.}}

\setlength{\csromangathapagewidth}{0pt}
\setlength{\csromangathawakwidth}{0pt}
\csromangathameasuregroup{}{\textbf{อสนฺตสํ ชีวิต}\textbf{สงฺขยมฺหิ,} \\ เอโก จเร ขคฺควิสาณกปฺโปติ.}
\csromangathameasurewak{\textbf{อสนฺตสํ ชีวิต}\textbf{สงฺขยมฺหิ,} \\ เอโก จเร ขคฺควิสาณกปฺโปติ.}
\setlength{\csromangathaleftcol}{0pt}
\csromangathagroup{\csromangathastanza{\csromangathawak{\textbf{อสนฺตสํ ชีวิต}\-\textbf{สงฺขยมฺหิ,} \\ เอโก จเร ขคฺควิสาณ\-กปฺโปติ.}}}

\proseitem{161}{\textbf{ภชนฺติ เสวนฺติ จ การณตฺถา,}}

\setlength{\csromangathapagewidth}{0pt}
\setlength{\csromangathawakwidth}{0pt}
\csromangathameasuregroup{}{\textbf{นิกฺการณา ทุลฺลภา อชฺช มิตฺตา,} \\ \textbf{อตฺตตฺถปญฺญา}\textsuperscript{1}\textbf{ อสุจี มนุสฺสา,}}
\csromangathameasurewak{\textbf{นิกฺการณา ทุลฺลภา อชฺช มิตฺตา,} \\ \textbf{อตฺตตฺถปญฺญา}\textsuperscript{1}\textbf{ อสุจี มนุสฺสา,}}
\setlength{\csromangathaleftcol}{0pt}
\csromangathagroup{\csromangathastanza{\csromangathawak{\textbf{นิกฺการณา ทุลฺลภา อชฺช มิตฺตา,} \\ \textbf{อตฺตตฺถปญฺญา}\csromansharedfootnote{dd304db529a1c150}{อตฺตฏฺฐปญฺญา (ก)}\textbf{ อสุจี มนุสฺสา,}}}}

\prose{\textbf{เอโก จเร ขคฺควิสาณ}\-\textbf{กปฺโป.} (10)}

\prose{\csromanfolio{306}\textbf{ภชนฺติ เสวนฺติ จ การณตฺถา}ติ อตฺตตฺถการณา ปรตฺถการณา อุภยตฺถ\-การณา ทิฏฺฐธมฺมิก\-ตฺถการณา สมฺปรายิกตฺถ\-การณา ปรมตฺถ\-การณา ภชนฺติ สมฺภชนฺติ เสวนฺติ นิเสวนฺติ สํเสวนฺติ ปฏิเสวนฺตีติ ภชนฺติ เสวนฺติ จ การณตฺถา.}

\prose{\textbf{นิกฺการณา ทุลฺลภา อชฺช มิตฺตา}ติ เทฺว มิตฺตา อคาริกมิตฺโต จ อนาคาริกมิตฺโต จ ฯเปฯ อยํ อคาริกมิตฺโต ฯเปฯ อยํ อนาคาริกมิตฺโต. \textbf{นิกฺการณา ทุลฺลภา อชฺช มิตฺตา}ติ อิเม เทฺว มิตฺตา อการณา นิกฺการณา อเหตู อปฺปจฺจยา ทุลฺลภา (ทุลฺลทฺธา สุทุลฺลทฺธา)\csromansharedfootnote{63e3590ed1ce5d75}{( ) นตฺถิ สฺยาโปตฺถเก.}ติ นิกฺการณา ทุลฺลภา อชฺช มิตฺตา.}

\prose{\textbf{อตฺตตฺถปญฺญา อสุจี มนุสฺสา}ติ \textbf{อตฺตตฺถ}\-\textbf{ปญฺญา}ติ อตฺตโน อตฺถาย อตฺตโน เหตุ อตฺตโน ปจฺจยา อตฺตโน การณา ภชนฺติ สมฺภชนฺติ เสวนฺติ นิเสวนฺติ สํเสวนฺติ ปฏิเสวนฺติ อาจรนฺติ สมาจรนฺติ ปยิรุปาสนฺติ ปริปุจฺฉนฺติ ปริปญฺหนฺตีติ อตฺตตฺถปญฺญา. \textbf{อสุจี มนุสฺสา}ติ อสุจินา กายกมฺเมน สมนฺนาคตาติ อสุจี มนุสฺสา, อสุจินา วจีกมฺเมน สมนฺนาคตาติ อสุจี มนุสฺสา, อสุจินา มโนกมฺเมน สมนฺนาคตาติ อสุจี มนุสฺสา, อสุจินา ปาณาติปาเตน. อสุจินา อทินฺนาทาเนน ฯเปฯ. อสุจินา กาเมสุ\-มิจฺฉาจาเรน. อสุจินา มุสาวาเทน. อสุจิยา ปิสุณาย วาจาย สมนฺนาคตา. อสุจิยา ผรุสาย วาจาย สมนฺนาคตา. อสุจินา สมฺผ\-ปฺปลาเปน สมนฺนาคตา. อสุจิยา อภิชฺฌาย สมนฺนาคตา อสุจินา พฺยาปาเทน สมนฺนาคตาติ อสุจี มนุสฺสา, อสุจิยา มิจฺฉาทิฏฺฐิยา สมนฺนาคตาติ อสุจี มนุสฺสา, อสุจิยา เจตนาย สมนฺนาคตาติ อสุจี มนุสฺสา, อสุจิยา ปตฺถนาย สมนฺนาคตาติ อสุจี มนุสฺสา, อสุจินา ปณิธินา สมนฺนาคตาติ อสุจี มนุสฺสา, อสุจี หีนา นิหีนา โอมกา ลามกา ฉตุกฺกา ปริตฺตาติ อตฺตตฺถปญฺญา อสุจี มนุสฺสา.}

\prose{\textbf{เอโก จเร ขคฺควิสาณ}\-\textbf{กปฺโป}ติ \textbf{เอโก}ติ โส ปจฺเจก\-สมฺพุทฺโธ ปพฺพชฺชา\-สงฺขาเตน เอโก ฯเปฯ. \textbf{จเร}ติ อฏฺฐ จริยาโย ฯเปฯ. \textbf{ขคฺควิสาณ}\-\textbf{กปฺโป}ติ ยถา ขคฺคสฺส นาม วิสาณํ เอกํ โหติ อทุติยํ ฯเปฯ เอโก จเร ขคฺควิสาณ\-กปฺโป. เตนาห โส ปจฺเจก\-สมฺพุทฺโธ\csromandash{}}

\setlength{\csromangathapagewidth}{0pt}
\setlength{\csromangathawakwidth}{0pt}
\csromangathameasuregroup{อชิโต ติสฺสเมตฺเตยฺโย, \\ โธตโก อุปสีโว จ, \\ โตเทยฺยกปฺปา ทุภโย, \\ ภทฺราวุโธ อุทโย จ, \\ โมฆราชา จ เมธาวี, \\ โสฬสานํ\textsuperscript{1} ปเนเตสํ, \\ ปารายนานํ นิทฺเทสา, \\ ขคฺควิสาณสุตฺตานํ, \\ นิทฺเทสา ทุวิธา เญยฺยา,}{ภชนฺติ เสวนฺติ จ การณตฺถา, \\ นิกฺการณา ทุลฺลภา อชฺช มิตฺตา. \\ อตฺตตฺถปญฺญา อสุจี มนุสฺสา, \\ เอโก จเร ขคฺควิสาณกปฺโปติ. \\ จตุตฺโถ วคฺโค. \\ ขคฺควิสาณสุตฺตนิทฺเทโส เอกูนวีสติโม. \\ \csromangathabat{อชิโต ติสฺสเมตฺเตยฺโย,}{ปุณฺณโก อถ เมตฺตคู.} \\ \csromangathabat{โธตโก อุปสีโว จ,}{นนฺโท จ อถ เหมโก.} \\ \csromangathabat{โตเทยฺยกปฺปา ทุภโย,}{ชตุกณฺณี จ ปณฺฑิโต.} \\ \csromangathabat{ภทฺราวุโธ อุทโย จ,}{โปสาโล จาปิ พฺราหฺมโณ.} \\ \csromangathabat{โมฆราชา จ เมธาวี,}{ปิงฺคิโย จ มหาอิสิ.} \\ \csromangathabat{โสฬสานํ\textsuperscript{1} ปเนเตสํ,}{พฺราหฺมณานํว สาสนํ.} \\ \csromangathabat{ปารายนานํ นิทฺเทสา,}{ตตฺถกา จ ภวนฺติ หิ\textsuperscript{1}.} \\ \csromangathabat{ขคฺควิสาณสุตฺตานํ,}{นิทฺเทสาปิ ตเถว จ.} \\ \csromangathabat{นิทฺเทสา ทุวิธา เญยฺยา,}{ปริปุณฺณา สุลกฺขิตาติ.}}
\csromangathameasurewak{ภชนฺติ เสวนฺติ จ การณตฺถา, \\ นิกฺการณา ทุลฺลภา อชฺช มิตฺตา. \\ อตฺตตฺถปญฺญา อสุจี มนุสฺสา, \\ เอโก จเร ขคฺควิสาณกปฺโปติ. \\ จตุตฺโถ วคฺโค. \\ ขคฺควิสาณสุตฺตนิทฺเทโส เอกูนวีสติโม.}
\csromangathasetleft{อชิโต ติสฺสเมตฺเตยฺโย, \\ โธตโก อุปสีโว จ, \\ โตเทยฺยกปฺปา ทุภโย, \\ ภทฺราวุโธ อุทโย จ, \\ โมฆราชา จ เมธาวี, \\ โสฬสานํ\textsuperscript{1} ปเนเตสํ, \\ ปารายนานํ นิทฺเทสา, \\ ขคฺควิสาณสุตฺตานํ, \\ นิทฺเทสา ทุวิธา เญยฺยา,}
\csromangathagroupwithclosermajor{\csromangathastanza{\csromangathawak{\csromanfolio{307}ภชนฺติ เสวนฺติ จ การณตฺถา, \\ นิกฺการณา ทุลฺลภา อชฺช มิตฺตา. \\ อตฺตตฺถปญฺญา อสุจี มนุสฺสา, \\ เอโก จเร ขคฺควิสาณ\-กปฺโปติ. \\ จตุตฺโถ วคฺโค. \\ ขคฺควิสาณสุตฺตนิทฺเทโส เอกูนวีสติโม.}}\csromangathastanza{\csromangathabat{อชิโต ติสฺสเมตฺเตยฺโย,}{ปุณฺณโก อถ เมตฺตคู.} \\ \csromangathabat{โธตโก อุปสีโว จ,}{นนฺโท จ อถ เหมโก.}}\csromangathastanza{\csromangathabat{โตเทยฺยกปฺปา ทุภโย,}{ชตุกณฺณี จ ปณฺฑิโต.} \\ \csromangathabat{ภทฺราวุโธ อุทโย จ,}{โปสาโล จาปิ พฺราหฺมโณ.}}\csromangathastanza{\csromangathabat{โมฆราชา จ เมธาวี,}{ปิงฺคิโย จ มหาอิสิ.} \\ \csromangathabat{โสฬสานํ\csromansharedfootnote{483895244a036fe8}{โสฬสนฺนํ (สฺยา, ก)} ปเนเตสํ,}{พฺราหฺมณานํว สาสนํ.}}\csromangathastanza{\csromangathabat{ปารายนานํ นิทฺเทสา,}{ตตฺถกา จ ภวนฺติ หิ\csromansharedfootnote{febd4a0c4e4d09a3}{วา (ก)}.} \\ \csromangathabat{ขคฺควิสาณสุตฺตานํ,}{นิทฺเทสาปิ ตเถว จ.} \\ \csromangathabat{นิทฺเทสา ทุวิธา เญยฺยา,}{ปริปุณฺณา สุลกฺขิตาติ.}}}{\textbf{จูฬนิทฺเทส}\-\textbf{ปาฬิ นิฏฺฐิตา.}}
\orphannote{สีมิกตํ มริยทิกตํ (สฺยา)}

